% Options for packages loaded elsewhere
\PassOptionsToPackage{unicode}{hyperref}
\PassOptionsToPackage{hyphens}{url}
\PassOptionsToPackage{dvipsnames,svgnames,x11names}{xcolor}
%
\documentclass[
  11pt,
  oneside]{book}
\usepackage{amsmath,amssymb}
\usepackage{iftex}
\ifPDFTeX
  \usepackage[T1]{fontenc}
  \usepackage[utf8]{inputenc}
  \usepackage{textcomp} % provide euro and other symbols
\else % if luatex or xetex
  \usepackage{mathtools}
\usepackage{unicode-math} % this also loads fontspec
  \defaultfontfeatures{Scale=MatchLowercase}
  \defaultfontfeatures[\rmfamily]{Ligatures=TeX,Scale=1}
\fi
\usepackage{lmodern}
\ifPDFTeX\else
  % xetex/luatex font selection
\fi
% Use upquote if available, for straight quotes in verbatim environments
\IfFileExists{upquote.sty}{\usepackage{upquote}}{}
\IfFileExists{microtype.sty}{% use microtype if available
  \usepackage[]{microtype}
  \UseMicrotypeSet[protrusion]{basicmath} % disable protrusion for tt fonts
}{}
\makeatletter
\@ifundefined{KOMAClassName}{% if non-KOMA class
  \IfFileExists{parskip.sty}{%
    \usepackage{parskip}
  }{% else
    \setlength{\parindent}{0pt}
    \setlength{\parskip}{6pt plus 2pt minus 1pt}}
}{% if KOMA class
  \KOMAoptions{parskip=half}}
\makeatother
\usepackage{xcolor}
\usepackage[a4paper,top=23mm,bottom=23mm,left=23mm,right=23mm]{geometry}
\usepackage{longtable,booktabs,array}
\usepackage{calc} % for calculating minipage widths
% Correct order of tables after \paragraph or \subparagraph
\usepackage{etoolbox}
\makeatletter
\patchcmd\longtable{\par}{\if@noskipsec\mbox{}\fi\par}{}{}
\makeatother
% Allow footnotes in longtable head/foot
\IfFileExists{footnotehyper.sty}{\usepackage{footnotehyper}}{\usepackage{footnote}}
\makesavenoteenv{longtable}
\setlength{\emergencystretch}{3em} % prevent overfull lines
\providecommand{\tightlist}{%
  \setlength{\itemsep}{0pt}\setlength{\parskip}{0pt}}
\setcounter{secnumdepth}{-\maxdimen} % remove section numbering
\usepackage{fontspec}
% Use Latin text metrics for the English edition. CJK occurs only in a few
% original source titles and is selected explicitly below.
\setmainfont{Latin Modern Roman}
\setsansfont{Latin Modern Sans}
\setmonofont{DejaVu Sans Mono}
\IfFileExists{UEOT_CJK.ttf}{
  \newfontfamily\UEOTCJK{UEOT_CJK.ttf}[BoldFont=UEOT_CJK.ttf,BoldFeatures={FakeBold=2}]
}{
  \newfontfamily\UEOTCJK{Noto Serif CJK SC}
}
\usepackage{mathtools}
\usepackage{fvextra}
\usepackage{etoolbox}
\usepackage{newunicodechar}
\newunicodechar{核}{{\UEOTCJK 核}}
\newunicodechar{心}{{\UEOTCJK 心}}
\newunicodechar{理}{{\UEOTCJK 理}}
\newunicodechar{论}{{\UEOTCJK 论}}
\newunicodechar{正}{{\UEOTCJK 正}}
\newunicodechar{式}{{\UEOTCJK 式}}
\newunicodechar{版}{{\UEOTCJK 版}}
\newunicodechar{数}{{\UEOTCJK 数}}
\newunicodechar{学}{{\UEOTCJK 学}}
\newunicodechar{基}{{\UEOTCJK 基}}
\newunicodechar{础}{{\UEOTCJK 础}}
\newunicodechar{统}{{\UEOTCJK 统}}
\newunicodechar{一}{{\UEOTCJK 一}}
\AtBeginDocument{\lefthyphenmin=3\righthyphenmin=3}
\newunicodechar{ℝ}{\ensuremath{\mathbb{R}}}
\newunicodechar{ℏ}{\ensuremath{\hbar}}
\newunicodechar{∎}{\ensuremath{\square}}
\newunicodechar{□}{\ensuremath{\square}}
\newunicodechar{⊗}{\ensuremath{\otimes}}
\usepackage{fancyhdr}
\pagestyle{fancy}
\fancyhf{}
\fancyhead[L]{\small UEOT PHYSICS}
\fancyhead[R]{\small Bilingual revision 3.11}
\fancyfoot[C]{\small\thepage}
\setlength{\headheight}{16pt}
\renewcommand{\headrulewidth}{0.3pt}
\setcounter{tocdepth}{0}
\setlength{\parindent}{0pt}
\setlength{\parskip}{5pt plus 1pt}
\setlength{\emergencystretch}{8em}
\setlength{\tabcolsep}{3pt}
\renewcommand{\arraystretch}{1.16}
\AtBeginEnvironment{longtable}{\small}
\allowdisplaybreaks[2]
\renewcommand{\contentsname}{Contents}
\makeatletter
% Reserve a real gap before a chapter page number. A zero-minimum \hfil
% can otherwise place a long final word directly against the number.
\patchcmd{\l@chapter}{#1\nobreak\hfil}{#1\nobreak\hspace{1em}\hfil}{}{\PackageError{UEOT}{Chapter TOC spacing patch failed}{Inspect the book class before publishing.}}
\renewcommand{\maketitle}{
  \begin{titlepage}
  \thispagestyle{empty}
  \vspace*{30mm}
  {\color{teal}\fontsize{46}{52}\selectfont UEOT\par}
  \vspace{5mm}
  {\fontsize{32}{40}\selectfont\bfseries PHYSICS\par}
  \vspace{11mm}
  {\Large From Object Mathematics to Quantum Physics\par}
  \vspace{4mm}
  {\large Derivations, Formal Proofs, and Research Boundaries\par}
  \vspace{10mm}
  {\color{teal}\rule{\textwidth}{1pt}\par}
  \vspace{7mm}
  {\large Third edition\quad Bilingual revision 3.11\par}
  \vspace{6mm}
  {Degui Qian\par}
  \vspace{5mm}
  {71 chapters\quad 6 appendices\quad UEOT-QM Kernel 1.14\par}
  \vfill
  {Objects and protocols\quad Quantum operations and fields\par}
  \vspace{3mm}
  {Proofs, countermodels, and empirical limits\par}
  \vspace{7mm}
  {14 November 2026\par}
  \end{titlepage}
}
\makeatother
\ifLuaTeX
  \usepackage{selnolig}  % disable illegal ligatures
\fi
\usepackage{bookmark}
\IfFileExists{xurl.sty}{\usepackage{xurl}}{} % add URL line breaks if available
\urlstyle{same}
\hypersetup{
  pdftitle={UEOT Physics},
  pdfauthor={Degui Qian},
  colorlinks=true,
  linkcolor={teal},
  filecolor={Maroon},
  citecolor={Blue},
  urlcolor={teal},
  pdfcreator={LaTeX via pandoc}}

\title{UEOT Physics}
\usepackage{etoolbox}
\makeatletter
\providecommand{\subtitle}[1]{% add subtitle to \maketitle
  \apptocmd{\@title}{\par {\large #1 \par}}{}{}
}
\makeatother
\subtitle{From Object Mathematics to Quantum Physics: Derivations,
Formal Proofs, and Research Boundaries}
\author{Degui Qian}
\date{2026-11-14}

\begin{document}
\frontmatter
\maketitle

{
\hypersetup{linkcolor=}
\setcounter{tocdepth}{0}
\tableofcontents
}
\mainmatter
\chapter{0. Scope, domains, and responsibilities of proof}\label{ch-0}

Reading this book begins by distinguishing mathematical validity of a
model from evidence that the model describes nature. This chapter
explains how to track those responsibilities and how UEOT organizes
standard physics, conditional derivations, and mechanisms awaiting
tests; these distinctions remain in force throughout.

Like the complete UEOT Core mathematical text, this book uses a
continuous, self-contained Markdown source with stable identifiers. It
contains definitions, derivations, the conditions under which standard
theorems apply, counterexamples, applications, interfaces, and
executable checks. It is not a collection of historical patches: the
reader should not have to recover an older edition to find an essential
formula.

The starting questions remain those of the original UEOT paper: how
objects persist, how openness acquires structure, and how phase
consistency, prediction, and resource constraints enter an account of
physical systems. The quantum branch supplies particular physical
realizations of these ideas. The book connects to the general
mathematical Core, but \textbf{does not treat general object mathematics
as an axiom system already sufficient to select quantum physics
uniquely}.

The quantum programme assembled here covers finite-dimensional objects
and measurement, P2 quantum selection, continuous variables, spin and
identical particles, many-body and open systems, free relativistic
fields, interactions and renormalization, QED, and candidate Standard
Model structures and generation-counting arguments. ``Complete'' refers
to a complete main text for this stated architecture. It does not mean
completion of every specialty in quantum physics, a rigorous
construction of interacting four-dimensional QFT, quantum gravity, or a
final unified theory.

\section{0.1 The role of a claim and its burden of proof}

\begin{longtable}[]{@{}
  >{\raggedright\arraybackslash}p{(\columnwidth - 4\tabcolsep) * \real{0.2500}}
  >{\raggedright\arraybackslash}p{(\columnwidth - 4\tabcolsep) * \real{0.3400}}
  >{\raggedright\arraybackslash}p{(\columnwidth - 4\tabcolsep) * \real{0.4100}}@{}}
\toprule\noalign{}
\begin{minipage}[b]{\linewidth}\raggedright
Label
\end{minipage} & \begin{minipage}[b]{\linewidth}\raggedright
Meaning
\end{minipage} & \begin{minipage}[b]{\linewidth}\raggedright
Required justification
\end{minipage} \\
\midrule\noalign{}
\endhead
\bottomrule\noalign{}
\endlastfoot
D & Definition, convention, or protocol & Specify types, support, and
units \\
P & Proof or explicit derivation supplied in this book & A conditional
mathematical result; the label does not itself assert originality \\
K & Standard mathematical or physical result & State its domain; do not
present an imported major theorem as a complete proof given here \\
H & Additional principle for a quantum or particle sector & Not an
automatic consequence of Core; requires independent justification \\
S & UEOT interpretation or connection between modules & Do not call it a
new law of nature when it leaves the physical predictions unchanged \\
E & Experimental or numerical evidence & Distinguish real observations,
synthetic data, and algebraic checks \\
N / O & Counterexample or impossibility result / open problem & Excluded
from the established positive dependency chain \\
\end{longtable}

The P-labelled results have proofs or derivations in the text. Major
imported K theorems are listed in the appendices. A numerical check is
neither a machine-checked proof nor empirical evidence about nature.

An independent \textbf{verification axis} distinguishes a proof in
prose, a proof checked by a formal kernel, exact finite computation,
floating-point checks, synthetic-data tests, and actual experimental
data. An item can be both a P result under additional physical
assumptions and a conditional theorem checked by the Lean kernel. Kernel
checking does not remove its physical assumptions. Conversely, an
experiment supports a specified model, apparatus, protocol, and error
range; finite empirical support does not become an unconditional
mathematical theorem. P/K/H labels identify argumentative roles, whereas
the verification axis records how an item has been checked. Chapter 48
adds the \texttt{dependency\_\allowbreak class} field for Core, bridge, sector,
standard-physics, effective-model, observable-consequence, and
open-conjecture roles. These three axes remain distinct. {[}S2, §§0.3,
28.2; S3, §§1--2; S4, §§2, 8--9{]}

\section{0.2 Precise realizations of the original ideas}

\begin{longtable}[]{@{}
  >{\raggedright\arraybackslash}p{(\columnwidth - 4\tabcolsep) * \real{0.2500}}
  >{\raggedright\arraybackslash}p{(\columnwidth - 4\tabcolsep) * \real{0.3400}}
  >{\raggedright\arraybackslash}p{(\columnwidth - 4\tabcolsep) * \real{0.4100}}@{}}
\toprule\noalign{}
\begin{minipage}[b]{\linewidth}\raggedright
Original idea
\end{minipage} & \begin{minipage}[b]{\linewidth}\raggedright
Valid quantum realization
\end{minipage} & \begin{minipage}[b]{\linewidth}\raggedright
Unlicensed inference
\end{minipage} \\
\midrule\noalign{}
\endhead
\bottomrule\noalign{}
\endlastfoot
An object is a persistent relational structure & Equivalence of protocol
futures, effective factors, transport along histories & A name or
wavefunction label automatically defines an object \\
Phase closure & Projective representations, holonomy, boundary
conditions, and gauge-anomaly consistency & A closed-orbit integral
directly establishes the CCR \\
Π/Φ & A conditional Fisher--action model, quantum-information costs, and
control objectives & Any two functions uniquely select QM \\
Parent objects and child interfaces & Tensor composition, purification,
and instrument dilation & Every general parent object is a pure quantum
parent system \\
Observation and resolution & POVMs, instruments, a public commutative
algebra, and process coarse-graining & Coarse-graining itself causes
objective collapse \\
GOD/GOA & Optimal control and long-term structure for a specified
quantum-control task & Every quantum state naturally optimizes the same
objective \\
A particle as a ``minimal closure'' & Irreducible representations and
anomaly-free physical realization packages & A Betti number equals the
number of generations \\
Unity across scales & Explicitly typed bridges between Core, processes,
and QFT & Transformations called RG are therefore the same
transformation \\
\end{longtable}

\section{0.3 How to use the main text}

Chapters 1--9 introduce foundations and measurement; Chapters 10--15
cover continuous variables and symmetry; Chapters 16--20 cover many-body
systems, environments, and control; Chapters 21--28 cover QFT, RG, and
QED; Chapters 29--32 address particle structure and research boundaries;
and Chapters 33--35 assemble interfaces, applications, and error
budgets. The appendices contain verification code and actual receipts
with version identities. Section 0.5 locates the additional modules.

The logical direction is \textbf{declare a physical realization → define
implementable protocols → calculate responses → establish object and
prediction certificates → solve control problems under additional task
specifications → test the sector assumptions}. A desired conclusion must
not be inserted into a definition of what an object ``must'' be and then
reported as a derivation.

\section{0.4 Three different senses of completion}

UEOT now has an assemblable architecture for research in quantum
physics. Once a quantum realization and its sector assumptions are
specified, many modules can be fully defined, derived, and computed.
General Core alone has not uniquely selected the quantum axioms,
nature's gauge group, or all physical parameters. Nor has this work
rigorously constructed interacting four-dimensional field theory or
quantum gravity. The completed deliverable is therefore an integrated
theoretical text and runnable benchmarks within the declared scope.
Completeness of a document is not completion of fundamental physics.

Each module is assessed by four questions. Are its definitions and
proofs complete? Are its physical inputs independently stated? Can its
outputs be calculated? Does a candidate admit a discriminating test? We
do not assign ``completion percentages'' without an objective
denominator.

\begin{longtable}[]{@{}
  >{\raggedright\arraybackslash}p{(\columnwidth - 4\tabcolsep) * \real{0.2500}}
  >{\raggedright\arraybackslash}p{(\columnwidth - 4\tabcolsep) * \real{0.3400}}
  >{\raggedright\arraybackslash}p{(\columnwidth - 4\tabcolsep) * \real{0.4100}}@{}}
\toprule\noalign{}
\begin{minipage}[b]{\linewidth}\raggedright
Module
\end{minipage} & \begin{minipage}[b]{\linewidth}\raggedright
Established position
\end{minipage} & \begin{minipage}[b]{\linewidth}\raggedright
Remaining conditions or work
\end{minipage} \\
\midrule\noalign{}
\endhead
\bottomrule\noalign{}
\endlastfoot
Core to quantum selection & Classical countermodels are given; the OPT
background and selection principles are stated & Major reconstruction
theorems remain external dependencies; the physical necessity of the
principles is unproved \\
States, effects, channels, and instruments & Finite-dimensional
definitions, constructive proofs, composition, and error bounds &
Availability and calibration of real preparations and measurements
remain experimental inputs \\
Objects, records, and prediction & Future equivalence, public records,
SRN, and interfaces with memory & Object partitions depend on the
allowed protocols and need not be unconditionally unique \\
Measurement and actualization & Coherent dilation, conditional states,
and absorption and conditional Born results for a specified SSE & Core
does not select that stochastic equation; arbitrary system channels
cannot distinguish the ontological interpretations \\
Continuous variables and global phase & Self-adjoint domains, Weyl
representations, energy truncation, and holonomy on a ring & Physical
boundary conditions and the action scale are inputs; integer
closed-orbit conditions do not establish the CCR in reverse \\
Spin, statistics, and many-body systems & Symmetry representations, Fock
spaces, finite models, and thermal-environment benchmarks & Relativistic
spin--statistics is an external theorem; this book provides no general
guarantee of efficient solution for arbitrary many-body problems \\
Open systems, control, and thermodynamics & GKSL/Davies conditions,
process memory, control gradients, and free energy & Strong coupling and
infinite systems require model-specific work \\
Topology and stability & Controlled, complete examples using Berry
bundles and SSH edge states & No general interacting topological
classification or selection of particle generations \\
Free fields and perturbative QFT & Fields, propagation, causality,
scattering, one-loop renormalization, and QED benchmarks & Higher
orders, infrared issues, and nonperturbative continuum limits are not
all completed \\
Lattice gauge theory & A finite pure-gauge measure, gauge invariance,
continuum expansion, and an exhaustively enumerable model & The
four-dimensional continuum limit, the QCD mass gap, chiral fermions, and
real-time error control require dedicated research \\
Candidate Standard Model structure & Gauge representations, local
anomalies, Higgs, Yukawa, and neutrino interfaces & The group, its
global form, couplings, masses, and mixing parameters are not uniquely
derived \\
Three generations and flavour structure & Conditional low-energy
selection, two high-energy candidates, and rank-one FCNC sum rules and
lower bounds & The origin of HGP, an independent exclusion of five
colours, a complete spectrum, and constraints from nature \\
Identification of new physics & A three-branch CP bound, a four-level
witness, and model relations and protocols for joint energy/decoherence
tests & No new real experiment was performed for this edition; excluding
a noise subclass does not establish a new law of nature \\
Gravity and cosmology & Chapter 47 provides controlled starting points
and necessary interfaces & No construction of quantum gravity or
derivation of dark matter, dark energy, or the cosmological initial
state \\
\end{longtable}

``Based on UEOT'' has a precise meaning here. A physical model supplies
intervention probabilities. Core organizes future equivalence, carriers,
persistence, resources, and control; sector bridges specify when these
structures can be composed. If removing the UEOT interpretation leaves
exactly the predictions of standard QM, that part is an S/K-level
organization or reconstruction. Only an additional assumption producing
different, identifiable probability constraints enters the testing
programme as a candidate H for new physics.

\section{0.5 Assembly of the text and successive additions}

Chapters 1--35 supply foundations, calculations, and interfaces.
Chapters 36--42 connect the high-energy generation audit, flavour
constraints, coherence witnesses, actualization proofs, and research
protocols to that foundation. Chapters 43--46 add global phase,
topology, lattice models, and statistics for correlated trials. Chapters
47--48 define research acceptance contracts and maintenance of the
integrated text. Chapters 49--56 develop the finite kernel, errors,
randomization, long-time averages, conditional survival,
quantum-selection separations, and the general finite quantum Doob
construction. Chapters 57--59 establish full finite complex
purification, genuine record memory, and conditional dynamics in a
specified noncommutative model. Chapters 60--62 develop separation by
all effects, complex local tomography for arbitrary joint states, full
convex purity, auxiliary-unitary relations, and action-by-action budgets
for a common fixed-depth controller. Chapters 63--65 add the
trace-distance expression and optimal binary discrimination, minimal
auxiliary dimension and directed isometries, and original finite
stopping trees with a budget for unstopped weight. The appendices
preserve notation, sources, all stable result identifiers, reproduction
code, and the complete audit of named assumptions.

The conditional three-generation theorem in Chapter 31 must be read with
Chapter 36's distinction between low and high energy. Chapter 40
supplies an absorption proof within a specific SSE for the conditional
martingale argument in Chapter 3. Chapters 38--39 turn the SRN and
tomography results of Chapters 7--8 into witnesses against noise
classes. Chapter 45 gives the cross-scale interface of Chapter 26 an
actual carrier in a finite gauge model. These connections are all
contained in this document.

The v2.0 appendices preserve the quantum-mainline, frontier,
global-phase, SSH-edge, finite-lattice, and correlated-trial checks
executed at that time. They are historical receipts tied to a version.
Subsequent kernel results and exact certificates have their own source
manifests, execution logs, and audit receipts. External experimental
evidence must be traced to its original papers and data descriptions;
any future hardware experiment by the project will require its own
acquisition and calibration records. A historical pass is not a fresh
check of the current version, and a new local proof does not
retrospectively formalize older chapters outside its scope. Finite
numerical checks neither prove imported major theorems nor replace
universal proofs or actual experimental data.

The research revision published as edition 3.6 connects, in Chapter 66,
the original finite quantum instrument, the complete conditioned record
and terminal state, randomization over allowed actions, initial-state
preparation, early-stopping costs, and errors. The original acceptance
tree, singular support, zero-probability events, visible labels, and
observation permissions remain explicit. General formal proofs, exact
finite examples, floating-point diagnostics, and physical empirical
evidence are recorded separately. The v31.9 source had already been
merged through PR \#356; its outdated pending wording is corrected using
the available CI records and source identities. Verification newly
performed for the v31.10 research revision retains its own execution
identity rather than borrowing that of earlier runs.

\section{0.6 From a research ambition to a falsifiable question}

Represent each candidate theory by a quadruple
\((\mathfrak M,\Theta,\mathcal U,\mathcal P)\): a model class, parameter
domain, allowed protocols, and response law. Given a record \(D\), an
``explanation'' must specify which parameters are fixed by training data
and which quantities are held-out predictions. If the model,
environment, or protocol may be changed arbitrarily for every possible
\(D\), the model class has no discriminating content.

For two candidates, define their separation over a specified protocol
family by \[
\Delta_{\mathcal U}(M_0,M_1)
=\sup_{u\in\mathcal U}\operatorname{TV}(p_0(\cdot|u),p_1(\cdot|u)).
\] Here \(\operatorname{TV}(P,Q)=\sup_B|P(B)-Q(B)|\); for a finite
outcome space it is \(\tfrac12\sum_y|p_y-q_y|\). The comparison is
between probability laws on the same record space under the same
protocol. If both candidates contain adjustable parameters, the relevant
objects are the distance and intersection between their response sets,
not merely a favoured example from each. The environmental-equivalence
obstruction in Chapter 41 concerns set inclusion; the witness in Chapter
39 supplies an explicit separating hyperplane. Every research contract
in Chapter 47 follows this distinction.

\section{0.7 The central philosophical question: identifying objects
through change}

UEOT begins with a question that can be made precise: when material
constituents, microscopic states, and environmental inputs continually
change, why can we still track an object and predict its next response?
The original paper locates an answer in persistent relations, feedback,
openness, and hierarchical composition. These ideas determine what this
book investigates first, but do not by themselves specify a Hamiltonian,
a set of quantum instruments, or a particle-representation spectrum.
{[}S1, §1.3, PDF pp.~2--4; §§2.2--2.4, pp.~4--8{]}

We adopt the mathematical reformulation completed in Core. First fix a
process, implementable protocols, and the candidate's own future, and
identify the information needed to determine the common responses. Then
ask separately whether this information admits recursive updating, what
physical apparatus realizes it, how it withstands interventions, and how
it is transported along histories. ``Relational identity'' now has
checkable content: a change of coordinates leaves the probabilities of
the same experiments invariant; retaining an appropriate state makes
information removed from the full history irrelevant to the declared
future. Existence certificates for objects may be strong or weak. Claims
of recovery, self-maintenance, inheritance, or optimal control require
the corresponding additional mechanisms. {[}S2, §§0.1, 3--4, 12.1{]}

This gives a specific meaning to the original phrase ``Being is
persistent Becoming.'' What persists is a verified relation between
response and realization; the state itself may change. A moving wave
packet need not remain at a fixed location. A logical qubit exposed to
noise and recovery need not occupy the same physical state throughout.
Different hardware may realize the same input--output responses.
Predictive identity, physical carrier, and historical instance are
therefore connected but distinct concepts.

The original paper associated general stable existence with nonzero
first homology, \(H_1\ne0\). This cannot serve as a universal theorem
here. For \(\dot x=-x\), the solution is \(x(t)=x(0)e^{-t}\): the origin
is globally asymptotically stable, but both the state space
\(\mathbb R\) and the attractor \(\{0\}\) are contractible, with
\(H_1=0\). Conversely, on the cylinder \(\mathbb R\times S^1\), let
\(\dot u=u,\ \dot\theta=1\). The state space has nontrivial first
homology and \(u=0\) is a closed orbit, yet transverse perturbations
grow as \(e^t\). Topological cycles, feedback edges in a
causal-dependence graph, and dynamical stability need separate
definitions and bridging proofs. {[}S1, §2.2, p.~4, printed lines
147--160; S2, §§9.4, 12.2{]}

Likewise, stationary action does not mean that a system monotonically
increases a benefit in physical time. A harmonic oscillator satisfies a
variational stationarity condition while its phase-space trajectory can
be periodic; that fact does not make it converge to a unique
maximum-payoff state. The scale parameter of an RG flow is also distinct
from physical time. Interpreting a scale flow as dynamics requires a
process bridge; interpreting a drift as optimal behaviour requires an
evaluation bridge. {[}S1, §2.3, pp.~6--7; S2, §§18--23, 28.3--4{]}

Unity in UEOT physics therefore begins with a common language for
identification, realization, transport, and certification. Standard
quantum models supply the actual probability laws. UEOT asks which state
distinctions matter to a task, how to prove that compression preserves
every allowed future, and when the memory, operation library, or object
boundary must be enlarged. Any added probability constraint beyond the
standard model becomes a physical hypothesis to be tested.

\section{0.8 Source hierarchy and stable texts}

The meaning of the theory follows an explicit source hierarchy. The
frozen definitions, assumptions, and proofs of Core v3 determine the
general mathematical semantics. Quantum chapters add their own physical
structures. Executable code and Lean declarations identify the
particular claims checked and the method of checking. The original paper
preserves the research motivations and conjectures that remain
incomplete. Its philosophical wording does not override subsequently
corrected mathematical definitions. {[}S2, §0 and Appendix E; S4,
§§1--2{]}

The foundational chapters use the following stable source identifiers so
that readers can locate the original responsibility directly.

\begin{longtable}[]{@{}
  >{\raggedright\arraybackslash}p{(\columnwidth - 4\tabcolsep) * \real{0.2500}}
  >{\raggedright\arraybackslash}p{(\columnwidth - 4\tabcolsep) * \real{0.3400}}
  >{\raggedright\arraybackslash}p{(\columnwidth - 4\tabcolsep) * \real{0.4100}}@{}}
\toprule\noalign{}
\begin{minipage}[b]{\linewidth}\raggedright
ID
\end{minipage} & \begin{minipage}[b]{\linewidth}\raggedright
Original material
\end{minipage} & \begin{minipage}[b]{\linewidth}\raggedright
Role in this book
\end{minipage} \\
\midrule\noalign{}
\endhead
\bottomrule\noalign{}
\endlastfoot
S1 & \texttt{UEOT\_\allowbreak v3\_\allowbreak top.\allowbreak pdf}, 43 pages, user-supplied copy &
Motivation for objects, dual drives, and physical conjectures; citations
use actual PDF pages and printed line numbers \\
S2 & \texttt{UEOT\_\allowbreak Core\_\allowbreak Mathematics\_\allowbreak v3.\allowbreak 0\_\allowbreak Complete.\allowbreak md}, 2026-09-05 &
Frozen general definitions, stable P-IDs, mathematical proofs, and
adaptation conditions \\
S3 & \texttt{CORE\_\allowbreak V3\_\allowbreak OPEN\_\allowbreak PORT\_\allowbreak MATRIX.\allowbreak md}, 2026-10-07 & Precise gaps
at seven ports, work packages, and historical source-to-Lean mappings \\
S4 & \texttt{CORE\_\allowbreak V3\_\allowbreak POST\_\allowbreak CORE\_\allowbreak SCIENTIFIC\_\allowbreak CLOSURE.\allowbreak md}, 2026-10-07 &
Independent review of the Core interpretation, four-generator
boundaries, and scientific closure \\
S5 & \texttt{THEORY\_\allowbreak COMPLETION\_\allowbreak ROADMAP\_\allowbreak V2.\allowbreak md}, 2026-10-07 &
Dependencies and acceptance rules for subsequent research; PLANNED
entries are not completion receipts \\
Q2 & \texttt{UEOT\_\allowbreak Physics\_\allowbreak Unified\_\allowbreak v2.\allowbreak 0.\allowbreak md}, 2026-09-06 & Direct
source of the continuous derivations inherited in Chapters 0--48 \\
\end{longtable}

Dates, branches, phases, and CI counts in S3--S5 describe the time those
records were written; they are not automatically the running status of
this edition. Formal revisions preserve existing identifiers and useful
proofs, adding assumptions or correcting formulas at the relevant
location. A new version number does not turn an open conjecture into a
theorem.

\section{0.9 Structure and reading paths of the integrated edition}

This edition carries forward all 49 chapters of \emph{UEOT PHYSICS:
Unified Text 2.0}. It does not abbreviate the inherited quantum
foundations, QFT/QED, generation-counting, actualization, topology,
lattice, or experimental-statistics derivations. Chapters 0--13 clarify
the original ideas, rigorous Core, the background to quantum selection,
and measurement/prediction interfaces. Chapters 14--32 and 36--46 retain
advanced physics and countermodels. Chapters 33--35 and 47--48 update
concrete interfaces, errors, research paths, and maintenance of the
integrated text. Chapters 49--50 develop the existing finite quantum
kernel and stability proofs, while Chapters 51--53 supply actual mixing,
Cesàro, and Core conditional-survival bridges.

The theory's author is Degui Qian. Appendix C and the source manifest
preserve attribution for the original ideas and project results. This
edition continues their integration, completion of derivations, and
collaborative review. Standard physical theorems remain attributed to
their original literature, and machine-verification and experimental
identities are recorded separately.

Readers interested in the theoretical core can start with Chapters 0--5
and 49--67, then turn to the relevant physical domain. An assessment of
the three-generation question must read Chapters 31, 36--37, and §49.7
together. Appendix A provides notation and responsibility for external
theorems; Appendix B preserves the stable result register; Appendix C
records sources; Appendix D contains all inherited verification code and
the actual v3.0 execution receipts. Editions 3.1--3.7 checked the
identity of unchanged programs without rerunning them. Appendix E gives
the compilation and maintenance contract. Appendix F lists the actual
formalization scope and all requests for named-axiom audits,
distinguishing executions in the relevant edition from reuse tied to
historical receipts.

\section{0.10 Seven evidence levels: from established mathematics to
tests of nature}

The argument begins with established mathematics and physics, passes
through conditional Core interfaces and independent physical-sector
choices, and reaches provable conclusions, computable models, and
hypotheses that might differ from existing theory. The levels below tell
the reader what to examine at each stage. They do not constitute seven
completed parts of a single unification proof.

\begin{longtable}[]{@{}
  >{\raggedright\arraybackslash}p{(\columnwidth - 4\tabcolsep) * \real{0.2500}}
  >{\raggedright\arraybackslash}p{(\columnwidth - 4\tabcolsep) * \real{0.3400}}
  >{\raggedright\arraybackslash}p{(\columnwidth - 4\tabcolsep) * \real{0.4100}}@{}}
\toprule\noalign{}
\begin{minipage}[b]{\linewidth}\raggedright
Level
\end{minipage} & \begin{minipage}[b]{\linewidth}\raggedright
Content and principal entry points
\end{minipage} & \begin{minipage}[b]{\linewidth}\raggedright
Responsibility to examine
\end{minipage} \\
\midrule\noalign{}
\endhead
\bottomrule\noalign{}
\endlastfoot
Level 1: standard mathematics and physics & Linear algebra, probability,
standard QM, QFT/EFT; Chapters 4, 9--30 and Appendices A, C & Original
precedents, theorem hypotheses, units, and domains; using an established
result does not confer priority \\
Level 2: UEOT Core & Future-response factors, valid updates, transport,
conditional survival, and control interfaces; Chapters 1, 5, 49--56 &
Frozen definitions, exact quantifiers, common support, and domain
neutrality \\
Level 3: physical-sector selection & Complex Hilbert space, state/effect
cones, Born pairing, composition, symmetry, action, and spectrum & Which
structures are independent inputs and which follow from a reconstruction
theorem with its complete background \\
Level 4: rigorous conditional derivation & Measurement, channels,
purification, separation, finite protocols, and spectral-model theorems
& Whether the actual model meets the premises and whether the prose
proof matches the Lean declaration \\
Level 5: models and numerical verification & Rational witnesses, finite
matrices, cutoffs, one-loop benchmarks, and synthetic sampling;
Appendices D, F & Inputs, errors, execution dates, and source
identities; numerical agreement does not identify a unique mechanism \\
Level 6: additional testable hypotheses & Actualization strength, record
activation, extra spectral selection, or independent dynamical relations
& Additional assumptions, parameter-fixing rules, differing predictions,
and failure conditions; an underived mechanism remains OPEN \\
Level 7: empirical standing & Original experiments, calibration, raw
counts or sufficient statistics, and competing models; Chapters 8,
39--42, 46 & What an observation actually supports, constrains, or
excludes; absent data cannot be recorded as experimental success \\
\end{longtable}

These levels supplement the argument markers and verification axis of
§0.1 and the \texttt{dependency\_\allowbreak class} of Chapter 48. They do not
replace \texttt{CHAPTER\_\allowbreak DOSSIER.\allowbreak json} or
\texttt{SOURCE\_\allowbreak MANIFEST.\allowbreak json}. A formula may use standard physical
inputs, implement a Core bridge, and undergo numerical checks without
thereby discharging the independent tasks of physical selection and
empirical validation. The 173 entries in the stable result registry are
likewise not 173 original discoveries or experimental confirmations.

\textbf{What this book currently supports.} Within the explicitly
specified finite complex quantum realization, predictive responses,
actual quantum operations, error transport, purification, and finite
history protocols form an inspectable mathematical structure. The
generation results depend on their representation spectra and effective
windows. Object-RG, Wilsonian RG, and parameter running retain distinct
types. Generic Core has not uniquely derived complex quantum theory
without additional premises, nor has this book selected the gauge group,
all masses and couplings, or quantum gravity. Reproduction of standard
physics is useful for checking and reconstruction, but does not alone
supply a new UEOT-specific prediction about nature. {[}Chapters 1--2,
26, 31, 36, 47, 54--65{]}

\section{0.11 A guiding example: local records and global reversal}

Consider two ideal qubits in standard QM: a system \(S\) and a record
\(R\). Starting from \(|00\rangle\), apply a Hadamard gate to the
system, then a controlled NOT with the system as control. The recorded
joint state and reduced system state are \[
|\Psi\rangle_{SR}=\frac{|00\rangle+|11\rangle}{\sqrt2},
\qquad
\rho_S=\operatorname{Tr}_R|\Psi\rangle\langle\Psi|=\frac{I_2}{2}.
\] An observer reading the record label can predict the corresponding
pointer outcome of the system. An observer restricted to the system
cannot access the joint state's phase coherence. If both qubits are
fully controlled, applying the controlled NOT first and then the
Hadamard to the system returns the joint state exactly to
\(|00\rangle\). This reversal protocol does not first measure the record
and retain the outcome in an external apparatus. If the label has been
copied to an external reader, that reader must also be included in the
inverse; control of the system and record qubits alone is no longer
sufficient. The public label, conditional state, and global pure state
answer different experimental questions without contradiction.

This example is a conditional calculation in standard QM. UEOT's role is
to specify the allowed future operations, which records must be
retained, and whether a compressed description preserves their
responses. An added actualization mechanism that changes the statistics
after complete reversal would be a new Level 6 hypothesis. It must
provide dynamics, strength, and activation conditions, and its Level 7
test must independently calibrate, control, or conservatively bound
ordinary gate errors, leakage, and uncontrolled environments. A local
mixed state does not itself identify that mechanism. Chapters 6--8 and
38--41 extend this distinction to general record structures and testable
protocols.

\section{0.12 Bilingual revision, review, and authorship}

Edition 3.7 retains Chapters 0--66, Appendices A--F, and all 173 stable
result IDs of the current research edition 3.6, including every one of
the initial 161 IDs and the 12 subsequently merged results. The Chinese
edition adds chapter-specific reading guides and corrects identified
terminology and outdated cross-references. The English edition retains
every derivation, counterexample, source entry, and verification
boundary in independently readable academic prose. Linguistic labels
within formulas may be translated; mathematical structure and numerical
responsibility do not change with language. The translation-fidelity
report records correspondence and exceptions.

Three review lines address theoretical derivation, empirical evidence
and prior art, and language and pedagogy. They are internal AI-assisted
cross-reviews, not external peer review. The theoretical author is Degui
Qian; established scholarly results remain attributed to their original
sources. Appendix E.4 and \texttt{PUBLICATION\_\allowbreak RIGHTS.\allowbreak md} state the
boundary between historical CC BY 4.0 expression and new
edition-specific contributions. Both languages and their typeset
editions carry the same notice.

This publication revision changes neither the physical kernel nor the
Lean source. Appendices C.8, F.2, and F.4 reconcile the observed remote
acceptance with the precise reuse of existing heavy proof evidence.
Newly executed book checks, bilingual comparisons, and PDF builds are
recorded separately. Data not obtained, mechanisms not established, and
unproved limits remain open.

\section{0.13 The present physical bridges and their scope}

Edition 3.9 continues from main
\texttt{aea2fceef3c6\allowbreak 2cc8dc749b6c\allowbreak bb1b40a7d89d\allowbreak f80d}, which already
contains v31.11. Chapter 68 adds six stable results: growing carriers
for separated blue pulses, complete rectangular operations, a
same-record pulse-omission bound, and counterexamples for simultaneous
evolution and flag-only pulse identification. The book now contains 69
chapters, six appendices, and 188 results; counts are not a percentage
of physics completed. The scientific version is v31.12 / Kernel 1.12.
All 182 earlier result objects and their execution identities remain;
one explicit correction concerns the conditional mean after low-spectral
filtering. The unpublished bytes from the interrupted work were
unavailable at recovery, so this increment was rederived, executed, and
reviewed without retrospectively accepting lost executions. New results
have paper proofs and independent computations. Identity checks reuse
v31.11 machine evidence for 158 unchanged proof/configuration inputs; no
new Lean proof is claimed. Protocol-matched device data remain absent,
the empirical axis remains untested, and complex quantum structure and
effective sideband control remain independent physical inputs.

\section{0.14 The current persistent-environment and
complete-certification increment}

v31.14 / Kernel 1.14 continues from the merged and CI-verified Chapter
69 basis. Nine Chapter 70 paper results and independent computational
evidence connect one persistent two-mode environment, complete
multi-time instruments, finite propagation, joint calibration, and
holdout. The book now has 71 chapters, six appendices, 203 stable
results, and 126 ledger claims; counts are not a percentage of physics
completed. All previous 194 result objects and 117 claims retain their
identities. No new Lean proof or matched-device certification is
claimed. The strong standard baseline is equivalent with no additional
gain, and global quantum selection remains open. Actual replay,
historical execution reuse, and publication verification are recorded
separately; CURRENT and auditable receipts determine release status.

\chapter{1. Bridging UEOT Core and quantum physics}\label{ch-1}

A quantum device can accumulate a long history even when predicting a
particular future task does not require retaining it all. We construct
predictive factors from allowed experimental responses, then identify
which Core structures carry over to quantum models and why this transfer
does not itself select the quantum axioms.

\section{1.1 Operational histories and quantum responses}

A protocol \(u\) is a sequence of implementable preparations,
transformations, measurements, and permitted conditional feedback. Write
\(h\) for the record history. For a fixed future target \(Y^+\), define
the canonical response \[
C(h)=\bigl(P(Y^+\in\cdot\mid h,u)\bigr)_{u\in\mathcal U}.
\] The ``future'' is a declared record target, which may include an
entire experimental path; it does not mean every future fact about the
universe. The general probabilistic formulation uses standard Borel
history and record spaces, a common reference law \(\mu\) for histories,
a finite or countable protocol family, and chosen measurable response
kernels \(K_u(h,\cdot)\). Protocols applied only subsequently share the
same reference history law. An intervention that changes how the history
is generated requires a common structural model and an explicit account
of support. {[}S2, §§1, 2.3--4, 4.1{]}

Suppose a statistic \(M\) is sufficient for each protocol: kernels
\(L_u\) satisfy \(K_u(H,\cdot)=L_u(M(H),\cdot)\) almost surely.
Intersect the probability-one sets for all protocols and the countable
determining event classes of their future spaces. On that common set, \[
C(H)=\bigl(L_u(M(H),\cdot)\bigr)_{u\in\mathcal U}=\varphi(M(H)).
\] The right-hand side is measurable. Conversely, each \(K_u\) is a
coordinate of \(C\), so \(C\) itself is sufficient. This establishes a
minimal predictive factor for specified protocols, target, and reference
law. It neither presupposes a classical or quantum state space nor
uniquely specifies a physical carrier. {[}S2, P-PRED-01{]}

Countability matters. Let \(H\) be uniform on \([0,1]\). For every
protocol \(u\in[0,1]\), change a kernel version only at \(h=u\). Each
version is still separately equal almost everywhere to the original, but
the union of exceptional points is the entire interval.
Protocol-by-protocol equality modulo null sets therefore does not imply
a common pointwise equality for uncountably many protocols. Pointwise
specified structural kernels can instead follow a pointwise argument. A
countable dense protocol set and a common continuous version can also
support an extension, but these conditions must be supplied explicitly.
{[}S2, N-PRED-01{]}

A quantum realization provides a concrete probability generator, for
example \[
p(a_1,\ldots,a_T)=
\operatorname{Tr}\bigl[
\mathcal M_{a_T}^{(T)}\circ\cdots\circ
\mathcal M_{a_1}^{(1)}(\rho_0)\bigr].
\] In an adaptive protocol, later instruments may depend on records
already observed, never on future outcomes.

For discrete records, if the protocol family is closed under action
prefixes and permitted continuations after positive-probability
observations, the current response also determines the updated response:
\[
P(B\mid h,a,o)=\frac{P(o,B\mid h,a)}{P(o\mid h,a)},\qquad P(o\mid h,a)>0.
\] Both numerator and denominator are retained by the current canonical
response. On a zero-probability branch the ratio does not determine a
conditional state; choosing an arbitrary mathematical extension adds no
empirical support. This denominator condition recurs in quantum
instruments, finite coordinates, and postselection errors. {[}S2,
P-PRED-03{]}

\section{1.2 General Core does not select a quantum sector}

\phantomsection\label{p-qbr-01}{}

\textbf{P-QBR-01: A classical realization countermodel.} General
predictive factors, state closure, access carriers, persistence, and
control structures can coexist in a finite classical Markov system.
These structures alone therefore do not logically entail noncommuting
observables, a complex Hilbert sector, or quantum purification.

\textbf{Proof.} Take a finite controlled chain, partition histories by
all their future responses, and realize them by state readout and
classical probability kernels. Sufficient access families, their minimal
members, and hitting sets can all be defined. An attracting recurrent
class supplies persistence, and a finite discounted task admits an
optimal policy. Yet the state space remains a probability simplex, and
all classical effects admit joint readout. A mixed classical state
cannot be the marginal of a pure classical joint state, since a pure
joint state is a deterministic point. The quantum-specific conclusions
are thus not logical consequences of the general premises. \(\square\)

This countermodel does not negate UEOT's unifying organizational role.
It identifies the physical content that requires sector-specific
principles.

\section{1.3 Representation, gauge, and observable equivalence}

A representation change \(T\) must transport preparations, channels,
effects, and composition rules together. If every allowed protocol
probability is preserved, the realizations belong to the same
operational equivalence class for those protocols. An invertible
coordinate change is not automatically a physical gauge symmetry.
Whether a gauge orbit represents redundancy is determined by the
particular physical model.

For a fixed target and protocol family, physical gauge invisibility can
allow a response to descend to a quotient. However, ``no experiment
performed so far has distinguished them'' is weaker than ``no physical
experiment can ever distinguish them.'' Finite tests supply finite
certificates.

\section{1.4 How the four counted generators enter physical derivations}

Section 6.2 of S4 gives the exact scope of the four-generator register:
of 106 items, 11 are \texttt{generated}, 89 are
\texttt{retained\_\allowbreak adapter}, and 6 are \texttt{retained\_\allowbreak boundary}. Their
nonredundancy means
\texttt{nonredundant\_\allowbreak under\_\allowbreak declared\_\allowbreak derivation\_\allowbreak system} within the
registered derivation system. It is not absolute independence in every
mathematical proof language, nor an unconditional derivation of 106
conclusions from four physical axioms. The four mechanisms have the
following concrete roles in quantum applications.

\textbf{M-QD: Quotient and descent.} Given maps \(q:X\to C\) and
\(F:X\to Y\), suppose \[
q(x)=q(x')\Longrightarrow F(x)=F(x'),
\] Then define \(\bar F(q(x))=F(x)\) on \(q(X)\). It is independent of
the representative, and \(F=\bar F\circ q\). For finite sets this
completes the construction; measurable spaces additionally require
suitable measurable-factorization conditions. Core predictive minimality
and structured internal--environmental decompositions use this pattern.
The quantum question is whether the Born responses of all allowed
records depend only on the retained finite coordinates. Proving descent
does not also select the complex scalar field or the Born pairing.
{[}S2, P-PRED-01, P-INT-02; S3, appendix mapping all 106 items{]}

\textbf{M-TC: Finite quantitative transport.} Let the reference
trajectory satisfy \(y_{t+1}=F_t(y_t)\) and define \(e_t=d(x_t,y_t)\) in
the common metric. If the defect between the actual evolution
\(x_{t+1}\) and a reference update \(F_t(x_t)\) is at most
\(\varepsilon_t\), and \(F_t\) has Lipschitz constant \(L_t\), inserting
the intermediate reference image \(F_t(x_t)\) gives \[
e_{t+1}\le L_te_t+\varepsilon_t,
\] Induction yields \[
e_n\le\left(\prod_{j<n}L_j\right)e_0+
\sum_{k<n}\varepsilon_k\prod_{j=k+1}^{n-1}L_j.
\] If quantum channels contract the chosen operational distance, then
\(L_t\le1\) and the error is bounded by the initial error plus the sum
of stepwise defects. But contraction in that distance must be proved
from quantum positivity, normalization, and effect pullback. It is not
supplied automatically by the transport template. Postselection
denominators, memory ancillas, and protocol support also remain
necessary. {[}S2, P-DYN-03--04, P-ID-01--02; §4.4 of this book{]}

\textbf{M-PE: Positive eigenstructure.} For a primitive nonnegative
substochastic matrix \(Q\), positive left and right Perron data satisfy
\[
qQ=rq,\quad Qh=rh,\quad q\mathbf1=1,\quad qh=1.
\] Consequently \(qQ^n=r^nq\), and conditioning on survival preserves
\(q\). Moreover, \[
Q^\uparrow_{ij}=\frac{Q_{ij}h_j}{rh_i},\qquad
\sum_jQ^\uparrow_{ij}=1,\qquad
\sum_iq_i h_iQ^\uparrow_{ij}=q_jh_j.
\] Positivity, spectral data, and normalization jointly produce a
quasi-stationary distribution and a conditioned process. Quantum CP maps
also act on positive cones, so analogous positive-operator spectral
questions can be studied. Irreducibility, primitivity, positive
eigenoperators, spectral gaps, and the type of conditioning must
nevertheless be established for the map at issue. Entrywise
nonnegativity of a classical matrix and complete positivity of a quantum
map are different input types. Chapter 53 proves a concrete bridge in a
diagonal measure-and-prepare model, directly using the Core spectral
theorem. General noncommutative spectral obligations remain separate.
{[}S2, §10.2, P-QSD-02; S4, §6.2{]}

\textbf{M-OI: Occupation limits and invariant closure.} For a fixed CPTP
channel \(\mathcal T\) on a finite quantum system, form the Cesàro
average of time-indexed states, \[
\bar\rho_N=\frac1N\sum_{t=0}^{N-1}\mathcal T^t(\rho_0).
\] The finite-dimensional density-matrix set is compact, hence a
convergent subsequence exists. Also, \[
\mathcal T(\bar\rho_N)-\bar\rho_N=
\frac{\mathcal T^N(\rho_0)-\rho_0}{N},\qquad
\|\mathcal T(\bar\rho_N)-\bar\rho_N\|_1\le\frac2N.
\] Continuity of the linear map implies that every subsequential limit
is invariant. This constructs invariant structure from occupation
averages. It does not prove uniqueness, convergence of the state at each
time, or an empirical limit along every random trajectory. Compactness
cannot be transferred unchanged to infinite dimension. A periodic
unitary evolution may have such an averaged limit without ever
converging to a single state. {[}S2, P-PER-02, P-GOA-01; the finite
quantum version here is a direct derivation in the declared CPTP
model{]}

Chapter 52 formalizes compactness of the actual density set, the Cesàro
matrix residual, the \(1/N\) bound for all effects, and existence of a
stationary subsequential limit in Lean. Its machine proof uses a \(2/N\)
bound in the Euclidean operator norm to show that the matrix residual
tends to zero. The trace-norm formula in this section retains its
analytical proof responsibility; equality of the constants does not make
it part of that formalization automatically.

These mechanisms explain why the kernel is reusable. The quantum layer
must still supply complex matrices, positive semidefiniteness,
instrument completeness, the actual trace pairing, and experimental
interfaces. A reusable proof is not another fundamental law of nature,
and retaining a specialized adapter is not a failure of Core.

\section{1.5 Objects, realizations, and the quantum meaning of Ω}

A quantum object certificate first records a candidate, its own future,
the permitted operations, and its environmental interfaces. It then
records predictive and recursive closure, intervention conditions,
transport along histories, and errors. Recovery, error correction,
self-maintenance, and functional objectives are additional modules when
the claim requires them. A pure-state label, a topological number, or a
finite matrix cannot replace this certificate. {[}S2, D-OBJ,
D-OBJ-STRONG{]}

For a finite set \(V\) of access interfaces, suppose adding accessible
interfaces only enlarges the set of permitted experiments. The
sufficient-access family is then upward closed. Let \(\mathcal M_0\) be
its minimal members. After deleting access permissions \(D\), the target
remains predictable exactly when \[
\exists M\in\mathcal M_0:\ M\cap D=\varnothing.
\] Negating this statement gives the condition that every minimal
realization is hit. Minimal hitting sets are access blockers. A quantum
realization must specify whether access means classical records, local
measurement, joint measurement, or coherent control. Stronger access
cannot simply be assumed implementable. {[}S2, §§13.1--2{]}

Physical deletion is defined separately. For a fixed certificate and
horizon \(T\), let \(\chi_{O,T}(D)=1\) mean that the certificate
survives removal of mechanism set \(D\). If failure is upward monotone,
then \[
\mathcal C_{\Omega,T}=\min_{\subseteq}\{D:\chi_{O,T}(D)=0\},\qquad
\chi_{O,T}(D)=0\iff\exists C\in\mathcal C_{\Omega,T},\ C\subseteq D.
\] The forward direction chooses a minimal member among finite failure
subsets; the reverse uses monotonicity. Here
\(\Omega_T=(\mathcal M,\mathcal C_{\Omega,T})\) is the failure
hypergraph for the specified mechanisms and certificate. It connects to
the original geometric-cycle intuition without making nonzero homology a
universal threshold for existence. If deleting an inhibitor restores the
system, failure is no longer monotone and the full intervention response
must be retained. {[}S2, P-OMG-01{]}

Interpreting an access blocker as physical failure additionally requires
a frame condition: surviving readouts and dynamics are preserved;
deletion semantics match the sets; recovery, power supply, and other
relevant mechanisms are included; and sufficiency of the remaining
access is actually equivalent to the target certificate. Disconnecting a
power port may also destroy a memory that is still nominally
``accessible.'' This is why the bridge is necessary. {[}S2, §13.3{]}

\section{1.6 The additional tasks behind GOD, GOA, and Π/Φ}

Specifying a quantum preparation and instruments does not specify a
payoff. An evaluation bridge must at least give the permitted policies,
reward or risk functional, horizon, resource budget, and evaluating
agent or level. On the same quantum memory, preserving an unknown state,
rapidly reading a label, estimating an external field, and cooling to
the ground state may call for different actions. Neither the object nor
the device need change. {[}S2, D-BRIDGE, N-TEL-01; Q2, §20{]}

A minimal counterexample uses the same process with available actions
\(a,b\). The objectives \(J_a=\mathbf1_{\{A_0=a\}}\) and
\(J_b=\mathbf1_{\{A_0=b\}}\) select opposite optimal actions. The
dependency used here is therefore \[
\begin{gathered}
(\text{quantum process},\text{object certificate},\text{evaluation bridge})\longrightarrow V^*\\
\longrightarrow\mathcal A^*\longrightarrow\text{typed long-term structure}.
\end{gathered}
\] For example, with finitely many actions, discount \(0<\beta<1\), and
bounded reward, the Bellman operator on conditional states is \[
(TV)(\rho)=\max_a\left\{r(\rho,a)+\beta\sum_o p(o\mid\rho,a)
V(\rho_{a,o})\right\}.
\] An arbitrary placeholder state may be assigned on a zero-probability
branch because its weighted contribution is zero. Since \(p_o\ge0\) and
the probabilities sum to one,
\(\|TV-TW\|_\infty\le\beta\|V-W\|_\infty\). Under the required
measurability and state-space conditions, the fixed point and optimal
selector follow from standard discounted-control constructions. GOD is
first the set of maximizing actions; it need not be single-valued or a
gradient in Euclidean space. {[}S2, §19; Q2, §20.2{]}

GOA must specify whether the subject is an invariant state, a Cesàro
limit, an orbit's \(\omega\)-limit set, a QSD, or a nonautonomous
tracking object. Even a unique invariant probability need not imply
mixing at individual times: deterministic alternation between two states
oscillates outside its unique uniform invariant law. Unitary conjugation
likewise permits periodic or quasiperiodic motion. Thus ``following an
optimal direction must converge to a unique optimal fixed point'' is not
a general theorem. {[}S1, §2.3.2, pp.~6--7; S2, §21; S4, §4{]}

If \(J=\Pi-\lambda\Phi-C\) is used, the dimensions, signs, and
independent measurements of \(\Pi,\Phi,C\) must first be specified. An
arbitrary \(\chi\) gives \[
\Pi'=\Pi+\lambda\chi,\quad \Phi'=\Phi+\chi,
\qquad \Pi'-\lambda\Phi'-C=J.
\] The same optimal behaviour therefore cannot uniquely identify a
physical axis and an informational axis. Energy costs of quantum
control, preparation time, memory requirements, predictive KL
divergence, or Bayesian risk can define concrete bridges. Fisher
information, entropy, information deficit, and information gain cannot
be added with the same sign merely because each is called
``information.'' {[}S2, §§23.1--3, 23.6{]}

Suppose one further proposes a two-dimensional drive
\(z(B)\in\mathbb R^2\) shared across environments, with responses
\(m_e(B)=g_e(z(B))\). At the same budget point, the chain rule gives
only \[
\begin{bmatrix}Dm_1(B)\\\vdots\\Dm_E(B)\end{bmatrix}
=\begin{bmatrix}Dg_1(z(B))\\\vdots\\Dg_E(z(B))\end{bmatrix}Dz(B),
\qquad\operatorname{rank}\le2.
\] This is a falsifiable necessary condition, not a universal
two-dimensional law derived from quantum formalism. Rank two in each
environment separately, or low rank after combining derivatives from
different budget points, does not replace the shared same-point test. A
scalar reward also does not force a two-dimensional response:
\(J(a,B)=-\|a-B\|^2/2\) has optimizer \(a=B\) in any dimension. {[}S2,
P-DDH-04; S3, C5 work package; S4, §5{]}

\section{1.7 Seven scientific ports and the quantum kernel}

The seven ports organize open research questions. They are neither new
axioms nor seven stages that must be completed sequentially. In the
quantum branch, they assign the following concrete responsibilities.
{[}S2, §31.2; S3, §§3--10; S5, §§2--7{]}

\begin{longtable}[]{@{}
  >{\raggedright\arraybackslash}p{(\columnwidth - 4\tabcolsep) * \real{0.2500}}
  >{\raggedright\arraybackslash}p{(\columnwidth - 4\tabcolsep) * \real{0.3400}}
  >{\raggedright\arraybackslash}p{(\columnwidth - 4\tabcolsep) * \real{0.4100}}@{}}
\toprule\noalign{}
\begin{minipage}[b]{\linewidth}\raggedright
Port
\end{minipage} & \begin{minipage}[b]{\linewidth}\raggedright
Precise connection to QM
\end{minipage} & \begin{minipage}[b]{\linewidth}\raggedright
Claim that the connection can support
\end{minipage} \\
\midrule\noalign{}
\endhead
\bottomrule\noalign{}
\endlastfoot
C1: existence and general states & Infinite-dimensional normal states,
measurable conditional versions, hidden memory, and noncompact limits;
explicit extension conditions for each finite-dimensional proof &
Existence, recursion, or limit theorems in a specified mathematical
model \\
C2: object and protocol discovery & Find sufficient response spaces from
permitted instruments and records; separate candidate search from
independent verification & A certificate within a given operation
library, not identification of every unknown device \\
C3: realization and damage & Access permissions, joint quantum readout,
correction/recovery mechanisms, and physical failure graphs & Minimal
realization or failure families under stated operational semantics \\
C4: identity and long-term transport & Encoding changes, drifting
devices, time-dependent channels, and transport of stationary or
metastable states & Transport bounds for a fixed distance, support, and
time range \\
C5: shared low-dimensional mechanisms & Independently anchor
energy/information budgets and test cross-environment response rank at
the same point & Rejection or support of a specified local
low-dimensional hypothesis, not a unique naming of drives \\
C6: identifying the purpose bridge & Independently determine the control
task or actual learning/selection mechanism; predict actions at budgets
not used for fitting & Behavioural support for a specified mechanism,
not an intention inferred from state stability \\
C7: statistics and reproduction & Freeze protocols, errors, and stopping
rules; preserve raw records, source versions, and rejection criteria &
Separate mathematical proof, model checking, experimental support, or
hypothesis rejection \\
\end{longtable}

An exact finite matrix can completely settle a computational instance.
Its Lean proof may also establish properties of that instance for every
finite record length. Whether a device realizes that matrix, its noise
belongs to the declared model, and training and testing are independent
remain physical and statistical responsibilities under C2/C7. Historical
work packages labelled PLANNED in S5 do not become complete merely
because a kernel now exists. {[}S5, status notes; §5.5 and the later
kernel-evidence chapters{]}

\chapter{2. Quantum selection: when complex quantum theory is
selected}\label{ch-2}

Classical probability, real quantum theory, and complex quantum theory
can each describe operational statistics; the question is which
additional principles distinguish them. After fixing the operational
background, purification and local-tomography countermodels clarify the
conditions of a reconstruction theorem without turning its use into an
unconditional derivation of QM from Core.

\section{2.1 A fixed operational-theory background}

We work in a finite-dimensional causal operational probabilistic theory.
The reconstruction background includes systems and finite-outcome tests,
the trivial probabilistic system, sequential and parallel composition
with compatibility laws, nonempty normalizable states, deterministic
effects, coarse-graining and conditional tests, finite-dimensional
closed normalized state sets, nontrivial distinguishable systems, and
equivalence under exchanging the order of composite systems. Equality of
transformations is defined by complete operational statistics
\textbf{including arbitrary auxiliary systems}. This is stronger than
comparing only single-system inputs and outputs. The full background is
fixed by the definitions and twelve conditions in Appendices A.1--A.2 of
the original paper cited below.

The background does not presuppose Hilbert spaces, purification, local
tomography, unrestricted measurement, or all CPTP channels. Defining
channel equality at the outset solely by local input--output equality
would make CIT a restatement of a definition, removing the substantive
condition required by the reconstruction theorem.

Quantum selection must be stated in a complete operational theory of
this kind. A state cone alone cannot reconstruct every physical
transformation.

\section{2.2 Two independent sector principles}

\textbf{H-QSEL-1: Local equivalence / CIT.} For channels
\(\Lambda,\mathcal E\) from system \(A\) to system \(B\), if
\(f(\Lambda\rho)=f(\mathcal E\rho)\) for every input state \(\rho\) and
every output effect \(f\), then \(\Lambda=\mathcal E\), including
operational equivalence in the presence of every auxiliary system.
Within the fixed background above, this is equivalent to local
distinguishability of composite states.

\textbf{H-QSEL-2: Equivalent-system purification (ESP).} For each system
\(A\), there is an operationally equivalent purifying system
\(\widetilde A\) such that every normalized state is a marginal of a
pure state of \(A\widetilde A\). Any two purifications of the same state
in this purifying system differ only by a reversible channel on
\(\widetilde A\).

\textbf{K-QSEL-01.} Under that fixed background and these two
principles, Nakahira's Theorem 1 characterizes finite-dimensional
complex operational quantum theory: all density matrices, POVMs, the
Born pairing, Hilbert-space tensor products, and CPTP channels for the
systems in question. This book invokes the external reconstruction; its
full proof is not included in the UEOT Lean kernel. See the
\href{https://arxiv.org/abs/2605.23217}{original preprint} and
\href{https://arxiv.org/html/2605.23217v1}{v1 text with Appendices
A--D}. Appendix A.3 explicitly distinguishes recovering quantum
structure for the systems present in a theory from proving that a
physical system of every natural-number dimension must exist.

To keep later chapters self-contained, we state and use these quantum
structures explicitly. The external paper is not a substitute for
explaining the calculation rules to the reader. It supplies the source
for the claim that this set of selection principles is sufficient.

\section{2.3 Independence and satisfaction in quantum theory}

\phantomsection\label{p-qsel-01}{}

\textbf{P-QSEL-01: Classical and real-quantum deletion witnesses.}
Finite classical probability theory satisfies local tomography but
cannot purify a nontrivial mixed state. Real quantum theory has
equivalent-system purification but does not satisfy general local
tomography.

\textbf{Proof.} A classical joint distribution is recovered from
products of local deterministic-event probabilities. A pure joint
distribution is a deterministic point, whose marginals cannot be mixed.
A real symmetric density matrix admits purification by real spectral
decomposition; two purification matrix factors differ by an orthogonal
transformation on the purifying space. For two rebits, let \[
\rho_\pm=\tfrac14(I\otimes I\pm Y\otimes Y).
\] The matrix \(Y\otimes Y\) is real symmetric, and \(\rho_\pm\) are
distinct positive semidefinite states. Single-rebit real symmetric
effects are spanned by \(I,X,Z\) and are orthogonal to \(Y\). Every
product effect therefore assigns the same probability to both states.
\(\square\)

The counterexample can be checked explicitly. Take
\(Y=\left(\begin{smallmatrix}0&-i\\i&0\end{smallmatrix}\right)\).
Although \(Y\) itself is not a real-quantum observable, \(Y\otimes Y\)
is real symmetric, squares to \(I\), and has trace zero. Its eigenvalues
are \(+1,-1\), each with multiplicity two. Both \(\rho_\pm\)
consequently have spectrum \(\{1/2,1/2,0,0\}\). For real symmetric
\(E,F\), \(\operatorname{Tr}(EY)=\operatorname{Tr}(FY)=0\), so \[
\operatorname{Tr}[(E\otimes F)(\rho_+-\rho_-)]=
\tfrac12\operatorname{Tr}(EY)\operatorname{Tr}(FY)=0.
\] The joint effect \((I+Y\otimes Y)/2\), however, gives probabilities 1
and 0. Local effects thus fail to separate all joint states. No
forbidden single-rebit effect has been inserted into the experiment: the
gap in local tomography is exhibited within the allowed real joint-state
space. {[}Q2, P-QSEL-01, with the explicit calculation added here{]}

In complex QM, local tomography follows because tensor products of
Hermitian bases span the full composite Hermitian space. Spectral
decomposition \(\rho=\sum_i p_i|i\rangle\langle i|\) gives the
purification \(|\Psi\rangle=\sum_i\sqrt{p_i}|i\rangle|i\rangle\).
Uniqueness of purification follows from Schmidt decomposition and
factorization with the same spectrum.

The relative indispensability of these two principles does not mean that
``QM requires exactly two axioms in every language.'' Combining them
into one conjunction or moving them into the background changes the
axiom count without changing the logic. Purification and
local-tomography mathematics in the historical CRP/CIT formulation
belong to existing theory; UEOT supplies an interpretation in terms of
parent systems and interfaces.

\section{2.4 Limits on the use of reconstruction}

There are two explicit routes to the formulas used later. The
reconstruction route accepts the fixed OPT background, CIT, and ESP, and
invokes K-QSEL-01. The direct modelling route instead takes
finite-dimensional complex Hilbert space, Born probabilities, tensor
composition, and CP instruments as quantum-sector inputs. Either
supports the subsequent mathematics. Inputs of the latter route must not
be relabelled as conclusions already derived from Core by the former.

A finite or countable collection of experiments constrains a device's
response set. Finite success alone cannot establish the universal
quantifiers over systems, states, and transformations in CIT or ESP. One
successful purification preparation also does not determine a unique
ontology of the actual environment. Approximate purification or
tomography at finite precision requires a separate stability theorem for
reconstruction. Replacing exact equalities by small numerical errors is
not sufficient to transfer an exact reconstruction theorem.

\chapter{3. Fisher action, complex amplitudes, and conditional Born
arguments}\label{ch-3}

How do a probability density and phase combine into Schrödinger
dynamics, and when does a statistical inequality become quantum
uncertainty? The selected action, momentum operator, and
branch-martingale assumptions carry different parts of the argument and
must not be collapsed into a single claim that an information principle
yields the Born rule.

\section{3.1 A continuous statistical action}

On a smooth region without nodes, take \(\rho>0\) and a real phase
action \(S\), with the required decay or boundary conditions. Given a
mass \(m>0\), potential \(V\), and action scale \(\hbar>0\), define \[
\mathcal A[\rho,S]=\int dt\,dx\left[
\rho\left(\partial_tS+\frac{|\nabla S|^2}{2m}+V\right)
+\frac{\hbar^2}{8m}\frac{|\nabla\rho|^2}{\rho}\right].
\] This local, second-order action with its particular Fisher
coefficient is an \textbf{H} input. Arbitrary object closure does not
select it automatically. Electromagnetic coupling additionally requires
an independent specification of \(\nabla S-qA\), \(q\phi\), and gauge
transformations.

\phantomsection\label{p-qact-01}{}

\textbf{P-QACT-01: Conditional Madelung--Schrödinger equivalence.}
Compactly supported variations of the action yield \[
\partial_t\rho+\nabla\cdot(\rho\nabla S/m)=0,
\] \[
\partial_tS+\frac{|\nabla S|^2}{2m}+V
-\frac{\hbar^2}{2m}\frac{\Delta\sqrt\rho}{\sqrt\rho}=0.
\] For \(\psi=\sqrt\rho\,e^{iS/\hbar}\), these are equivalent to the
real and imaginary parts of \[
i\hbar\partial_t\psi=
-\frac{\hbar^2}{2m}\Delta\psi+V\psi
\] on that smooth, node-free region.

\textbf{Proof.} Put \(S\mapsto S+\epsilon\eta\), where \(\eta\) has
compact support. Then \[
\left.\frac{d\mathcal A}{d\epsilon}\right|_0
=\int\rho\left(\partial_t\eta+\frac{\nabla S}{m}\cdot\nabla\eta\right)\,dt\,dx
=-\int\eta\left[\partial_t\rho+\nabla\cdot(\rho\nabla S/m)\right]dt\,dx.
\] Arbitrary \(\eta\) gives the continuity equation. Next vary \(\rho\),
taking the parameter sufficiently small to retain local \(\rho>0\). The
Euler derivative of the Fisher term is \[
\frac{\hbar^2}{8m}
\left(-2\frac{\Delta\rho}{\rho}
+\frac{|\nabla\rho|^2}{\rho^2}\right)
=-\frac{\hbar^2}{2m}\frac{\Delta\sqrt\rho}{\sqrt\rho}.
\] Writing \(R=\sqrt\rho\), direct differentiation gives \[
\partial_t\psi=e^{iS/\hbar}(\partial_tR+iR\partial_tS/\hbar),
\] \[
\Delta\psi=e^{iS/\hbar}\left[\Delta R+
\frac{2i}{\hbar}\nabla R\cdot\nabla S+
\frac{iR}{\hbar}\Delta S-\frac{R}{\hbar^2}|\nabla S|^2\right].
\] Substitute into the Schrödinger equation. Dividing its real part by
\(R>0\) gives the Hamilton--Jacobi equation; multiplying its imaginary
part by \(2R\) gives continuity. Each step is reversible on the declared
region. \(\square\)

Here \(|\psi|^2=\rho\) is part of the change of variables. It shows how
the specified statistical action is consistent with wavefunction
dynamics; it does not independently prove the POVM Born rule for every
observable or handle nodes, vortices, and globally multivalued phases. A
global wavefunction must satisfy suitable line-bundle or boundary
conditions. Integer phase winding on a closed path constrains only the
phase integral.

The coefficient and units are premises of the result. Both \(S\) and
\(\hbar\) have units of action, and \(|\nabla S|^2/(2m)\) has units of
energy. The Fisher coefficient \(\hbar^2/(8m)\) is what combines the two
equations into the stated linear Schrödinger equation. A different
coefficient, a different meaning of probability density, or additional
gradient terms require a fresh derivation. The original Π/Φ motivation
does not uniquely fix these choices. {[}S1, §4.2.1, pp.~30--32; S2, §23;
Q2, P-QACT-01{]}

\section{3.2 When a Fisher bound becomes position--momentum uncertainty}

Let a normalized one-dimensional density \(\rho\) have finite variance
and position Fisher information \(I_x=\int(\rho')^2/\rho\,dx\), with
vanishing boundary terms.

\phantomsection\label{p-qact-02}{}

\textbf{P-QACT-02.} \[
\operatorname{Var}(X)I_x\ge1.
\] If we also use the wavefunction above and quantum momentum
\(P=-i\hbar\partial_x\) on the required Sobolev domain, then \[
\operatorname{Var}(P)=\operatorname{Var}_\rho(S')
+\frac{\hbar^2}{4}I_x,
\quad \Delta X\Delta P\ge\hbar/2.
\]

\textbf{Proof.} Integration by parts gives \(\int(x-\bar x)\rho'=-1\),
hence \[
1=\left|\int(x-\bar x)\sqrt\rho\frac{\rho'}{\sqrt\rho}\,dx\right|^2
\le\left[\int(x-\bar x)^2\rho\,dx\right]
\left[\int\frac{(\rho')^2}{\rho}\,dx\right].
\] For \(\psi=\sqrt\rho e^{iS/\hbar}\), \[
P\psi=e^{iS/\hbar}\left(\sqrt\rho\,S'-i\hbar\frac{\rho'}{2\sqrt\rho}\right).
\] Vanishing boundary terms give \(\langle P\rangle=E_\rho S'\). The
real and imaginary parts have no cross term in the squared modulus, so
\(\|P\psi\|^2=E_\rho(S')^2+\hbar^2I_x/4\). As a second moment of the
self-adjoint momentum spectral measure, this is meaningful for
\(\psi\in D(P)\) with the relevant integrals finite; \(P^2\psi\) need
not already exist. Subtracting the squared mean yields the variance
decomposition. The first inequality then gives
\(\Delta X\Delta P\ge\hbar/2\). \(\square\)

Estimation theory and quantum uncertainty are thus explicitly connected,
but the momentum operator and amplitude--phase structure supply an
additional bridge. Finite information capacity alone does not establish
the CCR.

In general probability theory, the distributions of \(X\) and another
real random variable \(P\) can be independently specified with both
variances arbitrarily small. Fisher information does not identify them
as canonically conjugate variables. The momentum-variance identity above
is precisely the bridge from a statistical inequality to physical
uncertainty; ``finite information'' cannot replace it. {[}S1, §4.2.1,
p.~30, printed lines 1468--1487{]}

\section{3.3 Quadratic forms and the meaning of a ``two-phase test''}

\textbf{K-QBORN-01.} On a complex vector space, let a continuous
nonnegative function \(F\) satisfy \(F(\lambda x)=|\lambda|^2F(x)\) and
the parallelogram identity for all \(x,y\). Complex polarization gives a
unique positive semidefinite Hermitian form \(B\) with \(F(x)=B(x,x)\).
In finite dimension this is \(F(x)=x^\dagger Ax,A\succeq0\). This is the
standard characterization of quadratic forms and inner products, not a
global conclusion automatically proved by finitely many observations.

Using the convention that the inner product is linear in its second
argument, \[
B(x,y)=\frac14\{F(x+y)-F(x-y)
-iF(x+iy)+iF(x-iy)\}.
\] If each event \(a\) has such a form \(F_a(x)=x^\dagger E_ax\), and
\(\sum_aF_a(x)=\|x\|^2\) for every vector \(x\), then \[
x^\dagger\left(\sum_aE_a-I\right)x=0\quad\forall x.
\] Uniqueness under polarization implies \(\sum_aE_a=I\); together with
\(E_a\succeq0\), this gives \(E_a\preceq I\). We have therefore obtained
a POVM. If classical random mixtures are transported convex-linearly,
the probability for \(\rho=\sum_ip_i|\psi_i\rangle\langle\psi_i|\) is
\(\sum_ip_iF_a(\psi_i)=\operatorname{Tr}(E_a\rho)\). The derivation
explicitly depends on quadraticity, complex-vector structure, total
probability, and the mixing rule.

In a finite experiment, ``two phases'' generally means recovering a
declared quadratic model through real and imaginary interference
coordinates. It does not mean that two measured numbers exclude every
nonquadratic theory.

A response function can be changed outside a finite tested set while
retaining every measured value. The valid conclusion of the original
two-angle programme must therefore retain its model family, quantifiers,
and errors. A finite pass is not unconditional uniqueness of the Born
rule.

\section{3.4 Unbiased branch weights and actual outcomes}

\phantomsection\label{p-qborn-02}{}

\textbf{P-QBORN-02: Probabilities of an unbiased final branch.} Suppose
branch weights \(p_a(t)\in[0,1]\) are martingales with respect to a
common filtration and converge almost surely to a one-hot vector
\(1_{\{A=a\}}\). Then \[
P(A=a\mid\mathcal F_0)=p_a(0).
\]

\textbf{Proof.} For every \(B\in\mathcal F_0\), the martingale property
gives \(E[\mathbf1_Bp_a(t)]=E[\mathbf1_Bp_a(0)]\). Boundedness permits
dominated convergence, yielding
\(E[\mathbf1_B\mathbf1_{\{A=a\}}]=E[\mathbf1_Bp_a(0)]\), exactly the
conditional-probability identity. If the initial preparation is fixed so
that \(p_a(0)\) is constant, expectation gives \(P(A=a)=p_a(0)\). With a
random initial preparation, \(P(A=a)=Ep_a(0)\). \(\square\)

When the initial weights are \(|c_a|^2\), this gives Born outcome
probabilities for that actualization model. Unbiased martingales,
eventual selection of a single branch, and a concrete physical mechanism
are independent premises. Objectivity of public records has not
established them.

\chapter{4. Finite-dimensional quantum operations: states, effects,
channels, and instruments}\label{ch-4}

Every later finite quantum calculation uses the operational grammar
introduced here: states describe preparations, effects describe events,
and instruments determine both probabilities and postmeasurement states.
Choi--Kraus constructions make these objects checkable as matrices,
while distance bounds connect implementation errors to final
experimental probabilities.

\section{4.1 Basic conventions}

Let \(\mathcal H\cong\mathbb C^d\) be a finite-dimensional complex
Hilbert space, with inner product linear in the second argument. States
satisfy \(\rho\succeq0,\operatorname{Tr}\rho=1\), effects satisfy
\(0\preceq E\preceq I\), and probabilities are \[
p(E|\rho)=\operatorname{Tr}(\rho E).
\] This trace is real and belongs to \([0,1]\). Indeed,
\(\rho^{1/2}E\rho^{1/2}\succeq0\) has nonnegative trace, so cyclicity
gives \(\operatorname{Tr}(\rho E)\ge0\). Replacing \(E\) by \(I-E\)
gives \(1-\operatorname{Tr}(\rho E)\ge0\). This reasoning uses the
declared quantum states, effects, and pairing; an arbitrary complex
matrix expression is not thereby a probability model.

A POVM is a family \(\{E_a\}\) with \(\sum_aE_a=I\). Projective
measurement is the special case of mutually orthogonal projections. A
general POVM does not determine the postmeasurement state.

Composite systems use \(\mathcal H_A\otimes\mathcal H_B\). The local
state is uniquely determined by
\(\operatorname{Tr}(\rho_AE)=\operatorname{Tr}[\rho_{AB}(E\otimes I)]\).
A separable state has the form \(\sum_i p_i\rho_i^A\otimes\rho_i^B\);
states without such a representation are entangled.

\section{4.2 Complete positivity, Choi matrices, and Kraus operators}

A linear map
\(\mathcal E:\mathcal B(\mathcal H_A)\to\mathcal B(\mathcal H_B)\) is
completely positive if \(\operatorname{id}_n\otimes\mathcal E\)
preserves positivity for every auxiliary dimension \(n\). We use the
unnormalized Choi convention \[
J(\mathcal E)=\sum_{ij}|i\rangle\langle j|\otimes
\mathcal E(|i\rangle\langle j|),\qquad
\mathcal E(\rho)=\operatorname{Tr}_A[(\rho^\top\otimes I)J].
\]

\phantomsection\label{p-qop-01}{}

\textbf{P-QOP-01: Finite Choi--Kraus equivalence.} \[
\mathcal E\ {\rm CP}\iff J(\mathcal E)\succeq0
\iff\mathcal E(\rho)=\sum_\mu K_\mu\rho K_\mu^\dagger.
\] Trace preservation is equivalent to
\(\sum_\mu K_\mu^\dagger K_\mu=I_A\), and also to
\(\operatorname{Tr}_B J=I_A\).

\textbf{Proof.} Applying a CP map to the unnormalized maximally
entangled projector gives a positive Choi matrix. A positive Choi matrix
has a spectral decomposition
\(\sum_\mu|K_\mu\rangle\rangle\langle\langle K_\mu|\), where
\(|K\rangle\rangle=\sum_i|i\rangle\otimes K|i\rangle\). Substitution
into the reconstruction formula, entry by entry, gives the Kraus
representation. With any auxiliary space, the Kraus expression remains a
sum of positive terms. Cyclicity of the trace for arbitrary inputs gives
the trace-preservation condition. \(\square\)

More explicitly, write
\(J=\sum_\mu\lambda_\mu|v_\mu\rangle\langle v_\mu|\) and reshape
\(\sqrt{\lambda_\mu}v_\mu\) into \(K_\mu\) using the input and output
indices. Then \[
\operatorname{Tr}\mathcal E(X)
=\operatorname{Tr}\left[X\sum_\mu K_\mu^\dagger K_\mu\right]
\] equals \(\operatorname{Tr}X\) for every \(X\) exactly when the
bracketed operator is \(I\). Changing to a normalized maximally
entangled vector or reversing input/output order in the Choi convention
requires the corresponding change of coefficients or transpose placement
in the partial-trace and reconstruction formulas.

Positivity alone is insufficient. Transposition preserves positivity of
a single-system density matrix, but applying it on one side of a Bell
state gives a partial transpose with a negative eigenvalue. With
arbitrary entangled ancillas allowed, CP is therefore a physical
compatibility condition arising from the specified composite sector.

\section{4.3 Instruments and normalization}

Each branch of an instrument \(\{\mathcal M_a\}\) is CP and
trace-nonincreasing, and their sum is CPTP: \[
p_a=\operatorname{Tr}\mathcal M_a(\rho),\quad
E_a=\mathcal M_a^\dagger(I),\quad
\rho_a=\mathcal M_a(\rho)/p_a\quad(p_a>0).
\] At \(p_a=0\), division by zero cannot define a certified posterior.
The Lüders update \(\mathcal M_a(\rho)=P_a\rho P_a\) is a particular
instrument, not a consequence of the POVM alone.

A finite Kraus realization is
\(\mathcal M_o(X)=\sum_kK_{o,k}XK_{o,k}^\dagger\). Completeness of the
whole instrument requires \(\sum_o\sum_kK_{o,k}^\dagger K_{o,k}=I\),
giving \(E_o=\sum_kK_{o,k}^\dagger K_{o,k}\). Different branches may
have different numbers of Kraus terms. Padding a common array with zero
terms leaves the channel unchanged; omitting a branch or a term does not
establish completeness.

\phantomsection\label{p-qop-02}{}

\textbf{P-QOP-02: No signalling.} A local CPTP map \(\mathcal E_A\) does
not change the unconditioned state of \(B\).

\textbf{Proof.} For every effect \(F\) on \(B\), unitality of the dual
gives
\((\mathcal E_A^\dagger\otimes\operatorname{id})(I\otimes F)=I\otimes F\),
so all local probabilities remain unchanged. Selecting a measurement
outcome can change a conditional state, but a receiver without the
outcome label still observes the sum. \(\square\)

\section{4.4 Distances and errors}

The trace distance is \(D(\rho,\sigma)=\frac12\|\rho-\sigma\|_1\). For a
channel difference we use the diamond norm \[
\|\Phi\|_\diamond=\sup_{\|X\|_1=1}
\|(\operatorname{id}_d\otimes\Phi)(X)\|_1,
\] where auxiliary dimension equal to the finite input dimension \(d\)
suffices. The channel discrimination distance is half this norm.

\phantomsection\label{p-qop-03}{}

\textbf{P-QOP-03: Operational error and propagation.} Any
effect-probability difference is bounded by trace distance, and CPTP
maps contract state trace distance. If the channel diamond distance at
step \(t\) is at most \(\epsilon_t\), the same causal quantum controller
and initial state give a final distinguishable-probability error at most
\(\sum_t\epsilon_t\).

\textbf{Proof.} Write the positive--negative spectral decomposition
\(\Delta=\rho-\sigma=\Delta_+-\Delta_-\). Since
\(\operatorname{Tr}\Delta=0\), both parts have trace \(\|\Delta\|_1/2\).
For \(0\preceq E\preceq I\), \[
-\operatorname{Tr}\Delta_-\le\operatorname{Tr}(E\Delta)
\le\operatorname{Tr}\Delta_+.
\] Taking \(E\) to be the positive spectral-support projection attains
the upper bound, hence \[
\sup_{0\preceq E\preceq I}|\operatorname{Tr}[E(\rho-\sigma)]|
=\tfrac12\|\rho-\sigma\|_1.
\] For CPTP \(\mathcal T\), the dual is positive and unital, so
\(0\preceq\mathcal T^*(E)\preceq I\). Trace duality identifies an output
effect-probability difference with that of its pulled-back input effect.
Taking the supremum proves contraction.

For multiple steps, include the existing classical record, controller,
and permitted quantum memory in the circuit. Both circuits use the same
causal control rule. Replace their channels one step at a time to form
intermediate circuits. Each replacement has half-diamond-norm bound
\(\epsilon_t\); identical CPTP processing before and after it and the
final effect do not increase discrimination distance. The triangle
inequality gives the final probability bound
\(\min(1,\sum_t\epsilon_t)\). \(\square\)

Instrument comparisons must retain the outcome label in the channel:
\(\rho\mapsto\sum_o|o\rangle\langle o|\otimes\mathcal M_o(\rho)\).
Bounding only the total channel after forgetting labels does not
guarantee closeness of conditional-feedback circuits. The maximal
response difference over a finite protocol family also defines an
operational pseudodistance, but generally only a lower bound on the full
effect supremum. Tomographic completeness cannot be assumed for both
without justification.

Postselection amplifies error. If unnormalized posteriors
\(\sigma,\tau\succeq0\) satisfy \(\|\sigma-\tau\|_1\le\delta\), and
\(p=\operatorname{Tr}\sigma,q=\operatorname{Tr}\tau\ge p_{\min}>0\),
then \[
D(\sigma/p,\tau/q)\le\delta/p_{\min}.
\] Add and subtract \(\tau/p\) and use \(|p-q|\le\delta\) to obtain the
bound. Rare outcomes do not retain the unconditioned error radius.

Trace distance, the diamond norm, and Choi equivalence are presented
here as complete finite-dimensional mathematics. Their particular
coverage by the machine kernel must be checked declaration by
declaration against Lean sources and audit receipts. An exact numerical
check of one model does not automatically formalize all the spectral and
norm theory used here. {[}Q2, §4; S4, §§2, 9{]}

\chapter{5. Quantum FFIPS and realizations of quantum
objects}\label{ch-5}

Whether a density matrix contains all predictive information depends on
the allowed future experiments and on environmental memory. This chapter
uses that dependence to show how branchwise intertwining and finite
effect-space closure preserve every finite record, with a counterexample
to comparing total channels alone.

\section{5.1 States as representations of complete future predictions}

\phantomsection\label{p-qff-01}{}

\textbf{P-QFF-01: Finite quantum FFIPS.} Assume the current \(\rho_h\)
is sufficient for every future process, and future protocols contain a
family of effects separating Hermitian operators. Then \[
h\sim h'\iff\rho_h=\rho_{h'}.
\]

\textbf{Proof.} Equal states give equal probabilities under the same
future instrument sequence. Conversely, equality of all allowed future
probabilities makes every pairing of the separating effects with
\(\Delta=\rho_h-\rho_{h'}\) vanish, so \(\Delta=0\). If all effects are
allowed, the positive spectral-support projection \(P_+\) of a nonzero
traceless Hermitian \(\Delta\) directly distinguishes the states because
\(\operatorname{Tr}(P_+\Delta)>0\). A finite separating family need not
contain every such \(P_+\); its spanning property suffices. \(\square\)

If the actual protocols only measure \(Z\), the current responses of
\(|+\rangle,|-\rangle\) agree. Allowing a future Hadamard followed by
\(Z\) immediately distinguishes them. The canonical state is therefore
relative to the protocols, not to the analyst's preferred coordinates.

With environmental memory or initial system--environment correlations,
\(\rho_h\) need not be sufficient. Chapter 18 provides an explicit
causal-break counterexample. One must then use the process response
rather than insist that the reduced state represents every future.

\section{5.2 Quantum readout and general carriers}

Fix finite physical readout blocks \(V\) and a target response \(C\). An
access set \(S\subseteq V\) supplies implementable quantum or classical
interfaces; sufficiency requires recovery of all specified future
responses. Minimal sufficient families and access blockers retain their
Core definitions, but quantum interfaces must state whether joint
measurements, ancillas, and coherent access are permitted.

Logically deleting an output port, physically erasing a subsystem, and
measuring it while forgetting the outcome are different operations. They
may preserve, restrict, or destroy coherence in different ways; a damage
result without specified semantics cannot cover them indiscriminately.

\section{5.3 A quotient of quantum dynamics}

Let encoding \(\mathcal E:M\to X\) and decoding \(\mathcal D:X\to M\)
between logical and physical systems be CPTP, with
\(\mathcal D\mathcal E=\operatorname{id}\) on the declared logical state
class. For each permitted action \(a\) and outcome \(o\), let
\(\mathcal M_{a,o}\) and \(\overline{\mathcal M}_{a,o}\) be the physical
and logical instrument branches. Require branchwise commutation: \[
\mathcal M_{a,o}\mathcal E=\mathcal E\overline{\mathcal M}_{a,o}.
\] This must hold on the full logical operator space, or at least on a
closed class containing every reachable branch output. If a terminal
effect \(F\) is tested, its logical effect must be
\(\bar F=\mathcal E^*(F)\), or have the same pairing on the specified
input class.

\phantomsection\label{p-qff-02}{}

\textbf{P-QFF-02: Path preservation by a quantum quotient.} Under
branchwise commutation, a physical initial state
\(\mathcal E(\bar\rho)\) has, for every permitted finite record, an
unnormalized branch state exactly corresponding by encoding to the
logical one. Record probabilities agree, as do joint probabilities with
corresponding terminal effects.

\textbf{Proof.} For \(w=((a_1,o_1),\ldots,(a_t,o_t))\), successive
branchwise substitutions give \[
\mathcal M_{a_t,o_t}\cdots\mathcal M_{a_1,o_1}\mathcal E(\bar\rho)
=\mathcal E\bigl(\overline{\mathcal M}_{a_t,o_t}\cdots
\overline{\mathcal M}_{a_1,o_1}(\bar\rho)\bigr).
\] Trace preservation of the encoding makes the traces equal. A terminal
effect is handled by
\(\operatorname{Tr}[F\mathcal E(X)]=\operatorname{Tr}[\bar F X]\). An
adaptive protocol selects the same action on a shared observed record
path, so the branchwise equality remains valid. Summing over a finite
leaf set gives equal probabilities for every leaf event.
Zero-probability branches retain the zero unnormalized operator, without
division. \(\square\)

The requirement cannot be weakened to equal total channels and equal
single-step POVMs. On a qubit, let \[
\mathcal M_0(X)=\tfrac12\operatorname{Tr}(X)|0\rangle\langle0|,
\quad\mathcal M_1(X)=\tfrac12\operatorname{Tr}(X)|1\rangle\langle1|,
\] and \(\mathcal N_0=\mathcal M_1,\ \mathcal N_1=\mathcal M_0\). Both
instruments are CP, every branch effect is \(I/2\), and both total
channels are \(X\mapsto\operatorname{Tr}(X)I/2\). Yet a subsequent \(Z\)
readout always agrees with the first label for \(\mathcal M\) and always
disagrees for \(\mathcal N\). This exact counterexample shows that
one-step marginals and unconditioned evolution alone do not preserve
multitime predictions. {[}Q2, §5.3, revised premises; S2, §§7.5,
28.3--4{]}

A quotient valid on the encoding image does not cover arbitrary physical
initial states. Leakage, encoding error, and ancillary completeness need
separate treatment. Approximate commutation should use a restricted or
full diamond norm with its input class stated.

\section{5.4 Persistence and quantum existence}

Possible persistent structures include invariant subspaces, noiseless
subsystems, stationary states, metastable states, correctable logical
states, topologically protected sectors, and effective quasiparticles.
Each has its own mechanism conditions. A wavefunction need not exhibit a
classical closed loop to qualify as an object, and a stable spectral gap
does not automatically imply topological protection.

Chapter 20 adds control value to quantum objects under further task
specifications. An object certificate alone does not select a unique
reward function.

\section{5.5 From finite closure checks to every finite future}

A finite quantum system need not have its proof re-enumerated for every
future length. The key is closure of the effect space generated by the
actual instruments, not calling a finite-horizon test table the
``complete future.'' Fix a finite operation library. For event
\(e=(a,o)\), the Heisenberg pullback is \[
\mathcal M_e^*(E)=\sum_kK_{e,k}^\dagger E K_{e,k}.
\] For records ordered chronologically, define \[
E_{\varnothing}=I,\qquad E_{e::w}=\mathcal M_e^*(E_w).
\] Cyclicity gives
\(\operatorname{Tr}[E\mathcal M_e(X)]=\operatorname{Tr}[\mathcal M_e^*(E)X]\).
Induction on word length then gives \[
p(w\mid\rho)=\operatorname{Tr}(E_w\rho).
\] Define the real linear space
\(\mathcal O=\operatorname{span}_{\mathbb R}\{E_w\}\). Every \(E_w\) is
Hermitian, so a \(d\)-dimensional quantum system has
\(\dim\mathcal O\le d^2\). This counts real coordinates of Hermitian
matrices; it does not say that arbitrarily few measurements recover
every state.

A finite candidate list \(B_1,\ldots,B_r\) requires three kinds of
finite checks: some \(B_{j_0}=I\); every \(B_j\) is an actual word
effect \(E_{w_j}\); and, for every allowed event and candidate element,
real coefficients satisfy \[
\mathcal M_e^*(B_j)=\sum_iT_{e,ji}B_i.
\] The first two conditions prevent insertion of free operators
unrelated to the operation library. The third establishes one-step
pullback closure. Induction places every \(E_w\) in the candidate real
span: the empty word uses the identity, and the prefix step uses
linearity and the checked local equations. Thus the candidate span is
exactly \(\mathcal O\).

With unnormalized coordinates \(x_j(X)=\operatorname{Tr}(B_jX)\), \[
x_j(\mathcal M_e(X))=\sum_iT_{e,ji}x_i(X),\qquad
x(\mathcal M_e(X))=T_ex(X).
\] For \(w=(e_1,\ldots,e_t)\), event \(e_1\) occurs first, hence \[
p(w\mid\rho)=\bigl[T_{e_t}\cdots T_{e_1}x(\rho)\bigr]_{j_0}.
\] Two states have equal finite coordinates exactly when all finite-word
probabilities agree. The forward implication follows from closure and
recursion; the reverse follows because every \(B_j\) is itself a word
effect. The candidate list may be linearly redundant. Sufficiency
remains valid, but the certificate alone does not establish minimality
or independence of the coordinates.

The result also covers complete finite adaptive trees. Each leaf is an
actual record word with actions chosen from its history, and its
probability is computed from the same initial coordinates. Summing leaf
probabilities gives record-event probabilities. Conditional coordinates
at a positive-probability leaf are divided by its probability;
zero-probability leaves retain unnormalized semantics. Posterior closure
and all-word equivalence are conclusions of local pullback identities,
not assumptions.

The UEOT-QM files \texttt{Observables.\allowbreak lean},
\texttt{Finite\allowbreak Coordinates.\allowbreak lean}, and \texttt{Adaptive.\allowbreak lean} implement
this actual Kraus--trace-pairing chain. \texttt{Exact.\allowbreak lean} soundly
interprets pairs of matrices with rational real and imaginary parts as
complex matrices and proves preservation of addition, multiplication,
adjoint, real rational scaling, and finite sums.
\texttt{Exact\allowbreak Certificate.\allowbreak lean} additionally checks nonnegative weights,
completeness of every instrument, the empty word, an actual word witness
for each basis matrix, and local closure, then constructs ordinary
complex quantum instruments and a \texttt{Finite\allowbreak Observable\allowbreak Presentation}.
The finite data can be proved valid by \texttt{decide\ +kernel}; the
general soundness theorem then lifts the result to every finite record.
{[}Source interfaces: \texttt{UEOTQM.\allowbreak record\allowbreak Effect},
\texttt{observable\allowbreak Space},
\texttt{response\_\allowbreak eq\_\allowbreak iff\_\allowbreak pairing\_\allowbreak observable\allowbreak Space};
\texttt{Exact\allowbreak Certificate.\allowbreak Conditions},
\texttt{exact\allowbreak Record\allowbreak Effect\_\allowbreak to\allowbreak Matrix},
\texttt{coordinates\_\allowbreak eq\_\allowbreak iff\_\allowbreak response\_\allowbreak eq},
\texttt{adaptive\_\allowbreak leaf\allowbreak Probability\_\allowbreak eq}{]}

These guarantees concern the finite model bound to those sources. They
do not automatically verify external JSON parsers, search algorithms, or
all exporters; nor do they show that a real device realizes the input
numbers. They do not derive particle generations, a mass spectrum, or
full QFT from a finite quantum response space. Source binding and
negative controls answer which data were compiled. Experiments and
sector bridges answer how those data correspond to physics.

\chapter{6. Measurement: coherent parents, public quotients, and
conditional predictive states}\label{ch-6}

Measurement involves joint quantum evolution, publicly readable labels,
and conditional prediction once an outcome is known. Separating these
aspects shows how local records can be stable and mutually consistent
while the joint system retains coherence, setting up the dynamics and
inversion protocols of the next two chapters.

\section{6.1 Coherent completion of any finite instrument}

Let \(\mathcal M_a(\rho)=\sum_\mu A_{a\mu}\rho A_{a\mu}^\dagger\), with
trace-preserving sum. Define \[
V|\psi\rangle=\sum_{a,\mu}A_{a\mu}|\psi\rangle\otimes|a,\mu\rangle,\qquad
Q_a=\sum_\mu|a,\mu\rangle\langle a,\mu|.
\]

\phantomsection\label{p-qmeas-01}{}

\textbf{P-QMEAS-01: Parent-system completion.} \(V^\dagger V=I\), and \[
\mathcal M_a(\rho)=
\operatorname{Tr}_R[(I\otimes Q_a)V\rho V^\dagger(I\otimes Q_a)].
\] After a finite enlargement, \(V\) is realizable by a unitary
operation and a fixed initial ancilla.

\textbf{Proof.} Orthogonality of the auxiliary basis gives
\(V^\dagger V=\sum A_{a\mu}^\dagger A_{a\mu}=I\). The projection selects
label \(a\), and the partial trace retains only equal-\(\mu\) terms,
giving the instrument. Complete the image basis of the isometry to an
orthonormal basis of the full space to obtain a unitary extension.
\(\square\)

A pure input gives a pure joint parent state. Whether this parent can
actually be preserved reversibly depends on environmental access and
noise. A mathematical dilation alone does not establish that the inverse
can be implemented experimentally.

\section{6.2 Conditioning as the predictive state of an outcome-extended
history}

\phantomsection\label{p-qmeas-02}{}

\textbf{P-QMEAS-02.} After outcome \(a\) occurs, with \(p_a>0\), every
future quantum effect \(F\) has conditional probability \[
P(F|h,a)=\operatorname{Tr}(F\rho_a).
\] If future protocols separate states, \(\rho_a\) is the unique quantum
FFIPS for the outcome-extended history.

\textbf{Proof.} Divide the joint probability
\(\operatorname{Tr}[F\mathcal M_a(\rho)]\) by \(p_a\). Uniqueness
follows from P-QFF-01. \(\square\)

The standard measurement chain is thus organized as coherent parent
completion, access to a record algebra, and conditioning on an outcome.
This agrees with standard instrument predictions in QM. The UEOT
connection organizes those steps; it does not add collapse dynamics.

\section{6.3 What public records cannot see}

Let \(\Delta_R(\Omega)=\sum_a(I\otimes Q_a)\Omega(I\otimes Q_a)\).

\phantomsection\label{p-qmeas-03}{}

\textbf{P-QMEAS-03: Public blindness.} If \([B,I\otimes Q_a]=0\) for
every \(a\), then \[
\operatorname{Tr}(B\Omega)=\operatorname{Tr}[B\Delta_R(\Omega)].
\]

\textbf{Proof.} The observable \(B\) is block diagonal in the records.
Expanding the projection decomposition and using cyclicity makes every
cross term with distinct labels vanish. \(\square\)

The coherent parent and its dephased state therefore have the same
public-record statistics. Redundant records, stable pointers, and
observer agreement alone cannot decide whether global coherence has
fundamentally disappeared. Distinguishing the states requires coherent
tests across record blocks, such as complete erasure or inversion.

This blindness holds for all record-block-diagonal effects. The
block-diagonal operator algebra \(\bigoplus_a\mathcal B(\mathcal H_a)\)
can remain noncommutative when a branch has internal dimension greater
than one; the whole algebra is not classical. The next section uses only
\(\mathcal Z=\operatorname{span}\{P_a\}\), generated by label
projections. This smaller algebra is automatically commutative.
Classical public labels can coexist with quantum information inside
branches.

\section{6.4 Classical dynamics on a commutative algebra}

Public projections \(P_a\) span a commutative algebra \(\mathcal Z\). If
\(\mathcal T^\dagger(\mathcal Z)\subseteq\mathcal Z\), then \[
\mathcal T^\dagger(P_b)=\sum_a M_{ab}P_a,\quad
M_{ab}\ge0,\quad\sum_bM_{ab}=1.
\] The row vector of public probabilities consequently evolves as
\(p'=pM\).

\phantomsection\label{p-qmeas-04}{}

\textbf{P-QMEAS-04.} The matrix \(M\) is a classical stochastic matrix.

\textbf{Proof.} Positivity gives nonnegative expansion coefficients,
unitality gives row sums equal to one, and trace pairing gives the
probability update. \(\square\)

This directly connects quantum public classicality to the classical
predictive quotient of Core. The entire parent system need not become a
classical probabilistic mixture.

A multitime classical record process also requires an appropriate
measurement implementation or closure of the record branches on the
algebra. Invariance under one unconditioned \(\mathcal T^*\) directly
guarantees only Markov updates of public probabilities under its
repeated action. Other intermediate instruments or feedback require the
branchwise checks of §5.3.

\section{6.5 Objective record certificates (ORCs)}

Specify the labels, error \(\delta\), persistence window \(T\), number
of records \(m\), and confidence level in advance. A certificate
requires at least a commutative public algebra; \(m\) disjoint fragments
that can each identify the label; persistence and transport of the
records along histories; a positive margin to a defined failure class;
and remaining confidence margin in every statistical condition.

An ORC is operationally defined. A mutual-information plateau in
environmental fragments is relevant evidence, but does not replace
independent accessibility, nondisturbing readout, or a complete failure
model. Nor does it determine the strength of an additional actualization
channel considered below.

\textbf{Public records and total correlations.} Quantum mutual
information includes both classical and quantum correlations. A plateau
for an environment fragment does not by itself establish that an
observer can read the required pointer label without disturbing the
stipulated record. Strong Quantum Darwinism and spectrum broadcast
structure impose further conditions on accessible information,
distinguishable records, and conditional fragment structure. In
particular, pairwise conditional independence does not ensure full
factorization of the conditional joint state of all fragments; the 2021
Comment and Reply clarified this point in the 2019 equivalence
statement.
\href{https://doi.org/10.1103/PhysRevA.98.032103}{Le--Olaya-Castro
2018}, \href{https://doi.org/10.1103/PhysRevLett.122.010403}{2019},
\href{https://doi.org/10.1103/PhysRevLett.126.188901}{Feller et al.,
Comment}, \href{https://doi.org/10.1103/PhysRevLett.126.188902}{authors'
Reply}. The public quotient, strong record nondisturbance, and empirical
record criteria in this book are related but distinct conditions;
terminology alone does not make them interchangeable.

\chapter{7. Complete classification of strong record-nondisturbing
dynamics}\label{ch-7}

If a channel leaves every quantum state within each record branch
untouched, how can it still change coherence between branches? A
correlation matrix captures this freedom completely, and the semigroup
analysis derives its global compatibility conditions: individually
plausible pairwise behaviour need not assemble into a physical channel.

Later connections: Chapter 38 gives cycle-phase bounds, Chapter 39 a
four-level noise-class witness, and Chapter 46 tests for correlated
trials.

\section{7.1 SRN and the Schur form}

Let finite-dimensional \(\mathcal H=\bigoplus_{a=1}^m\mathcal H_a\) have
nonzero branches and \(\sum_aP_a=I\). A CPTP channel satisfies
\textbf{SRN} if it fixes every density matrix within each branch
pointwise, not merely each branch's total probability. Remove
zero-dimensional branches first to avoid assigning spurious independent
parameters to unobservable blocks.

\phantomsection\label{p-qsrn-01}{}

\textbf{P-QSRN-01: Complete SRN classification.} SRN holds exactly for
channels \[
\mathcal A_K(X)=\sum_{a,b}K_{ab}P_aXP_b,\qquad
K\succeq0,\quad K_{aa}=1.
\]

\textbf{Proof.} Choose Kraus operators \(L_\mu\). Every pure-state
projector in a branch is fixed. The positive terms
\(L_\mu|\psi\rangle\langle\psi|L_\mu^\dagger\) sum to a rank-one
operator, so every \(L_\mu|\psi\rangle\) is parallel to
\(|\psi\rangle\). Applying linearity to two independent vectors in the
same branch and their sum forces the same proportionality constant:
\(L_\mu P_a=c_{\mu a}P_a\). Hence
\(K_{ab}=\sum_\mu c_{\mu a}\overline{c_{\mu b}}\) is a Gram matrix, with
unit diagonal by trace preservation. Conversely, any such Gram
decomposition gives \(L_\mu=\sum_ac_{\mu a}P_a\), a CPTP channel fixing
every branch state pointwise. \(\square\)

This operator structure agrees with existing theory of Schur channels,
nondisturbed algebras, and genuinely incoherent operations. We re-prove
it to establish a complete parameterization for the measurement
extension; it is not claimed as an original UEOT operator theorem.

\section{7.2 Composition, blindness, and underdetermination}

\phantomsection\label{p-qsrn-02}{}

\textbf{P-QSRN-02.} \[
\mathcal A_K\mathcal A_L=\mathcal A_{K\odot L}.
\] Every block-diagonal effect has the same statistics under
\(\mathcal A_K\) as under the identity channel.

\textbf{Proof.} Each \((a,b)\) block is multiplied successively by
\(L_{ab}\) and \(K_{ab}\). Public effects read only the \(a=b\) blocks,
whose coefficients are one. \(\square\)

The all-ones matrix \(J\) represents the identity channel, while the
identity matrix \(I\) represents complete branch dephasing. They have
the same public records. An ORC therefore does not determine a unique
\(K=F(\mathrm{ORC})\). If standard fully coherent dynamics have already
been assumed, the corresponding channel has already been selected; it
has not then been independently derived from objectivity.

\section{7.3 Global generators of continuous semigroups}

Continuity and \(K(t+s)=K(t)\odot K(s),K(0)=J\) imply entrywise
\(K_{ab}(t)=e^{t\lambda_{ab}}\). A zero at finite time would, by
repeated halving, contradict continuity at zero. A nonzero continuous
multiplicative function admits a local logarithm and hence the
exponential form.

\phantomsection\label{p-qsrn-03}{}

\textbf{P-QSRN-03: Necessary and sufficient generator-cone conditions.}
These maps are CPTP for every \(t\ge0\) exactly when \[
\Lambda=\Lambda^\dagger,\quad\lambda_{aa}=0,\quad
z^\dagger\Lambda z\ge0\quad(\mathbf1^\dagger z=0).
\]

\textbf{Proof.} Necessity follows by writing \(K(t)=J+t\Lambda+o(t)\),
restricting to the zero-sum subspace to eliminate \(J\), and taking the
right derivative of positivity. Sufficiency has a direct construction,
without invoking Schoenberg's theorem as a black box. Fix an anchor
\(r\) and define \[
G_{ab}=\lambda_{ab}-\lambda_{ar}-\lambda_{rb}.
\] For any \(x\), \(x^\dagger Gx\) equals \(z^\dagger\Lambda z\) with
\(z=x-(\sum_ax_a)e_r\), hence \(G\succeq0\). Put
\(\theta_a=\operatorname{Im}\lambda_{ar}\). Then \[
\lambda_{ab}=G_{ab}-\tfrac12G_{aa}-\tfrac12G_{bb}
+i(\theta_a-\theta_b).
\] Schur products preserve positivity, so
\(\exp_\odot(tG)=J+tG+t^2G\odot G/2+\cdots\succeq0\). Multiplying on the
left by the diagonal matrix \(D_{aa}=e^{-tG_{aa}/2+it\theta_a}\) and on
the right by its adjoint gives exactly \(K(t)\). Its diagonal is one;
P-QSRN-01 gives CPTP, and the semigroup law holds entrywise. \(\square\)

This is an explicit global compatibility proof. Nonnegative pairwise
decay rates are insufficient. For example, \[
\Lambda=\begin{pmatrix}0&-1&-1\\-1&0&-10\\-1&-10&0\end{pmatrix}
\] gives \(z^\top\Lambda z=-12\) for \(z=(2,-1,-1)\), so its small-time
map is not physical.

In the real symmetric case, put \(D_{ab}=-\lambda_{ab}\). The Gram
construction gives \(D_{ab}=\|v_a-v_b\|^2\). General complex generators
also contain the complex geometry of the Gram matrix; arbitrary phase
differences alone do not capture it.

Explicitly, if \(G_{ab}=\langle u_a,u_b\rangle\) is a real Gram
representation, then \(D_{ab}=\|u_a-u_b\|^2/2\). Setting
\(v_a=u_a/\sqrt2\) gives the stated distance form. In the complex case,
imaginary cycle phases are constrained by the same positive Gram
structure. Arbitrary phases added to nonnegative pairwise decay rates
need not generate a CP semigroup. The classification assumes a fixed
branch decomposition, finite dimension, continuity, and a
time-homogeneous semigroup. Changing records or environmental memory
require a different model.

\chapter{8. QDRR: record writing, erasure by inversion, and finite
data}\label{ch-8}

QDRR writes a record and then reverses the writing to probe coherence
invisible to public labels. We derive what it reconstructs within SRN,
quantify finite-sample errors, and state the hardware-control
requirements; residual loss must first be distinguished from incomplete
reversal and ordinary noise.

\section{8.1 The interface for complete inversion}

Distinguish two types. The full unitary \(U_W\) acts on the system and
record register, initialized at \(|0\rangle_R\). The actual writing
isometry is \[
W:\mathcal H\to\mathcal H\otimes\mathcal R,\qquad
W|\psi\rangle=U_W(|\psi\rangle\otimes|0\rangle_R).
\] Its image branch projections are \(\widetilde P_a=WP_aW^\dagger\).
They sum to \(WW^\dagger\), not automatically to the identity on the
enlarged space. The intermediate channel is SRN with kernel \(K\) on
these orthogonal image branches and does not leak input outside the
image. Defining it on the whole space requires an explicit action on the
orthogonal complement.

\phantomsection\label{p-qqdr-01}{}

\textbf{P-QQDR-01: Coherent pullback.} \[
W^\dagger\widetilde{\mathcal A}_K(WXW^\dagger)W=\mathcal A_K(X).
\]

The corresponding full physical inverse satisfies \[
U_W^\dagger\widetilde{\mathcal A}_K
\bigl[U_W(X\otimes|0\rangle\langle0|)U_W^\dagger\bigr]U_W
=\mathcal A_K(X)\otimes|0\rangle\langle0|.
\]

\textbf{Proof.} Expand \(X=\sum_{a,b}P_aXP_b\). Since \(W^\dagger W=I\),
\[
\widetilde P_a(WXW^\dagger)\widetilde P_b=WP_aXP_bW^\dagger.
\] The intermediate SRN map multiplies this block by \(K_{ab}\),
yielding \(W\mathcal A_K(X)W^\dagger\). Multiplication by
\(W^\dagger,W\) gives the first identity. Using
\(U_W^\dagger W|\psi\rangle=|\psi\rangle\otimes|0\rangle\) gives the
second. \(\square\)

Inversion must include every relevant environment that actually retains
which-branch information. Residual decoherence after reversing only some
registers cannot be attributed directly to fundamental collapse.

\section{8.2 Tomography of the full SRN class}

Prepare \((|a\rangle+e^{i\phi}|b\rangle)/\sqrt2\) and measure
\(X_{ab}=|a\rangle\langle b|+|b\rangle\langle a|\).

\phantomsection\label{p-qqdr-02}{}

\textbf{P-QQDR-02.} \[
\langle X_{ab}\rangle_\phi
=\operatorname{Re}(e^{-i\phi}K_{ab}).
\] Phases \(\phi=0,\pi/2\) give the real and imaginary parts. The
\(m(m-1)\) real coordinates for all \(a<b\) recover the entire SRN
channel.

\textbf{Proof.} The input \((a,b)\) entry is \(e^{-i\phi}/2\) and is
multiplied by \(K_{ab}\). Adding its conjugate gives the expression. The
diagonal is already known to be one, so reconstruction is unique.
\(\square\)

Completeness here concerns a single finite-dimensional CPTP-SRN map. It
does not cover population transfer, leakage, disturbance within
branches, or arbitrary history-dependent new physics.

\section{8.3 Finite sampling and projection onto physical matrices}

Let \(m\ge2\). At each setting obtain \(n\ge1\) independent samples in
\([-1,1]\) with that setting's fixed mean. Set \(L=m(m-1)\) and
\(0<\alpha<1\). The union bound does not require independence between
settings, but the sampling assumptions for Hoeffding within each setting
must hold. When \(m=1\) there are no off-diagonal parameters; \(L=0\)
must not be inserted into the logarithmic formula below.

\phantomsection\label{p-qqdr-03}{}

\textbf{P-QQDR-03: Simultaneous error.} With probability at least
\(1-\alpha\), \[
\max_\ell|\widehat\mu_\ell-\mu_\ell|
\le\epsilon_n=\sqrt{\frac{2\log(2L/\alpha)}n},
\] and \[
\|\widehat K-K\|_{\rm op}\le\|\widehat K-K\|_F
\le\epsilon_n\sqrt{2m(m-1)}.
\]

\textbf{Proof.} Hoeffding gives \(2e^{-n\epsilon^2/2}\) for a variable
of range width two; take the union bound over \(L\) settings. Each
complex off-diagonal squared error is at most \(2\epsilon^2\), and
Hermitian symmetry counts each pair twice. \(\square\)

For \(m=3,\epsilon=.12,\alpha=.05\), each setting requires \(n\ge762\).
The coefficient for \([0,1]\) variables cannot be used without rescaling
the measurement outputs.

An empirical matrix need not be positive semidefinite. Projection onto
the nearest correlation matrix can improve physical usability, but the
projection residual must remain as a model-consistency diagnostic; it
cannot be discarded while reporting only a postprojection pass.
Euclidean projection onto a closed convex set containing the true \(K\)
does not increase its Frobenius error. It does not turn an incorrect SRN
model into an empirically validated model.

\section{8.4 Identifiability of hardware effects and new physics}

Only when hardware noise and candidate actualization are independent SRN
processes does \(K_{\rm obs}=K_{\rm hw}\odot K_{\rm act}\) hold.
Entrywise division is unstable near zeros of \(K_{\rm hw}\), and general
noise need not admit this factorization.

An experiment requires, at minimum, matched no-record and
immediate-inversion controls, scrambled redundancy, physical-qubit
permutations, randomized interleaving, and independent calibration.
Existing quantum-Darwinism data generally test public records. Without
coherent inversion and matched calibration, record redundancy cannot
retrospectively be interpreted as evidence for fundamental collapse. No
new hardware experiment was performed for this text.

\section{8.5 What existing experiments establish}

Quantum Darwinism has been investigated experimentally in controlled
settings. Each result should be cited with its actual model and readout;
these studies do not collectively constitute experimental confirmation
of UEOT.

\begin{longtable}[]{@{}
  >{\raggedright\arraybackslash}p{(\columnwidth - 4\tabcolsep) * \real{0.2500}}
  >{\raggedright\arraybackslash}p{(\columnwidth - 4\tabcolsep) * \real{0.3400}}
  >{\raggedright\arraybackslash}p{(\columnwidth - 4\tabcolsep) * \real{0.4100}}@{}}
\toprule\noalign{}
\begin{minipage}[b]{\linewidth}\raggedright
Original experiment
\end{minipage} & \begin{minipage}[b]{\linewidth}\raggedright
System and readout
\end{minipage} & \begin{minipage}[b]{\linewidth}\raggedright
Evidence relevant here
\end{minipage} \\
\midrule\noalign{}
\endhead
\bottomrule\noalign{}
\endlastfoot
Ciampini et al., 2018,
\href{https://doi.org/10.1103/PhysRevA.98.020101}{PRA 98, 020101(R)} &
Information distribution in photonic cluster states & Redundant-record
signatures in a specified model \\
Chen et al., 2019,
\href{https://doi.org/10.1016/j.scib.2019.03.032}{Science Bulletin 64,
580--585} & Six-photon quantum simulator & Classical records in
controlled environment fragments \\
Unden et al., 2019,
\href{https://doi.org/10.1103/PhysRevLett.123.140402}{PRL 123, 140402} &
NV centre and nuclear-spin environment & Redundant information in a
specified system--environment setting \\
Zhu et al., 2025, \href{https://doi.org/10.1126/sciadv.adx6857}{Science
Advances 11, eadx6857} & Branching structure and correlations in
superconducting circuits & Record formation under standard quantum
operations \\
Zheng et al., 2026, \href{https://doi.org/10.1103/t2yx-9bqh}{PRL 137,
150201} & Coherence revival of a mesoscopic cat state in a controlled
reservoir & Suppression of spontaneous revival with increasing reservoir
size over the measured time window \\
\end{longtable}

These are laboratory studies, including quantum simulations on physical
devices. They are distinct from the synthetic computer calculations in
Appendix D. Neither record formation nor revival suppression alone
establishes the elimination of global coherence by fundamental
nonunitary dynamics. The fully controlled inversion in §8.1 is also a
different operation from waiting for a spontaneous revival.

Data-processing conventions must accompany the results. For example, Zhu
et al.~reconstruct logical states using postselection that excludes
outcomes outside the logical subspace; this cannot directly certify an
unconditional channel. Experimental points, numerical curves, and the
definitions of error bars must also remain distinct. This book has not
obtained the studies' raw counts for a common reanalysis, inferred a new
actualization-strength bound from them, or performed a QDRR hardware
experiment.

\chapter{9. Quantum information, error correction, and
metrology}\label{ch-9}

Copyable public records and the impossibility of cloning arbitrary
unknown quantum states are compatible because they preserve different
information. Entanglement, correction, and metrology are placed in their
respective operational tasks, using standard quantum-information results
to assess maintenance and estimation without making one task's optimum a
universal purpose.

\section{9.1 Entanglement and information}

The von Neumann entropy is \(S(\rho)=-\operatorname{Tr}\rho\log\rho\).
Quantum relative entropy is \[
D(\rho\|\sigma)=\operatorname{Tr}\rho(\log\rho-\log\sigma)
\] with value infinity when the required support inclusion fails.
\textbf{K-QINFO-01:} relative-entropy data processing under CPTP maps,
strong subadditivity, and the Holevo bound on accessible information are
invoked as standard quantum-information theorems.

Schmidt coefficients \(p_i\) of a pure parent state give
\(S(\rho_A)=S(\rho_B)\). Entanglement can produce local mixedness, but
local mixedness alone does not show that the actual purifying
environment has been experimentally identified.

Quantum mutual information \(I(A:B)=S(A)+S(B)-S(AB)\) measures total
correlations. Entanglement, shared classical records, accessible mutual
information, and irreducibility under a parent-object partition answer
different questions.

\section{9.2 No cloning}

\phantomsection\label{p-qinfo-01}{}

\textbf{P-QINFO-01.} No single unitary implements
\(|\psi\rangle|0\rangle\mapsto|\psi\rangle|\psi\rangle\) for every
unknown pure state.

\textbf{Proof.} For two states with inner product \(c\), unitary
preservation of inner products would require \(c=c^2\), impossible for
\(0<|c|<1\). Allowing an input-dependent final environment requires
pure-state cloning to have joint output
\(|\psi\rangle|\psi\rangle|e_\psi\rangle\). Inner-product preservation
becomes \(c=c^2\langle e_\psi|e_\varphi\rangle\); absolute values and
Cauchy--Schwarz give \(|c|\le|c|^2\), again impossible for \(0<|c|<1\).
An isometric dilation reduces general deterministic CPTP cloning to the
same argument. \(\square\)

Orthogonal classical records can be copied redundantly without violating
no cloning. An ORC copies distinguishable public labels, not arbitrary
unknown quantum states.

\section{9.3 Error-correction conditions and object maintenance}

\textbf{K-QEC-01: Knill--Laflamme.} On a code projection \(P\), noise
Kraus operators \(E_a\) admit a common perfect recovery exactly when \[
PE_a^\dagger E_bP=c_{ab}P.
\] This requires coherence with a full auxiliary system, not merely
separate recoverability of two basis states.

For sufficiency, diagonalize the positive matrix \(c\) to obtain
mutually orthogonal error images of the code. Measure the error
syndrome, then reverse each isometric error without reading logical
information. Necessity follows because the complementary environment of
a recoverable channel cannot distinguish logical states. The full
general proof is invoked as a standard K result.

In UEOT terms, the code subspace is a candidate realization, recovery is
a maintenance protocol, and logical futures are the prediction target.
Code distance and thresholds depend on noise, geometry, and
implementable gates; general Ω does not determine their numerical
values.

\section{9.4 Quantum Fisher information and metrology}

For a parameterized state \(\rho_\theta\), the symmetric logarithmic
derivative satisfies \(\partial_\theta\rho=(L\rho+\rho L)/2\), and
quantum Fisher information is \(F_Q=\operatorname{Tr}\rho L^2\). The
classical Fisher information of a given POVM does not exceed \(F_Q\).
Locally unbiased estimators obey the corresponding Cramér--Rao bound,
subject to support and regularity conditions.

For a pure unitary family
\(|\psi_\theta\rangle=e^{-i\theta G}|\psi\rangle\):

\phantomsection\label{p-qinfo-02}{}

\textbf{P-QINFO-02.} \[
F_Q=4(\langle G^2\rangle-\langle G\rangle^2).
\]

\textbf{Proof.} Substitute \(|\dot\psi\rangle=-iG|\psi\rangle\) into the
pure-state tangent-projection formula
\(F_Q=4(\langle\dot\psi|\dot\psi\rangle-
|\langle\psi|\dot\psi\rangle|^2)\). This formula can also be derived
directly by solving for the SLD on the state support and its orthogonal
tangent subspace. \(\square\)

Metrological optimality is not GOD for every task. Optimal measurements
for several parameters may be incompatible, so individual
single-parameter Fisher bounds cannot automatically be attained
simultaneously.

\chapter{10. Infinite-dimensional quantum theory: domains and normal
states}\label{ch-10}

When finite matrices give way to continuous position or momentum,
familiar formulas may remain while operator domains and state types
become part of the proof. Normal states, self-adjoint Hamiltonians, and
energy truncation provide legitimate calculation interfaces, with finite
energy alone still insufficient to guarantee a finite-dimensional model.

\section{10.1 Infinite Hilbert dimension and infinitely many degrees of
freedom}

A single harmonic oscillator already has infinite-dimensional
\(\mathcal H=L^2(\mathbb R)\), but only one canonical pair. A quantum
field has infinitely many modes, so the finite-degree-of-freedom
Stone--von Neumann uniqueness theorem cannot be applied directly.
Finite-dimensional operational reconstruction also does not uniquely
specify an infinite-dimensional extension.

We use a separable complex Hilbert space. \textbf{H-QINF} selects the
normal operational sector: normal states on \(\mathcal B(\mathcal H)\)
are represented by positive trace-class density operators of trace one.
Schrödinger-picture operations act on trace-class operators, and the
dual picture uses the corresponding normal CP maps. General operator
algebras may also have singular states; these are not silently
identified with density matrices. This infinite-dimensional structure is
an additional domain, not something supplied directly by
finite-dimensional reconstruction or finite \texttt{Fin\ d} matrix
certificates.

Bounded effects retain the Born pairing. A POVM is a weakly countably
additive positive-operator-valued measure. Continuous position
probabilities are \(\int_B|\psi(x)|^2dx\). The generalized \(|x\rangle\)
is not a unit vector in the Hilbert space, and a zero-probability point
cannot be normalized as a discrete posterior.

\section{10.2 Self-adjointness is part of the physical specification}

An unbounded operator is specified by both its action and its dense
domain. Symmetry, \(A\subseteq A^\dagger\), is not self-adjointness,
\(A=A^\dagger\).

\textbf{K-QINF-01: Spectral theorem.} A self-adjoint \(A\) has a unique
PVM \(E_A\) with \[
A=\int\lambda\,dE_A(\lambda),\quad
D(A)=\{\psi:\int\lambda^2d\langle\psi,E_A(\lambda)\psi\rangle<\infty\}.
\] Measurement probabilities are defined through the PVM. Means and
variances additionally require finite corresponding spectral moments.

On \(L^2(\mathbb R)\), position has domain \(x\psi\in L^2\), and
momentum is \(-i\hbar\partial_x\) on the Sobolev domain
\(H^1(\mathbb R)\). On an interval, the same derivative expression with
\(\psi(L)=e^{i\vartheta}\psi(0)\) gives different self-adjoint
extensions and momentum spectra. The boundary phase belongs to the
model, not to the symbol \(-i\partial_x\) alone.

\textbf{K-QINF-02.} A closed symmetric operator has self-adjoint
extensions when its deficiency indices
\(n_\pm=\dim\ker(A^\dagger\mp i)\) agree. A semibounded operator admits
the Friedrichs extension. Kato--Rellich requires a perturbation
relatively bounded with respect to an already self-adjoint \(A\), with
relative bound less than one. A closed semibounded quadratic form can
also generate a self-adjoint Hamiltonian through the representation
theorem. These are standard routes to a legitimate Hamiltonian.

\section{10.3 Stone's theorem and the strong Schrödinger equation}

\textbf{K-QINF-03: Stone.} A strongly continuous one-parameter unitary
group is uniquely \(U(t)=e^{-itA}\) with self-adjoint \(A\). Defining
energy by \(H=\hbar A\), for \(\psi_0\in D(H)\) one has \[
i\hbar\partial_t\psi_t=H\psi_t.
\] An initial state outside the domain still has a continuous unitary
orbit, but need not have a strong derivative. Stone's theorem supplies
the time generator, not the numerical value of \(\hbar\).

\phantomsection\label{p-qinf-01}{}

\textbf{P-QINF-01: Trace-norm continuity of normal states.} For strongly
continuous \(U(t)\) and any trace-class \(\rho\),
\(U(t)\rho U(t)^\dagger\) is trace-norm continuous.

\textbf{Proof.} Start with a finite-rank expansion and use
\(\||x\rangle\langle y|\|_1=\|x\|\|y\|\). Approximate any trace-class
operator by finite-rank operators, use unitary invariance of the trace
norm, and apply the triangle inequality. \(\square\)

The von Neumann equation \(\dot\rho=-i[H,\rho]/\hbar\) is a strong
equation in trace norm only on the domain of the trace-class
conjugation-group generator. For an arbitrary trace-class state, two
potentially undefined products cannot be treated as an ordinary matrix
formula.

\section{10.4 Infinite-dimensional FFIPS and energy truncation}

For a countable orthonormal basis, expectations of projections onto
\(|i\rangle\), \((|i\rangle+|j\rangle)/\sqrt2\), and
\((|i\rangle+i|j\rangle)/\sqrt2\) determine every density-matrix entry.
Under separating protocols, a normal state therefore still uniquely
represents a sufficient future. Countable identifiability does not imply
efficient learning from finite samples.

\phantomsection\label{p-qinf-02}{}

\textbf{P-QINF-02: Energy-truncation error.} Let \(H\ge0\), let a
truncation projection \(\Pi\) commute with \(H\), and assume
\(H\ge E_{\rm cut}(I-\Pi)\) outside the cutoff. If
\(\operatorname{Tr}\rho H\le E\), then \[
\epsilon=\operatorname{Tr}[(I-\Pi)\rho]\le E/E_{\rm cut}.
\] When \(p=1-\epsilon>0\), the normalized truncation
\(\rho_\Pi=\Pi\rho\Pi/p\) satisfies \[
D(\rho,\rho_\Pi)\le\sqrt{\epsilon}.
\]

\textbf{Proof.} The first inequality follows by integrating the
operator-order inequality. Purify \(\rho\) as \(|\Psi\rangle\), project
its system factor, and normalize. The purification overlap is
\(\sqrt p\), so their pure-state trace distance is \(\sqrt{1-p}\).
Contraction under partial trace gives the result. \(\square\)

A finite-dimensional computable model additionally requires finite rank
of \(\Pi\). An energy bound alone does not guarantee finitely many
states below the cutoff. Multistep truncation must control the energy
tail and dynamical defect at each step.

\chapter{11. Phase space, Weyl--CCR structure, and the action
scale}\label{ch-11}

This chapter asks which physical inputs produce position--momentum
noncommutativity. The argument proceeds from projective phase-space
displacements to Weyl relations and then CCR on a common domain, keeping
irreducibility, regularity, and calibration of the action scale as
separate responsibilities.

\section{11.1 Starting from a projective action}

On classical phase space \(V=\mathbb R^{2n}\), define
\(\sigma((q,p),(q',p'))=q\cdot p'-p\cdot q'\). A projective unitary
action obeys \[
W(z)W(z')=m(z,z')W(z+z').
\] Associativity requires the 2-cocycle equation. Rephasing
\(W'(z)=f(z)W(z)\) changes \(m\) by a coboundary without changing the
physical ray action.

\textbf{K-QCCR-01.} The relevant continuous multiplier classes of a
finite-dimensional real vector group are represented by real alternating
bilinear forms. Full linear symplectic covariance further restricts that
form to \(\kappa\sigma\). These are representation-theoretic and
kinematic conditions, not consequences of quantizing an action integral.

To fix coefficients, after rephasing choose
\(m(z,z')=\exp[iB(z,z')/2]\). The exponent containing \(B\) must be
dimensionless. If \(q,p\) have physical position and momentum units,
\(\kappa\) has inverse-action units. A different sign or factor-of-two
convention for the multiplier requires a corresponding change in the
definition of \(\kappa\) below.

\phantomsection\label{p-qccr-01}{}

\textbf{P-QCCR-01: A single scale from symplectic covariance.} If every
symplectic transformation preserves the alternating form \(B\), then
\(B=\kappa\sigma\).

\textbf{Proof.} Write \(B(u,v)=\sigma(Au,v)\). Covariance makes \(A\)
commute with all symplectic transformations. In Darboux blocks,
commutation with \(\operatorname{diag}(\lambda I,\lambda^{-1}I)\)
eliminates off-diagonal blocks; commutation with
\(\operatorname{diag}(R,R^{-T})\) makes each diagonal block scalar; and
commutation with symplectic shears
\(\left(\begin{smallmatrix}I&C\\0&I\end{smallmatrix}\right)\),
\(C=C^\top\), makes the two scalars equal. \(\square\)

The possibility \(\kappa=0\) remains. Projectivity and symplectic
covariance alone therefore do not force quantum noncommutativity.

\section{11.2 Conditional nontriviality under UCA}

\textbf{H-UCA.} There is no nontrivial invariant orthogonal projection
for all canonical displacements, and the canonical factor has dimension
greater than one. This is the projection formulation of irreducibility,
imposed only on the orbital canonical factor; it does not exclude
additional internal multiplicities such as spin.

\phantomsection\label{p-qccr-02}{}

\textbf{P-QCCR-02.} In a complex Hilbert representation, UCA and
dimension greater than one exclude the trivial multiplier class.

\textbf{Proof.} A trivial class can be rephased to an ordinary Abelian
unitary representation, whose unitaries commute. If all are scalar,
every nontrivial projection is invariant. If one is nonscalar, a
nontrivial spectral projection of that unitary commutes with the whole
family. Either case violates UCA. \(\square\)

Together with symplectic covariance this gives \(\kappa\ne0\). With this
chapter's multiplier convention, choose the orientation so that
\(\kappa=-1/\hbar\), where \(\hbar=1/|\kappa|>0\). A direct
nonzero-cocycle assumption has been replaced by a specific
irreducible-sector assumption. General objecthood has not
unconditionally established UCA or the action scale.

\section{11.3 Weyl relations and rigorous CCR}

We use \[
W(q,p)=\exp[i(p\cdot Q-q\cdot P)/\hbar],\qquad
W(z)W(z')=e^{-i\sigma(z,z')/(2\hbar)}W(z+z').
\] In the Schrödinger representation, these bounded unitaries can be
defined directly by \[
[W(q,p)\psi](x)=e^{ip\cdot(x-q/2)/\hbar}\psi(x-q).
\] Successive application verifies the Weyl phase.

\textbf{K-QCCR-02: Stone--von Neumann.} For finitely many canonical
degrees of freedom, a nondegenerate symplectic form, a fixed nonzero
central character, and a regular irreducible Weyl representation, the
representation is unitarily equivalent to the Schrödinger
representation. Dropping irreducibility introduces a multiplicity
factor. Dropping regularity or passing to field theory with infinitely
many degrees of freedom removes this uniqueness conclusion.

\phantomsection\label{p-qccr-03}{}

\textbf{P-QCCR-03.} On the common invariant dense Schwartz domain, \[
[Q_j,P_k]\psi=i\hbar\delta_{jk}\psi.
\]

\textbf{Proof.} Apply the product rule to \(Q_j\psi=x_j\psi\) and
\(P_k\psi=-i\hbar\partial_k\psi\). Both compositions are defined on this
domain. \(\square\)

Finite matrices cannot satisfy \([Q,P]=i\hbar I\) exactly: the left-hand
trace vanishes and the right-hand trace does not. A truncated oscillator
must report its boundary correction, rather than treating finite
matrices as an exact CCR realization.

\section{11.4 Uncertainty and the limits of closed-orbit arguments}

For normalized \(\psi\) in the relevant product-operator domains,
Cauchy--Schwarz gives \[
\Delta A\,\Delta B\ge\tfrac12|\langle[A,B]\rangle|.
\] Take the imaginary part of the inner product of
\((A-\langle A\rangle)\psi\) and \((B-\langle B\rangle)\psi\). Retaining
the real part gives the stronger Robertson--Schrödinger bound.

A classical periodic orbit can also satisfy a chosen action integral
\(nh\) while its functions still commute under ordinary multiplication.
Thus \(\oint p\,dq=nh\) does not establish operator CCR. WKB
closed-orbit quantization can also contain a Maslov correction; it is
not an exact universal spectral formula.

\section{\texorpdfstring{11.5 A common \(\hbar\) across
sectors}{11.5 A common \textbackslash hbar across sectors}}

Two independent phase spaces may have
\(B=\kappa_1\sigma_1\oplus\kappa_2\sigma_2\); compositional consistency
does not force the coefficients to agree. If full symplectic covariance
mixing the sectors on the joint phase space is added, P-QCCR-01 gives a
common \(\kappa\) and hence one scale. Global physical mixing and
numerical calibration remain sector facts, not consequences of tensor
composition alone.

\chapter{12. Standard nonrelativistic calculation models}\label{ch-12}

An abstract quantum structure becomes usable when it calculates wave
packets, transmission probabilities, and bound energies. Free particles,
the oscillator, wells, and the Coulomb problem show the steps from a
Hamiltonian and boundary conditions to concrete results, while supplying
simple checks for numerical truncation.

\section{12.1 Free wave packets and interference}

The free Hamiltonian \(H=P^2/(2m)\) is self-adjoint on the appropriate
second-order Sobolev domain. Its Fourier components acquire phase
\(e^{-ip^2t/(2m\hbar)}\). An initially minimum-uncertainty Gaussian has
\[
\sigma_x^2(t)=\sigma_0^2+\frac{\hbar^2t^2}{4m^2\sigma_0^2},
\quad \langle X_t\rangle=x_0+p_0t/m.
\] This follows directly from \(X_t=X_0+tP_0/m\), zero initial
covariance, and \(\operatorname{Var}P=\hbar^2/(4\sigma_0^2)\).

For two paths the final amplitude is \(\psi_1+e^{i\phi}\psi_2\), giving
interference term
\(2\operatorname{Re}(e^{i\phi}\psi_2\overline{\psi_1})\). If the paths
are correlated with environmental states \(|E_1\rangle,|E_2\rangle\),
the cross term is multiplied by \(\langle E_1|E_2\rangle\). This is a
visibility factor after partial tracing, not automatically a fundamental
actualization rate.

\section{12.2 The harmonic oscillator}

\[
H=\frac{P^2}{2m}+\frac{m\omega^2Q^2}{2},\quad
a=\sqrt{\frac{m\omega}{2\hbar}}Q+\frac{iP}{\sqrt{2m\hbar\omega}}.
\] On a common domain, \([a,a^\dagger]=I\) and
\(H=\hbar\omega(a^\dagger a+1/2)\). The ground-state equation
\(a\psi_0=0\) gives a Gaussian; successive raising operators give \[
E_n=\hbar\omega(n+1/2),\quad
a|n\rangle=\sqrt n|n-1\rangle.
\] The standard spectral theorem for the self-adjoint harmonic
oscillator establishes the full spectrum and completeness of the basis.

Truncation to \(0,\ldots,N-1\) gives \[
[a_N,a_N^\dagger]=I-N|N-1\rangle\langle N-1|.
\] Entrywise evaluation of the two diagonal products proves this exact
diagnostic of the finite numerical boundary.

\section{12.3 Wells and barriers}

The Dirichlet infinite well on \((0,L)\) has \[
\psi_n=\sqrt{2/L}\sin(n\pi x/L),\quad
E_n=\frac{\hbar^2\pi^2n^2}{2mL^2}.
\] Solve the second-order eigenvalue equation and impose the two zero
boundary values to obtain this spectrum.

For the finite well \(V=-V_0\) on \(|x|<a\), bound energies obey
\(-V_0<E<0\). Write
\(k=\sqrt{2m(E+V_0)}/\hbar,\kappa=\sqrt{-2mE}/\hbar\). Matching the
wavefunction and its derivative gives \[
k\tan(ka)=\kappa\quad\text{even states},\qquad
-k\cot(ka)=\kappa\quad\text{odd states}.
\] Root searches must respect continuous branches of the trigonometric
functions. Crossing a pole can produce a false eigenvalue.

A rectangular barrier of height \(V_0>E\) and width \(a\) has
transmission \[
T(E)=\left[1+
\frac{V_0^2\sinh^2(\kappa a)}{4E(V_0-E)}
\right]^{-1},\quad
\kappa=\sqrt{2m(V_0-E)}/\hbar.
\] Match plane and exponential waves in the three regions and take the
current ratio. Different asymptotic potentials require an additional
velocity ratio. The current
\(j=(\hbar/m)\operatorname{Im}\bar\psi\nabla\psi\) checks conservation
of reflected and transmitted flux.

When both exterior potentials vanish, let \(k=\sqrt{2mE}/\hbar\). Inside
\(0<x<a\), the wavefunction and derivative obey \[
\begin{pmatrix}\psi(a)\\\psi'(a)\end{pmatrix}
=\begin{pmatrix}\cosh\kappa a&\sinh(\kappa a)/\kappa\\
\kappa\sinh(\kappa a)&\cosh\kappa a\end{pmatrix}
\begin{pmatrix}\psi(0)\\\psi'(0)\end{pmatrix}.
\] Insert boundary values from \(e^{ikx}+re^{-ikx}\) on the left and
\(te^{ikx}\) on the right, then eliminate \(r\): \[
|t|^{-2}=\cosh^2(\kappa a)+
\frac{(\kappa^2-k^2)^2}{4\kappa^2k^2}\sinh^2(\kappa a)
=1+\frac{(\kappa^2+k^2)^2}{4\kappa^2k^2}\sinh^2(\kappa a).
\] Using \(\kappa^2+k^2=2mV_0/\hbar^2\) gives the stated \(T=|t|^2\).
Energy, width, and mass all come from the specified model.

\section{12.4 The Coulomb problem and atoms}

Given a Coulomb potential and reduced mass \(\mu\), \[
H=-\frac{\hbar^2}{2\mu}\Delta-\frac{Ze^2}{4\pi\epsilon_0r}.
\] Separate spherical harmonics and require radial regularity at the
origin and integrability at infinity. Termination of the confluent
hypergeometric series gives \[
E_n=-\frac{\mu Z^2e^4}{2(4\pi\epsilon_0)^2\hbar^2n^2}.
\] With \(a_0=4\pi\epsilon_0\hbar^2/(\mu Ze^2)\), the ground state is
\(\psi_{100}=(\pi a_0^3)^{-1/2}e^{-r/a_0}\). This is a standard
calculation after choosing the electromagnetic potential, mass, spatial
dimension, and quantum sector. It does not count hydrogenic spectral
constants from Ω.

Put \(g=Ze^2/(4\pi\epsilon_0)\) and \(\psi=u(r)Y_{\ell m}/r\). The
bound-state radial equation is \[
u''+\left[-\kappa^2+\frac{2\mu g}{\hbar^2r}
-\frac{\ell(\ell+1)}{r^2}\right]u=0,
\qquad \kappa=\frac{\sqrt{-2\mu E}}{\hbar}.
\] Substitute \(s=2\kappa r\) and \(u=e^{-s/2}s^{\ell+1}v(s)\) to obtain
\[
s v''+(2\ell+2-s)v'+\left(\frac{\mu g}{\hbar^2\kappa}-\ell-1\right)v=0.
\] In the usual regular self-adjoint Coulomb realization, regularity at
the origin and square integrability at infinity require the series to
terminate: \(\mu g/(\hbar^2\kappa)=n_r+\ell+1=n\). Eliminating
\(\kappa\) yields \(E_n=-\mu g^2/(2\hbar^2n^2)\). The discrete energies
thus come from a concrete differential operator and global boundary
conditions. Local phase-closure intuition cannot replace the eigenvalue
problem.

Many-electron problems also require antisymmetry, electron interactions,
and approximation errors. Hartree--Fock, configuration interaction, and
density-functional theory can be connected to this framework; a one-body
solution does not automatically certify their accuracy.

\chapter{13. The finite-mode Gaussian sector}\label{ch-13}

Gaussian models compress a restricted family of infinite-dimensional
quantum states into means and covariances, making them especially
tractable. This relies on CCR, Gaussian assumptions, and complete
positivity; the chapter derives the covariance uncertainty condition and
explains why checking one input does not certify a channel.

Let \(R=(Q_1,P_1,\ldots,Q_n,P_n)\) satisfy
\([R_j,R_k]=i\hbar\Omega_{jk}\). Define the mean \(d=\langle R\rangle\)
and symmetric covariance \(V_{jk}=\langle\{R_j-d_j,R_k-d_k\}\rangle/2\).

\phantomsection\label{p-qgauss-01}{}

\textbf{P-QGAUSS-01: A necessary covariance condition.} \[
V+\frac{i\hbar}{2}\Omega\succeq0.
\]

\textbf{Proof.} For any complex vector \(z\), the operator
\(A=\sum_jz_j(R_j-d_j)\) satisfies \(\langle A^\dagger A\rangle\ge0\).
Expansion gives the stated quadratic form. \(\square\)

\textbf{K-QGAUSS-01.} A positive-definite covariance matrix has
Williamson form

\[
V=S\operatorname{diag}(\nu_1,\nu_1,\ldots,\nu_n,\nu_n)S^\top,
\]

with symplectic \(S\). A physical Gaussian state requires
\(\nu_j\ge\hbar/2\); a pure Gaussian state has every \(\nu_j=\hbar/2\).
First and second moments do not determine a non-Gaussian state.

A Gaussian channel acts by \[
d'=Xd+d_0,\quad V'=XVX^\top+Y,
\] where \(X,d_0\) are real and \(Y\) is real symmetric, with
input/output mode numbers and their symplectic forms fixed. For equal
input and output mode numbers with the same standard symplectic form,
its CP condition is \[
Y+\frac{i\hbar}{2}(\Omega-X\Omega X^\top)\succeq0.
\] For different input and output mode numbers, the first \(\Omega\) in
the parentheses is the output symplectic form, while the second, between
\(X\) and \(X^\top\), is the input symplectic form; \(X\) is a
rectangular real matrix from the input phase space to the output phase
space.

This standard condition is stronger than covariance positivity for one
output state: consistency with arbitrary auxiliary modes is essential.

A quadratic Hamiltonian \(H=R^\top GR/2+g^\top R\) gives
\(\dot d=\Omega(Gd+g)\) and \(\dot V=\Omega GV+VG^\top\Omega^\top\) from
the Heisenberg equation in the declared units. Damping and noise add
drift and diffusion terms subject to CP.

Here \(G\) is real symmetric, the quadratic term is symmetrically
ordered, and \(g\) is real. Using \([R_j,R_k]=i\hbar\Omega_{jk}\) term
by term gives \(\dot R=(i/\hbar)[H,R]=\Omega(GR+g)\). Expectations yield
the mean equation. The centred operator \(\delta R=R-d\) satisfies
\(\dot{\delta R}=A\delta R\), \(A=\Omega G\), so differentiating the
anticommutator covariance gives \(\dot V=AV+VA^\top\). These equations
depend on the established CCR and the chosen quadratic Hamiltonian;
information geometry alone does not supply them.

The physical covariance condition can be expanded as \[
\left\langle\delta R_j\delta R_k\right\rangle
=V_{jk}+\tfrac12\langle[R_j,R_k]\rangle
=V_{jk}+\frac{i\hbar}{2}\Omega_{jk}.
\] Thus \(\langle A^\dagger A\rangle=z^\dagger(V+i\hbar\Omega/2)z\ge0\),
fixing the sign and factor of one half. The Gaussian-channel CP
condition applies to every input and every auxiliary mode and is a
standard Gaussian-channel theorem. A single output covariance check
establishes only that this input's second moments violate no necessary
condition. Numerical truncation must also use the tail bound of §10.4;
approximate CCR of finite matrices do not replace the actual
infinite-dimensional modes.

\chapter{14. Symmetry, spin, and nonrelativistic particles}\label{ch-14}

Spin labels, free-particle energies, and time reversal arise from
concrete symmetry representations rather than geometric metaphors alone.
We track the distinct inputs of Wigner, rotation, and Galilei structure,
then reconnect representation theory to computable probabilities through
Rabi and CHSH benchmarks.

This chapter and Chapters 15--32 and 36--46 start from the standard
quantum sector and the additional physical inputs stated in each
chapter. Propositions proved in prose, imported major theorems,
exact-arithmetic regression checks, floating-point calculations, and
Lean theorems are registered as different evidence types. The finite
Kraus kernel makes some operational contracts machine-checkable. It does
not automatically prove the infinite-dimensional representations,
self-adjoint domains, field theory, or thermodynamic limits introduced
from this chapter onward. Formalized-proposition registers and frozen
receipts define the precise coverage.

\section{14.1 Wigner's theorem and generators}

\textbf{K-QSYM-01: Wigner.} A bijection preserving all ray transition
probabilities is implemented by a unitary or antiunitary operator,
unique up to phase. This standard theorem assumes a complex Hilbert ray
space already exists. It does not reconstruct quantum theory from a
general state set, object names, or tests at finitely many angles.

For a suitably continuous group action on rays and a continuous choice
of implementation near the identity, the connected identity component
has a unitary projective representation. The projective multiplier still
matters; this does not establish a globally single-valued ordinary
unitary representation. Where applicable, one passes to a covering group
or central extension. Discrete time reversal can be antiunitary.

Unbounded generators are not freely manipulable matrices. Stone's
theorem supplies a self-adjoint generator for each strongly continuous
one-parameter unitary subgroup. A strongly continuous unitary
representation of a finite-dimensional Lie group has a dense space of
smooth vectors on which a common infinitesimal representation and Lie
brackets can be defined. Commutators involving \(P,J,K,H\) below are
first used on such a common invariant domain; closures and domains of
the operators require subsequent treatment. Formal commutation relations
alone do not establish the existence of a global group representation.

A reversible finite-dimensional CPTP map whose inverse is also CPTP is
unitary conjugation. Differentiating continuous dynamics gives \[
\dot\rho=-i[H,\rho]/\hbar.
\] This standard quantum order-isomorphism result excludes
antiunitary-transpose channels and does not select \(H\). For example,
complex-matrix transposition preserves single-system positivity, but its
Choi operator is the swap, with a negative eigenvalue on the
antisymmetric subspace, so transposition is not CP. Antiunitarity
describes an implementation of a ray symmetry; it does not thereby
belong to an implementable CPTP control library on arbitrary systems.

\section{14.2 Rotations and spin}

Finite-dimensional irreducible unitary representations of \[
[J_i,J_j]=i\hbar\epsilon_{ijk}J_k
\] are labelled by \(s=0,\tfrac12,1,\ldots\) and satisfy \[
J^2=\hbar^2s(s+1),\qquad J_z=\hbar m,
\qquad m=-s,-s+1,\ldots,s.
\] With \(J_\pm=J_x\pm iJ_y\), the algebra gives \[
\|J_\pm|s,m\rangle\|^2
=\hbar^2[s(s+1)-m(m\pm1)].
\] Nonnegative norms, termination at highest and lowest weights, and
unit spacing of adjacent weights require \(2s\in\mathbb N\). This
explains the discrete representation labels. Single-valued orbital
scalar spherical harmonics have integer \(\ell\); intrinsic spin is a
separately specified internal representation.

A half-integer-spin vector changes sign under a \(2\pi\) rotation while
its ray is unchanged. The \(SU(2)\) cover handles the corresponding
projective representation of \(SO(3)\). A picture of a half twist in
space does not replace representation theory. An \(SU(2)\)
representation alone also does not fix a particle's mass, magnetic
moment, or \(g\) factor; these require a Hamiltonian, matter model, or
experimental input.

\section{14.3 The Galilei central extension and massive particles}

For generators \(P_i,J_i,K_i,H\), write on a common smooth domain \[
[K_i,P_j]=i\hbar\delta_{ij}M,\qquad
[K_i,H]=i\hbar P_i,\qquad [P_i,H]=0,
\] where \(M\) commutes with every generator. Irreducibility makes this
central operator scalar, \(M=mI\). We explicitly choose the free massive
branch \(m>0\). Massless and other singular Galilei representations are
not covered by the divisions by \(m\) below.

Define \[
Q_i=K_i/m,\qquad
\mathbf S=\mathbf J-\mathbf Q\times\mathbf P,
\qquad C_E=H-\frac{P^2}{2m}.
\] The first equation gives orbital CCR. The full Galilei commutation
relations make \(\mathbf S\) commute with \(\mathbf Q,\mathbf P\), while
satisfying its own angular-momentum algebra. Also, \[
[K_i,P^2/(2m)]=i\hbar P_i,
\] so \(C_E\) commutes with boosts, as well as spatial translations,
rotations, and time translations. A Schur argument in the specified
irreducible representation gives \(C_E=E_{\rm int}I\). For a regular
orbital representation and finite spin factor, \[
\mathcal H_1=L^2(\mathbb R^3)\otimes\mathbb C^{2s+1},\qquad
H=P^2/(2m)+E_{\rm int}.
\] The labels are \((m,s,E_{\rm int})\), not a mass spectrum computed
from winding numbers. Particle irreducibility refers to the centrally
extended symmetry group. Weyl displacements act only on the orbital
factor, so their reducibility on the full space including spin is
consistent.

The central extension also has a directly checkable consequence:
interchanging a boost and a spatial translation produces a phase
proportional to \(m\mathbf v\cdot\mathbf a/\hbar\), with sign fixed by
exponential conventions. The historical benchmark applied both operator
sequences independently to check this phase instead of inserting the
same target formula twice. It verifies the chosen representation, not
the numerical values of \(m\) or \(\hbar\).

\section{14.4 Time reversal and standard benchmarks}

Time reversal sends \(Q\to Q,P\to-P\). A unitary implementation would
change the sign of the CCR's left side without changing \(i\), a
contradiction. An antiunitary also sends \(i\to-i\), restoring
consistency. In the standard spin representation,
\(T=e^{-i\pi J_y/\hbar}K\) and \(T^2=(-1)^{2s}\). If \(T^2=-1\) and the
Hamiltonian is time-reversal invariant, orthogonality of \(\psi,T\psi\)
gives Kramers degeneracy.

For Stern--Gerlach measurement of spin one half along \(\mathbf n\),
effects \((I\pm\mathbf n\cdot\boldsymbol\sigma)/2\) give
\(p_\pm=(1\pm\mathbf r\cdot\mathbf n)/2\) for Bloch vector
\(\mathbf r\). The Rabi Hamiltonian \(H=\hbar(\Delta Z+\Omega X)/2\)
gives the transition probability between computational basis states \[
P_{0\to1}(t)=\frac{\Omega^2}{\Omega^2+\Delta^2}
\sin^2\left(\frac{\sqrt{\Omega^2+\Delta^2}\,t}{2}\right).
\] Expand the exponential using
\((\Delta Z+\Omega X)^2=(\Delta^2+\Omega^2)I\).

CHSH has absolute bound two for local deterministic \(\pm1\)
assignments. A quantum Bell state and appropriate Pauli directions
attain \(2\sqrt2\). The identity
\(\mathcal B^2=4I-[A_0,A_1]\otimes[B_0,B_1]\) gives the norm bound
\(2\sqrt2\). Experimental violation depends on actual
measurement-choice, detection, and no-signalling conditions. A
simulation is not a substitute for their realization.

\chapter{15. Identical particles, Fock space, and
superselection}\label{ch-15}

Indifference to particle labels is not enough to determine which
exchange statistics nature realizes. We first examine permutation
representations and their commutants, then construct Bose/Fermi Fock
spaces and explain why superselection depends on observables, reference
frames, and allowed operations.

\section{15.1 Indistinguishability alone does not select Bose or Fermi
statistics}

The labelled tensor product is an auxiliary representation. If
observables commute with the permutation representation of \(S_N\),
decompose the carrier as \[
\mathcal h^{\otimes N}\simeq\bigoplus_\lambda V_\lambda\otimes M_\lambda,
\qquad U_\pi=\bigoplus_\lambda U_\lambda(\pi)\otimes I.
\] Permutation-invariant observables belong to the representation's
commutant algebra and have the form
\(\bigoplus_\lambda I_{V_\lambda}\otimes A_\lambda\), where
\(A_\lambda\) acts on \(M_\lambda\). This algebra is generally
noncommutative. If the density matrix itself is permutation invariant, a
Schur argument gives \[
\rho=\bigoplus_\lambda\frac{I_{V_\lambda}}{\dim V_\lambda}\otimes\tau_\lambda,
\qquad \tau_\lambda\succeq0,\qquad\sum_\lambda\operatorname{Tr}\tau_\lambda=1.
\] For \(N\ge3\), \(S_N\) has irreducible representations beyond the
trivial and sign representations. Invariance of observables under
exchanging labels therefore does not alone exclude other sectors or
force states into only the Bose or Fermi subspace.

Standard Bose/Fermi selection adds a physical-sector condition using \[
P_\pm=\frac1{N!}\sum_{\pi\in S_N}
(\operatorname{sgn}\pi)^{(1\mp1)/2}U_\pi.
\] Group and sign multiplication give \(P_\pm^2=P_\pm=P_\pm^\dagger\).
State support in one of these subspaces must be explicit. The literature
also studies stronger permutation invariance under all auxiliary
extensions, quantum-controlled permutations, and related selection
conditions. Their quantifiers and permitted operations differ from
ordinary permutation invariance and from uniqueness of purification up
to an auxiliary reversible transformation. They cannot be compressed
into one assumption-free indistinguishability principle.

In three-dimensional nonrelativistic theory, spin and statistics remain
independent labels: one can formally construct spin-\(1/2\) Bose sectors
or spin-\(0\) Fermi sectors. These are countermodels to the strength of
a nonrelativistic derivation, not assertions that nature contains such
fundamental particles. The relativistic spin--statistics connection
requires locality, positive energy, positive metric, and appropriate
domain assumptions. In two dimensions, exchange-configuration topology
replaces the ordinary permutation group by the braid group and permits
anyons. This chapter's three-dimensional symmetrization formulas do not
exclude them.

\section{15.2 Fock construction}

\[
\mathcal F_\pm(\mathcal h)=\bigoplus_{N=0}^\infty
P_\pm\mathcal h^{\otimes N}.
\] Bosonic modes obey \([a_i,a_j^\dagger]=\delta_{ij}\); fermionic modes
obey \(\{c_i,c_j^\dagger\}=\delta_{ij}\) and \(\{c_i,c_j\}=0\). For
fermions, \((c_i^\dagger)^2=0\) restricts single-mode occupation to zero
or one. Bosonic occupation is any nonnegative integer.

Bosonic creation and annihilation operators are unbounded and are first
defined on the common domain of finite-particle vectors. A single
fermionic mode operator is bounded. The number operator is
\(N=\sum_i a_i^\dagger a_i\), and second quantization of a one-body
operator is \(d\Gamma(h)=\sum_{ij}h_{ij}a_i^\dagger a_j\). This
organizes a many-body Hilbert space; it does not by itself constitute
relativistic field theory.

For orthogonal orbitals \(u,v\), the two-particle wavefunction \[
\Psi_\pm(x_1,x_2)=
[u(x_1)v(x_2)\pm v(x_1)u(x_2)]/\sqrt2
\] contains exchange interference. For nonorthogonal orbitals, replace
the normalization denominator by
\(\sqrt{2(1\pm|\langle u,v\rangle|^2)}\).

\section{15.3 Operational types of superselection}

Superselection is defined relative to the observable algebra and allowed
operations. If central projections of \(\mathcal A\) define mutually
exclusive sectors and no allowed operation accesses intersector
coherence, the realization has that superselection structure.
Distinguish the centre \(Z(\mathcal A)\) from the commutant
\(\mathcal A'\). A projection in \(\mathcal A'\) may merely select an
unobservable multiplicity factor; it does not automatically define a new
fundamental charge. In algebraic QFT, superselection can also appear
through inequivalent representations and is not exhausted by central
projections on one fixed Hilbert space.

An enlarged reference frame, operation set, or boundary degrees of
freedom can change operational superselection. Bargmann mass
superselection depends on the chosen Galilei representation and
reference framework. It does not prohibit every formal superposition
involving different masses in every extended theory. Fock construction
itself does not force universal particle-number superselection. A
number-related conclusion requires, for example, a fixed species of mass
\(m\), the corresponding mass-superselection rule, and total mass
\(mN\).

Finite-mode Fock spaces and infinite-mode field representations also
differ. Stone--von Neumann uniqueness for finitely many canonical
degrees of freedom cannot directly control all representations of
infinite-mode CCR. Chapter 21 states separate contracts for vacuum,
thermal-state, and other phase representations in quantum fields.

\chapter{16. Many-body structure, locality, and phases}\label{ch-16}

Many-body difficulty is not merely larger matrices: locality,
thermodynamic limits, and controlled approximations must fit together.
Using Ising, Heisenberg, and Hubbard-type models, this chapter separates
single-truncation errors, long-time dynamical errors, and genuine phase
classification, and states what an object certificate must retain.

\section{16.1 Local Hamiltonians and infinite volume}

On a finite lattice \(\Lambda\) with space
\(\bigotimes_{x\in\Lambda}\mathcal H_x\), let \[
H_\Lambda=\sum_{X\subseteq\Lambda}\Phi(X),
\] where \(\Phi(X)\) is self-adjoint and supported on \(X\).
Finite-dimensional local spaces and finite-range or sufficiently rapidly
decaying interactions form an important controlled class.

\textbf{K-QMB-01: Lieb--Robinson interface.} Under bounded local terms,
controlled lattice growth, and suitable interaction decay, local
operators in separated regions obey a bound of the form \[
\|[A_X(t),B_Y]\|
\le C\|A_X\|\|B_Y\||X|
e^{-\mu(d(X,Y)-v|t|)}.
\] The constants depend on the interaction norm. This is not a
relativistic light cone and does not apply unconditionally to unbounded
bosonic modes or arbitrary long-range interactions.

Infinite volume is usually described by a quasilocal \(C^*\) algebra and
state functionals. One must not assume a single trace-one Gibbs matrix
for the entire infinite system. Infinite-volume, long-time, and
zero-external-field limits may fail to commute.

\section{16.2 Ising and Heisenberg models}

The transverse-field Ising Hamiltonian is \[
H=-J\sum_iZ_iZ_{i+1}-h\sum_iX_i.
\] For the uniform infinite one-dimensional chain, standard conventions
and Jordan--Wigner, Fourier, and Bogoliubov transformations give
dispersion \[
\epsilon(k)=2\sqrt{J^2+h^2-2Jh\cos k}.
\] The thermodynamic gap closes at \(|h|=|J|\). For a finite periodic
chain, fermionic boundary conditions depend on the parity sector. A
single momentum grid does not describe every finite state.

For two Heisenberg spins, \(H=(J/4)(X_1X_2+Y_1Y_2+Z_1Z_2)\), total-spin
decomposition gives singlet energy \(-3J/4\) and triplet energy \(J/4\),
directly from
\(\mathbf S_1\cdot\mathbf S_2=(\mathbf S_{\rm tot}^2-\mathbf S_1^2-\mathbf S_2^2)/2\).

The Bose--Hubbard model \[
H=-t\sum_{\langle ij\rangle}(a_i^\dagger a_j+\mathrm{h.c.})
+\frac U2\sum_i n_i(n_i-1)-\mu\sum_i n_i
\] also needs occupation-truncation error control. Fermionic hopping
must preserve CAR signs. Multiplying independent site matrices naively
loses anticommutation; use a fixed-order Jordan--Wigner map or an
explicit Fock basis.

\section{16.3 Correlations, entanglement, and phase classification}

Correlation length, spectral gap, order parameters, entanglement
entropy, and topological indices address different questions. A
finite-system level crossing or a peak in local fluctuations is not
automatically a thermodynamic phase-transition proof. Spontaneous
symmetry breaking generally selects infinite-volume states in the order
\[
\lim_{h\to0^\pm}\lim_{|\Lambda|\to\infty}
\] Reversing the limits may recover a symmetric state. Finite-size
extrapolation must therefore state boundaries, symmetry-breaking fields,
initial states, and observables. A finite-chain curve cannot replace a
limit theorem.

An MPS representation is \[
|\psi\rangle=\sum_{i_1\cdots i_L}
 v_L^\dagger A_1^{i_1}\cdots A_L^{i_L}v_R
 |i_1\cdots i_L\rangle,
\] The bond dimension bounds the Schmidt rank across each cut. For one
cut, discarding Schmidt weight \(\delta<1\) from a normalized pure state
gives overlap \(\sqrt{1-\delta}\) with the normalized truncation and
pure-state trace distance exactly \(\sqrt\delta\). This one-cut result
is not a uniform error bound for an entire chain at every time.
Sequential compression, time discretization, Trotter error, finite
volume, and initial approximation must still be combined.

In one dimension, finite local dimension, controlled local interactions,
and a uniform spectral gap can support area-law and MPS approximation
results. Existence of an approximation, computational complexity of
finding it, and low-bond-dimension representation of arbitrary highly
entangled real-time states are separate claims. An MPS description of a
low-energy ground state does not establish efficient general many-body
dynamics.

Benchmarks must discriminate between correct and incorrect
implementations. Single-particle hopping never exchanges two fermions. A
fermionic backend therefore needs genuinely multiparticle CAR checks or
an explicitly antisymmetric Fock basis, for example a fixed-\(N=2\)
chain compared with an independent second-quantized construction.

\section{16.4 Many-body object certificates}

Candidates may be localized excitations, logical subspaces, effective
order-parameter states, or spatial regions. First fix the target future,
initial-state/environment class, permitted interventions, persistence
window, and resource scale; then specify each certificate's distance and
estimator. One reusable register is \[
(\epsilon_{\rm pred},\epsilon_{\rm mem},\epsilon_{\rm leak},
\Pi_{\rm pers},R_\delta,C_{\rm glob}).
\] These record predictive error for the chosen protocols, residual
error after selecting a memory model, leakage, persistence, accessible
record redundancy, and global coherence. They are not undefined
platform-independent numbers. The quantity \(\epsilon_{\rm mem}\)
measures omitted memory; it does not require the true physical process
to be Markovian.

Shared correlations do not prove causal irreducibility of a parent
object. Deriving effective object dynamics from a many-body Hamiltonian
requires a projection, a reachable input class, and memory control. A
band diagram or correlation plot does not complete the Core interface. A
response-sufficient general object also need not be a strong \(\Omega\)
object with internal repair and resource maintenance; the latter needs
explicit contracts for legitimate damage, repair, and persistence.

Likewise, a Holevo-information plateau in environmental fragments, total
mutual information, the smallest fragment needed for a specified record
accuracy, and the number of independently accessible redundant fragments
are not interchangeable. Strong objectivity requires the relevant
operational conditions. Several fragments may carry the same public
record while their joint parent retains global coherence invisible to
public readout. The nonidentifiability bounds of Chapters 6--8 and 41
continue to apply to many-body systems.

\chapter{17. Open systems: GKSL dynamics, thermal baths, and
trajectories}\label{ch-17}

Tracing out an environment yields a channel under suitable initial-state
assumptions, but not necessarily a memoryless semigroup. This
distinction organizes the role of GKSL structure, the physical inputs of
the Davies weak-coupling limit, and the existence of different
conditional trajectories for the same master equation.

\section{17.1 Microscopic reduction and the semigroup assumption}

For initial state \(\rho_S\otimes\sigma_E\) and joint unitary evolution,
\[
\Lambda_t(\rho_S)=\operatorname{Tr}_E[
U_t(\rho_S\otimes\sigma_E)U_t^\dagger]
\] is CPTP. In general, however, \(\Lambda_{t+s}=\Lambda_t\Lambda_s\)
fails because the environment retains correlations. If changing the
input also changes the initial correlations, there need not be a common
CP assignment on every system state.

\textbf{K-QOPEN-01: Finite-dimensional GKSL.} The generator of a
continuous CPTP semigroup has exactly the form \[
\mathcal L(\rho)=-i[H,\rho]/\hbar+
\sum_\mu\left(L_\mu\rho L_\mu^\dagger-
\tfrac12\{L_\mu^\dagger L_\mu,\rho\}\right).
\] This characterizes physical semigroups. It does not derive the
microscopic origin of particular \(L_\mu\).

\section{17.2 Two explicit channels}

Amplitude damping with \(\eta=e^{-\gamma t}\) has Kraus operators \[
K_0=|0\rangle\langle0|+\sqrt\eta|1\rangle\langle1|,
\quad K_1=\sqrt{1-\eta}|0\rangle\langle1|.
\] The excited-state population becomes \(\eta\rho_{11}\) and coherence
becomes \(\sqrt\eta\rho_{01}\).

Pure dephasing with \(L=\sqrt{\gamma/2}Z\) gives
\(\dot\rho=(\gamma/2)(Z\rho Z-\rho)\). The diagonal is unchanged and
off-diagonal entries acquire \(e^{-\gamma t}\). This is also ordinary
two-branch SRN noise. Its SRN form does not make it evidence for a new
collapse mechanism.

\section{17.3 The Davies interface}

GKSL characterizes continuous CPTP semigroups in finite dimension.
Deriving one from a particular bath requires microscopic conditions. Let
\[
H_S=\sum_\epsilon\epsilon\Pi_\epsilon,\qquad
H_{
\rm int}=\lambda\sum_\alpha A_\alpha\otimes B_\alpha,
\] Specify the stationary bath state, correlation decay or mixing,
coupling-operator conditions, and the weak-coupling limit
\(\lambda\to0\) with macroscopic time scaled as \(\lambda^{-2}\). Finite
but long times, nonzero coupling, and error ranges require additional
estimates. The limiting generator is not an exact equation at every
coupling strength.

Define Bohr-frequency components by \[
A_\alpha(\omega)=
\sum_{\epsilon'-\epsilon=\hbar\omega}
\Pi_\epsilon A_\alpha\Pi_{\epsilon'},
\qquad [H_S,A_\alpha(\omega)]=-\hbar\omega A_\alpha(\omega).
\] With consistent Fourier conventions, the bath spectral matrix
\(\gamma_{\alpha\beta}(\omega)\) is positive semidefinite. Each
frequency block \[
\sum_{\alpha\beta}\gamma_{\alpha\beta}(\omega)
\left[A_\beta(\omega)\rho A_\alpha(\omega)^\dagger
-\tfrac12\{A_\alpha(\omega)^\dagger A_\beta(\omega),\rho\}\right]
\] then has GKSL form. The thermal KMS relation is \[
\gamma_{\alpha\beta}(-\omega)
=e^{-\beta_T\hbar\omega}\gamma_{\beta\alpha}(\omega),
\qquad \beta_T=(k_BT)^{-1}.
\] In this standard weak-coupling limit, the Lamb shift commutes with
\(H_S\). Degenerate or equal-frequency transitions must be treated
together; cross terms within one Bohr block cannot be dropped
arbitrarily.

Using
\(A_\alpha(\omega)\rho_T=e^{-\beta_T\hbar\omega}\rho_T A_\alpha(\omega)\)
and the thermal relation between positive and negative frequencies
verifies Gibbs stationarity block by block. Stationarity means only
\(\mathcal L(\rho_T)=0\). Uniqueness, mixing rates, and thermalization
from arbitrary initial states require further irreducibility or
ergodicity conditions. Extra conserved quantities or noiseless subspaces
can leave multiple stationary states.

For a two-level gap \(\Delta>0\), thermal Bose occupation \[
n_T=(e^{\beta_T\Delta}-1)^{-1}
\] gives \(\gamma_\downarrow=\kappa(n_T+1)\) and
\(\gamma_\uparrow=\kappa n_T\), with upward/downward ratio
\(e^{-\beta_T\Delta}\). Rates are thus tied to the specified bath rather
than fitted arbitrarily to a desired equilibrium state.

\section{17.4 Quantum trajectories and computation}

In a short interval \(dt\), the jump probability is
\(dp_\mu=dt\,\langle L_\mu^\dagger L_\mu\rangle\). After a jump,
normalize \(L_\mu|\psi\rangle\); no-jump evolution uses
\(H_{\rm eff}=H-i\hbar\sum L_\mu^\dagger L_\mu/2\). Averaging outcomes
to order \(dt\) gives GKSL.

Different measurements of the environment produce different unravellings
with the same unconditioned master equation. Mathematical existence of a
stochastic trajectory is therefore not experimental selection of an
ontological collapse account.

Numerical checks should include trace, Hermiticity, positivity,
step-size convergence, and stationary residual
\(\|\mathcal L(\rho_*)\|\). A small residual alone does not bound
distance to a unique stationary state; control of an inverse generator
or a mixing gap is needed.

\chapter{18. Non-Markovian processes and memory objects}\label{ch-18}

Two histories can leave the same current system state while their
environments retain information that changes the future. A two-qubit
causal-break example makes the omission explicit; process tensors and
memory kernels then explain when the state description must be enlarged
instead of refining an incorrect Markov approximation.

\section{18.1 Why one-time channels are insufficient}

CP divisibility requires, for \(t\ge s\), a CPTP map \(V_{t:s}\) with
\(\Lambda_t=V_{t:s}\Lambda_s\). It constrains a family of one-time maps,
usually measured under a fixed initial system--environment preparation.
Interventions change intermediate correlations. Divisibility of one-time
maps therefore does not determine all multitime intervention statistics.

A process tensor maps a sequence of CP interventions to a final state or
record probabilities. With explicit input/output leg order, Choi
normalization, and transpose convention, it is characterized by a
positive Choi operator and sequential causal trace constraints: tracing
out the final output yields the identity on the corresponding input
tensored with the preceding comb. For history-dependent instruments,
retain every outcome and only then select the required complete-record
event. Interchanging an input and output leg or omitting a transpose
changes the pairing; it is not a harmless notation change.

\textbf{K-QMEM-01.} Under complete instrument protocols, operational
Markovianity requires future conditional statistics after a current
causal break to be independent of earlier operations and outcomes. In a
common Choi convention, a Markov process has the temporal-product form
of an initial state and successive channel Choi operators. This
multitime operational criterion is used with the conditions of its
original theorem; CP divisibility cannot replace it.
\href{https://arxiv.org/abs/1512.00589}{Original research}

``Complete'' carries a quantifier. A small set of passive readouts can
hide environmental memory that an enlarged set of active interventions
reveals. Sufficiency proved by the finite Kraus kernel for all finite
words and complete finite adaptive trees of a fixed time-homogeneous
operation library does not automatically extend to new environmental
access, drift, a time-inhomogeneous library, or undeclared process
slots.

\section{18.2 A fully executable memory counterexample}

Let system \(S\) and environment \(E\) be qubits. The historical input
is \(S=|b\rangle\) and the environment starts at zero. First,
\(\mathrm{CNOT}_{S\to E}\) writes \(b\) to the environment. Discard and
reprepare the system as \(|0\rangle\), then apply
\(\mathrm{CNOT}_{E\to S}\) to write the memory back. Both histories have
the same system state at the break but futures \(|0\rangle\) and
\(|1\rangle\).

\phantomsection\label{p-qmem-01}{}

\textbf{P-QMEM-01: Irreducible error from omitted memory.} If the true
future states of two histories have trace distance \(d\), any predictor
using only their common current state has worst-case trace-distance
error at least \(d/2\).

\textbf{Proof.} The predictor returns the same \(\tau\), and
\(d\le D(\sigma_0,\tau)+D(\tau,\sigma_1)\). \(\square\)

Here \(d=1\), so the lower bound is \(1/2\). This standard quantum
process refutes an incorrect Markov compression; it does not violate
standard QM.

\section{18.3 Memory kernels and simulation choices}

Let the joint linear evolution be \(\dot R=\mathcal LR\), with
projection superoperator \(\mathcal P\) and \(\mathcal Q=1-\mathcal P\).
Where the operators and semigroups are well defined, solve \[
\mathcal QR(t)=e^{t\mathcal Q\mathcal L}\mathcal QR(0)
+\int_0^t e^{(t-s)\mathcal Q\mathcal L}
\mathcal Q\mathcal L\mathcal PR(s)\,ds.
\] Substitution into the \(\mathcal P\) equation gives the
Nakajima--Zwanzig identity \[
\frac d{dt}\mathcal PR(t)=\mathcal P\mathcal L\mathcal PR(t)
+\int_0^t\mathcal P\mathcal L e^{(t-s)\mathcal Q\mathcal L}
\mathcal Q\mathcal L\mathcal PR(s)\,ds
+\mathcal P\mathcal L e^{t\mathcal Q\mathcal L}\mathcal QR(0).
\] For a concrete system projection it can be written \[
\dot\rho_S(t)=\mathcal L_0\rho_S(t)
+\int_0^t\mathcal K(t-s)\rho_S(s)\,ds+\mathcal I(t).
\] The kernel and inhomogeneous term depend on the projection and
initial correlations. They are not a universal, unique number measuring
non-Markovianity. Nor does the exact identity guarantee CP for an
arbitrarily truncated approximate kernel; appropriate structural
conditions or actual protocol probabilities must establish it
separately.

Excitation exchange with a single environment gives \(G(t)=\cos gt\),
with revivals in excited-state distinguishability \(D(t)=\cos^2gt\). The
time-local rate \(2g\tan gt\) is singular where the map is noninvertible
and cannot be treated across its zeros as an ordinary smooth semigroup
generator. A distinguishability revival is evidence within this model,
not a replacement for the process criterion over all permitted
interventions.

HEOM, pseudomodes, chain mappings, and tensor networks are possible
simulation backends. Their errors arise from bath-spectrum
approximation, auxiliary-level truncation, time integration, and
initial-state assumptions. Compare them only after translating errors
into a common target response. Process tensors provide an operational
comparison language, not a guarantee of lower computational cost in
every problem.

\chapter{19. Thermal equilibrium, response, and entropy
production}\label{ch-19}

Free-energy minimality of a Gibbs state, stationarity under a thermal
bath, and actual thermalization are different questions. Relative
entropy, KMS, and linear response connect them here while retaining the
limits imposed by degeneracies, finite-system recurrence, and heat/work
definitions at strong coupling.

\section{19.1 The variational meaning of Gibbs equilibrium}

For a finite system, let \(\rho_T=e^{-\beta_TH}/Z\) and
\(F_T(\rho)=\operatorname{Tr}\rho H-\beta_T^{-1}S(\rho)\).

\phantomsection\label{p-qtherm-01}{}

\textbf{P-QTHERM-01: Free energy and relative entropy.} \[
F_T(\rho)-F_T(\rho_T)
=\beta_T^{-1}D(\rho\|\rho_T)\ge0.
\]

\textbf{Proof.} Substitute \(\log\rho_T=-\beta_TH-\log Z\) and
rearrange. Nonnegativity of relative entropy gives minimality.
\(\square\)

Given a Hamiltonian, temperature, and entropy definition, equilibrium
thus has an exact resource--information variational expression. This
does not prove that isolated unitary systems converge to Gibbs states
from arbitrary initial states. Unitary evolution preserves the global
spectrum and entropy.

\section{19.2 KMS and spectral detailed balance}

Define \(\tau_t(B)=e^{iHt/\hbar}Be^{-iHt/\hbar}\). The equilibrium KMS
boundary condition is \[
F(t)=\langle A\tau_t(B)\rangle_T,\quad
F(t+i\beta_T\hbar)=\langle\tau_t(B)A\rangle_T.
\]

\phantomsection\label{p-qtherm-02}{}

\textbf{P-QTHERM-02: KMS for finite Gibbs systems.} The stated boundary
condition holds.

\textbf{Proof.} Analytic continuation gives \(\tau_{t+i\beta_T\hbar}(B)=
e^{-\beta_TH}\tau_t(B)e^{\beta_TH}\). Insert this into the Gibbs trace
and cyclically move factors to obtain the other boundary. Finite
dimension guarantees analyticity. \(\square\)

For Hermitian \(B\), define
\(S_B(\omega)=\int dt\,e^{i\omega t}\langle B(t)B(0)\rangle_T\). The
Lehmann expansion gives \[
S_B(\omega)=2\pi\sum_{nm}p_n|B_{nm}|^2
\delta(\omega-(E_m-E_n)/\hbar),
\quad S_B(-\omega)=e^{-\beta_T\hbar\omega}S_B(\omega).
\] For infinite systems, KMS is defined through \(C^*\) dynamics and
analytic elements. It cannot require a nonexistent finite global
partition function.

\section{19.3 Kubo response and fluctuation--dissipation}

With perturbation \(H(t)=H_0-f(t)B\), the first-order
interaction-picture expansion yields \[
\delta\langle A(t)\rangle=
\int_{-\infty}^tds\,f(s)\chi_{AB}(t-s),\quad
\chi_{AB}(t)=\frac{i}{\hbar}\theta(t)\langle[A(t),B]\rangle_T.
\] The equilibrium spectral expansion gives \[
\chi_{BB}''(\omega)
=\frac{1-e^{-\beta_T\hbar\omega}}{2\hbar}S_B(\omega)
\] with the Fourier convention above. Response signs depend on the
perturbation convention; \(+fB\) and \(-fB\) cannot be interchanged
silently.

\section{19.4 Stationarity, equilibration, and thermalization}

The long-time average of a finite closed system retains complete
degenerate energy blocks: \[
\overline\rho=\lim_{T\to\infty}\frac1T\int_0^T
 e^{-iHt/\hbar}\rho_0e^{iHt/\hbar}\,dt
=\sum_E P_E\rho_0P_E.
\] In an energy basis, a coefficient with \(E\ne E'\) averages as
\(T^{-1}\int_0^T e^{-i(E-E')t/\hbar}dt\to0\), whereas within one energy
block it stays one. Coherence inside a degenerate block cannot be
removed arbitrarily. Retaining only the diagonal in a chosen eigenvector
basis would give the wrong result.

Existence of the average does not imply convergence of instantaneous
states. For \(H=\hbar\omega Z/2\) and initial \(|+x\rangle\),
\(\langle X\rangle_t=\cos\omega t\). Its mean is zero, but late-time
oscillations persist. Finite spectra also allow recurrence. Statements
that most late-time local observations lie near equilibrium, or bounds
on mean-square deviations, require additional effective-dimension,
gap-structure, observable-class, and initial-state assumptions.

ETH is a hypothesis about local matrix-element structure in particular
nonintegrable systems, not a theorem for every Hamiltonian.
Integrability and extra conserved quantities can require generalized
Gibbs or other stationary descriptions. Variational existence of
finite-system Gibbs states, stationarity in an open thermal bath,
dynamical equilibration, and Gibbs thermalization must be registered
separately.

\section{19.5 Spohn's inequality and physical entropy production}

\phantomsection\label{p-qtherm-03}{}

\textbf{P-QTHERM-03.} If a CPTP semigroup \(\Lambda_t\) has a full-rank
stationary state \(\sigma\), then \(D(\Lambda_t\rho\|\sigma)\) is
nonincreasing. Along differentiable full-rank trajectories, \[
-\operatorname{Tr}[\mathcal L(\rho)(\log\rho-\log\sigma)]\ge0.
\]

\textbf{Proof.} Relative-entropy data processing gives
\(D(\Lambda_{t+s}\rho\|\Lambda_s\sigma)
\le D(\Lambda_t\rho\|\sigma)\). Use stationarity of \(\sigma\) and take
the right derivative. \(\square\)

For \(\sigma=\rho_T\), a fixed Hamiltonian, and a valid heat/work split,
this becomes \(\dot S-\beta_T\dot Q\ge0\). For an arbitrary reference
state, strong coupling, or a non-Markovian reduction, declining relative
entropy is not automatically total thermodynamic entropy production.

\chapter{20. Quantum GOD/GOA: tasks, control, and long-term
structure}\label{ch-20}

Existence of a quantum object does not specify the task a controller
wants it to perform. Once rewards, resources, and feedback are fixed,
Bellman values and pulse gradients become meaningful; the resulting
long-term structure can be a fixed state, periodic motion, or
conditional survival rather than a compulsory unique endpoint.

\section{20.1 A task is needed to pass from objects to purposes}

Fix permitted actions, resources, measurement feedback, and reward
\(r(\rho,a)\). Target-state fidelity, logical survival, energy
consumption, and information acquisition may favour different actions.
Stability of a quantum object does not automatically maximize all these
criteria.

Without observed feedback, states update as CPTP
\(\rho'=\mathcal T_a(\rho)\). With observations, each action specifies
an instrument, posterior \(\rho_{a,o}\), and probability
\(p(o|\rho,a)\). A control state may also include a classical budget,
device mode, and belief about the environment.

\section{20.2 Bellman equations and a quantum goal direction}

For finite actions, bounded rewards, and discount \(0<\beta_C<1\),
define \[
(TV)(\rho)=\max_a\left[r(\rho,a)+
\beta_C\sum_op(o|\rho,a)V(\rho_{a,o})\right].
\]

\phantomsection\label{p-qctl-01}{}

\textbf{P-QCTL-01.} The operator \(T\) is a \(\beta_C\) contraction on
bounded functions and has a unique bounded fixed point. Under the
required measurability conditions, finite actions admit a measurable
optimal policy.

\textbf{Proof.} For each action, the difference of conditional
expectations is at most \(\|V-W\|_\infty\); taking a maximum preserves
the bound. Banach's fixed-point theorem gives the fixed point. A fixed
ordering of the finite actions breaks ties measurably. Iteration or
discounted telescoping proves optimality. \(\square\)

GOD is the set of optimal actions. If actions correspond to Hamiltonian
directions or Lindblad generators, map it to the corresponding set of
dynamics rather than forcing it to be a Euclidean vector. Changes to
future actions, instruments, or rewards require recomputing value.

\section{20.3 Finite-time quantum-control gradients}

For a closed system with \(H(t)=H_0+\sum_k u_k(t)H_k\) and objective
\(J=\langle\psi_T|O|\psi_T\rangle-\int c(u)dt\), define the adjoint
state \(|\chi_t\rangle=U(T,t)^\dagger O|\psi_T\rangle\). Duhamel
variation gives \[
\frac{\delta J}{\delta u_k(t)}
=\frac2\hbar\operatorname{Im}\langle\chi_t|H_k|\psi_t\rangle
-\partial_{u_k}c.
\] The sign follows directly from
\(\delta\psi_T=-(i/\hbar)U(T,t)H_k\psi_t\,\delta u_kdt\) and its
conjugate contribution. Discrete-pulse gradients must also state
approximations to derivatives of matrix exponentials. Nonconvex control
does not guarantee discovery of a global GOD.

\section{20.4 GOA and stability}

After optimal feedback is fixed, GOA may be a stationary-state set,
limit cycle, noiseless subspace, occupation distribution, or
conditional-survival structure. It is not a universally unique pure
state.

Finite-time approximate channels use the error composition of Chapter 4.
If a closed-loop channel has a uniform trace-distance contraction
coefficient \(\alpha<1\), it has a unique fixed state, and a perturbed
channel's fixed state is within \(\epsilon/(1-\alpha)\). As in Core,
insert a common channel output to obtain \(d\le\alpha d+\epsilon\).
Without contraction, near degeneracy, or at a bifurcation, a small
one-step error need not imply long-term identity stability.

Chapter 52 proves finite-channel stationary-state existence without
optimal control or strict contraction, and gives periodic-orbit and
multiple-stationary-state counterexamples. Chapter 53 first specifies
survival outcomes, then proves a QSD survival law and a different
invariant distribution for the Doob process. These results give actual
matrix certificates for persistence structures; they do not supply an
undeclared evaluation function.

\chapter{21. Relativistic quantum fields: domains and additional
assumptions}\label{ch-21}

Relativistic field theory does not arise simply by enlarging finite
matrices: local algebras, operator distributions, and representation
choices impose new obligations. We state these inputs before free-field
calculations and perturbation theory and explain why restricting a
region is not generally an ordinary partial trace.

Unless stated otherwise, from this chapter onward we use \(\hbar=c=1\),
metric \(\eta=\operatorname{diag}(1,-1,-1,-1)\), and
\(p x=p^0x^0-\mathbf p\cdot\mathbf x\). Mass, momentum, and energy have
the same dimension; length and time have inverse-energy dimension.
Uniqueness results for finitely many continuous-variable degrees of
freedom cannot be transferred directly to fields with infinitely many
degrees of freedom.

\section{21.1 Physical inputs for a field-theoretic realization}

\textbf{H-QFT-1.} Specify a Minkowski background, a strongly continuous
unitary implementation of Poincaré symmetry, a joint energy--momentum
spectrum in the closed forward light cone, locality, a positive-definite
physical state space, and the existence, cyclicity, regularity, and any
required uniqueness conditions of the vacuum. An interacting model also
requires field content, an action, a regulator, and renormalization
conditions. These physical and mathematical inputs select the
field-theory sector; they do not follow from general object closure. The
spectrum condition alone does not guarantee vacuum uniqueness.

A field \(\phi(x)\) is generally an operator-valued distribution. For a
specified test-function class, such as Schwartz functions, define \[
\phi(f)=\int d^4x\,f(x)\phi(x),
\] and discuss closability, products, and distributional matrix elements
on a common invariant dense domain. Point fields, coincident-point
products, and time-ordered products are not automatically
everywhere-bounded matrices; the latter two may require separate
extension or renormalization.

For a bounded region \(\mathcal O\), local observable algebras
\(\mathcal A(\mathcal O)\) satisfy isotony, covariance, and spacelike
commutativity. Bounded Weyl operators generate local algebras of free
bosonic fields. Fermion fields have graded locality, while physical even
observables obey ordinary spacelike commutativity. An object interface
can record the region, algebra, restricted state, and allowed local
instruments: \[
(\mathcal O,\mathcal A(\mathcal O),\omega|_{\mathcal A(\mathcal O)},\mathfrak I_{\mathcal O}).
\] Time/energy resolution and approximation errors are then added. A
sharp region is not assumed in advance to have an independent Hilbert
tensor factor.

\section{21.2 Infinite degrees of freedom, representations, and local
algebras}

Uniqueness of regular irreducible CCR representations for finitely many
canonical degrees of freedom does not extend to infinitely many modes.
In a bosonic Fock representation, unitary implementability of an
appropriate linear canonical Bogoliubov transformation is controlled by
the Hilbert--Schmidt condition on its creation--annihilation mixing part
\(\beta\): \[
\|\beta\|_{\rm HS}^2<\infty.
\] For compatible diagonal transformations of independent modes,
\(\beta_n=1/n\) is square summable whereas \(\beta_n=1/\sqrt n\) is not.
Transformations legitimate mode by mode can thus differ in
implementability on the infinite system. This diagnostic does not
truncate every mode and silently interchange an infinite limit. Vacua,
thermal states, and distinct phases may require inequivalent
representations.

Under standard conditions in usual continuum relativistic models,
regional local algebras have Type III structures rather than those of
finite matrix algebras. Even if \(\mathcal B(\mathcal H)\) on the full
Fock space is Type I, this does not make every local algebra
\(\mathcal B(\mathcal H_{\mathcal O})\otimes I\). A state restricted to
a sharp region need not have a density matrix relative to an ordinary
local trace, and ordinary reduced entanglement entropy need not be
finite.

``Tracing out the exterior'' therefore needs an implementation. One may
use a lattice or energy cutoff; alternatively, under additional
split/nuclearity conditions, separated inner and outer regions admit an
intermediate Type I factor. Such a split is conditional and generally
nonunique. The algebraic restriction \[
\omega\longmapsto\omega|_{\mathcal A(\mathcal O)},
\] remains available in general. Relative entropy or modular-theoretic
quantities may then be used in their proper framework, without asserting
a nonexistent finite local partition function or ordinary local trace.

The Reeh--Schlieder cyclicity result requires its spectrum, vacuum, and
suitable generating assumptions. Cyclicity does not allow arbitrary
remote control. If an observable \(B\) in a spacelike region commutes
with all local Kraus operators \(K_\mu\), and
\(\sum_\mu K_\mu^\dagger K_\mu=I\), then \[
\sum_\mu K_\mu^\dagger B K_\mu=B.
\] A nonselective local operation leaves the other region's expectation
unchanged. Conditional states depending on a particular measurement
result must include the outcome probability and communication of the
classical result.

Under its vacuum, covariance, and field assumptions, Haag's theorem
obstructs a naive global unitary equivalence between free and
interacting fields. It does not make interaction impossible. It requires
distinguishing finitely regulated models, perturbative
interaction-picture constructions, and rigorous theories after regulator
removal. Core's response and restriction language can accommodate these
realizations. Extensions of finite-dimensional density, entropy,
diamond-norm, or matrix methods still need explicit topology, energy
constraints, and limit contracts.

\section{21.3 Three distinct states of completion}

Free fields admit explicit constructions; finitely regulated models can
be defined and computed directly; perturbative series for interacting
four-dimensional theories yield formal or asymptotic predictions order
by order. These are not interchangeable with rigorous nonperturbative
completion. This book does not prove the four-dimensional Yang--Mills
mass gap, full nonperturbative existence of QED, or quantum gravity.

\chapter{22. The free scalar field: modes, propagation, and
causality}\label{ch-22}

The free scalar field provides an explicit starting point for modes,
propagators, and causality. Distinguishing vacuum correlations, Feynman
propagation, and causal response prevents nonzero spacelike correlations
from being mistaken for controllable faster-than-light signals.

\section{22.1 The Klein--Gordon field}

Take \[
\mathcal L=\tfrac12\partial_\mu\phi\,\partial^\mu\phi-\tfrac12m^2\phi^2,\qquad
(\Box+m^2)\phi=0,\qquad E_{\mathbf p}=\sqrt{\mathbf p^2+m^2}.
\] Define \[
[a(\mathbf p),a^\dagger(\mathbf q)]=(2\pi)^3\,2E_{\mathbf p}\,
\delta^3(\mathbf p-\mathbf q),\quad [a,a]=0,
\] \[
\phi(x)=\int\frac{d^3p}{(2\pi)^3\,2E_{\mathbf p}}
\left[a(\mathbf p)e^{-ipx}+a^\dagger(\mathbf p)e^{ipx}\right].
\] The vacuum obeys \(a(\mathbf p)|0\rangle=0\). The normal-ordered
Hamiltonian is \[
:H: =\int\frac{d^3p}{(2\pi)^3\,2E_{\mathbf p}}\,
E_{\mathbf p}a^\dagger(\mathbf p)a(\mathbf p).
\] This selects a positive-frequency decomposition in a flat background.
A general curved spacetime need not have a unique global vacuum.

\phantomsection\label{p-qkg-01}{}

\textbf{P-QKG-01: Canonical relations and propagation.}
Distributionally, these modes give \[
[\phi(t,\mathbf x),\pi(t,\mathbf y)]=i\delta^3(\mathbf x-\mathbf y),
\quad \pi=\partial_t\phi.
\] Substitute the two mode expansions and retain the only nonzero
\(a,a^\dagger\) commutators. The time-derivative factor
\(E_{\mathbf p}\) cancels the measure denominator, and the two terms
each supply half the Fourier delta. Other equal-time commutators vanish
by momentum reversal. \(\square\)

\section{22.2 Three different kernels}

\[
\Delta^+(x)=\langle0|\phi(x)\phi(0)|0\rangle
=\int\frac{d^3p}{(2\pi)^3\,2E_{\mathbf p}}e^{-ipx},
\] \[
[\phi(x),\phi(y)]=\Delta^+(x-y)-\Delta^+(y-x)=i\Delta(x-y),
\] \[
\Delta_F(x)=\langle0|T\phi(x)\phi(0)|0\rangle
=\int\frac{d^4p}{(2\pi)^4}\frac{i\,e^{-ipx}}{p^2-m^2+i0}.
\] Multiplying by the denominator gives
\((\Box+m^2)\Delta_F=-i\delta^4\). Causal response uses a retarded
kernel. The Feynman kernel is not a classical probability for signal
propagation.

\phantomsection\label{p-qkg-02}{}

\textbf{P-QKG-02: Spacelike causality.} For \(x^2<0\), an orthochronous
Lorentz transformation takes \(x\) to \(x^0=0\). The commutator integral
is then odd under momentum reversal and vanishes. Covariance returns the
result to the original coordinates. Suitably smeared observables with
spacelike-separated supports therefore commute. \(\square\)

The spacelike Wightman function need not vanish. For example, with
\(m>0\) and \(r=|\mathbf x|\), \[
\Delta^+(0,\mathbf x)=\frac{m}{4\pi^2r}K_1(mr).
\] Vacuum correlations are not controllable faster-than-light signals.
If a local operation commutes with an effect in another region, the
nonselective operation preserves that effect's expectation. Exploiting
selective conditional statistics still requires communication of the
classical result.

\section{22.3 Sources and detectors}

For an external source \(J\) coupled as \(\int J\phi\), the commutator
determines linear response. A localized two-level detector can be
modelled by \(H_I(\tau)=g\chi(\tau)\mu(\tau)\phi(x(\tau))\). Its
second-order transition probability is an integral of the smeared
\(\Delta^+\). Switching function, trajectory, gap, and regulator are
inputs. An instantaneous point-detector limit may diverge.

\chapter{23. Dirac and Maxwell fields and their physical degrees of
freedom}\label{ch-23}

Dirac and Maxwell fields show how spin, statistics, and gauge
constraints alter the physical degrees of freedom. Mode and constraint
calculations lead to antiparticles, transverse polarizations, and
Gauss's law while identifying which field variables serve only as
auxiliary descriptions.

\section{23.1 Spinors and antiparticles}

The Clifford relations and Dirac Lagrangian \[
\{\gamma^\mu,\gamma^\nu\}=2\eta^{\mu\nu}I,\quad
\bar\psi=\psi^\dagger\gamma^0,\quad
\mathcal L_D=\bar\psi(i\gamma^\mu\partial_\mu-m)\psi
\] give \((i\not\partial-m)\psi=0\). Momentum spinors satisfy \[
(\not p-m)u_s(p)=0,\quad(\not p+m)v_s(p)=0,
\] \[
\sum_su_s\bar u_s=\not p+m,\quad
\sum_sv_s\bar v_s=\not p-m
\] in the usual covariant normalization. The mode expansion is \[
\psi(x)=\sum_s\int\frac{d^3p}{(2\pi)^3\,2E_p}
\left[b_s(p)u_s(p)e^{-ipx}+d_s^\dagger(p)v_s(p)e^{ipx}\right],
\] with \(b,d\) obeying CAR in the same normalization. Normal-ordered
energy counts positive-energy particle and antiparticle excitations;
normal-ordered charge assigns opposite signs.

\phantomsection\label{p-qdir-01}{}

\textbf{P-QDIR-01.} Multiplication of the Dirac operator by its
conjugate factor gives
\((i\not\partial-m)(i\not\partial+m)=-(\Box+m^2)\). Each spinor
component therefore satisfies the KG equation, and \[
S_F(x)=(i\not\partial+m)\Delta_F(x),\qquad
S_F(p)=\frac{i(\not p+m)}{p^2-m^2+i0}.
\] The Clifford relation cancels the antisymmetric part of mixed
derivatives. \(\square\)

The same mode expansion verifies graded causality. Let \[
D_0(x)=\Delta^+(x)-\Delta^+(-x)=i\Delta(x),
\] With Chapter 22's normalization and signs, \[
\{\psi_\alpha(x),\bar\psi_\beta(y)\}
=\bigl[(i\not\partial_x+m)D_0(x-y)\bigr]_{\alpha\beta}.
\] Use the spin sums \(u\bar u=\not p+m\) and \(v\bar v=\not p-m\);
\(i\partial\) supplies opposite momentum signs on positive- and
negative-frequency terms. Differentiation preserves the causal support
of \(D_0\), so spacelike-separated fermion fields anticommute.
Appropriate local observables with even fermion parity obey ordinary
commutativity. A massive-model check must retain the \(m\) term; a
numerical massless substitute does not verify the entire massive
interface.

\textbf{K-QSTAT-01: Spin--statistics.} Under relativistic locality, a
positive-energy spectrum, a positive-definite state space, and the
relevant vacuum/domain assumptions, integer and half-integer spins obey
Bose and Fermi statistics, respectively. This is an imported standard
major theorem, not an independent Core proof. It is stronger than
ordinary indistinguishability in Chapter 15 and does not transfer
unchanged to unchecked two-dimensional braid statistics or other models
outside its domain.

\section{23.2 Maxwell constraints and transverse modes}

Take \[
\mathcal L_M=-\tfrac14F_{\mu\nu}F^{\mu\nu},\qquad
F_{\mu\nu}=\partial_\mu A_\nu-\partial_\nu A_\mu,
\qquad A_\mu\sim A_\mu+\partial_\mu\alpha.
\] The potential \(A_0\) has no independent kinetic term and its
canonical conjugate vanishes. Preserving the constraint in time yields
Gauss's constraint. Constraints and gauge redundancy reduce the four
potential components at each nonzero momentum to two physical transverse
polarizations, the Maxwell photon helicities \(\pm1\).
Polarization-vector phases and local frames are representation choices;
changing a frame does not create another physical polarization.

In Coulomb gauge the transverse projector is \[
P^T_{ij}(\mathbf k)=\delta_{ij}-k_i k_j/|\mathbf k|^2.
\] It is nonlocal in position space, but this does not establish
superluminal observable signals. The causal algebra of physical
electromagnetic fields \(F_{\mu\nu}\) is treated as in Chapter 21. An
auxiliary propagator of a gauge potential, a field-strength commutator,
and preparable or detectable photon states are different objects.

One cannot unconditionally require a nontrivial \(A_\mu\) to be
point-local, Lorentz covariant, subject to the strict Maxwell equations,
and represented on a positive-definite Hilbert space simultaneously.
Standard no-go results for Maxwell potentials explain the use of an
auxiliary indefinite-metric space in covariant gauges, followed by a
subsidiary condition or BRST quotient to obtain a positive physical
state space. An alternative works directly with transverse modes or
local gauge-invariant fields. Treating all four \(A_\mu\) modes as
positive-norm photons overcounts unphysical degrees of freedom.

\section{23.3 Charged objects, Gauss's law, and boundaries}

Gauss's law relates charge to flux through a surrounding surface,
qualifying the intuition that a charged particle is wholly localized in
an independent Hilbert factor. Regional partitions may require flux
centres, edge variables, or dressed operators. Different choices change
entropy and carrier decompositions and must be included in protocol
interfaces. Eliminating gauge redundancy and deleting physical boundary
degrees of freedom are different operations.

\chapter{24. Interactions, generating functionals, and
scattering}\label{ch-24}

Calculations with interactions first require a distinction between bare
and renormalized parameters, regulators, and perturbative order. On that
basis we introduce Dyson expansion, generating functionals, and
scattering, retaining the asymptotic-state and infrared conditions
needed to turn formal expressions into observable cross sections.

\section{\texorpdfstring{24.1 A regulated \(\lambda\phi^4\)
benchmark}{24.1 A regulated \textbackslash lambda\textbackslash phi\^{}4 benchmark}}

Fix notation for bare and renormalized parameters first. The bare
Lagrangian is \[
\mathcal L_0=\tfrac12(\partial\phi_0)^2
-\tfrac12m_0^2\phi_0^2-\frac{\lambda_0}{4!}\phi_0^4.
\] An additional independent set of counterterms must not be added to
this expression. With dimensional regularization in \(d=4-\epsilon\),
define a renormalized field and dimensionless coupling by \[
\phi_0=Z_\phi^{1/2}\phi,\qquad
Z_\phi m_0^2=m^2+\delta m^2,\qquad
Z_\phi^2\lambda_0=\mu^\epsilon(\lambda+\delta\lambda),
\qquad Z_\phi=1+\delta Z_\phi.
\] The same bare expression decomposes as \[
\begin{aligned}
\mathcal L={}&\tfrac12(\partial\phi)^2-\tfrac12m^2\phi^2
-\frac{\mu^\epsilon\lambda}{4!}\phi^4+\mathcal L_{\rm ct},\\
\mathcal L_{\rm ct}={}&\tfrac12\delta Z_\phi(\partial\phi)^2
-\tfrac12\delta m^2\phi^2-\frac{\mu^\epsilon\delta\lambda}{4!}\phi^4.
\end{aligned}
\] This is a reparameterization, not two independent sets of masses and
couplings. In the four-dimensional limit, the two-point counterterm
vertex is \(i(\delta Z_\phi p^2-\delta m^2)\) and the four-point vertex
is \(-i\delta\lambda\); before \(\epsilon\to0\), the latter retains
\(\mu^\epsilon\). QED in Chapter 27 uses \(d=4-2\epsilon_{\rm DR}\),
requiring a corresponding conversion of pole notation. The two epsilon
conventions cannot be mixed.

Global stability of the classical real-scalar potential requires an
appropriate positive quartic term. This is distinct from relations
between bare and renormalized parameters in a perturbative scheme.
Chapter 25's Euclidean hard-cutoff mass benchmark uses an explicit
subtraction scheme. Parameters obtained with different regulators must
be matched before comparing observables.

The interaction-picture notation is \[
U_I(t,t_0)=T\exp\left[-i\int_{t_0}^{t}H_I(s)\,ds\right].
\] Its \(n\)th term is a time-ordered integral; writing it over the full
integration region introduces \(1/n!\). With finite regulation and
suitable operator conditions, dynamics can first be defined directly. In
four-dimensional continuum theory the Dyson series is generally used
perturbatively. Regulator removal, asymptotic-state limits, and series
summation cannot be interchanged because an exponential symbol appears,
nor does that notation evade Haag's theorem in Chapter 21.

\section{24.2 Wick expansion and Feynman rules}

Even time-ordered free-vacuum correlation functions equal the sum of all
pairwise contractions; odd ones vanish (K: Gaussian/Wick theorem).
Combining this with the Dyson series gives scalar internal lines
\(i/(p^2-m^2+i0)\), vertices \(-i\lambda\), independent loop integrals
\(\int d^4\ell/(2\pi)^4\), and symmetry factors from graph automorphisms
and expansion counting.

The vacuum-normalized generating functionals are \[
\begin{gathered}
Z[J]=\frac{\int\mathcal D\phi\,
e^{iS[\phi]+i\int J\phi}}{\int\mathcal D\phi\,e^{iS[\phi]}},
\\ W[J]=-i\log Z[J],\\
\Gamma[\varphi]=W[J]-\int J\varphi,\\
\varphi=\frac{\delta W}{\delta J}.
\end{gathered}
\] These expressions are used with finite regulation or perturbatively.
A Lorentzian path integral is not automatically an ordinary probability
measure. The functional \(W\) generates connected diagrams; \(\Gamma\)
generates 1PI vertices.

\section{24.3 LSZ and cross sections}

\textbf{K-QLSZ-01.} With suitable stable one-particle mass shells,
nonzero pole residues, asymptotic states, and spectral assumptions,
amputating external propagators from connected correlation functions and
inserting external-leg residue factors gives scattering amplitudes.
Infrared problems for long-range charged fields prevent unqualified use
of the simplest charged Fock-asymptotic-state version.

Use \(S=I+iT\) and
\(\langle f|iT|i\rangle=i(2\pi)^4\delta^4(P_f-P_i)\mathcal M\). The
two-particle incident flux is \(4\sqrt{(p_1p_2)^2-m_1^2m_2^2}\), with
final-state measure \[
d\Phi_n=(2\pi)^4\delta^4(P-\sum k_i)
\prod_i\frac{d^3k_i}{(2\pi)^3\,2E_i}.
\] Divide by the appropriate factorials for identical final particles.
For centre-of-mass two-to-two scattering, \[
\frac{d\sigma}{d\Omega}=\frac1{64\pi^2s}
\frac{|\mathbf p_f|}{|\mathbf p_i|}\overline{|\mathcal M|^2}
\] with identical-particle factors and angular-domain conventions
handled consistently. Average initial spins only when the preparation
does not resolve them.

\phantomsection\label{p-qunit-01}{}

\textbf{P-QUNIT-01: Optical theorem.} From \(S^\dagger S=I\),
\(-i(T-T^\dagger)=T^\dagger T\). Take the initial-state diagonal element
and insert a complete final-state set. Then
\(2\operatorname{Im}\mathcal M_{ii}\) equals the phase-space integral of
squared intermediate-state amplitudes, including statistical factors.
\(\square\) This constrains perturbation theory order by order; it does
not make every truncated amplitude exactly unitary.

\chapter{25. Renormalization: explicit one-loop benchmarks and
schemes}\label{ch-25}

One-loop integrals make renormalization concrete: regulator dependence
enters parameter relations, while matched quantities enter physical
comparisons. Mass subtraction, bubbles, thresholds, and running are
checked step by step, showing why beta functions from different schemes
cannot be mixed without conversion.

\section{25.1 The two-point function}

With a Euclidean spherical hard cutoff, \[
I_\Lambda(m)=\int_{|q|<\Lambda}\frac{d^4q}{(2\pi)^4}
\frac1{q^2+m^2}
=\frac1{16\pi^2}\left[\Lambda^2-m^2\log(1+\Lambda^2/m^2)\right].
\] \phantomsection\label{p-qren-01}{}

\textbf{P-QREN-01.} The four-dimensional unit-sphere area is \(2\pi^2\).
Setting \(u=q^2\) and integrating \(u/(u+m^2)\) gives the expression
above. The quartic tadpole has combinatorial factor \(1/2\), hence
\(\Sigma^{(1)}=\lambda I_\Lambda/2\). Absorb this constant using a
zero-momentum mass condition:
\(m_0^2=m_R^2-\lambda I_\Lambda(m_R)/2+O(\lambda^2)\), which gives
\(\Gamma_R^{(2)}(p_E)=p_E^2+m_R^2+O(\lambda^2)\). The one-loop diagram
is momentum independent, so its wavefunction-renormalization
contribution at this order vanishes. \(\square\)

The quadratic cutoff term is a scheme-dependent parameter relation, not
an independent observable. A regulator change need not change physical
predictions.

\section{25.2 The four-point function and thresholds}

At a specified subtraction configuration with three equivalent Euclidean
channels, define \[
B_R(Q;\mu,m)=\frac1{16\pi^2}\int_0^1dx
\log\frac{m^2+x(1-x)\mu^2}{m^2+x(1-x)Q^2},
\] \[
F(Q)=\lambda-\tfrac32\lambda^2 B_R(Q;\mu,m)+O(\lambda^3).
\] Here \(F_{\rm tree}=\lambda\), whereas the Minkowski tree amplitude
is \(\mathcal M_{\rm tree}=-\lambda\). The sign difference is
conventional, not a physical inconsistency. General kinematics requires
three separate channels; \(s=t=u\) cannot be imposed arbitrarily.

\phantomsection\label{p-qren-02}{}

\textbf{P-QREN-02: Subtracted bubble.} Use
\(1/(AB)=\int_0^1dx/[xA+(1-x)B]^2\). After shifting the loop variable,
the denominator is \((\ell^2+m^2+x(1-x)Q^2)^2\). Subtract the \(\mu\)
configuration with the same regulator. Divergent constants cancel, and
radial integration gives the logarithm above. Each of the three channels
has factor \(1/2\). \(\square\)

At \((\lambda,m,\mu,Q)=(0.4,1,2,5)\),

\[
B_R=-0.00668304090527,\qquad F=0.401603929817.
\]

These are one-loop truncated values, not exact all-order results.

For a Minkowski channel \(z\), analytic continuation gives \[
B_R(z;z_0)=\frac1{16\pi^2}\int_0^1dx
\log\frac{m^2-x(1-x)z_0-i0}{m^2-x(1-x)z-i0},
\quad z_0=-\mu^2,
\] \[
\mathcal M(s,t,u)=-\lambda+
\frac{\lambda^2}{2}\bigl[B_R(s)+B_R(t)+B_R(u)\bigr]+O(\lambda^3).
\] \phantomsection\label{p-qren-03}{}

\textbf{P-QREN-03: Two-particle cut.} For \(s>4m^2\), the interval where
the logarithm's denominator is negative has length
\(\beta_s=\sqrt{1-4m^2/s}\), so \[
\operatorname{Im}\mathcal M^{(1)}=\frac{\lambda^2\beta_s}{32\pi}.
\] Meanwhile \(\int d\Phi_2=\beta_s/(8\pi)\), with division by \(2!\)
for identical scalars. The optical theorem's right side is
\(\lambda^2\beta_s/(16\pi)\), exactly twice the imaginary part.
\(\square\)

\section{25.3 Running and critical points}

\textbf{K-QREN-01.} In the massless/UV MS scheme for a four-dimensional
real scalar, \[
\beta_\lambda=\mu\frac{d\lambda}{d\mu}
=\frac{3\lambda^2}{16\pi^2}+O(\lambda^3).
\] The positive sign means growth toward high energy. A massive
momentum-subtraction scheme has a threshold function \[
\beta_{\rm MOM}=\frac{3\lambda^2}{16\pi^2}\,
T(\mu/m)+O(\lambda^3),\quad
T(r)=\int_0^1\frac{x(1-x)r^2}{1+x(1-x)r^2}\,dx.
\] It obeys \(T(r)\sim r^2/6\) in the infrared and tends to one in the
UV. A scheme's beta function at all scales cannot be transferred
unchanged to another scheme.

For a connected renormalized \(n\)-point function, let
\(\phi_0=Z_\phi^{1/2}\phi_R\) and
\(\gamma_\phi=\tfrac12\mu d\log Z_\phi/d\mu\). Holding bare quantities
fixed gives \[
[\mu\partial_\mu+\beta_\lambda\partial_\lambda+
\beta_{m^2}\partial_{m^2}+n\gamma_\phi]G_R^{(n)}=0
\] The anomalous-field-dimension term has the opposite sign for 1PI
vertices.

In the \(d=4-\epsilon\) expansion,
\(\beta_\lambda=-\epsilon\lambda+3\lambda^2/(16\pi^2)+\cdots\), with
nontrivial fixed point \(\lambda_*=16\pi^2\epsilon/3+O(\epsilon^2)\).
This is controlled for small \(\epsilon\). Using \(\epsilon=1\) requires
reported truncation and resummation uncertainties.

\section{25.4 Scheme matching, expansion parameters, and independent RG
checks}

Equality of leading beta coefficients has its own conditions. Suppose
two massless, mass-independent schemes differ only by an explicitly
scale-independent finite redefinition \[
\lambda'=\lambda+c\lambda^2+O(\lambda^3),
\] The chain rule gives \[
\beta_{\lambda'}=(1+2c\lambda+\cdots)\beta_\lambda
=b\lambda'^2+O(\lambda'^3)
\] when \(\beta_\lambda=b\lambda^2+\cdots\). If \(c\) depends on
\(\mu/m\), the explicit \(\mu\partial_\mu c\) adds a term at the same
order. This explains why a massive MOM threshold need not equal the MS
beta function at all scales. Physical comparisons must convert
couplings, masses, and external-leg normalizations together and fix
matching conditions.

For the Wilson--Fisher expansion, write \[
g=\frac{\lambda}{16\pi^2},\qquad
\beta_g=-\epsilon g+3g^2+O(g^3)
\] which displays the loop expansion parameter directly:
\(g_*=\epsilon/3+O(\epsilon^2)\). One must also tune to the critical
mass surface. The derivative \(\beta_g'(g_*)=\epsilon+O(\epsilon^2)\)
refers to the chosen direction of the RG parameter; stability cannot be
interpreted without specifying UV versus IR direction. Application at
\(\epsilon=1\) still needs higher-order and resummation analysis, not an
error-free three-dimensional number inferred from a small-\(\epsilon\)
proof.

Chapter 26's thin-shell integration toward the IR gives \[
\int_{\rm shell}\frac{d^4q}{(2\pi)^4q^4}=\frac{ds}{8\pi^2},
\qquad\delta\lambda=-\frac32\lambda^2\int_{\rm shell}\frac{d^4q}{(2\pi)^4q^4}.
\] This follows from three-channel combinatorics and an independent
radial integral. Differentiating the subtracted bubble of §25.2 with
respect to \(\log\mu\) obtains the corresponding UV coefficient by
another calculation. They agree in matched massless schemes with
opposite flow directions. The historical regression used independent
differentiation/integration and order-scaling checks, rather than
duplicating one known beta formula on both sides. Agreement at first
order does not identify the full Wilsonian action flow with every
nonperturbative datum of Callan--Symanzik evolution.

\chapter{26. Wilsonian RG, effective field theory, and bridges to
Core}\label{ch-26}

Low-energy observations need not retain a heavy field explicitly, but
its effects must enter effective coefficients. Schur complements and
scattering matching demonstrate this process before distinguishing the
data preserved by Wilsonian RG, process coarse-graining, and Core object
coarse-graining.

\section{26.1 Integrating out fast modes}

With a Euclidean regulator, split \(\phi=\phi_<+\phi_>\) and define, up
to normalization, \[
e^{-S_k[\phi_<]}=\int\mathcal D\phi_>\,e^{-S_\Lambda[\phi_<+\phi_>]}
\] The resulting \(S_k\) generally contains infinitely many allowed
operators. After rescaling, organize them as relevant, marginal, or
irrelevant. ``Irrelevant'' describes scaling near a specified fixed
point, not dispensability at every requested accuracy.

With IR-directed \(s=\log(\Lambda/k)\), a four-dimensional thin shell
gives \(\int_{\rm shell}d^4q/(2\pi)^4q^{-4}=ds/(8\pi^2)\). In a matched
massless normalization scheme,
\(d\lambda/ds=-3\lambda^2/(16\pi^2)+\cdots\), opposite to the
UV-directed parameter of Chapter 25.

\section{26.2 A fully checkable heavy-field matching example}

\phantomsection\label{p-qeft-01}{}

\textbf{P-QEFT-01: Quadratic Schur complement.} Eliminating the heavy
field from the two-real-field momentum kernel \[
K(p)=\begin{pmatrix}p^2-m^2&-\kappa\\-\kappa&p^2-M^2\end{pmatrix}
\] gives \[
K_{\rm eff}(p)=p^2-m^2+\frac{\kappa^2}{M^2-p^2}.
\] \textbf{Proof.} The heavy-field equation is
\(\chi=(p^2-M^2)^{-1}\kappa\phi\). Substituting into the quadratic form
yields the Schur complement. \(\square\)

For \(|p^2|<M^2\), \[
\frac{\kappa^2}{M^2-p^2}
=\frac{\kappa^2}{M^2}\sum_{j=0}^{n}(p^2/M^2)^j+R_n,\quad
|R_n|\le\frac{\kappa^2}{M^2}
\frac{|p^2/M^2|^{n+1}}{1-|p^2/M^2|}.
\] The parameter \(\kappa\) has mass dimension two; it is not
dimensionless. For \(m=0.7,M=10,\kappa=3,p^2=4\), the exact correction
is 0.09375; LO, NLO, and NNLO give 0.09, 0.0936, and 0.093744. Positive
definiteness of the original quadratic mass matrix also requires
\(m^2M^2>\kappa^2\). The derivative expansion fails near the heavy pole.

A general EFT has \[
\mathcal L_{\rm EFT}=\sum_{d,i}\frac{C_i^{(d)}(\mu)}
{\Lambda_*^{d-4}}\mathcal O_i^{(d)}.
\] The operator basis, redundancy relations, renormalization scheme,
matching scale, power counting, and fit data are all part of the model
specification.

\subsection{A scattering matching test at the same heavy scale}

The quadratic mixing coefficient \(\kappa\) has mass dimension two. A
separate, discriminating benchmark uses heavy exchange through
\(-g\chi\phi^2/2\), where \(g\) has mass dimension one. A locally stable
vacuum and the required complete potential must be declared. The
calculation below includes only tree-level heavy-\(\chi\) exchange and
does not prove global stability of a complete potential.

Standard vertices and propagators give \[
\mathcal M_\chi(s,t,u)=g^2\left[\frac1{M^2-s}+\frac1{M^2-t}+\frac1{M^2-u}\right].
\] When every \(|z|/M^2\le r<1\), channelwise expansion gives \[
\mathcal M_\chi=
\frac{3g^2}{M^2}+\frac{g^2(s+t+u)}{M^4}
+\frac{g^2(s^2+t^2+u^2)}{M^6}+\cdots.
\] Truncating each channel after \(z^n\) bounds the absolute total
remainder by \[
\frac{3g^2}{M^2}\frac{r^{n+1}}{1-r}.
\] For identical external particles of mass \(m\), on-shell kinematics
imposes \(s+t+u=4m^2\). Match lower-dimensional local and derivative
terms under this constraint rather than treating the three invariants as
arbitrary independent values. For \(M=10,g=2,m=1,s=9,t=u=-2.5\), the
exchange amplitude is approximately \(0.12200482444\). This compares a
full tree diagram with its EFT series, independently of the preceding
Schur complement, which changes only a two-point kernel.

Removing a heavy field from the low-energy active-field list requires
matching Wilson coefficients and running them in the appropriate theory.
Smooth threshold suppression of a full-theory MOM coupling is not
automatic EFT matching, nor does it make every heavy-field effect
exactly zero.

\section{26.3 Three types of RG}

\begin{longtable}[]{@{}
  >{\raggedright\arraybackslash}p{(\columnwidth - 4\tabcolsep) * \real{0.2500}}
  >{\raggedright\arraybackslash}p{(\columnwidth - 4\tabcolsep) * \real{0.3400}}
  >{\raggedright\arraybackslash}p{(\columnwidth - 4\tabcolsep) * \real{0.4100}}@{}}
\toprule\noalign{}
\begin{minipage}[b]{\linewidth}\raggedright
Transformation
\end{minipage} & \begin{minipage}[b]{\linewidth}\raggedright
Input and output
\end{minipage} & \begin{minipage}[b]{\linewidth}\raggedright
Preserved content
\end{minipage} \\
\midrule\noalign{}
\endhead
\bottomrule\noalign{}
\endlastfoot
Core object coarse-graining & Maps between declared protocol responses
and predictive factors & Responses, or bounded errors, for specified
future tasks \\
Quantum-process coarse-graining & Multitime instrument slots, initial
states, and memory processes & Record probabilities under allowed
interventions \\
Wilsonian RG & Regulated actions, fields, and couplings & Low-energy
correlation functions or matched amplitudes at stated accuracy \\
\end{longtable}

Bridges can connect them, but their domains, codomains, and preserved
data differ. Fixed points, spectral gaps, Markov order, and object
minimality are concepts belonging to their respective settings. No
general theorem identifies them unconditionally one to one.

\phantomsection\label{p-qeft-02}{}

\textbf{P-QEFT-02: A restricted real-time bridge.} Fix the same
system--environment partition, initial state, joint unitary evolution,
allowed slot interventions, and terminal readout. Direct partial tracing
and integration of the environment into a Schwinger--Keldysh influence
functional give the same protocol probabilities when both calculations
are exact and the required sums, integrals, and limits are legitimate.

\textbf{Proof.} In a finite-dimensional, finitely regulated realization,
the probability of a complete record \(\mathbf o\) is a linear pairing
of joint operators: insert successive joint unitaries and system CP
branches into the same initial state in time order, then trace over all
final states. Expanding matrix elements puts the amplitude and its
conjugate on the two closed-time-path legs. The initial environmental
state and final trace supply the same environmental index contractions.
Contracting the environment first gives the influence functional;
contracting system indices afterwards or first composing the joint
operator and then tracing are merely different orders of the same finite
sum. Passing this equality to a controlled continuum limit requires
additional convergence and interchange conditions. \(\square\)

The common responses can then be quotiented by Core predictive
equivalence. Approximate influence functionals, memory truncation, or
discarded interventions require probability-error bounds; exact equality
no longer applies. Euclidean fast-mode integration is not itself general
CPTP time evolution. Initial correlations, external insertions, gauge
constraints, and nonlocal memory can all defeat a naive channel
interpretation. A shared coarse-graining intuition does not replace
these interfaces.

\chapter{27. QED: gauge structure, BRST, and Ward
identities}\label{ch-27}

Gauge redundancy must remain consistent in quantum calculations, or
longitudinal terms and counterterms can produce spurious results. The
Ward--Takahashi identity connects gauge fixing, BRST, and one-loop UV
checks here, with the complete theorem distinguished from evidence at a
particular perturbative order.

\section{27.1 Action and symmetry}

\[
D_\mu=\partial_\mu+ieA_\mu,\quad
\mathcal L=-\tfrac14F^2+\bar\psi(i\not D-m)\psi,
\quad \psi\mapsto e^{-ie\alpha}\psi,\quad A_\mu\mapsto A_\mu+\partial_\mu\alpha.
\] The interaction is \(-e\bar\psi\gamma^\mu A_\mu\psi\), with vertex
\(-ie\gamma^\mu\). Linear covariant gauge fixing adds
\(-(\partial A)^2/(2\xi)\), giving free photon propagator \[
D_{\mu\nu}(q)=\frac{-i}{q^2+i0}
\left[\eta_{\mu\nu}-(1-\xi)\frac{q_\mu q_\nu}{q^2+i0}\right].
\] The local nonzero modes of the Abelian Faddeev--Popov determinant are
independent of \(A\), so the corresponding ghosts decouple in ordinary
perturbative QED. Global zero modes and boundaries require separate
treatment.

With auxiliary field \(B\), take
\(sA=\partial c,\ s\psi=-iec\psi,\ sc=0,\ s\bar c=B,\ sB=0\).
Gauge-fixing and ghost terms can be written
\(s[\bar c(\partial A+\xi B/2)]\). Eliminating \(B\) gives the stated
gauge-fixing term. The BRST quotient constructs a physical sector under
declared quantum-consistency assumptions; it does not treat arbitrary
indefinite-metric states as probability states.

\section{27.2 The Ward--Takahashi identity}

\textbf{K-QWARD-01.} With the overall charge factored out of the full
vertex \(\Gamma^\mu\), a regulator and renormalization prescription
preserving vector gauge symmetry give \[
q_\mu\Gamma^\mu(p+q,p)=S^{-1}(p+q)-S^{-1}(p).
\] The proof mechanism is an anomaly-free local phase change of
variables in the sourced path integral, differentiation with respect to
fermion sources, and external-leg amputation. Contact terms produce the
two inverse propagators. We invoke the standard complete identity
explicitly; a tree-level check is not an all-loop proof.

At tree level it is simply
\(q_\mu\gamma^\mu=\not p+\not q-m-(\not p-m)\). Vacuum polarization
obeys \(q_\mu\Pi^{\mu\nu}=0\). The identity constrains longitudinal
structure, not every finite transverse vertex form factor.

For \(d=4-2\epsilon_{\rm DR}\), \[
\psi_0=Z_2^{1/2}\psi,\quad A_0=Z_3^{1/2}A,\quad
e_0=\mu^{\epsilon_{\rm DR}}Z_e e,\quad
Z_1=Z_eZ_2Z_3^{1/2}.
\] The Ward identity gives \(Z_1=Z_2\) and therefore \(Z_e=Z_3^{-1/2}\).

\section{27.3 A one-loop UV consistency benchmark}

Fix Feynman gauge, one unit-charge Dirac fermion, and the one-loop
MS-bar UV-pole convention \[
a=\frac{\alpha}{4\pi\bar\epsilon},\qquad
\frac1{\bar\epsilon}=\frac1{\epsilon_{\rm DR}}-\gamma_E+\log4\pi.
\] The purpose is an independent consistency calculation of
counterterms, not all finite self-energy and vertex parts in an
arbitrary gauge.

First simplify the fermion self-energy numerator with
\(\gamma_\mu\gamma_\alpha\gamma^\mu=(2-d)\gamma_\alpha\). After Feynman
parameterization and the loop shift, the UV coefficient is \[
\Sigma_{\rm UV}(p)=a\int_0^1(-2x\not p+4m)dx
=a(-\not p+4m).
\] For inverse-propagator convention \[
\Gamma^{(2)}=\not p-m+\delta Z_2\not p
-\delta m_{\rm direct}\,m-\Sigma_{\rm loop},
\] cancelling the separate \(\not p\) and \(m\) UV residuals gives \[
\delta Z_2=-a,\qquad\delta m_{\rm direct}=-4a.
\] Multiplicative mass renormalization, \(m_0=Z_m m\), satisfies \[
1+\delta m_{\rm direct}=Z_2Z_m,
\qquad\delta Z_m=\delta m_{\rm direct}-\delta Z_2=-3a.
\] The direct mass coefficient is not \(Z_m-1\). Confusing them creates
a spurious Ward-identity conflict.

There is an independent check of the vertex's logarithmic UV part.
Symmetric loop-momentum tensor integration contracts with the metric,
and on the four-dimensional pole residue, \[
\frac14\gamma_\alpha\gamma_\rho\gamma^\mu
\gamma^\rho\gamma^\alpha=\gamma^\mu.
\] Apply the four-dimensional identity
\(\gamma_\rho\gamma^\mu\gamma^\rho=-2\gamma^\mu\) twice. Terms of order
\(d-4\) can affect finite parts, so this residue cannot replace the full
dimension-dependent vertex. It gives
\(\Lambda^\mu_{\rm UV}=a\gamma^\mu\), hence \(\delta Z_1=-a\).
Comparison with \(\delta Z_2\) from the independent self-energy
calculation confirms \(\delta Z_1=\delta Z_2\) in this gauge and at this
order.

The photon UV coefficient is \[
\frac{2\alpha}{\pi\bar\epsilon}\int_0^1x(1-x)dx
=\frac{\alpha}{3\pi\bar\epsilon},
\] so \(\delta Z_3=-\alpha/(3\pi\bar\epsilon)\). The full Ward relation
\(Z_e=Z_3^{-1/2}\) then gives the charge pole. The historical regression
separately integrated self-energy coefficients, contracted the vertex
Dirac tensor, and checked the photon coefficient. Assuming \(Z_1=Z_2\)
and then comparing them would not be an independent loop calculation.

This does not cover all finite parts at arbitrary \(\xi\), infrared
structure of on-shell \(Z_2\), or precision multiloop QED. The
tree-level Ward check, this one-loop UV check, and the complete
Ward--Takahashi theorem invoked in §27.2 are three distinct forms of
evidence.

\chapter{28. Computable QED applications: scattering, running, and
magnetic moments}\label{ch-28}

The preceding rules now calculate three quantities: a cross section, an
effective charge, and an anomalous magnetic moment. These reproduce
standard QED benchmarks; each number retains its mass approximation,
perturbative order, units, and constant version and does not thereby
become a new UEOT prediction.

\section{28.1 Electron--positron annihilation}

For \(e^-e^+\to\mu^-\mu^+\), neglect both lepton masses and retain only
one-photon tree exchange: \[
\mathcal M=\frac{e^2}{s}
[\bar v(p_2)\gamma^\mu u(p_1)]
[\bar u(k_1)\gamma_\mu v(k_2)].
\] \phantomsection\label{p-qqed-01}{}

\textbf{P-QQED-01.} Spin completeness and \[
\operatorname{Tr}(\not a\gamma^\mu\not b\gamma^\nu)
=4(a^\mu b^\nu+a^\nu b^\mu-\eta^{\mu\nu}a\cdot b)
\] give, after dividing by four for the initial-state average, \[
\overline{|\mathcal M|^2}=\frac{2e^4}{s^2}(t^2+u^2)
=e^4(1+\cos^2\theta).
\] Insert the two-body phase space and integrate: \[
\frac{d\sigma}{d\Omega}=\frac{\alpha^2}{4s}(1+\cos^2\theta),
\qquad \sigma=\frac{4\pi\alpha^2}{3s}.
\] The angular integral is \(16\pi/3\), directly checking the
coefficient. \(\square\)

The longitudinal propagator term vanishes by external-leg Dirac
equations and current conservation, so the tree cross section is
independent of \(\xi\). Finite masses, Z exchange, radiative
corrections, and detector acceptance are outside this benchmark. With
frozen reproduction input \(\alpha^{-1}=137.035999084\) and
\(\sqrt s=10\) energy units, \(\sigma=2.2305875\times10^{-6}\)
inverse-energy-squared units. This does not claim the latest constant
determination or a precision experimental prediction.

The fixed inverse fine-structure constant is the central value from
CODATA 2018, retained solely to preserve the reproducibility of this
benchmark.
\href{https://physics.nist.gov/cuu/Constants/archive2018.html}{NIST 2018
archive}, \href{https://doi.org/10.1103/RevModPhys.93.025010}{original
adjustment report}.

\section{28.2 Vacuum polarization and decoupling}

For spacelike \(Q^2>0\), define the positive screening function \[
\widehat\Pi(Q^2)=\frac{2\alpha}{\pi}\int_0^1dx\,x(1-x)
\log\!\left[1+\frac{Q^2}{m^2}x(1-x)\right],\quad
\alpha_{\rm eff}=\frac{\alpha}{1-\widehat\Pi}.
\] This fixes a subtraction and sign convention. For small
\(a_Q=Q^2/m^2\), \(\widehat\Pi=\alpha a_Q/(15\pi)+O(a_Q^2)\); for large
\(a_Q\), \[
\widehat\Pi=\frac{\alpha}{3\pi}\left(\log a_Q-\frac53\right)+O(a_Q^{-1}).
\] The first uses \(\int x^2(1-x)^2dx=1/30\); the second uses
\(\int x(1-x)dx=1/6\) and its logarithmic moment.

\phantomsection\label{p-qqed-02}{}

\textbf{P-QQED-02: Threshold beta function.} Differentiating with
respect to \(\log Q\), through \(O(\alpha^2)\), gives \[
\frac{d\alpha_{\rm eff}}{d\log Q}
=\frac{4\alpha^2}{\pi}\int_0^1dx\,
\frac{a_Qx^2(1-x)^2}{1+a_Qx(1-x)}.
\] It vanishes in the infrared and tends to \(2\alpha^2/(3\pi)\) in the
UV. For several active Dirac fields, multiply each contribution by its
\(N_cQ_f^2\) and sum, treating thresholds with the corresponding masses.
\(\square\) The one-loop Landau pole warns against extrapolation; it is
not a proved physical energy singularity.

\section{28.3 Schwinger's magnetic moment and infrared boundaries}

Decompose the on-shell vertex as \[
\bar u(p')\Gamma^\mu u(p)=\bar u(p')
\left[\gamma^\mu F_1(q^2)
+\frac{i\sigma^{\mu\nu}q_\nu}{2m}F_2(q^2)\right]u(p),
\quad a_\ell=F_2(0).
\] The standard one-loop starting point, with external legs on shell,
overall \(-ie\) factored out, and fixed momentum flow, is for example \[
\Lambda^\mu(p',p)=-ie^2\int\frac{d^d\ell}{(2\pi)^d}
\frac{\gamma_\rho(\not p'-\not\ell+m)\gamma^\mu
(\not p-\not\ell+m)\gamma^\rho}
{[\ell^2+i0][(p'-\ell)^2-m^2+i0][(p-\ell)^2-m^2+i0]}.
\] Dimensional regularization also retains \(\mu^{2\epsilon_{\rm DR}}\),
with consistent infrared treatment where needed. Feynman
parameterization, a loop shift, Gordon decomposition, and projection
onto the Pauli structure reduce the finite \(F_2\) contribution as
\(q^2\to0\) to the integral below. This analytical reduction is a
standard QED input. The subsequent two-dimensional numerical integral
does not machine-prove every interchange, regularization step, and limit
in the reduction.

\textbf{K-QQED-01.} Standard one-loop vertex parameterization and Pauli
projection give \[
F_2(0)=\frac{\alpha}{\pi}
\int_0^1dz\int_0^{1-z}dy\,\frac{z}{1-z}
=\frac{\alpha}{2\pi}.
\] Integrating over \(y\) cancels \(1-z\), after which the integral over
\(z\) gives the result. The independent checks here concern this
parameter integral and its numerical value, not a machine derivation of
all regularization and on-shell limits from the full loop graph.

Virtual photon loops and real soft radiation can separately diverge in
the infrared. Physical inclusive observables require fixed
energy/angular resolution and combined real and virtual terms.
Cancellation of a vector gauge anomaly does not eliminate axial-current
anomalies. The QED sector has not supplied a new independent UEOT
prediction of physical constants.

\chapter{29. Non-Abelian gauge fields and a candidate Standard Model
structure}\label{ch-29}

Gauge anomalies strongly constrain candidate particle tables but
generally leave several consistent solutions. After fixing non-Abelian
conventions, we calculate representation and hypercharge conditions and
distinguish evidence for local cancellation, global consistency, and
nature's actual choice.

\section{29.1 Yang--Mills definitions and calculation responsibilities}

Use Hermitian generators \([T^a,T^b]=if^{abc}T^c\) and
\(\operatorname{Tr}_R(T^aT^b)=T(R)\delta^{ab}\). Consistently with QED's
\(D=\partial+igA\), define \[
F_{\mu\nu}=\partial_\mu A_\nu-\partial_\nu A_\mu+ig[A_\mu,A_\nu],
\quad [D_\mu,D_\nu]=igF_{\mu\nu},
\quad\mathcal L_{\rm YM}=-\tfrac12\operatorname{Tr}(F_{\mu\nu}F^{\mu\nu})
\] with fundamental trace normalization \(T=1/2\). The nonlinear
component term is \(-g f^{abc}A_\mu^bA_\nu^c\). Changing to
\(D=\partial-igA\) requires the corresponding sign changes throughout.

Non-Abelian self-coupling gives three- and four-gauge-boson vertices.
Ghosts no longer decouple. Gauge fixing and BRST transformations use
covariant derivatives and Lie brackets; nilpotency depends on Jacobi.
Quantum constraints are Slavnov--Taylor identities, not QED's
\(Z_1=Z_2\) applied to every vertex. Perturbative BRST does not solve
the global Gribov problem.

\textbf{K-QYM-01.} For Dirac fermions and complex scalars, \[
\beta(g)=-\frac{g^3}{16\pi^2}b_0+O(g^5),\quad
b_0=\frac{11}{3}C_A-\frac43\sum_{\rm Dirac}T(R_f)
-\frac13\sum_{\rm complex}T(R_s).
\] A Weyl fermion contributes half the Dirac value. Six-flavour QCD in
the fundamental representation gives \(b_0=7\). Asymptotic freedom
proves neither confinement nor a mass gap. Strong-coupling low-energy
predictions generally need lattice methods, effective theories, or other
nonperturbative techniques with error control.

\section{29.2 The Standard Model as a concrete candidate}

Declare the gauge Lie algebra
\(\mathfrak{su}(3)_C\oplus\mathfrak{su}(2)_L\oplus\mathfrak u(1)_Y\).
The global group may have a representation-compatible discrete quotient.
The local Lagrangian does not determine line operators or every global
sector. The table uses exclusively left-handed Weyl fields and
\(Q=T_3+Y\).

\begin{longtable}[]{@{}
  >{\raggedright\arraybackslash}p{(\columnwidth - 6\tabcolsep) * \real{0.1800}}
  >{\raggedright\arraybackslash}p{(\columnwidth - 6\tabcolsep) * \real{0.3700}}
  >{\raggedleft\arraybackslash}p{(\columnwidth - 6\tabcolsep) * \real{0.1300}}
  >{\raggedleft\arraybackslash}p{(\columnwidth - 6\tabcolsep) * \real{0.3200}}@{}}
\toprule\noalign{}
\begin{minipage}[b]{\linewidth}\raggedright
Field
\end{minipage} & \begin{minipage}[b]{\linewidth}\raggedright
\(SU(3)_C,SU(2)_L\)
\end{minipage} & \begin{minipage}[b]{\linewidth}\raggedleft
\(Y\)
\end{minipage} & \begin{minipage}[b]{\linewidth}\raggedleft
Weyl components per generation
\end{minipage} \\
\midrule\noalign{}
\endhead
\bottomrule\noalign{}
\endlastfoot
\(Q\) & \(3,2\) & \(1/6\) & 6 \\
\(U^c\) & \(\bar3,1\) & \(-2/3\) & 3 \\
\(D^c\) & \(\bar3,1\) & \(1/3\) & 3 \\
\(L\) & \(1,2\) & \(-1/2\) & 2 \\
\(E^c\) & \(1,1\) & 1 & 1 \\
\(H\) & \(1,2\) & \(1/2\) & Complex scalar, not Weyl \\
\end{longtable}

Gauge and fermion kinetic terms, Higgs kinetic and potential terms, and
Yukawa couplings form the usual operator-dimension-at-most-four
renormalizable structure; topological theta terms may also occur.
``Minimal object'' does not automatically set the strong-CP parameter to
zero. This field table is a sector input, not the unique particle table
derived from Core.

\section{29.3 Exact anomaly equations}

For general \(N_c\) and the same multiplet types, local anomalies
require \[
2y_Q+y_U+y_D=0,\quad N_cy_Q+y_L=0,
\] \[
2N_cy_Q+N_cy_U+N_cy_D+2y_L+y_E=0,
\] \[
2N_cy_Q^3+N_cy_U^3+N_cy_D^3+2y_L^3+y_E^3=0.
\] The pure-colour cubic anomaly cancels between the doublet's two
fundamentals and two antifundamentals. \(SU(2)\) has no perturbative
cubic anomaly, but its Witten global anomaly requires an even total
number of fundamental doublets, here \(N_f(N_c+1)\).

\phantomsection\label{p-qsm-01}{}

\textbf{P-QSM-01: Anomaly solutions and nonuniqueness.} The three linear
equations give \(y_L=-N_cy_Q,\ y_D=-2y_Q-y_U,\ y_E=2N_cy_Q\). The cubic
equation factors as \[
6N_cy_Q[(N_c-1)y_Q-y_U][(N_c+1)y_Q+y_U].
\] Expansion or substitution verifies this directly. Besides the SM-type
branch and its up/down interchange, a degenerate \(y_Q=0\) branch
remains. Anomaly cancellation alone does not uniquely select
hypercharge. \(\square\)

Ordinary Yukawa gauge invariance of
\(QHU^c,QH^\dagger D^c,LH^\dagger E^c\) gives \(y_H=N_cy_Q\). Normalize
\(y_H=1/2\) to obtain \[
Y_Q=\frac1{2N_c},\quad
Y_U=-\frac{N_c+1}{2N_c},\quad
Y_D=\frac{N_c-1}{2N_c},\quad Y_L=-\frac12,\quad Y_E=1.
\] Any \(N_c\) satisfying the other conditions admits this local
solution branch; three has not been selected. Adding \(N^c\) can reopen
Abelian mixing directions, requiring all conditions to be solved again.

\section{29.4 Local and global anomalies and identity transport}

Vanishing perturbative anomaly coefficients are only the first layer.
The global gauge group, spin structure, discrete symmetries, background
bundles, boundaries, and possible inflow affect the complete anomaly
problem.

\textbf{S-QGOIT: Gauge-Orbit Identity Transport.} Once a physical
realization identifies a gauge transformation as representational
redundancy, complete physical responses must descend consistently to the
orbit quotient. An uncancelled curvature or holonomy obstruction in the
chiral-fermion determinant/Pfaffian line prevents the isolated candidate
from being a consistent realization of that dynamical gauge theory.
Legitimate additional matter, counterterms, or inflow may complete it.

This locates standard anomaly obstructions within UEOT identity
transport; it does not identify arbitrary Core holonomy with a gauge
anomaly. A 't Hooft anomaly of an ungauged global symmetry is not the
same as inconsistent dynamical gauge redundancy. Mathematical sources
include \href{https://arxiv.org/abs/1607.01873}{Yonekura's
determinant-line discussion} and the
\href{https://arxiv.org/abs/1808.00009}{Dai--Freed particle-physics
analysis}.

\chapter{30. Higgs, Yukawa, flavour, and neutrino
structures}\label{ch-30}

The Higgs mechanism explains how mass terms arise in a chosen model
without automatically fixing Yukawa values. Explicit mass matrices lead
to flavour mixing and neutrino realizations, separating conditional
relations, fitted parameters, and the mass-hierarchy questions that
still require a mechanism.

\section{30.1 An explicit realization of electroweak breaking}

For a single complex doublet, take \[
\begin{gathered}
V(H)=-\mu_H^2H^\dagger H+\lambda_H(H^\dagger H)^2,
\\ \lambda_H>0,\\ \mu_H^2>0,\\
H=\frac1{\sqrt2}\binom{0}{v+h},\\ v^2=\mu_H^2/\lambda_H.
\end{gathered}
\] This is unitary-gauge parameterization near a particular vacuum, not
a proof of global gauge fixing.

\phantomsection\label{p-qhig-01}{}

\textbf{P-QHIG-01: Mass matrices.} Expanding
\((D_\mu H)^\dagger D^\mu H\) gives \[
m_W^2=\frac{g^2v^2}{4},\qquad
M_{\rm neutral}^2=\frac{v^2}{4}
\begin{pmatrix}g^2&-gg'\\-gg'&g'^2\end{pmatrix}
\] in a compatible \(W^3,B\) neutral-field basis. Its determinant is
zero, and its other eigenvalue is \((g^2+g'^2)v^2/4\). Define \[
\tan\theta_W=g'/g,\quad
A=\sin\theta_W W^3+\cos\theta_W B,\quad
Z=\cos\theta_W W^3-\sin\theta_W B,
\] Then \(m_A=0,m_Z=v\sqrt{g^2+g'^2}/2\) and
\(e=g\sin\theta_W=g'\cos\theta_W\). The quadratic potential expansion
gives \(m_h^2=2\lambda_Hv^2\). At tree level,
\(\rho=m_W^2/(m_Z^2\cos^2\theta_W)=1\). \(\square\)

Several suitable doublets can also preserve tree-level \(\rho=1\), so
the relation does not uniquely select one Higgs. A ``topologically
trivial scalar'' does not determine a doublet, potential, or vacuum
scale either.

\section{30.2 Free flavour parameters and CP}

In four-component notation, \[
\mathcal L_Y=-\bar Q_LY_dd_RH-\bar Q_LY_uu_R\widetilde H
-\bar L_LY_ee_RH+\mathrm{h.c.},\quad \widetilde H=i\sigma_2H^*.
\] After symmetry breaking, \(M_f=vY_f/\sqrt2\). Singular-value
decomposition gives \[
U_{fL}^\dagger M_fU_{fR}=D_f\ge0,\qquad
V_{\rm CKM}=U_{uL}^\dagger U_{dL}.
\] Flavour diagonality of neutral currents, charged-current mixing, and
GIM cancellation follow from this structure. Degenerate masses enlarge
the residual basis-change group.

\phantomsection\label{p-qflav-01}{}

\textbf{P-QFLAV-01: Generic nondegenerate \(N\)-generation quark
counting.} Two complex Yukawa matrices contain \(4N^2\) real parameters.
Basis transformations \(U(N)_Q\times U(N)_u\times U(N)_d\) remove
\(3N^2-1\), since a common baryon phase leaves Yukawa matrices
unchanged. The remaining \(N^2+1\) consist of \[
2N\text{ masses}+\frac{N(N-1)}2\text{ angles}
+\frac{(N-1)(N-2)}2\text{ CP phases}.
\] Three generations give six masses, three angles, and one phase.
\(\square\)

This counts Yukawa flavour parameters only; the strong-CP parameter is
not among those ten. Nonzero KM-type CP violation requires at least
three generations, not exactly three. Physical CP violation also
requires a suitable nondegenerate structure and a nonzero rephasing
invariant, such as appropriate combinations of
\(J=\operatorname{Im}(V_{ij}V_{kl}V^*_{il}V^*_{kj})\).

\section{30.3 Neutrinos: distinct realizations}

Adding \(N_R\) permits Dirac masses. A Majorana EFT may use the explicit
convention \[
\mathcal L_5=-\frac{C_{ij}}{2\Lambda_*}(L_iH)(L_jH)+\mathrm{h.c.},
\qquad m_\nu=\frac{v^2}{2\Lambda_*}C
\] with spin and \(SU(2)\) contractions implicit and \(C=C^T\). The
factor of two depends on the operator definition and must agree with the
mass formula.

\phantomsection\label{p-qnu-01}{}

\textbf{P-QNU-01: Type-I seesaw.} With invertible \(M_R\) and
heavy--light separation \(\|M_D M_R^{-1}\|\ll1\), \[
\mathcal M=\begin{pmatrix}0&M_D\\M_D^T&M_R\end{pmatrix},
\quad m_\nu^{\rm eff}=-M_DM_R^{-1}M_D^T
\] is the Schur complement of the zero-momentum quadratic kernel.
Dynamics and mixing matrices have additional higher-order corrections.
Physical masses of a complex symmetric mass matrix come from Takagi
factorization, not ordinary eigenvalues of a general complex matrix.
\(\square\)

Ultrarelativistic vacuum oscillations are approximated by \[
P_{\alpha\to\beta}(L)=
\left|\sum_i U_{\beta i}U_{\alpha i}^*
e^{-im_i^2L/(2E)}\right|^2,
\] under suitable preparation, coherent-propagation, and detection
approximations. Matter effects, wave-packet decoherence, and nonunitary
extensions must be added separately. Oscillation probabilities do not
uniquely determine Dirac versus Majorana character or the high-energy
seesaw realization.

\section{30.4 Flavour models and the original nesting intuition}

A discrete depth \(d=1,2,3\) becomes testable only after an independent
map \(d\mapsto Y\) is specified. For example, equal-step exponential
masses \(m(d)=Me^{\kappa d}\) require \(m_2^2=m_1m_3\). This is a
falsifiable relation, not a restatement that deeper nesting means larger
mass.

An ansatz \(Y_{ij}=c_{ij}\epsilon^{n_{ij}}\) already has
Froggatt--Nielsen structure. To add physical content, UEOT must
independently select the integers, \(\epsilon\), and coefficient
constraints, then predict flavour observables not used to select the
realization. An arbitrarily fitted relabelling is not a mass theory.

\chapter{31. Three generations and three colours: conditional theorems
and simpler premises}\label{ch-31}

Selecting exactly three generations first faces replication: conditions
closed under copying families cannot determine a unique positive count.
A specified representation assignment and weak-running budget break that
freedom here, while countermodels preserve the role of each premise;
Chapter 36 supplies the distinct high-energy-window analysis.

Reading note: this chapter retains a low-energy weak-asymptotic-freedom
premise. Chapter 36 shows that a different set of high-energy
asymptotic-freedom conditions does not automatically imply it and
supplies a five-colour countermodel. Cite the two chapters together.

We retain the valid arithmetic of the available v30.11 mainline and
strengthen it by removing the ``primitive package'' assumption. The
result remains a conditional theorem, not a three-generation theorem of
general UEOT Core.

\section{31.1 The replication-invariance obstruction}

\phantomsection\label{p-qg3-01}{}

\textbf{P-QG3-01.} Suppose an axiom class \(\mathcal A\) contains a
model \(T\) and is closed under family replication by some \(r>1\): the
replicated model satisfies every axiom and \(N_f(R_rT)=rN_f(T)\). Then
the class cannot select a unique fixed positive integer generation
count.

\textbf{Proof.} If the unique value were \(N_0>0\), the replicated model
in the same class would have \(rN_0\ne N_0\), a contradiction.
\(\square\)

Independent anomaly cancellation in each generation, general Born
kinematics, and replication-closed object certificates therefore cannot
explain three by themselves. A successful selector needs a premise that
genuinely breaks replication freedom. An upper bound, an
intergenerational anomaly, or a coupling structure could serve that
role.

\section{31.2 A fixed representation class for a horizontal gauge
parent}

\textbf{H-QG3-1 (HGP and representation assignment).} There is a
dynamical chiral horizontal or extended-weak gauge carrier \(G_H\), with
a complex representation \(R\) whose cubic anomaly coefficient obeys
\(A(R)=a\ne0,A(\bar R)=-a\). Each lepton generation is one colourless
\(R\). Quarks have colour multiplicity \(N_c\), with \(k\) generations
in \(R\) and the remaining \(N_f-k\) in \(\bar R\). This specifies more
than the existence of a parent object: it fixes its chiral
representation assignment and assumes no other uncounted contribution to
that cubic anomaly.

\phantomsection\label{p-qg3-02}{}

\textbf{P-QG3-02: All integer solutions.} \[
A_H/a=N_f+N_c(2k-N_f)=0
\iff 2N_ck=(N_c-1)N_f.
\] Let \(d=\gcd(2N_c,N_c-1)=\gcd(2,N_c-1)\). All positive solutions are
\[
k=t\frac{N_c-1}{d},\qquad N_f=t\frac{2N_c}{d},\qquad
t\in\mathbb Z_{>0}.
\] \textbf{Proof.} Divide by the greatest common divisor. The remaining
coefficients are coprime, so divisibility forces the indicated common
multiple. For odd \(N_c\), \(d=2\); for even \(N_c\), \(d=1\).
\(\square\)

A primitive package means \(t=1\). General Core atomicity does not force
the universe to realize exactly one package. The anomaly assignment has
a clear precedent in the classic 331 mechanism and is not original to
this book: see
\href{https://arxiv.org/abs/hep-ph/9206242}{Pisano--Pleitez} and the
\href{https://link.aps.org/doi/10.1103/PhysRevD.50.R34}{extension with
right-handed neutrinos}.

\section{31.3 Conditional selection without a primitive-package premise}

Add the following explicit assumptions.

\begin{itemize}
\tightlist
\item
  \textbf{H-QG3-2:} The colour group is \(SU(N_c)\) with a genuinely
  complex fundamental representation, so \(N_c\ge3\). Excluding the
  pseudoreal fundamental of \(SU(2)\) is a physical-class choice.
\item
  \textbf{H-QG3-3:} In a declared perturbative window with every
  ordinary generation active and fixed gauge structure and matter
  content, require the one-loop weak \(SU(2)_L\) coefficient to point
  toward asymptotic freedom. Ordinary fermions give the necessary
  condition \[
  N_f(N_c+1)<22.
  \] Scalars and additional ordinary weakly charged matter add
  screening. A changed gauge group or new vector structure requires
  recalculation; the bound is not unconditional.
\end{itemize}

\phantomsection\label{p-qg3-03}{}

\textbf{P-QG3-03: Selection without primitivity.} Within the
intersection of H-QG3-1, 2, and 3, the unique positive integer anomaly
solution is \[
(N_c,N_f,k)=(3,3,1).
\] No separate \(t=1\) premise is needed.

\textbf{Proof.} For odd \(N_c\ge3\), \(N_f=tN_c\), so
\(tN_c(N_c+1)<22\). For \(N_c\ge5\), the left side is at least 30,
impossible. For \(N_c=3\), \(12t<22\) forces the positive integer \(t\)
to be one, hence \(N_f=3,k=1\). For even \(N_c\ge4\), \(N_f=2tN_c\)
makes the left side at least 40, also impossible. \(\square\)

This removes a redundant premise within a fixed model class. It does not
establish HGP, uniqueness of the representation class, or the necessity
of asymptotic freedom for every weak effective theory in nature.
Relaxing complex colour admits \((2,4,1)\) under the arithmetic bound;
removing the weak bound restores \(t>1\) and larger odd colour counts.
Additional horizontal anomaly-cancelling matter changes the integer
equation itself.

Adding one Higgs doublet at the candidate point gives the complete
ordinary weak coefficient \[
b_2=\frac{22}{3}-\frac{N_f(N_c+1)}3-\frac16
=\frac{19}{6}>0.
\] This checks one permitted weak-matter configuration. It is not an
existence proof for every high-energy HGP completion.

\section{31.4 An explicit 331 realization and anomaly inventory}

One RHN-type branch of \(SU(3)_C\times SU(3)_L\times U(1)_X\) has the
all-left-handed representations

\begin{longtable}[]{@{}
  >{\raggedleft\arraybackslash}p{(\columnwidth - 2\tabcolsep) * \real{0.3000}}
  >{\raggedright\arraybackslash}p{(\columnwidth - 2\tabcolsep) * \real{0.7000}}@{}}
\toprule\noalign{}
\begin{minipage}[b]{\linewidth}\raggedleft
Multiplicity
\end{minipage} & \begin{minipage}[b]{\linewidth}\raggedright
Representation \((C,L,X)\)
\end{minipage} \\
\midrule\noalign{}
\endhead
\bottomrule\noalign{}
\endlastfoot
3 & \((1,3,-1/3)\) \\
3 & \((1,1,1)\) \\
1 & \((3,3,1/3)\) \\
2 & \((3,\bar3,0)\) \\
4 & \((\bar3,1,-2/3)\) \\
5 & \((\bar3,1,1/3)\) \\
\end{longtable}

The pure \(C^3,L^3\) anomalies, mixed \(C^2X,L^2X\), mixed gravitational
\(X\), and \(X^3\) all vanish. Fundamental and antifundamental
representations have opposite cubic non-Abelian anomaly signs but the
same quadratic Dynkin index \(1/2\). The \(U(1)\) contributions include
the dimensions of the other representations. These rules reproduce the
appendix's rational checks. The table demonstrates that a locally
anomaly-free carrier exists. It is neither evidence that nature selects
it nor a complete global-anomaly or phenomenological assessment.

\section{31.5 Another conditional branch with fixed ordinary charges}

For general \(N_c\), take \(Q=T_3+\beta_8T_8+X\) and fix ordinary
charges \(Q_\nu=0,Q_e=-1,Q_u=2/3,Q_d=-1/3\). Then \[
X_L=-\tfrac12-\frac{\beta_8}{2\sqrt3},\quad
X_{Q+}=\tfrac16-\frac{\beta_8}{2\sqrt3},\quad
X_{Q-}=\tfrac16+\frac{\beta_8}{2\sqrt3},
\] and third-component charges are \[
q_E=-\tfrac12-\frac{\sqrt3\beta_8}{2},\quad
q_{J+}=\tfrac16-\frac{\sqrt3\beta_8}{2},\quad
q_{J-}=\tfrac16+\frac{\sqrt3\beta_8}{2}.
\] A complete chiral table for arbitrary \(\beta_8\) includes the
left-handed conjugate singlets of every corresponding right-handed
field, \textbf{including charged \(E^c\)}, not just \(e^c,u^c,d^c,J^c\).
For \(q_E=0\), the neutral singlet contributes nothing to these local
anomalies.

\phantomsection\label{p-qg3-04}{}

\textbf{P-QG3-04.} For this full representation table, imposing the
\(L^3\) anomaly equation and summing directly gives \[
\begin{gathered}
A_{L^2X}=\frac{N_f(N_c-3)}{12},\\
A_{X^3}=-\frac{N_f(N_c-3)(1+\beta_8^2)}4,
\\ A_{\mathrm{grav}^2X}=A_{C^2X}=A_{C^3}=0.
\end{gathered}
\] To verify this, the triplet \(X^3\) contributions are \(3X_L^3\) and
\(3N_cX_{Q\pm}^3\). Add the cubes of the singlets' negative physical
charges, substitute \(k=N_f(N_c-1)/(2N_c)\), and expand to obtain the
factors above. Mixed terms use \(T(3)=1/2\) directly. \(\square\)

For \(N_f>0\), local cancellation within this class forces \(N_c=3\),
then \(N_f=3k\). An upper bound or another replication-breaking
principle is still needed to select \(k=1\). Replacing the ordinary
quark charges by the \(N_c\)-dependent solution of Chapter 29 removes
the mixed-anomaly argument forcing three colours. Ordinary charges and
three colours have not both been derived from nothing.

\section{31.6 Exact kernel correspondence and permanent countermodels}

The integer content of the §31.3 selector is implemented in the finite
UEOT-QM kernel as \texttt{Selector.\allowbreak selects\_\allowbreak three\_\allowbreak without\_\allowbreak atomicity}.
Its inputs explicitly include natural numbers \(N_c,N_f,k\),
\(N_c\ge3\), \(N_f>0\), \(k\le N_f\), the stated cubic equation, and the
strict weak budget. It does not generate these physical premises from
general Core.

\texttt{Spectrum.\allowbreak spectrum\_\allowbreak selects\_\allowbreak three} further connects finite
spectrum-row sums to the selector. Its contract requires homogeneous
lepton orientation, a nonzero anomaly unit, ordinary lepton and quark
doublets in every row active in the same weak window, nonnegative
additional screening, and positivity of the complete weak coefficient.
Any anomaly compensation outside the ordinary rows must separately sum
to zero. These are substantive conditions. The theorem does not
construct an \(SU(3)_L\) quantum field theory, prove an additional
compensation mechanism effective, or preserve the original window after
thresholds change.

For example, relax homogeneous leptons by placing \(l\) generations in
\(R\) and the others in \(\bar R\). The cubic equation becomes \[
(2l-N_f)+N_c(2k-N_f)=0.
\] The point \((N_c,N_f,l,k)=(3,2,1,1)\) satisfies it and has
\(N_f(N_c+1)=8<22\). This explicitly refutes the claim that spectrum
orientation is immaterial. The kernel verifies the corresponding row
counts and weak coefficient. Under a separately specified generalized
331 charge table, \(\beta_8=0\), and further conditions, Python exact
arithmetic also checks six local anomalies. That calculation is not
thereby a complete six-anomaly theorem in Lean or an already viable full
particle model.

Similarly, at the three-generation point with ordinary fermions, a total
of twenty complex Higgs doublets reduces the weak coefficient to zero.
Passing the ordinary necessary budget therefore does not guarantee
positivity after arbitrary scalar additions. Allowing additional cubic
compensation for a nonzero ordinary-row anomaly can admit one-generation
and other points in a wider class. Counterexamples must remain alongside
positive results so that summaries do not erase physical premises.

The complete six-local-anomaly sums of §§31.4--31.5, the framing/index
route of Chapter 32, and the classification on the correct high-energy
gauge group in Chapter 36 are distinct mathematical objects. Each
retains its own analytical and exact-regression evidence. A count of
formalized theorems cannot replace these precise coverage relations.

\chapter{32. Further results, topological routes, and open
boundaries}\label{ch-32}

Topological and framing ideas yield testable selection only after their
realizations, maps, and completion conditions are specified. This
chapter retains checkable conditional classifications and removes
invalid shortcuts such as deriving generations directly from Betti
numbers, keeping open research distinct from the established chain.

\section{32.1 Conditional classification of framing completions}

The following candidate route is independent of HGP. The two premise
sets cannot be interchanged without explanation.

\textbf{H-QFR-1 (FAD).} Specify a dimensional-reduction or defect
realization whose descendant must forget a chosen framing autonomously,
without additional inflow. Assume completion contributions are additive
per generation and recorded as
\(\chi=2c_-^{\mathcal C}\in\mathbb R/48\mathbb Z\), with the designated
descent of the ordinary \(4N_c+3\) Weyl components per generation
contributing the same number of units. The resulting condition is \[
N_f(4N_c+3+\chi)=0\quad\text{in }\mathbb R/48\mathbb Z.
\] This is one compatibility equation within a complete physical
assumption package, not a general four-dimensional anomaly-cancellation
formula.

Complex colour \(N_c\ge3\), KM-type CP requiring \(N_f\ge3\), and the
necessary weak-AF bound leave seven candidates: \[
(3,3),(3,4),(3,5),(4,3),(4,4),(5,3),(6,3).
\]

\phantomsection\label{p-qfr-01}{}

\textbf{P-QFR-01.} On this seven-point set, the framing equation selects
exactly \((3,3)\) if and only if \(\chi=1\bmod48\).

\textbf{Proof.} Acceptance of \((3,3)\) requires
\(3(15+\chi)=0\bmod48\), leaving residues 1, 17, or 33 modulo 48. For 1,
the condition is \(12\mid N_f(N_c+1)\), and only \((3,3)\) among the
seven has product 12. Residue 17 also admits \((4,4)\), while 33 also
admits \((3,4),(3,5)\). \(\square\)

For \(\chi=0\), oddness of \(4N_c+3\) requires \(16\mid N_f\),
contradicting the upper bound. If completions are restricted to zero or
one singlet Weyl field per generation, the narrow class selects one.
This does not exclude coupled TQFTs or completions not additive
generation by generation.

The same four-dimensional anomaly class need not yield the same
descendant \(\chi\). A concrete natural descent map and a proof of
framing response are needed. That universal bridge has not been
established.

\section{32.2 Autonomy, object partitions, and external completions}

\phantomsection\label{p-qfr-02}{}

\textbf{P-QFR-02.} If a joint process factorizes exactly as
\(\mathbf T_{FC}=\mathbf T_F\otimes\mathbf T_C\), and a legitimate cut
deletes only cross-couplings, the cut leaves the process unchanged. Its
distance is zero in every operational distance.

\textbf{Proof.} No cross term exists to delete, so the compared
responses are equal. \(\square\)

This prevents a wholly decoupled external system from being called a
causally integrated parent merely by naming it so. Coupled topological
sectors, shared gauge variables, and line/surface protocols can have
real cross-structure and are not excluded. Parent autonomy also does not
imply autonomy of every descendant with a boundary: bulk inflow can
cancel a boundary anomaly.

A flavour-space dimension \(N_f\) does not automatically supply a
gauge-independent basis of generation subobjects. Masses, interactions,
and implementable protocols may select structures, but they need
certificates. An arbitrary coordinate basis is not nature's
decomposition.

\section{32.3 The retained topological research route}

\textbf{K-QIDX-01.} For a closed spin manifold and an appropriate
twisting bundle, \[
\operatorname{Index}D_E=\int_M\widehat A(TM)\operatorname{ch}(E)
\] taking the top-degree component. On \(S^2\), a line bundle of degree
\(n=(2\pi)^{-1}\int F\) gives \(\operatorname{Index}D_E=n\), while
\(b_1(S^2)=0\). The index counts the difference between positive- and
negative-chirality zero modes, not an arbitrary Betti number.

The legitimate chain is: independently select \(M,E,D\); calculate the
index; connect it to net chiral zero modes; then check paired extra zero
modes and the low-energy spectrum. Choosing a degree-three bundle and
computing three is a conditional realization, not generation selection.
The original implication \(b_1=3\Rightarrow\operatorname{Index}D=3\) is
excluded from this book's proof system.

\section{32.4 Claims removed from the established dependency chain}

\begin{longtable}[]{@{}
  >{\raggedright\arraybackslash}p{(\columnwidth - 4\tabcolsep) * \real{0.2500}}
  >{\raggedright\arraybackslash}p{(\columnwidth - 4\tabcolsep) * \real{0.3400}}
  >{\raggedright\arraybackslash}p{(\columnwidth - 4\tabcolsep) * \real{0.4100}}@{}}
\toprule\noalign{}
\begin{minipage}[b]{\linewidth}\raggedright
Claim
\end{minipage} & \begin{minipage}[b]{\linewidth}\raggedright
Status here
\end{minipage} & \begin{minipage}[b]{\linewidth}\raggedright
What would be required to restore it
\end{minipage} \\
\midrule\noalign{}
\endhead
\bottomrule\noalign{}
\endlastfoot
A closed-orbit phase integral directly establishes CCR & N & Regular
operations, representations, continuity, and a nontrivial central
character \\
Finite information automatically yields the standard uncertainty
constant & N & Physical operators, state domains, and calibration of
\(\hbar\) \\
Public records imply a unique objective collapse & O & Dynamics and
experiments distinguishable from open-system accounts \\
Arbitrary two-angle tests prove the global Born rule & N & Global
quantifiers or a fixed finite-parameter model \\
\(b_1\) directly determines three generations & N & An index realization
and independent selection mechanism \\
Common-\(q\) or CP capacity alone selects three & N & Independent
conditions excluding other consistent candidates \\
An autonomous parent forces an autonomous descendant & N & Additional
no-inflow conditions \\
Gauge-anomaly cancellation automatically selects the full SM & N &
Selection of the gauge group, representations, global structure, and
parameters \\
General object mathematics has proved HGP/FAD & O & An explicit
realization theorem and exclusion of countermodels \\
\end{longtable}

QBX or other quantum-foundations breakthrough candidates can advance
only through mathematical model, simulation evidence, hardware protocol,
real blinded testing, and independent reproduction. This book contains
no new hardware data and announces no discovery of objective collapse,
supraquantum performance, or particles beyond the Standard Model.

\chapter{33. Interface contracts for the complete
architecture}\label{ch-33}

Composition across chapters can easily lose types, normalization, system
identity, or evidential scope. Interface contracts make finite quantum
programs, memory processes, field models, and data exchange checkable
results, while distinguishing specifications from implemented features.

\section{33.1 Layers and responsibilities}

\begin{longtable}[]{@{}
  >{\raggedright\arraybackslash}p{(\columnwidth - 6\tabcolsep) * \real{0.2100}}
  >{\raggedright\arraybackslash}p{(\columnwidth - 6\tabcolsep) * \real{0.2000}}
  >{\raggedright\arraybackslash}p{(\columnwidth - 6\tabcolsep) * \real{0.2800}}
  >{\raggedright\arraybackslash}p{(\columnwidth - 6\tabcolsep) * \real{0.3100}}@{}}
\toprule\noalign{}
\begin{minipage}[b]{\linewidth}\raggedright
Layer
\end{minipage} & \begin{minipage}[b]{\linewidth}\raggedright
Main objects
\end{minipage} & \begin{minipage}[b]{\linewidth}\raggedright
Outputs
\end{minipage} & \begin{minipage}[b]{\linewidth}\raggedright
Unlicensed upgrade
\end{minipage} \\
\midrule\noalign{}
\endhead
\bottomrule\noalign{}
\endlastfoot
Core & Protocols, histories, response equivalence, resources, and
control & Predictive factors, carriers, persistence certificates &
Assuming complex Hilbert space, Born probabilities, or a particle
table \\
Q-Selector & Finite causal OPT and additional principles & Conditional
selection of a quantum realization & Calling sector principles Core
theorems \\
Q-Operations & States, channels, instruments, and effects & Sequential
probabilities, posteriors, error bounds & Assuming a one-time channel is
automatically multitime memoryless \\
Q-Objects & Encoding/decoding, access, coherent parents & Effective
factors, public records, SRN kernels & Equating protocol invisibility
with ontological nonexistence \\
Q-Dynamics & Self-adjoint operators, CPTP evolution, process tensors &
Time responses, control, stationary states & Hiding domains, initial
correlations, or truncation \\
Q-Fields & Local algebras, field distributions, regulators, actions &
Correlators, amplitudes, EFT coefficients & Equating a formal series
with a nonperturbative construction \\
Q-Particles & Gauge groups, representations, anomalies, Higgs/flavour &
Candidate consistency and conditional selection & Equating consistency
with a unique realization in nature \\
Q-Evidence & Data, calibration, blinded tests, statistical plans &
Intervals, model comparisons, failure explanations & Equating synthetic
validation with experimental evidence \\
\end{longtable}

Every layer carries model identity, units, domain, assumptions, and
certificates. Interfaces exchange typed results, not merely shared
abbreviations.

\section{33.2 A minimal API for finite quantum operations}

The following specification covers the full quantum architecture.
Section 33.8 maps the implemented finite exact kernel to actual names.
Completion of that kernel does not automatically implement the
specification's general Choi, continuous-domain, or process-tensor
interfaces. Matrices are two-dimensional complex arrays with globally
fixed vectorization and tensor-factor order.

\begin{Verbatim}[breaklines,breakanywhere,fontsize=\scriptsize]
SystemSpec:
  labels: ordered subsystem labels
  dims: positive integers
  basis_id: immutable string
  units: explicit energy/time/action convention

State:
  system: SystemSpec
  rho: Hermitian positive matrix
  trace: 1, or separately tagged subnormalized
  validity: tolerance + minimum eigenvalue + trace residual

Channel:
  input, output: SystemSpec
  choi: input-first, UNNORMALIZED Choi
  validate: J >= 0; Tr_output J = I_input
  apply(rho): Tr_input[(rho^T tensor I_output) J]
  compose(after, before): match system identities, not just dimensions

Instrument:
  maps: outcome -> CP trace-nonincreasing map
  validate: sum of maps is TP
  observe(rho): outcome probabilities + subnormalized branches
  condition(outcome): normalized state only if probability > 0

Protocol:
  initialization
  causal history -> permitted instrument/action
  terminal effects
  resource/time budget
  target and stopping rule

Certificate:
  model_id, protocol_id, target_id
  assumptions, theorem_ids, estimator_spec
  domain, tolerances, confidence, failure_modes
  evidence_kind: proof | exact_arithmetic | synthetic | experiment
\end{Verbatim}

\texttt{Channel.\allowbreak compose} is sequential composition, not a Kronecker
product; tensor composition is parallel. Our vectorization convention is
\(|A\rangle\rangle=\sum_i|i\rangle_{\rm in}\otimes A|i\rangle_{\rm out}\),
so \(J=\sum_\mu|A_\mu\rangle\rangle\langle\langle A_\mu|\). An
unnormalized Choi matrix has trace equal to the input dimension. A
state-normalized Choi matrix is a different, explicitly labelled object.

Invalid inputs should return interpretable errors such as
INVALID\_\allowbreak STATE, NOT\_\allowbreak CP, NOT\_\allowbreak TP, BASIS\_\allowbreak MISMATCH, or
ZERO\_\allowbreak PROBABILITY\_\allowbreak BRANCH. A nonpositive matrix must not be silently
clipped into a valid model. Requested projection-based repair should
also return the original residual and the constraint problem solved.

\section{33.3 Memory processes and combs}

In temporal order, let the input leg of process step \(j\) be \(2j-2\)
and its output leg \(2j-1\). A deterministic \(n\)-comb has Choi
operator \(R^{(n)}\) satisfying \[
R^{(n)}\succeq0,\qquad
\operatorname{Tr}_{2n-1}R^{(n)}
=I_{2n-2}\otimes R^{(n-1)},\quad R^{(0)}=1.
\] Compare each recursive relation after the specified leg permutation,
not an array's default ordering. Insert instruments using compatible
Choi link contractions; transposes on shared legs must agree with §33.2.
Only the sum over selectable-outcome combs has deterministic
normalization.

A process tensor absorbs initialization, free system propagation, and
environmental memory into a comb-like object assigning probabilities or
final states to slot operations. Its API includes slot times,
input/output legs, the controllable operation set, and the allowed
initial correlations. Compare two processes only on a common tester set.
Complete positivity, CP divisibility, Markov factorization, and future
sufficiency must be reported separately.

\section{33.4 Object interfaces and algorithms}

\begin{Verbatim}[breaklines,breakanywhere,fontsize=\scriptsize]
QuantumObjectSpec:
  carrier: physical system or accessible algebra
  protocol_family: allowed future interventions/readouts
  target: predicted observable/history
  encoder, decoder: optional CP maps
  lifecycle: initialization, survival, recovery, termination
  cuts: legal access restrictions or physical interventions
  records: pointer projectors + fragment specification
  tolerance: metric, horizon, confidence

certify_predictive_factor(spec):
  verify admissibility and common protocol identity
  compare future response laws or derive exact intertwining
  distinguish exact equality from finite-test agreement
  return factor certificate, unresolved directions, witness protocols

certify_SRN(channel, projectors):
  check each branch's entire operator algebra, not only diagonal populations
  recover block coefficients K_ab
  test Hermiticity, diagonal=1, PSD and model residual
  return K, generator status if time data available, deviations

discover_access_carriers(spec):
  enumerate permitted access subsets or use a declared search algorithm
  test target sufficiency for each
  output minimal certified subsets relative to this finite universe
  do not claim global minimality beyond searched cuts
\end{Verbatim}

Restricted access reduces available effects; physical cutting or erasure
changes the process. Keeping these interfaces separate prevents
inaccessible coherence from being mistaken for vanished coherence.

\section{33.5 Continuous variables, operators, and numerical backends}

An \texttt{Operator\allowbreak Spec} records at least the Hilbert space, dense
domain, boundary conditions, symmetric/self-adjoint status, spectral
lower bound, common invariant core, and chosen extension. A differential
expression without boundary conditions is not a self-adjoint
Hamiltonian.

A \texttt{CVBackend} records phase-space ordering, \(\hbar\), symplectic
form, Weyl multiplier, Gaussian mean/covariance conventions, and
truncation. Finite Fock truncation approximates low-energy behaviour
only; it cannot claim exact CCR. Separate tail probability,
truncated-dynamics error, integration error, and statistical error. A
static energy-tail bound does not automatically replace the latter two.

\section{33.6 Field and particle backends}

\begin{Verbatim}[breaklines,breakanywhere,fontsize=\scriptsize]
QFTModel:
  spacetime, metric, units
  local_fields + representations + statistics
  global_gauge_group and permitted bundles
  regulated_action and gauge_fixing
  scheme, scale, renormalization_conditions
  perturbative_order or nonperturbative_definition
  operator_basis, boundary_conditions, infrared_resolution

AmplitudeRequest:
  external states, mass/charge convention
  spin preparation or averaging
  kinematic domain and phase-space cuts
  inclusive/exclusive definition
  requested order, EFT power and uncertainty

AnomalyCertificate:
  all-left-handed representation table
  multiplicities, Dynkin and cubic normalizations
  local anomaly polynomial coefficients
  global/discrete/boundary/inflow checks: done or unresolved
  assumptions under which quotient is defined
\end{Verbatim}

Checking six vanishing local coefficients for a field table should
return LOCAL\_\allowbreak ANOMALIES\_\allowbreak CANCELLED, not
FULL\_\allowbreak QUANTUM\_\allowbreak THEORY\_\allowbreak CONSISTENT. Perturbative computability, global
gauge questions, a stable vacuum, unitarity, the actual spectrum, and
experimental exclusions require further checks.

\section{33.7 Data, versions, and evidence interfaces}

Every experimental datum must point to raw counts or a justified
sufficient statistic, preparation/measurement calibration, acquisition
times, randomization, exclusion rules, and a data freeze. Preserve noisy
response-matrix estimates, not just their PSD-projected ``physical''
versions.

A version change records the previous claim, reason for change, added or
removed premises, affected downstream results, and regression outcomes.
Changes in protocol, target, basis, units, or action-operator basis are
semantic changes; a file version alone is insufficient without an
updated model identity.

\section{33.8 Implemented finite-kernel interfaces}

\begin{longtable}[]{@{}
  >{\raggedright\arraybackslash}p{(\columnwidth - 4\tabcolsep) * \real{0.2500}}
  >{\raggedright\arraybackslash}p{(\columnwidth - 4\tabcolsep) * \real{0.3400}}
  >{\raggedright\arraybackslash}p{(\columnwidth - 4\tabcolsep) * \real{0.4100}}@{}}
\toprule\noalign{}
\begin{minipage}[b]{\linewidth}\raggedright
Research interface
\end{minipage} & \begin{minipage}[b]{\linewidth}\raggedright
Actual entry point
\end{minipage} & \begin{minipage}[b]{\linewidth}\raggedright
Implemented scope
\end{minipage} \\
\midrule\noalign{}
\endhead
\bottomrule\noalign{}
\endlastfoot
Exact states and operations & \texttt{QComplex}, \texttt{Matrix},
\texttt{Kraus\allowbreak Term}, \texttt{Instrument} & Gaussian-rational matrices and
nonnegative rational weights; Python PSD checking in dimensions 1--4 \\
Allowed library and finite trees & \texttt{Instrument\allowbreak Library},
\texttt{Protocol\allowbreak Node} & Context fingerprints, complete outcomes, exact
finite leaf laws \\
Response representation & \texttt{Predictive\allowbreak Model},
\texttt{Encoded\allowbreak State} & Word-effect closure search, exact coordinates,
branch and posterior updates \\
Portable representation certificates &
\texttt{export\_\allowbreak representation\_\allowbreak certificate},
\texttt{verify\_\allowbreak representation\_\allowbreak certificate} & Independent checks of
numerical bases, generating words, every local identity, and context \\
Numerical Lean entry & \texttt{export\_\allowbreak lean\_\allowbreak representation},
\texttt{Exact\allowbreak Certificate.\allowbreak check} & Generate a concrete finite model, then
kernel-check its truth proof and consequences \\
Actual effects and error mathematics & \texttt{Effect},
\texttt{effect\allowbreak Distance}, \texttt{channel\allowbreak Defect} & Supremum over all
effects, complete-channel contraction, positive-event posteriors, and
trajectory theorems; the supremum is not a runtime optimizer \\
Particle spectra & \texttt{Spectrum}, \texttt{inspect\_\allowbreak spectrum},
\texttt{local\_\allowbreak 331\_\allowbreak anomalies} & Exact arithmetic on explicit spectrum
rows and windows; separate scope labels for six local anomalies and Lean
conditional selection \\
\end{longtable}

Python exact arithmetic, general Lean theorems, and compilation of a
particular exported model are three independent evidence layers. The
representation checker accepts dimensions 1--4, at most \(d^2\) word
effects, and witness words of length at most \(d^2-1\). These limits
constrain the certificate format, not the length of all future responses
certified by it. The Lean exporter additionally limits input rationals
to 4096 bits, padded scalar storage to 65536 entries, and source size to
8 MiB. These software boundaries are not physical constants.

The minimal reproduction path is: construct an instrument library with
fixed context; generate a representation certificate; independently
check the local matrix conditions; export concrete Lean data; actually
compile its \texttt{checked} proof; read axiom and source-hash receipts;
and only then cite the corresponding all-word, posterior, or error
result in Chapters 49--50. Export alone retains status
\texttt{LEAN\_\allowbreak SOURCE\_\allowbreak EXPORTED\_\allowbreak NOT\_\allowbreak COMPILED}, not formal verification.

\chapter{34. Complete application paths}\label{ch-34}

An application must connect preparation, calculation, measurement,
error, and controls rather than stop at a result formula. Calibration,
record reversal, correction, many-body models, and particle candidates
become executable paths, with fitting data distinguished from
observations reserved for testing.

\section{34.1 Single-qubit and channel calibration}

Inputs are an unknown preparation, three calibrated Pauli measurement
settings, and counts. Linear reconstruction uses
\(\hat r_j=(n_{j,+}-n_{j,-})/n_j\) and
\(\hat\rho=(I+\sum_j\hat r_j\sigma_j)/2\). Retain the raw estimate and
confidence intervals. A value \(|\hat r|>1\) may result from sampling
noise, drift, or model mismatch; clipping alone does not establish
experimental agreement with the model.

Channel tomography reconstructs \(J\) from a linear input basis and
informationally complete output measurements, then checks CP, TP, and
held-out inputs. Without independent SPAM calibration, gate error cannot
be uniquely separated from preparation and measurement error. Use joint
estimation or an explicit gauge-equivalence report. Metrics can include
average fidelity, worst-case protocol error, leakage, and drift; one
high-fidelity number does not replace all certificates.

\section{34.2 Quantum records, QDRR, and actualization tests}

The full sequence is coherent preparation, pointer/environment record
writing, a specified waiting period and public readout, inverse writing
within a controllable boundary, and phase tomography of coherence.
Record access, leakage to uncontrolled environments, and
inverse-operation errors are recorded separately.

Within SRN for \(m\) branches, reconstruct \(K\) from the \(m(m-1)/2\)
branch pairs, with real and imaginary phase settings for each. Chapter
8's bound assumes \(n\) independent bounded outcomes per phase. For
\(m=3,\epsilon=0.12,\alpha=0.05\), this requires 762 shots per setting,
six settings and 4572 shots in total, excluding calibration and
controls. Test SRN itself as well; exact SRN cannot simply be assumed
before claiming discovery or exclusion of new physics.

Important controls vary storage duration, environmental access, echo
accuracy, record redundancy, randomized order, and the empty protocol.
If additional irreversible loss appears, first compare frozen
open-system and device-noise models. A new actualization model needs
independent parameters, scale dependence, and falsifiable held-out
predictions. Without real counts, the result is a protocol design.

\section{34.3 Distinct uses of two-phase Gram and Choi tests}

Internal labels such as Paper-A/B/C are not proof categories. Their
applications reduce to three kinds of objects:

\begin{longtable}[]{@{}
  >{\raggedright\arraybackslash}p{(\columnwidth - 4\tabcolsep) * \real{0.2500}}
  >{\raggedright\arraybackslash}p{(\columnwidth - 4\tabcolsep) * \real{0.3400}}
  >{\raggedright\arraybackslash}p{(\columnwidth - 4\tabcolsep) * \real{0.4100}}@{}}
\toprule\noalign{}
\begin{minipage}[b]{\linewidth}\raggedright
Application
\end{minipage} & \begin{minipage}[b]{\linewidth}\raggedright
Data and reconstruction
\end{minipage} & \begin{minipage}[b]{\linewidth}\raggedright
Supported conclusion
\end{minipage} \\
\midrule\noalign{}
\endhead
\bottomrule\noalign{}
\endlastfoot
Quadratic interference structure & Basis states and pairwise responses
at 0 and \(\pi/2\) & A Gram estimate within a declared finite quadratic
model \\
Operational process & Input basis, complete output readout, and Choi
constraints & CP/TP compatibility at specified dimension and
calibration \\
Coherence after public records & Write--store--reverse and \(K_{ab}\) &
SRN parameters and differences from control models \\
\end{longtable}

Freeze the model and test statistic, retain raw residuals, perform any
needed constrained estimation, and test on unused data. Projecting data
onto a quantum set and observing that the projection is quantum does not
prove that every probability law must be quantum.

\section{34.4 Quantum memory, correction, and control}

For encoding \(V:\mathcal H_L\to\mathcal H_P\), noise \(\mathcal E\),
and recovery \(\mathcal R\), the logical channel is
\(\mathcal N_L=\mathcal D\mathcal R\mathcal E\mathcal V\).
Knill--Laflamme specifies ideal correctable classes. Approximate
correction must directly assess logical errors and leakage.

A three-qubit repetition code corrects a single \(X\) error, not
arbitrary single-qubit quantum errors. Ignoring this difference
misreports classical-bit protection as full quantum error correction.
Fault-tolerance thresholds require a particular code, noise-correlation
model, decoder, and gate schedule. No model-independent threshold is
given here.

Closed-loop control uses Chapter 20's value or pulse-gradient
calculations. A real controller may depend only on records already
acquired; computation latency and bandwidth are resources. GOA is
reported as logical survival probability, a stationary-state set, or a
recoverable domain, without presupposing a universal ultimate state.

\section{34.5 Many-body quenches and thermal environments}

For a finite TFIM, specify boundaries, \(J,h\), initial state, and
quench time, then calculate magnetization, correlations, and entropy.
Exact diagonalization of small systems checks tensor networks;
short-time series check integrators. Converge finite volume, bond
dimension, Trotter step, and observation error separately.

With an environment, first check the scale separation needed for GKSL.
If distinct preparation histories yield different futures after the same
causal break, the system state is not a sufficient predictive state: add
memory variables or use a process tensor. Finer time steps cannot repair
an incorrect memoryless model.

\section{34.6 Interfaces for semiconductor and superconducting devices}

A quantum dot may start from a finite-orbital Anderson-type Hamiltonian
with declared Coulomb energies, tunnelling, temperature, bias, and
readout. A common effective BdG model for a proximitized semiconductor
nanowire is \[
H_{\rm BdG}(k)=
\left(\frac{k^2}{2m_*}-\mu\right)\tau_z
+\alpha_R k\,\sigma_y\tau_z+V_Z\sigma_x+\Delta\tau_x,
\] with fixed Nambu/spin basis and phase conventions. This is an
effective device-model input, not a Core theorem. Disorder, multiple
bands, orbital effects, self-energies, and contacts can alter the
simplified predictions. See the theoretical precedents of
\href{https://arxiv.org/abs/1002.4033}{Lutchyn--Sau--Das Sarma} and
\href{https://arxiv.org/abs/1003.1145}{Oreg--Refael--von Oppen}.

The gap-closing condition \(V_Z^2=\mu^2+\Delta^2\) holds within the
assumptions of the closed, uniform, single-band model. A local zero-bias
peak or an Andreev-bound-state signal alone does not prove a topological
Majorana mode. Trivial mechanisms must be excluded through more
discriminating protocols. No experimental discovery in such a device is
claimed for this project.

\section{34.7 Comparing scattering, EFTs, and particle candidates}

For QED, calculate the Born cross section first, then add masses,
radiation, acceptance, and their uncertainties in controlled stages. For
EFTs, specify \(E/\Lambda_*<1\) and operator truncation, match
coefficients, and compare held-out processes. Do not force a low-order
local expansion to fit near a heavy pole.

For HGP/331 candidates, anomaly freedom starts further work: scalar
breaking, spectra and mixing, additional gauge bosons and exotic
particles, flavour-changing neutral currents, precision observables,
vacuum stability, global anomalies, and running windows. Framing
candidates first require descent and completion classifications. Freeze
predictions, fitted parameters, and exclusion criteria separately.
Passing three-generation arithmetic does not establish the complete
model's success.

\chapter{35. Assembling the quantum framework: errors and
extensions}\label{ch-35}

After modules are assembled, the framework must show whether local
errors still control final records. This chapter combines channel and
confidence bounds under common protocols and targets, states what
historical verification receipts support, and provides precise entry
points to the later finite-kernel proofs.

\section{35.1 A common assembly rule}

In a finite memoryless realization, \[
p(\mathbf a\mid\pi)=
\operatorname{Tr}\!\left[
\mathcal M_{a_T|\mathbf a_{<T}}\circ\cdots\circ
\mathcal M_{a_1}(\rho_0)\right].
\] With memory, replace this by pairing a legitimate comb with its
tester. Continuous variables require energy and operator domains. Field
theory connects preparation, detection, and phase-space resolution to
correlators or amplitudes. Core then uses these probability responses to
define objects, future sufficiency, and control tasks.

The common structure is therefore the composition of realization,
intervention, response, and certificate. It makes each physical layer's
independent content traceable; it does not collapse all physics into one
untyped formula.

\section{35.2 Composable error certificates}

First state the analytical contract including ancillary systems. An
instrument retains every label in a CPTP channel, rather than only a
selected branch \(\mathcal M_o\): \[
\widehat{\mathcal M}(X)=\sum_o|o\rangle\langle o|\otimes\mathcal M_o(X).
\] Adaptive control stores the prior classical history explicitly and
selects the next instrument by history block. Stopped paths retain an
absorbing stop marker. Completeness makes every full step CPTP on a
common classical--quantum carrier. Compared models need matching labels,
registers, and allowed protocols. Channels equal only after outcomes are
discarded can still have different multitime record laws, as §5.3 shows.

For a common protocol and target, suppose the initial-state
trace-distance error, including the compared initial memory, is at most
\(\epsilon_0\). Suppose each complete step satisfies \[
\tfrac12\|\widehat{\mathcal E}_j-\widehat{\widetilde{\mathcal E}}_j\|_\diamond
\le\epsilon_j,
\] with arbitrary quantum ancillas, and the terminal-readout
total-variation defect is uniformly at most \(\epsilon_M\). Then \[
\operatorname{TV}(p,\widetilde p)
\le\min\left\{1,\epsilon_0+\sum_j\epsilon_j+\epsilon_M\right\}.
\] Replace channels successively in the actual composition. Use the
diamond defect on the common ancilla, contraction by subsequent complete
channels, the triangle inequality, and measurement contraction. This
argument does not treat an unnormalized selected branch as TP or omit
adaptive memory from a model without a record register.

Chapter 50 supplies a \textbf{Lean-checked finite portion} of this
analytical contract: actual complete Kraus channels contract the
all-effect distance \(d_E\); their actual output defect \(\Delta_E\)
gives a one-step recursion; and existing Core error-transport theorems
propagate it. That portion does not formalize the diamond norm,
arbitrary-ancilla stability, or general history-register assembly. The
phrase ``channel error'' does not supply these properties by
implication.

Sector-model approximation, truncation, and estimation errors can be
added only after each is bounded in the same target-probability metric.
Energy error, fidelity error, log likelihood, and Wilson-coefficient
error first need proved conversions. Postselection separately requires a
positive-event lower bound. The actual-Born-branch result developed in
Chapter 50 is \(d_E(\rho_o,\sigma_o)\le2d_E(\rho,\sigma)/p_o(\rho)\). At
small probabilities it may be uninformative, and zero-probability events
have no normalized posterior.

A union bound combines confidence events with total failure probability
at most the sum. Independence cannot be assumed without checking.
Estimates of unknown higher orders in an asymptotic series are model
uncertainty, not automatically frequentist confidence intervals with
guaranteed coverage.

\section{35.3 Scope of the inherited independent checks}

Appendix D contains the complete code and actual historical execution
results. The three inherited programs were rerun for edition 3.0,
comprising 757 finite checks. This bilingual publication reuses those
receipts for unchanged programs and does not claim a fresh execution.
The NumPy/SciPy and rational-arithmetic programs cover:

\begin{itemize}
\tightlist
\item
  CPTP/Choi relations, non-CP counterexamples, local rebit blindness,
  coherent records, and inverse writing;
\item
  SRN block structure, complex Gram kernels, generator positivity and
  counterexamples, and phase recovery;
\item
  truncated CCR, Gaussian positivity, Rabi dynamics, spin, CHSH, and
  control gradients;
\item
  open dynamics, KMS, free energy, and causal-break memory;
\item
  Clifford relations, tree-level QED Ward identities, scattering traces,
  loop parameter integrals, and EFT remainders;
\item
  Higgs/seesaw benchmarks, rational SM and 331 anomalies, the
  three-generation search without primitive packages, and framing
  classification.
\end{itemize}

These are finite numerical and exact-arithmetic regression checks, not
formal verification of every theorem. In particular, they do not
independently machine-prove the reconstruction theorem, Stone--von
Neumann, GKSL, LSZ, spin--statistics, or a four-dimensional
nonperturbative construction. They contain no real experimental data.
Finite integer searches are cross-checks; the all-integer
three-generation result rests on Chapter 31's analytical proof, with
separate formal support from \texttt{Selector}'s integer theorem with
full premises and \texttt{Spectrum}'s conditional spectral bridge. The
Lean results in Chapters 49--50 cover their exact definitions and
theorem inventories, not every physical calculation listed here.

\section{35.4 Completed work, improvements, and open questions}

The book maintains a continuous formal text covering every declared
module and integrates the finite quantum kernel and error proofs in
Chapters 49--50. Substantive improvements preserve Core's neutrality
across domains; give coherent-parent, public-record, and posterior
aspects of measurement distinct interfaces; state SRN complex-generator
conditions, continuous-variable domains, process memory, the three types
of RG, and anomaly responsibilities explicitly; remove the redundant
primitive-package premise from the conditional three-generation theorem;
and complete the general 331 representation inventory.

Open questions include an assumption-free derivation of nature's quantum
physics from Core, an objective actualization mechanism, a complete
nonperturbative quantum-field construction, selection of gauge groups
and all particle parameters, independent physical necessity of HGP/FAD,
and an unconditional explanation of mass hierarchies and three
generations. They remain open; a higher version number is not a proof.

The later constructions in Chapters 58, 62, and 65 complete the
specified finite history-memory and common-controller cases mentioned in
§35.2. They retain their explicit hypotheses and do not supply a general
diamond-distance theorem.

\chapter{36. The generation problem: calculations must use the correct
gauge group}\label{ch-36}

The same field content has different gauge carriers and running
coefficients in different energy regimes, so generation selection first
requires a model and threshold specification. This chapter recalculates
anomaly and asymptotic-freedom conditions for the correct high-energy
group and retains the five-colour counterexample; read it alongside the
low-energy conditional selection in Chapter 31.

\section{36.1 The fixed representation content}

Consider \[
G=SU(N_c)_C\times SU(3)_L\times U(1)_X,\qquad N_c\ge3.
\] Here \(SU(3)_L\) is a physical input; it does not follow from the
general UEOT Core. Each lepton generation occupies a colourless \(3_L\);
\(k\) quark generations occupy \((N_c,3_L)\), and the remaining
\(N_f-k\) occupy \((N_c,\bar3_L)\). Each quark triplet is accompanied by
three corresponding colour-antifundamental, right-handed conjugate
singlets. A general charge embedding also requires the appropriate
singlet for the third lepton component. All fermions are written as
left-handed Weyl fields.

The pure \(SU(3)_L^3\) anomaly is \[
A_L=N_f+N_c(2k-N_f).
\] Cancellation therefore requires \[
2N_ck=(N_c-1)N_f.
\] Writing \(d=\gcd(2,N_c-1)\), all positive-integer solutions are \[
N_f=\frac{2N_c}{d}t,\qquad
k=\frac{N_c-1}{d}t,\qquad t=1,2,\ldots
\tag{36.1}
\] This follows by divisibility of the coprime coefficients. Odd colour
number gives \(N_f=tN_c\), while even colour number gives \(N_f=2tN_c\).
The simplicity of the smallest solution is not a reason to assert that
nature must choose \(t=1\).

This representation pattern has precedents in 331 research. Flipped
models can instead place generation nonuniversality in the lepton
sector, so it is not the unique form of every extended gauge theory.
\href{https://arxiv.org/abs/hep-ph/9206242}{Original 331 paper},
\href{https://arxiv.org/abs/1606.01109}{flipped 331}.

\section{36.2 Why the earlier weak-sector bound is not a high-energy
requirement}

The conditional theorem in the quantum treatment uses \[
N_f(N_c+1)<22,
\tag{36.2}
\] a necessary one-loop asymptotic-freedom bound for a particular
low-energy \(SU(2)_L\) field content. The arithmetic selecting
\((3,3,1)\) under this premise is correct.

Above a scale \(M_B\), however, the weak group may already have enlarged
to \(SU(3)_L\), changing the gauge-boson self-interaction contribution.
Requiring a well-behaved high-energy theory does not automatically make
low-energy \(b_2>0\) necessary. An effective low-energy coupling can
grow over a finite window and then match to a different high-energy
group; a negative effective coefficient \(b_2\) alone does not rule out
that matching.

The coefficients must therefore be recalculated. Stronger wording cannot
preserve a three-generation conclusion that the revised premises do not
imply.

\section{36.3 One-loop coefficients for the full active field content}

Use the convention \[
\mu\frac{dg}{d\mu}=-\frac{b\,g^3}{16\pi^2}+O(g^5),\quad
b=\frac{11}{3}C_A-\frac23\sum_{\rm Weyl}T(R)
-\frac13\sum_{\rm complex\ scalar}T(R).
\tag{36.3}
\] Assume that all exotic fermions are active, that there are no
additional coloured scalars, and that the scalar sector contains \(n_s\)
complex \(SU(3)_L\) fundamentals with no other matter carrying \(L\)
charge.

For the colour group, each generation has three fundamental and three
antifundamental Weyl fields, equivalent to three Dirac colour flavours.
Thus \[
b_C=\frac{11N_c}{3}-2N_f.
\tag{36.4}
\] For \(SU(3)_L\), the sum of Weyl quadratic indices is
\(N_f(N_c+1)/2\), and each complex scalar has index \(1/2\). Hence \[
b_L=11-\frac{N_f(N_c+1)}3-\frac{n_s}{6}.
\tag{36.5}
\] This is consistent with the known contribution of extra quarks to QCD
running in extended weak-group models, but the counting here is
performed afresh for the stated \(N_c,N_f\) content.
\href{https://arxiv.org/abs/1205.4483}{Original research on extended
gauge groups}.

The conditions \(b_C,b_L>0\) establish only the one-loop asymptotically
free direction of the two non-Abelian factors. The \(U(1)_X\) coupling,
Yukawa interactions, scalar potential and higher orders require separate
analysis; two coefficient signs are not a complete UV completion.

\section{36.4 High-energy classification and a five-colour
counterexample}

\phantomsection\label{g-01}{}

\textbf{G-01: All integer solutions for the three-scalar content.} Under
(36.1), \(N_c\ge3\), \(n_s=3\) and \(b_C>0,b_L>0\), \[
(N_c,N_f,k)=(3,3,1)\quad\text{or}\quad(5,5,2).
\tag{36.6}
\]

\textbf{Proof.} The inequality \(b_L>0\) gives \(N_f(N_c+1)<31.5\).

For even \(N_c\ge4\), \(N_f\ge2N_c\), so the product is at least 40 and
is excluded. For odd \(N_c\ge7\), the product is at least 56 and is
excluded. Only 3 and 5 remain.

For \(N_c=3\), \(N_f=3t\) and \(b_C=11-6t>0\) force \(t=1\). For
\(N_c=5\), \(N_f=5t\) and \(b_L>0\) force \(30t<31.5\), hence \(t=1\);
also \(b_C=25/3>0\). Both points satisfy the conditions. \(\square\)

Explicitly, \[
\begin{array}{c|ccc}
(N_c,N_f,k)&b_C&b_L&b_2\text{ (low-energy, one doublet)}\\\hline
(3,3,1)&5&13/2&19/6\\
(5,5,2)&25/3&1/2&-17/6
\end{array}
\tag{36.7}
\] The final column shows precisely that asymptotically free directions
for both high-energy non-Abelian factors do not imply low-energy
\(b_2>0\).

This does not assert that nature has five colours. It is a countermodel
to the claim that high-energy physical consistency has already
independently selected three colours and three generations. When the
five-colour branch has charges different from those of observed quarks,
it cannot be presented as a fit to nature.

\section{36.5 All local anomalies can also cancel in the five-colour
counterexample}

To avoid checking only one anomaly, introduce \[
Q_u=\frac{N_c+1}{2N_c},\qquad
Q_d=-\frac{N_c-1}{2N_c},\quad Q_\nu=0,\quad Q_e=-1.
\] With \(Q=T_3+\beta_8T_8+X\), \[
X_L=-\frac12-\frac{\beta_8}{2\sqrt3},\quad
X_{Q+}=\frac1{2N_c}-\frac{\beta_8}{2\sqrt3},\quad
X_{Q-}=\frac1{2N_c}+\frac{\beta_8}{2\sqrt3}.
\] The third-component charges are \[
q_E=-\frac12-\frac{\sqrt3\beta_8}{2},\quad
q_{J+}=\frac1{2N_c}-\frac{\sqrt3\beta_8}{2},\quad
q_{J-}=\frac1{2N_c}+\frac{\sqrt3\beta_8}{2}.
\] Include \(e^c,E^c,u^c,d^c,J^c\) with the required multiplicities.

\phantomsection\label{g-02}{}

\textbf{G-02: Local completion for general colour number.} Under (36.1),
all six local anomalies, \(C^3,L^3,C^2X,L^2X,\mathrm{grav}^2X,X^3\),
vanish.

\textbf{Proof.} The three colour fundamentals and three antifundamentals
cancel the cubic colour anomaly generation by generation. The sum of
physical charges in each triplet is \(3X\), so the addition of
oppositely charged singlets cancels the mixed gravitational and colour
anomalies. The \(L^2X\) anomaly is \[
\tfrac12[N_fX_L+N_c kX_{Q+}+N_c(N_f-k)X_{Q-}]=0
\] because the constant and \(\beta_8\) terms vanish separately by the
chosen charges and (36.1).

For one triplet, write its charges as \(X+t_i\), with \(\sum_i t_i=0\).
After adding oppositely charged singlets, its cubic contribution is
\(3X^3-\sum_i(X+t_i)^3=-3X\sum_i t_i^2-\sum_i t_i^3\). The quadratic
sums are equal for the fundamental and antifundamental, while the cubic
sums have opposite signs. In the total, the first term vanishes by
\(L^2X=0\) and the second by \(L^3=0\). \(\square\)

This construction establishes compatibility of local anomaly
cancellation with the one-loop counting. It does not establish
cancellation of global anomalies, stable symmetry breaking, all masses,
or experimental viability. Fixing \(Q_u=2/3,Q_d=-1/3\) directly restores
the three-colour selection condition of Chapter 31; the measured charges
are then inputs.

\section{36.6 Which premises remain unresolved}

Within this fixed family, both points survive for \(0\le n_s\le5\); only
\((3,3,1)\) survives for \(6\le n_s\le41\); and no solution survives for
\(n_s\ge42\). The same inequality proves this: \(b_L=7-n_s/6\) at the
three-colour point and \(b_L=1-n_s/6\) at the five-colour point.
Introducing six scalars solely to exclude five colours would require an
independent justification.

If three colours are already fixed by experiment or another independent
principle, (36.1) and the correctly computed high-energy colour
coefficient give \(N_f=3t,\ b_C=11-6t>0\Rightarrow t=1\). This is a
conditional three-generation route with a clearer physical premise. Why
the extended weak carrier exists, and why it has this representation
content, remain open.

\chapter{37. Testable flavour constraints from generation
nonuniversality}\label{ch-37}

For a specified neutral-vector candidate with generation-dependent
couplings, unknown flavour rotations do not make every observation
freely adjustable. Rank-one structure yields sum rules and bounds, while
matching them to measured low-energy processes still requires masses,
couplings, running, and other contributions; no new particle observation
is reported.

This chapter does not require a complete solution to the origin of three
generations. It asks a more direct question: once the preceding
three-generation carrier is adopted, what structure must it leave at low
energy?

\section{37.1 Fixing a neutral-vector candidate}

Suppose a heavy neutral vector \(Z'\) couples to ordinary left-handed
quark doublets in the weak basis through \[
G_{\rm weak}=g_0I+\Delta g\,P,\qquad P^2=P=P^\dagger,\quad
\operatorname{rank}P=1,\quad \Delta g\ne0.
\tag{37.1}
\] The projector \(P\) selects the generation that differs from the
other two, and \(g_0,\Delta g\) are real. This structure arises from a
particular difference in generation representations; not every HGP or
flipped model has it.

Assume unitary mass rotations for the three ordinary generations,
neglect significant mixing with exotic states, and compare gauge
couplings of the same left-handed doublets. Set \[
V=U_{uL}^\dagger U_{dL},\quad
P_u=U_{uL}^\dagger PU_{uL}=uu^\dagger,\quad
P_d=V^\dagger P_uV=dd^\dagger,\quad d=V^\dagger u.
\tag{37.2}
\] Here \(u^\dagger u=d^\dagger d=1\). Tree-level FCNCs in general 331
models have long been studied. The specific algebraic relations and
bounds below are derived step by step; the existence of FCNCs is not
claimed as an original discovery.
\href{https://arxiv.org/abs/hep-ph/9312312}{Liu's original analysis}.

\section{37.2 Rank-one sum rules constrain the unknown parameters}

\phantomsection\label{f-01}{}

\textbf{F-01.} For \(B=G-g_0I=\Delta g\,uu^\dagger\), \[
B_{ij}B_{kl}=B_{il}B_{kj},\qquad
|B_{ij}|^2=B_{ii}B_{jj},\qquad
B_{12}B_{23}B_{31}=(\Delta g)^3|u_1u_2u_3|^2\in\mathbb R.
\tag{37.3}
\] Moreover, \[
B_d=V^\dagger B_uV.
\tag{37.4}
\]

\textbf{Proof.} Substitute \(B_{ij}=\Delta g\,u_i\bar u_j\) term by
term; the common factors give each equality directly. The phase around a
closed triangle is invariant under rephasing of the individual quark
fields. \(\square\)

These are three potentially falsifiable constraints: matrix rank,
magnitudes, and cyclic phase. If couplings of the same heavy vector
extracted from different processes violate them, the fixed rank-one
candidate is excluded. Refitting a mutually incompatible \(u\)
separately for each process would not preserve that candidate.

\section{37.3 FCNCs cannot vanish in both up and down sectors}

\phantomsection\label{f-02}{}

\textbf{F-02.} In a nondegenerate mass basis, if both \(P_u\) and
\(P_d\) are diagonal, at least one entry of \(V\) has modulus one, and
every other entry in its row and column vanishes. Hence, when
\(c=\max_{ij}|V_{ij}|<1\), at least one sector has a nonzero
off-diagonal gauge coupling.

\textbf{Proof.} A diagonal rank-one projection must be
\(|e_i\rangle\langle e_i|\). If
\(V^\dagger|e_i\rangle\langle e_i|V=|e_j\rangle\langle e_j|\), then
\(V^\dagger e_i=e^{i\theta}e_j\), so \(|V_{ij}|=1\). Unitarity forces
every other entry in that row and column to vanish. \(\square\)

The statement is stronger than the possibility of FCNCs: under the
stated assumptions, simultaneously aligning the distinguished generation
with both mass bases cannot hide all of them. Degenerate masses or a
nonunitary light sector require a new analysis.

\section{37.4 A quantitative bound independent of the unknown
orientation}

Let \(\operatorname{off}(A)\) denote the matrix obtained by removing the
diagonal, and let \(p_i=|u_i|^2,q_j=|(V^\dagger u)_j|^2\).

\phantomsection\label{f-03}{}

\textbf{F-03.} \[
\|\operatorname{off}G_u\|_F^2+
\|\operatorname{off}G_d\|_F^2
\ge |\Delta g|^2(1-c),\qquad c=\max_{ij}|V_{ij}|.
\tag{37.5}
\]

\textbf{Proof.} A rank-one projection has squared Frobenius norm one.
The left-hand side divided by \(|\Delta g|^2\) is therefore
\(2-\sum_i p_i^2-\sum_jq_j^2\). Choose \(p_i,q_j\) at their respective
maxima. Normalization gives \(\sum p_i^2\le p_i,\sum q_j^2\le q_j\). The
sum of the two rank-one projections \[
|e_i\rangle\langle e_i|+
V|e_j\rangle\langle e_j|V^\dagger
\] has largest eigenvalue \(1+|V_{ij}|\), obtained directly from the
quadratic characteristic equation on the span of the two vectors. Thus
\(p_i+q_j\le1+|V_{ij}|\le1+c\). Rearrangement proves (37.5). \(\square\)

The bound is generally conservative; it is not the exact minimum for
every \(V\). It is small when \(c\) is close to one, so strict
positivity does not imply easy experimental discovery. More precise
candidate comparisons can optimize over the four real degrees of freedom
of projective \(u\), but a numerical optimum does not replace the proved
lower bound.

\section{37.5 EFT matching and observable meaning}

With \(\mathcal L\supset \tfrac12M_{Z'}^2Z'_\mu Z'^\mu+Z'_\mu J^\mu\),
integrating out \(Z'\) below its mass gives \[
\mathcal L_{\rm EFT}=-\frac{J_\mu J^\mu}{2M_{Z'}^2}
+O(\partial^2/M_{Z'}^4).
\tag{37.6}
\] For the fixed chiral operator
\(O_{ij}^{u,d}=(\bar q_{Li}\gamma_\mu q_{Lj})
(\bar q_{Li}\gamma^\mu q_{Lj})\), \[
C_{ij}^{u,d}=-\frac{(G_{u,d})_{ij}^{\,2}}{2M_{Z'}^2}.
\tag{37.7}
\] Recording both the operator and its conjugate with directed indices
\(i\ne j\) gives \[
\sum_{i\ne j}(|C_{ij}^u|+|C_{ij}^d|)
\ge\frac{|\Delta g|^2}{2M_{Z'}^2}(1-c).
\tag{37.8}
\] If only independent unordered pairs are listed, divide both sides by
two.

This is a relation among coefficients at the matching scale, not a
direct sum rule for measured meson mixing. Running, hadronic matrix
elements, interference with the SM, right-handed couplings and other
heavy-particle contributions require separate calculations; multiple
vectors can cancel. These equations do not predict \(M_{Z'}\) or a
specific decay rate here. They address whether this candidate can leave
no flavour structure at all; they do not identify a new particle already
found in nature.

\section{37.6 Lepton universality is a condition on the full matrix}

If a heavy neutral vector has the same representation coupling to all
three ordinary lepton generations, and their mass rotation remains
unitary after exotic mixing is neglected, its weak-basis coupling must
be \[
G_\ell=g_\ell I.
\] For any mass-basis rotation \(U\), \(U^\dagger G_\ell U=g_\ell I\).
This term therefore generates no tree-level flavour-changing neutral
current through that rotation. It is a useful way to distinguish
candidates, not a conclusion valid for every HGP, flipped-331, or
substantial heavy--light mixing model.

Equal diagonal entries alone are insufficient. In a two-dimensional
subblock, take \(G=\operatorname{diag}(g_1,g_2)\) and rotate with \[
U=\frac1{\sqrt2}\begin{pmatrix}1&1\\1&-1\end{pmatrix}
\] The result is \[
U^\dagger GU=
\begin{pmatrix}
(g_1+g_2)/2&(g_1-g_2)/2\\
(g_1-g_2)/2&(g_1+g_2)/2
\end{pmatrix}.
\] The diagonal entries coincide, but the off-diagonal coupling remains
whenever \(g_1\ne g_2\). Embedding this block in three dimensions gives
the same kind of counterexample. Universality therefore requires testing
that the full matrix is a scalar multiple of the identity, or checking
its spectrum and all relevant entries. A diagonal-only fit cannot
establish the absence of flavour change. Running, operator mixing and
other meson/heavy-vector contributions to low-energy processes still
follow the contract in §37.5.

\chapter{38. Coherence phases impose constraints on the
dynamics}\label{ch-38}

Coherent phases cannot be chosen independently of decay: complete
positivity constrains their joint behaviour. A sharp three-branch bound
quantifies this restriction, and a classical Poisson phase-jump
counterexample shows why even saturation does not identify a new
physical mechanism.

\section{38.1 The minimal physical definition of SRN}

Take a complete family of nonzero orthogonal pointer projections
\(P_a\). Strong record nondisturbance (SRN) requires a channel to fix
every state within each branch, rather than merely preserve the total
population. In finite dimension, the CPTP channels with this property
are precisely \[
\mathcal A_K(\rho)=\sum_{ab}K_{ab}P_a\rho P_b,\qquad
K\succeq0,\quad K_{aa}=1.
\tag{38.1}
\] To prove this, each Kraus operator must act as a scalar on every pure
state fixed pointwise. Linearity forces that scalar to be constant
throughout a branch. Hence \(L_\mu=\sum_a l_{\mu a}P_a\), and
\(K_{ab}=\sum_\mu l_{\mu a}\bar l_{\mu b}\) is a Gram matrix.
Conversely, a Gram decomposition constructs the Kraus operators.

Public readout of a coherent parent system can be completely insensitive
to \(K_{ab},a\ne b\). New measurement dynamics must therefore be probed
through coherence, not solely by counting public records.

\section{38.2 Continuous generators and cycle phases}

A continuous time-homogeneous semigroup has \[
K_{ab}(t)=e^{t\lambda_{ab}},\qquad
\lambda_{ab}=-\gamma_{ab}+i\omega_{ab}.
\] Here \(\gamma_{ab}=\gamma_{ba}\ge0,\omega_{ab}=-\omega_{ba}\), and
the diagonal entries vanish. Complete positivity at all times requires
\(\Lambda\) to be positive semidefinite on the zero-sum subspace.

For three branches, define \[
\Omega_{abc}=\omega_{ab}+\omega_{bc}+\omega_{ca}.
\tag{38.2}
\] Changing branch energies or choosing a diagonal rotating frame linear
in time only adds a difference \(h_a-h_b\) to \(\omega_{ab}\); its cycle
sum is zero. Such a Hamiltonian redefinition therefore cannot remove
\(\Omega\). It is also unaffected by redundancies of the Lindblad
representation, since \(\lambda_{ab}\) is a coefficient of the fixed
physical channel.

This is the same type of Hamiltonian obstruction studied by Fagnola,
Gough, Nurdin and Viola; it is not a first discovery by UEOT.
\href{https://arxiv.org/abs/1811.11784}{Original research}.

\section{38.3 A sharp CP bound for three branches}

\phantomsection\label{m-01}{}

\textbf{M-01.} Every three branches must satisfy \[
\boxed{\ 
\Omega_{abc}^{\,2}
\le
4\gamma_{ac}\gamma_{bc}
-(\gamma_{ac}+\gamma_{bc}-\gamma_{ab})^2.\ }
\tag{38.3}
\] For exactly three branches, this condition together with
\(\gamma\ge0\) for each rate is necessary and sufficient for the
generator to yield CP evolution.

\textbf{Proof.} Anchor at \(c\) and set
\(G_{ij}=\lambda_{ij}-\lambda_{ic}-\lambda_{cj}\). Its nonzero
\(2\times2\) block is \[
G=
\begin{pmatrix}
2\gamma_{ac}&
\gamma_{ac}+\gamma_{bc}-\gamma_{ab}+i\Omega_{abc}\\
\gamma_{ac}+\gamma_{bc}-\gamma_{ab}-i\Omega_{abc}&
2\gamma_{bc}
\end{pmatrix}.
\] Positivity on the zero-sum subspace is equivalent to positive
semidefiniteness of this matrix. Nonnegativity of its second-order
principal minor is exactly (38.3), and the two nonnegative diagonal
entries supply the remaining conditions. Conversely, this positive Gram
matrix constructs
\(\lambda_{ij}=G_{ij}-(G_{ii}+G_{jj})/2+i(\theta_i-\theta_j)\). Schur
positivity shows that its entrywise exponential is a correlation matrix,
giving CPTP evolution for every \(t\ge0\). \(\square\)

With more branches, all triangle conditions remain necessary but do not
replace the higher-order principal-minor conditions of the full Gram
matrix.

\section{38.4 Phase accumulation entails coherence loss}

Write \(S=\gamma_{ab}+\gamma_{bc}+\gamma_{ca}\). The right-hand side can
be expressed as \[
S^2-2(\gamma_{ab}^2+\gamma_{bc}^2+\gamma_{ca}^2)\le S^2/3.
\] Consequently:

\phantomsection\label{m-02}{}

\textbf{M-02.} \[
S\ge\sqrt3\,|\Omega_{abc}|.
\tag{38.4}
\] For the cyclic coherence product
\(B(t)=K_{ab}(t)K_{bc}(t)K_{ca}(t)\), \[
|B(t)|=e^{-St}
\le e^{-\sqrt3|\Omega_{abc}|t}
=e^{-\sqrt3|\arg_{\rm unwrap} B(t)|}.
\tag{38.5}
\] The proof follows directly from (38.3) and the lower bound on the sum
of squares. The phase must be unwrapped continuously; using a principal
value modulo \(2\pi\) cannot justify a spurious long-time conclusion.
\(\square\)

For three coherences at a fixed finite time, complete positivity also
provides a test that requires no Markov assumption: \[
1-|K_{ab}|^2-|K_{bc}|^2-|K_{ca}|^2
+2\operatorname{Re}(K_{ab}K_{bc}K_{ca})\ge0.
\tag{38.6}
\] This is the determinant of the \(3\times3\) correlation matrix. The
semigroup and single-time bounds have different hypotheses. Violating
the former may indicate memory or time-dependent dynamics, rather than a
failure of QM.

\section{38.5 Why a nonzero cycle phase still does not establish quantum
noise}

For real diagonal Hermitian noise operators
\(L_\mu=\operatorname{diag}(\ell_{\mu a}),\ell_{\mu a}\in\mathbb R\), \[
\lambda_{ab}=-\tfrac12\sum_\mu(\ell_{\mu a}-\ell_{\mu b})^2
-i(h_a-h_b),
\] so \(\Omega=0\). Conversely, if every cycle vanishes, \(\omega_{ab}\)
is a branch-potential difference. The remaining anchored Gram matrix is
real and positive, so it can be factored into real vectors, yielding a
Gaussian diffusion representation of this form.

Classical Poisson phase jumps, however, give \[
\mathcal L(\rho)=r(U\rho U^\dagger-\rho),\qquad
U=\operatorname{diag}(e^{i\theta_a}),
\] \[
\lambda_{ab}=r(e^{i(\theta_a-\theta_b)}-1).
\tag{38.7}
\] For \(\theta=(0,2\pi/3,4\pi/3)\), every \(\gamma=3r/2\) and
\(|\Omega|=3\sqrt3r/2\), exactly saturating (38.4).

A nonzero cycle phase therefore excludes a model consisting solely of
diagonal Gaussian white noise plus a Hamiltonian, but even saturation of
the sharp bound can arise from classical Poisson noise. This
counterexample motivates the stronger four-level criterion in the next
chapter.

\chapter{39. A four-level criterion: excluding the entire random-phase
noise class}\label{ch-39}

This chapter turns the question of which noise class is excluded into a
concrete four-level test, with a channel, witness constants, and a
zero-event protocol. Pointer preservation and calibration are essential:
the witness excludes classical random phases, while the constructed
channel remains realizable by a standard quantum environment.

\section{39.1 The question and the minimum dimension}

At a fixed time, arbitrary classical diagonal phase noise can be written
as \[
\mathcal R(\rho)=\mathbb E[U_\theta\rho U_\theta^\dagger],
\quad U_\theta=\operatorname{diag}(z_1,\ldots,z_m),\quad |z_a|=1,
\] and therefore \[
K_{\rm cl}=\mathbb E[zz^\dagger].
\tag{39.1}
\] The phase distribution may have arbitrary correlations and may be
non-Gaussian or non-Markovian; it is not restricted to exponential
decay.

\phantomsection\label{m-03}{}

\textbf{M-03: Four is the minimum number of branches for single-time
separation.} For \(m\le3\), every correlation matrix
\(K\succeq0,K_{aa}=1\) has form (39.1). Even complete single-time SRN
tomography on two or three branches therefore cannot exclude all
classical random-phase noise.

\textbf{Proof.} Write a rank-\(r\) matrix as \(K=VV^\dagger\). If
\(r^2>m\), the space of \(r\times r\) Hermitian matrices contains a
nonzero \(A\) such that \(v_aAv_a^\dagger=0\) for every row. For
sufficiently small \(\epsilon>0\), \(V(I\pm\epsilon A)V^\dagger\) are
two distinct correlation matrices whose mean is \(K\). Thus \(K\) is not
extreme. An extreme point must therefore satisfy \(r^2\le m\). For
\(m\le3\), only \(r=1\) is possible, and unit diagonal forces the form
\(zz^\dagger\). A compact convex set in finite dimension is the convex
hull of its extreme points, proving the claim. \(\square\)

This dimension bound and the four-dimensional matrix below have explicit
precedents, particularly the 2005 work of Buscemi, Chiribella and
D'Ariano. This chapter supplies an explicit analytical witness,
calibration bounds and a counting protocol for use in tests; it does not
rename a known counterexample as a new law of nature.
\href{https://doi.org/10.1103/PhysRevLett.95.090501}{Original paper}.

\section{39.2 A physical channel requiring only a two-dimensional
environment}

Take \[
V=\begin{pmatrix}
1&0\\0&1\\1/\sqrt2&1/\sqrt2\\1/\sqrt2&i/\sqrt2
\end{pmatrix},\qquad K_*=VV^\dagger.
\tag{39.2}
\] Every row has unit norm, so \(K_*\) is a valid correlation matrix of
rank two. Kraus operators can be read off directly: \[
L_1=\operatorname{diag}(1,0,1/\sqrt2,1/\sqrt2),\quad
L_2=\operatorname{diag}(0,1,1/\sqrt2,i/\sqrt2).
\tag{39.3}
\] Since \(\sum L_\mu^\dagger L_\mu=I\), this is a standard CPTP-SRN
channel.

An isometric implementation is \[
|a\rangle|0\rangle_E\longmapsto |a\rangle|e_a\rangle_E,
\] where the two components of \(|e_a\rangle\) are the entries of row
\(a\) of \(V\). The environmental Gram relation is
\(\langle e_b|e_a\rangle=K_{*,ab}\). Four controlled environmental
unitaries can each send \(|0\rangle_E\) to its corresponding
\(|e_a\rangle_E\), completing the joint evolution unitarily. This is a
finite model consisting of a four-level system and one environmental
qubit.

Because \(K_{*,12}=0\), it is not the kernel of a continuous
time-homogeneous SRN semigroup with finite rates at a finite time. This
chapter makes no Markov assumption.

\section{39.3 An analytical separating witness with rigorous constants}

Define \[
B=\begin{pmatrix}
-1/\sqrt2&-1/\sqrt2&1&0\\
-1/\sqrt2&-i/\sqrt2&0&1
\end{pmatrix},\qquad W=B^\dagger B.
\tag{39.4}
\] Since \(BV=0\), \[
\operatorname{Tr}(WK_*)=0.
\tag{39.5}
\]

\phantomsection\label{m-04}{}

\textbf{M-04: A sharp classical lower bound.} \[
\boxed{\operatorname{Tr}(WK_{\rm cl})\ge
\kappa:=3-2\sqrt2=0.171572875\ldots}
\tag{39.6}
\] The bound holds throughout class (39.1), and the constant is
attainable.

\textbf{Proof.} First consider any phase vector \(z\). Put \[
r_1=\left|z_3-\frac{z_1+z_2}{\sqrt2}\right|,\qquad
r_2=\left|z_4-\frac{z_1+iz_2}{\sqrt2}\right|.
\] Then \(z^\dagger Wz=r_1^2+r_2^2\). Write \(w=z_1\bar z_2\), with
\(|w|=1\). The reverse triangle inequality gives \[
|\operatorname{Re}w|
=\left|\left|\frac{z_1+z_2}{\sqrt2}\right|^2-1\right|
\le(\sqrt2+1)r_1,
\] \[
|\operatorname{Im}w|
=\left|\left|\frac{z_1+iz_2}{\sqrt2}\right|^2-1\right|
\le(\sqrt2+1)r_2.
\] Squaring and adding, and using
\((\operatorname{Re}w)^2+(\operatorname{Im}w)^2=1\), yields
\(r_1^2+r_2^2\ge1/(\sqrt2+1)^2=3-2\sqrt2\). Averaging over a probability
distribution proves (39.6). The choice \(z=(1,1,1,(1+i)/\sqrt2)\) gives
\(r_1=\sqrt2-1,r_2=0\) and attains the bound. \(\square\)

Thus no classical random distribution of diagonal phases produces
\(K_*\). Moreover, \[
\|W\|_F=3,\qquad \|W\|_{\rm op}=2+1/\sqrt2=:w_{\max}.
\] By Cauchy--Schwarz, \[
\inf_{K_{\rm cl}}\|K_*-K_{\rm cl}\|_F\ge\kappa/3.
\tag{39.7}
\] This is a proved conservative distance bound, not a conjecture
obtained from an optimizer.

\section{39.4 Tolerance to additional dephasing}

\phantomsection\label{m-05}{}

\textbf{M-05.} For \(K_\epsilon=(1-\epsilon)K_*+\epsilon I\), \[
\operatorname{Tr}(WK_\epsilon)=4\epsilon.
\] The witness therefore still strictly excludes (39.1) whenever \[
\epsilon<\frac{3-2\sqrt2}{4}=0.0428932188\ldots
\tag{39.8}
\]

Here \(I\) is the kernel of complete pointer dephasing, not the kernel
of the identity channel; the latter is the all-ones matrix. Other noise
forms require separate calculations. The value 4.29\% is not a universal
hardware error threshold. Above it, this witness no longer guarantees
detection; that alone does not show that the channel has become
random-unitary.

If a raw tomographic estimate satisfies \(\|\widehat K-K\|_F\le\eta\),
then \[
\operatorname{Tr}(WK)\le \operatorname{Tr}(W\widehat K)+3\eta.
\] An exclusion allowing for error requires this upper bound to lie
below \(\kappa\). A PSD projection cannot replace that confidence
guarantee.

\section{39.5 A zero-event protocol with one input and one effect}

Prepare \(|+\!_4\rangle=(1,1,1,1)^T/2\). The SRN output is
\(\rho_K=K/4\). Choose the valid binary measurement effect \[
E=\frac{W}{w_{\max}},\qquad 0\preceq E\preceq I.
\tag{39.9}
\] Then \[
p_E(K_*)=0,\qquad
p_E(K_{\rm cl})\ge
p_{\rm cl}:=\frac{\kappa}{4w_{\max}}
=0.0158446719\ldots
\tag{39.10}
\] One can measure \(W\) spectrally and accept with probability
\(\lambda/w_{\max}\) conditional on its eigenvalue, or implement the
binary POVM with an ancillary qubit. The complexity has moved into
trusted preparation and effect calibration; it has not disappeared.

\phantomsection\label{e-01}{}

\textbf{E-01: A finite-sample rejection test.} For independent
identically distributed trials with the model and measurement fixed in
advance, the probability of no accepted events in \(n\) trials under any
classical hypothesis is at most \((1-p_{\rm cl})^n\). Thus, with ideal
calibration, zero events in \(n=188\) trials exclude the entire
classical phase class at significance level 0.05.

\textbf{Proof.} The Bernoulli zero-event probability decreases
monotonically with \(p\). Substitute the smallest allowed value
\(p_{\rm cl}\), require \(n\ge\log(0.05)/\log(1-p_{\rm cl})\), and round
up. \(\square\)

If \(k>0\) events are observed, use the prespecified \[
p\text{-value}=\sum_{j=0}^k{n\choose j}p_{\rm cl}^{\,j}(1-p_{\rm cl})^{n-j}.
\tag{39.11}
\] For a conservative error allowance \(\delta\), replace \(p_{\rm cl}\)
with \(\max(0,p_{\rm cl}-\delta)\). The count 188 excludes calibration
trials, SRN validation, drift control and reproducibility checks; it is
not a performance promise for existing devices.

\section{39.6 Ordinary nondiagonal unitaries must not imitate witness
success}

An arbitrary unitary can send the single chosen input into \(\ker W\).
Equation (39.10) is therefore not an unconditional test of all
random-unitary channels. The physical class of pointer-state-preserving
channels must be validated independently.

If a random-unitary channel preserves every \(|a\rangle\langle a|\)
exactly, its positive rank-one summands force every constituent unitary
to be diagonal in this basis. Only then does it belong to (39.1).

\phantomsection\label{e-02}{}

\textbf{E-02: A robust bound for approximate population preservation.}
If a general random-unitary channel
\(\sum_\nu p_\nu U_\nu(\cdot)U_\nu^\dagger\) satisfies
\(\Pr(a|a)\ge1-\ell\) for every basis state, then, for any fixed input,
its output is within trace distance \[
\sqrt{2m\ell}.
\tag{39.12}
\] of the output of some diagonal random-unitary channel. For the
four-level protocol, the event-probability lower bound over this broader
class is therefore \(p_{\rm cl}-\sqrt{8\ell}\), before subtracting
preparation and measurement errors.

\textbf{Proof.} For every \(U_\nu\), choose the diagonal unitary
\(D_{\nu,aa}=U_{\nu,aa}/|U_{\nu,aa}|\), assigning an arbitrary phase
when the entry is zero. \[
\|U_\nu-D_\nu\|_F^2=2\sum_a(1-|U_{\nu,aa}|)
\le2\sum_a(1-|U_{\nu,aa}|^2).
\] Its average over \(\nu\) is at most \(2m\ell\). The trace distance
between pure outputs is bounded by their vector distance; the operator
norm is bounded by the Frobenius norm; Jensen's inequality then bounds
the average. Convexity extends the result to mixed inputs. \(\square\)

This conservative bound exposes a substantial calibration requirement:
99\% population preservation is far from sufficient for the ideal
188-trial conclusion. Even with no additional error, it requires
\(\ell<p_{\rm cl}^2/8\) merely to make the lower bound positive. Tighter
calibration bounds are a legitimate research goal; omitting model
validation is not.

\section{39.7 What the criterion establishes and what remains open}

The criterion resolves a task that single-time two- or three-level SRN
data cannot accomplish: excluding all classical random-phase noise. Four
levels suffice, with an explicit physical channel, a sharp witness, a
noise radius and a finite-data protocol.

It does not exclude ordinary quantum environments. Such an environment
produces \(K_*\) itself. A successful witness shows that the specified
classical noise class is inadequate; it does not show that standard QM
is wrong or that objective collapse has been discovered. This
construction is a candidate tool for quantum-noise spectroscopy and
QCVV, not an established new law of nature.

\chapter{40. Actualization dynamics: from assumed single-branch
selection to a conditional theorem}\label{ch-40}

For an explicitly specified stochastic state equation, can absorption
into spectral branches be proved rather than assumed? This chapter
answers that mathematical question within the stated finite model, while
keeping the choice of equation, strength, record activation, and energy
cost as separate physical obligations.

\section{40.1 First exclude an apparently plausible signalling rule}

Suppose the intuition that stronger superposition should produce faster
actualization motivates the nonlinear single-qubit rule \[
\dot x=-\gamma(1-z^2)x,\quad
\dot y=-\gamma(1-z^2)y,\quad \dot z=0
\tag{40.1}
\] and suppose this rule acts separately on each pure preparation branch
before averaging with the usual Born weights. This is not an equivalent
rewriting of a linear quantum channel.

\phantomsection\label{c-01}{}

\textbf{C-01: An explicit contradiction from remote preparation.} The
same density matrix \(\rho=(I+0.6X)/2\) has two decompositions: \[
\mathcal E_1=\{0.8,|+x\rangle;\ 0.2,|-x\rangle\},
\] \[
\mathcal E_2=\{1/2,\mathbf r_+=(0.6,0,0.8);\
1/2,\mathbf r_-=(0.6,0,-0.8)\}.
\] Applying (40.1) branch by branch gives the averages \[
x_1(t)=0.6e^{-\gamma t},\qquad
x_2(t)=0.6e^{-0.36\gamma t}.
\tag{40.2}
\] They differ.

\textbf{Proof and operational implementation.} Share the pure state \[
|\Psi\rangle=\sqrt{0.8}|0\rangle_A|+x\rangle_B
+\sqrt{0.2}|1\rangle_A|-x\rangle_B.
\] Alice chooses to measure either \(Z_A\) or \(X_A\) without
communicating the outcome. These choices prepare the two decompositions
on Bob's side, both with marginal \(\rho\). After \(\gamma t=1\), Bob
measures \(X\); the difference between the probabilities of its positive
outcome is \[
\Delta p=0.3(e^{-0.36}-e^{-1})=0.0989390655\ldots
\] It is nonzero. Combined with standard remote preparation and the
stated branch-averaging rule, the dynamics would therefore allow
signalling. \(\square\)

This is not an unconditional assertion that every mathematical
nonlinearity entails faster-than-light signalling. The contradiction
uses specified pure-state dynamics, composite states and a probability
interpretation. Original no-signalling theorems address these
constraints.
\href{https://arxiv.org/abs/quant-ph/0102125}{Simon--Bužek--Gisin}.

Directly inserting ORC redundancy, state entropy or branch weights into
a deterministic collapse rate must pass the same
decomposition-independence test. Without a declared composite-system
rule and stochastic structure, an appealing single-system trajectory
does not define viable new physics.

\section{40.2 A norm-preserving stochastic candidate with linear
unconditional dynamics}

In finite dimension, take Hermitian \(A\) and \(\gamma>0\), and use the
Itō equation \[
d|\psi\rangle=
\left[-\frac{i}{\hbar}Hdt
-\frac{\gamma}{2}(A-\bar a)^2dt
+\sqrt\gamma(A-\bar a)dW_t\right]|\psi\rangle,
\quad \bar a=\langle\psi|A|\psi\rangle.
\tag{40.3}
\] Here \(W_t\) is a standard Wiener process under the physical
probability measure. If \(A\) is dimensionless, \(\gamma\) has units of
inverse time; in general, \([\gamma]=[A]^{-2}[t]^{-1}\).

This is a standard class of continuous-measurement/stochastic-reduction
models. The equation itself is not claimed as original. The aim is to
derive fully which parts of the UEOT measurement problem it resolves and
which it leaves open. \href{https://arxiv.org/abs/1301.3626}{Continuous
quantum measurement},
\href{https://arxiv.org/abs/quant-ph/0107153}{original
martingale-reduction research}.

\phantomsection\label{c-02}{}

\textbf{C-02: Norm preservation and a linear ensemble average.} A
normalized initial state remains normalized, and \[
\bar\rho=\mathbb E|\psi\rangle\langle\psi|
\] satisfies \[
\frac{d\bar\rho}{dt}
=-\frac{i}{\hbar}[H,\bar\rho]
-\frac{\gamma}{2}[A,[A,\bar\rho]].
\tag{40.4}
\]

\textbf{Proof.} Apply the Itō product rule to
\(\langle\psi|\psi\rangle\). The real drift is
\(-\gamma\langle(A-\bar a)^2\rangle dt\), canceled by the quadratic
variation. The stochastic term is
\(2\sqrt\gamma\langle A-\bar a\rangle dW=0\), so the norm is unchanged.
The same calculation for the pure-state projection gives \[
d\rho=-i[H,\rho]dt/\hbar-\tfrac\gamma2[A,[A,\rho]]dt
+\sqrt\gamma\{A-\bar a,\rho\}dW.
\] The stochastic integral has an adapted bounded integrand and hence
zero mean, yielding the linear equation (40.4). In finite dimension the
coefficients are smooth and bounded on the unit sphere, providing a
global process. No decomposition-dependent term is silently retained by
the averaging. \(\square\)

For local \(A\), the averaged channel is CPTP and compatible with
ordinary composition, so unconditional statistics do not allow
signalling. Nonlinearity of individual trajectories does not by itself
produce the contradiction in C-01.

\section{40.3 Proving spectral-branch absorption}

Now add the crucial condition \([H,A]=0\). Write \(A=\sum_a aP_a\), with
distinct labels for distinct eigenvalues. The projections \(P_a\) may be
degenerate; the conclusion concerns spectral subspaces.

Define \(p_a(t)=\langle P_a\rangle_t\) and
\(V(t)=\langle A^2\rangle_t-\bar a_t^2\).

\phantomsection\label{c-03}{}

\textbf{C-03: A single spectral branch with Born weights.} Under (40.3)
and the conditions just stated, \[
dp_a=2\sqrt\gamma(a-\bar a)p_a\,dW_t,
\tag{40.5}
\] each \(p_a\) is a bounded martingale, and the weight vector converges
almost surely to a one-hot vector: \[
p_a(t)\longrightarrow 1_{\{A_\infty=a\}},\qquad
\Pr(A_\infty=a)=p_a(0).
\tag{40.6}
\]

\textbf{Full proof.} Taking the trace against \(P_a\) in the
trajectory-density equation, commutation removes the Hamiltonian and
double-commutator drifts and gives (40.5). Each \(p_a\) is bounded, so
the martingale convergence theorem gives almost-sure and \(L^1\)
convergence.

Equation (40.5) implies \(d\bar a=2\sqrt\gamma VdW\). Applying Itō's
rule to \(\langle A^2\rangle-\bar a^2\) gives \[
dV=2\sqrt\gamma\langle(A-\bar a)^3\rangle dW
-4\gamma V^2dt.
\tag{40.7}
\] The bounded nonnegative supermartingale \(V\) therefore converges.
Its mean \(v(t)=\mathbb EV(t)\) satisfies \(v'\le-4\gamma v^2\), hence
\[
\mathbb EV(t)\le\frac{V(0)}{1+4\gamma V(0)t}\longrightarrow0.
\tag{40.8}
\] Bounded convergence shows that the nonnegative limit \(V_\infty\) has
zero mean and is thus zero almost surely. A probability distribution on
finitely many distinct eigenvalues has zero variance exactly when all
weight is concentrated at one eigenvalue. The martingale limit is
consequently one-hot. Finally, \(\mathbb E p_a(\infty)=p_a(0)\) proves
(40.6). \(\square\)

This advances beyond Chapter 3's conditional statement that an unbiased
martingale has Born weights if it ultimately selects one branch: for a
specified dynamics, single-branch selection is now also proved. Core
still does not select these \(A,\gamma\) or this stochastic equation. If
\([H,A]\ne0\), the Hamiltonian can continually drive the state away from
eigenbranches, so the absorption proof cannot be transferred unchanged.

Spectral-branch absorption here is an almost-sure limit as
\(t\to\infty\), not an exact projection at a finite time. Indeed, when
\(p_a(0)>0\), writing \(b_a(t)=2\sqrt\gamma[a-\bar a_t]\) gives \[
p_a(t)=p_a(0)\exp\left[
\int_0^t b_a(s)dW_s-\frac12\int_0^t b_a(s)^2ds\right]>0
\] almost surely at every finite time, because the finite spectrum makes
\(b_a\) bounded. An initially zero weight remains zero, whereas multiple
initially nonzero weights generally separate only in the infinite-time
limit. A finite experiment needs an error or confidence bound of the
following kind; C-03 does not assert exact finite-time collapse.

\section{40.4 A computable finite-time concentration rate}

\phantomsection\label{c-04}{}

\textbf{C-04: Decay of pairwise uncertainty.} Retain the finite-spectrum
and \([H,A]=0\) conditions of §40.3. Then \[
\mathbb E\sqrt{p_a(t)p_b(t)}
=\sqrt{p_a(0)p_b(0)}
 e^{-\frac\gamma2(a-b)^2t}.
\tag{40.9}
\] Consequently, \[
\mathbb E[1-\max_a p_a(t)]
\le\sum_{a<b}\sqrt{p_a(0)p_b(0)}
 e^{-\frac\gamma2(a-b)^2t}.
\tag{40.10}
\]

\textbf{Proof.} Apply (40.5) and Itō's formula to
\(q_{ab}=\sqrt{p_ap_b}\). The difference of the two multiplicative-noise
coefficients is \(2\sqrt\gamma(a-b)\), and the second-order term gives
drift \(-\frac\gamma2(a-b)^2q_{ab}\). The stochastic integral has a
bounded integrand and hence zero mean. Its expectation therefore solves
a first-order linear equation, giving (40.9). Positive initial
probabilities stay positive at finite times; initially zero
probabilities stay zero, so the boundary cases are handled separately.
For any probability vector, let its largest entry be \(p_{\max}\). The
sum of pairwise minima involving this entry already equals
\(1-p_{\max}\). Thus \[
1-p_{\max}\le\sum_{a<b}\min(p_a,p_b)
\le\sum_{a<b}\sqrt{p_ap_b}.
\] Taking expectations gives (40.10). \(\square\)

If \(A\) is degenerate, choose a common eigenbasis \(|a,\alpha\rangle\)
such that \[
A|a,\alpha\rangle=a|a,\alpha\rangle,
\qquad H|a,\alpha\rangle=E_{a\alpha}|a,\alpha\rangle.
\] The unconditional master equation then gives \[
\bar\rho_{a\alpha,b\beta}(t)=
 e^{-i(E_{a\alpha}-E_{b\beta})t/\hbar}
 e^{-\frac\gamma2(a-b)^2t}
\bar\rho_{a\alpha,b\beta}(0).
\] An eigenvalue label of the chosen observable need not specify a
unique energy of the Hamiltonian; coherence within one spectral subspace
may continue to evolve unitarily. The model links selection between
different spectral branches to the rate of nonselective coherence loss.
It does not select a unique pure state within a degenerate subspace.

These weights can be simulated with an exact filtering representation
that requires no Euler time step. First draw a label \(A_*\) according
to \(p_a(0)\), take an independent standard Brownian process \(B_t\) on
the simulation probability space, and generate \[
Y_t=2\sqrt\gamma\,a_{A_*}t+B_t,
\] Bayes' rule gives \[
p_a(t)=
\frac{p_a(0)e^{2\sqrt\gamma aY_t-2\gamma a^2t}}
{\sum_b p_b(0)e^{2\sqrt\gamma bY_t-2\gamma b^2t}}.
\tag{40.11}
\] In the filtration generated by the observation record, the innovation
process \[
W_t=Y_t-2\sqrt\gamma\int_0^t\bar a_s\,ds
\] is Wiener under the physical mixture measure, and the weights satisfy
(40.5). The processes \(B_t\) and \(W_t\) have different roles. Drawing
a label is a simulation representation of the same filtering law; it
does not replace C-03's dynamical absorption proof by assuming the
outcome in advance.

\section{40.5 A new physical candidate must account for its energy cost}

Equation (40.4) and valid cyclicity of the trace give the extra noise
contribution to energy: \[
\frac{d\langle H\rangle}{dt}\bigg|_{\rm noise}
=-\frac\gamma2\langle[A,[A,H]]\rangle.
\tag{40.12}
\] This identity does not itself require \([H,A]=0\). In a
one-dimensional nonrelativistic extension with \(A=X\) and
\(H=P^2/(2m)+V(X)\), one must separately supply common domains for the
unbounded operators, energy conditions on states, and justification for
interchanging the trace and generator. Under those conditions, \[
[X,[X,H]]=-\hbar^2/m,
\qquad\frac{d\langle H\rangle}{dt}=\frac{\gamma\hbar^2}{2m}.
\tag{40.13}
\] For two positions separated by \(d\ne0\), the coherence decay rate
contributed by this additional double-commutator dissipator is
\(\Gamma(d)=\gamma d^2/2\). Therefore \[
\boxed{\frac{d\langle H\rangle}{dt}
=\frac{\hbar^2}{md^2}\Gamma(d).}
\tag{40.14}
\] This is a constraint linking two observables within the same linear
position-noise model, not a universal heating formula for all CSL,
dissipative or coloured-noise models. The \([H,A]=0\) condition from
§40.3 is not used here: position noise for a free particle can cause
continuing heating and cannot simultaneously be described as permanent
absorption into fixed position eigenstates. Exact global CCR do not hold
for finite matrices, so a numerical position truncation cannot inherit
(40.13) without error analysis.

If an independent experiment supplies an upper bound \(\dot E_{\max}\)
on extra heating, this particular model must satisfy
\(\gamma\le2m\dot E_{\max}/\hbar^2\), also constraining the
corresponding additional interference decay. A record-triggered model
that fits a larger decay without the associated energy effect must
explicitly change \(A\), the noise spectrum or the dissipative structure
and submit that new model to testing. The second consequence of the same
noise parameter cannot simply be removed from the contract.

\section{40.6 Record dependence can be introduced through physical
registers}

If a record \(R\) is an actual local register with fixed projections
\(Q_r\), define on the joint space \[
L_r=\sqrt{\gamma_r}\,A\otimes Q_r,\qquad\gamma_r\ge0,
\] and use a standard Lindblad equation. In a branch already classically
conditioned on record value \(r\), the system acquires the corresponding
rate. Averaging over all record outcomes leaves the overall evolution
linear and CPTP.

Here \(\gamma_r\) is a fixed model parameter and \(Q_r\) is an
implementable physical variable. If the register remains in a coherent
superposition, the calculation must retain the joint model rather than
first imposing a classical label. This differs from inserting an
analyst's entropy, redundancy or ``degree of objecthood,'' computed from
an unobserved global state, arbitrarily into a nonlinear single-system
rate.

Independently selecting \(A,Q_r,\gamma_r\) from resources,
self-maintenance, actually accessible records and specified tasks
remains the responsibility of an additional physical bridge. Standard
quantum models with record feedback, fundamental stochastic-field
candidates, and calibrated device control are distinct realizations.
Their comparison requires common responses and experimentally
distinguishable conditions.

\chapter{41. Why coherence loss alone still does not establish objective
collapse}\label{ch-41}

Different system--environment realizations can produce the same system
channel, limiting what system-only statistics can identify about
collapse. This chapter separates exclusion of a noise subclass,
controlled inversion of the full environment, and tests of fundamental
nonunitarity, so that failures are assessed against an explicit model
and control specification.

\section{41.1 Identifiability limits of system-only statistics}

\phantomsection\label{c-05}{}

\textbf{C-05: Equivalence with an unrestricted quantum environment.}
Every finite CPTP channel
\(\mathcal E(\rho)=\sum_\mu L_\mu\rho L_\mu^\dagger\) can be written as
\[
\mathcal E(\rho)=\operatorname{Tr}_E[V\rho V^\dagger],
\quad V=\sum_\mu L_\mu\otimes|\mu\rangle,\quad V^\dagger V=I.
\tag{41.1}
\] Protocols observing only the system therefore cannot uniquely
attribute this channel to fundamental collapse rather than an unobserved
quantum environment.

\textbf{Proof.} Expanding the partial trace gives the Kraus sum. Trace
preservation makes the map an isometry, which can be completed to a
unitary on an extended space. The Born pairing with every system effect
is identical. For finitely many interventions at multiple times, a
sequential environment with memory dilates the causal comb, giving the
same conclusion; the entire process must be matched, not merely its
single-time states. \(\square\)

This obstruction does not make collapse models experimentally
untestable. It means that ordinary-environment hypotheses must be
constrained, or that an additional relation must link systems, energy,
space and time, or parameters shared across different experiments,
before two physical theories can be compared. The unrestricted class of
all finite CPTP explanations already contains every individual such
channel.

\section{41.2 What a complete reversal measures}

If all standard environmental degrees of freedom, writing operations and
dynamics have been independently characterized, and the complete
\(U^\dagger\) can be implemented, the standard model predicts a return
to the initial state. A failed return can reject this complete
model-and-control package, but it can also reflect omitted environmental
degrees of freedom or inversion error. Evidence for a specified
additional noise model must come from distinctive scale dependence or
parameter relations predicted in preregistered controls.

Chapter 39 addresses a narrower, explicit hypothesis test: rejecting all
diagonal classical random-phase models. It does not overcome C-05's
environmental equivalence. Three-branch cycle tests, four-branch
random-unitary tests, inversion of the complete environment, and tests
of fundamental collapse belong to different experimental levels and
should be labelled accordingly.

\section{41.3 Existing experimental constraints on collapse models}

Specified collapse models can be constrained through the additional
force noise, diffusion or radiation that they predict. Vinante et
al.~used low-temperature nanocantilever noise measurements to constrain
a specified CSL model; Donadi et al.~tested their Diósi--Penrose
formulation against underground radiation data; and Altamura et
al.~analyzed existing LISA Pathfinder angular-motion data to obtain CSL
constraints dependent on the model and apparatus geometry. These are
published experiments or analyses of measured data, not this book's
synthetic benchmarks or experimental confirmation of UEOT.
\href{https://doi.org/10.1103/PhysRevLett.116.090402}{Vinante et al.,
2016}, \href{https://doi.org/10.1038/s41567-020-1008-4}{Donadi et al.,
2021}, \href{https://doi.org/10.1103/PhysRevA.111.L020203}{Altamura et
al., 2025}.

An upper bound and a discovery require different evidence.
Conservatively allowing the candidate mechanism to account for all
measured noise can constrain it without asserting its detection. A
stronger limit based on a smaller residual requires an explicit
ordinary-background model, independent calibration and propagated
uncertainty. The LISA study distinguishes these approaches. Noise or
coherence loss alone identifies neither its source nor environmental
parameters that have not been independently measured.

These limits cannot be transferred directly to an undefined UEOT
actualization strength. Chapter 40's position-noise parameter is not
automatically the CSL collapse rate, spatial correlation length, or a
parameter of another model. The actualization operator, noise spectrum,
composite-system rule and testable parameter correspondence must first
be specified. This book has neither refitted UEOT parameters from these
papers nor performed a QDRR hardware experiment. QDRR remains a proposed
test with explicit controls, calibration and failure conditions.

\chapter{42. Executable research protocols and the status of the
arguments}\label{ch-42}

After deriving a candidate curve, one must still specify which
parameters may be fitted, which data are held out, and what would count
as failure. This chapter consolidates experimental, numerical, and
attribution obligations; historical computational receipts are not
relabelled as new executions or experimental confirmations.

\section{42.1 Experimental and fitting contracts}

\begin{Verbatim}[breaklines,breakanywhere,fontsize=\scriptsize]
ExperimentContract:
  model_family:
    classical_diagonal_phase | approximately_pointer_preserving_random_unitary
    | Gaussian_diffusion | calibrated_quantum_environment | specified_collapse
  system: 4 levels, fixed pointer basis, fixed preparation
  channel: duration, control sequence, environment access
  measurement: E=W/(2+1/sqrt(2)), calibration bound
  assumptions:
    conditional per-shot null lower bound (chapter 46), or IID
    no post-hoc choice of witness or exclusion rules
    pointer survival and SPAM calibration data kept separate
  data:
    raw accept/reject counts
    time order, randomization, rejected-shot reasons
    calibration runs and their confidence allocation
  test:
    predeclared n and alpha
    lower null probability after systematic-error allowance
    lower-tail binomial test, no optional stopping without valid adjustment
  outputs:
    rejected family, unresolved larger families
    parameter intervals, calibration failures, reproducibility record
\end{Verbatim}

Using data to choose \(W\) and then computing an unadjusted p-value from
the same data introduces selection bias. One can fix the chapter's \(W\)
in advance or use training and held-out data. Postselection losses must
be represented in the instrument; trials containing an event cannot
simply be removed.

The particle-physics contract fixes the complete representation table,
global group, active mass thresholds, scalar count, beta-function
scheme, \(Z'\) mixing and right-handed couplings. Flavor data require
matching to Wilson coefficients before running and conversion through
matrix elements. No unprovided experimental spectrum, meson dataset or
device dataset has been used as ``validation'' here.

\section{42.2 Independent checks and formalization boundaries of these
advanced derivations}

The appendices preserve independently written numerical/rational
verification programs, fixed inputs and results. The principal checks
compare expressions obtained in different ways:

\begin{itemize}
\tightlist
\item
  Compute all six anomalies from the complete left-handed representation
  table and compare them with the analytical factors.
\item
  Enumerate a finite integer domain and compare it with the all-integer
  classification proved in the text.
\item
  Construct matrices directly from unitary mass rotations and check the
  rank-one sum rules, CKM conjugation and FCNC lower bound.
\item
  Construct generators from complex Lindblad vectors and independently
  compute triangle positivity and Poisson saturation.
\item
  Construct the Kraus operators, environmental dilation,
  four-dimensional \(W\) and actual statistical effect directly.
\item
  Check the sharp witness's saturation point, random phases, mixtures
  and calibration-distance bounds.
\item
  Recheck the decomposition-dependent counterexample through two remote
  measurements on the same joint pure state.
\item
  Use exact QND filtering samples to check martingale means, pairwise
  decay and the finite-time absorption-error bound.
\end{itemize}

Historically retained results and executions in the present edition are
recorded separately, bound to program SHA, inputs, runtime environment
and actual exit status. A success flag in old JSON does not certify a
new execution. Integer/rational equalities, finite-dimensional
floating-point residuals and Monte Carlo coverage are also different
kinds of evidence.

The universal mathematical claims in these advanced sections rest on the
analytical arguments and cited external theorems; sampling and finite
searches do not replace them. The finite Kraus kernel does use Lean and
has separate build, axiom-audit, positive/negative instance checks and
remote-CI receipts. This does not mean that G-01, F-03, M-04, Itō
absorption, Berry/SSH results, history-dependent statistics, or every
other advanced proposition has been formalized. Selectors and
spectral-row contracts with actual source mappings are handled according
to §31.6 and the formalization registry.

These programs presently provide model regression and arithmetic checks,
not new hardware data. A real device must execute the preregistered
physical and statistical contract. Passing a simulation is not empirical
confirmation of a classical-noise, fundamental-collapse, five-colour or
new-particle model.

\section{42.3 Explicit answers about new physics and earlier questions}

\begin{longtable}[]{@{}
  >{\raggedright\arraybackslash}p{(\columnwidth - 4\tabcolsep) * \real{0.2500}}
  >{\raggedright\arraybackslash}p{(\columnwidth - 4\tabcolsep) * \real{0.3400}}
  >{\raggedright\arraybackslash}p{(\columnwidth - 4\tabcolsep) * \real{0.4100}}@{}}
\toprule\noalign{}
\begin{minipage}[b]{\linewidth}\raggedright
Earlier question
\end{minipage} & \begin{minipage}[b]{\linewidth}\raggedright
Progress in these advanced derivations
\end{minipage} & \begin{minipage}[b]{\linewidth}\raggedright
Unresolved responsibility
\end{minipage} \\
\midrule\noalign{}
\endhead
\bottomrule\noalign{}
\endlastfoot
Are three generations and three colours derived unconditionally? &
Correct high-energy counting gives the complete two-candidate
classification and exposes the difference between low- and high-energy
premises. & HGP, representation selection and an independent physical
principle excluding five colours. \\
Does the candidate three-generation parent structure have low-energy
consequences? & Rank-one sum rules, impossibility of simultaneous
up/down flavour alignment, and a total FCNC lower bound. & Heavy masses
and couplings, actual signals and complete competing models. \\
Can coherence cycles be assigned arbitrarily? & A sharp CP phase--decay
bound and a Poisson counterexample. & A nonzero cycle does not establish
fundamentally new physics. \\
Can two quadratures or few levels identify every classical noise model?
& Single-time three-level data remain insufficient; an explicit sharp
four-level witness and finite-sample protocol are provided. & Hardware
execution and experimental comparison with existing methods. \\
Does Born actualization merely assume eventual absorption? & Absorption
and Born weights are proved for a specified stochastic model, with
finite-time bounds. & Whether nature uses this as a fundamental
stochastic mechanism. \\
Does extra dephasing establish collapse? & An unrestricted-environment
equivalence obstruction and a two-observable energy constraint are
given. & Identification of a specified fundamental model using
independent calibration and parameters shared across experiments. \\
\end{longtable}

These are calculable research outputs, but the evidence does not justify
announcing a new law of nature. At UEOT's present stage, substantive
progress consists in reducing an untestable interpretation to fewer,
clearer physical premises and deriving the quantitative consequences
they must satisfy. Those derivations can also disprove a preferred
conjecture.

\section{42.4 Sources and originality boundaries}

These advanced derivations build on the original research programme, the
existing complete quantum treatment and the integrated v2 monograph,
including their SRN, HGP and QFT interfaces. The maintenance work also
checked accessible historical v23--v30.11 texts, two complementary v28.1
sources, the formal Core--Physics contract and v31-series kernel
reports. Appendix C records material identities and the scope actually
obtained. Cached readable text, complete original files and executed
verification are separately identified; none is substituted for another,
and access to every historical conversation is not claimed. External
technical dependencies retain their original-paper or author-source
attribution.

\begin{itemize}
\tightlist
\item
  \textbf{331 representations, intergeneration anomalies and FCNCs:}
  Pisano--Pleitez 1992; Liu 1993; Frampton--Ho--Kephart 2012;
  Fonseca--Hirsch 2016. Links are given in Chapters 36 and 37.
\item
  \textbf{Four-dimensional decoherence outside random-unitary models and
  its precedent matrix:} Buscemi--Chiribella--D'Ariano 2005; see also
  the geometric study by
  \href{https://arxiv.org/abs/1508.04027}{Kayser--Luoma--Strunz}. This
  book claims no priority for those underlying phenomena.
\item
  \textbf{Cycle phases and the Hamiltonian obstruction:} Fagnola et al.,
  2018 preprint, revised in 2019. A dynamically displayed date
  encountered during webpage retrieval is not treated as a new
  publication date.
\item
  \textbf{No-signalling and linear CP dynamics:} Simon--Bužek--Gisin
  2001.
\item
  \textbf{Continuous measurement and stochastic reduction:} the
  martingale work of Adler et al.~and the stochastic-equation work of
  Barchielli. The finite-time expressions here are derived for the
  explicitly stated model; stochastic reduction is not presented as a
  newly proposed mechanism.
\end{itemize}

The specific conservative bound F-03, sharp witness M-04, experimental
assembly E-01/E-02 and their combination receive complete independent
derivations here. The literature search does not prove that these exact
expressions have never appeared elsewhere. A submission would require
further focused prior-art searches, independent checking and
experimental performance comparisons. Deriving something independently
does not by itself establish scholarly priority.

\chapter{43. Global phase, boundary conditions, and topological
quantization}\label{ch-43}

Local amplitude and phase equations do not replace the global boundary
conditions of a wavefunction. The ring with twisted boundary conditions
places holonomy, operator domains, and observable spectra in one
calculation, showing both a valid use of phase closure and why it does
not determine the action scale.

\section{43.1 What a local Madelung representation leaves unspecified}

Chapter 3 uses \(\psi=\sqrt\rho e^{iS/\hbar}\) in regions where
\(\rho>0\). Phases on different patches may differ by constants or gauge
transition functions; in the presence of nodes, \(S\) need not be a
globally single-valued real function. A complete model should specify
the configuration space \(X\), a Hermitian line bundle \(L\to X\), a
connection, the Hilbert space \(L^2(X,L)\), and the domain of a
particular self-adjoint Hamiltonian. Models with the same local
differential expression can still have different spectra because of
their boundaries or holonomy.

For a fixed ordinary complex scalar wavefunction with no nodes on a
closed curve \(C\), single-valuedness implies \[
\oint_C dS=2\pi\hbar n,\qquad n\in\mathbb Z.
\] This constrains the phase winding of a given wavefunction; it does
not construct all position and momentum operators from an integer
condition. For a charged system, the observable mechanical momentum is
\(dS-qA\). Its circulation includes a magnetic-flux term and is not
unconditionally \(2\pi\hbar n\).

\section{43.2 The full spectrum on a ring with twisted boundary
conditions}

Take \(X=[0,L]\) with its endpoints identified, mass \(m>0\), and a
constant tangential potential \(A\). On \(L^2([0,L])\), define \[
P_\theta=-i\hbar\partial_x,
\quad D(P_\theta)=\{f\in H^1:f(L)=e^{i\theta}f(0)\}.
\] The boundary form is \(-i\hbar[\bar f g]_0^L\), and this domain makes
it vanish. Requiring the boundary form to vanish against every \(f\)
satisfying the boundary condition imposes the same condition on the
adjoint domain. Thus the domain is self-adjoint, not merely formally
symmetric.

\phantomsection\label{p-ugl-01}{}

\textbf{P-UGL-01: Ring spectrum and gauge invariance.} \[
f_n(x)=L^{-1/2}e^{i(2\pi n+\theta)x/L},\quad
E_n=\frac{[\hbar(2\pi n+\theta)/L-qA]^2}{2m}.
\] The domain of \(H=(P_\theta-qA)^2/(2m)\) consists of \(H^2\)
functions whose values and first derivatives satisfy the same twisted
boundary condition. Multiplying the ordinary Fourier basis by
\(e^{i\theta x/L}\) produces this basis, proving completeness. The
spectrum depends only on \[
\eta=\theta-\frac{qAL}{\hbar}\pmod {2\pi}.
\] For a gauge transformation that preserves the constant-potential
representation, with \(A'=A+\partial_x\chi\) and
\(f'=e^{iq\chi/\hbar}f\), one has
\(\theta'=\theta+q[\chi(L)-\chi(0)]/\hbar\), so \(\eta\) is invariant.
The shift \(\theta\mapsto\theta+2\pi\) merely relabels \(n\). If a
periodic domain was fixed initially, changing a boundary label must be
accompanied by transporting the operator domain in the Hilbert space.
\(\square\)

For example, set \(\hbar=m=L=q=1\). At \(\theta=0\), \(A=0\), the
ground-state energy is zero. At \(A=\pi\), the lowest energy is
\(\pi^2/2\) and is twofold degenerate. If \(\theta\) is also changed to
\(\pi\), the spectrum returns to that for \(\eta=0\). This illustrates
gauge transport; it does not imply that magnetic flux has no physical
effect.

\section{43.3 Refining the original phase-closure route}

Phase closure should be separated into four ingredients: local curvature
of the connection, global holonomy, transition functions of the
wavefunction bundle, and operator domains. The original closed-curve
intuition now has a calculable realization, but the quantum action
scale, complex state space and physical choice of boundary conditions
remain inputs. Core can receive gauge-invariant spectra and responses to
allowed interference protocols as a response object; naming the
protocols does not derive the numerical value of \(\hbar\).

\chapter{44. Topological quantum objects: Berry bundles and SSH edge
states}\label{ch-44}

Topological protection becomes an inspectable physical property only
with a specified Hamiltonian, symmetry, and boundary termination. The
Berry and SSH examples derive geometric quantities, spectra, and edge
states while retaining finite-chain errors, lifetimes, and
symmetry-breaking perturbations in the certificate.

\section{44.1 A complete two-level Berry-curvature calculation}

Take \(H(\mathbf n)=\Delta\mathbf n\cdot\boldsymbol\sigma\), with
\(\Delta>0\) and \(\mathbf n\in S^2\), and consider the instantaneous
eigenline of energy \(+\Delta\). In spherical coordinates, a
northern-patch choice is \[
|u_N\rangle=(\cos(\vartheta/2),e^{i\varphi}\sin(\vartheta/2))^T.
\] Define the connection by \(\mathcal A=i\langle u|du\rangle\). On the
southern patch take \(u_S=e^{-i\varphi}u_N\). Then \[
\mathcal A_N=-\frac{1-\cos\vartheta}{2}d\varphi,
\quad \mathcal A_S=\frac{1+\cos\vartheta}{2}d\varphi.
\]

\phantomsection\label{p-utp-01}{}

\textbf{P-UTP-01: Chern number of this eigenline.} Direct
differentiation gives \[
\mathcal F=-\tfrac12\sin\vartheta\,d\vartheta\wedge d\varphi,
\qquad C_1=\frac1{2\pi}\int_{S^2}\mathcal F=-1.
\] The difference between the local connections,
\(\mathcal A_S-\mathcal A_N=d\varphi\), agrees with the transition
function. The nonzero total flux prevents an everywhere smooth choice of
eigenvector over the sphere. The negative-energy band has the opposite
sign. The sign here depends on the explicitly stated connection, band
and orientation of the sphere. \(\square\)

The geometric phase in adiabatic evolution is \(\oint\mathcal A\); the
dynamical phase must be separated from the total phase. Finite-speed
errors require separate control in terms of the gap, path regularity and
speed. This example uses Berry's established geometric-phase structure;
the calculation here is not its first discovery.
\href{https://doi.org/10.1098/rspa.1984.0023}{Berry's original paper}.

\section{44.2 SSH bulk spectrum and winding number}

Fix the unit cell, boundary termination, and real positive hopping
amplitudes \(t_1,t_2\). The single-particle Hamiltonian is \[
H=\sum_{n\ge1}(t_1|n,A\rangle\langle n,B|
+t_2|n+1,A\rangle\langle n,B|+\mathrm{h.c.}).
\] For the infinite translation-invariant system, the chosen Fourier
convention gives the Bloch matrix \[
H(k)=\begin{pmatrix}0&t_1+t_2e^{-ik}\\t_1+t_2e^{ik}&0\end{pmatrix},
\quad E_\pm(k)=\pm\sqrt{t_1^2+t_2^2+2t_1t_2\cos k}.
\] The chiral operator \(\Gamma=\operatorname{diag}(1,-1)\) satisfies
\(\Gamma H\Gamma=-H\). With \(q(k)=t_1+t_2e^{-ik}\), the winding number
\[
\nu=\frac{1}{2\pi i}\int_{-\pi}^{\pi}q^{-1}\partial_kq\,dk
\] is \(-1\) for \(t_2>t_1\) and zero for \(t_2<t_1\). Indeed, \(q(k)\)
traces a clockwise circle of radius \(t_2\) centred at \(t_1\). At
\(t_1=t_2\) the gap closes and the winding number is undefined. Changing
the Fourier or sublattice convention can reverse the winding sign.
Reselecting the unit cell must be accompanied by the corresponding
change in boundary termination; changing a number alone does not remove
the physical edge state.

\section{44.3 A semi-infinite boundary and finite-chain errors}

\phantomsection\label{p-utp-02}{}

\textbf{P-UTP-02: A semi-infinite zero mode.} For a chain terminated on
the left at \(A_1\), if \(t_1<t_2\), \[
|\psi_L\rangle=\sqrt{1-r^2}\sum_{n=1}^\infty(-r)^{n-1}|n,A\rangle,
\qquad r=t_1/t_2
\] is a normalized zero-energy eigenstate. To prove this, every \(B_n\)
component requires \(t_1a_n+t_2a_{n+1}=0\). The geometric series
converges exactly when \(r<1\), with the normalization factor shown.
Every \(A_n\) equation holds automatically because \(b_n=0\).
\(\square\)

On a finite chain of \(N\) unit cells, truncate and normalize this
vector to obtain \(\psi_{L,N}\), with coefficient
\(c_N=\sqrt{(1-r^2)/(1-r^{2N})}\). Only the final \(B_N\) component has
a residual, so \[
\|H_N\psi_{L,N}\|=c_Nt_1r^{N-1}.
\] The spectral theorem gives
\(\operatorname{dist}(0,\operatorname{spec}H_N)\le\|H_N\psi_{L,N}\|\). A
near-zero mode in a finite chain with equal sublattice sizes is not
generally an exact zero mode: coupling between its ends produces a small
energy splitting. Sublattice imbalance and special boundaries require
separate treatment. The established SSH model originates in the
\href{https://link.aps.org/doi/10.1103/PhysRevLett.42.1698}{Su--Schrieffer--Heeger
paper}.

\section{44.4 The UEOT interface and possible experiments}

Take the boundary subspace as a candidate carrier, with protocols for
local injection, finite-time propagation, boundary survival and
perturbation response. A certificate should record the bulk gap, degree
of chiral-symmetry preservation, boundary termination, finite size and
lifetime. A localized peak alone is not a topological certificate.

A practical sequence is to verify the residual by exact diagonalization
on a small chain, vary \(t_1/t_2\) to observe gap closing and the
boundary distribution, and then introduce known symmetry-preserving or
symmetry-breaking perturbations and compare preregistered spectral and
survival curves. The original paper's idea of a persistent structure
becomes a concrete question about a protected low-energy subspace. It
does not directly derive the number of electron generations, fundamental
particles or an absolutely everlasting object.

\chapter{45. A nonperturbative starting point: finite lattice gauge
models and interfaces across scales}\label{ch-45}

A finite lattice provides a well-defined, computable starting point for
nonperturbative work, but existence of its measure does not construct
continuum QCD. This chapter establishes gauge invariance, action
normalization, and an enumerable example before separating the
obligations for lattice spacing, volume, fermions, and real-time
extrapolation.

\section{45.1 An exact definition of a finite pure-gauge model}

Take a finite four-dimensional Euclidean hypercubic lattice, with
\(U_{x\mu}\in SU(N)\) on oriented links and \(U^\dagger\) on the reverse
links. The oriented plaquette is \[
U_p=U_{x\mu}U_{x+\hat\mu,\nu}U^\dagger_{x+\hat\nu,\mu}U^\dagger_{x\nu}.
\] Define the Wilson action and normalized measure by \[
\begin{gathered}
S_W=\beta\sum_p\left(1-\frac1N\operatorname{ReTr}U_p\right),
\\ Z=\int\prod_l dU_l\,e^{-S_W},\\ d\mu=Z^{-1}e^{-S_W}\prod_l dU_l,
\\\beta=\frac{2N}{g^2}>0.
\end{gathered}
\] Here \(dU_l\) is normalized Haar measure, and one orientation is
chosen for each unoriented plaquette. There is no fermion determinant,
chemical potential or real-time phase factor in this section.

\phantomsection\label{p-ulg-01}{}

\textbf{P-ULG-01: Existence and gauge invariance of the finite model.}
On the finite product of compact groups, \(0\le S_W\le2\beta N_p\),
hence \(e^{-2\beta N_p}\le Z\le1\), so the measure is well defined. A
local transformation \(U_{xy}\mapsto G_xU_{xy}G_y^\dagger\) conjugates
each closed-loop product at its starting point. The trace and left/right
Haar measure are invariant, so all Wilson-loop expectations are gauge
invariant. \(\square\)

This existence statement requires no global gauge choice. It establishes
only a finite Euclidean statistical model. Reconstructing a
positive-energy real-time quantum theory requires reflection positivity,
continuity and suitable limits, among other conditions. The positivity
condition \(e^{-S}>0\) alone does not prove them all. Wilson's original
lattice construction and strong-coupling study appear in
\href{https://link.aps.org/doi/10.1103/PhysRevD.10.2445}{\emph{Confinement
of quarks}}.

\section{45.2 Checking the continuum gauge-action coefficient}

Use Chapter 29's convention \(D_\mu=\partial_\mu+igA_\mu\), with
Hermitian generators normalized by
\(\operatorname{Tr}T^aT^b=\delta^{ab}/2\). Then \[
F_{\mu\nu}=\partial_\mu A_\nu-\partial_\nu A_\mu+ig[A_\mu,A_\nu].
\] For a smooth field, regard link variables as path-ordered
exponentials consistent with this transport convention. The local
expansion gives \[
U_p=\exp\{iga^2F_{\mu\nu}+O(a^3)\},\quad
1-\frac1N\operatorname{ReTr}U_p
=\frac{g^2a^4}{2N}\operatorname{Tr}F_{\mu\nu}^2+O(a^5).
\] The linear term vanishes because \(F\) is traceless, and the
quadratic term comes from the exponential expansion. For a smooth field
at fixed volume, with controlled boundaries, multiplication by \(\beta\)
and summation give \[
S_W\longrightarrow\int d^4x\sum_{\mu<\nu}\operatorname{Tr}F_{\mu\nu}^2
=\frac14\int d^4x\sum_{\mu,\nu,a}(F_{\mu\nu}^a)^2.
\] This establishes consistency of the action on classical smooth
fields. Typical quantum fields are not smooth, so it does not establish
a continuum limit of the quantum measure. Symmetric summation permits
the study of higher-order lattice-spacing improvements, but boundary and
smoothness assumptions must be specified. The expansion order is not
automatically a strong-coupling quantum error bound.

\section{45.3 An exactly enumerable gauge-invariant example}

To make the finite computation explicit, take the \(\mathbb Z_2\)
submodel on the four links of a single square: \(u_l=\pm1\), plaquette
\(P=\prod_{l=1}^4u_l\), and action \(S=-\beta P\). With unnormalized
counting measure, there are 16 configurations, eight each with
\(P=\pm1\).

\phantomsection\label{p-ulg-02}{}

\textbf{P-ULG-02: Exact partition function and response.} \[
Z=16\cosh\beta,\qquad \langle P\rangle=\tanh\beta,
\qquad\partial_\beta\langle P\rangle=1-\langle P\rangle^2.
\] Summing the two eightfold sectors proves the expressions. A vertex
transformation multiplies each of its two adjacent links by the same
\(g_x=\pm1\), so \(P\) is unchanged. Every single-link expectation is
zero, as follows by flipping a transformation at one endpoint and using
invariance of the measure. Normalized Haar counting divides \(Z\) by 16
but leaves normalized expectations unchanged. \(\square\)

This gives a concrete UEOT readout: if all observables \(f(P)\) are
allowed, \(P\) is a complete carrier of invariant information. Orbit
completeness follows by successively gauge-fixing the first three links
to \(+1\); the final link is then exactly \(P\). Configurations with the
same \(P\) therefore lie on the same gauge orbit. This example concerns
only the stated finite group and single loop; it is neither \(SU(3)\)
QCD nor a proof of four-dimensional confinement.

\section{45.4 Lattice coarse-graining and a valid connection to Core}

Define each coarse link as an ordered product along a fine path,
obtaining a map \(B:U_f\mapsto U_c\). The coarse measure in the finite
model is directly defined as the pushforward \(\mu_c=B_*\mu_f\). For
every bounded coarse observable, \[
\mathbb E_{\mu_c}f(U_c)=\mathbb E_{\mu_f}f(B(U_f)).
\] This is the defining identity of the pushforward measure and gives
exact readout consistency. It does not imply that the coarse action
retains only single-plaquette terms: multi-loop and nonlocal terms
generally appear. Truncation to a finite operator family must report the
error in discarded responses rather than assert closure in the same
functional form.

A deterministic map of Euclidean configurations is classical statistical
coarse-graining; it is not automatically a real-time CPTP channel.
Connecting it to the quantum-process interface of Chapter 26 requires a
transfer matrix or controlled real-time model, an initial state and the
same intervention slots, after which responses can be compared. Three
distinct questions can then be studied: exact readout pushforward,
approximate effective actions, and quantum-intervention coarse-graining.

\section{45.5 Explicit tasks toward QCD}

First validate local updates, gauge invariance, thermalization and
autocorrelation on finite \(SU(2)\)/\(SU(3)\) lattices. Then add a
specified fermion discretization and sea-quark weight, checking the
treatment of doubling, chirality and anomalies. Finally, extrapolate the
same dimensionless observable jointly in volume, lattice spacing, mass
and statistical error. Noisy Euclidean data do not provide an
unconditionally stable inversion to real-time response. Neither
existence of a finite pure-gauge measure nor local one-loop asymptotic
freedom proves the mass gap of four-dimensional continuum Yang--Mills
theory.

\chapter{46. Experimental statistics with history dependence and
optional stopping}\label{ch-46}

Device drift and history-dependent choices can invalidate a simple
independent-trial analysis. Conditional null hypotheses and
supermartingales incorporate temporal order and optional stopping into
the statistical guarantee; inadequate calibration can make the resulting
bound explicitly powerless.

\section{46.1 Replacing independent trials with a conditional null
hypothesis}

With the input and effect fixed, the classical diagonal random-phase
model in Chapter 39 gives the four-level witness's event-probability
lower bound \[
p_* =\frac{3-2\sqrt2}{4(2+1/\sqrt2)}.
\] Let \(X_t\in\{0,1\}\) denote the event on trial \(t\), and let
\(\mathcal F_{t-1}\) contain the entire preceding record and allowed
feedback. The new null hypothesis is explicitly \[
\Pr(X_t=1\mid\mathcal F_{t-1})\ge p_0\quad\text{almost surely, for every }t.
\] Here \(0<p_0<1\), and \(p_0\) already accounts for calibration,
leakage and identifiable systematic errors. The condition holds if every
history-conditioned channel belongs to the tested class. Knowing only
that time-averaged noise belongs to that class is insufficient.

\phantomsection\label{p-ust-01}{}

\textbf{P-UST-01: A fixed-sample test for correlated trials.} Under
these conditions, for fixed \(n\), \[
\Pr\left(\sum_{t=1}^nX_t\le k\right)
\le\sum_{j=0}^k{n\choose j}p_0^j(1-p_0)^{n-j}.
\] A sequential coupling proves the result. On an enlarged probability
space, take independent \(U_t\sim\mathrm{Unif}[0,1]\) and realize
\(X_t=1\{U_t\le p_t\}\) with the corresponding history-conditioned
probability \(p_t\). Set \(Y_t=1\{U_t\le p_0\}\). Pathwise,
\(X_t\ge Y_t\), while the \(Y_t\) are independent Bernoulli variables,
giving lower-tail domination. Additional history variables are generated
sequentially using their conditional kernels; the marginal joint law of
the binary trials is unchanged. \(\square\)

The zero-event case also follows directly by iterated conditional
expectation, giving \((1-p_0)^n\). Thus, with ideal \(p_0=p_*\), 188
trials with zero accepted events still provide a fixed-sample 5\% test,
even with history dependence satisfying the conditional lower bound. The
conclusion does not cover arbitrary drift or memory.

\section{46.2 Why single-trial marginal bounds are insufficient}

Let all trials equal the same hidden Bernoulli variable \(Z\), with
\(\Pr(Z=1)=p_0\). Each trial has marginal probability \(p_0\), but the
probability that all trials are zero remains \(1-p_0\), however many
trials are taken; it does not decay as \((1-p_0)^n\). After the first
zero, the next conditional event probability is zero, violating
precisely the condition in §46.1. This counterexample must accompany the
statistical certificate.

\section{46.3 A rigorous version allowing continuous inspection}

For a prefixed \(q\in[0,p_0)\), define \[
L_n=\prod_{t=1}^n
\left(\frac q{p_0}\right)^{X_t}
\left(\frac{1-q}{1-p_0}\right)^{1-X_t},\quad L_0=1.
\]

\phantomsection\label{p-ust-02}{}

\textbf{P-UST-02: An optional-stopping certificate.} Under the null
hypothesis, \(L_n\) is a nonnegative supermartingale, because the
conditional expectation of its one-step factor is \[
1+\frac{(p_t-p_0)(q-p_0)}{p_0(1-p_0)}\le1.
\] Let \(\tau=\inf\{n:L_n\ge1/\alpha\}\). Applying iterated conditional
expectation to the nonnegative supermartingale at the truncated stopping
time \(\tau\wedge N\) gives
\(1\ge E L_{\tau\wedge N}\ge\alpha^{-1}\Pr(\tau\le N)\). Letting
\(N\to\infty\) yields \(\Pr(\tau<\infty)\le\alpha\). This permits
inspection at any time and stopping at the stated threshold while
retaining a valid test. \(\square\)

For \(q=0\), the no-event factor is defined as \(1/(1-p_0)\), and the
entire product becomes zero as soon as an event occurs. This choice is
best suited to an explicit zero-event target. To allow a few events, fix
\(q>0\) or mix several \(q\) values with weights fixed in advance.
Maximizing \(L_n(q)\) after seeing the data requires an adjustment.

\section{46.4 Calibration, leakage and the physical contract}

Suppose a simultaneous calibration event with probability at least
\(1-\alpha_{\rm cal}\) guarantees the lower bound for every trial used.
The total error probability is then at most
\(\alpha_{\rm cal}+\alpha_{\rm test}\). The null must be shown to remain
valid under that good event and the relevant conditioning information,
or one can use an independently calibrated, frozen confidence set of
physically consistent models. A calibration event selected afterward
cannot simply be treated as independent of the test. Taking the worst
lower bound over the entire confidence set is more direct.

Execution should fix \(W\), preparation, pointer basis, randomization,
failure rate and stopping rule, and preserve the time series. Chapter
39's approximate pointer-preservation bound
\(p_0\ge p_*-\sqrt{8\ell}-\delta_{\rm SPAM}\) is very conservative. If
the right-hand side is nonpositive, this witness does not establish
exclusion capability for the current experiment. A useful improvement
would optimize over a complete calibration confidence set with certified
bounds, rather than arbitrarily shrink error bars.

This chapter removes the IID requirement under explicit conditions and
provides a test valid under optional stopping. No hardware experiment
has been performed, and no new decoherence mechanism in nature has been
established.

\chapter{47. Research programme: from closed modules to distinguishable
physics}\label{ch-47}

A research roadmap should be organized around propositions or
experiments that can be delivered, checked, and rejected, rather than
around version numbers. This chapter specifies inputs, outputs, and
failure conditions across quantum selection, dynamics, fields,
particles, and gravity, distinguishing completed finite results from
larger open goals.

\section{47.1 Dependencies and priorities}

The research programme advances according to testable dependencies,
without promising dates for solving open problems. P0 denotes
foundations required before any new-physics claim. P1 denotes directions
already supported by theorems or calculable models and most likely to
produce definite outputs. P2 denotes longer-term directions requiring a
new physical selection principle or a substantial mathematical advance.

\begin{longtable}[]{@{}
  >{\raggedright\arraybackslash}p{(\columnwidth - 6\tabcolsep) * \real{0.2100}}
  >{\raggedright\arraybackslash}p{(\columnwidth - 6\tabcolsep) * \real{0.2000}}
  >{\raggedright\arraybackslash}p{(\columnwidth - 6\tabcolsep) * \real{0.2800}}
  >{\raggedright\arraybackslash}p{(\columnwidth - 6\tabcolsep) * \real{0.3100}}@{}}
\toprule\noalign{}
\begin{minipage}[b]{\linewidth}\raggedright
Work package
\end{minipage} & \begin{minipage}[b]{\linewidth}\raggedright
Priority
\end{minipage} & \begin{minipage}[b]{\linewidth}\raggedright
Prerequisites
\end{minipage} & \begin{minipage}[b]{\linewidth}\raggedright
Deliverable
\end{minipage} \\
\midrule\noalign{}
\endhead
\bottomrule\noalign{}
\endlastfoot
R1 Quantum selection and minimal premises & P0 & OPT background;
classical and real-quantum countermodels & Dependency graph, individual
assumption-deletion certificates, and an auditable reconstruction
chain \\
R2 Global phase and continuous domains & P0 & Self-adjoint domains, line
bundles and energy bounds & Models with explicit boundaries and
convergence certificates \\
R3 Record witnesses and QCVV & P1 & SRN, witnesses, calibration and
statistical contracts & Valid confidence bounds, preregistered
experiments and raw data \\
R4 Actualization and new dynamics & P1/P2 & CP, no-signalling,
absorption and heating constraints & Excludable models and predictions
sharing parameters across observables \\
R5 Memory processes and object compression & P1 & Combs, intervention
sufficiency and resource constraints & Prediction-compression errors and
upper/lower memory-dimension bounds \\
R6 Many-body topology and control & P1 & Finite models, spectral gaps,
stability and noise & Reproducible edge-state lifetimes, logical errors
and control gains \\
R7 Precise QFT/EFT interfaces & P0/P1 & Actions, regulators, schemes and
infrared definitions & Scattering and matching benchmarks with order and
error statements \\
R8 Lattice and nonperturbative physics & P1/P2 & Finite gauge measures
and update algorithms & Finite-volume observables and controlled
extrapolation; continuum limits remain a separate goal \\
R9 Gauge and generation selection & P2 & Complete field tables,
thresholds, global groups and countermodels & An independent mechanism
excluding five colours, or an explicit set of surviving candidates \\
R10 Flavor, masses and neutrinos & P1/P2 & Symmetry-breaking spectra,
rank-one structure and EFT matching & Held-out sum rules, correlations
between processes and competing-model comparisons \\
R11 Global anomalies and geometry & P2 & Specified bundles, operators
and completion & Calculable anomaly maps; retention or withdrawal of
conditional framing conclusions \\
R12 Gravity and cosmology interfaces & P2 & Quantum fields, locality,
energy and geometric inputs & Controlled-background models and
consistency checks; no claim of completed gravitational unification \\
R13 Proof and software reproducibility & P0 & The book's status IDs and
code & Type checking, counterexample collections, a small formal core
and version tracking \\
\end{longtable}

The first stage fixes definitions and benchmarks for R1/R2/R7/R13. The
second prioritizes parallel finite-model work on R3/R5/R6/R8 that does
not require new fundamental axioms. Only then do the candidates in
R4/R9/R10/R11 compete with experiments and mathematical counterexamples.
R12 begins as interface research; a unification goal cannot bypass
preceding unresolved steps. Parallelism here describes dependencies
between research tasks. Every result still requires a concrete proof,
source files and independent review.

\section{47.2 R1: Quantum selection and minimality}

\textbf{Goal.} Distinguish UEOT's ability to organize quantum theory
from a necessity for UEOT to select quantum theory, and seek reductions
in the genuinely independent quantum-domain inputs.

\textbf{Mathematical starting point.} Form a background package
\(\mathcal B\) from the objects, tests, composition, equivalence, convex
closure and finite-dimensional conditions of a complete causal OPT.
Choose a set of selection hypotheses \(\mathcal H\), and state the
conclusion as
\(\mathcal B\cup\mathcal H\Rightarrow\mathrm{CQM}_{\rm fd}\). Every
change of background counts toward the cost of minimality. Chapter 2's
external reconstruction theorem is not a general Core theorem.

\textbf{Steps.} First compare each lemma with the original
reconstruction paper and tabulate the background it actually uses. Next
perform deletion tests in classical, real-quantum,
restricted-effect/transformation quantum subtheories and other concrete
OPTs. Check whether a quantum-looking state space also supplies all
quantum processes and composition rules. Finally investigate whether
UEOT's original resource/access conditions independently imply a domain
principle. Every proposed implication must also face an attempt to
construct a classical countermodel.

\textbf{Acceptance.} Produce a noncircular proof chain traceable to its
basic quantifiers, or an explicit counterexample showing that the
proposed weaker package is insufficient. A counterexample is a valid
deliverable. Renaming purification ``parent-object completeness'' is not
progress; hiding the target quantum structure in the meaning of
``realizable'' is likewise insufficient. If a revision of the cited
preprint changes its background, update the K dependency instead of
silently retaining a stronger conclusion.

\section{47.3 R2: Global phase, operator domains and numerical
convergence}

\textbf{Goal.} Extend Chapter 3's local statistical action safely to
actual boundaries, nodes and multiply connected configuration spaces.

\textbf{Steps.} Begin with Chapter 43's ring and a punctured plane. Fix
the line bundle and connection, then enumerate the Hamiltonian's
self-adjoint extensions. Use ring spectra, flux periodicity and
interference phases as benchmarks. At \(\rho=0\), use weak-solution or
wavefunction-domain formulations; do not divide by zero density.
Numerical truncation must control the initial-state tail, generator
approximation and time integration together. Small
\(D(\rho,\rho_{\rm cut})\) alone does not establish a small long-time
dynamical error.

If \(\|H-\widetilde H\|\le\delta_H\) on a common invariant finite
domain, the Duhamel identity \[
U(t)-\widetilde U(t)
=-\frac{i}{\hbar}\int_0^tU(t-s)(H-\widetilde H)\widetilde U(s)\,ds
\] gives \(\|U-\widetilde U\|\le |t|\delta_H/\hbar\). For unbounded
problems, estimate the integrated norms on the states being used; a
nonexistent global operator norm cannot be retained. This is a minimal
prototype for a dynamical-error contract.

\textbf{Acceptance.} Spectra and responses agree before and after
consistent boundary transport; grid or energy-truncation errors converge
over the specified time window; and the treatment of global nodes no
longer relies on singular Madelung division. If distinct self-adjoint
extensions predict different outcomes without an independent physical
selection, retain several models rather than choose the spectrum closest
to the desired answer.

\section{47.4 R3: SRN witnesses, QDRR and experimental statistics}

\textbf{Goal.} Turn the question of coherence loss associated with
records into experiments that exclude explicit model classes, without
jumping directly to ontological collapse.

\textbf{Order of candidate tests.} Test the full SRN class first, then
time-homogeneous semigroups. Next use three-branch cycles to distinguish
Gaussian diagonal diffusion from broader classical jumps, and finally
use a four-branch witness to exclude the entire classical diagonal
random-phase class. Each level leaves a broader class unexcluded.
Chapter 39's \(K_*\) has zero entries, so its exact channel is not a
finite-time member of a homogeneous SRN semigroup with finite rates.
Implement it through a controlled environmental channel rather than fit
it incorrectly to that semigroup.

\textbf{Steps.} Prepare \(|+_4\rangle\) in a known four-level system or
a two-qubit encoding; implement the specified Kraus dilation with an
environmental qubit; measure \(E=W/(2+1/\sqrt2)\); independently
calibrate pointer survival and SPAM and freeze \(p_0\); acquire
time-ordered data using Chapter 46's fixed-sample or supermartingale
procedure; and reserve different storage times and environment controls
for model checks. Leakage must be recorded as an output or probability
loss, not discarded to improve the signal.

\textbf{Optimization.} Let \(\mathcal C(D_{\rm cal})\) be a
channel/instrument set with a coverage guarantee, and compute \[
p_{\min}=\inf_{M\in\mathcal C(D_{\rm cal})\cap\mathfrak M_0}p_M(E).
\] The convex hull of random phases does not automatically have an exact
small SDP representation. Certified outer approximations can provide
conservative lower bounds. Rejection requires that this lower bound
exceed the measured held-out upper bound; one local minimum returned by
an optimizer is insufficient.

\textbf{Acceptance and stopping.} Calibrate error rates on simulated
null, drifting and approximately SRN data. A real experiment must report
raw counts and worst-case errors together. If \(p_{\min}\le0\), or if
preparation and the effect cannot be independently calibrated, improve
the experimental contract before publishing a claim that classical noise
has been excluded. Recalculate the sample size from calibrated \(p_0\);
188 is not a universal hardware promise.

\section{47.5 R4: Actualization and new dynamics}

\textbf{Goal.} Develop Chapter 40's standard stochastic candidate into
physical predictions that differ from a calibrated environment and share
parameters across observables.

\textbf{Fix the model first.} Specify the probability space, observable
\(A\), rate \(\gamma\), temporal and spatial correlations,
composite-system rule, energy budget, and how records enter the
dynamics. The finite-dimensional QND model is a closed starting point.
It can represent a continuous-measurement trajectory or a fundamental
stochastic hypothesis; identical unconditional channels do not
experimentally distinguish those interpretations.

\textbf{Steps.} Prove norm preservation, averaged CP/TP and local
no-signalling; then determine absorption or stationary behaviour. When
\([H,A]\ne0\), repeat the analysis rather than import the QND theorem.
Derive common relations among heating, momentum diffusion, interference
contrast and spectra. Finally apply the same parameter domain across
experimental classes, independently constrain ordinary environmental
models, and reserve a set of observables for prediction. A
mass-density-field model needs a spatial smoothing scale to control the
UV, together with checks of amplification for composites and
particle-permutation symmetry.

\textbf{Critical comparison.} Chapter 40's position noise gives
\(\dot E=\gamma\hbar^2/(2m)\) and \(\Gamma(d)=\gamma d^2/2\), hence
\(\dot E=\hbar^2\Gamma(d)/(md^2)\). A candidate asserting this dephasing
law without the corresponding energy cost must specify the additional
dissipation or feedback that changes the model. It cannot retain only
the favourable half of the prediction.

\textbf{Acceptance.} Obtain reproducible held-out predictions from one
fixed model, including independent bounds on allowed environments. If
arbitrary environmental dilation can always absorb the candidate
completely and no cross-observable constraint is supplied, only an
interpretive difference has been established. A relativistically
covariant extension also requires local consistency and energy control;
this book has not completed that extension.

\section{47.6 R5: Non-Markovian memory and predictive-object
compression}

\textbf{Goal.} Use UEOT's future sufficiency to determine when a density
matrix is enough and when environmental memory must be retained.

\textbf{Steps.} Construct a finite system with a controllable memory
environment and enumerate a tomographically complete intervention set.
Use a causal break to create the same present state with different
future responses. Fit a finite-memory comb and test predictive error
with independent testers. For an allowed protocol family \(\mathcal U\),
define
\(d_{\mathcal U}(R,\widetilde R)=\sup_{T\in\mathcal U}\operatorname{TV}(p_R^T,p_{\widetilde R}^T)\).
Compression should minimize this error under the resource constraints,
rather than merely minimize single-time state-reconstruction error.

\textbf{Mathematical tasks.} Establish composition bounds for
approximate sufficient-factor errors; derive upper and lower
memory-dimension bounds for different access sets; and use numbers of
distinguishable histories or certified matrix ranks to obtain
compression lower bounds. A rank of a process Choi matrix does not
unconditionally equal the minimum physical environment dimension. The
correspondence between representation and implementation needs proof.

\textbf{Acceptance.} Verify Choi positivity and causal recursion of the
compressed model; obtain certified errors for adaptive protocols
excluded from training; and outperform a model retaining only \(\rho_t\)
at the same resource budget. A good fit to trajectories without
interventions does not establish complete future sufficiency.

\section{47.7 R6: Many-body topology, thermodynamics and control
applications}

\textbf{Goal.} Starting from known concrete Hamiltonians, assess whether
object persistence, encoding and control yield comparable performance
improvements.

\textbf{Steps.} Use SSH, TFIM, Hubbard and small codes as four
benchmarks. Solve small volumes exactly before extending them with
MPS/tensor-network or other stated methods. Converge the time step, bond
dimension, truncation and volume separately. Topological tasks must
retain symmetry and boundary specifications; thermal tasks must retain
bath and detailed-balance assumptions; error-correction tasks must
retain noise correlations and decoder latency.

\textbf{Objective function.} For example, given storage time \(T\),
control-energy cost \(C[u]\) and a logical-channel error, define
\(J[u]=1-F_e(\mathcal N_L^u)+\lambda C[u]\). The application specifies
its weights and task; no claim follows that every natural system
minimizes this \(J\). Check dynamical gradients independently by finite
differences or an adjoint method, then compare with a baseline using the
same resources.

\textbf{Acceptance.} The SSH finite-chain residual satisfies Chapter
44's formula; many-body observables have a stated convergence window;
control gains survive held-out noise, temperature and initial states;
and logical errors are actually below the unencoded baseline. If the
gain comes solely from increased resources, attribute it to those
resources. Infinite-volume phase transitions, universal fault-tolerance
thresholds and complete classifications of interacting topological
systems remain separate problems.

\section{47.8 R7: QFT, gauge fixing, EFT and the infrared}

\textbf{Goal.} Make the field-theory modules into computational
interfaces with defined quantities, orders and errors.

\textbf{Steps.} Fix the action, metric, generators, gauge parameter,
regulator and subtraction conditions. Extend the one-loop benchmarks in
Chapters 25/28 to complete amplitudes at the same order. Verify
Ward/Slavnov--Taylor identities and gauge-parameter independence of
physical outputs. For soft/collinear divergences, specify inclusive
observables and resolution, and display cancellation between real and
virtual contributions term by term. A formal LSZ expression cannot
dispose of the problem of charged asymptotic states.

\textbf{EFT tasks.} Fix the operator basis and treatment of
redundancies, retaining cancellation of light-field loops between the
two sides of matching. Run coefficients and convert across thresholds in
the same scheme. Separate loop-order errors, powers of \(E/\Lambda_*\),
parameter errors and input matrix-element errors. The heavy-field Schur
complement in Chapter 26 is an exact algebraic test, not a convergence
theorem for all EFTs.

\textbf{Acceptance.} Two compatible representations give the same
physical cross section; arbitrary renormalization-scale dependence
remains beyond the retained order; and the optical theorem and symmetry
identities hold to that accuracy. A numerical difference caused by
changing schemes without rematching parameters is not new physics. Local
perturbative BRST closure does not automatically solve Gribov or
global-gauge problems.

\section{47.9 R8: Nonperturbative physics across scales}

\textbf{Goal.} Express unification across scales as provable consistency
of the same specified observable.

\textbf{First step: finite models.} Use Chapter 45's exact single-loop
model to test code and normalization. Then implement finite \(SU(N)\)
link updates and test detailed balance, thermalization and
autocorrelation, recording effective sample size and error. A correct
stationary distribution without sufficient mixing does not make a
sampling estimate reliable.

\textbf{Second step: controlled limits.} Select a definite Euclidean
observable, keeping the same physical scale across several volumes,
lattice spacings and bare parameters. Separate statistical,
discretization, finite-volume and parameter-tuning errors. Investigate
the conditions for reflection positivity and transfer matrices;
existence of a finite measure does not by itself establish a complete
continuum QFT.

\textbf{Third step: the Core bridge.} First prove the
readout-pushforward identity for the coarse map. Intervention responses
additionally require a real-time realization and operation slots. When
comparing truncated effective actions, bound the influence of omitted
multi-loop terms on the specified observable. A Core quotient and
Wilsonian integration are not interchangeable names for one
construction.

\textbf{Acceptance.} Obtain checks by independent algorithms and
traceable limits for the selected quantity. If no continuum limit is
established, publish only the finite-lattice result. Computing a finite
spectrum does not solve the four-dimensional mass gap, general chiral
gauge theory or real-time analytic continuation.

\section{47.10 R9: Gauge groups, generations and high-energy
consistency}

\textbf{Goal.} Treat Chapter 31's conditional three-generation theorem
and Chapter 36's high-energy two-candidate classification as comparable
candidates, with their distinct premises retained.

\textbf{Fixed table.} Specify the global group \(G\), all left-handed
Weyl representations, scalars, allowed Yukawa terms, mass thresholds and
gauge embeddings. Beta functions must count every field active in the
relevant energy range. UV asymptotic freedom of \(SU(3)_L\) is not the
same as asymptotic freedom of the broken \(SU(2)_L\) sector under every
hypothetical high-energy extrapolation.

\textbf{Route A.} With three colours fixed by independent observations,
determine whether horizontal anomaly constraints reduce the independent
information compared with directly assuming three generations, and
derive extra-particle and flavour predictions. Three colours remain an
input. A conditional result can still have research value.

\textbf{Route B.} To select the colour number as well, start from an
independent physical selection rule that does not use the desired
\((3,3)\) answer, and ask whether it excludes \((5,5,2)\): for example,
vacuum stability, realizable breaking, or a rigorous global-anomaly
constraint. Compute each rule for the complete candidates. Do not add
scalars or change charges simply because five colours are unwanted.

\textbf{Route C.} Enlarge the representation class and investigate
countermodels. Nonminimal spectra, vectorlike partners or multiple
horizontal factors may change the arithmetic. The existing
classification applies only to its fixed class; an enlarged class
requires a new classification and cannot retain an unsupported
uniqueness title.

\textbf{Acceptance.} Supply local and required global anomaly
certificates for the complete field table, a symmetry-breaking minimum,
mass matrices and a running window, plus at least one low-energy
prediction not used in selection. If every extra condition merely
re-encodes known colour/generation counts, the result is a compatibility
reformulation.

\section{47.11 R10: Flavor, masses, neutrinos and observable matching}

\textbf{Goal.} Turn Chapter 37's rank-one coupling restrictions into
correlations between actual processes using common parameters, and seek
a physically motivated reduction of Yukawa freedom.

\textbf{Steps.} With specified scalars and symmetries, list every
renormalizable mass term and derive mass matrices, unitary
transformations and heavy--light mixing. Then derive the left-handed,
right-handed and mixing couplings of \(Z'\). Chapter 37 assumes a
unitary light sector, common doublets and one rank-one splitting.
Significant exotic mixing requires extra nonunitarity corrections; the
original lower bound cannot simply be reused.

\textbf{From couplings to data.} Match the \(\Delta F=2\) Wilson
coefficients, run them in the same operator basis to the hadronic scale,
and use matrix elements with traceable sources and covariance. Fit
multiple neutral-meson systems and rare decays jointly. Include other
bosons, right-handed currents and cancellation regions in the competing
models. A conservative total FCNC lower bound is not a positive lower
bound on any one channel amplitude.

\textbf{Mass route.} For a proposed geometric/family structure, first
count the free parameters in the allowed Yukawa sector and compare that
dimension with the observables. Arbitrary complex matrices do not
predict mixing. If genuine restrictions remain, fit part of the spectrum
and hold out CKM/PMNS and CP quantities. Count Dirac, Weinberg and
seesaw implementations separately. The light-neutrino mass matrix alone
does not uniquely determine heavy seesaw parameters.

\textbf{Acceptance.} Produce at least one sum rule or held-out
prediction that the remaining parameters cannot adjust arbitrarily, with
complete uncertainty propagation and explicit exclusion conditions. An
attractive fit point is not a proof of the origin of mass.

\section{47.12 R11: Global anomalies, framing and geometric indices}

\textbf{Goal.} Determine whether the original phase/topology ideas can
yield a concrete and natural particle selection, without conflating
distinct topological invariants.

\textbf{Starting point.} Specify a manifold with the required
tangent-bundle structure and gauge bundle, a concrete Dirac-type
operator, boundary conditions and a completion theory. Distinguish the
local anomaly polynomial, global phase, eta invariant, chiral index and
framing response. A Betti number cannot replace an unconstructed Dirac
index.

\textbf{Steps.} Compute spectral flow or boundary response in a
tractable low-dimensional model; construct an explicit inflow/descent
map and check gluing, naturality and background transport; then feed the
resulting response class into Chapter 32's arithmetic conditions. The
value of \(\chi\) must actually be derived, rather than setting
\(\chi=1\) to make three generations unique and seeking a justification
afterward.

\textbf{Acceptance.} Give the same completion's full calculation on
generating backgrounds, prove the necessary natural map, and test
whether an external theory really cannot change the object's autonomy
certificate. If \(\chi\) is not unique, or distinct completions produce
several classes, retain the corresponding candidates. Local anomaly
cancellation neither automatically excludes global anomalies nor proves
that this topological route exists.

\section{47.13 R12: Controlled entry points for gravity and cosmology}

\textbf{Scope.} This book does not construct quantum gravity. A
unification programme involving gravity must first add spacetime
geometry dynamics and a quantum-measurement interface, rather than
rename the existing Hilbert-space theory.

\textbf{First entry point.} Study free fields and local detectors on a
given globally hyperbolic background, specifying states, renormalized
stress--energy and local observables. Compare observer-dependent
particle notions through the same detector response. A nonunique
particle decomposition is not by itself a contradiction.

\textbf{Second entry point.} The semiclassical equation
\(G_{\mu\nu}+\Lambda g_{\mu\nu}=8\pi G_N\langle T_{\mu\nu}\rangle_{\rm ren}\)
is an additional candidate interface, requiring counterterms,
conservation and a suitable fluctuation scale. It is not a complete
solution to quantum gravity and does not automatically handle nonlinear
feedback after conditioning on a measurement. Coupling any stochastic
actualization mechanism to geometry requires checking compatibility of
the Bianchi identity with energy exchange.

\textbf{Third entry point.} For a cosmological candidate, specify the
background initial state, perturbations, observable transfer functions
and parameter priors before claiming held-out power-spectrum
predictions. Dark matter, dark energy, the cosmological constant and
black-hole information each require new concrete models and criteria;
they do not follow from a statement that objects exist.

\textbf{Acceptance.} Establish at least one calculable,
conservation-consistent background model within declared errors before
extending the geometric dynamics. Without that starting point, a broad
unification diagram is not part of a completed theory.

\section{47.14 R13: Formal proofs, code and literature responsibility}

Begin with finite matrices and integer classifications. Choi--Kraus
structure, SRN Gram matrices, the triangle CP bound, the
four-dimensional witness, rank-one FCNCs and global integer inequalities
are possible formalization targets. Encode types,
conjugation/transposition conventions and premises before proving
theorems; random numerical testing is no substitute.

Separate code into model, computation and certificate layers. The model
layer rejects incompatible units or bases; computation returns
approximations with errors; certificates preserve assumptions and
failure states. Non-CP data may be accompanied by an explicitly labelled
fitted model, while retaining the original residual. The proof registry
must distinguish standard prior art from new combinations made by this
project. Before submission, search for prior appearances of the exact
expressions and obtain independent review.

\textbf{Additional results in research edition 3.6.} Chapter 66 proves
the full conditioned classical--quantum (CQ) law of the same original
experimental tree, constructs spectral data for singular support, and
establishes a Choi convex-hull criterion that retains the allowed action
library and observation fibres. It also treats an actual initial filter,
retry costs computed from the original leaf probabilities, and a sharp
conditioning-error bound, while preserving both positive realization
certificates and certificates of impossibility. The source, separate
internal review, and execution receipts for that research revision are
recorded in \texttt{physics/\allowbreak qm/\allowbreak kernel/\allowbreak evidence/\allowbreak v31.\allowbreak 10}. No actual
hardware validation has been performed.

\textbf{Stage already reached in edition 3.5 (historical execution
counts).} Chapters 49--62 retain finite operations, responses,
representations, certificates, errors, mixture averaging, conditional
survival, full complex purification, physical record memory, specified
conditional dynamics, tomography of full complex states, full purity and
common fixed-depth-controller budgets. Chapter 63 identifies the
full-effect distance with half the trace of the matrix absolute value of
the density difference, constructs a spectral effect attaining the
bound, proves optimal equal-prior binary discrimination, and supplies
noncircular matrix certificates. Chapter 64 proves that purification on
a given ancillary carrier exists exactly when its dimension is at least
the density matrix's complex rank, and constructs an actual ancillary
isometry under an explicit dimension ordering. Chapter 65 constructs an
absorbing completion of a finite stopping tree, proves isometric
recovery of the joint state at the original leaves, and gives a
per-action budget weighted by the reference realization's unstopped
mass. At that stage, the main package contained 41 modules, 854 explicit
theorems and 1123 unique named axiom-audit targets. The three new
witness groups were actually executed in the recorded kernel revision.
The ten unchanged older witness groups, six certificates, 100 runtime
methods and four demonstrations reuse genuine earlier gates on the basis
of file identities and dependency-closure identities. The 757 checks
from the three monograph programs retain their edition-3.0 execution
identity. None of these counts denotes new physical laws or experiments.

\textbf{Next verifiable obligations.} Finite quantum properties continue
to assume complete complex matrices, valid effects, finite dimension and
the stated control processes. Quantum selection requires checking the
candidate class and all background premises of the reconstruction
theorem, and independently explaining why a UEOT physical realization
satisfies the selection principles. A trace-distance expression is not a
general effective numerical algorithm. Arbitrary priors, stabilized or
diamond distances, coherent controllers, arbitrary external memory and
unbounded stopping times require new contracts. Dynamical work should
pursue checkable finite contraction or primitivity conditions, existence
of general positive spectral objects and infinite conditioned processes.
Exact program refinement, actual QDRR data, complete spectra and
thresholds, and continuous-variable energy bridges each have separate
acceptance gates.

Real hardware data require their own source and raw-count labels;
synthetic outputs cannot be relabelled as empirical evidence. A
declaration receives a FORMAL label only after its proof script actually
succeeds and its input identity and domain are explicit.

\section{47.15 Common completion criteria and conditions for stopping a
research direction}

Every completed work package should provide a problem definition,
dependencies, model/parameters, proof or counterexample, algorithm and
error, observable contract, independent verification, and uncovered
domains. Integrate new results into the same monograph and update their
dependencies instead of producing several incomplete documents each
called final.

Stop a direction and record the reason if it only adds untestable
assumptions, feeds data back into selection premises, or repeatedly fits
without fixing a model. Retain counterexamples that defeat the original
conjecture and weaken the conclusion accordingly. Progress toward a
unified theory includes rigorously ruling out incorrect routes, not only
obtaining desired numbers.

\section{47.16 The dependency ledger toward a complete physical theory}

The overarching project goal is a complete quantum physical theory based
on UEOT. To make progress assessable, this edition tracks four separate
levels: mathematical definitions and logical closure; the ability to
select physical realizations; composable prediction across models and
scales; and independent experimental constraints. Progress at one level
does not replace progress at another.

\begin{longtable}[]{@{}
  >{\raggedright\arraybackslash}p{(\columnwidth - 6\tabcolsep) * \real{0.2100}}
  >{\raggedright\arraybackslash}p{(\columnwidth - 6\tabcolsep) * \real{0.2000}}
  >{\raggedright\arraybackslash}p{(\columnwidth - 6\tabcolsep) * \real{0.2800}}
  >{\raggedright\arraybackslash}p{(\columnwidth - 6\tabcolsep) * \real{0.3100}}@{}}
\toprule\noalign{}
\begin{minipage}[b]{\linewidth}\raggedright
Current interface
\end{minipage} & \begin{minipage}[b]{\linewidth}\raggedright
Substantive results obtained
\end{minipage} & \begin{minipage}[b]{\linewidth}\raggedright
Remaining obligations
\end{minipage} & \begin{minipage}[b]{\linewidth}\raggedright
Next outcome that can succeed or fail
\end{minipage} \\
\midrule\noalign{}
\endhead
\bottomrule\noalign{}
\endlastfoot
Core → quantum sector & Separation of three models; tomography of full
complex states, full purity, ancillary unitaries and minimum dimension,
full-effect trace distance & Independent origin of selection principles,
all reconstruction backgrounds, and alternative candidate models & Check
reconstruction premises within an explicit candidate class; provide a
testable UEOT physical bridge \\
Quantum sector → kernel & Valid operations, all words and finite
representations, certificate soundness & General runtime refinement,
continuous domains and infinite protocols & A general refinement theorem
and error bound for the same input \\
Kernel → persistence and error & Physical record memory,
unstopped-weight budget under a common finite stopping rule, exact
convergence of specified conditional states & Coherent or unbounded
stopping, external environments, general positive spectral objects and
mixing & Prove the control and distance contracts required by broader
quantifiers; check general sufficient dynamical conditions \\
Known QFT → independent physical selection & Actions, QED/EFT
derivations, anomalies and candidate spectra & Gauge group, observed
spectrum, masses, vacuum and thresholds & One prediction sharing
parameters in a complete fixed model, not used for fitting \\
Candidate → nature & Witnesses, calibration and sequential statistical
procedures & Actual raw data and constraints on competing models & An
auditable experimental conclusion that can reject the candidate \\
\end{longtable}

Acceptance in this ledger consists of a concrete proposition or
counterexample, not an adjective implying proximity to ultimate truth.
New derivations return to their appropriate chapters, with countermodels
and applicability conditions preserved. Evidence closes open entries,
and failed directions retain traceable reasons.

\chapter{48. The integrated monograph, status interfaces, and continuing
maintenance}\label{ch-48}

A maintainable monograph must keep its text, stable IDs, sources, and
verification status aligned. This chapter specifies how those interfaces
are updated and explains why inserting physical selection conditions
into generic Core would make the claimed derivation circular.

\section{48.1 One canonical theoretical source}

The book's fixed title is \emph{UEOT Physics}, its stable document
identity is \texttt{UEOT-\allowbreak PHYSICS}, and its current bilingual revision is
3.11. The canonical text is \texttt{physics/\allowbreak qm/\allowbreak book/\allowbreak UEOT\_\allowbreak PHYSICS.\allowbreak md};
the Chinese PDF is generated from that text. The complete English
counterpart is \texttt{UEOT\_\allowbreak PHYSICS.\allowbreak en.\allowbreak md}, and its PDF is compiled
from that English source. Chapters 0--70 and appendices A--F form one
continuous, complete document. All earlier physical topics are retained,
and new proofs are integrated into the corresponding arguments. Readers
should not have to reconstruct a proof from several old versions.

Version reports, PRs, logs and mathematical modules each retain their
own identity, without competing with the monograph as the authoritative
theoretical text. The source manifest, chapter dossiers and evidence
receipts record the basis of a statement and the part actually checked.
Historical files preserve provenance. Their old statuses do not override
a more recent verified kernel, and a new kernel does not retroactively
confer machine proof on earlier unformalized physical propositions.

\section{48.2 Three independent status axes}

The text's argument-status axis remains D/P/K/H/S/E/N/O: definition,
handwritten proposition/derivation, substantial external theorem,
additional hypothesis, interpretation, empirical input, negative result,
and open problem. It answers what the passage is doing. Physical
dependencies use the repository contract C/B/S/P/E/O/X: Core result,
domain bridge, sector selection, standard-physics input, effective-model
assumption, observable consequence, and open scientific conjecture. This
axis records the dependence and physical role of a conclusion. An
observable consequence labelled O does not by itself indicate existing
experimental data. The two alphabets have different meanings at P, S, E,
O and other shared letters; the field name must accompany a label.

Edition 3.1 explicitly corrected edition 3.0's conflicting use of E for
empirical, O for open and X for counterexample in the dependency field,
preserving the migration reason in the chapter dossiers. Empirical
sources continue to be recorded through argument status and data
identity. A counterexample has argument status N and must not be
migrated to an open conjecture X. A standard mathematical result not yet
formalized by this project is a development obligation; that alone does
not make it X. The upstream contract was not rewritten, and the label
correction neither added nor revoked a mathematical theorem.

The third axis, \texttt{verification\_\allowbreak kind}, records analytical
derivation, exact arithmetic, synthetic numerics, Lean checking or
actual experiment. A handwritten theorem with argument status P is not
thereby Lean-checked; a formalized conclusion supported by
standard-physics input P is not thereby an unconditional Core theorem.
Statement, domain and checking evidence must be retrieved together.

\begin{Verbatim}[breaklines,breakanywhere,fontsize=\scriptsize]
document:
  id: UEOT-PHYSICS
  edition: "3.11"
  canonical_source: physics/qm/book/UEOT_PHYSICS.md
  derivative_pdf: physics/qm/book/UEOT_PHYSICS.pdf
  english_source: physics/qm/book/UEOT_PHYSICS.en.md
  english_pdf: physics/qm/book/UEOT_PHYSICS.en.pdf
  chapter_range: [0, 70]
  fundamental_unique_derivation_complete: false
  four_dimensional_nonperturbative_qft_complete: false
  quantum_gravity_complete: false
  new_experimental_data_in_this_edition: false
  formal_proof_assistant_used_for_listed_kernel_claims: true
claim_schema:
  required:
    - id
    - argument_status
    - dependency_class
    - statement_and_hypotheses
    - domain_units_protocol
    - proof_or_primary_source
    - countermodels_and_exclusions
    - verification_kind_and_receipt
  argument_status: [D, P, K, H, S, E, N, O]
  dependency_class: [C, B, S, P, E, O, X]
  verification_kind: [analytic, exact_arithmetic, synthetic, lean, experiment]
maintenance:
  preserve_existing_result_ids: true
  update_affected_chapters_and_dossier: true
  bind_verification_to_actual_source_hashes: true
  require_separate_evidence_for_stronger_scope: true
\end{Verbatim}

This schema organizes the document and its evidence; a full research
knowledge graph still requires item-by-item development. Counts use
their actual counting units. Multiple parameter points within a test are
not multiple universal physical theorems.

\section{48.3 High-risk conclusions must carry their premises}

\begin{longtable}[]{@{}
  >{\raggedright\arraybackslash}p{(\columnwidth - 2\tabcolsep) * \real{0.3000}}
  >{\raggedright\arraybackslash}p{(\columnwidth - 2\tabcolsep) * \real{0.7000}}@{}}
\toprule\noalign{}
\begin{minipage}[b]{\linewidth}\raggedright
Retrievable conclusion
\end{minipage} & \begin{minipage}[b]{\linewidth}\raggedright
Conditions that must accompany it
\end{minipage} \\
\midrule\noalign{}
\endhead
\bottomrule\noalign{}
\endlastfoot
Conditional uniqueness of three generations and three colours & All
spectral and low-energy weak-window assumptions in Chapter 31; the
five-colour counterexample from the correct high-energy calculation in
Chapter 36 \\
Born-weight absorption & Finite spectrum, specified SSE and QND
condition; degeneracy permits selection only of an eigensubspace \\
Nonzero cycle phase & Excludes a specific diagonal Gaussian diffusion
class; broader classical jumps can still explain it \\
Four-level witness & Fixed pointers, SRN and calibration; excludes
classical diagonal random phases, not quantum environments \\
188 zero-event trials & Ideal event lower bound, fixed sample or valid
supermartingale; actual calibration can increase the required sample
size \\
Finite coordinates cover every future & Actual generating words and all
local pullback identities; fixed allowed library and context \\
Channel-error accumulation & Complete Kraus operators, common carrier
and actual defect bounds; general memory/ancillary-system contracts need
separate proofs \\
Bounded posterior error & Positive conditioning events on both sides and
explicit lower bounds; rare events may make the bound uninformative \\
Existence of a lattice model & The specified finite Euclidean gauge
model; continuum limits and the four-dimensional mass gap remain
separate \\
Existence of a formalized kernel & Source, toolchain and actual audit
for the listed theorems and instances; not machine proof of the entire
physical theory \\
\end{longtable}

\section{48.4 Review, revision and feedback to Core}

Before each change, identify the proposition and its affected domain. Is
the edit correcting an error, adding an omitted premise, or proving a
stronger result? A correction should retain the defect in the earlier
wording and a minimal counterexample. An extension should supply its new
proof and the affected downstream chapters. Check the complete text,
build scripts and evidence locally before the necessary coordinated
submission and remote acceptance.

The revisions recorded here add a branchwise intertwining condition and
counterexample to the former §5.3; distinguish the writing isometry from
the full unitary in Chapter 8; repair parameter decompositions and
restore the vertex UV derivation in Chapters 24 and 27; specify the full
cq channel in Chapter 35; and correct the norm and degeneracy labels in
Chapter 40. Together with the additions of Chapters 49--50, these are
substantive theoretical improvements: eliminating invalid inferences
makes the kernel more reliable.

Results suitable for feedback to the general mathematical Core should
first be abstracted as general theorems whose premises do not presuppose
the desired physical outcome---for example, theorems about response
factors, valid recursion, error composition, persistence events or
nonidentifiability. Complex Hilbert structure, Born probabilities,
particular gauge groups, the action scale and mass spectra are domain
inputs. Inserting them into object-mathematical axioms does not yield an
unconditional derivation of the target. Changes to the frozen Core
continue through its governance process; the developments recorded here
concern QM-domain implementation and integration into the monograph.

\chapter{49. The verified finite quantum kernel of UEOT-QM}\label{ch-49}

The finite quantum kernel begins here: matrices generate record
probabilities, states may agree on every allowed future, and finite
coordinates can preserve all finite-word responses. Formal proofs
establish universal statements and exact certificates check specified
instances; both presuppose the chosen complex quantum realization.

\section{49.1 The scientific role of the kernel}

This chapter connects the object and quantum-operation framework of
Chapters 1, 4, 5, and 33 to mathematical implementations that have
actually been checked. The kernel consists of finite quantum structures
with shared types, dependencies, and composable theorems. Valid states
and operations generate record probabilities; all future records
determine a response object relative to a protocol library; finite
representation certificates make that object computable;
positive-probability posterior updates remain within the space of
realizable representations; and complete finite adaptive trees preserve
normalization. Chapter 50 develops error contracts for these quantum
objects.

The construction implements a central UEOT method: specify the allowed
interventions and future targets before asking which distinctions can
affect future responses. It neither identifies a physical object in
advance with a particular kind of particle nor identifies equal response
coordinates with identical microscopic states. Its physical inputs are
finite complex Hilbert spaces, density operators, Kraus operations, and
the Born pairing. These belong to the quantum sector and its standard
physical realization. The reconstruction framework and additional
principles in Chapter 2 bear the separate responsibility of selecting
that sector.

The formal source is in \texttt{formalization/\allowbreak ueot-\allowbreak qm}; executable
rational models and exporters are in \texttt{physics/\allowbreak qm/\allowbreak kernel}. The
book supplies a continuous theoretical account, while the code and
itemized logs provide inspectable proof artifacts for the corresponding
propositions. These physical realizations do not alter the frozen
mathematics or counts of Core v3.

\section{49.2 Matrices, records, and the response quotient}

Fix a finite index set \(\iota\) and write
\(\mathcal H=\mathbb C^{\iota}\). A state is a matrix satisfying
\(\rho\succeq0\) and \(\operatorname{Tr}\rho=1\). For an allowed action
\(a\), a complete instrument supplies a finite family \(K_{aok}\)
satisfying \[
\sum_{o,k}K_{aok}^{\dagger}K_{aok}=I,
\qquad \Phi_{a,o}(X)=\sum_kK_{aok}XK_{aok}^{\dagger}.
\] The validity fields of \texttt{Density} and \texttt{Instrument} state
these local mathematical conditions. They do not assume in advance that
every future response is correct. Positivity of each branch follows from
\(X\succeq0\Rightarrow KXK^\dagger\succeq0\) and finite addition;
preservation of total trace follows from cyclicity of the trace and
completeness. In particular, \[
p_{a,o}(\rho)=\operatorname{Tr}\Phi_{a,o}(\rho)\ge0,
\qquad \sum_op_{a,o}(\rho)=1.
\]

A word \(w=((a_1,o_1),\ldots,(a_n,o_n))\) records operations in their
actual temporal order. Define \[
R_\rho(\varnothing)=1,\qquad
R_\rho(w)=\operatorname{Tr}\bigl[\Phi_{a_n,o_n}\circ\cdots\circ\Phi_{a_1,o_1}(\rho)\bigr].
\] Once the protocol library is fixed, \(\rho\sim\sigma\) means that
\(R_\rho(w)=R_\sigma(w)\) for \textbf{every finite word}. This
quantifier is not restricted to the finite depths listed in a
certificate. For a target factor \(S\), if \(S(\rho)=S(\sigma)\) implies
equality of responses, the response function factors uniquely through
the reachable image of \(S\). Arbitrary extensions to unreachable
elements are outside the uniqueness claim. This is the quantum instance
of protocol-relative sufficiency developed in Chapter 1.

\phantomsection\label{p-qkr-01}{}

\textbf{P-QKR-01 (Posterior closure of the response object).} For the
same action and outcome, suppose the probability is \(p>0\). The actual
posterior is \(\rho_{a,o}=\Phi_{a,o}(\rho)/p\), and \[
R_{\rho_{a,o}}(w)=\frac{R_\rho((a,o)::w)}{R_\rho([(a,o)])}.
\] Thus response-equivalent states assign the same probability to a
branch, and their posteriors remain response-equivalent on every
positive-probability branch. The update is determined by the response
class itself and remains in the response image of an actual density
operator.

\textbf{Proof.} The branch map and all subsequent operations are linear.
Pulling the scalar \(1/p\) out of the composition and trace gives the
formula. Both numerator and denominator are determined by the current
response function, so changing the representative state within a
response class leaves the posterior response unchanged. The posterior
matrix is positive semidefinite and has trace one, hence the updated
response is realizable. If the branch probability is zero, positivity
and zero trace imply that the branch matrix is zero. No normalized
posterior then exists, and the interface returns \texttt{none}.
\(\square\)

The Lean results include \texttt{response\_\allowbreak posterior},
\texttt{posterior\_\allowbreak response\_\allowbreak eq}, \texttt{next\allowbreak Response\_\allowbreak eq\_\allowbreak none\_\allowbreak iff},
and \texttt{exists\allowbreak Unique\_\allowbreak closed\allowbreak Reachable\allowbreak Update}. They establish more
than the existence of a set-theoretic quotient: the actual update is
independent of the chosen representative and preserves the realizable
domain.

\section{49.3 Finite generating spaces and coordinates for all words}

Define the Heisenberg pullback and word effects by \[
\Phi_{a,o}^*(E)=\sum_kK_{aok}^{\dagger}EK_{aok},\qquad
E_{\varnothing}=I,\quad E_{(a,o)::w}=\Phi_{a,o}^*(E_w).
\] Repeated use of trace cyclicity gives
\(R_\rho(w)=\operatorname{Re}\operatorname{Tr}(E_w\rho)\). Every word
effect is Hermitian. Let \[
W=\operatorname{span}_{\mathbb R}\{E_w:w\text{ is a finite word}\}.
\] This is the smallest real subspace containing \(I\) and closed under
every allowed branch pullback. Containment of the identity and closure
put every word effect in the subspace by induction; conversely, the
linear span of word effects is closed under pullback. Finite
dimensionality of the matrix space guarantees that \(W\) has a finite
generating family consisting of actual word effects.

\phantomsection\label{p-qkr-02}{}

\textbf{P-QKR-02 (A local finite representation determines every future
response).} Choose a finite generating family \(B_j=E_{w_j}\) with
\(B_u=I\). Suppose every branch satisfies the matrix identity \[
\Phi_{a,o}^*(B_j)=\sum_i A_{a,o;ji}B_i,
\] Define \(x_j(\rho)=\operatorname{Re}\operatorname{Tr}(B_j\rho)\).
Then \[
x(\rho)=x(\sigma)\ \Longleftrightarrow\ R_\rho=R_\sigma.
\] The unnormalized branch coordinates are \(y=A_{a,o}x\), with
probability \(p=y_u\). When \(p>0\), the posterior coordinates are
\(y/p\) and belong to the reachable image of coordinates of actual
states.

\textbf{Proof.} The generating words ensure that every \(B_j\) belongs
to \(W\). The identity and local closure ensure that every word effect
belongs to \(\operatorname{span}\{B_j\}\), so the two spaces coincide.
Equal coordinates therefore give equal pairings with every word effect.
The converse follows by reading the responses to the generating words.
Pairing the pullback identity with \(\rho\) gives the branch-coordinate
formula. The posterior formula follows from the actual normalization in
P-QKR-01. \(\square\)

Here row \(j\) of the coefficient matrix expands the pullback of effect
\(j\), so the update is \(Ax\). Changing the row/column convention
requires corresponding changes to the certificate, interpreter, and
proof. Lean's \texttt{Finite\allowbreak Observable\allowbreak Presentation} permits redundant
generating families; it does not require an unproved claim that the
supplied basis is minimal. The exact Gram independence check in Python
provides additional information, but the search and elimination
procedure has no general Lean refinement proof.

In a standard \(d\)-dimensional Hilbert space, diagonal matrix units
together with the real and imaginary Hermitian units for each
off-diagonal pair form a real basis of size \(d+2\binom d2=d^2\).
Ordinary linear algebra therefore gives \(\dim_{\mathbb R}W\le d^2\).
The present Lean interface proves finite dimensionality and the
existence of finitely many generating words; this specific dimension
count and termination of the Python search are not included in its
machine-checked scope.

A qubit illustrates the meaning of protocol relativity. With only a
Z-projective readout, \(W=\operatorname{span}\{I,Z\}\). Adding a
Hadamard operation provides the X direction; adding the phase gate S
also provides the Y direction, giving four real coordinates for the full
Hermitian response space. The same two-level carrier thus has different
predictive representations under different access libraries. Fewer
response coordinates do not by themselves supply a realization in a
smaller Hilbert space, fewer physical controls, or a lower information
cost.

\section{49.4 The complete record law of a finite adaptive protocol}

An allowed finite deterministic adaptive protocol is a finite tree that
selects an action at every nonterminal node and supplies a continuation
subtree for \textbf{every outcome} of that action. The tree may be
unbalanced: different paths need not have the same length. Branch
choices depend only on records already obtained. A leaf probability is
the trace after composing the branch maps along its path.

\phantomsection\label{p-qkr-03}{}

\textbf{P-QKR-03 (All-word coordinates determine a complete finite
adaptive tree).} For the same protocol described above, leaf
probabilities are nonnegative and sum to one. If the initial states have
equal responses to all words, or equal coordinates as in P-QKR-02, every
corresponding leaf has equal probability.

\textbf{Proof.} At a leaf, the leaf mass is the trace of the current
positive matrix. At an internal node, apply the induction hypothesis
separately to each outcome and then sum using the total-trace identity
of the complete instrument. The result is the incoming mass. For a
density operator, the mass at the root is one. Every leaf corresponds to
a unique finite path word, and its probability equals the response to
that word. Equal responses therefore give equal probabilities for all
corresponding leaves. \(\square\)

\texttt{Adaptive.\allowbreak lean} formalizes this induction, injectivity of leaf
words, and the coordinate representation of the leaf law. Deleting
unfavourable leaves, changing the stopping rule after observing the
data, or retaining only surviving outcomes changes this normalization
contract. A new conditional protocol must specify its event probability
and postselection error. Termination and limits for an infinite stopping
time require separate proofs. The statistical supermartingales of
Chapter 46 address a different conditional sampling problem; the shared
word ``adaptive'' does not make the theorems interchangeable.

\section{49.5 Rational weights and variable finite index sets}

The exact program stores \(q_{aok}\in\mathbb Q_{\ge0}\) and
\(M_{aok}\in M_d(\mathbb Q[i])\) and defines branches by \[
\Phi_{a,o}(X)=\sum_kq_{aok}M_{aok}XM_{aok}^{\dagger}
\] Take the complex matrices \(K_{aok}=\sqrt{q_{aok}}M_{aok}\). Termwise
multiplication proves \(KXK^\dagger=qMXM^\dagger\), with corresponding
identities for pullback and completeness. Existence of the square root
is established over the reals; the numerical certificate does not
approximate it. For example, an integer Hadamard matrix with weight
\(1/2\) represents the standard Hadamard channel exactly.

Different outcomes may have different numbers of Kraus terms, and
different actions may have different numbers of outcomes.
\texttt{Ragged\allowbreak Instrument} first allows an outcome-dependent finite type
of terms, then pads with zeros along an explicit injection into a common
type. Reindexing the finite sums along that injection preserves each
branch, its pullback, and Gram completeness, because the added terms
vanish. Padding outcomes likewise adds only zero-probability outcomes.
The theorem requires an actual injection and a specified zero extension.
Vague compatibility of array lengths is insufficient, and out-of-range
word indices are not repaired by taking them modulo an array length.

\section{49.6 Mathematical soundness of the computable checker}

A complex rational matrix is represented as \(A+iB\), with
multiplication \((AC-BD)+i(AD+BC)\) and conjugate transpose
\(A^T-iB^T\). \texttt{Exact.\allowbreak Matrix\allowbreak Data.\allowbreak to\allowbreak Matrix} interprets the data as
a complex matrix. Its injectivity and preservation of these algebraic
operations, finite sums, and rational real scaling are proved.

A numerical certificate \(D\) stores Kraus weights and matrices,
generating words \(w_j\), basis matrices \(B_j\), an identity index
\(u\), and local coefficients \(A_{a,o;ji}\). The independent Lean
checker verifies only five finite conditions: nonnegative weights;
completeness of each action; an empty identity word; equality of each
basis matrix to its actual word effect; and equality of every local
pullback to the supplied finite expansion. It contains no assumption
field asserting correctness of predictions for all futures.

\phantomsection\label{p-qkr-04}{}

\textbf{P-QKR-04 (Soundness of an exact certificate).} If the actual
computable check returns \(\operatorname{check}(D)=\mathrm{true}\), the
interpreted complex matrices form a valid instrument and a finite
observable presentation. The response, posterior-closure, and complete
finite-tree conclusions of P-QKR-01 through P-QKR-03 then hold.

\textbf{Proof.} \texttt{check\_\allowbreak sound} converts Boolean truth into the
five finite mathematical conditions. The algebraic interpretation turns
them into identities of complex matrices. Recursion on words proves that
interpretation preserves actual word effects, binding the numerical
generating words to the actual basis. Apply the local closure theorem
P-QKR-02 and then P-QKR-03. \(\square\)

The six checked instances are Z, H+Z, H+S+Z, a ragged example with
variable outcome and Kraus counts, a complex example with \(d=1\), and a
coarse-projection example with \(d=4\). They cover rational weights,
nonzero imaginary entries, zero-probability outcomes, and response
spaces of less than full dimension. Lean's kernel accepted the
\texttt{checked} declaration for each data set, after which the all-word
and tree consequences were derived. Four classes of post-export
corruption were also examined: weights, numerical basis matrices, local
pullback coefficients, and word order. Each class required both
rejection of the original true-value proof and successful proof of a
false checker value for \textbf{the same corrupted data}. This guards
against spurious negative controls caused by import or resource
failures.

These are a soundness theorem for a mathematical checker and acceptance
records for concrete numerical instances. All behaviours of the Python
implementation, JSON parser, search procedure, and output format have
not been jointly formalized. Acceptance of six instances cannot
establish that stronger claim. Chapter 33 distinguishes implemented
interfaces from future interfaces in the specification.

\section{49.7 Connections to generation spectra and other physical
layers}

\texttt{Spectrum.\allowbreak lean} sums explicitly specified finite family rows and
additional matter rows under effective-window assumptions. A homogeneous
lepton template, a nonzero anomaly unit, and conditions on compensating
terms imply the selector's cubic equation; nonnegative additional
screening contributions enter the weak-sector budget.
\texttt{Selector.\allowbreak lean} then proves conditional uniqueness of
\((N_c,N_f,k)=(3,3,1)\) under the complete integer premises. This goes
beyond finding one triple in a bounded search box, but the premises have
not been derived unconditionally from physical representation theory,
thresholds, and a vacuum mechanism.

The distinction between Chapters 31 and 36 remains essential.
Conditional uniqueness under a low-energy effective weak-sector budget
does not exclude the five-colour candidate under the appropriate
high-energy \(SU(3)_L\) assumptions. The kernel does not silently remove
that independent high-energy candidate from the model class or identify
a topological index with three generations. The physical selection of
masses, mixing, global anomalies, and the observed particle spectrum
remains the responsibility of the relevant chapters.

Continuous variables, general process tensors, relativistic quantum
fields, and the stochastic actualization model of Chapter 40 retain
their own analytic derivations and domain assumptions. The finite matrix
kernel can check controlled finite models within those settings. Passing
from a finite truncation to unbounded operators, infinite volume, or a
continuum limit requires the corresponding error and existence theorems.

\chapter{50. From Born probabilities to quantum error
propagation}\label{ch-50}

An error can be composed along dynamics only after the compared objects
and metric are specified. Born probabilities for all valid effects yield
channel contraction and trajectory budgets, while rare events explain
why normalized posteriors need not contract; Chapters 60 and 63 complete
separation and trace-distance identification.

\section{50.1 All effects and bounds on probabilities}

This chapter establishes a specific mathematical part of the QMU-04
programme: for complete Kraus channels on a common finite quantum
carrier, it constructs an observable probability distance from the
matrices themselves, then proves one-step contraction, model-defect
bounds, trajectory estimates, and postselection bounds for
positive-probability events. The quantities are defined from quantum
objects; contraction is proved rather than supplied as a certificate
field.

An effect \(E\) is a matrix satisfying \(E\succeq0\) and
\(I-E\succeq0\). Write \[
P_E(\rho)=\operatorname{Re}\operatorname{Tr}(E\rho).
\] This is the actual definition of \texttt{Effect}, covering the entire
matrix-order interval \([0,I]\).

\phantomsection\label{p-qst-01}{}

\textbf{P-QST-01 (Positive pairing and effect bounds).} For every
\(X\succeq0\) and effect \(E\), \[
0\le\operatorname{Re}\operatorname{Tr}(EX)\le\operatorname{Tr}X.
\] In particular, \(0\le P_E(\rho)\le1\), and
\(P_{I-E}(\rho)=1-P_E(\rho)\).

\textbf{Proof.} Positivity of a finite complex matrix gives a
factorization \(E=B^\dagger B\). Trace cyclicity yields
\(\operatorname{Tr}(EX)=\operatorname{Tr}(BXB^\dagger)\). The right-hand
side is the trace of a positive semidefinite matrix and is therefore
nonnegative. Apply the same argument to \(I-E\) and expand the linear
pairing to obtain the upper bound. For a density operator, use trace
one. \(\square\)

In Lean, \texttt{pairing\_\allowbreak nonneg} obtains \(B\) from the existing PSD
predicate through Mathlib's matrix-order factorization. The user need
not supply a separate Gram factor. \texttt{Effect.\allowbreak probability\_\allowbreak nonneg},
\texttt{probability\_\allowbreak le\_\allowbreak one}, and \texttt{probability\_\allowbreak complement}
consequently apply to any valid density operator.

\section{50.2 Outcome-dependent pullbacks and all finite records}

For a complete instrument and a continuation effect \(E_o\) for each
outcome, define \[
F=\sum_o\Phi_o^*(E_o)=\sum_{o,k}K_{ok}^{\dagger}E_oK_{ok}.
\]

\phantomsection\label{p-qst-02}{}

\textbf{P-QST-02 (Complete pullback of continuation effects).} \(F\) is
an effect, and \[
\operatorname{Re}\operatorname{Tr}(FX)
=\sum_o\operatorname{Re}\operatorname{Tr}(E_o\Phi_o(X)).
\]

\textbf{Proof.} Every summand is positive semidefinite, so
\(F\succeq0\). Completeness gives \[
I-F=\sum_{o,k}K_{ok}^{\dagger}(I-E_o)K_{ok}\succeq0.
\] The pairing identity follows term by term from trace cyclicity.
\(\square\)

Setting all \(E_o=E\) gives the effect pullback for the complete channel
\(\Phi=\sum_o\Phi_o\) that forgets the outcome. Taking \(E\) at one
selected outcome and zero at the others gives the branch effect
\(\Phi_o^*(E)\in[0,I]\). Starting from \(E_\varnothing=I\) and applying
this branch result recursively proves that every word effect \(E_w\) is
PSD and that \(I-E_w\) is also PSD. Every finite-word response is
therefore an actual effect probability, automatically in \([0,1]\).

The corresponding interfaces are \texttt{Instrument.\allowbreak pullback\allowbreak Effects},
\texttt{branch\allowbreak Effect}, \texttt{record\allowbreak Test}, and
\texttt{record\allowbreak Effect\_\allowbreak complement\_\allowbreak pos\allowbreak Semidef}. They connect all-word
responses to the distance construction. The result allows
outcome-dependent continuation effects. A proof only for the total
channel, with outcome labels discarded, would not establish this
stronger interface.

\section{50.3 The all-effect response distance and its scope}

Define \[
d_E(\rho,\sigma)=\sup_{0\preceq E\preceq I}|P_E(\rho)-P_E(\sigma)|.
\] The zero effect makes the set nonempty, and P-QST-01 bounds it. This
is a supremum over all effects. Neither finitely many tomographic
measurements nor a training or test set is substituted for that
supremum.

\phantomsection\label{p-qst-03}{}

\textbf{P-QST-03 (Basic contract for the observable distance).} The
distance satisfies \(0\le d_E\le1\), symmetry, \(d_E(\rho,\rho)=0\), and
the triangle inequality. Moreover, \[
d_E(\rho,\sigma)=0
\iff \forall E\in[0,I],\ P_E(\rho)=P_E(\sigma).
\] Every finite record in any allowed library satisfies \[
|R_\rho(w)-R_\sigma(w)|\le d_E(\rho,\sigma).
\]

\textbf{Proof.} Every absolute probability difference lies in \([0,1]\)
and obeys the ordinary triangle inequality. Apply it effect by effect
and take suprema to obtain the basic bounds, symmetry, and triangle
inequality. The supremum is zero exactly when all the nonnegative
probability differences vanish. For the record bound, use
\(E_w\in[0,I]\) from P-QST-02. \(\square\)

The standard spectral argument in Chapter 4 identifies
\(d_E=\tfrac12\|\rho-\sigma\|_1\) in finite dimensions. The positive
spectral projector of the difference attains its positive mass, and zero
trace makes the positive and negative masses equal. The interface
originally introduced in this chapter stopped at the probability
pseudometric contract above. Chapter 60 subsequently proved that zero
distance is equivalent to equality of density operators, and Chapter 63
proved the half-trace-absolute-value formula and constructed an optimal
effect. Registration as a Lean \texttt{Metric\allowbreak Space} instance is a
separate implementation question. A restricted access library may still
fail to distinguish different density operators; its protocol-relative
space must not be confused with the set of all effects.

\section{50.4 Channel contraction and model defects}

\phantomsection\label{p-qst-04}{}

\textbf{P-QST-04 (One-step stability of a complete Kraus channel).} The
channel \(\Phi\) defined by a complete instrument satisfies \[
d_E(\Phi(\rho),\Phi(\sigma))\le d_E(\rho,\sigma).
\] For any two state maps on the same carrier, define the worst output
discrepancy by \[
\Delta_E(\Phi,\Psi)=\sup\bigl(\{0\}\cup
\{d_E(\Phi(\tau),\Psi(\tau)):\tau\text{ is a density operator}\}\bigr).
\] If \(\Phi,\Psi\) are complete channels as above, then \[
d_E(\Phi(\rho),\Psi(\sigma))
\le d_E(\rho,\sigma)+\Delta_E(\Phi,\Psi).
\]

\textbf{Proof.} Pullback by \(\Phi^*\) sends every output effect to an
input effect and preserves the Born pairing. Each output probability
difference is therefore bounded by the input \(d_E\); taking the
supremum proves contraction. Insert the intermediate state
\(\Phi(\sigma)\) and apply the triangle inequality, contraction, and the
upper-bound property of \(\Delta_E\). Including zero in the defining set
makes it nonempty even when the density type is empty. It does not
change the value for an ordinary positive-dimensional physical system.
\(\square\)

\texttt{channel\allowbreak State\_\allowbreak effect\allowbreak Distance\_\allowbreak le} and
\texttt{channel\allowbreak State\_\allowbreak perturbation\_\allowbreak bound} prove these relations between
the actual objects. \(\Delta_E\) is a noncomputable supremum, not an
implemented efficient error estimator. Chapter 51 gives a rational upper
bound in the specific model of randomization over a common library. Any
numerical upper bound \(\widehat\Delta\) supplied at runtime still
requires an independent proof that it covers \(\Delta_E\). A maximum
over finitely sampled inputs gives only a lower bound on the worst-case
discrepancy; it cannot be used as a certified upper bound.

\section{50.5 Connecting the one-step quantum theorem to Core transport
certificates}

Given two sequences of complete channels \(\Phi_j,\Psi_j\), define them
in their actual temporal order by \[
\rho_{j+1}=\Phi_j(\rho_j),\qquad
\sigma_{j+1}=\Psi_j(\sigma_j),\qquad
\delta_j=\Delta_E(\Phi_j,\Psi_j).
\]

\phantomsection\label{p-qst-05}{}

\textbf{P-QST-05 (Composable error for quantum trajectories).} For every
finite \(n\), \[
d_E(\rho_n,\sigma_n)
\le d_E(\rho_0,\sigma_0)+\sum_{j=0}^{n-1}\delta_j.
\] For every finite word \(w\) in the same future instrument library, \[
|R_{\rho_n}(w)-R_{\sigma_n}(w)|
\le d_E(\rho_0,\sigma_0)+\sum_{j=0}^{n-1}\delta_j.
\] Since a probability difference cannot exceed one, the bound may be
capped at one in applications.

\textbf{Proof.} Set \(e_j=d_E(\rho_j,\sigma_j)\). P-QST-04 gives the
recurrence \(e_{j+1}\le e_j+\delta_j\), with nonnegative probability
distance and defect. Apply the existing
\texttt{Transport.\allowbreak scalar\_\allowbreak transport\_\allowbreak certificate} as the case \(L_j=1\)
to obtain its recursive budget bound. Induction on \(n\) identifies that
budget with \(e_0+\sum_{j<n}\delta_j\). Apply P-QST-03 for the second
inequality. \(\square\)

The substantive step is that the recurrence follows from the definition
of Kraus channels: the caller no longer has to supply the one-step
quantum inequality that is to be proved. The Core composition theorem is
unchanged; its quantum premises are discharged. The Lean declarations
are \texttt{channel\allowbreak Trajectory\_\allowbreak transport\_\allowbreak certificate},
\texttt{channel\allowbreak Trajectory\_\allowbreak effect\allowbreak Distance\_\allowbreak le\_\allowbreak sum}, and
\texttt{channel\allowbreak Trajectory\_\allowbreak response\_\allowbreak sub\_\allowbreak le\_\allowbreak sum}.

This formal trajectory contract uses a common finite carrier. The two
sequences may have different outcome and Kraus index types, but the
types are fixed within each sequence; the matrices may change from step
to step. Changing index sets can be handled by explicit zero-padding
into a common finite type and the preservation theorem of Section 49.5.
The contract does not include a supremum over inputs entangled with
arbitrary external ancillary systems, so it is not a diamond-distance
theorem. General adaptive instrument errors require explicitly retained
classical histories, outcomes, and quantum memories, followed by
construction of a common complete channel. Section 35.2 gives the
corresponding analytic assembly contract; Chapters 58, 62, and 65
formalize the stated finite cases. Comparing only channels with their
outcomes discarded can miss label information used by a controller.

\section{50.6 Conditioning on positive events and rare branches}

For the same instrument and outcome \(o\), suppose \(p=p_o(\rho)>0\) and
\(q=p_o(\sigma)>0\), with posteriors \(\rho_o,\sigma_o\).

\phantomsection\label{p-qst-06}{}

\textbf{P-QST-06 (Postselection error and preservation of positive
support).} We have \[
|p-q|\le d_E(\rho,\sigma),\qquad
d_E(\rho_o,\sigma_o)\le \frac{2d_E(\rho,\sigma)}{p}.
\] If \(p\ge\eta>0\), the right-hand side can be replaced by
\(2d_E(\rho,\sigma)/\eta\). If \(d_E(\rho,\sigma)<p\) is known, then
\(q>0\), so the approximate model also has a valid posterior.

\textbf{Proof.} For an output effect \(E\), write
\(a=P_{\Phi_o^*(E)}(\rho)\) and \(b=P_{\Phi_o^*(E)}(\sigma)\). P-QST-02
makes the pullback an effect, so \(|a-b|\le d_E(\rho,\sigma)\). Taking
\(E=I\) gives the bound on \(|p-q|\). The posterior pairings are
\(a/p,b/q\), with \(0\le b/q\le1\). The identity \[
\frac ap-\frac bq=
\frac{a-b+(q-p)(b/q)}p
\] and the triangle inequality give \[
\left|\frac ap-\frac bq\right|
\le\frac{|a-b|+|p-q|}{p}.
\] Substitute the bounds for each effect and take the supremum. A lower
bound on the probability controls the denominator; \(q\ge p-d_E>0\)
preserves positive support. \(\square\)

This is not a contraction theorem for selective normalization, and no
optimality claim is made for the constant two. The estimate in Chapter
4, based on the trace norm of subnormalized operators, uses a different
input-error definition. The estimates are not interchangeable merely
because their numerators look similar. If a branch has probability zero
in one model, its posterior remains undefined; assigning an arbitrary
state would not continue the proof legitimately.

Consider a purely illustrative application of the contract. Suppose
independent certificates give initial error \(10^{-3}\) and a defect
bound \(2\times10^{-4}\) at each of twenty steps. The final error is at
most \(0.005\). Postselecting an event with probability at least \(0.2\)
gives posterior-error bound \(0.05\). If the event probability in the
first model is guaranteed only to be \(10^{-3}\), the error bound
\(0.005\) no longer guarantees a positive probability in the second
model, whose posterior may therefore be undefined. The conditioning
theorem applies only after positivity in both models is separately
established. A bound greater than one is merely trivial and indicates
that the certificate is inadequate for this rare-event task. These
numbers illustrate the contract; they are neither measurements nor
performance claims for an existing device.

\section{50.7 Proof, execution, and physical claims}

The first layer of this chapter consists of analytic propositions and
their Lean counterparts. The second consists of actual builds, axiom
checks, and compatibility checks for certificates under a pinned
toolchain. The third asks whether a physical device or domain model
satisfies the quantum realization and numerical error bounds supplied as
inputs. The first two provide a reusable mathematical kernel. The third
requires evidence about preparation, control, noise, thresholds, and
data.

QMU-04 therefore includes a completed slice for complete Kraus channels
on a common finite carrier and for positive-event posteriors. The later
finite-memory and common-controller results are developed in Chapters
58, 62, and 65. General coherent or unequal controllers,
ancillary-system-stabilized norm identification, executable worst-case
error certification, unbounded-energy control, and continuum limits
retain separate obligations in Chapter 47. Further progress must
discharge those explicit obligations; stronger titles for existing
theorems cannot do so.

\chapter{51. Randomized quantum realizations and certified error
budgets}\label{ch-51}

A worst-case supremum over unknown channels is not generally an
executable algorithm; probability mixtures over a common operation
library provide an exactly tractable case. This chapter derives the
factor of one half in the total-variation budget from normalization and
separates weight changes from errors in library members.

Chapter 50 defines quantum errors through suprema over all effects and
all input states. This specifies what is compared but does not provide a
computable upper bound for every device. Here we establish a common,
exactly tractable case: a classical random choice selects a complete
quantum operation from a fixed finite library. Models may differ in
their selection weights or in individually controlled errors of the
library members. Randomization is an explicitly specified operational
mechanism; arbitrary unknown noise is not assumed to have this form.

\section{51.1 Constructing quantum objects from probability weights}

Fix a finite index set \(B\), a common finite-dimensional complex
Hilbert space, and probability weights \[
p_b\ge0,\qquad \sum_{b\in B}p_b=1.
\] Each library member is a complete instrument
\(\{K^{(b)}_{ok}\}_{o,k}\) satisfying
\(\sum_{o,k}K^{(b)\dagger}_{ok}K^{(b)}_{ok}=I\). We use a common finite
outcome set and Kraus index set. If the index sets differ, the verified
zero-padding construction of Section 49.5 may be applied first.

\phantomsection\label{p-qmx-01}{}

\textbf{P-QMX-01 (An actual realization of a probability mixture).} For
any library of density operators \(\rho_b\), \(\rho_p=\sum_b p_b\rho_b\)
is a density operator. With the new index \((b,k)\), define \[
L^{(p)}_{o,(b,k)}=\sqrt{p_b}\,K^{(b)}_{ok},
\] This gives a complete quantum instrument. Its branch maps, Heisenberg
pullbacks, and outcome-forgetting channel are, respectively, \[
\begin{aligned}
\Phi^{(p)}_o(X)&=\sum_b p_b\Phi^{(b)}_o(X),\\
\Phi^{(p)*}_o(E)&=\sum_b p_b\Phi^{(b)*}_o(E),\\
\Phi^{(p)}(X)&=\sum_b p_b\Phi^{(b)}(X).
\end{aligned}
\]

\textbf{Proof.} The positive semidefinite cone is closed under
nonnegative weighted sums. Trace linearity and normalized weights give
\(\operatorname{Tr}\rho_p=1\). The new instrument satisfies \[
\sum_{o,b,k}L^{(p)\dagger}_{o,(b,k)}L^{(p)}_{o,(b,k)}
=\sum_b p_b\sum_{o,k}K^{(b)\dagger}_{ok}K^{(b)}_{ok}
=\Bigl(\sum_b p_b\Bigr)I=I.
\] Expanding \(LXL^\dagger\) or \(L^\dagger EL\) term by term gives the
three linear identities. \(\square\)

Lean's \texttt{Mixture.\allowbreak Weights} stores only the weights, nonnegativity,
and normalization. \texttt{density} and \texttt{weighted\allowbreak Instrument}
implement the formulas above, deriving completeness from the existing
instruments. \texttt{instrument} obtains the square roots through the
verified conversion from weighted operations to ordinary Kraus
operators. A zero weight has zero square root and remains a valid term;
no outcome need be deleted and the protocol identity need not change.
The runtime kernel stores rational weights and unscaled matrices. It can
therefore represent operations such as Hadamard gates using rational
weights without pretending that \(\sqrt2\) is rational.

\section{51.2 Normalization supplies the factor of one half}

For two distributions on the same index set, define \[
\operatorname{TV}(p,q)=\frac12\sum_b|p_b-q_b|.
\] Then \(0\le\operatorname{TV}\le1\). The key estimate applies to all
functions taking values in \([0,1]\), with a representative function
denoted by \(f_b\): \[
\left|\sum_b(p_b-q_b)f_b\right|
=\left|\sum_b(p_b-q_b)(f_b-\tfrac12)\right|
\le\frac12\sum_b|p_b-q_b|.
\] The first step uses \(\sum_b(p_b-q_b)=0\); the second uses
\(|f_b-1/2|\le1/2\). Bounding \(|f_b|\le1\) separately would give the
looser \(\ell^1\) estimate. The formal proof includes the factor of one
half obtained from normalization.

\phantomsection\label{p-qmx-02}{}

\textbf{P-QMX-02 (A uniform bound for changed weights in a common
library).} For a common density-operator library and a common
complete-channel library, respectively, \[
d_E\left(\sum_b p_b\rho_b,\sum_b q_b\rho_b\right)
\le\operatorname{TV}(p,q),
\] \[
\Delta_E\left(\sum_b p_b\Phi_b,\sum_b q_b\Phi_b\right)
\le\operatorname{TV}(p,q).
\]

\textbf{Proof.} Fix an arbitrary effect \(E\) and set
\(f_b=P_E(\rho_b)\). Chapter 50 proves that every \(f_b\in[0,1]\), so
the centring estimate bounds the probability difference for every
effect. Taking the supremum over effects gives the first inequality. For
the second, take any input state \(\tau\), apply the first result to the
density-operator library \(\Phi_b(\tau)\), and then take the supremum
over input states. \(\square\)

This is not a maximum over finitely sampled states. The quantifiers
assert that the same budget controls every valid input and every effect;
this is what permits both suprema. The Lean entry points are
\texttt{expectation\_\allowbreak sub\_\allowbreak le\_\allowbreak tv},
\texttt{density\_\allowbreak weight\_\allowbreak change\_\allowbreak le\_\allowbreak tv}, and
\texttt{channel\allowbreak Defect\_\allowbreak weight\_\allowbreak change\_\allowbreak le\_\allowbreak tv}.

\section{51.3 Decomposing the budget when library members also change}

\phantomsection\label{p-qmx-03}{}

\textbf{P-QMX-03 (A joint bound for weight and member errors).} Given
two complete-channel libraries \(\Phi_b,\Psi_b\) on a common carrier,
write \(\delta_b=\Delta_E(\Phi_b,\Psi_b)\). Then \[
\Delta_E\left(\sum_b p_b\Phi_b,\sum_b q_b\Psi_b\right)
\le\operatorname{TV}(p,q)+\sum_b q_b\delta_b.
\] For fixed weights \(p=q\), the right-hand side becomes
\(\sum_b p_b\delta_b\). The two libraries may have different outcome and
Kraus index types because the comparison concerns their respective
outcome-forgetting channels on the common carrier.

\textbf{Proof.} Fix an input and an effect, and insert the intermediate
channel \(\sum_b q_b\Phi_b\). P-QMX-02 bounds the first difference.
Expand the second as a nonnegative weighted sum of member probability
differences, apply the triangle inequality, and use the true uniform
defect bound for each member. Both suprema preserve this right-hand
side, which is independent of the input and effect. \(\square\)

Exchanging the models gives the symmetric version weighted by \(p_b\).
An implementation may use the smaller of the proved budgets. In either
direction, member bounds require evidence. Equal weights do not
establish equivalence of different libraries: the single-member
distribution \(p=q=(1)\) has zero TV distance from itself, but the
identity channel and the Pauli \(X\) channel send \(|0\rangle\langle0|\)
to outputs distinguishable by the \(|0\rangle\langle0|\) effect with
probability difference one. Omitting the common-library premise would
therefore make the conclusion false.

\section{51.4 Rational budgets and composition over time}

\phantomsection\label{p-qmx-04}{}

\textbf{P-QMX-04 (Interpretation of an exact rational budget).} Suppose
the weights are rational, nonnegative, and normalized. Define, using
only finite rational arithmetic, \[
b_{\rm TV}=\frac12\sum_b|p_b-q_b|\in\mathbb Q.
\] Its real interpretation is exactly the TV distance of the real-valued
distributions. If \(b_{\rm TV}\le b\), the common-library channel defect
satisfies \(\Delta_E\le b\). More generally, if each member also has a
proved rational upper bound \(\delta_b\le\varepsilon_b\), then \[
b_{\rm joint}=b_{\rm TV}+\sum_b q_b\varepsilon_b
\] is a valid upper bound after interpretation over the reals.

\textbf{Proof.} The embedding of the rationals into the reals preserves
finite sums, multiplication, subtraction, absolute values, and order.
Convert each term and apply P-QMX-02 or P-QMX-03. \(\square\)

The corresponding declarations are \texttt{Rational\allowbreak Weights.\allowbreak to\allowbreak Weights},
\texttt{tv\_\allowbreak to\allowbreak Weights},
\texttt{channel\allowbreak Defect\_\allowbreak weight\_\allowbreak change\_\allowbreak le\_\allowbreak rational\_\allowbreak budget}, and
\texttt{channel\allowbreak Defect\_\allowbreak joint\_\allowbreak change\_\allowbreak le\_\allowbreak rational\_\allowbreak budget}. They prove
the mathematical interpretation of finite rational expressions. They do
not automatically verify the Python parser, loops, sequence indexing, or
external file transfer.

If step \(j\) has a budget \(b_j\) of this form, the trajectory
certificate of Section 50.5 gives \[
d_E(\rho_n,\sigma_n)\le d_E(\rho_0,\sigma_0)+\sum_{j<n}b_j.
\] This supplies a concrete connection from actual quantum evolution
through a computable local budget to the Core composition certificate. A
time-dependent library may be specified at each step, provided the
models share the declared correspondence between components at that
step. The existing Lean trajectory interface fixes the finite outcome
and Kraus index types within each sequence; left and right sequences may
differ but share a matrix carrier. Index changes over a finite horizon
require explicit zero-padding or reindexing adaptations.

\section{51.5 Sharp witnesses, implementation obligations, and
extensions}

Take \(\Phi_0=\mathrm{id}\), \(\Phi_1(\rho)=X\rho X\), \(p=(3/4,1/4)\),
and \(q=(1/4,3/4)\). For input \(\rho=|0\rangle\langle0|\), the outputs
are \(\operatorname{diag}(3/4,1/4)\) and
\(\operatorname{diag}(1/4,3/4)\). Measuring \(E=|0\rangle\langle0|\)
gives probability difference \(1/2\), attaining the TV bound. If both
components are instead the identity channel, the same weights still give
budget \(1/2\), while the actual channel defect is zero. The budget is
thus a proved upper bound, generally not an exact-distance algorithm.

The runtime entry \texttt{Mixture\allowbreak Weight\allowbreak Budget} binds two weight tables
to the same collection of verified complete instruments. It computes the
budget exactly from the inputs and does not accept an arbitrarily
supplied smaller number. It records fingerprints of the actual Kraus
representations; different but physically equivalent representations can
have different fingerprints. Tests also cover negative weights, failure
of normalization, floating-point inputs, dimension or outcome-label
mismatches, and context checks for zero-weight components. Runtime tests
provide independent finite witnesses; the universal quantifiers are
discharged by the Lean theorems.

The chapter assumes a specified quantum realization and randomization
model. Its conditional results depend on standard physical input P,
effective-model assumptions E, and domain bridges B; testable
probability constraints belong to dependency class O. It neither selects
complex quantum theory nor supplies a universal algorithm for unknown
channels. It also does not automatically bound instrument distances
retaining outcomes or diamond distances on arbitrary ancillary systems.
The finite history-and-memory developments in Chapters 58, 62, and 65
explicitly retain random branches and control histories as classical
registers before comparing the complete joint channels.

\chapter{52. Long-time averages, compactness, and stationary quantum
states}\label{ch-52}

Stationarity of long-time averages and convergence of the state at every
step are different assertions. Finite-dimensional compactness yields
invariant subsequential Cesàro limits; periodic orbits and multiple
stationary states explain why mixing and uniqueness need additional
conditions.

UEOT's notion of object persistence needs a dynamical meaning that can
be checked. Given a complete quantum channel, a basic question is
whether an invariant average structure exists even when the state keeps
changing. This chapter proves such a result in the actual topology of
finite complex matrices and distinguishes it from pointwise convergence,
uniqueness, conditional survival, and optimization.

\section{52.1 Iterates and average states}

Fix a complete Kraus channel \(\Phi\) and an initial density operator
\(\rho_0\). Define \[
\rho_{n+1}=\Phi(\rho_n),\qquad
\bar\rho_N=\frac1N\sum_{n=0}^{N-1}\rho_n,\qquad N>0.
\] Every iterate is a density operator. Positivity, trace linearity, and
\(N>0\) also make the average a density operator. The average includes
the initial state and ends at \(N-1\), fixing the endpoints in the
telescoping identity below.

\phantomsection\label{p-qce-01}{}

\textbf{P-QCE-01 (Exact residual and all-effect budget for an average
state).} We have \[
\Phi(\bar\rho_N)-\bar\rho_N=\frac{\rho_N-\rho_0}{N},
\qquad
d_E\bigl(\Phi(\bar\rho_N),\bar\rho_N\bigr)\le\frac1N.
\]

\textbf{Proof.} By channel linearity, \[
\Phi\left(\sum_{n=0}^{N-1}\rho_n\right)-\sum_{n=0}^{N-1}\rho_n
=\sum_{n=0}^{N-1}(\rho_{n+1}-\rho_n)=\rho_N-\rho_0.
\] Divide by \(N\) for the first identity. For every effect \(E\), both
endpoint probabilities lie in \([0,1]\), so the absolute probability
residual is at most \(1/N\). Taking the supremum over all effects gives
the second result. \(\square\)

This is not a distance bound to a proposed stationary state. It compares
the average after one additional channel application with the average
itself. Without a stability constant or a spectral gap, a small residual
does not imply proximity to a specified stationary state.
\texttt{matrix\_\allowbreak residual}, \texttt{probability\_\allowbreak residual}, and
\texttt{effect\allowbreak Distance\_\allowbreak residual\_\allowbreak le} formalize precisely this contract.

\section{52.2 Why the density-operator set is compact}

In the usual finite-dimensional topology on complex \(d\times d\)
matrices, let \[
\mathcal D_d=\{\rho:\rho\succeq0,\ \operatorname{Tr}\rho=1\}.
\] Here \(\|\cdot\|_{\rm op}\) is the operator norm induced by the
complex Euclidean vector norm, not the trace norm. Lean uses Mathlib's
\texttt{Matrix.\allowbreak Norms.\allowbreak L2\allowbreak Operator} scope. This structure preserves the
usual matrix topology; one cannot replace it by an artificial topology
that makes every sequence converge.

\phantomsection\label{p-qce-02}{}

\textbf{P-QCE-02 (Compactness of the actual density-operator set).}
\(\mathcal D_d\) is closed and bounded, hence compact in finite
dimensions. Every \(\rho\in\mathcal D_d\) satisfies
\(\|\rho\|_{\rm op}\le1\), and \[
\|\Phi(\bar\rho_N)-\bar\rho_N\|_{\rm op}\le\frac2N.
\]

\textbf{Proof.} The positive semidefinite cone is closed. Equivalently,
for every vector \(v\), nonnegativity of the continuous function
\(\rho\mapsto v^\dagger\rho v\) is preserved under limits. The trace is
continuous and linear, so the trace-one set is closed. Their
intersection is closed. The finite-dimensional Hermitian spectral
theorem gives nonnegative eigenvalues \(\lambda_i\) with
\(\sum_i\lambda_i=1\). Thus every \(\lambda_i\le1\). Unitary conjugation
preserves the operator norm, giving the norm bound. Closed and bounded
subsets of a finite-dimensional normed space are compact. Finally, take
norms in the matrix residual of P-QCE-01 and use
\(\|\rho_N-\rho_0\|_{\rm op}\le2\). \(\square\)

Compactness, the norm bound, and continuity are established for the
matrices themselves; they are not supplied as ``persistence
certificate'' fields that a caller must first accept. The Lean
declarations are \texttt{density\_\allowbreak operator\allowbreak Norm\_\allowbreak le\_\allowbreak one},
\texttt{density\allowbreak Matrices\_\allowbreak closed}, \texttt{density\allowbreak Matrices\_\allowbreak compact},
\texttt{continuous\_\allowbreak channel}, and \texttt{matrix\_\allowbreak residual\_\allowbreak norm\_\allowbreak le}.
Mathlib supplies general results in spectral theory, C*-order, and
finite-dimensional topology; the project proves that density operators
and Kraus channels satisfy their premises.

\section{52.3 Invariant subsequential limits and existence of stationary
states}

\phantomsection\label{p-qce-03}{}

\textbf{P-QCE-03 (A stationary density operator for a finite complete
channel).} For every initial state \(\rho_0\), there exist a density
operator \(\sigma\) and a strictly increasing subsequence of natural
numbers \(N_k\to\infty\) such that \[
\bar\rho_{N_k}\longrightarrow\sigma,\qquad \Phi(\sigma)=\sigma.
\] The limit of any convergent average-state subsequence of this kind is
stationary. In particular, when \(d>0\), every complete finite Kraus
channel has at least one stationary density operator.

\textbf{Proof.} All average states lie in the compact set
\(\mathcal D_d\), so a convergent subsequence exists. Closedness makes
its limit positive semidefinite with trace one. The linear map \(\Phi\)
is continuous in finite dimensions, giving \[
\Phi(\sigma)-\sigma
=\lim_k\bigl(\Phi(\bar\rho_{N_k})-\bar\rho_{N_k}\bigr)=0,
\] The final term tends to zero by the \(2/N_k\) bound in P-QCE-02. The
same argument applies to any convergent subsequence whose indices tend
to infinity. If a channel is given without an initial state, positive
dimension supplies \(I/d\) as a seed, and the result follows.
\(\square\)

\texttt{exists\_\allowbreak stationary\_\allowbreak subsequence} returns the density operator,
its stationarity, and a strictly increasing subsequence.
\texttt{subsequence\_\allowbreak limit\_\allowbreak density} and
\texttt{subsequence\_\allowbreak limit\_\allowbreak invariant} establish physical validity and
invariance of the limit separately. \texttt{exists\_\allowbreak stationary}
constructs the normalized identity on a nonempty finite carrier. This is
an existence theorem, not an implemented stationary-state solver, and it
supplies no general complexity bound for finding a sufficiently long
subsequence.

Core's P-PER-02/M-OI and the associated persistence chapters provide a
general interface linking averages, residuals, and invariant structures.
This chapter independently establishes the required premises and
finite-matrix proof for quantum density operators and channels. The
source does not directly invoke Core's occupation-measure theorem, so a
conceptual correspondence is not reported as a code dependency. The QSD
bridge in Chapter 53 does directly use a frozen Core theorem; the two
kinds of connection are recorded separately. {[}S2, Sections 7--8; S3,
persistence interface{]}

\section{52.4 Counterexamples: periodic orbits and multiple stationary
states}

\phantomsection\label{n-qce-01}{}

\textbf{N-QCE-01 (Average invariance does not imply orbit convergence or
uniqueness).} The premises of P-QCE-01 through P-QCE-03 imply neither
convergence of the original sequence \(\rho_n\) nor uniqueness of the
stationary state.

\textbf{Proof and exact witnesses.} Take \(\Phi(\rho)=X\rho X\) and
\(\rho_0=|0\rangle\langle0|\). Even iterates are \(\rho_0\), and odd
iterates are \(\rho_1=|1\rangle\langle1|\). The \(P_0\) effect
distinguishes them with probability difference one, so the orbit does
not converge. Nevertheless, \[
\bar\rho_N=\operatorname{diag}\left(\frac{\lceil N/2\rceil}{N},
\frac{\lfloor N/2\rfloor}{N}\right)\longrightarrow I/2.
\] For odd \(N\), the one-step residual of \(P_0\) is exactly \(1/N\);
for even \(N\), it is zero. Alternatively, take the identity channel.
Every density operator is stationary, and at least \(\rho_0\) and
\(\rho_1\) are distinct, ruling out uniqueness. \(\square\)

Both counterexamples have the algebraic proofs above and rational
runtime witnesses. Runtime checks are not counted as universal Lean
theorems. The implementation rejects averages with \(N=0\), noninteger
step counts, and invalid initial states. Its averaging endpoints agree
exactly with P-QCE-01.

\section{52.5 The place of long-time averages in the physical framework}

Existence of a stationary state, uniqueness of that state, mixing from
every initial state, and almost-sure empirical convergence of quantum
trajectories are different statements. The finite-dimensional channel
literature distinguishes ergodicity from mixing. The present existence
theorem is consistent with that distinction; a subsequence proof does
not automatically establish the stronger spectral or mean-ergodic
conclusions. \href{https://arxiv.org/abs/1210.5625}{Original study by
Burgarth et al.}

The chapter gives object persistence a concrete certificate without a
circular definition: within a declared physical realization, there is a
valid quantum state invariant in the average-limit sense proved above.
This does not determine which object should persist, choose a value
function, or imply optimal control. GOD/GOA requires separately
specified actions, costs, constraints, and evaluation. A lifetime claim
also requires a specified survival event, which is the subject of the
next chapter.

A further step could prove uniqueness and suitable spectral conditions
for a particular finite channel, then derive a convergence rate. Another
could formalize the projection limit of the entire Cesàro sequence.
Stronger standard tools are available for the latter in
finite-dimensional linear theory, but only the actual subsequence and
existence results are recorded here as machine-checked. In infinite
dimensions, the density-operator set generally lacks the compactness
used above. Energy compactness, tail control, and operator-domain
contracts must then be supplied.

\chapter{53. Conditional survival, diagonal quantum realizations, and
the Core spectral bridge}\label{ch-53}

Survival must first correspond to an actual instrument outcome before
lifetime or conditional state is discussed. A diagonal quantum
realization connects the same survival record to Core quasi-stationary
and Doob formulas, with an exact example distinguishing the
quasi-stationary distribution, Doob stationary law, and initial-state
dependence.

``The object still exists'' becomes a probability statement only after a
survival event is specified. A complete quantum channel preserves total
probability, yet it may include both a continuing-survival outcome and
an outcome that leaves the task domain. The survival probability
generally decays, and the state conditioned on survival need not equal
the unconditional state of the complete channel. This chapter constructs
a finite realization that connects the frozen Core substochastic and QSD
theorems to actual quantum instruments and record probabilities.

\section{53.1 Conventions, outcomes, and completion by failure}

Fix a finite state index set \(i=1,\ldots,d\) and use the
row-distribution convention \[
(pQ)_j=\sum_i p_iQ_{ij}.
\] Let a nonnegative substochastic matrix be given: \[
Q_{ij}\ge0,\qquad s_i:=\sum_jQ_{ij}\le1.
\] The diagonal embedding is
\(D(p)=\operatorname{diag}(p_1,\ldots,p_d)\). A nonnegative normalized
\(p\) gives a complex density operator. The matrix unit
\(U_{ji}=|j\rangle\langle i|\) measures the source \(i\) and prepares
the destination \(j\). Interchanging these indices while retaining the
same \(Q\) would change the convention and the model.

\phantomsection\label{p-qsv-01}{}

\textbf{P-QSV-01 (A complete quantum instrument for a substochastic
kernel).} Take survival and failure Kraus matrices \[
A_{ji}=\sqrt{Q_{ij}}\,|j\rangle\langle i|,
\qquad B_i=\sqrt{1-s_i}\,|i\rangle\langle i|.
\] Together they form a complete instrument, whose survival branch is \[
\mathcal K_Q(X)=\sum_{i,j}Q_{ij}U_{ji}XU_{ji}^{\dagger}.
\]

\textbf{Proof.} Since \(U_{ji}^{\dagger}U_{ji}=|i\rangle\langle i|\),
the Gram sum of the survival terms is \(\operatorname{diag}(s_i)\), and
that of the failure terms is \(\operatorname{diag}(1-s_i)\). Their sum
is \(I\). Nonnegativity and the row-sum bound make all weights valid, so
both branches are completely positive and their total channel preserves
trace. \(\square\)

\texttt{Survival.\allowbreak weighted\allowbreak Instrument} and \texttt{instrument} construct
this instrument. The two Lean outcomes are \texttt{true} and
\texttt{false}; the runtime names them \texttt{survive} and
\texttt{lost}. Lean uses the full \((i,j)\) index set even for failure,
with nonzero weights only on the diagonal. Python omits the zero
off-diagonal failure terms. The operations coincide, while their
representation fingerprints need not.

Using the same matrix carrier after failure is an encoding of the
complete outcome set. ``Dead'' is marked by the classical record, not by
an undeclared absorbing Hilbert-space state. The lifetime task ends at
the first failure, so applying a survival operation again after failure
is not interpreted as resurrection.

\section{53.2 Quantum and classical descriptions of the same survival
record}

\phantomsection\label{p-qsv-02}{}

\textbf{P-QSV-02 (Intertwining diagonal dynamics with finite success
words).} For every complex operator \(X\), \[
\mathcal K_Q(X)=\operatorname{diag}\left(\sum_iQ_{ij}X_{ii}\right)_j.
\] Consequently, for every real vector \(p\) and every \(n\ge0\), \[
\mathcal K_Q(D(p))=D(pQ),\qquad
\mathcal K_Q^n(D(p))=D(pQ^n).
\] If \(p\) is a probability distribution, the probability of \(n\)
consecutive survival outcomes is \(pQ^n\mathbf1\). When this mass is
positive, the conditional state is \[
D\left(\frac{pQ^n}{pQ^n\mathbf1}\right).
\]

\textbf{Proof.} Compute
\(U_{ji}XU_{ji}^{\dagger}=X_{ii}|j\rangle\langle j|\) directly and sum
to obtain the first identity. Substitute \(D(p)\) to obtain one-step
intertwining, then induct in the temporal order of the records using
\((pQ)Q^n=pQ^{n+1}\). The trace of a diagonal state is the sum of its
components, giving the record probability. Normalize by the positive
mass to obtain the conditional state. \(\square\)

Here \(\mathcal K_Q^n\) is the existing \texttt{record\allowbreak Operator} for the
all-success word of length \(n\). \texttt{survival\allowbreak Mass\_\allowbreak eq\_\allowbreak response}
identifies it with \texttt{response} for the same operation library; it
is not an unrelated scalar process. For every \(p,n\), Lean supplies
\texttt{killed\allowbreak Iterate\_\allowbreak diagonal} and \texttt{survival\allowbreak Mass\_\allowbreak diagonal}.
The explicit one-step normalization identity is
\texttt{success\_\allowbreak posterior\_\allowbreak matrix}. The general multistep conditional
formula follows analytically by the same normalization construction; the
QSD case below has its own fully named formal theorems.

The first identity also displays the physical restriction of this model:
a single survival measure-and-prepare step removes any input coherence.
It is a valid quantum realization of a classical process, not a model of
every noncommuting survival channel. Use of complex matrices alone does
not establish a general quantum Perron theory.

\section{53.3 The QSD survival law and lifetime}

Suppose a nonnegative normalized distribution \(q\) and a positive
number \(r\) satisfy the explicitly required left-eigenvector relation
\[
qQ=rq,\qquad q_i\ge0,\quad \sum_iq_i=1,\quad r>0.
\] These spectral data are inputs. This chapter does not first prove
that every \(Q\) has the required positive spectral data, nor does it
assume convergence from every initial state to \(q\). Here \(r\) is a
survival eigenfactor satisfying the relation. Its identification with
the global spectral radius and uniqueness as a dominant eigenvalue are
not outputs of the theorem below.

\phantomsection\label{p-qsv-03}{}

\textbf{P-QSV-03 (The geometric survival law of the actual all-success
record).} Under these premises, \(0<r\le1\), and \[
\mathcal K_Q^n(D(q))=r^nD(q),\qquad
\Pr_q(\text{survival for the first }n\text{ steps})=r^n.
\] Survival has positive probability for every finite \(n\), and the
actual quantum state conditioned on that event remains \(D(q)\).

\textbf{Proof.} Core's frozen \texttt{left\allowbreak Perron\_\allowbreak pow} obtains
\(qQ^n=r^nq\) from the left-eigenvector relation. Transporting it to the
quantum record through P-QSV-02 gives the first identity. Normalization
of \(\sum q_i=1\) and taking the trace give the second; this step
directly invokes Core's \texttt{qsd\_\allowbreak survival\_\allowbreak mass\_\allowbreak pow}. Positivity
\(r>0\) makes every finite-step mass positive. Core's
\texttt{qsd\_\allowbreak all\_\allowbreak steps} says that the normalized distribution remains
\(q\). Diagonal embedding and extensional equality of density operators
then establish equality of the quantum conditional states. Finally, \[
r=\sum_j(qQ)_j=\sum_iq_is_i\le\sum_iq_i=1,
\] This is consumed formally through Core's
\texttt{perron\_\allowbreak root\_\allowbreak le\_\allowbreak one\_\allowbreak of\_\allowbreak substochastic}. \(\square\)

The quantum entry points are \texttt{qsd\_\allowbreak killed\allowbreak Iterate},
\texttt{qsd\_\allowbreak survival\_\allowbreak pow}, \texttt{qsd\_\allowbreak survival\_\allowbreak pos},
\texttt{qsd\_\allowbreak conditional\allowbreak State}, and \texttt{qsd\_\allowbreak rate\_\allowbreak le\_\allowbreak one}. The
dependency chain first constructs the quantum model, then invokes the
frozen Core results, and finally returns to actual records and
conditional states. Quantum--classical equality is proved, not supplied
as a structure field.

A direct analytic consequence is a lifetime formula. Let
\(\tau\in\{1,2,\ldots\}\cup\{\infty\}\) be the step of the first
failure, with infinity assigned if failure never occurs. Then
\(\Pr(\tau>n)=r^n\). When \(r<1\), \(\tau\) is almost surely finite and
\[
\Pr(\tau=m)=(1-r)r^{m-1},\qquad
\mathbb E\tau=\sum_{n=0}^{\infty}\Pr(\tau>n)=\frac1{1-r}.
\] If each step represents a calibrated duration \(\Delta t\), the mean
physical lifetime is \(\Delta t/(1-r)\). The Lean results recorded here
establish the finite-record law \(r^n\) and the conditional state. The
infinite tail-sum expectation is an analytic consequence, and time
calibration is a physical input; neither is reported as an additional
machine proof. The ideal case \(r=1\) has no finite-step failure.

\section{53.4 The Doob realization and its different invariant
distribution}

Suppose also that a strictly positive right eigenvector \(h\) satisfies
\[
Qh=rh,\qquad h_i>0,\qquad \sum_iq_ih_i=1.
\] Define \[
\widetilde Q_{ij}=\frac{Q_{ij}h_j}{rh_i},\qquad \pi_i=q_ih_i.
\]

\phantomsection\label{p-qsv-04}{}

\textbf{P-QSV-04 (A complete realization of the Doob kernel and
invariant success records).} \(\widetilde Q\) is nonnegative with each
row summing to one; \(\pi\) is a probability distribution satisfying
\(\pi\widetilde Q=\pi\). Realizing \(\widetilde Q\) by P-QSV-01 gives a
quantum instrument. For every normalized diagonal input, one-step
success has probability one. Starting from \(D(\pi)\), the success
branch and every finite all-success word preserve \(D(\pi)\), their
record probability is one, and the conditional state is \(D(\pi)\).

\textbf{Proof.} Strict positivity of the denominators gives
nonnegativity. The right-eigenvector relation gives \[
\sum_j\widetilde Q_{ij}=\frac{\sum_jQ_{ij}h_j}{rh_i}=1.
\] Normalization ensures \(\sum_i\pi_i=1\). The left-eigenvector
relation yields \[
(\pi\widetilde Q)_j
=\sum_iq_ih_i\frac{Q_{ij}h_j}{rh_i}
=\frac{h_j}{r}(qQ)_j=q_jh_j=\pi_j.
\] Apply P-QSV-01 and P-QSV-02 to these Core results for row sums,
normalization of the invariant distribution, and stationarity. This
produces the actual instrument, probability, and success-branch
identity. Induction gives every finite success word; normalization by
mass one gives the conditional state. \(\square\)

The implementation directly uses \texttt{doob\allowbreak Transform\_\allowbreak row\_\allowbreak sum},
\texttt{doob\allowbreak Invariant\_\allowbreak sum}, and \texttt{doob\allowbreak Invariant\_\allowbreak stationary}. The
quantum declarations are \texttt{doob\allowbreak Instrument}, \texttt{doob\allowbreak Density},
\texttt{doob\_\allowbreak success\_\allowbreak probability}, \texttt{doob\_\allowbreak success\_\allowbreak branch},
\texttt{doob\_\allowbreak killed\allowbreak Iterate}, \texttt{doob\_\allowbreak survival\allowbreak Mass}, and
\texttt{doob\_\allowbreak conditional\allowbreak State}. These names cover branches, records,
and conditional states. No separate stationarity corollary for the
outcome-forgetting \texttt{channel\allowbreak State} is registered here as an
additional named machine theorem.

Doob transforms are often used to study conditioned trajectories or
alternative dynamics for rare events. General quantum versions involve
noncommuting spectral data and suitable trajectory limits.
\href{https://arxiv.org/abs/1711.10951}{Original study by Carollo et
al.} The present chapter proves only the finite construction under the
stated spectral premises and the invariant relations above. It does not
establish a process theorem equating this transform with the limit of
conditioning on survival to an increasingly distant future. In
particular, the finite-time survival-conditioned law \(q\) cannot be
interchanged with the Doob stationary law \(\pi=qh\).

\section{53.5 An exact model separating the QSD, Doob law, and
initial-state premise}

Take \[
Q=\begin{pmatrix}1/4&1/2\\1/8&1/4\end{pmatrix},\quad
q=(1/3,2/3),\quad r=1/2,\quad h=(3/2,3/4).
\] The two row survival probabilities are \(3/4\) and \(3/8\), so losses
are unequal. Direct calculation gives \[
qQ=(1/6,1/3)=rq,\qquad
Qh=(3/4,3/8)=rh,\qquad \sum_iq_ih_i=1.
\] Starting from \(D(q)\), the all-success record of length \(n\) has
probability \(2^{-n}\), and every conditional state remains \(D(q)\).
The unconditional state also contains the failure-branch contribution.

\phantomsection\label{n-qsv-01}{}

\textbf{N-QSV-01 (The QSD, Doob invariant law, and an arbitrary initial
state are distinct).} In this example, \[
\widetilde Q=\begin{pmatrix}1/2&1/2\\1/2&1/2\end{pmatrix},\qquad
\pi=(1/2,1/2)\ne q.
\] Starting from the uniform state \(p=(1/2,1/2)\) gives first-step
survival probability \(9/16\ne r\). Incorrectly transposing \(Q\) while
retaining the row-distribution interpretation means that, starting from
\(q\), the first-step survival probability instead becomes \(5/8\ne r\).

\textbf{Proof.} Substitution into the Doob formula gives the displayed
rows, and each product \(q_ih_i\) is \(1/2\). Also, \(pQ=(3/16,3/8)\)
has sum \(9/16\), whereas \(qQ^T=(5/12,5/24)\) has sum \(5/8\). Each of
the three substitutions therefore changes the proposition or the actual
probability. \(\square\)

The rational runtime witness computes the first six success words,
obtaining \(1/2,1/4,\ldots,1/64\) and the same conditional state. It
also checks that the Doob dynamics fixes \(\pi\) but does not fix \(q\).
The model has analytic mean failure step two, but six runtime checks are
not a proof of an infinite expectation. Negative controls include row
sums exceeding one, negative entries, nonsquare matrices, and inexact
inputs. Deleting the failure branch breaks completeness. For \(Q=0\),
survival probability is zero and the runtime returns an undefined
posterior; \(q\) cannot be assigned by fiat.

\section{53.6 The contribution to the UEOT programme}

The chapter separates three checkable meanings of persistence:
probability conservation under a complete operation; probability and
conditional state for a specified survival event; and an invariant
structure of alternative dynamics constructed from given left and right
spectral data. The abstract language of object persistence thereby
acquires calculable quantum-record content, together with explicit
diagnoses of which missing premise makes a conclusion fail.

In the dependency classification, frozen QSD algebra belongs to C;
complex density operators and Kraus instruments are standard physical
input P; the choice of diagonal measure-and-prepare dynamics and a
termination event is effective model E; actual intertwining is bridge B;
and lifetime and record-probability constraints are O. General
noncommuting Perron existence, irreducibility or primitivity conditions,
conditional limits, and spectral gaps retain obligations in standard
mathematics, domain bridges, or specific models. Being unformalized does
not by itself make them scientific conjectures X. Only an additional,
unproved claim that a UEOT mechanism is necessary or physically selected
is separately classified as X. The original dual-drive philosophy may
guide which persistence mechanisms to study, but these finite theorems
do not directly determine an objective function or natural optimization.
{[}S1; S2, QSD/Perron and GOA sections; S4, the separation of
persistence and evaluation{]}

Chapter 56 supplies a coherence-preserving finite CP survival map,
explicit left and right spectral objects, a noncommuting Doob
realization, and a finite conditioning connection for arbitrary initial
states. Existence and uniqueness of the required spectral data,
conditional limits, and rates remain separate problems under suitable
assumptions. Only after combining such results with perturbation budgets
of the kind in Chapter 51 can one ask whether lifetime and conditional
predictions transfer stably under errors in preparation, transitions, or
survival classification. This develops reusable physical contracts while
preserving the independent question of why UEOT should select quantum
structure.

\chapter{54. The classical purification boundary on the same Core
process}\label{ch-54}

If generic Core interfaces admit a complete classical model in which a
mixed state cannot be purified, those interfaces do not automatically
imply universal purification. This chapter binds the countermodel to one
actual process and terminal quotient, and explains how deleting readout
changes the proposition being tested.

Chapters 1 and 2 distinguish object mathematics from selection of a
quantum sector. Realizability of a general object interface by a quantum
system does not establish that only quantum systems can realize it. This
chapter makes the distinction concrete using an actual frozen Core
finite stochastic kernel, terminal stable quotient, control quotient,
and infinite trajectory. The classical interpretation is explicit.
Failure of purification is proved using convex extremality in the full
state space, for the actual quotient marginal of the same trajectory.

\section{54.1 The complete classical state space and its operational
interpretation}

For a finite carrier \(Z\), define \[
\Delta(Z)=\{p:Z\to\mathbb R\mid p_z\ge0,\ \sum_zp_z=1\}.
\] The allowed states form the entire simplex; effects are all
\(e:Z\to[0,1]\), with probability pairing \[
P_e(p)=\sum_zp_ze_z.
\] Composite systems use \(\Delta(Z\times B)\), and discarding the right
system is the actual sum marginal \[
(\operatorname{marg}_Z\omega)_z=\sum_b\omega_{zb}.
\] These choices specify a classical domain interpretation B. The H1--H6
interface group listed below does not itself select or exclude this
composition rule; it is not a necessary physical consequence of Core.

\phantomsection\label{p-qcl-01}{}

\textbf{P-QCL-01 (States, responses, and transport on the full
simplex).} The pairing above lies in \([0,1]\), and all effects separate
states. Dirac preparation, valid convex combinations, pushforward along
any finite function, and the stated marginal preserve normalized states.
Pushforward commutes with an equal mixture of two states, and the
marginal is pushforward along the projection.

\textbf{Proof.} Termwise nonnegativity gives \(P_e(p)\ge0\). From
\(e_z\le1\) we obtain \(P_e(p)\le\sum_zp_z=1\). If two states agree on
all effects, the point-indicator effect \(e^{(x)}_z=\mathbf1_{z=x}\)
makes their individual coordinates equal. For a function \(f:Z\to W\),
\[
(f_*p)_w=\sum_{z:f(z)=w}p_z.
\] Interchanging finite sums proves normalization: each source point
occurs in exactly one fibre. The same interchange proves compatibility
with the equal mixture and the marginal identity. \(\square\)

Lean's \texttt{Classical.\allowbreak State} reuses the three purely scalar fields of
\texttt{Mixture.\allowbreak Weights}. It does not assume complex density matrices or
the Born structure in these classical theorems. Quantum modules occur in
the transitive import closure, which is a software dependency fact; the
new state definition and its proof premises concern finite nonnegative
real sums. \texttt{all\_\allowbreak effects\_\allowbreak separate} quantifies over the entire
\texttt{Effect} type. A few sampled measurements do not establish the
same statement.

For any actual finite Core stochastic kernel \(K\), any
\(p\in\Delta(Z)\), and a fixed action \(a\), \texttt{step\allowbreak State} defines
the stochastic action on the complete simplex: \[
(T_ap)_y=\sum_xp_xK(x,a,y).
\] Nonnegativity is termwise; normalization follows from unit row sums
and interchange of finite sums. \texttt{to\allowbreak PMF} uses Core's existing
construction from real probability rows. \texttt{to\allowbreak PMF\_\allowbreak step\allowbreak State}
proves \[
\operatorname{toPMF}(T_ap)
=\operatorname{toPMF}(p)\mathbin{\mathrm{bind}}
\bigl(x\mapsto\operatorname{toPMF}(K(x,a,\cdot))\bigr).
\] Thus the classical state action agrees with the actual PMF transition
for every input state. The trajectory connection below chooses an
explicit policy for the same model. This fixed-action identity alone
does not formalize arbitrary history-dependent feedback.

A state \(p\) is pure exactly when it is an extreme point of the entire
allowed state space: for every \(q,r\in\Delta(Z)\) and \(0<t<1\),
\(p=tq+(1-t)r\) implies \(q=r=p\). \texttt{Is\allowbreak Pure} uses precisely these
quantifiers. Purity of a Dirac state is proved separately. At each zero
coordinate, two nonnegative weighted terms sum to zero, so both vanish.
Normalization forces the remaining coordinates to be one. Purity is not
defined to mean ``Dirac'' and then presented as a classification
theorem.

\section{54.2 Failure of purification for every finite ancillary system}

\phantomsection\label{n-qcl-01}{}

\textbf{N-QCL-01 (A two-point mixed state has no finite classical
purification).} Suppose \(z_0\ne z_1\) and
\(p=\tfrac12\delta_{z_0}+\tfrac12\delta_{z_1}\). For every finite type
\(B\), no \(\omega\in\Delta(Z\times B)\) is both a convex extreme point
and satisfies \(\operatorname{marg}_Z\omega=p\).

\textbf{Proof.} Suppose such a joint state exists. Its marginal has
support at \(z_0,z_1\), with mass \(1/2\) at each point and zero mass in
every other row. Nonnegativity makes every entry in the other rows zero.
Define \[
\omega_i(z,b)=2\mathbf1_{z=z_i}\omega(z,b),\quad i=0,1.
\] Both states are nonnegative. In row \(z_i\), mass \(1/2\) gives
\(\sum_{z,b}\omega_i(z,b)=1\). Coordinatewise, \[
\omega=\tfrac12\omega_0+\tfrac12\omega_1,
\qquad \operatorname{marg}_Z\omega_i=\delta_{z_i}.
\] Extremality forces \(\omega_0=\omega\). Taking marginal coordinate
\(z_1\) gives \(0=1/2\), a contradiction. \(\square\)

The proof does not require \(B\) to be nonempty. If \(B\) is empty, no
normalized joint state exists in the first place. Lean's
\texttt{no\_\allowbreak pure\_\allowbreak extension\_\allowbreak half\allowbreak Mix} quantifies over every finite
ancillary type, not just one selected dimension. It refutes the
existence part of purification, so reversible uniqueness need not first
be formalized.

Removing purity changes the claim entirely. Given an ancillary point
\(b\), \(p\otimes\delta_b\) is an ordinary joint extension of any \(p\).
For a Dirac preparation, \(\delta_{(z,b)}\) is a pure joint extension.
\texttt{marginal\allowbreak Left\_\allowbreak attach} and \texttt{dirac\_\allowbreak has\_\allowbreak pure\_\allowbreak extension}
establish these two facts separately. Existence of a parent carrier, a
record extension, and a pure joint state imposes different mathematical
requirements.

\section{54.3 An actual terminal quotient that retains operational
mixing}

Use the same frozen Core kernel \texttt{survival\allowbreak Micro\allowbreak Kernel}, with
states \(x=(s,h)\in\mathrm{Bool}^2\) and actions \(a\in\mathrm{Bool}\):
\[
K(x,a,z)=\begin{cases}1/2,&z_1=a,\\0,&z_1\ne a.\end{cases}
\] Write \(x_0=(\mathrm{true},\mathrm{false})\) and
\(x_1=(\mathrm{true},\mathrm{true})\). The first bit is the registered
safety predicate. The new instance \texttt{fully\allowbreak Observed\allowbreak Survival\allowbreak Model}
registers the full readout \(x\mapsto x\), together with the safety bit,
trivial program/resource fields, and zero reward. Full readout is this
model's B interpretation; it does not assert that every conceivable
physical label is measurable.

\phantomsection\label{p-qcl-02}{}

\textbf{P-QCL-02 (The actual terminal quotient retains two-point
support).} The model's actual \texttt{terminal\allowbreak Setoid} satisfies
\(x\sim y\iff x=y\). Consequently the control-quotient map \(f\)
constructed by \texttt{organizational\allowbreak Terminal\allowbreak Control} is injective. The
first row of the original kernel, \[
p_* = K((\mathrm{true},\mathrm{false}),\mathrm{true},\cdot)
=\tfrac12\delta_{x_0}+\tfrac12\delta_{x_1}
\] is pushed forward to
\(\bar p=\tfrac12\delta_{f(x_0)}+\tfrac12\delta_{f(x_1)}\), where
\(f(x_0)\ne f(x_1)\).

\textbf{Proof.} Terminal equivalence preserves registered outputs. The
first output is the full source state, so terminal equivalence implies
equality of source states; reflexivity gives the converse. If two images
under the actual control quotient are equal, extract the quotient-class
component and use quotient-equality elimination to obtain terminal
equivalence, then equality of source states. The two-point formula for
the original kernel row is checked from its four entries. Apply
compatibility with equal mixtures from P-QCL-01 and the Dirac
pushforward formula. \(\square\)

An identity quotient has not been substituted by hand for the algorithm.
The carrier is \texttt{Reachable\allowbreak State} of the actual stable relation,
and the map is the actual \texttt{terminal\allowbreak Control.\allowbreak f}. Full observation
means that this example gives no strict reduction in state count, which
is a valid algorithmic output. Compression ratio and retention of
operationally relevant mixing are distinct questions.

\section{54.4 The same control quotient, actual trajectories, and the
first-time marginal}

Fix initial state \(\delta_{x_0}\) and the explicit always-true policy
\(\pi_*(x)=\mathrm{true}\). The transition PMF comes directly from
\texttt{micro.\allowbreak transition\allowbreak PMF} of the same control quotient. Let
\(\mathbb P_*\) be its actual \texttt{stationary\allowbreak Traj\allowbreak Measure}. No
separate trajectory is invented merely to satisfy the desired
conclusion.

\phantomsection\label{p-qcl-03}{}

\textbf{P-QCL-03 (A common-source safe trajectory and a
purification-separation certificate).} This model satisfies the finite
Core interface group listed below, and \[
\mathbb P_*\{\omega:\forall n\in\mathbb N,\ (\omega_n)_1=\mathrm{true}\}=1,
\] \[
(\omega\mapsto f(\omega_1))_*\mathbb P_*
=\operatorname{Measure}(\operatorname{toPMF}\bar p).
\] The state \(\bar p\) on the right has no purification using any
finite classical ancillary system.

\begin{longtable}[]{@{}
  >{\raggedright\arraybackslash}p{(\columnwidth - 2\tabcolsep) * \real{0.3000}}
  >{\raggedright\arraybackslash}p{(\columnwidth - 2\tabcolsep) * \real{0.7000}}@{}}
\toprule\noalign{}
\begin{minipage}[b]{\linewidth}\raggedright
Interface
\end{minipage} & \begin{minipage}[b]{\linewidth}\raggedright
What the example establishes
\end{minipage} \\
\midrule\noalign{}
\endhead
\bottomrule\noalign{}
\endlastfoot
H1: normalized controlled process & Every row of the same \(K\) is
nonnegative and sums to one. \\
H2: registered interface & The transition is exactly \(K\); outputs are
the full readout and the declared attributes. \\
H3: terminal stable quotient & The actual \texttt{terminal\allowbreak Setoid} is
stable, and terminal equivalence is exactly equality. \\
H4: exact control & The actual control quotient, discount \(1/2\), and
zero reward are used; microscopic and macroscopic optimal values are
transported by \(f\). \\
H5: nonempty safety & \(x_0\) belongs to the actual viability set at
every finite horizon, with the required initial PMF support. \\
H6: actual safe trajectory & An explicit always-true policy, equality to
the actual source PMF, safety at all discrete times, and the displayed
actual measure marginal are supplied. \\
\end{longtable}

\textbf{Proof.} Kernel normalization follows from its frozen definition.
Preservation of the registered interface, terminal stability, and
optimal-value transport use the actual theorems of the frozen
interfaces. Prove viability by induction on the horizon. After action
\texttt{true}, every positive-probability destination has first bit
\texttt{true}, giving recursive safety. For trajectories, the initial
distribution is supported on the safe set, and the always-true policy
maps that support back to the same set. The actual
coordinate-distribution theorem gives safety probability one at every
time. Then \texttt{all\_\allowbreak times\_\allowbreak safe\_\allowbreak of\_\allowbreak coordinate\_\allowbreak one} gives
probability one for the countable intersection over all times.

The first-time marginal is not inferred from mere existence of some safe
policy. First use the actual coordinate theorem to identify the
distribution of \(\omega_1\) with one transition from \(\delta_{x_0}\)
under the always-true policy, namely \(K(x_0,\mathrm{true},\cdot)\).
Push this measure forward by the measurable \(f\). Composition of
measure pushforwards and \texttt{to\allowbreak PMF\_\allowbreak pushforward} give the displayed
measure identity. Finally instantiate N-QCL-01 at the two distinct
quotient points supplied by P-QCL-02. \(\square\)

\texttt{Finite\allowbreak Core\allowbreak Slice} expresses these concrete facts as a
conjunction, and \texttt{finite\allowbreak Core\allowbreak Slice\_\allowbreak holds} proves them
individually from the actual theorems.
\texttt{classical\_\allowbreak purification\_\allowbreak separation} then connects the actual
trajectory marginal to the nonexistence quantifier over every finite
ancillary space. Neither failure of purification nor marginal equality
is assumed as an input when constructing the object.

Zero reward is a valid control instance; optimal-value transport here
selects no nontrivial preference. The actions nevertheless have
different safety consequences. For the same \(K\), action \texttt{false}
sets the first bit to \texttt{false}, and the formal witness gives an
unsafe destination probability \(1/2\). Safety at all times means
preservation in probability of this registered predicate. It does not
establish material repair, self-production, lineage persistence, or
endogenous goal selection. {[}S2, S4; QM316{]}

\section{54.5 Why removing the readout must remove this counterexample}

\phantomsection\label{n-qcl-02}{}

\textbf{N-QCL-02 (An operational quotient can remove microscopic
mixing).} Keep exactly the same \(K\), but use the existing
\texttt{registered\allowbreak Survival\allowbreak Model}, which reads only the first bit. Its
actual terminal control quotient merges \(x_0,x_1\). The actual quotient
pushforward of \(p_*\) is therefore a Dirac state, pure according to the
convex-extreme-point definition.

\textbf{Proof.} The frozen theorem
\texttt{registered\_\allowbreak survival\_\allowbreak two\_\allowbreak hidden\_\allowbreak states\_\allowbreak merge} gives terminal
equivalence of the two points. Quotient introduction then gives equal
images under the actual control quotient. Hence \[
f_*p_* = \tfrac12\delta_{f(x_0)}+\tfrac12\delta_{f(x_0)}
=\delta_{f(x_0)}.
\] Apply the independently proved purity of a Dirac state. \(\square\)

This deletion witness shows that the identity of the operational object
changes which mixing is visible. The microscopic second bit remains
random, but the specified readout and actual quotient identify it
completely. Continuing to call this quotient state operationally mixed
would change the object of the purification argument. Conversely, the
main model retains the distinction by full registered readout and bears
responsibility for that domain interpretation. Failure of this
particular mixed-state counterexample under restricted readout does not
prove universal purification for every state of the restricted model.

\section{54.6 The established boundary and the next quantum-selection
problem}

\begin{longtable}[]{@{}
  >{\raggedright\arraybackslash}p{(\columnwidth - 2\tabcolsep) * \real{0.3000}}
  >{\raggedright\arraybackslash}p{(\columnwidth - 2\tabcolsep) * \real{0.7000}}@{}}
\toprule\noalign{}
\begin{minipage}[b]{\linewidth}\raggedright
Result in the text
\end{minipage} & \begin{minipage}[b]{\linewidth}\raggedright
Main Lean entry points
\end{minipage} \\
\midrule\noalign{}
\endhead
\bottomrule\noalign{}
\endlastfoot
All effects, pushforwards, and marginals &
\texttt{all\_\allowbreak effects\_\allowbreak separate}, \texttt{pushforward\_\allowbreak half\allowbreak Mix},
\texttt{marginal\allowbreak Left\_\allowbreak eq\_\allowbreak pushforward} \\
Common-source stochastic action on arbitrary states &
\texttt{step\allowbreak State}, \texttt{to\allowbreak PMF\_\allowbreak step\allowbreak State},
\texttt{to\allowbreak PMF\_\allowbreak pushforward} \\
Extremality and nonexistence for every finite ancilla &
\texttt{dirac\_\allowbreak is\allowbreak Pure}, \texttt{no\_\allowbreak pure\_\allowbreak extension\_\allowbreak half\allowbreak Mix} \\
Actual terminal quotient and separation of states &
\texttt{terminal\_\allowbreak eq\_\allowbreak iff}, \texttt{terminal\allowbreak Control\_\allowbreak injective},
\texttt{first\allowbreak Quotient\allowbreak State\_\allowbreak half\allowbreak Mix} \\
Actual common-source trajectory & \texttt{actual\allowbreak Path\_\allowbreak all\_\allowbreak times\_\allowbreak safe},
\texttt{actual\allowbreak Path\_\allowbreak first\_\allowbreak quotient\_\allowbreak marginal} \\
Interface conjunction and separation & \texttt{finite\allowbreak Core\allowbreak Slice\_\allowbreak holds},
\texttt{classical\_\allowbreak purification\_\allowbreak separation} \\
Removal of the readout & \texttt{restricted\allowbreak First\allowbreak Quotient\allowbreak State\_\allowbreak dirac},
\texttt{restricted\allowbreak First\allowbreak Quotient\allowbreak State\_\allowbreak pure} \\
\end{longtable}

The precise conclusion is that this finite Core interface group is
compatible with the explicit full classical interpretation, while a
reachable operationally mixed state of the same model violates
purification existence. Obtaining universal purification from those
interfaces therefore requires an independent principle that excludes the
example, or an additional physical bridge proved with its actual
premises. A finite certificate does not replace the entire
reconstruction problem of Chapter 1. This is not a model of all 106 Core
statements and every physical interpretation, nor does it rule out
quantum reconstruction from a richer UEOT theory.

Chapter 57 makes the candidate principle's system types, complete state
sets, composition, and discarding into actual interfaces. It proves that
the adapted classical example still violates purification existence and
establishes purification for every finite complex density operator in
the full density space. Exclusion of the real sector in Chapter 55
concerns local tomography instead. Purification existence and essential
uniqueness remain separate properties. Further work must not describe
complex Hilbert, Born, or Kraus structure already supplied as input as
an origin derived from Core.

\chapter{55. A local-tomography separation for the real quantum
sector}\label{ch-55}

Local tomography asks whether all local product tests determine an
arbitrary joint state; a finite sample of tests cannot replace that
quantifier. Two joint states give an exact counterexample in the real
quantum sector, while allowing complex local effects makes the
difference observable.

Purification and local tomography impose different selection
requirements. Chapter 54 refutes purification existence in a complete
classical interpretation. Here the usual restriction to real states and
real effects is fixed, and two joint states are constructed directly:
every real local product test gives the same statistics, while one valid
real joint test distinguishes them perfectly. The proof uses concrete
matrices and quantification over all effects; failure of one measurement
is not presented as failure of every local measurement.

The composition of two rebits and the additional \(Y\otimes Y\)
direction have established precedents; see pages 4--5 and equations
(1.6)--(1.12) of the original paper by Caves, Fuchs, and Rungta. The
contribution here is to connect the specified state pair, universal
effect quantifiers, and deletion witnesses to the formal kernel. The
standard mathematical phenomenon is not classified as a new law of
nature. \href{https://arxiv.org/abs/quant-ph/0009063}{Original paper}

\section{55.1 Reality is an additional restriction; composition must be
specified}

Within the finite complex-matrix realization already used, a matrix is
real when every entry has zero imaginary part. Local density operators
and effects must satisfy their existing positivity, trace, or
upper-bound conditions, together with this reality condition. A real
Hermitian matrix is real symmetric. The joint carrier of two local
systems is a four-dimensional real Hilbert space, with real symmetric
\(4\times4\) density matrices as joint states.

Fix \[
Y=\begin{pmatrix}0&-i\\i&0\end{pmatrix},\qquad H=Y\otimes Y.
\] In basis order \((00,01,10,11)\), \[
H=\begin{pmatrix}
0&0&0&-1\\
0&0&1&0\\
0&1&0&0\\
-1&0&0&0
\end{pmatrix}.
\] A single \(Y\) is Hermitian but not real, whereas \(Y\otimes Y\) is
real symmetric. Excluding the local \(Y\) does not automatically exclude
its tensor product from the full real joint system. This is precisely
why the composition rule must be explicit.

Define \[
P_\pm=\frac{I\pm H}{2},\qquad
\rho_\pm=\frac{P_\pm}{2}=\frac{I\pm H}{4}.
\]

\phantomsection\label{p-qrt-01}{}

\textbf{P-QRT-01 (Two valid, distinct real joint states).} Both
\(\rho_+,\rho_-\) are real, Hermitian, positive semidefinite density
matrices of trace one, and \(\rho_+\ne\rho_-\).

\textbf{Proof.} We have \(Y^\dagger=Y\) and \(Y^2=I\), hence
\(H^\dagger=H\) and \(H^2=I\). Therefore \[
P_\pm^2=P_\pm,\quad P_+P_-=0,\quad P_++P_-=I.
\] A Hermitian idempotent satisfies \(P=P^\dagger P\) and is positive
semidefinite. From \(\operatorname{Tr}H=0\) and joint dimension four,
\(\operatorname{Tr}P_\pm=2\) follows; division by two gives trace one.
Reality follows from the displayed entries of \(H\). The difference of
the two states is \(H/2\), whose \((00,11)\) entry equals \(-1/2\), so
the states are distinct. \(\square\)

Lean's \texttt{involution\allowbreak Effect} constructs an actual effect from a
Hermitian involution. Positivity is proved through a Gram expression,
and the upper bound through the complementary projection.
\texttt{rho\allowbreak Plus} and \texttt{rho\allowbreak Minus} are actual \texttt{Density}
objects, not names for matrices whose validity remains unproved.

\section{55.2 Every real local product effect gives identical
statistics}

Arbitrary local effects \(E,F\) satisfy \(0\le E,F\le I\). Their tensor
product is a valid joint effect: tensor products preserve positive
semidefiniteness, and \[
I-E\otimes F=(I-E)\otimes I+E\otimes(I-F)\ge0.
\] \texttt{product\allowbreak Effect} proves this upper bound; validity of a local
product is not supplied as an unchecked additional input.

Every real symmetric local matrix satisfies \[
E=\begin{pmatrix}a&b\\b&d\end{pmatrix},
\qquad \operatorname{Tr}(EY)=ib-ib=0.
\] This identity holds for all real Hermitian local matrices, not merely
for a few commonly used projections.

\phantomsection\label{n-qrt-01}{}

\textbf{N-QRT-01 (Real local product tests do not determine the joint
state).} For every pair of real local effects \(E,F\), \[
\operatorname{Tr}[(E\otimes F)\rho_+]
=\operatorname{Tr}[(E\otimes F)\rho_-].
\] Nevertheless, the joint states are distinct by P-QRT-01.

\textbf{Proof.} Multiplication and trace identities for finite Kronecker
products give \[
\begin{aligned}
\operatorname{Tr}[(E\otimes F)(\rho_+-\rho_-)]
&=\tfrac12\operatorname{Tr}[(E\otimes F)(Y\otimes Y)]\\
&=\tfrac12\operatorname{Tr}(EY)\operatorname{Tr}(FY)=0.
\end{aligned}
\] The two states therefore have equal Born responses to these tests.
Analytically, the common value is
\(\tfrac14\operatorname{Tr}E\operatorname{Tr}F\). \(\square\)

\texttt{all\_\allowbreak real\_\allowbreak product\_\allowbreak probabilities\_\allowbreak equal} quantifies over actual
\texttt{Effect} objects, with \texttt{Is\allowbreak Real} specifying the allowed
local subset. \texttt{real\_\allowbreak local\_\allowbreak tomography\_\allowbreak fails} combines universal
equality of responses with inequality of states. Local tomography here
means determination of a joint state by all local product responses. The
claim is not extended to a classification of arbitrary unmodelled
communication, adaptive protocols, or external ancillary resources.

A dimension argument explains the obstruction. A single two-dimensional
real symmetric matrix space has dimension three. Products of two such
local spaces span nine dimensions, whereas the four-dimensional real
symmetric joint matrix space has dimension ten. The additional direction
can be chosen as \(Y\otimes Y\). The main Lean result proves the
universal response equality directly rather than substituting an
unformalized dimension count for it.

\section{55.3 A valid real joint test distinguishes the states
perfectly}

\phantomsection\label{p-qrt-02}{}

\textbf{P-QRT-02 (Deterministic discrimination by a joint effect).}
\(P_+\) is a valid real joint effect, and \[
\operatorname{Tr}(P_+\rho_+)=1,\qquad
\operatorname{Tr}(P_+\rho_-)=0.
\]

\textbf{Proof.} Positivity of the projection in P-QRT-01 gives
\(P_+\ge0\), while \(I-P_+=P_-\ge0\), so the full effect conditions
hold. Idempotence gives \(P_+\rho_+=P_+/2=\rho_+\), with trace one;
orthogonality gives \(P_+\rho_-=P_+P_-/2=0\). \(\square\)

The example is therefore not a pair of different mathematical labels for
the same operational state. Joint allowed effects yield different
statistics; only the local product-effect set fails to reconstruct the
difference. \texttt{joint\allowbreak Test} and its two probability theorems
establish this distinction.

\section{55.4 Removing local reality restores a distinguishing test}

\phantomsection\label{n-qrt-02}{}

\textbf{N-QRT-02 (Discrimination when complex local effects are
allowed).} Take the complex local effect \[
E_Y=\frac{I+Y}{2}.
\] It is a valid Hermitian projection although it is not real. Its
product test satisfies \[
\operatorname{Tr}[(E_Y\otimes E_Y)\rho_+]=\tfrac12,
\qquad \operatorname{Tr}[(E_Y\otimes E_Y)\rho_-]=0.
\]

\textbf{Proof.} \(Y\) is a Hermitian involution, so \(E_Y\) and
\(I-E_Y\) are orthogonal positive semidefinite projections. Using
\(\operatorname{Tr}E_Y=1\) and \(\operatorname{Tr}(E_YY)=1\) in the
tensor-trace expression for \(\rho_\pm\) gives \((1\pm1)/4\).
\(\square\)

This deletion control distinguishes the particular state pair above. It
is not the complete theorem of complex quantum local tomography. The
joint states remain unchanged; only the allowed local test set changes.
The witness identifies the actual role of the reality restriction and
prevents N-QRT-01 from being misread as a failure of local tomography in
complex quantum theory itself.

\section{55.5 Linear postprocessing cannot recover the missing joint
effect}

\phantomsection\label{n-qrt-03}{}

\textbf{N-QRT-03 (The distinguishing effect is outside the finite real
span of product effects).} There is no finite index set \(J\), no
collection of real coefficients \(c_j\), and no real local effects
\(E_j,F_j\) such that \[
P_+=\sum_{j\in J}c_j(E_j\otimes F_j).
\]

\textbf{Proof.} If such an expansion existed, real linearity of the
trace pairing and N-QRT-01 would make the right-hand side respond
equally to \(\rho_+,\rho_-\). P-QRT-02 instead gives left-hand-side
responses one and zero, a contradiction. \(\square\)

The coefficients may be negative, so this is stronger than exclusion of
a convex combination of product effects. It concerns the specified
finite linear span, not every physical implementation or infinite
operational protocol. The dedicated \texttt{Real\allowbreak Tomography\allowbreak Two\allowbreak Rebit}
witness connects this linear consequence to the main module's states,
effects, and universal equality.

\section{55.6 Its exact position in the UEOT quantum-selection chain}

\begin{longtable}[]{@{}
  >{\raggedright\arraybackslash}p{(\columnwidth - 4\tabcolsep) * \real{0.2500}}
  >{\raggedright\arraybackslash}p{(\columnwidth - 4\tabcolsep) * \real{0.3400}}
  >{\raggedright\arraybackslash}p{(\columnwidth - 4\tabcolsep) * \real{0.4100}}@{}}
\toprule\noalign{}
\begin{minipage}[b]{\linewidth}\raggedright
Selection question
\end{minipage} & \begin{minipage}[b]{\linewidth}\raggedright
Established result
\end{minipage} & \begin{minipage}[b]{\linewidth}\raggedright
Remaining responsibility
\end{minipage} \\
\midrule\noalign{}
\endhead
\bottomrule\noalign{}
\endlastfoot
Excluding the full classical interpretation & Chapter 54 proves failure
of purification on the specified common-source Core interfaces. & Why
nature must satisfy the proposed purification principle. \\
Excluding this real composition rule & This chapter proves failure of
local tomography for jointly distinguishable states. & Why a physical
realization of UEOT must satisfy local tomography. \\
Uniquely selecting complex quantum structure & A necessary separation
witness is supplied here. & The candidate-theory domain, other
reconstruction premises, and a complete reconstruction theorem. \\
Describing nature & Exact probability separation is proved. & Actual
preparation and measurement capabilities, with calibrated data. \\
\end{longtable}

The machine proof represents the real-entry subspace inside complex
matrices. Chapter 57 supplies an operational presentation of all finite
real states, real effects, real Kraus processes, composition, and
partial trace, and proves seven basic laws. This still does not
formalize every OPT reconstruction premise, categorical coherence, or
essential uniqueness of purification. Nor is the present example a
second complete certificate for the H1--H6 Core interface group of
Chapter 54.

In the dependency classification, standard finite states and effects are
P. The reality restriction and selected composition interpretation carry
B/E responsibilities. Requiring physical systems to satisfy local
tomography is an additional sector-selection principle S, and the
displayed response difference is O. Any further unproved assertion that
a UEOT mechanism necessarily generates these selection principles must
be separately classified as X. Comparing independence and physical
meaning of candidate principles turns the question of quantum selection
into a mathematical problem with explicit exclusion obligations. {[}S2,
S3; QM316{]}

\chapter{56. Coherence-preserving conditional evolution and quantum Doob
channels}\label{ch-56}

Discarding off-diagonal entries while conditioning on survival would
lose quantum coherence. Explicit left and right spectral objects
construct a coherence-preserving Doob channel, separating finite-time
intertwining from stronger long-time convergence and providing an
exactly computable noncommuting example.

Chapter 53 realizes a classical substochastic process diagonally,
removing input coherence after every successful step. This chapter
starts instead from an actual success branch of a quantum instrument.
Under explicit positive spectral and invertible-filter premises, it
constructs a complete quantum Doob channel. The main connection is not
confined to a specially chosen stationary state: for every initial
state, the actual finite-time conditional state is related, by a
specified normalization transform, to an iterate of the Doob channel.

Quantum Doob transforms have precedents in open-system research. Carollo
et al.~relate rare dynamics to an alternative Markov dynamics. The
finite discrete-instrument propositions below are proved here; citation
of that work does not replace the required finite proofs or establish an
infinite-time conditioning limit.
\href{https://arxiv.org/abs/1711.10951}{Original paper}

\section{56.1 The success branch, actual records, and left/right
spectral premises}

Fix a finite complex-matrix carrier and a complete instrument, and
select a success outcome \(o\). Its branch and Heisenberg adjoint are \[
\Phi(X)=\sum_kK_kXK_k^\dagger,
\qquad \Phi^*(A)=\sum_kK_k^\dagger AK_k.
\] The complete instrument also has other outcomes, so this branch
generally preserves or decreases the trace of positive operators.
\texttt{success\allowbreak Operator} uses the existing \texttt{record\allowbreak Operator} for
the actual record in which the same success outcome occurs \(n\) times.
It is not an unrelated iteration proxy. Define \[
m_n(\rho)=\operatorname{Tr}\Phi^n(\rho),\qquad
\rho_n^{\mathrm{cond}}=\frac{\Phi^n(\rho)}{m_n(\rho)}
\] The second expression requires a strictly positive actual mass.
\texttt{success\allowbreak Mass\_\allowbreak eq\_\allowbreak response} connects the first expression to the
record probability in the same instrument library.

The right-eigenstate condition is \(\Phi(\rho_*)=r\rho_*\). The
left-eigenoperator condition concerns the observable matrix \(H\) and
reads \(\Phi^*(H)=rH\). They serve different purposes. The former gives
the geometric survival law of a quasi-stationary state; the latter gives
Doob completeness and the finite conditioning relation for arbitrary
initial states. Existence of one spectral object does not establish
existence of the other.

\section{56.2 Invertible conjugation and normalized state transforms}

Let \(S,T\) be inverse matrices satisfying \(TS=ST=I\). Write \[
F(X)=SXS^\dagger,\quad F^{-1}(X)=TXT^\dagger,
\quad H=S^\dagger S.
\] For a density operator \(\rho\), set \[
c_\rho=\operatorname{Tr}F(\rho),\qquad
\mathcal F(\rho)=F(\rho)/c_\rho.
\]

\phantomsection\label{p-qqd-01}{}

\textbf{P-QQD-01 (Validity of filter normalization).} \(F\) and
\(F^{-1}\) are inverses and preserve positive semidefiniteness. For
every density operator, \(c_\rho>0\), so \(\mathcal F(\rho)\) is a valid
density operator; furthermore, \(c_\rho=\operatorname{Tr}(H\rho)\).

\textbf{Proof.} Conjugation preserves positive semidefiniteness.
Substitution of \(TS=ST=I\) gives the two inverse identities. If the
positive operator \(F(\rho)\) had zero trace, the PSD zero-trace
criterion would make it zero. Applying \(F^{-1}\) would then give
\(\rho=0\), contradicting trace one. Finally, finite-matrix trace
cyclicity gives
\(\operatorname{Tr}(S\rho S^\dagger)=\operatorname{Tr}(S^\dagger S\rho)\).
\(\square\)

\texttt{Invertible\allowbreak Filter} assumes only the two matrices and their actual
inverse identities. \texttt{mass\_\allowbreak apply\_\allowbreak pos} derives positive mass;
positive mass is not inserted as a field of the input structure. There
is no requirement that \(S\) be a previously chosen positive square root
of \(H\), avoiding an additional unproved choice of canonical form.

This mathematical transform must be distinguished from a physical
operation. \(F\) is linear and completely positive, but need not
preserve or decrease trace; normalization makes \(\mathcal F\) generally
nonlinear. The factor \(c_\rho\) can exceed one and therefore is not
directly a filtering success probability. The complete quantum channel
actually proved in this chapter is \(\widetilde\Phi\) in the next
section.

\section{56.3 Deriving Kraus completeness from the left spectral
condition}

Explicitly require \[
r>0,\qquad \Phi^*(H)=rH.
\]

\phantomsection\label{p-qqd-02}{}

\textbf{P-QQD-02 (An actual quantum Doob instrument).} Define \[
\widetilde K_k=r^{-1/2}SK_kT.
\] These matrices form a complete single-outcome instrument, whose
channel satisfies \[
\widetilde\Phi(X)=r^{-1}F\bigl(\Phi(F^{-1}(X))\bigr),
\qquad \widetilde\Phi(F(X))=r^{-1}F(\Phi(X)).
\]

\textbf{Proof.} Compute the Gram sum directly: \[
\begin{aligned}
\sum_k\widetilde K_k^\dagger\widetilde K_k
&=r^{-1}T^\dagger\Bigl(\sum_kK_k^\dagger S^\dagger SK_k\Bigr)T\\
&=r^{-1}T^\dagger\Phi^*(H)T
=T^\dagger HT=(ST)^\dagger(ST)=I.
\end{aligned}
\] Expand \(\widetilde K_kX\widetilde K_k^\dagger\) term by term to
obtain the channel formula, then use \(F^{-1}F=\mathrm{id}\) for
one-step intertwining. \(\square\)

Lean first constructs \texttt{doob\allowbreak Weighted} with explicit real weight
\(r^{-1}\) and matrices \(SK_kT\), then uses the existing verified
conversion to obtain \texttt{doob\allowbreak Instrument} with square-root Kraus
operators. \texttt{doob\_\allowbreak complete} is a derived completeness theorem.
The visible outcome type is \texttt{Unit}; Kraus indices are summed
within the branch. These indices have not been made visible monitoring
outcomes, and equivalence of arbitrary fine-grained record processes is
not proved here.

\section{56.4 Finite-time conditioning for an arbitrary initial state}

\phantomsection\label{p-qqd-03}{}

\textbf{P-QQD-03 (Finite intertwining on all operators and the actual
conditional state).} Under P-QQD-02, for every complex operator \(X\)
and every finite \(n\), \[
\widetilde\Phi^n(F(X))=r^{-n}F(\Phi^n(X)),
\qquad \operatorname{Re}\operatorname{Tr}F(\Phi^n(X))=r^n\operatorname{Re}\operatorname{Tr}F(X).
\] For every initial density operator \(\rho\) and every finite \(n\),
\(m_n(\rho)>0\), and \[
\mathcal F\bigl(\rho_n^{\mathrm{cond}}\bigr)
=\widetilde\Phi^n\bigl(\mathcal F(\rho)\bigr).
\]

\textbf{Proof.} Induction on \(n\) using one-step intertwining gives the
first identity. The left spectral relation and adjoint pairing give \[
\operatorname{Tr}F(\Phi(X))
=\operatorname{Tr}(H\Phi(X))
=\operatorname{Tr}(\Phi^*(H)X)
=r\operatorname{Tr}F(X),
\] Iteration gives the real-part trace relation in the second identity.
The same cyclic-trace calculation also gives the corresponding full
complex-trace identity analytically; the named Lean theorem has the
real-valued \texttt{mass} type. For a density operator, the relevant
operators are PSD, their traces are real, and the right-hand side is
\(r^nc_\rho>0\). If \(m_n(\rho)=0\), the PSD operator \(\Phi^n(\rho)\)
would vanish by the zero-trace criterion, forcing the left-hand side to
vanish as well, a contradiction. All finite conditional states are
therefore well defined.

Finally substitute the actual normalizations: \[
\begin{aligned}
\mathcal F\!\left(\frac{\Phi^n(\rho)}{m_n(\rho)}\right)
&=\frac{F(\Phi^n(\rho))}{\operatorname{Tr}F(\Phi^n(\rho))}\\
&=\frac{r^{-n}}{c_\rho}F(\Phi^n(\rho))
=\widetilde\Phi^n\!\left(\frac{F(\rho)}{c_\rho}\right).
\end{aligned}
\] Cancellation of \(m_n\) in the first step requires strict positivity
and cannot be used for a zero-probability event. \(\square\)

The Lean declarations are \texttt{doob\_\allowbreak success\allowbreak Operator\_\allowbreak apply},
\texttt{eigenmetric\_\allowbreak success\allowbreak Mass},
\texttt{eigenmetric\_\allowbreak success\allowbreak Mass\_\allowbreak pos}, and
\texttt{doob\_\allowbreak finite\_\allowbreak conditioning}. The first two hold on all
operators; the conditional-state conclusion explicitly uses actual
\texttt{Density} objects. Normalization refers to the originally defined
success record, not to similar notation in a different evolution.

The left-hand side still contains \(\mathcal F\). The theorem does not
equate the original conditional state directly with a Doob orbit
starting from the untransformed input. Nor does it assert survival rate
\(r^n\) for every initial state. Inverting the first identity
analytically also gives \[
m_n(\rho)=r^nc_\rho\operatorname{Tr}\!\left[
T\,\widetilde\Phi^n(\mathcal F(\rho))T^\dagger\right].
\] Thus actual survival mass for a general input has an additional
state-dependent factor. This analytic consequence is not counted as a
separate named Lean theorem.

\section{56.5 A right eigenstate supplies the quasi-stationary and
Doob-stationary states}

\phantomsection\label{p-qqd-04}{}

\textbf{P-QQD-04 (Quasi-stationarity, geometric survival, and
stationarity of the actual channel).} Suppose \(r>0\). If a density
operator \(\rho_*\) is additionally supplied with
\(\Phi(\rho_*)=r\rho_*\), then \[
\Phi^n(\rho_*)=r^n\rho_*,\quad m_n(\rho_*)=r^n,
\quad \rho_{*,n}^{\mathrm{cond}}=\rho_*.
\] The probability bound of the actual instrument gives \(r\le1\). If
the preceding left spectral condition also holds, then \[
\widetilde\Phi(\mathcal F(\rho_*))=\mathcal F(\rho_*).
\] Hence \(\mathcal F(\rho_*)\) is a stationary density operator of the
actual \texttt{channel\allowbreak State}. Every finite repeated single-outcome
record preserves this state and has mass one.

\textbf{Proof.} Induction on the number of steps gives the first
identity from the right-eigenstate relation. Taking traces and using
\(\operatorname{Tr}\rho_*=1\) gives geometric mass. Strict positivity of
\(r^n\) makes normalization return \(\rho_*\). The one-step success
probability is \(r\), which is at most one by completeness of the
instrument. Substitute the right-eigenstate relation into one-step
intertwining and normalize by \(c_{\rho_*}\) to obtain the Doob fixed
point. The single branch is the whole channel, so the conclusion applies
to the actual \texttt{channel\allowbreak State}. \(\square\)

Unit mass of the Doob single-outcome record follows from trace
preservation. The substantive construction is its all-operator
intertwining with the original success branch, the conditioning relation
for every input, and the explicit stationary matrix. No theorem here
asserts that every CP branch has the required positive invertible left
spectral matrix, identifies the spectral radius, proves simplicity of
the leading eigenvalue, or proves long-time attraction.

\section{56.6 An exact rational example retaining input coherence}

Take a two-dimensional source instrument with success Kraus operators \[
K_0=\tfrac14I,\quad
K_1=\begin{pmatrix}1/4&0\\1/2&0\end{pmatrix},\quad
K_2=\begin{pmatrix}0&1/8\\0&1/4\end{pmatrix},\quad K_3=0.
\] and failure Kraus operators \[
L_0=\operatorname{diag}(3/4,7/8),\quad
L_1=\operatorname{diag}(1/4,1/4),\quad
L_2=L_3=\operatorname{diag}(0,1/8).
\] Direct calculation gives \[
\sum_kK_k^\dagger K_k=\operatorname{diag}(3/8,9/64),
\quad \sum_kL_k^\dagger L_k=\operatorname{diag}(5/8,55/64),
\] Their sum is the identity. For every complex matrix \(X\), the
success branch is \[
\Phi(X)=\begin{pmatrix}
X_{00}/8+X_{11}/64&X_{01}/16+X_{00}/8+X_{11}/32\\
X_{10}/16+X_{00}/8+X_{11}/32&X_{00}/4+X_{11}/8
\end{pmatrix}.
\] Thus the off-diagonal input \(|0\rangle\langle1|\) is multiplied by
\(1/16\) rather than erased; also, \(K_1K_2\ne K_2K_1\). These exact
witnesses show that the example extends beyond the one-step diagonal
measure-and-prepare model of Chapter 53. They do not constitute a
classification of every nonclassical resource.

Choose \[
\begin{aligned}
\rho_* &=\tfrac15\begin{pmatrix}1&2\\2&4\end{pmatrix},\qquad r=3/16,\\
S&=\operatorname{diag}(2,1),\quad T=\operatorname{diag}(1/2,1),\quad H=\operatorname{diag}(4,1).
\end{aligned}
\] The exact identities are \[
\Phi(\rho_*)=r\rho_*,\qquad \Phi^*(H)=rH,\qquad
c_{\rho_*}=8/5,
\] \[
\sigma_*:=\mathcal F(\rho_*)=\tfrac12\begin{pmatrix}1&1\\1&1\end{pmatrix}
=|+\rangle\langle+|.
\] The normalization factor \(8/5\) illustrates the interpretive
boundary of Section 56.2: it exceeds one and is therefore not an outcome
probability of the source instrument.

\phantomsection\label{p-qqd-05}{}

\textbf{P-QQD-05 (A nondiagonal Doob witness computable on all
matrices).} For every complex operator \(X\), the actual Doob channel in
this example satisfies \[
\widetilde\Phi(X)=\tfrac13X+\tfrac23\operatorname{Tr}(X)\,\sigma_*.
\] Starting from the source quasi-stationary state, the probability of
\(n\) consecutive successes is \((3/16)^n\), and the conditional state
remains \(\rho_*\). The transformed channel preserves \(\sigma_*\). The
arbitrary-input finite-conditioning theorem P-QQD-03 also applies to
this example.

\textbf{Proof.} Substitute \(S,T,r\) and the source-branch formula on
all matrices. Simplifying the four entries gives the displayed
linear-operator identity. The remaining conclusions instantiate P-QQD-03
and P-QQD-04. \(\square\)

The dedicated \texttt{Quantum\allowbreak Doob\allowbreak Examples} proves
\texttt{source\_\allowbreak branch\_\allowbreak formula} and
\texttt{transformed\_\allowbreak branch\_\allowbreak formula} for arbitrary complex \(X\). The
latter uses the complex trace explicitly. The real-valued \texttt{mass},
which agrees with trace on Hermitian inputs, must not be substituted for
trace on all complex operators. All matrices and coefficients are exact
rational inputs. Square roots enter only through the verified conversion
to an ordinary instrument; floating-point approximations do not enter
the proof.

\phantomsection\label{n-qqd-01}{}

\textbf{N-QQD-01 (Three distinctions that must be retained).} In this
example, removing the failure Kraus operators breaks completeness of the
source instrument; the original quasi-stationary state \(\rho_*\)
differs from the filtered stationary state \(\sigma_*\); and from the
uniform initial state \(I/2\), the one-step success probability is
\(33/128\), not \(r=3/16\).

\textbf{Proof.} The success Gram sum has diagonal entries \(3/8,9/64\),
neither equal to one. The two states have off-diagonal entries \(2/5\)
and \(1/2\), respectively. Finally, \[
\operatorname{Tr}\Phi(I/2)=\tfrac12(3/8+9/64)=33/128\ne3/16.
\] All three conclusions have dedicated Lean witnesses. \(\square\)

\section{56.7 The next dynamical obligations}

This construction extends the diagonal success embedding of Chapter 53
to a conditional construction for general finite Kraus branches and
completes the all-matrix calculation in a nondiagonal example. It
complements Chapters 54 and 55: those chapters examine exclusion
obligations for quantum-selection principles, while this one strengthens
usable dynamics after quantum structure is supplied. The tasks do not
replace one another.

The results are conditional on standard quantum structure P and actual
operational bridges B. The selected success task and source instrument
are E, and record probabilities are O. Core is not directly invoked to
select complex quantum structure, so thematic connections to QSD or
persistence do not place every result in C. Spectral existence,
irreducibility, uniqueness, and conditional limits remain obligations in
standard mathematics or domain bridging; additional claims of physical
necessity specific to UEOT would be classified separately as X.

For the same rational nondiagonal model, Chapter 59 proves exact
iterates, uniqueness of the Doob invariant state, and a geometric
convergence rate for the actual conditional state. A valid identity Doob
channel also shows that a positive spectral metric does not by itself
guarantee mixing. General progress requires verifiable finite
irreducibility, primitivity, or contraction conditions before
establishing spectral objects and long-time conditioned processes. The
original dual-drive idea can suggest persistence mechanisms to study,
but success events, inputs, control costs, and actual observations still
require separate definitions. {[}S1, S2, S4; QM316; QM317{]}

\chapter{57. Purification and operational principles in the complete
state space}\label{ch-57}

A purification principle concerns the full state space and an actual
discard operation, not only a selected class of vector states. Positive
semidefinite matrices yield purifications of arbitrary finite complex
states; common operational interfaces then compare classical, real, and
complex presentations while separating existence from a full
reconstruction principle.

Chapter 54 showed that a collection of concrete finite Core process
interfaces admits a complete classical interpretation, while mixed
states in that interpretation have no finite classical pure extension.
This result creates a constructive obligation. If purification is
proposed as a principle for selecting a quantum sector, systems, all
allowed states, composition, and discarding must first be specified in a
common operational language. One must then prove that the complex
quantum model satisfies the principle as stated. This chapter constructs
purifications of finite-dimensional complex density states and
identifies the remaining distance between that construction and a
quantum reconstruction.

\section{57.1 A Complete Operational Presentation Before Selection
Principles}

An operational presentation of a system must at least specify its
states, effects, effect probabilities, convex mixtures, processes with
explicit input and output types, and the action of those processes on
states. Composite systems also require product states, product effects,
and concrete discard maps. A separate package contains seven basic laws:
probability bounds; affinity of probability under mixtures;
multiplication of product probabilities; marginal identities for product
states; affinity of processes; the state action of identity processes;
and the state action of sequential composition. These laws must be
proved in each concrete model. A structure field having the required
type does not establish all the model's physical properties.

System types must genuinely represent composition. For finite carriers
\(\iota\) and \(\alpha\), the joint matrix carrier in this chapter is
\(\iota\times\alpha\). This avoids equating equal dimensions with equal
types and makes the summation index in a partial trace correspond
directly to the discarded system. Processes with different input and
output spaces use the rectangular Kraus channels of Chapter 58: their
Kraus matrices are actual linear maps from the input vector space to the
output vector space.

For a finite complex quantum system, the state space remains

\[
\mathcal D(\iota)=\{\rho\in M_\iota(\mathbb C):\rho\succeq0,\ \operatorname{Tr}\rho=1\}.
\]

Effects comprise all matrices satisfying \(0\preceq E\preceq I\), with
probability \(\operatorname{Re}\operatorname{Tr}(E\rho)\). The matrix
product \(E\rho\) is generally not Hermitian. Its trace is real by the
PSD Gram factorization, cyclicity of trace, and the reality of the trace
of a positive matrix; one cannot replace that argument by asserting that
the product remains positive. Taking the real part makes the probability
function explicitly real-valued in its type. For \(0\le t\le1\), a
mixture is the actual matrix \(t\rho+(1-t)\sigma\); positivity and unit
trace follow from the corresponding linear identities.

Purity is defined as extremality in this complete convex set:

\[
\begin{aligned}
\operatorname{Pure}(\rho)\iff{}&
\forall\sigma,\tau\in\mathcal D(\iota),\ \forall 0<t<1,\\
&\rho=t\sigma+(1-t)\tau
\ \Longrightarrow\ \sigma=\rho\ \land\ \tau=\rho.
\end{aligned}
\]

This property quantifies over every valid density state. Merely calling
rank-one matrices ``pure states,'' or checking extremality in a selected
subset of states, does not discharge the proof obligation here.

\section{57.2 Why Unit-Vector States Are Extreme Points of the Complete
Convex Set}

\phantomsection\label{p-pur-01}{}

\textbf{P-PUR-01 (Vector purity in the complete density-state space).}
Let \(P=|v\rangle\langle v|\) with \(\operatorname{Tr}P=1\). Then \(P\)
is an extreme point of the complete set \(\mathcal D(\iota)\).

\textbf{Proof.} Normalization gives \(P^2=P\). Set \(Q=I-P\); then
\(Q=Q^\dagger\) and \(QP=0\). Suppose \(P=t\sigma+(1-t)\tau\), where
\(0<t<1\) and \(\sigma,\tau\) are arbitrary density states. Compression
on both sides gives

\[
0=QPQ=tQ\sigma Q+(1-t)Q\tau Q.
\]

Both terms on the right are positive semidefinite. Antisymmetry of the
positive cone, together with the two strictly positive coefficients,
forces \(Q\sigma Q=Q\tau Q=0\).

The vanishing positive compression must still be converted into a
statement about the entire support. For any positive semidefinite \(X\),
finite-matrix \(C^*\) positivity supplies \(X=B^\dagger B\). If
\(QXQ=0\), then

\[
\operatorname{Tr}\bigl((BQ)^\dagger(BQ)\bigr)=0,
\]

so \(BQ=0\), and hence \(XQ=B^\dagger BQ=0\). Hermiticity also gives
\(QX=0\). Therefore \(XP=PX=X\). For rank-one \(P\), direct matrix
multiplication yields

\[
PXP=\operatorname{Tr}(XP)P.
\]

Substituting \(X=\sigma\) and using \(\operatorname{Tr}\sigma=1\) gives
\(\sigma=P\); the same argument applies to \(\tau\). At no point are
\(\sigma,\tau\) restricted to rank-one states, diagonal states, or a
finite sample. \(\square\)

In Lean, \texttt{Is\allowbreak Pure} quantifies over the actual \texttt{Density}
type. The declarations \texttt{positive\_\allowbreak compression\_\allowbreak zero},
\texttt{density\_\allowbreak eq\_\allowbreak vector\_\allowbreak of\_\allowbreak compression\_\allowbreak zero}, and
\texttt{vector\allowbreak State\_\allowbreak is\allowbreak Pure} establish, respectively, the
compression-support step, identification using unit trace, and
extremality in the complete set. The invoked matrix-positivity theorems
are explicit mathematical dependencies, not new physical selection
principles. {[}QM317{]}

\section{57.3 The Actual Partial Trace and Gram Purification of Any
Finite State}

Let the discarded auxiliary system have carrier \(\alpha\). For each
\(a\in\alpha\), define the rectangular Kraus matrix

\[
L_a(i,(j,b))=\delta_{ij}\delta_{ab}.
\]

Direct summation gives \(\sum_aL_a^\dagger L_a=I_{\iota\times\alpha}\),
so these matrices define a concrete channel. Its action on an arbitrary
complex matrix is

\[
\operatorname{Tr}_{\alpha}(X)_{ij}
=\sum_aX_{(i,a),(j,a)}.
\]

Thus the partial trace preserves positivity and the complex trace. Every
finite ancillary amplification also preserves positivity on all positive
joint inputs; the check is not restricted to product states. For product
density states, \(\operatorname{Tr}\sigma=1\) further gives
\(\operatorname{Tr}_{\alpha}(\rho\otimes\sigma)=\rho\).

\phantomsection\label{p-pur-02}{}

\textbf{P-PUR-02 (Existence of a purification for every finite complex
density state).} For every \(\rho\in\mathcal D(\iota)\), there exists
\(\Psi\in\mathcal D(\iota\times\iota)\) such that \(\Psi\) is pure in
the complete joint density-state space and satisfies
\(\operatorname{Tr}_{2}\Psi=\rho\).

\textbf{Proof.} From \(\rho\succeq0\), obtain a matrix \(A\) such that
\(\rho=AA^\dagger\), and define the joint vector \(\psi(i,a)=A_{ia}\).
Its outer product \(\Psi=|\psi\rangle\langle\psi|\) is positive, and

\[
\operatorname{Tr}\Psi
=\sum_{i,a}|A_{ia}|^2
=\operatorname{Tr}(AA^\dagger)=1.
\]

P-PUR-01 establishes that \(\Psi\) is an extreme point of the complete
joint state space. Entrywise, the actual partial trace gives

\[
(\operatorname{Tr}_{2}\Psi)_{ij}
=\sum_a A_{ia}\overline{A_{ja}}
=(AA^\dagger)_{ij}=\rho_{ij}.
\]

The existence of \(A\) follows from positivity of the original density
state; it is not an additional assumption that a caller must supply.
\(\square\)

The formal construction \texttt{purification} selects a Gram factor
supplied by PSD existence. The declarations
\texttt{purification\_\allowbreak is\allowbreak Pure} and \texttt{purification\_\allowbreak marginal} prove
purity and the marginal identity, and \texttt{exists\_\allowbreak purification}
combines them. This selection uses \texttt{Classical.\allowbreak choice}, already
included in the foundational audit. It establishes existence and
provides a noncomputable choice in Lean; it does not include a certified
numerical matrix-square-root algorithm. On an empty carrier, there is no
unit-trace density state in the first place. The theorem does not
manufacture states through a hidden nonemptiness assumption. {[}QM317{]}

The Bell vector \(\psi=(|00\rangle+|11\rangle)/\sqrt2\) supplies the
familiar two-dimensional positive control, with marginal \(I/2\). The
concrete Lean witness uses \((1+i)/2\) at its two nonzero positions to
retain rational complex coefficients. This differs from the standard
vector only by a global phase and gives the same Bell projector. The
general theorem covers every finite dimension and every valid initial
state. The small example checks indices and marginals; it cannot replace
the general proof.

\section{57.4 Three Concrete Models and Separate Candidate Principles}

\phantomsection\label{p-opt-01}{}

\textbf{P-OPT-01 (Three concrete presentations in one operational
language).} The complete classical simplex, complete real density-matrix
model, and complete complex density-matrix model each have concrete
states, effects, convex mixtures, products, discards, and their declared
processes. The seven basic laws above are proved from their respective
data as \texttt{classical\_\allowbreak laws}, \texttt{real\_\allowbreak laws}, and
\texttt{complex\_\allowbreak laws}. Purification existence and local tomography are
defined separately as candidate properties, not required fields of a
presentation.

In the classical model, states are all finite probability distributions
and effects are all \([0,1]\)-valued functions. Composite systems
contain all distributions on Cartesian products, discarding is marginal
summation, and processes are concrete stochastic kernels. The real model
uses all real symmetric PSD unit-trace matrices, real effects, and real
Kraus processes. The complex model uses the complex matrix objects
above. Multiplication of product probabilities follows from classical
product summation or the Kronecker-product trace identity; process
affinity follows from the actual linear action. When real matrices are
embedded into complex matrices, their real-valuedness and joint carrier
must be preserved. The full complex state set must not be mistaken for
the real model's state space.

In this language, purification existence means

\[
\forall A\ \forall\rho\in\operatorname{State}(A),\
\exists B\ \exists\Psi\in\operatorname{State}(A\otimes B),\
\operatorname{Pure}(\Psi)\ \land\ \operatorname{discard}_{B}\Psi=\rho.
\]

Local tomography instead requires that two joint states with identical
probabilities for every product effect be the same joint state. The
former concerns pure extensions; the latter concerns identification of
joint states from local-test statistics. They impose different logical
obligations.

\section{57.5 How Countermodels Constrain Quantum Selection Without
Presupposing It}

\phantomsection\label{n-opt-01}{}

\textbf{N-OPT-01 (Basic operational laws do not automatically select a
quantum theory).} The complete classical model satisfies the declared
basic laws, but the nontrivial half-mixture of Chapter 54 has no pure
extension with any finite auxiliary system. In the real quantum model,
the two states of Chapter 55 give identical probabilities for all real
local product effects, yet a valid real joint effect distinguishes them.
Once these witnesses are connected to the same operational presentation,
the first refutes purification existence for that classical model and
the second refutes local tomography for that real model.

The classical mixed state in Chapter 54 still has ordinary
extensions---for example, its product with a deterministic auxiliary
state. What fails is the additional requirement that the extension be
pure. Conversely, if the allowed state set is deliberately reduced to a
singleton containing one mixed state, that state becomes an extreme
point of the reduced set. This negative control shows why quantification
over the complete state space is itself substantive scientific content
and cannot be silently weakened.

Nor may the real quantum counterexample be rewritten as ``real quantum
theory does not admit purification.'' The property rejected here is
local tomography: the two real joint states of Chapter 55 have identical
statistics for every real product test. Once the specified complex local
effects are allowed, their difference becomes visible. This
removal-of-restriction control identifies precisely the local reality
assumption on which the counterexample depends.

A rigorous research path follows: formulate selection principles
independently, test their satisfaction or failure model by model, and
investigate whether UEOT itself or empirical evidence provides an
independent justification for them. The fact that the complex model
performs appropriately among three examined models does not exclude
every alternative theory. Nor does this chapter exhaust all
interpretations of the complete 106 Core propositions by those three
models.

\section{57.6 From Purification Existence to Full Quantum
Reconstruction}

Purification principles in quantum-reconstruction work usually include
essential uniqueness up to reversible transformations on the auxiliary
system. The reconstruction of Chiribella, D'Ariano, and Perinotti also
uses background principles including causality, perfect
distinguishability, ideal compression, local distinguishability, and
pure conditioning. Selecting only the names ``purification existence''
and ``local tomography'' does not make its conclusion a proved
consequence of this project. {[}CDP2010; CDP2011{]}

This chapter supplies the existence construction and the extremality
proof for vector states in a finite complex matrix realization. The
converse-purity and uniqueness gaps left here in edition 3.3 are closed
in Chapter 61: every fully convex-pure state has a normalized vector
representation; arbitrary purifications on the same auxiliary carrier
are related by an actual auxiliary unitary; and different auxiliary
carriers are handled separately through isometric embedding and common
extension. Chapter 60 proves local tomography for arbitrary joint states
in the present complex operational presentation. These results are bound
to their new actual source files and acceptance records; they do not
retrospectively enlarge this chapter's original proof. The complete
reconstruction background, and why UEOT must satisfy its principles,
remain separate bridging obligations.

The original idea that a parent object can support a mixed description
of a child now has a concrete construction and a clear boundary: the
actual partial trace of a pure parent state gives any finite complex
density state, whereas the same assertion does not hold universally in
the complete classical model. Advancing this difference into a physical
selection requires proof or empirical constraints on the source of the
additional principles themselves.

\chapter{58. Joint records, quantum memory, and finite adaptive
protocols}\label{ch-58}

When later actions depend on earlier outcomes, comparing
outcome-forgetting quantum channels can miss real predictive
differences. Joint classical records and quantum memory make a complete
finite adaptive tree a comparable physical object and allow exact
control of record distance over all events.

An instrument maps an input state to several unnormalized branches
carrying outcome labels. Summing the branches produces an unconditional
quantum channel, but discards the outcome information. The complete
record laws of Chapter 49 and the effect distance of Chapter 50
therefore need a common carrier: a joint state retaining both the actual
classical record and the quantum memory associated with each record.
This chapter constructs that carrier from actual Kraus matrices and
extends it to the existing class of arbitrary finite adaptive protocols.

\section{58.1 Kraus Channels with Different Input and Output Types}

Let the input carrier be \(\iota\) and the output carrier \(\omega\). A
finite rectangular Kraus channel specifies
\(K_k\in M_{\omega\times\iota}(\mathbb C)\) satisfying completeness on
the input space:

\[
\sum_kK_k^\dagger K_k=I_\iota.
\]

Its actual operator map is \(\mathcal C(X)=\sum_kK_kXK_k^\dagger\).
Input and output dimensions may differ; the trace identities remain
meaningful on their respective spaces.

\phantomsection\label{p-cq-01}{}

\textbf{P-CQ-01 (Complete positivity, trace preservation, and effect
contraction for rectangular channels).} This map preserves the complex
trace of every complex matrix. For every finite auxiliary carrier
\(\nu\), the actual amplification defined by \(K_k\otimes I_\nu\)
preserves positivity of every \(X\succeq0\). For every output effect
\(E\), the pullback \(\mathcal C^*(E)=\sum_kK_k^\dagger E K_k\) is a
valid input effect, and

\[
\operatorname{Tr}(E\mathcal C(X))=\operatorname{Tr}(\mathcal C^*(E)X),\qquad
d_{\rm eff}(\mathcal C(\rho),\mathcal C(\sigma))\le d_{\rm eff}(\rho,\sigma).
\]

\textbf{Proof.} Each term \(K_kXK_k^\dagger\) is positive for positive
input. Cyclicity of trace and completeness give

\[
\operatorname{Tr}\mathcal C(X)
=\operatorname{Tr}\!\left(\left(\sum_kK_k^\dagger K_k\right)X\right)
=\operatorname{Tr}X.
\]

The amplified Kraus completeness relation is
\(\sum_k(K_k^\dagger K_k)\otimes I=I\), so the same positivity proof
applies to every positive joint input. An entrywise expansion also
identifies this map with the actual identity amplification, rather than
a different map agreeing only on product states.

For effects, \(\mathcal C^*(E)\succeq0\), while
\(I-\mathcal C^*(E)=\mathcal C^*(I-E)\succeq0\). The trace-pairing
identity turns each output test into an input test. Taking the supremum
over all output effects proves contraction. Identity processes and
sequential composition also have actual Kraus realizations: the Kraus
matrices for \(\mathcal C\) followed by \(\mathcal D\) are \(D_lC_k\).
\(\square\)

Complete positivity here includes positive inputs on every finite
auxiliary system. The error measure remains the full-effect distance on
the stated input space. This positivity result does not, by itself, turn
the later defect supremum into a diamond distance. {[}QM317{]}

\section{58.2 A Single Instrument's Record Register and Two Marginals}

Given a complete instrument \(I\) with Kraus matrices \(K_{o,k}\), take
the outcome set \(O\) as the record carrier. Define

\[
V_{o,k}((p,i),j)=\delta_{po}K_{o,k}(i,j).
\]

These are actual rectangular matrices from \(\iota\) to
\(O\times\iota\). Summing over the label yields
\(V_{o,k}^\dagger V_{o,k}=K_{o,k}^\dagger K_{o,k}\), so completeness
follows directly from the original instrument.

\phantomsection\label{p-cq-02}{}

\textbf{P-CQ-02 (Joint block state and marginals of an actual
instrument).} These Kraus matrices define the channel

\[
\mathcal R_I(X)=\bigoplus_{o\in O}\Phi_{I,o}(X),\qquad
\Phi_{I,o}(X)=\sum_kK_{o,k}XK_{o,k}^\dagger.
\]

The block identity holds for every complex input \(X\). For a density
state \(\rho\), \(\Omega_{I,\rho}=\mathcal R_I(\rho)\) is a valid joint
density state, with quantum and record marginals

\[
\operatorname{Tr}_{O}\Omega_{I,\rho}=\sum_o\Phi_{I,o}(\rho),\qquad
\operatorname{Tr}_{\iota}\Omega_{I,\rho}
=\operatorname{diag}\bigl(p_I(o\mid\rho)\bigr).
\]

\textbf{Proof.} Compute the output entry indexed by \((p,i),(q,j)\). The
two label constraints make blocks with \(p\ne q\) vanish; the block with
\(p=q=o\) is precisely the actual branch \(\Phi_{I,o}(X)\). The two
partial traces sum over the label or memory indices, respectively,
giving the displayed identities. The trace of each record block is the
existing instrument probability. \(\square\)

If \(p_I(o\mid\rho)>0\), dividing that block by its trace gives the
existing \texttt{posterior}; no different conditional update is
introduced. If the probability is zero, the entire positive block is
zero, since a positive semidefinite matrix with zero trace vanishes. A
zero-probability event is therefore not assigned an arbitrary ``actual
posterior state.''

This is an operationally explicit representation of a measurement
record. It establishes consistency between the declared instrument and
the public record; it does not derive an objective-collapse mechanism
from discarding or block diagonalization alone.

\section{58.3 Actual Leaf Blocks in Any Finite Adaptive Tree}

The existing protocol type is defined recursively by a stop node or an
action node \(\operatorname{node}(a,\pi_o)\). Each leaf \(\ell\) has its
prescribed chronological word \(\operatorname{word}_\pi(\ell)\);
different outcomes may enter subtrees of different depths. The joint
output uses the actual carrier \(\operatorname{Leaf}(\pi)\times\iota\),
preserving the complete record and terminal memory.

Define a Kraus-path type for each leaf. A stop leaf has one empty path,
with matrix \(I\). At a leaf \((o,\ell')\), a path consists of the
current Kraus index \(k\) and a subtree path. Its chronological matrix
product is

\[
K_{\text{path}}=K_{\text{tail}}K_{a,o,k}.
\]

The total Kraus index is
\(\sum_{\ell:\operatorname{Leaf}(\pi)}\operatorname{Path}(\pi,\ell)\).
Embedding each path matrix in its actual leaf-label block gives the
protocol's rectangular channel. This dependent index retains the leaf
label, path length, and chosen subtree simultaneously; it does not
require all leaves to have equal depth.

\phantomsection\label{p-cq-03}{}

\textbf{P-CQ-03 (Record-and-memory realization of every finite
protocol).} For any current complete finite protocol \(\pi\), actual
instrument family \(I\), and density state \(\rho\), the constructed
joint density state satisfies

\[
\Omega_{\pi,\rho}[\ell,\ell]
=\operatorname{recordOperator}\bigl(I,\operatorname{word}_\pi(\ell),\rho\bigr),
\qquad
\Omega_{\pi,\rho}[\ell,m]=0\quad(\ell\ne m).
\]

The trace of this block is therefore the existing
\texttt{leaf\allowbreak Probability}; the record marginal is the complete leaf
distribution, and the quantum marginal is the sum of all actual terminal
unnormalized memories.

\textbf{Proof.} Induct on the protocol tree. At a stop node, the
identity Kraus matrix preserves the input state. At an action node,
apply subtree completeness for each outcome, then sum over the current
Kraus matrices:

\[
\sum_{o,k}K_{a,o,k}^\dagger
\left(\sum_{\ell',\text{tail}}K_{\text{tail}}^\dagger K_{\text{tail}}\right)
K_{a,o,k}
=\sum_{o,k}K_{a,o,k}^\dagger K_{a,o,k}=I.
\]

For a fixed leaf, the sum of subtree-path conjugations equals its record
operator. Substitution of the current branch then agrees with the
existing chronologically defined \texttt{record\allowbreak Operator}. The label
embedding gives the joint block identities, after which the partial
traces can be taken. This leaf-block relation follows from path products
and tree induction; it is not an input field of the channel
construction. \(\square\)

Retaining the entire unnormalized branch matrix makes this result
stronger than retaining leaf probabilities alone. A subsequent effect
may depend on the complete leaf record and act on its corresponding
quantum memory. The model still concerns the stated finite protocol and
memory carrier, rather than an arbitrary unmodelled external
environment.

\section{58.4 Actual Record Defect, Event Effects, and the Exact TV
Constant}

For rectangular channels with common input and output carriers, define
the actual defect

\[
\delta(\mathcal C,\mathcal D)=
\sup\left(\{0\}\cup\{d_{\rm eff}(\mathcal C(\tau),\mathcal D(\tau)):
\tau\in\mathcal D(\iota)\}\right).
\]

Including \(0\) gives the boundary case of an empty input-state set an
explicit meaning. All actual distances lie in \([0,1]\), so the supremum
is bounded. For record channels, set
\(\delta_{\rm rec}(I,J)=\delta(\mathcal R_I,\mathcal R_J)\). This is
defined from the actual channels, rather than by renaming a
user-supplied tolerance ``the defect.''

For an outcome event \(S\subseteq O\), the joint effect is

\[
E_S=\operatorname{diag}(\mathbf1_S)\otimes I_\iota.
\]

Both it and its complement are positive, and its probability on the
joint block state is exactly the actual record probability of \(S\).

\phantomsection\label{p-cq-04}{}

\textbf{P-CQ-04 (Joint-distance and perturbation control of record
distributions).} Let \(p_{I,\rho}\) be the actual record distribution.
Then

\[
\operatorname{TV}(p_{I,\rho},p_{J,\sigma})
\le d_{\rm eff}(\mathcal R_I(\rho),\mathcal R_J(\sigma))
\le d_{\rm eff}(\rho,\sigma)+\delta_{\rm rec}(I,J).
\]

The same statement applies to the complete finite-protocol joint channel
constructed in §58.3 and its actual leaf distribution. In that case, the
defect is the actual defect of the entire protocol channel.

\textbf{Proof.} For the second inequality, split the triangle inequality
at \(\mathcal R_I(\sigma)\), then use the input-distance contraction of
P-CQ-01 and the supremum property of the defect. For the first
inequality, apply \(E_S\) to every event to obtain

\[
|p(S)-q(S)|\le d_{\rm eff}(\Omega_p,\Omega_q).
\]

For finite normalized distributions, let \(S_+=\{o:p(o)\ge q(o)\}\).
Since \(\sum_o(p(o)-q(o))=0\), the total positive and negative
differences agree, so

\[
p(S_+)-q(S_+)=\frac12\sum_o|p(o)-q(o)|=\operatorname{TV}(p,q).
\]

This actual event attains TV. Thus ``every event difference is at most
\(\varepsilon\)'' is equivalent to ``the TV of the complete
distributions is at most \(\varepsilon\).'' Taking \(S_+\) proves the
first inequality. \(\square\)

\texttt{Record\allowbreak Probability} proves this event equivalence separately; the
joint quantum module constructs the actual effects and verifies the
matching probabilities. This preserves the constant in the bound.
Bounding individual outcomes and then summing directly gives a different
bound that may grow with the number of outcomes; it cannot retain the
constant above without proof. {[}QM317{]}

\section{58.5 Discarding Labels Erases Actual Distinguishability}

\phantomsection\label{n-cq-01}{}

\textbf{N-CQ-01 (Equal unconditional channels do not control record
error).} Use the same two-dimensional quantum memory. Instrument \(I\)
deterministically outputs label \(0\) while preserving the input state;
instrument \(J\) deterministically outputs label \(1\) while preserving
it. Both have complete Kraus realizations. After labels are discarded,
both channels are the identity and their earlier unlabelled channel
defect is zero. Their actual record distributions have TV \(1\), and
their joint record defect is \(1\).

\textbf{Proof.} For \(I\), use the identity Kraus matrix in the
label-\(0\) branch and a zero matrix in the label-\(1\) branch; exchange
the labels for \(J\). Both completeness identities are
\(I^\dagger I=I\), and both unconditional maps return the input matrix.
For any common initial state, the record laws are two distinct Dirac
distributions. The effect \(E_{\{0\}}\) has probabilities \(1\) and
\(0\), so the joint distance is at least \(1\). The full-effect distance
is at most \(1\), and hence equals \(1\). Taking the supremum over
initial states gives the defect. \(\square\)

If a common controller selects the second action from the first label,
the record difference enters the subsequent actual protocol. For
example, after reading label \(o\), it can choose an action that
deterministically outputs \(o\); the second output again distinguishes
the instrument families. This two-step control has been checked using
actual protocol leaf words, the subsequent event effect, and complete
record probabilities: the event that the second outcome is zero has
probabilities 1 and 0, and complete-record TV is 1. Quantum memory
remains the same in this label-deletion control; the proved difference
lies in the actual records and their subsequent selections.

\section{58.6 Established Interfaces and the Next Error Obligation}

This chapter places the existing single-step instruments, complete
finite record laws, and quantum-state errors on one concrete joint
system. That is enough to discuss each finite protocol's entire output
and all-event error, although its actual defect may still be difficult
to compute.

Chapter 62 completes the fixed-depth common causal-controller task posed
here. It constructs actual history-label--quantum-memory channels step
by step, connects them to the original \texttt{Protocol} leaf and word
order, and derives action-record budgets from common-input PSD blocks,
normalized weights, zero blocks, and valid effects. Chapter 65 further
supplies actual absorbing padding for finite stopping trees, the
connection to the original joint leaf state, and budgets weighted by the
probability of remaining active. Those new proofs carry the extension
obligations; this chapter's whole-protocol defect has not simply been
renamed a per-action conclusion. Coherent controllers, arbitrary
additional environments, and unbounded stopping times remain open.

This also identifies a structural requirement for UEOT predictive
objects. When allowed future actions depend on a record, discarding that
record generally changes the available future tests and object
equivalence. The restricted-readout counterexample of Chapter 54 and the
label counterexample here show, from two directions, that the selected
interface determines distinguishability. A complete physical theory must
retain allowed observations and allowed controls as substantive model
content.

\chapter{59. Exact dynamics of noncommutative conditional
persistence}\label{ch-59}

An exactly solvable noncommuting model tests whether the objects used in
conditioned dynamics actually fit together. Explicit Doob trajectories,
success probabilities, and conditional states establish geometric
approach in this model; a positive spectral metric alone does not ensure
mixing in general.

Chapter 56 constructed an actual Doob channel from a given positive left
eigenmetric and proved its intertwining relation with finite conditioned
evolution. The general construction does not guarantee a unique
stationary state or attraction. Here we fix the same rational,
nondiagonal model with noncommuting Kraus matrices used in Chapter 56,
derive closed forms and convergence rates for its actual iterates,
record probabilities, and conditional states, and use a valid
identity-Doob model to test which conclusions do not generalize.

\section{59.1 The Fixed Model and Two State Coordinates}

Retain

\[
r=\frac3{16},\qquad S=\operatorname{diag}(2,1),\qquad T=S^{-1},\qquad
H=S^\dagger S=\operatorname{diag}(4,1),
\]

\[
\rho_* =\frac15\begin{pmatrix}1&2\\2&4\end{pmatrix},\qquad
\sigma_+=\frac12\begin{pmatrix}1&1\\1&1\end{pmatrix}.
\]

Write the source's actual success branch as \(\Phi\). Its verified right
eigendensity is \(\rho_*\), with \(\Phi(\rho_*)=r\rho_*\), while the
left spectral relation is \(\Phi^*(H)=rH\). The complete source
instrument also contains an actual failure branch. This chapter leaves
its Kraus matrices unchanged and does not replace the example by a
simpler diagonal classical model.

The Doob channel is

\[
D(X)=\frac13X+\frac23\operatorname{Tr}(X)\sigma_+.
\]

This formula holds for every complex matrix. For non-Hermitian inputs
the trace may be complex, so its imaginary part cannot be discarded in
advance. For density states, the normalized coordinate transformation is

\[
F(\tau)=\frac{S\tau S^\dagger}{\operatorname{Tr}(H\tau)}.
\]

Since its denominator depends on the input, \(F\) is generally not a
linear channel. The success probability of \(\Phi\), the normalization
factor of \(F\), and trace preservation of \(D\) have different
meanings.

The public module \texttt{Quantum\allowbreak Doob\allowbreak Model} extracts the same model from
the preceding edition's concrete witness, so the general dynamics source
no longer depends on a witness script. The extracted definitions and
proofs retain their provenance; moving them into a public module does
not count as another new physical result. {[}QM317{]}

\section{59.2 Closed Form, Uniqueness, and Exact Contraction of the Full
Doob Orbit}

Let \(a_n=3^{-n}\).

\phantomsection\label{p-qdd-01}{}

\textbf{P-QDD-01 (Exact dynamics of the present Doob channel).} For
every complex matrix \(X\) and every \(n\ge0\),

\[
D^n(X)=a_nX+(1-a_n)\operatorname{Tr}(X)\sigma_+.
\]

The unique invariant density state is \(\sigma_+\). For any density
states \(\tau,\upsilon\),

\[
d_{\rm eff}(D^n\tau,D^n\upsilon)
=a_n d_{\rm eff}(\tau,\upsilon),
\qquad D^n\tau\longrightarrow\sigma_+.
\]

The final limit is convergence of the entire orbit in the usual topology
of the actual finite matrix space.

\textbf{Proof.} Trace preservation gives
\(\operatorname{Tr}D(X)=\operatorname{Tr}X\). The formula holds at
\(n=0\). Applying it to \(D(X)\) and combining coefficients using
\(a_{n+1}=a_n/3\) proves the closed form for all complex operators by
induction. If \(D(\tau)=\tau\), then
\(\tau=\frac13\tau+\frac23\sigma_+\). Rearranging each matrix entry
gives \(\tau=\sigma_+\). This uniqueness argument directly rearranges
matrix entries. Chapter 60 separately proves separation by the
full-effect distance, without changing the present proof.

Subtracting two density-state orbits gives
\(D^n\tau-D^n\upsilon=a_n(\tau-\upsilon)\). For every effect, the
absolute probability difference scales by \(a_n\). When the coefficient
is strictly positive, division by it proves the reverse supremum
inequality; a zero coefficient is handled directly. The full-effect
supremum therefore scales exactly, rather than merely satisfying an
upper bound. Since \(a_n\to0\), continuity of matrix addition and real
scalar multiplication gives the limit of the whole orbit. \(\square\)

This full-effect scaling argument uses only probability pairings and
suprema. It does not hide an identification with another matrix norm
among its assumptions, and it does not relabel convergence of a Cesàro
subsequence as convergence of the whole orbit.

\section{59.3 From Intertwining Back to Actual Success Records}

For any input density state \(\tau\), define

\[
c_\tau=\operatorname{Tr}(H\tau),\qquad
b_\tau=\frac58c_\tau,\qquad
z_n(\tau)=a_n+(1-a_n)b_\tau.
\]

Since \(\tau_{00},\tau_{11}\ge0\) and \(\tau_{00}+\tau_{11}=1\),

\[
c_\tau=4\tau_{00}+\tau_{11},\qquad
1\le c_\tau\le4,\qquad
\frac58\le b_\tau\le\frac52.
\]

Also \(0<a_n\le1\), so \(z_n\) is a convex combination of \(1\) and
\(b_\tau\), giving

\[
\frac58\le z_n(\tau)\le\frac52.
\]

These bounds follow from positivity and unit trace in the complete
density-state space; they are not checks only on diagonal states or
finitely sampled inputs. The filter normalization factor \(c_\tau\) may
exceed \(1\) and is not a physical success probability.

\phantomsection\label{p-qdd-02}{}

\textbf{P-QDD-02 (Closed forms for the actual source success operator
and its mass).} For the same source model,

\[
\Phi^n(\tau)
=r^n\left[a_n\tau+(1-a_n)b_\tau\rho_*\right],
\]

\[
m_n(\tau)=\operatorname{Tr}\Phi^n(\tau)
=r^nz_n(\tau).
\]

Here \(\Phi^n\) is the existing \texttt{success\allowbreak Operator}, and \(m_n\) is
the existing \texttt{success\allowbreak Mass}, also equal to the \texttt{response}
of the actual instrument family to the repeated-success word.

\textbf{Proof.} Chapter 56's intertwining relation gives, for every
complex matrix,

\[
D^n(SXS^\dagger)=r^{-n}S\Phi^n(X)S^\dagger.
\]

Multiply on the left and right by \(T,T^\dagger\), substitute P-QDD-01,
and use

\[
T\sigma_+T^\dagger=\frac58\rho_*.
\]

This yields the pullback formula for every complex \(X\). Taking
\(X=\tau\) and using \(\operatorname{Tr}(S\tau S^\dagger)=c_\tau\) gives
the density-state version. Taking its trace and using unit trace of both
states yields the actual mass identity. \(\square\)

The proof first connects actual operator iterates and only then takes
the trace. It does not define a separate, similar-looking sequence and
treat it as the original record probability. The success mass is
strictly positive everywhere, and

\[
\frac58r^n\le m_n(\tau)\le\frac52r^n.
\]

These are uniform bounds; they do not assert \(m_n=r^n\) for a general
initial state.

\section{59.4 Exact Geometric Approach of the Actual Conditional State}

Since the success mass is strictly positive at every finite \(n\), the
existing \texttt{conditional\allowbreak State} is well defined for every density
input.

\phantomsection\label{p-qdd-03}{}

\textbf{P-QDD-03 (The actual conditional state and its quantitative
convergence).} Write \(\tau_{\mid n}=\Phi^n(\tau)/m_n(\tau)\). Then

\[
\tau_{\mid n}
=\frac{a_n}{z_n(\tau)}\tau
+\frac{(1-a_n)b_\tau}{z_n(\tau)}\rho_*.
\]

Both coefficients are nonnegative and sum to \(1\). Moreover,

\[
d_{\rm eff}(\tau_{\mid n},\rho_*)
=\frac{a_n}{z_n(\tau)}d_{\rm eff}(\tau,\rho_*)
\le\frac85\,3^{-n}d_{\rm eff}(\tau,\rho_*).
\]

The actual conditional-state matrix converges to \(\rho_*\), and its
full-effect distance from that state tends to zero.

\textbf{Proof.} Divide P-QDD-02 by the strictly positive \(r^nz_n\) and
cancel the common factor \(r^n\) to obtain the conditional-state
formula. By the definition of the normalization factor, the second
coefficient is \(1-a_n/z_n\), so

\[
\tau_{\mid n}-\rho_*
=\frac{a_n}{z_n(\tau)}(\tau-\rho_*).
\]

Repeating the full-effect supremum-scaling argument of P-QDD-01 gives
the exact equality. The bound \(z_n\ge5/8\) implies
\(a_n/z_n\le(8/5)a_n\). Since \(z_n\to b_\tau>0\), the coefficient
\(a_n/z_n\to0\), giving both actual limits: the matrix limit and the
effect-distance limit. \(\square\)

In this specified model, normalized source conditioning has a nonlinear
denominator, while the Doob orbit is generated by a linear
trace-preserving channel. Intertwining makes the two comparable. It does
not make their physical probabilities or stationary states
interchangeable.

\section{59.5 Different Limits for the Conditional Steady State and
Unconditional Success Mass}

\phantomsection\label{p-qdd-04}{}

\textbf{P-QDD-04 (Decay of finite-success probabilities and the role of
the initial state).} For every density state \(\tau\) and each finite
\(n\), one has \(m_n(\tau)>0\), but

\[
m_n(\tau)\longrightarrow0.
\]

For \(\tau=\rho_*\), \(c_{\rho_*}=8/5\) and \(b_{\rho_*}=1\), so
\(m_n(\rho_*)=r^n\). For the uniform initial state \(\tau_0=I/2\),

\[
c_{\tau_0}=\frac52,\qquad b_{\tau_0}=\frac{25}{16},\qquad
m_1(\tau_0)=\frac{33}{128}\ne\frac3{16}.
\]

\textbf{Proof.} The sequence \(z_n(\tau)\) converges to the finite
positive number \(b_\tau\), while \(0<r<1\); hence \(r^nz_n\to0\). The
special initial-state identities follow by substituting the actual
matrix entries into the general mass formula. \(\square\)

A conditional distribution can therefore approach a steady state while
finite records satisfying the condition become increasingly rare. What
is proved here is the limit of a sequence of finite-event probabilities.
This chapter does not construct an infinite-path measure or an
infinite-time conditioning event, and it does not describe the
conditional limit as unconditional permanent existence.

This distinction is essential to UEOT's original interest in
persistence. Increasing stability of an organization among surviving
samples does not, by itself, show that the overall survival probability
retains a positive lower bound. To ask which controls can increase
lifetime or change success rates, one must also specify the available
actions, success event, constraints, and costs. The present solvable
uncontrolled model does not select an objective function that nature
must optimize. {[}S1; S2; S4; QM317{]}

\section{59.6 A Positive Spectral Metric Does Not Guarantee a Unique
Attractor}

\phantomsection\label{n-qdd-01}{}

\textbf{N-QDD-01 (A valid identity-Doob negative control).} There exists
an actual complete source instrument whose success branch strictly
decreases the trace of every nonzero PSD input and has a positive left
eigenmetric, but whose Doob channel is the identity. Every density state
is stationary, so there is no common unique attracting state. The
success probability of a normalized density input is always \(1/4\).

\textbf{Construction and proof.} In two dimensions, take one success
Kraus matrix \(I/2\) and three failure Kraus matrices \(I/2\), filling
any remaining required indices with zero matrices. The sum of
\(K^\dagger K\) for the four nonzero matrices is \(I\), so this is a
complete instrument. Its success branch is \(\Phi(X)=X/4\). Choose
\(H=I\), \(S=T=I\), and \(r=1/4\). The left spectral relation follows
immediately, and the actual Doob construction of Chapter 56 gives
\(D(X)=X\).

Every density state's entire Doob orbit consequently remains at that
state. Two distinct density states give two distinct stationary states.
If all initial states converged to one matrix, uniqueness of the limit
of a constant orbit would force those two initial states to coincide, a
contradiction. The source success mass is still \((1/4)^n\), and the
actual finite conditional state always equals the input state.
\(\square\)

The counterexample retains the success branch, complete instrument,
strict decay, positive eigenmetric, and actual Doob construction. It
refutes precisely an inference of uniqueness or mixing from those
conditions alone. A general convergence theorem needs explicit
irreducibility, primitivity, contraction, or spectral conditions and
their proofs. The original GOA picture of a unique fixed point cannot be
made a universal conclusion about every object.

\section{59.7 From a Concrete Dynamics Kernel to Broader Physical
Connections}

This chapter establishes finite dynamics for repeated application of the
specified success branch in one noncommutative model, together with
long-time limits of its Doob orbit and conditional state. Existence of
general positive spectral objects, identification of spectral radii,
uniqueness of general quasistationary states, and infinite conditioned
processes remain separate obligations. They cannot be extrapolated from
the formulas of this solvable model. Physical implementation of the
finite model, experimental readout of its success event, and parameter
calibration also require separate verification.

Together with the preceding two chapters, the kernel gains three
complementary capabilities: representing any finite complex mixed state
by an actual pure state and partial trace; retaining protocol history
and quantum memory in a joint record register; and quantitatively
explaining conditional persistence in a specified model through actual
spectral intertwining. All three hold within the declared quantum
realization. The resulting research priorities are rigorous logical
connections between selection principles, composable record-error
budgets, and dynamical extensions under explicit general conditions,
followed by experimentally implementable observables. The finite
common-controller budgets are developed later in Chapters 62 and 65; the
broader physical and dynamical obligations remain. In this way, model
counterexamples, formal proofs, and empirical requirements jointly
constrain UEOT's explanatory aims.

\chapter{60. Separation by all effects and universal local tomography in
the complex model}\label{ch-60}

When does the all-effect distance actually separate density matrices?
Valid effects first span all Hermitian operators; real and imaginary
matrix entries then recover the joint state, establishing local
tomography for arbitrary joint states in the finite complex model.

Chapter 50 constructed the full-effect distance from actual Born
probabilities and proved probability bounds, symmetry, the triangle
inequality, and channel contraction. Its zero-distance conclusion
stopped precisely at ``all effects give the same probabilities.''
Chapter 55 showed that, in the complete real quantum model, even the
statistics of all local product effects may fail to determine the joint
state. Chapter 57 presented the complete complex model in the same
operational language, leaving the constructive obligation to establish
local tomography for all its joint states.

This chapter closes both obligations in the original complete finite
complex density-state space. The essential step is not to treat matrix
entries directly as measurable quantities. We first express any
Hermitian test as a real linear combination of valid effects, then
recover the matrix from actual effect probabilities. The joint result
holds for arbitrary density states, including nonproduct inputs. Its
information consists of joint probabilities for effects on both sides,
including correlations, rather than merely the two marginals.

\section{60.1 Complete States, Complete Effects, and the Real-Valued
Pairing}

Fix an arbitrary finite carrier \(\iota\). The matrix, state, and effect
spaces are

\[
\begin{aligned}
M_\iota&=M_\iota(\mathbb C),\\
\mathcal D(\iota)&=\{\rho\in M_\iota:\rho\succeq0,\ \operatorname{Tr}\rho=1\},\\
\mathcal E(\iota)&=\{E\in M_\iota:E\succeq0,\ I-E\succeq0\}.
\end{aligned}
\]

Positivity uses the full positive-semidefinite cone of complex matrices.
States are not restricted to diagonal states, vector states, supplied
Gram factors, or finite samples; effects are not restricted to tests
realizable by one chosen instrument library. Lean's \texttt{Density} and
\texttt{Effect} are the actual types already in use. The new proofs do
not change their definitions.

For any matrices \(H,X\), write the real-valued trace pairing as

\[
\langle H,X\rangle_{\mathbb R}
=\operatorname{Re}\operatorname{Tr}(HX).
\]

The effect probability is
\(p_\rho(E)=\langle E,\rho\rangle_{\mathbb R}\). This pairing is
real-linear in both arguments. Multiplying complex matrices by complex
scalars and then taking the real part does not produce a complex-linear,
complex-valued probability function. All probability extensions below
therefore use real coefficients.

The positive-operator order \(H\preceq K\) means that \(K-H\) is
positive semidefinite. For Hermitian \(H\), the finite-matrix operator
norm gives

\[
-\|H\|_{\rm op}I\preceq H\preceq\|H\|_{\rm op}I.
\]

This is the matrix operator norm on complex Euclidean space; the
relevant Lean proof section explicitly uses
\texttt{Matrix.\allowbreak Norms.\allowbreak L2\allowbreak Operator}. These order bounds follow from the
pinned finite-matrix \(C^*\) dependencies. They are not additional
physical spectral assumptions, and they do not require the caller to
supply a separating effect in advance. {[}QM318{]}

\section{60.2 Every Hermitian Operator Lies in the Real Span of Valid
Effects}

\phantomsection\label{p-sep-01}{}

\textbf{P-SEP-01 (Representation of a Hermitian operator by valid
effects).} For every finite carrier \(\iota\) and every Hermitian matrix
\(H\), there exist a real number \(r>0\) and an actual effect
\(E\in\mathcal E(\iota)\) such that

\[
H=2rE-rI.
\]

Consequently, every Hermitian matrix lies in the real linear span of all
effect matrices.

\textbf{Proof.} Set \(r=\|H\|_{\rm op}+1\). Nonnegativity of the norm
gives \(r>0\). Adding the positive-semidefinite \(I\) to each of the
preceding order bounds yields

\[
H+rI\succeq0,\qquad rI-H\succeq0.
\]

Define

\[
E=\frac{H+rI}{2r}.
\]

The coefficient \((2r)^{-1}\) is strictly positive, so \(E\succeq0\).
The effect's upper bound must also be verified; positivity alone is
insufficient. The actual matrix identity gives

\[
I-E=\frac{rI-H}{2r}\succeq0.
\]

This establishes both positivity fields required by \texttt{Effect}.
Since \(2r\ne0\), rearranging the definition gives \(H=2rE-rI\). The
identity matrix is itself a valid effect, so \(H\) is a real linear
combination of two effect matrices. \(\square\)

``Span'' does not mean that every Hermitian matrix is a convex mixture
of effects. The coefficient \(-r\) may be negative. The reconstruction
formula allows signed real arithmetic on probabilities from different
effects; it does not define that arithmetic as a new random-mixture
experiment.

In the formalization, \texttt{bounded\allowbreak Hermitian\allowbreak Effect} first handles any
supplied two-sided order bounds. \texttt{bounded\allowbreak Hermitian\allowbreak Effect\_\allowbreak affine}
proves the actual matrix identity;
\texttt{exists\_\allowbreak effect\_\allowbreak affine\_\allowbreak of\_\allowbreak hermitian} derives the necessary
bounds from \(H\) itself; and
\texttt{hermitian\_\allowbreak mem\_\allowbreak real\_\allowbreak span\_\allowbreak effects} states the conclusion
explicitly in \texttt{Submodule.\allowbreak span\ ℝ}. The use of real coefficients
is part of the declaration. {[}QM318{]}

The construction remains meaningful for an empty carrier. Its only
matrix is zero, its identity matrix is also zero, and one can still
choose \(r=1\). Later assumptions about density states automatically
exclude fictitious normalized states on an empty carrier; see §60.6.

\section{60.3 Agreement on Every Effect Forces Equality of the Actual
Density States}

First obtain a reusable intermediate step from P-SEP-01. If any two
matrices \(X,Y\) satisfy

\[
\langle E,X\rangle_{\mathbb R}
=\langle E,Y\rangle_{\mathbb R}
\quad\text{for all }E\in\mathcal E(\iota),
\]

that is, for every valid effect, then their pairings with any Hermitian
\(H\) also agree. Indeed, write \(H=2rE-rI\). Real linearity gives

\[
\langle H,X\rangle_{\mathbb R}
=2r\langle E,X\rangle_{\mathbb R}
-r\langle I,X\rangle_{\mathbb R},
\]

and similarly for \(Y\). The hypothesis covers both the actual effect
\(E\) and the identity effect \(I\), so the two expressions coincide.
This intermediate step does not require \(X,Y\) to be density states.
Hermiticity is needed for the valid-effect representation of the test
\(H\).

\phantomsection\label{p-sep-02}{}

\textbf{P-SEP-02 (Separation by all effects and the matrix meaning of
zero distance).} For any \(\rho,\sigma\in\mathcal D(\iota)\),

\[
\rho=\sigma
\quad\Longleftrightarrow\quad
\forall E\in\mathcal E(\iota),\ p_\rho(E)=p_\sigma(E).
\]

For the full-effect distance actually defined in Chapter 50,

\[
d_{\rm eff}(\rho,\sigma)
=\sup_{E\in\mathcal E(\iota)}
|p_\rho(E)-p_\sigma(E)|,
\]

one accordingly has \(d_{\rm eff}(\rho,\sigma)=0\) if and only if
\(\rho=\sigma\).

\textbf{Proof.} Equal density states plainly give equal probabilities.
Conversely, set \(D=\rho-\sigma\). Positivity of both states implies
their Hermiticity, so \(D\) is Hermitian. The real-linear extension
above can be applied with \(H=D\), giving

\[
\langle D,\rho\rangle_{\mathbb R}
=\langle D,\sigma\rangle_{\mathbb R}.
\]

Subtracting in the right argument of the pairing yields

\[
\operatorname{Re}\operatorname{Tr}(D^\dagger D)
=\operatorname{Re}\operatorname{Tr}(DD)=0.
\]

The matrix \(D^\dagger D\) is positive semidefinite, and its trace has
zero imaginary part. Thus its full complex trace is zero. The
finite-matrix identity

\[
\operatorname{Tr}(D^\dagger D)
=\sum_{i,j}|D_{ij}|^2
\]

shows that every entry vanishes. Hence \(D=0\), or \(\rho=\sigma\). The
formal proof uses
\texttt{Matrix.\allowbreak trace\_\allowbreak conj\allowbreak Transpose\_\allowbreak mul\_\allowbreak self\_\allowbreak eq\_\allowbreak zero\_\allowbreak iff} for this
step, after separately justifying the passage from zero real part to
zero full trace.

The distance conclusion uses the supremum property already proved in
Chapter 50: distance is zero exactly when every effect has equal
probabilities. Combining it with the separation result gives the
equivalence. \(\square\)

The new public declaration
\texttt{Quantum\allowbreak Separation.\allowbreak effect\allowbreak Distance\_\allowbreak eq\_\allowbreak zero\_\allowbreak iff\_\allowbreak density\_\allowbreak eq} is
named separately from the earlier
\texttt{UEOTQM.\allowbreak effect\allowbreak Distance\_\allowbreak eq\_\allowbreak zero\_\allowbreak iff}. The earlier theorem still
states agreement of all effect probabilities precisely; the new theorem
supplies the previously missing connection to equality of the actual
densities. It requires no rewriting of earlier contraction,
conditioning, or perturbation proofs, and does not itself assert
identification with another matrix norm. {[}QM318{]}

\section{60.4 From Four Valid Product Tests to Hermitian Products}

Now fix two arbitrary finite carriers \(\iota,\nu\), with actual joint
carrier \(\iota\times\nu\). For local effects \(E,F\), the matrix of the
existing \texttt{Real\allowbreak Local\allowbreak Tomography.\allowbreak product\allowbreak Effect} is \(E\otimes F\).
That historical module name does not restrict the inputs to real
matrices: the definition accepts arbitrary complex effects. The
real-model restriction is carried separately by a real-valuedness field.

Validity of the product effect follows from Kronecker positivity and the
identity

\[
I_{\iota\times\nu}-E\otimes F
=(I_\iota-E)\otimes I_\nu
+E\otimes(I_\nu-F).
\]

Both terms on the right are positive semidefinite. The joint test is
therefore an actual effect whose positivity and upper bound have both
been proved.

Suppose two arbitrary joint density states \(\rho,\sigma\) give the same
probabilities for all local product effects. For any local Hermitian
matrices \(H,K\), apply P-SEP-01 separately to write

\[
H=2rE-rI_\iota,\qquad K=2sF-sI_\nu,
\qquad r,s>0.
\]

Distributivity of the actual Kronecker product gives a four-term real
expansion:

\[
\begin{aligned}
H\otimes K={}&4rs\,E\otimes F
-2rs\,E\otimes I_\nu\\
&-2rs\,I_\iota\otimes F
+rs\,I_\iota\otimes I_\nu.
\end{aligned}
\]

The hypothesis covers the four actual local-test combinations \((E,F)\),
\((E,I_\nu)\), \((I_\iota,F)\), and \((I_\iota,I_\nu)\). Real linearity
of the real-valued pairing applied to those terms gives

\[
\langle H\otimes K,\rho\rangle_{\mathbb R}
=\langle H\otimes K,\sigma\rangle_{\mathbb R}.
\]

This is the actual conclusion of
\texttt{hermitian\_\allowbreak product\_\allowbreak pairing\_\allowbreak eq}. The four tests depend on the
representations of \(H,K\). The proof does not assume that \(\rho\) or
\(\sigma\) can be written as a probabilistic mixture of product states.
It expands test operators without imposing separability on the inputs.

\section{60.5 Recovering Real Parts, Imaginary Parts, and Every Joint
Matrix Entry}

Let \(e_{ij}\) be the matrix unit whose \((i,j)\) entry is one and whose
other entries are zero. Define the Hermitian matrices

\[
S_{ij}=e_{ij}+e_{ji},\qquad
T_{ij}=\mathrm i(e_{ij}-e_{ji}).
\]

Here \(\mathrm i\) is the imaginary unit. Conjugate transposition gives
\(S_{ij}^\dagger=S_{ij}\) and \(T_{ij}^\dagger=T_{ij}\). For any
Hermitian \(X\), write \(X_{ij}=u+\mathrm i v\); then
\(X_{ji}=u-\mathrm i v\). The trace pairings of matrix units give

\[
\begin{aligned}
\langle S_{ij},X\rangle_{\mathbb R}
&=2\operatorname{Re}X_{ij},\\
\langle T_{ij},X\rangle_{\mathbb R}
&=2\operatorname{Im}X_{ij}.
\end{aligned}
\]

The second sign is positive:
\(\operatorname{Tr}(T_{ij}X)=\mathrm i(X_{ji}-X_{ij})=2v\). If one
instead used \(\mathrm i(e_{ji}-e_{ij})\), the reconstruction sign would
have to reverse. The declaration \texttt{pairing\_\allowbreak imaginary\allowbreak Unit} fixes
the present convention and prevents the two conventions from being
mixed.

These matrix units on the two local carriers also satisfy

\[
S_{(i,a),(j,b)}
=\frac12\bigl(S_{ij}\otimes S_{ab}
-T_{ij}\otimes T_{ab}\bigr),
\]

\[
T_{(i,a),(j,b)}
=\frac12\bigl(T_{ij}\otimes S_{ab}
+S_{ij}\otimes T_{ab}\bigr).
\]

They can be checked by expansion. In the first identity, the cross terms
\(e_{ij}\otimes e_{ba}\) and \(e_{ji}\otimes e_{ab}\) cancel, leaving
the joint matrix unit and its transpose. In the second, what remains is
\(\mathrm i\) times their difference. All external coefficients are
still real.

These identities do not require \(i\ne j\) or \(a\ne b\). For example,
\(S_{ii}=2e_{ii}\) and \(T_{ii}=0\). When both index pairs repeat, the
first identity is still exactly twice the joint diagonal matrix unit,
and both sides of the second are zero. For Hermitian \(X\), the diagonal
imaginary parts vanish in any case. These boundaries are included in the
actual general declarations.

\phantomsection\label{p-clt-01}{}

\textbf{P-CLT-01 (Universal local tomography of the complete finite
complex model).} For any finite carriers \(\iota,\nu\) and any
\(\rho,\sigma\in\mathcal D(\iota\times\nu)\), if

\[
\forall E\in\mathcal E(\iota),\ 
\forall F\in\mathcal E(\nu),\qquad
p_\rho(E\otimes F)=p_\sigma(E\otimes F),
\]

then \(\rho=\sigma\). Hence \texttt{Operational\allowbreak Systems.\allowbreak complex\allowbreak Model},
already constructed in Chapter 57, satisfies the complete
\texttt{Local\allowbreak Tomography} principle of the original interface.

\textbf{Proof.} Section 60.4 extends agreement on product-effect
probabilities to agreement on all local Hermitian-product pairings. For
arbitrary joint row and column indices \((i,a),(j,b)\), use the product
identity for \(S\) and real linearity to obtain

\[
\langle S_{(i,a),(j,b)},\rho\rangle_{\mathbb R}
=\langle S_{(i,a),(j,b)},\sigma\rangle_{\mathbb R}.
\]

Both states are Hermitian, so the matrix-unit reconstruction formula
equates the real parts of the corresponding joint entries. Applying the
product identity for \(T\) in the same way equates their imaginary
parts. Every complex entry therefore agrees, and the two actual
\texttt{Density} objects are equal.

The joint-state type of the original \texttt{complex\allowbreak Model} is exactly
this complete joint density space; its local effects and product
probabilities are also the definitions used here. Substitution of the
theorem into the universal quantifiers of the original
\texttt{Local\allowbreak Tomography} establishes the declared model principle. No
replacement presentation with a smaller state space is introduced.
\(\square\)

Thus, in the complete complex matrix model, the joint statistics of all
product effects identify the entire joint state. The result requires
neither prior access to a global entangled effect nor an invalid
strengthening from knowledge of the two marginals alone. Designing a
finite test set for a particular device, estimating sampling error,
controlling ill-conditioned reconstruction, and proving completeness of
a measurement library remain tasks with additional inputs.

\section{60.6 Empty Carriers, Singleton Carriers, and Unequal Local
Dimensions}

A finite type may be empty. The trace sum on an empty carrier is zero,
so no unit-trace density state exists there; the concrete witness
\texttt{empty\_\allowbreak carrier\_\allowbreak has\_\allowbreak no\_\allowbreak density} retains this fact. The general
theorems accept arbitrary finite carriers. Statements about density
inputs remain logically valid over an empty state domain, without a
hidden \texttt{Nonempty} assumption that would manufacture a state.

On a singleton carrier, the only matrix entry is the trace. Unit trace
therefore forces that entry to equal one, so any two density states
coincide. The witness \texttt{singleton\_\allowbreak densities\_\allowbreak equal} follows from
actual unit trace, rather than from an undeclared nontrivial
two-dimensional test.

The two local dimensions may differ. Their indices are separated from
the outset as \(\iota\) and \(\nu\), with product carrier
\(\iota\times\nu\); every distributive and matrix-unit identity
preserves both index sets. The witness
\texttt{unequal\_\allowbreak local\_\allowbreak dimensions\_\allowbreak tomography} instantiates the general
theorem at \(\operatorname{Fin}2\times\operatorname{Fin}3\). It checks
the type boundary; the general conclusion still holds for any two finite
dimensions.

These witnesses are counted separately from the general theorems.
Empty-state exclusion inherits an existing result; the singleton
calculation and \(2\times3\) instance check their respective boundaries.
They are not three new laws of nature. {[}QM318{]}

\section{60.7 Diagonal-Effect Negative Controls and the Inherited
Obstruction in the Complete Real Model}

The quantifier over all effects cannot be reduced to ``some commonly
used effects.'' To check both real and imaginary coherence information,
take the two-dimensional Hermitian involutions

\[
X=\begin{pmatrix}0&1\\1&0\end{pmatrix},\qquad
Y=\begin{pmatrix}0&-\mathrm i\\\mathrm i&0\end{pmatrix}.
\]

They satisfy \(X^2=Y^2=I\) and
\(\operatorname{Tr}X=\operatorname{Tr}Y=0\). In the present convention,
\(T_{01}\) is \(-Y\), so using standard \(Y\) here does not alter the
reconstruction sign fixed in §60.5.

\phantomsection\label{n-sep-01}{}

\textbf{N-SEP-01 (Even all diagonal effects fail to identify coherent
states).} Define

\[
\rho_{x,\pm}=\frac12(I\pm X),\qquad
\rho_{y,\pm}=\frac12(I\pm Y).
\]

All four matrices are actual density states. For every valid diagonal
effect \(E\),

\[
p_{\rho_{x,+}}(E)=p_{\rho_{x,-}}(E),\qquad
p_{\rho_{y,+}}(E)=p_{\rho_{y,-}}(E).
\]

Nevertheless, the two states in each pair are distinct. Actual effects
give probabilities one and zero for each pair, and

\[
d_{\rm eff}(\rho_{x,+},\rho_{x,-})
=d_{\rm eff}(\rho_{y,+},\rho_{y,-})=1.
\]

\textbf{Construction and proof.} For any Hermitian involution \(A\), the
operators \(P_\pm=(I\pm A)/2\) are complementary orthogonal projections.
Direct multiplication gives \(P_\pm^2=P_\pm\) and \(P_+P_-=0\);
Hermiticity and \(P_\pm=P_\pm^\dagger P_\pm\) give positivity. For
\(A=X,Y\), zero trace and the trace two of the two-dimensional identity
further imply \(\operatorname{Tr}P_\pm=1\), yielding the four density
states above.

Both diagonal entries of each state are \(1/2\). Write an arbitrary
diagonal effect as \(E=\operatorname{diag}(e_0,e_1)\). Its probability
depends only on those diagonal entries:

\[
\operatorname{Re}\operatorname{Tr}(E\rho)
=\frac12(e_0+e_1).
\]

Hermiticity of a valid effect makes \(e_0,e_1\) real, while positivity
of the effect and its complement gives \(0\le e_0,e_1\le1\). The
equal-probability conclusion therefore covers every valid diagonal
effect, not merely two selected examples.

Now take \(E_x=(I+X)/2\) and \(E_y=(I+Y)/2\). Their
complementary-positive-projection properties also prove they are valid
effects. From \(P_+^2=P_+\) and \(P_+P_-=0\), the respective
probabilities are one and zero, directly distinguishing both pairs. Each
full-effect distance is at least the probability difference one of this
valid test; Chapter 50 bounds it above by one. The distance is therefore
exactly one. \(\square\)

The diagonal measurement family is entirely valid within its declared
scope, but a statistical equivalence class for that family can contain
distinct complete density states. This counterexample retains the loss
of both real and imaginary off-diagonal information; testing only one
coherence direction would not establish completeness of the measurement
library.

The complete real-model obstruction of Chapter 55 also remains valid,
and is not reduced to a diagonal library. Its two valid real joint
states are

\[
\omega_\pm=\frac14(I\pm Y\otimes Y).
\]

Although \(Y\otimes Y\) is formed from two local matrices with imaginary
entries, the joint matrix itself is a real Hermitian involution. For any
real local effects \(E,F\), real Hermiticity gives symmetry and thus
\(\operatorname{Tr}(EY)=\operatorname{Tr}(FY)=0\). Consequently,

\[
p_{\omega_+}(E\otimes F)
=p_{\omega_-}(E\otimes F)
=\frac14\operatorname{Tr}E\operatorname{Tr}F.
\]

The valid joint effect \((I+Y\otimes Y)/2\) still distinguishes the
states with probabilities one and zero. If the complex local effect
\((I+Y)/2\) is allowed, its product on the two sides has probabilities
\(1/2\) and zero. These are already numbered and proved results from
Chapter 55, used here only for inherited cross-checking.
\texttt{full\_\allowbreak real\_\allowbreak model\_\allowbreak not\_\allowbreak local\allowbreak Tomography} calls the original
\texttt{Operational\allowbreak Systems.\allowbreak real\_\allowbreak not\_\allowbreak local\allowbreak Tomography} directly. It
neither changes the real-model fields nor counts the inherited result as
a new numbered counterexample in this chapter. {[}QM316; QM317; QM318{]}

This identifies the concrete role of imaginary-part tests in complex
local tomography. The local Hermitian tests of type \(T\) in §60.5 can
be expressed as real combinations of valid effects within the complex
model. The real model does not thereby permit the silent addition of
local effects with imaginary matrix entries. The models share actual
composition and trace operations, but their allowed tests differ.

\section{60.8 Public Declarations, Actual Evidence, and the Remaining
Physical Connection}

The four numbered results correspond to actual code as follows. The
first two groups of declarations are in
\texttt{UEOTQM.\allowbreak Quantum\allowbreak Separation}, the third in
\texttt{UEOTQM.\allowbreak Complex\allowbreak Local\allowbreak Tomography}, and the final witness group in
\texttt{UEOTQM.\allowbreak Quantum\allowbreak Separation\allowbreak Examples}. The namespace and local
declaration together determine each stable name.

\begin{longtable}[]{@{}
  >{\raggedright\arraybackslash}p{(\columnwidth - 2\tabcolsep) * \real{0.3000}}
  >{\raggedright\arraybackslash}p{(\columnwidth - 2\tabcolsep) * \real{0.7000}}@{}}
\toprule\noalign{}
\begin{minipage}[b]{\linewidth}\raggedright
Numbered result or proof step
\end{minipage} & \begin{minipage}[b]{\linewidth}\raggedright
Public declarations and actual responsibility
\end{minipage} \\
\midrule\noalign{}
\endhead
\bottomrule\noalign{}
\endlastfoot
P-SEP-01: two-sided positivity and effect representation &
\texttt{bounded\allowbreak Hermitian\allowbreak Effect};
\texttt{bounded\allowbreak Hermitian\allowbreak Effect\_\allowbreak matrix};
\texttt{bounded\allowbreak Hermitian\allowbreak Effect\_\allowbreak affine};
\texttt{exists\_\allowbreak effect\_\allowbreak affine\_\allowbreak of\_\allowbreak hermitian} \\
P-SEP-01: span over the real field &
\texttt{hermitian\_\allowbreak mem\_\allowbreak real\_\allowbreak span\_\allowbreak effects} \\
P-SEP-02: real extension and density separation &
\texttt{pairing\_\allowbreak eq\_\allowbreak of\_\allowbreak all\_\allowbreak effect\_\allowbreak pairings\_\allowbreak eq};
\texttt{pairing\_\allowbreak sub\_\allowbreak right};
\texttt{density\_\allowbreak eq\_\allowbreak of\_\allowbreak all\_\allowbreak effect\_\allowbreak probabilities\_\allowbreak eq};
\texttt{density\_\allowbreak eq\_\allowbreak iff\_\allowbreak all\_\allowbreak effect\_\allowbreak probabilities\_\allowbreak eq} \\
P-SEP-02: actual meaning of zero distance &
\texttt{effect\allowbreak Distance\_\allowbreak eq\_\allowbreak zero\_\allowbreak iff\_\allowbreak density\_\allowbreak eq} \\
P-CLT-01: local matrix units & \texttt{symmetric\allowbreak Unit};
\texttt{imaginary\allowbreak Unit}; their \texttt{\_\allowbreak hermitian} theorems;
\texttt{pairing\_\allowbreak symmetric\allowbreak Unit}; \texttt{pairing\_\allowbreak imaginary\allowbreak Unit} \\
P-CLT-01: real product expansion & \texttt{sub\_\allowbreak kronecker};
\texttt{kronecker\_\allowbreak sub}; \texttt{symmetric\allowbreak Unit\_\allowbreak product};
\texttt{imaginary\allowbreak Unit\_\allowbreak product};
\texttt{hermitian\_\allowbreak product\_\allowbreak pairing\_\allowbreak eq} \\
P-CLT-01: arbitrary joint states and the original model &
\texttt{density\_\allowbreak eq\_\allowbreak of\_\allowbreak all\_\allowbreak product\_\allowbreak effect\_\allowbreak probabilities\_\allowbreak eq};
\texttt{complex\_\allowbreak local\allowbreak Tomography} \\
N-SEP-01: equal statistics for all diagonal effects &
\texttt{diagonal\_\allowbreak effect\_\allowbreak pairing\_\allowbreak eq};
\texttt{all\_\allowbreak diagonal\_\allowbreak effects\_\allowbreak agree\_\allowbreak x};
\texttt{all\_\allowbreak diagonal\_\allowbreak effects\_\allowbreak agree\_\allowbreak y} \\
N-SEP-01: actual distinction and distance &
\texttt{x\allowbreak Test\_\allowbreak probability\_\allowbreak plus}, \texttt{x\allowbreak Test\_\allowbreak probability\_\allowbreak minus};
the corresponding \texttt{y\allowbreak Test} probability theorems;
\texttt{x\allowbreak Plus\_\allowbreak ne\_\allowbreak x\allowbreak Minus}, \texttt{y\allowbreak Plus\_\allowbreak ne\_\allowbreak y\allowbreak Minus};
\texttt{x\_\allowbreak effect\allowbreak Distance\_\allowbreak exact}, \texttt{y\_\allowbreak effect\allowbreak Distance\_\allowbreak exact} \\
Dimension boundaries and inherited-principle cross-checks &
\texttt{empty\_\allowbreak carrier\_\allowbreak has\_\allowbreak no\_\allowbreak density};
\texttt{singleton\_\allowbreak densities\_\allowbreak equal};
\texttt{unequal\_\allowbreak local\_\allowbreak dimensions\_\allowbreak tomography};
\texttt{full\_\allowbreak complex\_\allowbreak model\_\allowbreak local\allowbreak Tomography};
\texttt{full\_\allowbreak real\_\allowbreak model\_\allowbreak not\_\allowbreak local\allowbreak Tomography} \\
\end{longtable}

The two new public modules contain 20 explicit \texttt{theorem}
declarations and 3 definitions. The separate concrete witness contains
19 \texttt{theorem} declarations, 8 definitions, and 1 abbreviation. The
actual focused receipts for that stage check the named axiom outputs for
all 23 public declarations and all 28 witness declarations, binding them
to the source bytes that compiled successfully. The permitted
foundations are only \texttt{propext}, \texttt{Classical.\allowbreak choice}, and
\texttt{Quot.\allowbreak sound}. Ordinary linter warnings remain in the actual
compilation. The records are stored in
\texttt{physics/\allowbreak qm/\allowbreak kernel/\allowbreak evidence/\allowbreak v31.\allowbreak 8/\allowbreak A}; integrated acceptance,
public-source identity, and remote-release status are recorded
separately in the unified v31.8 receipts. Focused success, integrated
acceptance, and actual remote identity are distinct evidence facts.
Definitions, witnesses, and numbered results are also distinct counts.
{[}QM318{]}

Scientifically, this chapter completes identifiability inside the
supplied complex quantum model: all effects distinguish complete density
states, and complete joint statistics of local products distinguish
complete joint states. In the predictive-object language, equal
full-effect responses can now identify the same density object within an
explicit complex matrix realization. For a restricted instrument
library, the quotient must still be defined through the responses
actually available; N-SEP-01 retains the possibility of a strictly
coarser quotient.

The finite complex vector spaces, complete PSD state cone, complete
effect cone, tensor composition, and trace probabilities are the already
declared mathematical and quantum-sector inputs. The results belong
principally to standard quantum structure P and actual operational
connections B, with concrete probabilities carrying O. Requiring nature
to satisfy local tomography would additionally require registering and
justifying that selection principle and its background under S. The
frozen 106 Core items and four generators have not thereby acquired a
``must be complex'' axiom. This chapter does not show that the candidate
principles uniquely exclude every other operational model, establish an
infinite-dimensional or continuum limit, or supply new experimental
data.

The next physical connection must identify which UEOT mechanisms select
these states, effects, and composition rules; which real devices
implement the required tests; and which empirical predictions
distinguish the candidate theories. This chapter provides a complete,
checkable finite realization for that work. Completeness of the finite
realization is not final completion of a theory of nature.

\chapter{61. Complete purity and unitary equivalence of finite
purifications}\label{ch-61}

The assertion that a pure state is a vector state must follow from
extremality in the full density-state convex set, rather than being
assumed. This chapter proves that characterization and constructs actual
ancilla unitaries or isometries from equal Gram data, carefully checking
transposes, conjugation, and matrix direction.

Chapter 57 derived a Gram factor from positivity of an arbitrary finite
complex density state and proved purification existence using the actual
partial trace. Proving that certain vector states are pure, however,
does not yet compare all purifications defined by complete convex
extremality. Nor does constructing one Gram purification explain its
relation to other purifications of the same marginal. This chapter
supplies both obligations, treating equal and unequal auxiliary carriers
separately.

The conclusions concern the explicitly selected finite complex density
model. Its complete convex extreme points are exactly normalized vector
outer products. Any two purifications on the same auxiliary carrier are
connected by an actual auxiliary unitary. Purifications on different
carriers can first be embedded by actual isometric channels into a
common direct sum and then compared there. These last two statements
have different matrix types. They cannot be merged into a dimension-free
assertion about a ``rectangular unitary matrix.'' {[}QM318{]}

\section{61.1 Complete State Quantifiers and Normalization}

Continue to use

\[
\mathcal D(\iota)=\{\rho\in M_\iota(\mathbb C):\rho\succeq0,\ \operatorname{Tr}\rho=1\}.
\]

Here \(\iota\) is any finite carrier. The two conditions in
\texttt{Density} are positivity of the actual matrix and complex trace
one. \texttt{Is\allowbreak Pure} retains Chapter 57's extreme-point definition over
all density states and all \(0<t<1\). The allowed state set is not
reduced, and vector representation is not inserted as a defining field.

For a vector \(v:\iota\to\mathbb C\), write

\[
P_v=|v\rangle\langle v|,
\qquad (P_v)_{ij}=v_i\overline{v_j}.
\]

In this chapter's Lean declarations, normalization is written
\(\operatorname{Tr}P_v=1\). Expanding it gives \(\sum_i|v_i|^2=1\), the
unit-vector condition in finite complex Euclidean space. The vector norm
here is the two-norm; a different default norm on a function type must
not be substituted silently.

For a positive matrix \(X\), retain the mass notation

\[
m(X)=\operatorname{Re}\operatorname{Tr}X\ge0.
\]

The inherited zero-mass theorem for positive matrices gives
\(m(X)=0\iff X=0\). Only when \(m(X)>0\) do we use the valid normalized
density \(\widehat X=X/m(X)\). This condition is checked explicitly in
the convex decompositions below. A zero remainder is never assigned a
fictitious normalized state. {[}QM317; QM318{]}

\section{61.2 Every Positive-Mass Positive Component of an Extreme Point
Gives the Same State}

First prove an intermediate statement that directly reflects the
complete extremality quantifier. Suppose \(\rho\) is a pure density
state and

\[
\rho=X+Y,\qquad X,Y\succeq0,\qquad m(X)>0.
\]

Then \(\widehat X=\rho\). Unit trace gives \(m(X)+m(Y)=1\). The proof
separates two genuinely different cases.

If \(m(Y)=0\), positivity forces \(Y=0\), so \(m(X)=1\) and \(X=\rho\);
normalization leaves \(X\) unchanged. If \(m(Y)>0\), then \(0<m(X)<1\)
and

\[
\rho=m(X)\widehat X+m(Y)\widehat Y
=m(X)\widehat X+(1-m(X))\widehat Y.
\]

Both normalized states lie in the complete \(\mathcal D(\iota)\), so
they can be substituted directly into the original \texttt{Is\allowbreak Pure}
definition, giving \(\widehat X=\rho\) and \(\widehat Y=\rho\). No
additional rank, support, or diagonal-form assumption is imposed on
\(X\) or \(Y\). The Lean theorem
\texttt{Is\allowbreak Pure.\allowbreak normalize\_\allowbreak eq\_\allowbreak of\_\allowbreak positive\_\allowbreak summand} handles positivity,
strictly positive mass, normalized recombination, and the zero-remainder
branch explicitly. {[}QM318{]}

The result converts extremality into a constraint usable in a
Gram-column decomposition. It does not assume the desired rank-one
conclusion. Whether a component comes from a vector must still be
established by an actual positive-matrix decomposition in the next step.

\section{61.3 Deriving a Normalized Vector from Any Convex Extreme
Point}

\phantomsection\label{p-puc-01}{}

\textbf{P-PUC-01 (Purity characterization in the complete finite complex
density space).} For every \(\rho\in\mathcal D(\iota)\),

\[
\operatorname{Pure}(\rho)
\iff
\exists v:\iota\to\mathbb C,\quad
\operatorname{Tr}P_v=1\ \land\ \rho=P_v.
\]

\textbf{Proof.} The reverse implication---that a normalized vector state
is an extreme point of the complete set---was proved in P-PUR-01. Now
let \(\rho\) be any extreme point under the original definition.
Positivity supplies an actual matrix \(A\) with \(\rho=AA^\dagger\). For
each column \(a\), write \(u_a(i)=A_{ia}\). Entrywise multiplication
gives

\[
\rho=AA^\dagger=\sum_a |u_a\rangle\langle u_a|
=\sum_a P_a,\qquad P_a\succeq0.
\]

Taking traces gives \(\sum_a m(P_a)=1\), so at least one column has
\(m(P_a)>0\). If no column had positive mass, their finite sum would be
at most zero, contradicting unit trace. For a chosen such column, define
the actual remainder

\[
R=\sum_{b\ne a}P_b\succeq0.
\]

Deletion and recombination of the finite sum give \(\rho=P_a+R\). The
two-case result of the preceding section then yields
\(\rho=P_a/m(P_a)\). Set

\[
v=\frac{u_a}{\sqrt{m(P_a)}}.
\]

Since \(m(P_a)>0\), the square of its real square root is \(m(P_a)\).
Real scaling contributes one coefficient on each side of the outer
product. Thus

\[
P_v=\frac1{m(P_a)}P_a=\rho,
\qquad\operatorname{Tr}P_v=1.
\]

A normalized vector has therefore been derived from any extreme point of
the complete set, rather than supplied by the caller as a witness.
\(\square\)

In Lean, \texttt{gram\_\allowbreak eq\_\allowbreak sum\_\allowbreak outer} proves the column-outer-product
sum for the actual Gram matrix; \texttt{outer\_\allowbreak real\_\allowbreak smul} proves the
squared scaling law for real scalars; \texttt{Is\allowbreak Pure.\allowbreak exists\_\allowbreak vector}
completes the construction; and \texttt{is\allowbreak Pure\_\allowbreak iff\_\allowbreak exists\_\allowbreak vector}
reconnects it to the earlier implication. Choosing a positive-mass
column and a Gram factor uses the permitted classical foundations,
without a new rank-one axiom or physical spectral assumption.

An empty carrier is not silently excluded. If a unit-trace density
exists, the total-mass-one argument itself supplies a positive-mass
column. This is impossible on an empty carrier, so there is no valid
\(\rho\) to input in that case. No extra nonemptiness field is needed to
hide this boundary. {[}QM318{]}

\section{61.4 How Equal Gram Data Produce an Actual Unitary}

Next consider a purely linear-algebraic problem. Let
\(F,G:\mathbb C^{\iota}\to\mathbb C^{\alpha}\) be arbitrary finite
complex matrices satisfying

\[
F^\dagger F=G^\dagger G.
\]

For every input \(x\), the actual Gram quadratic forms give the squared
Euclidean norms:

\[
\|Fx\|_2^2=x^\dagger F^\dagger Fx
=x^\dagger G^\dagger Gx=\|Gx\|_2^2.
\]

Hence \(Fx=0\) implies \(Gx=0\). This makes

\[
L:\operatorname{ran}F\longrightarrow\mathbb C^\alpha,
\qquad L(Fx)=Gx
\]

a well-defined complex-linear map. Two representations of the same
\(Fx\) differ by an element of \(\ker F\), and therefore of \(\ker G\).
The norm equality further makes \(L\) a linear isometry on its actual
generated support.

The Lean construction does not leave well-definedness of \(L\) as an
unproved comment. It first defines an induced map through the quotient
by \(\ker F\), then uses the linear equivalence from the quotient by the
kernel to the range to return to \(\operatorname{ran}F\). It next proves
norm preservation for each range vector. \texttt{Linear\allowbreak Isometry.\allowbreak extend}
extends this subspace isometry to a linear isometry of the same full
finite auxiliary space. Injectivity of a finite-dimensional endomorphism
implies surjectivity, producing a linear isometric equivalence.

Taking its actual matrix \(U\) in the standard orthonormal basis gives

\[
U^\dagger U=I_\alpha,\qquad UU^\dagger=I_\alpha,
\qquad UF=G.
\]

Both unitary identities are conclusions. Neither \(F\) nor \(G\) is
assumed full rank or required to fill the auxiliary space.
Finite-dimensional isometric extension handles directions outside the
generated support. The only input to
\texttt{exists\_\allowbreak unitary\_\allowbreak of\_\allowbreak column\_\allowbreak gram} is the actual Gram equality:
its inputs contain no desired unitary matrix, extendibility field, or
unitary-equivalence conclusion. {[}QM318{]}

This is a mathematical existence construction. It uses the
quotient-space and finite isometric-extension theorems proved in the
pinned Mathlib and checks its foundations through exact named axiom
auditing. It does not additionally claim a certified
matrix-decomposition algorithm with numerical stability bounds.

\section{61.5 Row Grams, Transposes, and the Direction of Auxiliary
Action}

Use the fixed vectorization convention for purification vectors:

\[
\psi(i,a)=A_{ia},\qquad\phi(i,a)=B_{ia}.
\]

Their actual system marginals are \(AA^\dagger\) and \(BB^\dagger\).
Equal marginals therefore give a \textbf{row-Gram} equality:

\[
AA^\dagger=BB^\dagger.
\]

The source implementation transposes this equality and takes
\(F=A^\mathsf T\), \(G=B^\mathsf T\). Since

\[
(A^\mathsf T)^\dagger A^\mathsf T
=\overline A A^\mathsf T=(AA^\dagger)^\mathsf T,
\]

the preceding section directly supplies an auxiliary unitary satisfying
\(UA^\mathsf T=B^\mathsf T\). Entrywise, this is

\[
B_{ia}=\sum_bU_{ab}A_{ib},
\qquad
\phi(i,a)=\sum_bU_{ab}\psi(i,b).
\]

Thus the matrix acting on the second vector index is \(U\), while right
multiplication of the matrix factor is written \(B=AU^\mathsf T\). The
corresponding joint-vector relation is

\[
\phi=(I_\iota\otimes U)\psi.
\]

\texttt{exists\_\allowbreak auxiliary\_\allowbreak unitary\_\allowbreak of\_\allowbreak gram} derives the required \(U\)
from the actual transposed Gram equality.
\texttt{auxiliary\_\allowbreak action\_\allowbreak apply} computes every entry of the
tensor-matrix action, and \texttt{auxiliary\_\allowbreak action\_\allowbreak vectorize} proves
the transpose in the right-hand factor explicitly. The matrices \(U\),
\(U^\dagger\), \(U^\mathsf T\), and \(\overline U\) are therefore not
silently interchanged between representations.

An alternative mathematical route can begin with \(F=A^\dagger\), but
the conjugation must then be tracked when converting the
support-isometry matrix back into the vector's auxiliary action. The
implementation here uses the transpose route above. Alternative routes
discussed during preparation are not presented as proof executions
performed by this source. {[}QM318{]}

\section{61.6 Arbitrary Purifications on the Same Auxiliary Carrier}

\phantomsection\label{p-pun-01}{}

\textbf{P-PUN-01 (Unitary equivalence of all purifications on the same
auxiliary carrier).} Suppose
\(\Psi,\Phi\in\mathcal D(\iota\times\alpha)\) are pure in the complete
joint density space and their actual partial traces satisfy

\[
\operatorname{Tr}_\alpha\Psi=\operatorname{Tr}_\alpha\Phi.
\]

Then there is an actual matrix \(U\in M_\alpha(\mathbb C)\) such that

\[
U^\dagger U=UU^\dagger=I_\alpha,
\qquad
\Phi=(I_\iota\otimes U)\Psi(I_\iota\otimes U)^\dagger.
\]

\textbf{Proof.} Apply P-PUC-01 to \(\Psi\) and \(\Phi\) to obtain
normalized vectors \(\psi\) and \(\phi\) whose actual outer products are
the two joint densities. Substituting these outer products into the
partial-trace equality gives the row-Gram equality for their matrix
factors. The preceding section supplies \(U\) with
\(\phi=(I\otimes U)\psi\). Finally, the actual outer-product identity

\[
|K\psi\rangle\langle K\psi|
=K|\psi\rangle\langle\psi|K^\dagger
\]

gives the unitary conjugation of joint densities. \(\square\)

The intermediate vector theorem
\texttt{exists\_\allowbreak auxiliary\_\allowbreak unitary\_\allowbreak of\_\allowbreak equal\_\allowbreak marginal} does not even
require normalization. Any two joint vectors on the same carriers with
equal actual outer-product partial traces satisfy the exact
vector-action relation. The final theorem
\texttt{pure\_\allowbreak states\_\allowbreak unitarily\_\allowbreak equivalent} uses P-PUC-01 to handle all
pure densities defined by complete convex extremality. ``Arbitrary
purifications'' is therefore not reduced to ``purifications for which
the caller has already supplied vector representatives.'' {[}QM318{]}

``Uniqueness'' here means equivalence under auxiliary unitary action. It
asserts neither \(\Psi=\Phi\) nor uniqueness of the connecting matrix
\(U\). Unused auxiliary directions allow different unitary extensions,
and a vector representative's global phase can be absorbed into the
auxiliary unitary. Together, existence in Chapter 57 and equivalence
here provide the full chain within this finite complex model:
purifications exist, and all purifications on a common auxiliary carrier
can be compared.

\section{61.7 Different Carriers Require Actual Rectangular Isometric
Channels}

\phantomsection\label{p-aux-01}{}

\textbf{P-AUX-01 (Actual channel, marginal, and purity of an auxiliary
isometric embedding).} Let \(J:\mathbb C^\alpha\to\mathbb C^\beta\) be
an actual rectangular matrix satisfying \(J^\dagger J=I_\alpha\). Then

\[
\mathcal E_J(X)=(I_\iota\otimes J)X(I_\iota\otimes J)^\dagger
\]

is an actual single-Kraus CPTP channel from \(\iota\times\alpha\) to
\(\iota\times\beta\), and for \textbf{every complex input matrix} \(X\),

\[
\operatorname{Tr}_\beta\mathcal E_J(X)
=\operatorname{Tr}_\alpha X.
\]

If the input is pure in the complete joint state space, the output is
pure in the complete enlarged state space.

\textbf{Proof.} Put \(K=I_\iota\otimes J\). Tensor multiplication and
conjugate transposition give

\[
K^\dagger K=I_\iota\otimes(J^\dagger J)
=I_{\iota\times\alpha},
\]

so the actual Kraus matrix satisfies rectangular-channel completeness.
Chapter 58's rectangular-channel theory applies, including complete
positivity under every finite ancillary amplification.

Fix system indices \(i,j\), and set \(X^{ij}_{cd}=X_{(i,c),(j,d)}\), an
actual complex matrix on the auxiliary carrier. Direct calculation of
the corresponding output block gives

\[
(\mathcal E_JX)^{ij}=JX^{ij}J^\dagger.
\]

The actual partial trace takes the trace of that block. By cyclicity,

\[
\operatorname{Tr}(JX^{ij}J^\dagger)
=\operatorname{Tr}(J^\dagger JX^{ij})
=\operatorname{Tr}X^{ij}.
\]

Applying this at every \(i,j\) proves the marginal identity. No
positivity, unit-trace, purity, or product-form restriction is imposed
on \(X\) in this calculation.

Finally, if \(\Psi\) is pure in the complete convex set, P-PUC-01 gives
\(\Psi=P_\psi\). The channel output is \(P_{K\psi}\). Trace preservation
gives normalization, and P-PUR-01---or the reverse implication of
P-PUC-01---makes the output an extreme point with respect to all density
states in the enlarged space. \(\square\)

\texttt{auxiliary\allowbreak Embedding\allowbreak Channel} constructs the channel with actual
input and output types.
\texttt{auxiliary\allowbreak Embedding\allowbreak Channel\_\allowbreak operator\_\allowbreak apply} gives the block
formula; \texttt{auxiliary\allowbreak Embedding\allowbreak Channel\_\allowbreak discard\_\allowbreak operator}
quantifies over all complex matrices; and
\texttt{auxiliary\allowbreak Embedding\allowbreak Channel\_\allowbreak discard\_\allowbreak state} and
\texttt{auxiliary\allowbreak Embedding\allowbreak Channel\_\allowbreak preserves\_\allowbreak purity} reconnect the
result to actual densities and complete purity. The only matrix
assumption is \(J^\dagger J=I_\alpha\). When the output space is larger,
\(JJ^\dagger=I_\beta\) generally fails and must not be added to the
conclusion. {[}QM318{]}

\section{61.8 Comparing Different Auxiliary Dimensions After a Common
Direct-Sum Extension}

Take the common auxiliary carrier to be the disjoint sum
\(\alpha\sqcup\beta\), whose complex vector space is the direct sum of
the two auxiliary spaces. Define the actual rectangular matrices

\[
J_\alpha=\begin{bmatrix}I_\alpha\\0\end{bmatrix},
\qquad
J_\beta=\begin{bmatrix}0\\I_\beta\end{bmatrix}.
\]

They satisfy \(J_\alpha^\dagger J_\alpha=I_\alpha\) and
\(J_\beta^\dagger J_\beta=I_\beta\), respectively. The right-hand
identities belong to different source carriers. Construct

\[
\widetilde\Psi=\mathcal E_{J_\alpha}(\Psi),
\qquad
\widetilde\Phi=\mathcal E_{J_\beta}(\Phi).
\]

\phantomsection\label{p-pun-02}{}

\textbf{P-PUN-02 (Unitary equivalence after common extension for
different finite auxiliary carriers).} Suppose
\(\Psi\in\mathcal D(\iota\times\alpha)\) and
\(\Phi\in\mathcal D(\iota\times\beta)\) are pure in their respective
complete joint density spaces and have the same actual system marginal.
Then there is \(U\in M_{\alpha\sqcup\beta}(\mathbb C)\) satisfying both
unitary identities and

\[
\widetilde\Phi
=(I_\iota\otimes U)\widetilde\Psi(I_\iota\otimes U)^\dagger.
\]

\textbf{Proof.} Block-matrix multiplication directly establishes the
isometry identities of the canonical embeddings. P-AUX-01 makes
\(\widetilde\Psi\) and \(\widetilde\Phi\) pure densities in the same
complete common state space and preserves their original system
marginals. Those marginals are equal, so P-PUN-01 applies on the common
auxiliary carrier \(\alpha\sqcup\beta\). \(\square\)

The formal definitions \texttt{auxiliary\allowbreak Inl} and \texttt{auxiliary\allowbreak Inr}
give the actual rectangular embeddings; \texttt{extend\allowbreak Auxiliary\allowbreak Left} and
\texttt{extend\allowbreak Auxiliary\allowbreak Right} give their actual channel outputs. Four
purity and marginal theorems first establish the required premises
independently. Only then does
\texttt{pure\_\allowbreak states\_\allowbreak common\_\allowbreak extension\_\allowbreak unitary} invoke the
same-carrier unitary-equivalence theorem.

The result requires no comparison between \(|\alpha|\) and \(|\beta|\),
and does not treat the carriers as the same type. The common dimension
here is \(|\alpha|+|\beta|\); the result does not claim it is minimal or
infer a two-sided unitary directly between unequal-dimensional spaces.
Chapter 64 supplies new proofs of minimum auxiliary dimension and, when
\(|\alpha|\le|\beta|\), an actual isometry from \(\alpha\) to \(\beta\).
Strict enlargement still has no two-sided inverse, and a reverse
full-space isometry may still fail to exist. Those later results do not
change the quantifiers of the present proposition.
{[}QM318{]}{[}QM319{]}

\section{61.9 An Exact Positive Control: Equal Marginals, Different
Joint States}

The chosen auxiliary matrix has mixed real and imaginary entries, making
the directions of conjugation and transposition testable:

\[
U=\begin{bmatrix}3/5&4i/5\\4/5&-3i/5\end{bmatrix}.
\]

Its columns are \((3,4)^\mathsf T/5\) and \(i(4,-3)^\mathsf T/5\). Each
has squared norm one, and their inner product is zero. The actual
witness verifies both \(U^\dagger U=I\) and \(UU^\dagger=I\).

Take

\[
\psi_1=\tfrac35|00\rangle+\tfrac45|11\rangle,
\]

and compute using the fixed second-index action:

\[
\begin{aligned}
\psi_2=(I\otimes U)\psi_1
={}&\tfrac9{25}|00\rangle+\tfrac{12}{25}|01\rangle\\
&+\tfrac{16i}{25}|10\rangle-\tfrac{12i}{25}|11\rangle.
\end{aligned}
\]

The corresponding vectorization matrices are

\[
A=\begin{bmatrix}3/5&0\\0&4/5\end{bmatrix},
\qquad
B=\begin{bmatrix}9/25&12/25\\16i/25&-12i/25\end{bmatrix}.
\]

Exact multiplication gives

\[
AA^\dagger=BB^\dagger
=\begin{bmatrix}9/25&0\\0&16/25\end{bmatrix}=\rho.
\]

For example, the off-diagonal entry of \(BB^\dagger\) is
\(-144i/625+144i/625=0\). Both vectors have total squared norm one, so
their joint densities are pure in the complete convex sense and both
have actual partial trace \(\rho\).

The two joint vectors and densities are indeed different: their
\(|00\rangle\langle00|\) entries are \(9/25\) and \(81/625\). Equal
marginals thus do not imply equal joint densities; the correct
conclusion is that an auxiliary unitary connects them.
\texttt{exact\_\allowbreak auxiliary\_\allowbreak vector\_\allowbreak action} checks the exact vector
equality, \texttt{same\_\allowbreak actual\_\allowbreak marginal} checks the original actual
discard channel, and \texttt{general\_\allowbreak uniqueness\_\allowbreak instantiated} actually
invokes the general complete-pure-state uniqueness theorem. This avoids
checking only the small matrices calculated specifically for this
example. {[}QM318{]}

\section{61.10 Joint-Density Negative Controls for Three Incorrect
Action Directions}

\phantomsection\label{n-pun-01}{}

\textbf{N-PUN-01 (The adjoint, transpose, and entrywise conjugate cannot
replace the actual auxiliary action).} For the fixed \(U\) and
\(\psi_1\) above, replacing the correct auxiliary matrix by
\(U^\dagger\), \(U^\mathsf T\), or \(\overline U\) produces a joint
density different from \(P_{\psi_2}\) in each case.

\textbf{Proof.} Define \(\chi_T=(I\otimes T)\psi_1\). The following
three actual complex-matrix entries have each been calculated exactly:

\begin{longtable}[]{@{}
  >{\raggedright\arraybackslash}p{(\columnwidth - 6\tabcolsep) * \real{0.2100}}
  >{\raggedright\arraybackslash}p{(\columnwidth - 6\tabcolsep) * \real{0.2000}}
  >{\raggedright\arraybackslash}p{(\columnwidth - 6\tabcolsep) * \real{0.2800}}
  >{\raggedright\arraybackslash}p{(\columnwidth - 6\tabcolsep) * \real{0.3100}}@{}}
\toprule\noalign{}
\begin{minipage}[b]{\linewidth}\raggedright
Auxiliary matrix \(T\)
\end{minipage} & \begin{minipage}[b]{\linewidth}\raggedright
Compared joint-density entry
\end{minipage} & \begin{minipage}[b]{\linewidth}\raggedright
Value in \(P_{\chi_T}\)
\end{minipage} & \begin{minipage}[b]{\linewidth}\raggedright
Correct value in \(P_{\psi_2}\)
\end{minipage} \\
\midrule\noalign{}
\endhead
\bottomrule\noalign{}
\endlastfoot
\(U^\dagger\) & \(((0,0),(0,1))\) & \(108i/625\) & \(108/625\) \\
\(U^\mathsf T\) & \(((0,0),(0,1))\) & \(-108i/625\) & \(108/625\) \\
\(\overline U\) & \(((0,0),(1,0))\) & \(144i/625\) & \(-144i/625\) \\
\end{longtable}

The imaginary parts differ in every row, proving inequality of each
actual joint density. \(\square\)

\texttt{wrong\_\allowbreak adjoint\_\allowbreak joint\_\allowbreak density},
\texttt{wrong\_\allowbreak transpose\_\allowbreak joint\_\allowbreak density}, and
\texttt{wrong\_\allowbreak conjugate\_\allowbreak joint\_\allowbreak density} are three separately
axiom-audited negative-control theorems. Their entry values have
corresponding exact-calculation theorems. They compare joint densities,
so a wrong vector differing only by a global phase would not be mistaken
for a physical distinction merely because its expression looked
different. Actual density inequality in this example excludes that
ambiguity.

Testing only the system marginal cannot detect these direction errors.
The unitary identities of \(U\) and the transpose/conjugation relations
show that all three substitute matrices are also auxiliary unitaries, so
their action still preserves the system marginal. The negative controls
ask whether the specified joint purification is obtained. Their evidence
must therefore concern exact joint vectors or densities. They do not
claim that marginal tests identify a unitary choice on the discarded
auxiliary system. {[}QM318{]}

\section{61.11 An Actual Embedding from Two to Three Dimensions and a
Nonunitarity Control}

Take \(\alpha=\mathrm{Fin}\,2\) and \(\beta=\mathrm{Unit}\), so the
common carrier \(\alpha\sqcup\beta\) has dimension three. The actual
left canonical embedding is

\[
J=\begin{bmatrix}1&0\\0&1\\0&0\end{bmatrix}.
\]

\phantomsection\label{n-aux-01}{}

\textbf{N-AUX-01 (Strict auxiliary enlargement is an isometric
embedding, not a two-sided unitary equivalence).} This actual
\(3\times2\) matrix satisfies \(J^\dagger J=I_2\), but
\(JJ^\dagger\ne I_3\).

\textbf{Proof.} Column orthogonality gives the first identity directly.
The second product is \(\operatorname{diag}(1,1,0)\): its third diagonal
entry is zero, whereas the corresponding entry of \(I_3\) is one.
\(\square\)

The witnesses \texttt{strict\_\allowbreak embedding\_\allowbreak isometry} and
\texttt{strict\_\allowbreak embedding\_\allowbreak not\_\allowbreak coisometry} prove the two claims
separately, and \texttt{auxiliary\_\allowbreak dimensions\_\allowbreak differ} proves that the
carrier cardinalities really differ. Embedding the second joint pure
state of §61.9 along this \(J\) gives \texttt{larger\allowbreak State}. It remains
pure in the complete six-dimensional joint space, has a
three-dimensional auxiliary carrier, and preserves the original system
marginal. Together with \texttt{first\allowbreak State} on the original
two-dimensional auxiliary carrier, it forms a pair of purifications with
genuinely different auxiliary dimensions.

\texttt{different\_\allowbreak auxiliary\_\allowbreak common\_\allowbreak unitary} applies P-PUN-02 to these
actual densities. Their common auxiliary carrier is
\(\mathrm{Fin}\,2\sqcup(\mathrm{Fin}\,2\sqcup\mathrm{Unit})\), of
dimension five. The general theorem supplies a matrix on that carrier
satisfying both unitary identities. This five-dimensional matrix is
obtained by the proved existence construction. The witness does not
additionally display a numerically solved five-dimensional unitary, and
it does not conflate the two original auxiliary dimensions. {[}QM318{]}

\section{61.12 Source, Declaration-Level Auditing, and Evidence Scope}

The new public sources for this chapter are
\texttt{UEOTQM/\allowbreak Pure\allowbreak State\allowbreak Characterization.\allowbreak lean} and
\texttt{UEOTQM/\allowbreak Purification\allowbreak Uniqueness.\allowbreak lean}; the concrete witness is
\texttt{scripts/\allowbreak witnesses/\allowbreak Purification\allowbreak Uniqueness\allowbreak Examples.\allowbreak lean}. Their
counts are:

\begin{longtable}[]{@{}
  >{\raggedright\arraybackslash}p{(\columnwidth - 8\tabcolsep) * \real{0.4000}}
  >{\raggedleft\arraybackslash}p{(\columnwidth - 8\tabcolsep) * \real{0.1800}}
  >{\raggedleft\arraybackslash}p{(\columnwidth - 8\tabcolsep) * \real{0.1000}}
  >{\raggedleft\arraybackslash}p{(\columnwidth - 8\tabcolsep) * \real{0.1000}}
  >{\raggedleft\arraybackslash}p{(\columnwidth - 8\tabcolsep) * \real{0.2200}}@{}}
\toprule\noalign{}
\begin{minipage}[b]{\linewidth}\raggedright
Source file
\end{minipage} & \begin{minipage}[b]{\linewidth}\raggedleft
Explicit \texttt{theorem}
\end{minipage} & \begin{minipage}[b]{\linewidth}\raggedleft
\texttt{def}
\end{minipage} & \begin{minipage}[b]{\linewidth}\raggedleft
\texttt{abbrev}
\end{minipage} & \begin{minipage}[b]{\linewidth}\raggedleft
Exact named axiom audits
\end{minipage} \\
\midrule\noalign{}
\endhead
\bottomrule\noalign{}
\endlastfoot
\texttt{Pure\allowbreak State\allowbreak Characterization.\allowbreak lean} & 5 & 0 & 0 & 5 \\
\texttt{Purification\allowbreak Uniqueness.\allowbreak lean} & 22 & 5 & 0 & 27 \\
\texttt{Purification\allowbreak Uniqueness\allowbreak Examples.\allowbreak lean} & 28 & 10 & 0 & 38 \\
\end{longtable}

The two public modules contain 27 explicit theorems and 5 definitions in
total; the witness adds 28 theorems and 10 definitions. All 70 distinct
stable named public and witness declarations were checked in the actual
focused audit for this workstream and depend only on \texttt{propext},
\texttt{Classical.\allowbreak choice}, and \texttt{Quot.\allowbreak sound}. Every witness
\texttt{theorem}, \texttt{def}, and \texttt{abbrev} has its axiom
dependencies printed in source order; the audit does not merely sample
the final combined theorem.

Actual acceptance logs, final source hashes, source-to-olean bindings,
and the named inventory are recorded in
\texttt{physics/\allowbreak qm/\allowbreak kernel/\allowbreak evidence/\allowbreak v31.\allowbreak 8/\allowbreak B/\allowbreak focused-\allowbreak verification.\allowbreak json}
and \texttt{declaration\_\allowbreak inventory.\allowbreak json}. Focused module validation
corresponds to \texttt{pure-\allowbreak 005} and \texttt{uniqueness-\allowbreak 005}; the final
concrete witness to \texttt{witness-\allowbreak 003}; and the exact public audit to
\texttt{public-\allowbreak audit-\allowbreak 001}. These are actual invocations with exit status
zero. Earlier informative failures and their available source snapshots
are preserved separately and are not counted as accepted results.

This workstream receipt establishes focused proof acceptance. The
version-wide build, full axiom audit, inherited witnesses, and remote CI
have their own completion and identity records in the main version
receipts. Local results cannot substitute for a remote run not actually
read back. Ordinary linter notices must also be distinguished from
logical acceptance; permitted ordinary warnings must not be rewritten as
``zero warnings.'' No new real experiment is supplied here. Rational
complex-matrix positive and negative controls are formal witnesses, not
empirical confirmation by nature. {[}QM318{]}

\section{61.13 Closed Mathematical Obligations and Open Theory
Selection}

Chapter 57's existence construction and this chapter's converse purity,
same-auxiliary-carrier unitary equivalence, and different-carrier common
extension now connect within one finite complex density model.
Pure-state definition, Gram factor, joint vector, actual partial trace,
and actual rectangular channel each have an explicit interface. The
conclusions no longer rely on assuming all pure states are vector states
or requiring the caller to provide the desired auxiliary unitary.

Uniqueness of the auxiliary unitary itself, minimum dimension of a
common extension, direct isometric extension under additional
dimensional conditions, and certified numerical construction are
distinct questions. Chapter 64 subsequently addresses minimum
purification dimension and direct isometry under the required dimension
ordering; those are separate new proofs. Closed ranges, nonfinite rank,
and extendibility in infinite dimensions also cannot simply use this
chapter's finite-dimensional surjectivity argument. If a reconstruction
principle imposes more specific requirements on auxiliary completion,
allowed reversible processes, or system isomorphisms, its fields must be
connected individually to the actual proved types. The name ``essential
uniqueness'' alone does not establish every literature contract.

These results principally concern mathematics under standard finite
complex quantum structure P and actual operational connections B. They
strengthen the candidate model's internal completeness and test
interfaces, but do not establish a bridge uniquely selecting complex
quantum structure from the frozen 106 Core items and four generators.
The complete classical model's purification obstruction and the real
model's local-tomography counterexample continue to constrain selection
principles. Excluding these two examined models does not exclude every
possible theory of nature.

The generation-count conclusions continue to depend on their spectral,
representation, threshold, and effective-window assumptions. This
chapter neither changes those assumptions nor supplies new empirical
data. UEOT's aim remains to connect theoretical ideas, actual
derivations, counterexamples, and empirical obligations into a
reviewable whole. This chapter completes one finite quantum-mathematical
segment of that task. Further claims about nature still require
independent premises, actual proofs, and testable physical
correspondences. {[}QM318{]}

\chapter{62. Per-action record-error budgets with retained
history}\label{ch-62}

Resolving protocol error into individual actions requires causal
histories and quantum memory. Under a common fixed-depth controller,
actual PSD history blocks and weights from the realization chosen for
the common-input comparison yield actionwise budgets; a label-deletion
counterexample explains why the common-control contract matters.

Chapter 58 constructed a joint channel for actual instruments, all
finite adaptive leaf records, and quantum memory, and bounded
complete-record TV by its whole-channel defect. That defect defines how
two protocols differ, but usually still needs to be calculated or
estimated. Under a fixed finite depth and a common causal controller,
this chapter traces the error to the instrument actually selected at
each step. It first constructs a history-reading rectangular Kraus
channel, then derives action budgets from a common input's PSD principal
blocks, actual block weights, and valid block effects.

The controller actually reads and dephases its history register. The
formal channel can act on any history--memory density matrix, and the
proof uses all input principal blocks. This does not turn
classical-history control into coherent quantum control. The original
finite-tree interpreter retains its more general variable-length leaf
semantics; the present construction is connected to it through an
explicit fixed-depth subclass. {[}QM318{]}

\section{62.1 Finite Time, Actual Histories, and a Common Controller}

Fix a finite outcome carrier \(O\), a finite complex quantum-memory
carrier \(\iota\), and a finite depth \(N\in\mathbb N\). The complete
outcome history before step \(t\) is

\[
\mathsf H_t=\{h:\operatorname{Fin}(t)\to O\}.
\]

The subscript counts recorded outcomes; \(\mathsf H_0\) contains only
the empty history. Let \(A\) be the action carrier. Finite depth
requires the controller to choose actions at the relevant times and
histories, but does not require \(A\) itself to be finite. Its actual
type is

\[
C:\prod_{t:\operatorname{Fin}(N)}(\mathsf H_t\to A).
\]

The two implementations use \(I_a\) and \(J_a\), respectively, but both
select the same \(C_t(h)\) at that history. A common outcome alphabet
alone does not imply a common controller; here the same function value
is used. Each instrument has a complete Kraus family, for example

\[
\Phi^I_{a,o}(X)=\sum_{k:\kappa}K^I_{a,o,k}X(K^I_{a,o,k})^\dagger,
\qquad
\sum_{o,k}(K^I_{a,o,k})^\dagger K^I_{a,o,k}=I_\iota.
\]

The Kraus indices for \(J\) may have a different finite carrier
\(\eta\). Actual channels are compared; identical Kraus representations
are not required. The outcome-retaining joint channel for a single
action remains the one from Chapter 58:

\[
\mathcal R^I_a(X)=\sum_o|o\rangle\langle o|\otimes\Phi^I_{a,o}(X).
\]

Write the full-effect distance as \(d_{\rm eff}\) and define the actual
action-record defect

\[
\delta_{\rm rec}(I_a,J_a)
=\sup\left(\{0\}\cup
 \{d_{\rm eff}(\mathcal R^I_a(\rho),\mathcal R^J_a(\rho)):
   \rho\in\mathcal D(\iota)\}\right).
\]

The common-input quantifier in the supremum is essential, and record
labels must not be discarded before the defect is defined. Inserting
zero makes the set explicitly nonempty even when the input-state type is
empty. When a uniform upper bound is imposed on a history-channel
defect, the source explicitly retains nonnegativity of that upper bound.

\section{62.2 PSD History Blocks, Actual Weights, and Valid Effects}

For any finite history carrier \(H\) and any
\(X\in M_{H\times\iota}(\mathbb C)\), define the principal block

\[
X_h(i,j)=X((h,i),(h,j)).
\]

This is not a selection of arbitrary entries subsequently declared to be
a state. Writing the embedding \(V_h:e_i\mapsto e_{(h,i)}\) gives
\(X_h=V_h^\dagger XV_h\). The same principal-submatrix construction is
used for output effects.

\phantomsection\label{p-arb-01}{}

\textbf{P-ARB-01 (History blocks, normalized weights, and valid block
effects).} If \(\zeta\in\mathcal D(H\times\iota)\), every principal
block \(\zeta_h\) is PSD, and

\[
p_h=\operatorname{Tr}\zeta_h\ge0,
\qquad \sum_h p_h=1,
\qquad p_h=0\ \Longleftrightarrow\ \zeta_h=0.
\]

If \(E\in\mathcal E(H\times\nu)\) is any valid output effect, its
principal block \(E_h\) is a valid member of \(\mathcal E(\nu)\). For
any block family \(B_h\), the actual trace-pairing identity is

\[
\left\langle E,\operatorname{diag}_hB_h\right\rangle_{\mathbb R}
=\sum_h\langle E_h,B_h\rangle_{\mathbb R}.
\]

\textbf{Proof.} Principal submatrices preserve PSD, so
\(\zeta_h\succeq0\). Summing the full trace over product indices gives
\(\sum_h\operatorname{Tr}\zeta_h=\operatorname{Tr}\zeta=1\). A PSD
matrix has zero trace exactly when it is zero. This existing theorem
applies to every block without forcing a posterior state onto a
zero-weight history.

For an effect, restrict both \(E\succeq0\) and \(I-E\succeq0\) to
history \(h\). Since \(i\mapsto(h,i)\) is injective, the principal block
of the identity is \(I_\nu\), giving \(E_h\succeq0\) and
\(I_\nu-E_h\succeq0\). Finally, expand matrix multiplication and the
trace. The block output has zero off-diagonal history entries, leaving
precisely the sum of principal-block pairings. \(\square\)

The input \(\zeta\) may have nonzero coherences between histories.
Principal-block positivity and trace normalization do not require it to
be block diagonal in advance. The channel constructed next actually
removes those coherences; this explicit action determines the proof's
scope.

\section{62.3 From Actual Reading Kraus Operators to the Weighted Action
Bound}

First consider general history-dependent rectangular channels
\(\mathcal C_h:M_\iota\to M_\nu\). If their Kraus matrices are
\(L_{h,k}\), define on joint history--memory spaces

\[
\widetilde L_{h,k}((g,a),(g',i))
=\mathbf 1_{g=h}\mathbf 1_{g'=h}\,L_{h,k}(a,i).
\]

This matrix reads input history \(h\), applies \(\mathcal C_h\) to
memory, and writes the same \(h\) to the output. Summing Kraus terms
over different \(h\) is an incoherent channel sum, implementing actual
history dephasing.

\phantomsection\label{p-arb-02}{}

\textbf{P-ARB-02 (Action-weighted bound for the actual history-reading
channel).} This Kraus family forms a trace-preserving completely
positive channel \(\mathcal S_{\mathcal C}\) satisfying, for every
actual matrix input,

\[
\mathcal S_{\mathcal C}(X)
=\operatorname{diag}_h\mathcal C_h(X_h).
\]

If \(\mathcal D_h\) is another channel family with the same input and
output carriers, every \(\zeta\in\mathcal D(H\times\iota)\) satisfies

\[
d_{\rm eff}(\mathcal S_{\mathcal C}(\zeta),
             \mathcal S_{\mathcal D}(\zeta))
\le \sum_h p_h\,\delta(\mathcal C_h,\mathcal D_h),
\qquad p_h=\operatorname{Tr}\zeta_h.
\]

In particular, if the actual defect at every \(h\) is at most
\(\varepsilon\), this common-input distance is at most \(\varepsilon\).
If \(\varepsilon\ge0\) is also explicit, the actual history-channel
defect is at most \(\varepsilon\), including the empty-input-carrier
boundary of the supremum.

\textbf{Proof.} For fixed input indices \((g,i),(g',j)\), the matrix
product \(\widetilde L_{h,k}^\dagger\widetilde L_{h,k}\) is nonzero only
if \(g=g'=h\). First sum over \(k\) using completeness of each
\(\mathcal C_h\), then sum over \(h\), obtaining the joint identity
matrix. Complete positivity follows from the proved finite-ancilla
amplification theorem for rectangular channels. Similarly expanding
\(\widetilde L_{h,k}X\widetilde L_{h,k}^\dagger\) shows that only
\(X_h\) enters the corresponding output block, giving the action
formula.

For each PSD block \(X_h\), distinguish two cases. If \(p_h=0\),
P-ARB-01 gives \(X_h=0\), and both linear channel outputs vanish. If
\(p_h>0\), set \(\rho_h=X_h/p_h\), with actual PSD and unit-trace
proofs. For every valid effect \(F\), real linearity and the defect
definition give

\[
\begin{aligned}
&\left|\langle F,\mathcal C_h(X_h)\rangle_{\mathbb R}
 -\langle F,\mathcal D_h(X_h)\rangle_{\mathbb R}\right|\\
&\quad=p_h\left|p_{\mathcal C_h(\rho_h)}(F)
                 -p_{\mathcal D_h(\rho_h)}(F)\right|\\
&\quad\le p_h\delta(\mathcal C_h,\mathcal D_h).
\end{aligned}
\]

Now take any joint output effect \(E\). P-ARB-01 makes every \(E_h\)
valid and expands the total probability difference into the sum of
principal-block differences. The triangle inequality for absolute values
gives the weighted bound, followed by the supremum over all effects. The
uniform bound uses \(\sum_hp_h=1\). For the history-channel defect, take
the supremum over all input states and also check the inserted zero.
\(\square\)

The crucial weights come from the same \(\zeta\). The argument does not
assume that two different evolution histories have identical
distributions. It supplies a legitimate common input for the triangle
decomposition below.

\section{62.4 Chronological Outcome Appending and the Original Protocol}

For a length-\(t\) history \(h\) and a new outcome \(o\), write
\(h\mathbin{\|}o\) for the history obtained by appending \(o\), of
length \(t+1\). The actual carrier bijection is

\[
\mathsf H_t\times(O\times\iota)
\longleftrightarrow\mathsf H_{t+1}\times\iota,
\qquad(h,(o,i))\longmapsto(h\mathbin{\|}o,i).
\]

Its inverse reads the first \(t\) entries and the last entry. An actual
rectangular permutation Kraus matrix implements the bijection, and the
inverse bijection also gives a channel. Relabelling consequently
preserves complete states and full-effect distance, not merely trace.

At time \(t\), the actual step \(S^I_t\) first reads the history and
uses \(\mathcal C_h=\mathcal R^I_{C_t(h)}\), then appends the outcome by
that permutation. Thus

\[
S^I_t(X)_{h\|o,h\|o}
=\Phi^I_{C_t(h),o}(X_h),
\]

and output blocks between distinct complete histories vanish. An actual
permutation embeds the initial state through
\(\iota\simeq\mathsf H_0\times\iota\). Successive steps produce the
joint density \(\zeta^I_{C,\rho}\). At depth zero no action is called,
so nonemptiness of the action carrier is unnecessary.

\phantomsection\label{p-arb-03}{}

\textbf{P-ARB-03 (Matrix agreement between finite-history evolution and
the original complete protocol).} For any fixed \(N\), controller \(C\),
actual instrument family \(I\), and initial state \(\rho\), the
constructed actual output satisfies

\[
\zeta^I_{C,\rho}
=\operatorname{diag}_{h\in\mathsf H_N}
 \operatorname{recordOperator}(I,\operatorname{word}_C(h),\rho).
\]

From the same \(C\), construct the protocol \(\pi_C\) in the original
type. There is an actual leaf bijection
\(e_C:\operatorname{Leaf}(\pi_C)\simeq\mathsf H_N\) with

\[
\operatorname{word}_{\pi_C}(\ell)
=\operatorname{word}_C(e_C\ell).
\]

After the actual permutation \(e_C\times\operatorname{id}_\iota\), the
two joint density states coincide. The associated leaf probabilities,
complete-record TV, and terminal quantum marginals all agree. In
particular, the quantum marginal is

\[
\operatorname{Tr}_{\mathsf H_N}\zeta^I_{C,\rho}
=\sum_{\ell:\operatorname{Leaf}(\pi_C)}
 \operatorname{recordOperator}(I,\operatorname{word}_{\pi_C}(\ell),\rho).
\]

\textbf{Proof.} The original protocol recurses from the first letter of
a word: stopping corresponds to the empty history, and an action node
supplies a successor controller for every \(o\in O\). The
leaf-to-history map places the current outcome first and recursively
handles the tail. Its inverse reads the first entry and recursively
recovers the tail leaf. Both composites are identities, establishing an
actual bijection.

Channel execution instead appends at the end, so agreement of these
recursion orders must also be proved. The source first establishes

\[
\operatorname{word}_{C}(h\|o)
=\operatorname{word}_{C_{<t}}(h)
 \mathbin{+\!+}[(C_t(h),o)],
\]

where \(C_{<t}\) is the same controller with its last action time
removed. For the original interpreter it separately proves

\[
\operatorname{recordOperator}(I,w\mathbin{+\!+}v,X)
=\operatorname{recordOperator}
  (I,v,\operatorname{recordOperator}(I,w,X)).
\]

At depth zero, the permutation embedding preserves the original matrix.
In the inductive step, combine the principal-block formula for the
actual \(S^I_t\), the induction hypothesis, and the two word-order
identities. This identifies the full output matrix, not just
probabilities. Reindexing the finite sum by the leaf bijection then
gives the original protocol's complete-state permutation identity and
quantum marginal. \(\square\)

The original \texttt{Protocol} permits leaves of different lengths,
whereas \(\pi_C\) is the complete subclass constructed from one fixed
\(N\). Even zero-probability outcomes retain valid successor actions.
This chapter does not prove padding semantics for an arbitrary stopping
tree, and a leaf-type conversion alone cannot supply that broader
result.

\section{62.5 Per-Action Recursion, the Finite Sum, and Complete-Record
TV}

For the same \(C\), let the two initial states be \(\rho,\sigma\). Write
\(\zeta^I_t\) for the actual joint state reached from \(\rho\) through
\(I\) and the first \(t\) controller steps, and \(\zeta^J_t\) for the
state reached from \(\sigma\) through \(J\). Define

\[
\begin{aligned}
q^J_t(h)&=\operatorname{Tr}(\zeta^J_t)_h,\\
a_t&=\sum_{h\in\mathsf H_t}q^J_t(h)
 \delta_{\rm rec}(I_{C_t(h)},J_{C_t(h)}).
\end{aligned}
\]

This is \texttt{action\allowbreak Cost} in the source. Its prefix state is actually
constructed as \texttt{run\ J\ (C.\allowbreak truncate\ t)\ σ}. The value \(a_t\) is
nonnegative, and a zero-mass history contributes zero by the zero-block
theorem. It is not the whole-protocol defect under a different name.

\phantomsection\label{p-arb-04}{}

\textbf{P-ARB-04 (Actual per-action sum budget and complete-record TV).}
Under the stated fixed finite depth, common controller, and common
record and memory carriers, any initial states \(\rho,\sigma\) satisfy

\[
\begin{aligned}
d_{\rm eff}(\zeta^I_N,\zeta^J_N)
&\le d_{\rm eff}(\rho,\sigma)+\sum_{t=0}^{N-1}a_t,\\
\operatorname{TV}(p^I_N,p^J_N)
&\le d_{\rm eff}(\rho,\sigma)+\sum_{t=0}^{N-1}a_t.
\end{aligned}
\]

If, for every time \(t\) and every length-\(t\) history \(h\), one has
actually proved

\[
\delta_{\rm rec}(I_{C_t(h)},J_{C_t(h)})\le\varepsilon_t,
\]

both right-hand sides may be replaced by
\(d_{\rm eff}(\rho,\sigma)+\sum_{t<N}\varepsilon_t\). Here
\(p^I_N,p^J_N\) are complete-history distributions. After the leaf
bijection of P-ARB-03, the same bounds apply to the original protocol's
complete leaf laws.

\textbf{Proof.} Put \(D_t=d_{\rm eff}(\zeta^I_t,\zeta^J_t)\). Apply the
triangle inequality at the common intermediate state
\(S^I_t(\zeta^J_t)\):

\[
\begin{aligned}
D_{t+1}
&\le d_{\rm eff}(S^I_t(\zeta^I_t),S^I_t(\zeta^J_t))\\
&\quad+d_{\rm eff}(S^I_t(\zeta^J_t),S^J_t(\zeta^J_t))\\
&\le D_t+a_t.
\end{aligned}
\]

Actual CPTP effect pullback contracts the first term. In the second,
both instrument implementations are compared on the same \(\zeta^J_t\).
P-ARB-02 gives the action bound with common weights \(q^J_t(h)\), and
the outcome-appending permutation preserves distance. Induction and
finite summation give \(D_N\le D_0+\sum_ta_t\). The empty-history
initial permutation yields \(D_0=d_{\rm eff}(\rho,\sigma)\). If each
action defect is bounded by \(\varepsilon_t\), weight normalization in
P-ARB-01 gives \(a_t\le\varepsilon_t\), and summing proves the stated
alternative.

For any finite record event \(S\subseteq\mathsf H_N\), construct the
actual effect

\[
E_S=\sum_{h\in S}|h\rangle\langle h|\otimes I_\iota.
\]

Both it and its complement are PSD. Its probability on any joint state
is exactly the sum of principal-block masses over \(S\). Thus
\(|p^I_N(S)-p^J_N(S)|\le D_N\). The all-event characterization of finite
normalized distributions gives
\(\operatorname{TV}(p^I_N,p^J_N)\le D_N\), retaining coefficient one
instead of multiplying an individual-outcome bound by the number of
histories. Finally, use the proved TV relabelling identity in P-ARB-03
to return to the original protocol. \(\square\)

Here \(J\) specifies the chosen direction of the triangle decomposition
and its common input. The formula requires neither
\(\zeta^I_t=\zeta^J_t\) nor equal history weights for the two
implementations. Actual \(a_t\) can reflect histories rarely visited by
the \(J\) prefix. A uniform \(\varepsilon_t\) only requires checking
every underlying action selected by the controller, giving a more
conservative bound that may be easier to estimate externally. It still
needs a genuine instrument-error proof, calibration, or other
independent error input. Defining a number as a supremum does not
automatically compute it numerically.

\section{62.6 An Exact Rational Two-Step Positive Control}

Fix \(O=\iota=\operatorname{Fin}(2)\), with four Kraus indices per
instrument. Let \(\mathsf D_b\) deterministically output
\(b\in\{0,1\}\). Its sole nonzero Kraus matrix is the identity at index
zero of that label, so every matrix input satisfies

\[
\Phi_{\mathsf D_b,o}(X)=\mathbf1_{o=b}X.
\]

Define a biased instrument \(\mathsf B\) with Kraus quadruple
\((I/2,I/2,I/2,0)\) at label zero and \((0,0,0,I/2)\) at label one. The
four nonzero Gram terms sum to \(I\), making it an actual complete
instrument with

\[
\Phi_{\mathsf B,0}(X)=\tfrac34X,
\qquad \Phi_{\mathsf B,1}(X)=\tfrac14X.
\]

All matrix coefficients are exact rationals. Neither floating-point
approximation nor an estimated square root is used to set the Kraus
matrices.

The controller chooses action \(\mathrm{none}\) first. After reading
outcome \(o\), it chooses \(\mathrm{some}(o)\) on the second step. The
reference implementation takes \(I_{\mathrm{none}}=\mathsf D_0\); the
approximate implementation takes \(J_{\mathrm{none}}=\mathsf B\); both
take \(I_{\mathrm{some}(o)}=J_{\mathrm{some}(o)}=\mathsf D_o\). The
controller thus actually reads the previous record and changes the
subsequent instrument accordingly.

\phantomsection\label{p-arb-05}{}

\textbf{P-ARB-05 (An exact rational witness saturating the stepwise
budget).} For every two-dimensional input density state, the first
action's actual record defect is \(1/4\), and the second action's
corresponding defect is zero. The actual weighted costs are \(a_0=1/4\)
and \(a_1=0\), and complete two-step record TV is exactly \(1/4\). For
different initial states,

\[
d_{\rm eff}(\zeta^I_{2,\rho},\zeta^J_{2,\sigma})
\le d_{\rm eff}(\rho,\sigma)+\tfrac14.
\]

For a common input state \(\rho\), both terminal quantum marginals are
exactly \(\rho\). Complete-record TV in the original protocol's leaf
representation is also exactly \(1/4\).

\textbf{Proof.} For any valid joint effect \(E\), set \(x=p_\rho(E_0)\)
and \(y=p_\rho(E_1)\). P-ARB-01 gives \(0\le x,y\le1\). The
probabilities of the two first-step outputs are

\[
x,\qquad \tfrac34x+\tfrac14y.
\]

Their absolute difference is \(|x-y|/4\le1/4\), an upper bound over all
input states and all valid joint effects. The record event ``output
zero'' has probabilities \(1\) and \(3/4\), supplying a matching lower
bound. Taking the input-state supremum gives the actual action defect
\(1/4\), without mistaking a single tested input for a supremum proof.

On the second step, at the common history \(h\), both parties use the
same instrument \(\mathsf D_{h(0)}\), so the action defect is zero.
Normalized prefix weights give \(a_0=1/4,a_1=0\), and P-ARB-04 gives the
full-process upper bound. The complete history \((0,0)\) has actual
probabilities \(1\) and \(3/4\), whose event difference supplies TV
lower bound \(1/4\), equal to the upper bound.

For the terminal quantum marginal, sum the actual \texttt{run\_\allowbreak matrix}
over all four length-two histories. The reference sum is \(\rho\), and
the approximate sum is \(3\rho/4+\rho/4=\rho\). The leaf-TV identity
then gives the original-protocol conclusion. \(\square\)

The example checks the actual order: the current outcome enters the
history, the controller reads that existing history, and the next action
depends on it. It also checks weights, summation, and the TV
coefficient. This exact calculation is a finite formal witness, not
empirical confirmation. {[}QM318{]}

\section{62.7 A Label-Deletion Counterexample Under the Same Feedback
Rule}

Keep the controller and reference implementation \(I\) above. Instead
take \(\widehat J_{\mathrm{none}}=\mathsf D_1\) and
\(\widehat J_{\mathrm{some}(o)}=\mathsf D_o\). For every action, the two
implementations have the same actual quantum channel after labels are
discarded: the identity channel.

\phantomsection\label{n-arb-01}{}

\textbf{N-ARB-01 (Zero unlabelled action defects can coexist with
maximum complete-record error).} Every unlabelled action defect of these
actual instrument families is zero, but their two-step complete-record
TV is \(1\) for any common initial state. The event that the second
outcome is zero has probabilities \(1\) and \(0\). Both terminal quantum
density states equal the input \(\rho\), and hence each other.

\textbf{Proof.} The unlabelled channels of \(\mathsf D_0\) and
\(\mathsf D_1\) both send every input matrix \(X\) to \(X\); later
actions even use the same instrument on both sides. Their unlabelled
channels are equal as functions, and the existing
\texttt{channel\allowbreak Defect\_\allowbreak self} gives zero defect.

The reference implementation deterministically records \((0,0)\); the
other records \((1,1)\). These probabilities follow from the original
chronological words and actual branch formulas. The history event
\(\{(0,0)\}\) has probability difference \(1\), giving the same lower
bound on complete TV. TV between finite probability distributions is at
most \(1\), so it is exactly \(1\). Expanding the actual event
\(\{h:h(1)=0\}\) over the four histories gives second-outcome
probabilities \(1\) and \(0\). Summing the actual record blocks on
either side gives terminal quantum marginal \(\rho\). The actual leaf
bijection preserves original-protocol leaf TV. \(\square\)

The counterexample connects Chapter 58's label obligation to the new
stepwise realization and the old protocol; it does not reinterpret the
result as unequal terminal quantum states. Replacing
\(\delta_{\rm rec}\) in P-ARB-04 by the zero unlabelled action defects
would incorrectly predict zero record error for a common initial state.
The actual labelled action defects do not permit that replacement.

\section{62.8 Public Interfaces, Actual Evidence, and Physical
Obligations}

The new public modules are \texttt{UEOTQM.\allowbreak History\allowbreak Channels} and
\texttt{UEOTQM.\allowbreak Action\allowbreak Record\allowbreak Budget}. The table lists stable entry points
for their responsibilities; concrete positive and negative controls are
in \texttt{UEOTQM.\allowbreak Action\allowbreak Record\allowbreak Budget\allowbreak Examples}. {[}QM318{]}

\begin{longtable}[]{@{}
  >{\raggedright\arraybackslash}p{(\columnwidth - 2\tabcolsep) * \real{0.3000}}
  >{\raggedright\arraybackslash}p{(\columnwidth - 2\tabcolsep) * \real{0.7000}}@{}}
\toprule\noalign{}
\begin{minipage}[b]{\linewidth}\raggedright
Numbered result or interface
\end{minipage} & \begin{minipage}[b]{\linewidth}\raggedright
Actual source responsibility
\end{minipage} \\
\midrule\noalign{}
\endhead
\bottomrule\noalign{}
\endlastfoot
P-ARB-01: PSD, mass, and zero blocks & \texttt{principal\_\allowbreak pos\allowbreak Semidef},
\texttt{sum\_\allowbreak mass\_\allowbreak principal}, \texttt{weight\_\allowbreak nonneg},
\texttt{sum\_\allowbreak weight}, \texttt{zero\_\allowbreak weight\_\allowbreak iff\_\allowbreak zero\_\allowbreak principal} \\
P-ARB-01: valid block effects and actual pairing &
\texttt{restrict\allowbreak Effect}, \texttt{restrict\allowbreak Effect\_\allowbreak matrix},
\texttt{pairing\_\allowbreak block} \\
P-ARB-02: reading Kraus matrices and actual block action &
\texttt{read\allowbreak Kraus}, \texttt{read\allowbreak Kraus\_\allowbreak gram},
\texttt{read\allowbreak Kraus\_\allowbreak conjugation}, \texttt{controlled},
\texttt{controlled\_\allowbreak operator\_\allowbreak eq\_\allowbreak block} \\
P-ARB-02: from PSD blocks to the action bound &
\texttt{pairing\_\allowbreak operator\_\allowbreak difference\_\allowbreak le\_\allowbreak mass\_\allowbreak mul\_\allowbreak defect},
\texttt{controlled\_\allowbreak state\_\allowbreak distance\_\allowbreak le\_\allowbreak weighted\_\allowbreak defect},
\texttt{controlled\_\allowbreak defect\_\allowbreak le\_\allowbreak uniform} \\
P-ARB-03: finite controller and chronology &
\texttt{Controller.\allowbreak leaf\allowbreak Equiv}, \texttt{Controller.\allowbreak protocol\_\allowbreak word\_\allowbreak eq},
\texttt{Controller.\allowbreak word\_\allowbreak snoc}, \texttt{record\allowbreak Operator\_\allowbreak append} \\
P-ARB-03: actual steps and the original complete state & \texttt{step},
\texttt{step\_\allowbreak operator\_\allowbreak eq\_\allowbreak block}, \texttt{run\_\allowbreak matrix},
\texttt{run\_\allowbreak eq\_\allowbreak protocol\_\allowbreak state}, \texttt{terminal\_\allowbreak memory\_\allowbreak matrix},
\texttt{weight\_\allowbreak run\_\allowbreak leaf} \\
P-ARB-04: actual per-action sum & \texttt{action\allowbreak Cost},
\texttt{run\_\allowbreak succ\_\allowbreak distance\_\allowbreak le},
\texttt{run\_\allowbreak distance\_\allowbreak le\_\allowbreak action\allowbreak Cost\_\allowbreak sum}, \texttt{action\allowbreak Cost\_\allowbreak le},
\texttt{run\_\allowbreak distance\_\allowbreak le\_\allowbreak budget\_\allowbreak sum} \\
P-ARB-04: complete records and original leaf laws &
\texttt{record\_\allowbreak tv\_\allowbreak le\_\allowbreak action\allowbreak Cost\_\allowbreak sum},
\texttt{record\_\allowbreak tv\_\allowbreak le\_\allowbreak budget\_\allowbreak sum}, \texttt{protocol\_\allowbreak record\_\allowbreak tv\_\allowbreak eq},
\texttt{protocol\_\allowbreak record\_\allowbreak tv\_\allowbreak le\_\allowbreak budget\_\allowbreak sum} \\
P-ARB-05: exact rational positive control &
\texttt{biased\_\allowbreak record\_\allowbreak defect\_\allowbreak exact},
\texttt{first\_\allowbreak action\allowbreak Cost\_\allowbreak exact}, \texttt{second\_\allowbreak action\allowbreak Cost\_\allowbreak exact},
\texttt{two\_\allowbreak step\_\allowbreak record\_\allowbreak tv\_\allowbreak exact},
\texttt{original\_\allowbreak protocol\_\allowbreak record\_\allowbreak tv\_\allowbreak exact} \\
N-ARB-01: deleted labels and quantum marginals &
\texttt{negative\_\allowbreak unlabelled\_\allowbreak defect\_\allowbreak zero},
\texttt{negative\_\allowbreak two\_\allowbreak step\_\allowbreak record\_\allowbreak tv},
\texttt{exact\_\allowbreak second\_\allowbreak record\_\allowbreak zero},
\texttt{flipped\_\allowbreak second\_\allowbreak record\_\allowbreak zero},
\texttt{negative\_\allowbreak terminal\_\allowbreak memories\_\allowbreak equal} \\
\end{longtable}

The two public modules contain 54 explicit \texttt{theorem}
declarations, 20 definitions, and 1 abbreviation, giving 75 actual named
axiom-audit targets. The separate witness contains 36 \texttt{theorem}
declarations, 14 definitions, and 2 abbreviations, giving 52 targets.
All three source files have actual successful focused-compilation
receipts. The two complete audit sequences follow actual source order
and use only \texttt{propext}, \texttt{Classical.\allowbreak choice}, and
\texttt{Quot.\allowbreak sound}, or subsets thereof. Development failures and final
success logs are preserved; ordinary linter warnings remain in the final
logs. An independent internal B→C review covered the public proofs and
complete concrete witness, respectively. Actual evidence is in
\texttt{physics/\allowbreak qm/\allowbreak kernel/\allowbreak evidence/\allowbreak v31.\allowbreak 8/\allowbreak C}. Integrated acceptance and
remote publication are recorded separately by the coordinator's complete
receipts; local success does not substitute for those identities.

The closed mathematical obligation is precise: given explicit finite
complex quantum instruments and a common causal action rule, prediction
error can be decomposed through actual history channels into
action-record defects and passed to the complete record distribution. A
protocol that is difficult to assess as a whole thus has individually
reviewable error inputs. Memory dimension, completeness of each action,
record scope, the actual controller, and input-state error remain
premises.

Chapter 65 supplies the finite-stopping-tree responsibility left by the
fixed-depth result: the common stopping rule, uniform finite bound,
actual absorbing padding, and original joint leaf state all receive
explicit proofs. Coherent controllers, arbitrary environmental memory
outside the present carrier, diamond distance, unbounded stopping, and
continuous-time limits still require their own definitions and proofs.
Real-device calibration of action defects, leakage outside the model,
and empirical applicability also remain separate responsibilities. This
chapter adds no experiment, does not derive a complete natural theory
from a finite kernel, and does not change the spectral, representation,
threshold, or effective-window conditions of the generation-count
results. The frozen 106 Core items and four generators retain their
identity.

\chapter{63. Full-effect distance, trace absolute value, and optimal
binary discrimination}\label{ch-63}

The probability distance defined by a supremum has a concrete spectral
expression in the finite complex model, and an optimal test attains it.
A valid positive-spectral effect proves the half-trace-absolute-value
formula and optimal equal-prior binary discrimination: formal proofs of
established quantum-information results in this project's matrix
interfaces.

Chapter 50 constructed the full-effect distance from actual Born
probabilities and proved symmetry, the triangle inequality, and channel
contraction. Chapter 60 then proved that all effects separate complete
density states. This chapter determines its exact numerical expression
and whether its supremum is attained. For all finite complex density
states, the original full-effect supremum equals half the trace of the
standard matrix absolute value, and an actual valid positive-spectral
projection attains it.

This connection has two uses. Theoretically, it relates the greatest
distinction available to valid effects to the spectrum of a concrete
matrix. In applications, it permits certain distance questions to be
reduced to finite matrix certificates. The mathematical library's
spectral decomposition, however, is a noncomputable existence choice.
Its existence proof is not a general numerical eigensolver. {[}QM319{]}

\section{63.1 The Original Distance and the Standard Matrix Absolute
Value}

Fix any finite carrier \(\iota\) and retain the complete density space
\(\mathcal D(\iota)\) and effect space \(\mathcal E(\iota)\). The
probability of an actual effect \(E\) in density state \(\rho\) is

\[
p_\rho(E)=\operatorname{Re}\operatorname{Tr}(E\rho).
\]

The distance from Chapter 50 is

\[
d_{\rm eff}(\rho,\sigma)
=\sup_{0\preceq E\preceq I}|p_\rho(E)-p_\sigma(E)|.
\]

It quantifies over all valid effects. A maximum probability difference
over a fixed instrument library, finite measurement sample, or diagonal
effects alone is not a substitute for this definition. This chapter
neither changes the supremum nor places the desired conclusion into a
new alias.

The standard matrix absolute value is defined by the positive square
root:

\[
|X|=\sqrt{X^\dagger X},
\qquad
\tau(X)=\operatorname{Re}\operatorname{Tr}|X|.
\]

Since \(|X|\succeq0\), its trace is real and nonnegative. Lean's
\texttt{CFC.\allowbreak abs\ X} is this square root, and \texttt{trace\allowbreak Absolute\ X}
here is its actual real trace. Although \(\tau\) is the usual numerical
expression for the trace norm, this work does not establish all norm
axioms on the arbitrary matrix space or add a new \texttt{Normed\allowbreak Space}
instance. The actual distance theorem uses Hermiticity and zero trace of
a density difference.

\section{63.2 Actual Finite Spectral Transport and Positive and Negative
Parts}

Let \(H=H^\dagger\). The pinned finite-matrix spectral theorem supplies
an actual unitary \(U\) and real eigenvalue family
\((\lambda_i)_{i\in\iota}\) such that

\[
U^\dagger U=UU^\dagger=I,
\qquad
H=U\operatorname{diag}(\lambda_i)U^\dagger.
\]

These data follow from the spectral theorem. No physical axiom asserting
existence of the required spectral decomposition is added to the quantum
model, and the caller need not supply an optimal effect in advance.

To distinguish actual matrix multiplication from scalar operations on
the spectrum, write

\[
T_U(d)=U\operatorname{diag}(d_i)U^\dagger,
\qquad d_i\in\mathbb R.
\]

Entrywise matrix algebra gives

\[
\begin{aligned}
T_U(d+e)&=T_U(d)+T_U(e),\\
T_U(d-e)&=T_U(d)-T_U(e),\\
T_U(d)T_U(e)&=T_U(de),\\
T_U(0)&=0,\qquad T_U(1)=I.
\end{aligned}
\]

The third identity uses \(U^\dagger U=I\) in the middle. An unverified
similarity matrix cannot replace the unitary. If every \(d_i\ge0\), the
diagonal matrix is PSD and unitary conjugation preserves PSD, giving
\(T_U(d)\succeq0\). Cyclicity of trace also gives

\[
\operatorname{Re}\operatorname{Tr}T_U(d)=\sum_i d_i.
\]

\phantomsection\label{p-tdist-01}{}

\textbf{P-TDIST-01 (Positive and negative spectral parts and the actual
absolute value of a finite Hermitian matrix).} For any finite Hermitian
matrix \(H\), define

\[
P=T_U\bigl(\max(\lambda_i,0)\bigr),
\qquad
N=T_U\bigl(\max(-\lambda_i,0)\bigr).
\]

Then \(P,N\succeq0\), and

\[
H=P-N,\qquad |H|=P+N,
\qquad
\tau(H)=\sum_i|\lambda_i|.
\]

\textbf{Proof.} Positivity follows from nonnegativity of the two scalar
coefficient families. For each eigenvalue,

\[
\max(\lambda,0)-\max(-\lambda,0)=\lambda,
\]

so the subtraction identity of \(T_U\) gives \(P-N=H\). Meanwhile,
matrix continuous functional calculus identifies the absolute value of
the Hermitian matrix exactly as

\[
|H|=T_U(|\lambda_i|).
\]

This uses the actual agreement of the defined \(\sqrt{H^\dagger H}\)
with the functional-calculus absolute value, not a new
spectral-absolute-value definition for this chapter. The scalar identity

\[
|\lambda|=\max(\lambda,0)+\max(-\lambda,0)
\]

and additivity of \(T_U\) give \(|H|=P+N\). Taking the actual trace and
using its invariance under unitary conjugation gives the finite sum of
absolute eigenvalues. \(\square\)

The formal counterparts are \texttt{positive\allowbreak Part},
\texttt{negative\allowbreak Part}, \texttt{positive\allowbreak Part\_\allowbreak sub\_\allowbreak negative\allowbreak Part},
\texttt{abs\_\allowbreak eq\_\allowbreak positive\allowbreak Part\_\allowbreak add\_\allowbreak negative\allowbreak Part}, and
\texttt{trace\allowbreak Absolute\_\allowbreak eq\_\allowbreak sum\_\allowbreak abs\_\allowbreak eigenvalues}. Zero eigenvalues
contribute zero to both parts. No eigenvalue is divided by, and no
full-rank assumption is added. {[}QM319{]}

\section{63.3 First Establishing the Positive-Spectral Projection as a
Valid Effect}

Define the real indicator coefficients

\[
e_i=\begin{cases}
1,&\lambda_i>0,\\
0,&\lambda_i\le0,
\end{cases}
\qquad E_+=T_U(e).
\]

Every \(e_i\) and \(1-e_i\) is nonnegative, so

\[
E_+\succeq0,
\qquad
I-E_+=T_U(1-e)\succeq0.
\]

Together, these two positivity statements make \(E_+\) an actual effect
in the original \texttt{Effect} type. PSD alone is insufficient: without
the upper bound, pairing with a positive operator can exceed the
probability range.

Since \(e_i^2=e_i\), one also has \(E_+^2=E_+\). The support relations
needed for optimality are actual matrix identities:

\[
E_+P=P,
\qquad
E_+N=0.
\]

They follow by checking scalar coefficients. If \(\lambda_i>0\), then
\(e_i=1\) and the negative-part coefficient vanishes. If
\(\lambda_i\le0\), then \(e_i=0\) and the positive-part coefficient
vanishes. Zero eigenvalues are assigned to the complementary effect
without changing any pairing value.

Consequently,

\[
\operatorname{Re}\operatorname{Tr}(E_+H)
=\operatorname{Re}\operatorname{Tr}(E_+P-E_+N)
=\operatorname{Tr}P.
\]

This evaluates the actual performance of a concrete candidate effect. An
upper bound for all effects is still needed to prove optimality. The
formalization proves validity, idempotence, both support identities, and
the pairing separately, rather than hiding optimality in the effect
structure's input fields. {[}QM319{]}

\section{63.4 From the All-Effect Bound to Exact Distance and
Attainment}

First consider any actual PSD decomposition \(H=P-N\), without requiring
a spectral decomposition. For every valid effect \(E\), Chapter 50
already proved

\[
0\le\operatorname{Re}\operatorname{Tr}(EP)\le\operatorname{Tr}P,
\qquad
0\le\operatorname{Re}\operatorname{Tr}(EN)\le\operatorname{Tr}N.
\]

Thus

\[
-\operatorname{Tr}N
\le\operatorname{Re}\operatorname{Tr}(EH)
\le\operatorname{Tr}P.
\]

Only now use normalization. If \(H=\rho-\sigma\), then

\[
\operatorname{Tr}H=1-1=0,
\qquad
\operatorname{Tr}P=\operatorname{Tr}N.
\]

Every effect therefore satisfies

\[
|p_\rho(E)-p_\sigma(E)|\le\operatorname{Tr}P.
\]

Any PSD decomposition gives this valid upper bound. Equality requires a
further effect that attains it.

\phantomsection\label{p-tdist-02}{}

\textbf{P-TDIST-02 (Half-trace-absolute-value formula and an optimal
effect for full-effect distance).} For any complete finite complex
density states \(\rho,\sigma\), set \(H=\rho-\sigma\) and construct
\(E_+\) as above. Then

\[
\begin{aligned}
d_{\rm eff}(\rho,\sigma)
&=\operatorname{Tr}P\\
&=\frac12\operatorname{Tr}|\rho-\sigma|\\
&=\frac12\sum_i|\lambda_i(\rho-\sigma)|.
\end{aligned}
\]

The actual signed probability difference also satisfies

\[
p_\rho(E_+)-p_\sigma(E_+)
=d_{\rm eff}(\rho,\sigma).
\]

\textbf{Proof.} Applying the all-effect bound to the spectral positive
and negative parts gives an upper bound \(\operatorname{Tr}P\) on the
original supremum. The actual valid \(E_+\) has pairing
\(\operatorname{Tr}P\ge0\). Its value belongs to the set defining the
supremum, so that supremum is at least \(\operatorname{Tr}P\). Both
bounds coincide and are actually attained.

From \(|H|=P+N\) and equality of the two part traces implied by zero
trace,

\[
\operatorname{Tr}|H|
=\operatorname{Tr}P+\operatorname{Tr}N
=2\operatorname{Tr}P.
\]

This supplies the factor \(1/2\). P-TDIST-01 gives the sum of absolute
eigenvalues. \(\square\)

The corresponding Lean results are
\texttt{effect\allowbreak Distance\_\allowbreak eq\_\allowbreak positive\allowbreak Part\_\allowbreak mass},
\texttt{positive\allowbreak Spectrum\allowbreak Effect\_\allowbreak attains},
\texttt{effect\allowbreak Distance\_\allowbreak eq\_\allowbreak half\_\allowbreak trace\allowbreak Absolute}, and
\texttt{effect\allowbreak Distance\_\allowbreak eq\_\allowbreak half\_\allowbreak sum\_\allowbreak abs\_\allowbreak eigenvalues}. They all use
the original \texttt{Density}, \texttt{Effect}, and
\texttt{effect\allowbreak Distance}; the input pair is not restricted to pure,
diagonal, or simultaneously diagonalizable states. {[}QM319{]}

\section{63.5 Equal-Prior Binary Discrimination}

Suppose each candidate state has prior probability \(1/2\). Decide for
the first state on the outcome of effect \(E\) and for the second on its
complement \(I-E\). The actual success probability is

\[
s(E;\rho,\sigma)
=\frac12\bigl(p_\rho(E)+p_\sigma(I-E)\bigr).
\]

The complementary probability is \(p_\sigma(I-E)=1-p_\sigma(E)\), so

\[
s(E;\rho,\sigma)
=\frac{1+p_\rho(E)-p_\sigma(E)}2.
\]

\phantomsection\label{p-dec-01}{}

\textbf{P-DEC-01 (Actual optimal success probability for equal-prior
binary discrimination).} Every valid effect satisfies

\[
s(E;\rho,\sigma)\le\frac{1+d_{\rm eff}(\rho,\sigma)}2,
\]

and the actual \(E_+\) constructed above attains the bound. Hence

\[
s_{\max}
=\frac{1+d_{\rm eff}(\rho,\sigma)}2
=\frac12+\frac14\operatorname{Tr}|\rho-\sigma|.
\]

\textbf{Proof.} The signed probability difference is at most its
absolute value, which is at most the original full-effect supremum.
Substitute this into the success-probability identity to obtain the
upper bound. The attainment statement in P-TDIST-02 makes the success
probability of \(E_+\) equal to that bound. An actual optimal decision
therefore exists. \(\square\)

The proof concerns binary effects and their probabilities, without
imposing a unique postmeasurement state-update rule. Different
instruments can have the same effect probabilities and different
postmeasurement dynamics. Discrimination optimality alone does not imply
uniqueness of those dynamics.

The contract is limited to equal priors. With priors \(p\) and \(1-p\),
the weighted difference is \(p\rho-(1-p)\sigma\), whose trace is
\(2p-1\). When \(p\ne1/2\), this chapter's zero-trace half-absolute
pairing formula cannot be applied directly to that matrix. The
general-prior optimal discrimination formula requires a separate
nonzero-trace argument and is not counted among the completed results
here. {[}QM319{]}

\section{63.6 Finite Matrix Certificates and Supplied Spectral Data}

A noncomputable existence proof and a finite certificate checkable
component by component serve different purposes. The certificate here
takes neither the final distance nor an assertion that all effects
satisfy the final bound as input. Its inputs are matrix positivity and
multiplication identities.

\phantomsection\label{p-tcert-01}{}

\textbf{P-TCERT-01 (Distance certificate from a PSD support
decomposition).} Given actual density states \(\rho,\sigma\), matrices
\(P,N\), and a valid effect \(E\), suppose

\[
P,N\succeq0,\qquad
\rho-\sigma=P-N,\qquad
EP=P,\qquad EN=0,
\]

Then

\[
d_{\rm eff}(\rho,\sigma)=\operatorname{Tr}P,
\qquad
s(E;\rho,\sigma)=\frac{1+d_{\rm eff}(\rho,\sigma)}2.
\]

\textbf{Proof.} Normalization gives
\(\operatorname{Tr}P=\operatorname{Tr}N\). The
arbitrary-PSD-decomposition bound of §63.4 gives
\(d_{\rm eff}\le\operatorname{Tr}P\). The two support identities make
the actual probability difference \(\operatorname{Tr}P\). This is
nonnegative, so the original full-effect supremum is at least that
value. Both bounds coincide. The actual identity of §63.5 gives the
success probability. \(\square\)

An external tool need not certify that \(P,N\) are outputs of a named
eigensolver. The certificate requires actual positivity, an actual
matrix difference, and actual support relations. Whenever they hold, the
proved theorem guarantees the conclusion. Conversely, positivity and a
difference decomposition without the support relations generally provide
only an upper bound; see §63.8.

If the input already includes a concrete finite unitary decomposition

\[
\rho-\sigma=U\operatorname{diag}(d_i)U^\dagger,
\qquad U^\dagger U=UU^\dagger=I,
\qquad d_i\in\mathbb R,
\]

one can construct \(P,N,E\) from those actual data and verify the
certificate, obtaining

\[
d_{\rm eff}(\rho,\sigma)
=\sum_i\max(d_i,0)
=\frac12\sum_i|d_i|.
\]

This corollary accepts any actual unitary spectral data satisfying the
displayed identities; their eigenvector order need not match the
mathematical library's chosen order. The formal declarations are
\texttt{effect\allowbreak Distance\_\allowbreak eq\_\allowbreak mass\_\allowbreak of\_\allowbreak support\_\allowbreak certificate},
\texttt{success\_\allowbreak eq\_\allowbreak optimum\_\allowbreak of\_\allowbreak support\_\allowbreak certificate},
\texttt{effect\allowbreak Distance\_\allowbreak eq\_\allowbreak sum\_\allowbreak positive\_\allowbreak of\_\allowbreak unitary\_\allowbreak data}, and
\texttt{effect\allowbreak Distance\_\allowbreak eq\_\allowbreak half\_\allowbreak sum\_\allowbreak abs\_\allowbreak of\_\allowbreak unitary\_\allowbreak data}.
{[}QM319{]}

\section{63.7 Noncommuting Mixed States, Complex Phase, and Actual
Constants}

The first exact witness uses

\[
\rho=\begin{pmatrix}4/5&1/5\\1/5&1/5\end{pmatrix},
\qquad
\sigma=\begin{pmatrix}1/2&-1/5\\-1/5&1/2\end{pmatrix}.
\]

Positivity is not inferred backwards from the desired distance. The
state \(\rho\) is a nonnegative sum of a basis projection and a rank-one
positive matrix; \(\sigma\) is a nonnegative sum of a positive scalar
identity and another rank-one positive matrix. Both have trace one.
Their commutator is

\[
\rho\sigma-\sigma\rho
=\begin{pmatrix}0&-3/25\\3/25&0\end{pmatrix}\ne0,
\]

so this is not a test confined to a common diagonal family.

Set

\[
E=\begin{pmatrix}4/5&2/5\\2/5&1/5\end{pmatrix}.
\]

It is the positive projection of a Hermitian involution and satisfies
\(E^2=E\), \(E\succeq0\), and \(I-E\succeq0\). Take

\[
P=\frac12E,\qquad N=\frac12(I-E).
\]

The actual matrix difference is \(\rho-\sigma=P-N\), with \(EP=P\) and
\(EN=0\). The certificate therefore gives

\[
d_{\rm eff}(\rho,\sigma)=\frac12.
\]

Direct matrix-trace calculation also gives

\[
p_\rho(E)=\frac{21}{25},
\qquad
p_\sigma(E)=\frac{17}{50},
\qquad
s(E;\rho,\sigma)=\frac34.
\]

P-TDIST-02 yields \(\operatorname{Tr}|\rho-\sigma|=1\). This witness
cross-checks positivity, noncommutativity, actual effect probabilities,
the support certificate, the absolute-value trace, and the constants in
the success probability.

The second witness uses the two spectral pure states of Pauli \(Y\):

\[
\rho_{+i}=\frac{I+Y}{2},
\qquad
\rho_{-i}=\frac{I-Y}{2}.
\]

The corresponding complex-phase effect \(E=(I+Y)/2\) gives distance
\(1\), trace absolute value \(2\), and success probability \(1\). Yet
every real valid effect assigns equal probabilities to the two states,
as does every diagonal valid effect. Discarding imaginary parts or
restricting measurement to diagonal entries therefore loses a fully
distinguishable difference.

A further concrete complex unitary and two rational diagonal-density
coefficient families instantiate the supplied-spectrum certificate and
give distance \(1/2\). This checks the actual unitary direction and
matrix-difference decomposition. Approximate eigenvalues returned by a
numerical library are not treated as exact input. {[}QM319{]}

\section{63.8 Negative Controls: Zero Trace and Support Conditions Are
Both Necessary}

For the first negative control, take \(H=I\) on a one-dimensional
carrier. Its absolute value is still \(I\), so

\[
\frac12\operatorname{Tr}|H|=\frac12.
\]

But the valid effect \(E=I\) already gives

\[
|\operatorname{Tr}(EH)|=1>\frac12.
\]

Without zero trace, even the assertion that half the trace absolute
value bounds every valid-effect pairing fails. Section 63.4 first uses
\(\operatorname{Tr}\rho=\operatorname{Tr}\sigma=1\) and only then
obtains the factor \(1/2\); that step has substantive logical force.

For the second control, again in one dimension, compare a state with
itself, whose actual distance is zero. Nevertheless one can write

\[
0=P-N,\qquad P=N=\frac12I,
\]

with both parts PSD and \(\operatorname{Tr}P=1/2\). A PSD difference
decomposition therefore gives a general upper bound, not automatically
the optimum. Requiring \(EP=P\) and \(EN=0\) would immediately
contradict \(P=N\ne0\), so the support certificate correctly rejects
this decomposition.

For equal states the actual matrix difference and its standard
matrix-absolute-value trace are zero. For the orthogonal phase states,
the distance is one and that trace is two. These endpoint controls check
zero-eigenvalue handling and the normalization constant. {[}QM319{]}

\section{63.9 Scope and Remaining Obligations}

Within the complete finite complex matrix model, this chapter closes
three precise obligations: the original full-effect distance has the
standard matrix-absolute-value expression; an actual valid effect
attains the original supremum; and equal-prior binary discrimination and
concrete finite matrix certificates have checkable optimal values. No
new physical selection principle is added, and neither the complete
state space nor the complete effect space is redefined.

The remaining boundaries are explicit. Mathematical existence of
spectral data is not an executable numerical algorithm. Approximate
matrices, rounding error, and automatic certificate generation need new
error contracts. A state-distance formula does not automatically
establish a channel distance on arbitrary auxiliary environments or a
diamond norm. Coherent controllers, infinite-dimensional spectra, and
unbounded operators are also outside this chapter.

A more complete internal derivation in the complex quantum model
clarifies which observational differences follow from the model and
which come from apparatus errors or a restricted test library. A
completed finite mathematical segment does not establish that the model
is nature's unique complete theory. There are no new real experiments
here; every concrete example is exact arithmetic or a formal witness.
{[}QM319{]}

\chapter{64. Minimum auxiliary dimension and direct isometries for
purification}\label{ch-64}

The ancilla dimension needed for purification, and direct comparison of
different ancillas, depend on exact rank and isometry direction. This
chapter proves attainment and the lower bound for the minimum dimension;
experimental questions involving small eigenvalues require a separate
approximation tolerance rather than a silent change to exact rank.

Chapter 57 proved that every finite complex density state has a
purification. Chapter 61 further characterized complete convex extreme
points as normalized vector states and compared all purifications on the
same auxiliary carrier. The earlier existence construction, however,
usually uses an auxiliary space as large as the system itself. It does
not determine the auxiliary dimension needed for a particular state. Nor
does comparison through a common direct sum say when one auxiliary space
can be embedded isometrically into the other directly.

This chapter resolves the two questions separately. In the declared
complete finite complex density model, a purification on a specified
auxiliary carrier exists exactly when its dimension is at least the
original state's complex matrix rank. For any two actual purifications,
an actual auxiliary isometry connects them directly whenever the target
auxiliary space is no smaller than the source. The minimum dimension is
neither the number of terms in a prespecified decomposition nor a
``minimum'' named after an auxiliary type selected in a file. It is
determined by the original density matrix. {[}QM319{]}

\section{64.1 The Original State Space, Actual Rank, and the Problem to
Be Proved}

Continue to write

\[
\mathcal D(\iota)=\{\rho\in M_\iota(\mathbb C):\rho\succeq0,\quad
\operatorname{Tr}\rho=1\}.
\]

The carriers \(\iota,\alpha,\beta\) are arbitrary finite types, with
auxiliary complex dimensions \(|\alpha|,|\beta|\). The rank
\(\operatorname{rank}_{\mathbb C}\rho\) is the dimension of the actual
linear map's image. Every purity assertion uses convex extremality in
the original complete density-state space. A normalized vector
representative is obtained through the direction of P-PUC-01 that takes
complete convex purity to a vector representation.

For a given \(\rho\in\mathcal D(\iota)\), purification on a specified
\(\alpha\) means finding

\[
\Psi\in\mathcal D(\iota\times\alpha),\qquad
\operatorname{Pure}(\Psi),\qquad
\operatorname{Tr}_\alpha\Psi=\rho.
\]

The partial trace remains the rectangular Kraus channel implemented in
Chapter 58. No different partial trace valid only on vector states is
introduced, and the allowed joint-state set is not reduced. The main
theorem does not additionally assume a nonempty input carrier: existence
of a unit-trace density already excludes the relevant empty-state case.

\section{64.2 The Actual Coordinate-Inclusion Matrix and the Necessary
Direction}

First handle finite-dimensional geometry independently of any input
quantum state. For an injection \(e:\alpha\hookrightarrow\beta\), define

\[
(J_e)_{ba}=\begin{cases}1,&b=e(a),\\0,&b\ne e(a).\end{cases}
\]

For every \(a,a'\in\alpha\),

\[
(J_e^\dagger J_e)_{aa'}
=\sum_{b\in\beta}\overline{(J_e)_{ba}}(J_e)_{ba'}
=\mathbf1_{e(a)=e(a')}=\mathbf1_{a=a'}.
\]

Thus \(J_e^\dagger J_e=I_\alpha\). The finite-cardinality inequality
\(|\alpha|\le|\beta|\) supplies an actual injection and therefore an
actual rectangular isometry. This step does not invoke purification
existence or place the desired matrix into an assumption field.

\phantomsection\label{p-pdm-01}{}

\textbf{P-PDM-01 (Exact dimension ordering for a rectangular isometry).}
For finite carriers \(\alpha,\beta\),

\[
\exists J\in M_{\beta\times\alpha}(\mathbb C),\quad J^\dagger J=I_\alpha
\quad\Longleftrightarrow\quad |\alpha|\le|\beta|.
\]

\textbf{Proof.} Coordinate inclusion proves sufficiency. Necessity does
not depend on constructing the matrix that way. For any
\(J^\dagger J=I_\alpha\),

\[
|\alpha|=\operatorname{rank}(I_\alpha)
=\operatorname{rank}(J^\dagger J)
\le\operatorname{rank}(J)\le|\beta|.
\]

The product-rank and row-count bounds are actual complex-linear-algebra
theorems. \(\square\)

If \(|\alpha|<|\beta|\), one cannot also have \(JJ^\dagger=I_\beta\).
Applying the same necessity argument to \(J^\dagger\) would then imply
\(|\beta|\le|\alpha|\). The one-sided isometry condition for strict
enlargement is therefore the exact dimension-forced scope, not a
notational preference.

\section{64.3 From Same-Carrier Unitary Equivalence to a Direct Isometry
Between Different Carriers}

First let \(\psi:\iota\times\alpha\to\mathbb C\) and
\(\phi:\iota\times\beta\to\mathbb C\) have outer products with the same
system marginal, and assume \(|\alpha|\le|\beta|\). Choose
\(J_0:\mathbb C^\alpha\to\mathbb C^\beta\) as above and set

\[
\chi=(I_\iota\otimes J_0)\psi.
\]

The actual rectangular-channel theorem P-AUX-01 guarantees

\[
\operatorname{Tr}_\beta |\chi\rangle\langle\chi|
=\operatorname{Tr}_\alpha|\psi\rangle\langle\psi|
=\operatorname{Tr}_\beta|\phi\rangle\langle\phi|.
\]

Now \(\chi,\phi\) use the same auxiliary carrier \(\beta\). Chapter 61's
actual vector-unitary comparison therefore gives
\(U^\dagger U=UU^\dagger=I_\beta\) and \(\phi=(I_\iota\otimes U)\chi\).
Put

\[
J=UJ_0.
\]

Then \(J^\dagger J=J_0^\dagger U^\dagger UJ_0=I_\alpha\), and Kronecker
multiplication gives

\[
\phi=(I_\iota\otimes J)\psi,
\qquad
\phi(i,b)=\sum_{a\in\alpha}J_{ba}\psi(i,a).
\]

If \(\psi\) is rearranged into the matrix \(A_{ia}=\psi(i,a)\), the
transformed factor is \(AJ^{\mathsf T}\). This is an ordinary transpose
because \(J\) acts on the vector's auxiliary index; the factor is not
\(AJ^\dagger\). Conjugation belongs in outer products and isometry
identities and cannot be moved to the wrong position.

\phantomsection\label{p-pdm-02}{}

\textbf{P-PDM-02 (A dimension-directed direct isometry for arbitrary
actual purifications).} Suppose \(\Psi\in\mathcal D(\iota\times\alpha)\)
and \(\Phi\in\mathcal D(\iota\times\beta)\) are pure in the complete
convex sense, have equal actual system marginals, and satisfy
\(|\alpha|\le|\beta|\). There is then
\(J\in M_{\beta\times\alpha}(\mathbb C)\) such that

\[
J^\dagger J=I_\alpha,\qquad
\Phi=(I_\iota\otimes J)\Psi(I_\iota\otimes J)^\dagger.
\]

\textbf{Proof.} P-PUC-01 supplies normalized vectors \(\psi,\phi\) from
the two complete convex-pure states. The actual marginal equality
becomes the equality of vector-outer-product marginals used above,
yielding an actual \(J\) and its vector action. The identity
\(|Tv\rangle\langle Tv|=T|v\rangle\langle v|T^\dagger\) gives the
required density-matrix relation. \(\square\)

The rank may be smaller than either auxiliary dimension; neither
auxiliary support is assumed to fill its space. When the dimension
direction is reversed, the common-direct-sum comparison of Chapter 61
remains valid. A correspondence on lower-dimensional supports has not
been silently extended into a full-space isometry from a larger
auxiliary space to a smaller one.

\section{64.4 An Auxiliary-Dimension Lower Bound for Every Purification}

Let \(\Psi\in\mathcal D(\iota\times\alpha)\) be any pure state in the
complete sense. P-PUC-01 gives \(\Psi=|\psi\rangle\langle\psi|\). Define
the rectangular factor \(A\in M_{\iota\times\alpha}(\mathbb C)\) by
\(A_{ia}=\psi(i,a)\). Entrywise, the actual partial trace is

\[
\rho_{ij}=
\sum_{a\in\alpha}\psi(i,a)\overline{\psi(j,a)}
=(AA^\dagger)_{ij}.
\]

\phantomsection\label{p-pdm-03}{}

\textbf{P-PDM-03 (Actual rank lower bound for all convex-pure
extensions).} For every such pure state and actual marginal,

\[
\operatorname{rank}_{\mathbb C}(\operatorname{Tr}_\alpha\Psi)
=\operatorname{rank}_{\mathbb C}(AA^\dagger)
=\operatorname{rank}_{\mathbb C}A
\le |\alpha|.
\]

\textbf{Proof.} For any complex matrix \(A\), the equality
\(A^\dagger A x=0\) implies \(\|Ax\|_2^2=0\), hence \(Ax=0\); the
reverse implication is immediate. Equality of kernels and
finite-dimensional rank--nullity give
\(\operatorname{rank}(A^\dagger A)=\operatorname{rank}A\). Applying this
to \(A^\dagger\) yields the row-Gram rank identity. The image dimension
is at most the domain dimension, giving the last inequality. The actual
Lean proof invokes the corresponding matrix-rank theorems in the pinned
mathematical library. \(\square\)

The bound covers all purifications defined by complete convex
extremality. Different vector representatives, nonorthogonal columns, or
complex phases cannot evade it: every such representation has the same
actual system marginal.

\section{64.5 Actual Gram Compression onto Nonzero Spectral Support}

Let \(\rho\in\mathcal D(\iota)\). The finite Hermitian spectral theorem
supplies an actual unitary \(U\) and real eigenvalues \(\lambda_a\)
satisfying

\[
\rho=U\operatorname{diag}(\lambda_a)U^\dagger,
\qquad \lambda_a\ge0,
\qquad \sum_a\lambda_a=1.
\]

Spectral existence follows from the actual positive matrix and the
pinned mathematical library; it is not an additional physical
assumption. Define the nonzero-eigenvalue index set and factor

\[
S_\rho=\{a\in\iota:\lambda_a\ne0\},
\qquad
G_{ia}=U_{ia}\sqrt{\lambda_a}\quad(a\in S_\rho).
\]

The spectral decomposition also gives
\(|S_\rho|=\operatorname{rank}_{\mathbb C}\rho\). For every \(i,j\),

\[
\begin{aligned}
(GG^\dagger)_{ij}
&=\sum_{a\in S_\rho}U_{ia}\sqrt{\lambda_a}
  \overline{U_{ja}\sqrt{\lambda_a}}\\
&=\sum_{a\in S_\rho}U_{ia}\lambda_a\overline{U_{ja}}\\
&=\sum_{a\in\iota}U_{ia}\lambda_a\overline{U_{ja}}
=\rho_{ij}.
\end{aligned}
\]

The second line uses \(\lambda_a\ge0\), ensuring that the square of the
real square root equals the eigenvalue. The third uses, term by term,
that every discarded eigenvalue is exactly zero. The chosen subtype is
not a ``numerical support'' selected by an arbitrary approximate
threshold. Very small but nonzero eigenvalues are retained.

After vectorizing as \(g(i,a)=G_{ia}\), the trace and partial-trace
identities already proved in Chapter 57 give

\[
\operatorname{Tr}|g\rangle\langle g|
=\operatorname{Tr}(GG^\dagger)=1,
\qquad
\operatorname{Tr}_{S_\rho}|g\rangle\langle g|=\rho.
\]

Thus \(|g\rangle\langle g|\) is an actual pure state of the complete
joint state space, with auxiliary dimension equal to the actual rank.
After compression, normalization, purity, and the original matrix
marginal are all established from the actual construction; they are not
verbally inherited from the uncompressed vector.

\section{64.6 Attainment of the Minimum and the Criterion for a
Specified Carrier}

\phantomsection\label{p-pdm-04}{}

\textbf{P-PDM-04 (Exact minimum auxiliary dimension of finite complex
purification).} For every finite complex density state \(\rho\) and
every specified finite auxiliary carrier \(\alpha\),

\[
\begin{aligned}
&\exists\Psi\in\mathcal D(\iota\times\alpha),\quad
  \operatorname{Pure}(\Psi)\ \land\ \operatorname{Tr}_\alpha\Psi=\rho\\
&\hspace{18mm}\Longleftrightarrow\quad
  \operatorname{rank}_{\mathbb C}\rho\le|\alpha|.
\end{aligned}
\]

\textbf{Proof.} P-PDM-03 proves necessity. For sufficiency, use the
actual purification on \(S_\rho\) from the preceding section. From
\(|S_\rho|=\operatorname{rank}\rho\le|\alpha|\) and P-PDM-01, obtain an
actual coordinate isometry associated with
\(S_\rho\hookrightarrow\alpha\). The corresponding single-Kraus CPTP
channel of P-AUX-01 preserves purity and the system marginal, giving the
required purification on the specified \(\alpha\). \(\square\)

Take \(\alpha=\operatorname{Fin}(\operatorname{rank}\rho)\). The right
side becomes a tautological rank inequality, so the lower bound is
actually attained. If \(|\alpha|<\operatorname{rank}\rho\), necessity
rules out purification. Unit trace also excludes rank zero: at rank zero
the construction would produce a density state on an empty auxiliary
carrier, but all matrices on the empty product carrier have trace zero.
Thus every input density state has \(\operatorname{rank}\rho\ge1\), and
an empty auxiliary system cannot perform the purification task here.

\section{64.7 Exact Matrix Witnesses, Dimension Counterexamples, and
Actual Auditing}

The finite witnesses use exact rational complex numbers. The first
vector comes from the factor

\[
A=\begin{pmatrix}3/5&0\\0&4/5\end{pmatrix},
\qquad
AA^\dagger=\begin{pmatrix}9/25&0\\0&16/25\end{pmatrix}=\rho.
\]

Its two nonzero diagonal entries give \(\det\rho=144/625\ne0\), so
\(\operatorname{rank}\rho=2\). Vectorizing \(A\) gives a normalized
purification with a two-dimensional auxiliary space and actual partial
trace \(\rho\). P-PDM-04 also proves that a one-dimensional
\texttt{Unit} auxiliary cannot purify this state. This is not the
classical purification obstruction; it is the exact resource lower bound
within the current complete complex model.

Now take \(v=(1,0)\). Its normalized outer product
\(\rho_0=|v\rangle\langle v|=\operatorname{diag}(1,0)\) is pure in the
complete convex sense. A rank-one outer product has rank at most one,
while positivity of density-state rank gives the opposite bound. Hence
the actual rank is \(\operatorname{rank}\rho_0=1\). P-PDM-04 supplies an
actual purification on the auxiliary carrier \texttt{Unit} and excludes
the empty carrier \texttt{Fin\ 0}. A one-dimensional ancillary Hilbert
space is the tensor unit and adds no nontrivial ancillary degree of
freedom. It nevertheless differs from the zero-dimensional space, on
which no normalized state exists.

Take the matrix from two to three dimensions

\[
J=\begin{pmatrix}
3/5&4i/5\\
4/5&-3i/5\\
0&0
\end{pmatrix}.
\]

Direct calculation gives \(J^\dagger J=I_2\), but \(JJ^\dagger\ne I_3\),
since its third diagonal entry is zero. The correct right-factor
direction gives

\[
B=AJ^{\mathsf T}
=\begin{pmatrix}
9/25&12/25&0\\
16i/25&-12i/25&0
\end{pmatrix}.
\]

Actual multiplication gives \(BB^\dagger=AA^\dagger=\rho\), and the two
normalized vectors satisfy
\(\operatorname{vec}B=(I\otimes J)\operatorname{vec}A\). This checks the
direction of the actual auxiliary matrix and the complete joint-density
relation. Its complex phases cannot be discarded without conditions.
Chapter 61's negative controls for three wrong action directions remain
in force.

\phantomsection\label{n-pdm-01}{}

\textbf{N-PDM-01 (Equal marginals do not remove the dimension direction
required by a direct isometry).} The pure states with two- and
three-dimensional auxiliaries above have the same marginal. Yet no
\(K\in M_{2\times3}(\mathbb C)\) satisfies \(K^\dagger K=I_3\). Their
reverse comparison therefore cannot be interpreted as an isometry from
the entire three-dimensional space into the entire two-dimensional
space.

\textbf{Proof.} Such a \(K\) would imply \(3\le2\) by necessity in
P-PDM-01, a contradiction. A common-direct-sum comparison or a
correspondence on supports can still exist; neither removes the
dimensional obstruction for a full-space isometry. \(\square\)

The actual source is concentrated in \texttt{Purification\allowbreak Isometry} and
\texttt{Purification\allowbreak Dimension}. Together, the two modules contain 23
explicit theorems, 3 definitions, 1 abbreviation, and 2 named typeclass
instances: 29 named audit targets, each with actual successful focused
compilation and axiom auditing. The separate witness
\texttt{Purification\allowbreak Dimension\allowbreak Examples} contains 25 theorems and 9
definitions, giving 34 named audits actually passed. All these
declarations use only the permitted foundations \texttt{propext},
\texttt{Classical.\allowbreak choice}, and \texttt{Quot.\allowbreak sound}. Ordinary linter
warnings remain in the public modules. Focused execution, package-wide
acceptance, official book typesetting, and remote publication each
retain their own actual identities and cannot substitute for one
another. {[}QM319{]}

\section{64.8 Closed Obligations and Remaining Physical Meaning}

Within the complete finite complex model, this chapter closes exact
purification dimension and dimension-directed direct isometry. It
completes two finite obligations expressly left by Chapter 61 at that
stage, without retrospectively rewriting that chapter's historical
machine receipts. Spectral-support selection and finite-unitary
existence use noncomputable mathematical choice. No eigenvalue decision
procedure with rounding-error bounds, stable-rank estimator, or optimal
preparation circuit is provided.

A tiny nonzero eigenvalue and an exactly zero eigenvalue have different
consequences for the exact minimum dimension. If a real apparatus needs
only approximate purification, its state distance, allowed error, and
preparation operations must first be specified before approximate
resource costs are studied. A numerical threshold cannot be inserted
silently into this rank definition. Infinite dimensions, infinite rank,
and finite-energy approximations also require new analytic conditions.

Finally, the mathematically complete finite complex purification
properties still presuppose the complex matrix state space and allowed
operational model. They cannot alone establish that UEOT Core must
select complex quantum theory, and a successful formal-kernel check is
not a new experiment confirming nature. The complete reconstruction
background, independent physical-selection grounds, applicability to
actual systems, and testable predictions retain their separate open
obligations.

\chapter{65. Variable-length stopping protocols, original joint leaf
states, and budgets weighted by continued activity}\label{ch-65}

A protocol that stops early should not continue charging action error to
branches that have already stopped. The original finite stopping tree is
completed to an actual absorbing process, then joint states and effects
are transported back to the original leaves; reference active weights
give the budget, without extending the theorem to unbounded stopping or
coherent control.

Chapter 58 constructed the actual leaf records and quantum memory of
arbitrary complete finite adaptive trees. Chapter 62 then proved
per-action error sums under a common fixed-depth controller. This
chapter connects the two. It retains the stopping time on each original
leaf path, pads already stopped paths with actual identity operations,
and proves that padding is related to the original joint leaf state by
an actual isometric embedding. The fixed-depth budget therefore returns
rigorously to the original variable-length stopping tree, and stopped
paths incur no later action error. {[}QM319{]}

A common stopping rule is a mathematical input. The two
action-instrument families may differ, but use the same complete history
to select an action or stop. These histories remain classical records
that are read and dephased. They are not upgraded to a coherent quantum
controller, and stopping times without a uniform finite upper bound are
not treated.

\section{65.1 The Original Complete Tree and Actual Stopping Times}

Write the original protocol as \(\pi\). It either stops immediately or
performs action \(a\) at a node and for every possible outcome
\(o\in O\), supplies a continuation \(\pi_o\). Every outcome requires a
definition, including outcomes whose actual probability is zero. The
original leaf type \(L(\pi)\) is built recursively from the tree. Each
leaf \(l\) has an actual action--outcome word

\[
w_\pi(l)=((a_0,o_0),\ldots,(a_{\tau(l)-1},o_{\tau(l)-1})),
\qquad \tau(l)=|w_\pi(l)|.
\]

The quantity \(\tau(l)\) is the actual stopping time; padding does not
change it into a uniform depth \(N\). Since the outcome type is finite,
the tree's maximum height can be defined recursively:

\[
H(\mathrm{stop})=0,
\qquad H(\mathrm{node}(a,\pi_\cdot))=1+\max_{o\in O}H(\pi_o),
\]

where an empty finite maximum takes the natural-number bottom value
\(0\). Induction through the tree gives \(\tau(l)\le H(\pi)\). Fixing
\(N\ge H(\pi)\) allows every original leaf to fit completely into a
length-\(N\) history.

\section{65.2 An Absorbing Stop Symbol and Actual Kraus Completeness}

Extend the action, outcome, and Kraus-index types to
\(A_\bot=\mathrm{Option}(A)\), \(O_\bot=\mathrm{Option}(O)\), and
\(K_\bot=\mathrm{Option}(K)\). The values \texttt{some\ a} and
\texttt{some\ o} denote original actions and outcomes; \texttt{none} is
an explicit idle symbol after stopping.

For an original action \(a\), the extended instrument takes a nonzero
value only at \((\mathrm{some}\ o,\mathrm{some}\ k)\), where it uses the
original Kraus matrix \(K_{a,o,k}\); all other terms vanish. The idle
action has exactly one nonzero Kraus matrix: \(I\) at
\texttt{(none,\ none)}. Thus

\[
\begin{aligned}
\widehat\Phi_{\mathrm{some}\ a,\mathrm{some}\ o}(X)&=\Phi_{a,o}(X),&
\widehat\Phi_{\mathrm{some}\ a,\mathrm{none}}(X)&=0,\\
\widehat\Phi_{\mathrm{none},\mathrm{none}}(X)&=X,&
\widehat\Phi_{\mathrm{none},\mathrm{some}\ o}(X)&=0.
\end{aligned}
\]

These four identities quantify over arbitrary complex matrices \(X\),
not only normalized states or their traces. Completeness of active
actions follows from the original Kraus sum, and idle completeness from
\(I^\dagger I=I\). Every step is therefore an actual valid instrument;
idle identity operations preserve quantum memory.

Starting from the original tree, the controller follows recorded
\texttt{some} outcomes and continues to select \texttt{none} after
reaching a stop leaf. Illegal histories also receive continuations so
that the controller is a total function on all finite histories.
However, an outcome \texttt{none} at an active node, or \texttt{some\ o}
after stopping, immediately makes the branch matrix zero. Every later
linear record operator keeps it zero.

\section{65.3 Original-Leaf Encoding, Legal Histories, and Blockwise
Identities}

Encode an original leaf \(l\) by wrapping each outcome of its word in
\texttt{some}, then appending \(N-\tau(l)\) copies of \texttt{none}.
This gives \(e_N(l)\in O_\bot^N\):

\[
e_N(l)=\bigl(\mathrm{some}\ o_0,\ldots,
\mathrm{some}\ o_{\tau(l)-1},
\underbrace{\mathrm{none},\ldots,\mathrm{none}}_{N-\tau(l)}\bigr).
\]

The decoder requires the history to follow the original tree exactly to
a leaf, with all remaining positions equal to \texttt{none}. A complete
history is illegal if it has not stopped by depth \(N\), contains
\texttt{none} prematurely, or contains a real outcome after stopping.
Decoding an encoded leaf is proved to return that leaf. Conversely, any
history decoding to \(l\) is proved equal to \(e_N(l)\). Thus encoding
is injective, and the legal complete histories are exactly the image of
original-leaf encoding.

\phantomsection\label{p-stop-01}{}

\textbf{P-STOP-01 (Actual absorbing padding of an original finite
stopping tree).} For any \(N\ge H(\pi)\), any original input matrix
\(X\), and any padded history \(h\), the record operator generated by
the extended action family under the common padded controller satisfies

\[
\widehat\Phi_h(X)=
\begin{cases}
\Phi_{w_\pi(l)}(X),& h=e_N(l),\\
0,&h\notin e_N(L(\pi)).
\end{cases}
\]

\textbf{Proof.} Induct on remaining depth and the original tree. At an
immediate stop, the only legal continuation is entirely \texttt{none},
and every step is the identity; any \texttt{some} makes the matrix zero.
At an active node, the \texttt{none} branch vanishes, while a
\texttt{some\ o} branch first applies the original \(\Phi_{a,o}\) and
then recursively follows the subtree. The subtree height is at most the
remaining depth, so induction applies. The encoding and decoding inverse
identities exclude two different leaves producing the same legal
complete history. The proof compares matrices throughout; it never
infers matrix equality from equal traces. □

The original and padded joint states are respectively

\[
\Omega_I^\pi(\rho)
=\sum_{l\in L(\pi)}|l\rangle\langle l|\otimes
\Phi_{w_\pi(l)}^I(\rho),
\]

\[
\widehat\Omega_{I,N}^\pi(\rho)
=\sum_{h\in O_\bot^N}|h\rangle\langle h|\otimes
\widehat\Phi_h^I(\rho).
\]

P-STOP-01 therefore identifies the complete joint matrices: each legal
block equals its original block, and every illegal block is zero. The
encoding preserves stopping times and original labels.

\section{65.4 Isometric Extension of Valid Effects and Distance
Preservation}

To return a padded error bound to the original leaf space, one cannot
merely name a ``decoder'' and assume contraction. Let \(J\) be any
actual rectangular isometry satisfying \(J^\dagger J=I\). It defines the
actual single-Kraus CPTP channel \(\rho\mapsto J\rho J^\dagger\).

For an original valid effect \(0\preceq E\preceq I\), define its
extension \(\widetilde E=JEJ^\dagger\). Conjugation preserves PSD,
giving positivity. Its complement is

\[
I-JEJ^\dagger
=(I-JJ^\dagger)+J(I-E)J^\dagger.
\]

The matrix \(JJ^\dagger\) is a self-adjoint idempotent projection. Thus
\(Q=I-JJ^\dagger\) is also self-adjoint and idempotent, with
\(Q=QQ^\dagger\succeq0\). Both terms on the right are PSD, so
\(\widetilde E\) is valid on the entire enlarged space. Its actual
effect pullback is

\[
J^\dagger\widetilde E J
=(J^\dagger J)E(J^\dagger J)=E.
\]

\phantomsection\label{p-stop-02}{}

\textbf{P-STOP-02 (Full-effect distance preservation and the connection
to the original joint leaf state).} For any finite rectangular isometry
\(J\) and complete density states \(\rho,\sigma\),

\[
d_{\rm eff}(J\rho J^\dagger,J\sigma J^\dagger)
=d_{\rm eff}(\rho,\sigma).
\]

For the coordinate isometry
\(V_N:|l\rangle\otimes|i\rangle\mapsto|e_N(l)\rangle\otimes|i\rangle\),

\[
\widehat\Omega_{I,N}^\pi(\rho)
=V_N\Omega_I^\pi(\rho)V_N^\dagger,
\]

The original and padded joint-state distances therefore agree, as do TV
of the original complete-leaf distributions and TV of the padded-record
distributions.

\textbf{Proof.} CPTP contraction gives one distance inequality. Extend
each original effect by the valid construction above; its
Born-probability difference is unchanged. Taking the supremum over all
effects gives the reverse inequality. The coordinate-isometry columns
are orthonormal. Its joint-state matrix equals the original matrix on
the encoded image and vanishes outside it, so P-STOP-01 gives the
displayed actual-state identity. For TV, reindex by injective leaf
encoding: probabilities agree term by term on legal histories, and both
probabilities vanish on illegal histories. No default leaf is chosen,
and no object is silently added to a possibly empty type. □

\section{65.5 Actual Active-Action Defects and Zero Idle Cost}

The original record defect \(\delta_{\rm rec}(I_a,J_a)\) is the actual
joint-output distance maximized over all input density states. For the
extended action \texttt{some\ a}, its record--memory state is exactly
the isometric image of the original output under the coordinate
inclusion \(O\hookrightarrow O_\bot\). P-STOP-02 gives equality of
distances for each input. Taking the supremum over the same original
input set yields

\[
\delta_{\rm rec}(\widehat I_{\mathrm{some}\ a},
\widehat J_{\mathrm{some}\ a})
=\delta_{\rm rec}(I_a,J_a).
\]

The two implementations may use different Kraus-index carriers. Their
idle output nevertheless contains the same deterministic stop label and
original quantum memory, so

\[
\delta_{\rm rec}(\widehat I_{\mathrm{none}},
\widehat J_{\mathrm{none}})=0.
\]

Actual states are compared first, the zero supremum is proved next, and
only then are idle terms removed from the cost sum. ``Stopping has no
charge'' is not an additional assumption.

\section{65.6 The Reference Implementation's Probability of Remaining
Active}

Write the common padded controller as \(C_t(h)\). Let
\(\widehat\omega_t^J(\sigma)\) be the actual joint history state
obtained by reference implementation \(J\) from initial state \(\sigma\)
after the first \(t\) steps. Define

\[
q_t^J(h;\sigma)
=\operatorname{Tr}\bigl[\widehat\omega_t^J(\sigma)\bigr]_{h,h}.
\]

Chapter 62's PSD-block, zero-trace-if-and-only-if-zero-block, and
normalization results give \(q_t^J(h;\sigma)\ge0\) and
\(\sum_hq_t^J(h;\sigma)=1\). The controller still executes an original
action at history \(h\) exactly when \(C_t(h)=\mathrm{some}\ a\). Define

\[
\mathcal A_t^J
=\sum_{\substack{h\in O_\bot^t\\ C_t(h)=\mathrm{some}\ a}}
q_t^J(h;\sigma)\,\delta_{\rm rec}(I_a,J_a),
\qquad
s_t^J=\sum_{h:\ C_t(h)\text{ is active}}q_t^J(h;\sigma).
\]

The quantity \(s_t^J\) is the actual probability that an original action
will still be executed at that time. It is not a posterior-renormalized
survival distribution: stopped mass remains in the complete record
state.

\phantomsection\label{p-stop-03}{}

\textbf{P-STOP-03 (Per-action error budget under a common finite
stopping rule).} For the same original complete finite tree \(\pi\), two
action families \(I,J\), any input density states \(\rho,\sigma\), and
\(N\ge H(\pi)\),

\[
d_{\rm eff}(\Omega_I^\pi(\rho),\Omega_J^\pi(\sigma))
\le d_{\rm eff}(\rho,\sigma)+\sum_{t=0}^{N-1}\mathcal A_t^J.
\]

Complete original-leaf record TV and terminal quantum-memory distance
are each at most the same right-hand side. This main weighted budget
needs no uniform action bound. Suppose additionally that, at time \(t\),
for every history \(h\), including zero-weight histories, whenever the
common controller chooses \(C_t(h)=\mathrm{some}\ a\), one has
\(\delta_{\rm rec}(I_a,J_a)\le\varepsilon_t\). Then

\[
d_{\rm eff}(\Omega_I^\pi(\rho),\Omega_J^\pi(\sigma))
\le d_{\rm eff}(\rho,\sigma)+\sum_{t=0}^{N-1}s_t^J\varepsilon_t.
\]

\textbf{Proof.} Apply the proved recursion of Chapter 62 in the padded
model. First propagate the preceding input error through I's actual CPTP
step, then compare both actions on the same reference-prefix state of J.
Normalized PSD principal blocks and valid block effects give defects
weighted by \(q_t^J\). After summing the time recursion, active extended
defects equal the original defects and idle defects vanish, leaving
exactly \(\sum_t\mathcal A_t^J\). P-STOP-02 replaces the padded joint
distance by the original joint-leaf distance. Valid event effects bound
leaf-record TV, and actual partial-trace contraction bounds
terminal-memory distance. Finally, multiply each active-action bound by
its nonnegative \(q_t^J\) and sum to obtain \(s_t^J\varepsilon_t\). □

The reference implementation and initial state cannot be exchanged
without a separate argument: the displayed weights belong to J's actual
process started at \(\sigma\). Nor can the uniform finite bound N be
replaced by the verbal condition ``eventually stops almost surely.'' No
infinite stopped process or unbounded optional-stopping theorem is
constructed here.

\section{65.7 Exact Rational Positive Control: Charging Only Continuing
Paths}

On the first step, both implementations use the same rational Kraus
instrument, which preserves quantum memory and records outcomes 0 and 1
with probabilities \(3/4,1/4\). Outcome 0 stops immediately; outcome 1
proceeds to a second action and then stops. I's second action
deterministically records 0; J's records 0 and 1 with probabilities
\(3/4,1/4\). The second action's actual record defect over all inputs is
\(1/4\).

The three actual original leaves are the early stop \((0)\) and late
stops \((1,0)\) and \((1,1)\). Their probabilities are

\[
p_I=(3/4,1/4,0),\qquad p_J=(3/4,3/16,1/16).
\]

\phantomsection\label{p-stop-04}{}

\textbf{P-STOP-04 (Exact \(1/16\) record difference in a stopping
protocol).} For these actual qubit instruments and their common stopping
tree, complete original-leaf record TV and padded-record TV are both
\(1/16\), and the sum of per-action costs is also \(1/16\).

\textbf{Proof.} The first action is identical, so its cost is 0. At the
second step, only the path whose first outcome is 1 remains active; its
reference-J weight is \(1/4\). Multiplication by the original action
defect \(1/4\) gives \(1/16\). This sum bounds TV above. The event ``the
final leaf is \((1,1)\)'' has probability difference \(1/16\), giving a
matching lower bound. The padded records are
\((\mathrm{some}\ 0,\mathrm{none})\),
\((\mathrm{some}\ 1,\mathrm{some}\ 0)\), and
\((\mathrm{some}\ 1,\mathrm{some}\ 1)\), so P-STOP-02 preserves their
TV. □

This example checks that the actual probability of remaining active
enters the theorem. Charging the worst-case error at each time level
would give a looser bound and cannot replace the exact accounting here.

\section{65.8 A Boundary Control for the Common Stopping Rule}

\phantomsection\label{p-stop-05}{}

\textbf{P-STOP-05 (Identical quantum actions do not control all records
under different stopping rules).} Use identical identity quantum
operations and an original instrument that deterministically outputs 0.
One controller stops immediately, while another performs one step before
stopping. At common padded depth 1, their records are deterministically
\texttt{none} and \texttt{some\ 0}, so TV is 1. Their terminal quantum
memories coincide and both equal the original input state.

\textbf{Proof.} The immediately stopping controller gives probability 1
to ``record is none.'' The one-step controller has a zero branch matrix
for none during its active action, giving that event probability 0.
Probability-distribution TV is at most 1, so it is exactly 1. The first
controller makes no quantum change; the second's actual sole nonzero
quantum branch is also the identity. Their terminal quantum states are
therefore equal. □

This negative control uses different stopping rules, explicitly outside
the common-controller premise of P-STOP-03. Comparing only the defects
of identical original action families, or deleting the stopping record,
is insufficient to compare the complete outputs of different stopping
rules. Replacing the equal-memory conclusion with a claim of unequal
quantum memory would also destroy the role of the counterexample.

\section{65.9 Completed Scope and Open Obligations}

Actual focused Lean acceptance closes the following scope: original
complete finite stopping trees; finite classical-history controllers;
actual \texttt{Option} absorbing padding; all joint leaf matrices; valid
effect transport; actual original-action defects; and the action budget
weighted by reference J's probability of remaining active. The two
public modules contain 50 explicit \texttt{theorem} declarations, with
69 named declarations passing actual axiom auditing. The concrete
witness contains 45 \texttt{theorem} declarations, 15 definitions, and 2
abbreviations, giving 62 targets actually executed and audited by name.
The allowed foundations are only \texttt{propext},
\texttt{Classical.\allowbreak choice}, and \texttt{Quot.\allowbreak sound}; ordinary linter
warnings remain. This is local mathematical acceptance of the listed
finite-model obligations. Integrated and remote-release identities are
recorded separately in their actual receipts. {[}QM319{]}

Coherent controllers, unmodelled external quantum environments,
stabilized or diamond distance, unbounded stopping times, infinite
conditioned processes, and empirical applicability remain open. Exact
rational witnesses built from quantum instruments are mathematical
positive and negative controls; they add no real experiment. They
neither alter the frozen Core responsibilities nor promote the finite
quantum model into a complete theory of nature.

\chapter{66. Finite conditional processes of the original instrument,
allowed implementations, and full costs}\label{ch-66}

\textbf{Reading guide.} Constructing a new process in which an original
experiment's rare accepted records become typical requires preserving
the full conditioned joint state of the visible record and terminal
quantum system, and proving that the new operations belong to the
allowed library. Building on Chapter 56's conditioning constructions and
the record-error budgets of Chapters 62 and 65, this chapter connects
backward acceptance effects, singular support, labelled Choi criteria,
initialization, and retry costs. The conclusions apply to the stated
finite complex models; realizability, no resource advantage, and
impossibility under the specified access constraints can each be valid
outcomes.

This chapter completes a finite but general chain of physical reasoning.
An original complete quantum instrument and a prespecified acceptance
event determine the full conditioned record and terminal state. The
allowed actions and observation permissions determine whether that
process can be implemented. Initialization, failures, resets, and
numerical errors determine its actual cost. The conclusions are
conditional theorems within finite-dimensional complex quantum theory.
Quantum Doob constructions already provide precedents for changing the
dynamics so that rare trajectories of the original process become
typical trajectories of another process, including finite-time,
time-dependent constructions {[}CFW-R1{]}. The task here is to connect
the original experimental process, visible records, singular support,
the actual action library, resources, and formal proofs. These
connections do not establish a new law of nature.

\section{66.1 The original experiment and its distinct outputs}

Fix a finite-dimensional complex Hilbert space and an initial state
\(\rho\). Each internal node \(h\) of a finite classical control tree
carries a complete instrument \[
\mathcal I_{h,o}(X)=\sum_\mu K_{h,o,\mu}XK_{h,o,\mu}^{\dagger},\qquad
\sum_{o,\mu}K_{h,o,\mu}^{\dagger}K_{h,o,\mu}=I.
\] Here \(o\) is an outcome that is actually visible in the original
experiment; \(\mu\) is a hidden Kraus index. The tree specifies the
action after each outcome, the stopping rule, and the set \(G\) of
accepted leaves. Every rejected leaf remains part of the original
experiment. Write \(X_\ell(\rho)\) for the unnormalized quantum block at
leaf \(\ell\), after summing all hidden paths compatible with the same
visible record. The output of the original experiment is \[
\Omega(\rho)=\sum_\ell |\ell\rangle\langle\ell|\otimes X_\ell(\rho).
\] When \(p_G>0\), the object of study is the full joint state
conditioned on acceptance, \[
\Omega_G(\rho)=p_G^{-1}\sum_{\ell\in G}
 |\ell\rangle\langle\ell|\otimes X_\ell(\rho),\qquad
p_G=\sum_{\ell\in G}\operatorname{Tr}X_\ell(\rho).
\] This state retains correlations between the record and the terminal
state. Matching only the final quantum marginal, or only the success
probability, is insufficient. Leaves may have different depths. An
operation's internal random seed does not automatically belong to the
target visible record either. If the experiment requires that seed to be
retained, the record space must be enlarged and the object of comparison
specified again.

\section{66.2 Backward acceptance effects come from the original
experiment}

Set \(H_\ell=I\) at an accepted leaf and \(H_\ell=0\) at a rejected
leaf. Define the effects backwards, node by node: \[
H_h=\sum_{o,\mu}K_{h,o,\mu}^{\dagger}H_{ho}K_{h,o,\mu}.
\]

\phantomsection\label{p-fcond-01}{} \textbf{P-FCOND-01 (Acceptance
effects of the same finite original experiment).} At every node,
\(0\preceq H_h\preceq I\). For any unnormalized positive input \(X\),
the total mass of all accepted leaves in that subtree is exactly
\(\operatorname{Tr}(H_hX)\). In particular,
\(p_G=\operatorname{Tr}(H_0\rho)\).

\textbf{Proof.} The claim holds directly at the leaves. Suppose it holds
at the child nodes. The CP dual preserves positivity, and completeness
gives \[
I-H_h=\sum_{o,\mu}K_{h,o,\mu}^{\dagger}(I-H_{ho})K_{h,o,\mu}\succeq0.
\] Sum the accepted masses of the branches and use cyclicity of the
trace to obtain \(\operatorname{Tr}(H_hX)\). Induction on the finite
tree completes the proof. Here \(p_G\) is the original record
probability, not an independently supplied spectral parameter.
\(\square\)

A special case accepts exactly those paths whose outcomes lie in \(S\)
for \(T\) consecutive uses of a fixed instrument: \[
H_T=I,\qquad H_t=\Phi_S^*(H_{t+1}),\qquad
\Phi_S=\sum_{o\in S}\mathcal I_o.
\] A general adaptive tree requires the nodewise recurrence. It cannot
be replaced by a power of one fixed \(\Phi_S^*\).

\section{66.3 Singular support cannot be discarded as a zero-row detail}

Let \(S_h=H_h^{1/2}\), let \(P_h\) be the support projection of \(H_h\),
and let \(T_h\) be the inverse of \(S_h\) on its support, extended by
zero on its kernel. The finite-dimensional spectral theorem gives \[
T_hS_h=S_hT_h=P_h,\quad P_h^2=P_h=P_h^{\dagger},\quad P_hS_h=S_h.
\]

\phantomsection\label{p-fcond-02}{} \textbf{P-FCOND-02 (Kernel
annihilation from the backward Gram sum).} The original backward
recurrence implies, for every hidden branch, \[
S_{ho}K_{h,o,\mu}(I-P_h)=0.
\] Consequently, \(S_{ho}K_{h,o,\mu}P_h=S_{ho}K_{h,o,\mu}\). This is a
derived identity, not an additional assumption about paths.

\textbf{Proof.} For any \(v\in\ker H_h\), \[
0=\langle v,H_hv\rangle
=\sum_{o,\mu}\|S_{ho}K_{h,o,\mu}v\|^2.
\] A finite sum of nonnegative numbers is zero only if every term is
zero. Take \(v=(I-P_h)u\). \(\square\)

The formal interface allows matrices \(S,T,P\) and local Gram identities
to be supplied and checked individually. It does not require the desired
conditioned record law to be included in the input structure. The
canonical spectral construction produces these matrices from the
original \(H\). A noncomputable spectral existence proof and the
reliability of a floating-point eigendecomposition algorithm are
separate obligations.

\section{66.4 A complete conditioned instrument with the original
visible labels}

First take the principal branches \[
L_{h,o,\mu}=S_{ho}K_{h,o,\mu}T_h.
\] Choose any complete instrument \(J_{h,o,\nu}\) with the same visible
label set, and add the hidden completion branches
\(J_{h,o,\nu}(I-P_h)\). The label remains \(o\), while the hidden index
is enlarged to the disjoint sum of the two types of Kraus index.

\phantomsection\label{p-fcond-03}{} \textbf{P-FCOND-03 (Support
transformation and a valid completion on the full space).} The Gram sum
of the principal branches is \(P_h\), and that of the completion
branches is \(I-P_h\). The resulting \(\widetilde{\mathcal I}_{h,o}\) is
a complete CP instrument and satisfies, for every complex matrix \(X\),
\[
\widetilde{\mathcal I}_{h,o}(S_hXS_h)
=S_{ho}\mathcal I_{h,o}(X)S_{ho}.
\]

\textbf{Proof.} The backward recurrence gives \[
\sum_{o,\mu}L_{h,o,\mu}^{\dagger}L_{h,o,\mu}
=T_hH_hT_h=P_h.
\] The Gram sum of the completion branches is \((I-P_h)I(I-P_h)=I-P_h\).
By P-FCOND-02, \(L_{h,o,\mu}S_h=S_{ho}K_{h,o,\mu}\); moreover,
\((I-P_h)S_h=0\). Summing over the hidden indices yields the branch
transport identity. The proof holds for all complex matrices and does
not depend on taking traces first. \(\square\)

At a rejected leaf, \(S_\ell=0\), so the conditioned mass reaching it is
zero. Completion at a node with zero support must still be valid. An
incomplete or undefined operation cannot substitute for a complete
instrument.

\section{66.5 The full conditioned joint state is independent of the
completion}

For \(p_G>0\), define the actual initial state \[
\widetilde\rho_0=S_0\rho S_0/p_G.
\] Positivity and P-FCOND-01 ensure that its trace is one.

\phantomsection\label{p-fcond-04}{} \textbf{P-FCOND-04 (Conditioning the
full leaf record and terminal state from the same original experiment).}
Running the complete instrument tree constructed above gives, at every
original leaf, \[
\widetilde X_\ell(\widetilde\rho_0)=
\begin{cases}X_\ell(\rho)/p_G,&\ell\in G,\\0,&\ell\notin G.\end{cases}
\] The full classical--quantum (CQ) joint state therefore equals
\(\Omega_G(\rho)\) and is independent of the completion instrument \(J\)
chosen at each node. If \(p_G=0\), every original accepted block is zero
and the conditioned state is undefined.

\textbf{Proof.} Apply P-FCOND-03 repeatedly along the tree, transporting
the filtering relation through the successive levels. At an accepted
leaf, \(S_\ell=I\); at a rejected leaf, \(S_\ell=0\). Normalization is
performed only once, at the start, with the original \(p_G\) as
denominator. For positive inputs, zero total accepted trace implies that
every positive accepted block has zero trace. A positive matrix with
zero trace is zero. Different completions yield the same formula for
each block, hence the same full joint state. \(\square\)

At an intermediate time, the state incorporates the future acceptance
effect \(H_h\). It generally differs from the original experiment's
state conditioned only on the prefix observed so far. Inserting a new
measurement, changing future actions, or changing the acceptance event
changes the problem itself and the effects \(H_h\). This theorem does
not assert that a single fixed conditioned process preserves every
counterfactual intervention simultaneously.

Independence from the completion also requires exact confinement to the
reachable support. If the input leaks into \(I-P_h\), different
completions can produce different records; on an input entirely in that
complement, their TV distance can even be one. Error analysis must
retain both valid complete operations and leakage terms.

\section{66.6 An exact criterion for realization by allowed actions}

Fix a finite allowed action library \(\{\mathcal A_{h,a,o}\}_a\). An
allowed controller chooses a common set of classical randomization
weights \(w_y(a)\) from its observation \(y=q(h)\); it may not depend on
an unobserved node. Each action is itself a complete instrument. The
actual Kraus representation of the randomization is
\(\sqrt{w_y(a)}A_{h,a,o,\nu}\).

Use the following Choi coordinates for a complex-linear map: \[
C(F)_{(i,r),(j,s)}=F(|i\rangle\langle j|)_{r,s}.
\] Matrix units span all complex matrices, so \(C(F)=C(G)\) if and only
if \(F=G\). The vector representing an instrument must retain the Choi
block for every visible \(o\).

\phantomsection\label{p-fcond-05}{} \textbf{P-FCOND-05 (A convex-hull
equivalence preserving outcomes and observation permissions).} Given
complex-linear input restrictions \(R_h\) at the nodes, an allowed
randomization satisfying \[
\sum_a w_{q(h)}(a)\,\mathcal A_{h,a,o}\circ R_h
=\widetilde{\mathcal I}_{h,o}\circ R_h\quad(\forall h,o),
\] exists if and only if, on every observation fibre, the entire target
Choi family lies in the real convex hull of the action Choi families on
that fibre.

\textbf{Proof.} Summing the actual Kraus terms with square-root weights
yields a convex combination of maps, and their Choi coordinates are
linear term by term. Conversely, a finite convex combination supplies
weights and constructs a complete randomized instrument. Injectivity of
the Choi representation then gives the branch identities for every input
matrix. Treating all nodes and labels within one observation fibre as
coordinates of a single vector forces all of them to share the same
weights. This uses Core's general convex-hull theorem for real modules
directly, taking a family of complex matrices as the real module. It
does not identify classical state matrices with general quantum CP maps.
\(\square\)

Three tasks must be distinguished:

\begin{longtable}[]{@{}
  >{\raggedright\arraybackslash}p{(\columnwidth - 2\tabcolsep) * \real{0.3000}}
  >{\raggedright\arraybackslash}p{(\columnwidth - 2\tabcolsep) * \real{0.7000}}@{}}
\toprule\noalign{}
\begin{minipage}[b]{\linewidth}\raggedright
Input restriction and target
\end{minipage} & \begin{minipage}[b]{\linewidth}\raggedright
Exact meaning of the criterion
\end{minipage} \\
\midrule\noalign{}
\endhead
\bottomrule\noalign{}
\endlastfoot
\(R_h=\operatorname{id}\), with a specified complete target instrument &
Reproduce that specified completion on the full space. \\
\(R_h(X)=P_hXP_h\) & Reproduce the conditioned process on all declared
support inputs; the allowed mixture found by the criterion itself
supplies a valid completion on the full space. \\
The actually reachable set for a fixed initial state, or a requirement
only on the final output & A weaker task than matching all support
inputs at every node; it requires a separately specified criterion. \\
\end{longtable}

Failure in the first row does not rule out a different completion.
Failure in the second does not automatically rule out the third.
Matching at each node guarantees realization of the whole protocol, but
is not a universal necessary condition for matching only its final
output. The allowed class in this criterion is one round of classical
randomization over a finite library. The result is not a controllability
theorem for quantum circuits of arbitrary length.

\section{66.7 Positive certificates, separation certificates, and
observation constraints}

\phantomsection\label{p-fcond-06}{} \textbf{P-FCOND-06 (A finite
negative separation certificate for the allowed class).} Suppose a
real-linear functional \(f\) and a real number \(b\) satisfy
\(f(v_a)\le b\) for every allowed action vector \(v_a\), whereas
\(f(z)>b\) for the target \(z\). Then the target is not in that allowed
convex hull.

\textbf{Proof.} Any nonnegative normalized weights give
\(f(\sum_aw_av_a)=\sum_aw_af(v_a)\le b\), contradicting the target
inequality. \(\square\)

An actual positive certificate consists of nonnegative normalized
weights, complete action instruments, and all Choi equalities with
visible labels retained. A status message from an LP solver cannot
replace those equalities. A negative certificate likewise requires
independently checkable separating inequalities.

Two counterexamples are essential. First, deterministically outputting
label 0 and applying the identity channel is identical, after discarding
the label, to deterministically outputting label 1 and applying the
identity channel. Yet the records have TV distance one. Second, suppose
two hidden nodes share the same observation. The target at the first
node requires action 0 with probability one; at the second, it requires
action 1 with probability one. Each node is feasible separately, but no
common randomization works. Checking convex-hull membership only node by
node therefore gives an incorrect conclusion about the physically
permitted controller.

\section{66.8 When initialization can be deterministic}

Consider the map \[
F_H(\rho)=\frac{\sqrt H\rho\sqrt H}{\operatorname{Tr}(H\rho)},\qquad
\operatorname{Tr}(H\rho)>0,
\] and ask whether it admits a deterministic channel valid for all
unknown inputs. Checking a single known input cannot establish such a
channel.

\phantomsection\label{p-fcond-07}{} \textbf{P-FCOND-07 (Classification
of deterministic initialization on all unknown inputs).} Let
\(H\succeq0\), \(H\ne0\) be a finite-dimensional matrix. If only the
system output marginal is required, one deterministic CPTP map realizes
\(F_H\) on every density state with positive success probability if and
only if either \(H=cI\), \(c>0\), or \(\operatorname{rank}H=1\). If the
conditioned joint state with an arbitrary external reference must also
be preserved, only the nonzero scalar case \(H=cI\) remains.

\textbf{Proof.} In the scalar case, \(F_H\) is the identity map. In the
rank-one case, the output is a fixed pure state, which a replacement
channel can prepare. Conversely, if there are two distinct positive
eigenvalues, take the corresponding eigenstate projectors. Each is
unchanged by the filter, but their equal mixture is reweighted according
to the eigenvalues after normalized filtering, violating affinity. If
all positive eigenvalues are equal and \(2\le\operatorname{rank}H<d\),
choose two distinct pure states \(u,v\) in the support and a pure state
\(k\) in the kernel. The inputs \(u\) and \(t v+(1-t)k\), \(0<t<1\),
have success probabilities \(c,ct\) and conditioned outputs \(u,v\),
respectively. Filtering their equal mixture gives conditional weights
\(1/(1+t),t/(1+t)\), again violating affinity. This covers every case
other than the two listed classes.

If a reference system is retained, a local deterministic channel cannot
change its marginal. Classically correlate eigenstate branches of the
system having different success weights with orthogonal reference
labels. Filtering with a nonscalar \(H\) changes those label weights.
Thus, even when the system output is a fixed rank-one state, a local
replacement channel cannot reproduce the full conditioned joint state.
In the scalar case, local normalized filtering is the identity and
preserves all reference correlations. \(\square\)

This result classifies only the specified initialization map \(F_H\). It
does not directly rule out other algorithms required to match only the
final CQ output, nor does rank-one replacement automatically reproduce
the original success probability and filtering flag. The classification
is a paper proof in this chapter. Specific failures of affinity, the
rank-one marginal exception, and counterexamples involving reference
correlations are covered separately by finite checks. If \(\rho\) is
known and can be prepared repeatedly, one may also prepare \(F_H(\rho)\)
directly, but its preparation cost must be recorded. With only one copy
of an unknown state, a failed attempt does not provide a free identical
unknown input for another attempt.

\section{66.9 A valid initial filter and the cost of scaling}

This section retains the assumption of positive acceptance probability.
Choose \(c>0\) with \(c\ge\|H_0\|\), and take the success Kraus operator
\[
A=\sqrt{H_0/c}.
\] Since \(I-A^{\dagger}A\succeq0\), a failure Kraus operator can be
added to form an actual complete instrument.

\phantomsection\label{p-fcond-08}{} \textbf{P-FCOND-08 (Actual success
probability and posterior state of the scaled filter).} This complete
filtering instrument succeeds with probability \(p_G/c\). Its state
after success is \(\widetilde\rho_0\), independently of the scale \(c\).
If independent preparation of the same input is available, the mean
number of filtering attempts is \(c/p_G\).

\textbf{Proof.} The success probability is
\(\operatorname{Tr}(A\rho A^{\dagger})=\operatorname{Tr}(H_0\rho)/c\).
The factors \(1/c\) in the numerator and denominator cancel on
normalization. The waiting time for independent attempts has a geometric
distribution with success probability \(p_G/c\). \(\square\)

Within the implementation family in which this scaling is the only
permitted change to the success Kraus operator, \(c=\|H_0\|\) maximizes
the success probability. This is not a theorem of globally optimal cost
over all implementations. In particular, the scalar or rank-one marginal
tasks in P-FCOND-07 admit deterministic schemes; directly preparing a
known state is another resource option.

\section{66.10 Early stopping on failure must enter the actual cost}

Let \(R\) be the full record of one original attempt, let its success
probability be \(p>0\), and let \(C(R)\) be its actual nonnegative cost,
with mean \(\mu\). The cost includes the actions, preparation, failure
detection, and reset actually carried out during that attempt. Within a
single attempt, cost may be correlated with success or failure.

\phantomsection\label{p-fcond-09}{} \textbf{P-FCOND-09 (Full cost of
independent resets until success).} If each reset produces the same
independent trial, the expected total cost until the first success is
\(\mu/p\). Truncating at no more than \(n\) attempts gives expected cost
\[
\frac{\mu}{p}\bigl[1-(1-p)^n\bigr].
\]

\textbf{Proof.} The first attempt incurs unconditional expected cost
\(\mu\). It fails with probability \(1-p\), after which the new trial is
independent of the preceding outcome. Thus \(V_{n+1}=\mu+(1-p)V_n\),
with \(V_0=0\). Solving this finite recurrence gives the closed form.
Nonnegativity of the costs permits monotone convergence to identify the
limit with the expected total cost. Independence between cost and
success within a single trial is not required. \(\square\)

A comparison budget for the new process can be written as \[
C_{\rm new}=C_{\rm init}+\sum_{h\,\mathrm{active}}
 \widetilde q_h\,C_h+C_{\rm record}+C_{\rm calibration},
\] where \(\widetilde q_h\) is the probability that the actual
conditioned process visits the corresponding node. The term
\(C_{\rm init}\) should include all failed initial filtering attempts
and repeated preparations. Compare \(C_{\rm new}\) with \(\mu/p\) only
after specifying the allowed actions and common resource units. The
conclusion may be that there is no advantage, or that implementation is
impossible. Frequent production of accepted records by the new process
does not increase the success probability of the unchanged original
experiment.

\section{66.11 A sharp error bound for conditioning on acceptance}

Let \(\Omega,\Omega'\) be two full CQ outputs of original processes,
with the same acceptance extractor and positive success probabilities
\(p,q>0\). Write \(D\) for the trace distance, with the one-half
normalization used in this book, equivalently the distance defined by
all effects.

\phantomsection\label{p-fcond-10}{} \textbf{P-FCOND-10 (Propagation of
error from full outputs to conditioned outputs).} If
\(D(\Omega,\Omega')\le\varepsilon\), then \[
|p-q|\le\varepsilon,\qquad
D(\Omega_G,\Omega'_G)\le
\min\left\{1,\frac{\varepsilon}{\max(p,q)}\right\}.
\]

\textbf{Proof.} The acceptance effect gives the probability-difference
bound. For any effect \(E\) on the conditioned output, write its two
conditional probabilities as \(a,b\in[0,1]\). Placing \(E\) and its
complementary effect within the accepted branch and pulling each back to
the original input gives valid effects. Therefore \[
|pa-qb|\le\varepsilon,\quad
|p(1-a)-q(1-b)|\le\varepsilon.
\] If \(p\ge q\), the first inequality and \((p-q)b\ge0\) give
\(p(a-b)\le\varepsilon\); the second inequality and \((p-q)(1-b)\ge0\)
give \(p(b-a)\le\varepsilon\). Hence \(|a-b|\le\varepsilon/p\).
Exchanging the two states yields the denominator \(\max(p,q)\). Take the
supremum over all effects and combine this with the bound of one on the
distance. \(\square\)

This statement applies the same acceptance rule to two full outputs. The
premise \(\varepsilon\) cannot be estimated by retaining only attractive
samples after filtering. A classical three-outcome distribution already
attains the amplification factor. For \(p>0\), \(0\le\delta\le p\le1\),
take the distributions \((p,0,1-p)\) and \((p-\delta,\delta,1-p)\),
accepting the first two outcomes. Their original distance is \(\delta\),
and their conditioned distance is \(\delta/p\). Reporting the
unamplified source error when acceptance is rare can substantially
overstate predictive precision.

\section{66.12 From errors in allowed operations to the full-record
budget}

Chapters 62 and 65 already provide budgets for individual actions under
a common classical controller, weighted by the actual active-prefix
masses of the reference implementation. Here each step must be compared
as a complete CP instrument retaining its visible labels, with the
initial-state error recorded separately.

\phantomsection\label{p-fcond-11}{} \textbf{P-FCOND-11 (A usable local
error bound for valid Kraus approximations).} Consider two actual
complete Kraus families \(\{L_j\},\{M_j\}\) with the same visible labels
and with their hidden indices aligned by zero padding. For every
normalized input, the distance between their outputs with visible labels
retained is bounded by \[
\delta_{\rm step}\le\min\left\{1,
\frac12\sum_j(\|L_j\|+\|M_j\|)\|L_j-M_j\|\right\}.
\] If valid budgets of this form are available at every node of the same
fixed controller, the full-output distance is bounded by the
initial-state error plus the local budgets weighted by the reference
prefix masses.

\textbf{Proof.} Insert and subtract \(M_j\rho L_j^{\dagger}\) in each
Kraus term: \[
L_j\rho L_j^{\dagger}-M_j\rho M_j^{\dagger}
=(L_j-M_j)\rho L_j^{\dagger}
+M_j\rho(L_j-M_j)^{\dagger}.
\] Apply \(\|AXB\|_1\le\|A\|\|X\|_1\|B\|\), the triangle inequality, and
\(\|\rho\|_1=1\). Sum over the hidden indices, then over the visible
direct-sum blocks, and divide by two. The additional bound of one is
valid because both outputs are normalized. Propagation through the
finite tree follows either from the existing common-controller budgets
or from replacement one level at a time and CPTP contraction.
\(\square\)

The same operator-norm derivation remains valid when each local
operation is extended by the identity on an explicitly retained
reference system. It does not control unmodelled additional couplings or
physical errors that change the controller. This Kraus-norm estimate is
a paper-level sufficient bound in the present chapter, not a definition
of the optimal distance. Numerical implementations must provide evidence
for their matrix-error and norm upper bounds; machine residuals cannot
simply be treated as rigorous real-number bounds. Nor may an incomplete
floating-point Kraus family be treated as an actual CPTP action without
a valid correction.

\section{66.13 Small eigenvalues, exact certificates, and numerical
failure}

For two certified positive-definite effects \(H,H'\succeq\eta I>0\),
standard sufficient bounds for the square root and inverse square root
are \[
\|\sqrt H-\sqrt{H'}\|\le\frac{\|H-H'\|}{2\sqrt\eta},\qquad
\|H^{-1/2}-{H'}^{-1/2}\|\le
\frac{\|H-H'\|}{2\eta^{3/2}}.
\] The first follows from the Sylvester equation for the difference of
the square roots and its semigroup integral; the second additionally
uses the identity for the difference of inverse matrices. Substitution
into the product difference for \(S_{ho}KT_h\) yields local Kraus
errors. In the singular case, the support must be fixed and certified
before applying the corresponding lower bound on that support.
Floating-point sorting alone cannot determine the exact rank.

\phantomsection\label{c-fcond-01}{} \textbf{C-FCOND-01 (Small positive
eigenvalues cannot be treated as zero).} Let
\(H_\epsilon=\operatorname{diag}(1,\epsilon)\) and
\(H_0=\operatorname{diag}(1,0)\), with \(\epsilon>0\). Although
\(\|H_\epsilon-H_0\|=\epsilon\to0\), the distance between its inverse
square root on the support and that of \(H_0\) is \(\epsilon^{-1/2}\),
which diverges.

\textbf{Proof.} The two inverse square roots on the supports are
\(\operatorname{diag}(1,\epsilon^{-1/2})\) and
\(\operatorname{diag}(1,0)\). Subtract them and take the operator norm.
\(\square\)

When the smallest positive support eigenvalue is too small relative to
the largest positive eigenvalue and reaches the current
inverse-square-root precision threshold, this chapter's reference
program explicitly reports unreliable numerical conditioning. It does
not convert a positive-probability event into a zero event merely
because its overall scale is small. Rigorous finite examples use
rational or algebraic matrices and checks of the actual backward Gram
sums and support identities. Floating-point completeness and recordwise
residuals are recorded separately as numerical diagnostics. Exact finite
SymPy computations and Lean kernel checks are also recorded separately,
without treating either as the other.

\section{66.14 A realizable positive example connecting the whole chain}

On a two-level system, take the Pauli matrices \(Z,X\). The original
visible outcomes are \(0,1,f\), with hidden Kraus operators \[
K_0=Z/2,\quad K_1=X/2,\quad
K_{f,0}=I/2,\quad K_{f,1}=I/2.
\] The complete Gram sum is \(I\). The two hidden paths for \(f\) retain
the same visible failure label. Accept a three-step run if every outcome
lies in \(\{0,1\}\). Then \[
H_t=2^{-(3-t)}I,\quad p_G=1/8,\quad\widetilde\rho_0=\rho.
\] At each step, the conditioned instrument outputs 0 and applies \(Z\)
with probability one half, and outputs 1 and applies \(X\) with
probability one half. Choose the allowed library to consist of exactly
these two unitary actions with deterministic labels. The weights
\((1/2,1/2)\) give an actual positive Choi certificate, matching the
quantum blocks of all eight successful records simultaneously. Each
record retains the full quantum state after its corresponding unitary
evolution. Although \(XZ=-ZX\), the two conjugation channels commute
because the global phase cancels. This example therefore does not test
noncommuting CP superoperators.

Suppose one call to the original complete instrument and one call to an
allowed Pauli action each cost one unit, with no additional preparation
or recording charge. The mean action cost of one original attempt with
early stopping is \(1+1/2+1/4=7/4\), giving cost 14 per successful
sample. The new process requires three actions, or \(3+3r\) if each
randomization costs \(r\ge0\). Under this explicit resource convention,
an advantage requires \(3+3r<14\). Any conclusion about real devices
additionally requires measured operation, randomization, preparation,
calibration, and reset costs. No hardware data are supplied here.

\section{66.15 A nonscalar, noncommuting example and implementation
obstructions}

In another original instrument, the two success Kraus operators are \[
K_0=\begin{pmatrix}1/2&0\\0&1/4\end{pmatrix},\qquad
K_1=\frac14\begin{pmatrix}1&1\\1&-1\end{pmatrix}.
\] The failure label contains the two hidden Kraus operators
\(\operatorname{diag}(3/4,3/4)\) and \(\operatorname{diag}(1/4,1/2)\),
so the instrument is again exactly complete. For
\(\rho=\frac12\begin{pmatrix}1&-i\\i&1\end{pmatrix}\), the backward
effect and probability of three consecutive successes are \[
H_0=\frac1{4096}\begin{pmatrix}183&27\\27&51\end{pmatrix},\qquad
p_G=\frac{117}{4096}.
\] All eight visible success words retain their full complex quantum
blocks. With \(P_0=\operatorname{diag}(1,0)\), reversing the order of
the two success CP branches gives the difference \[
\mathcal I_0(\mathcal I_1(P_0))-\mathcal I_1(\mathcal I_0(P_0))
=\begin{pmatrix}0&-1/128\\-1/128&-3/256\end{pmatrix}\ne0,
\] so this example tests genuinely noncommuting branch maps. One
original attempt has action cost \(701/512\); the expected cost until
success is \(5608/117\), approximately 47.93. The best valid scale for
this single-Kraus filtering family,
\(c=\|H_0\|=(117+3\sqrt{565})/4096\), gives initial filtering success
probability \(39/(39+\sqrt{565})\). The original \(117/4096\) is the
filtering success probability only when the unscaled choice \(c=1\) is
used.

This conditioned process does not automatically belong to the original
action library. If the library contains only the original instrument,
its failure effect has probability \(23/32\) on the maximally mixed
state. At a positive-definite node, the target conditioned instrument
instead has zero failure probability. The negative of the failure
probability is a real-linear separating functional that excludes exact
realization by this library. Furthermore, the original failure effect is
\(\operatorname{diag}(5/8,13/16)\succeq(5/8)I\), so the original
instrument still has failure probability at least \(5/8\) on the actual
filtered initial state. In this particular case, that bound also
directly excludes matching the full record for the fixed initial state.
It does not exclude broader implementation classes containing additional
filters or new actions. A formally complete Doob instrument must not be
assumed to be available on a device.

Singular examples, zero events, alternative hidden Kraus decompositions,
conflicts within an observation fibre, leakage into the support
complement, rare-event error amplification, and changes of controller
are retained separately in the corresponding execution records. They are
mathematical and program controls, not experimental support from nature.

\section{66.16 What this chapter establishes}

The original experimental tree determines the backward effects. Spectral
support yields valid conditioned instruments. Matrix identities at every
leaf give the full conditioned CQ law. Choi families and observation
fibres give a necessary and sufficient criterion for allowed
randomization. Initialization and early stopping determine the resource
account. Errors in valid complete operations, together with the
acceptance extractor, determine the final error. Each step has its own
input objects, proof, or counterexample. The chain may terminate with a
precise negative answer: the specified library cannot realize the
process.

The formal sources are \texttt{Quantum\allowbreak Finite\allowbreak Window},
\texttt{Quantum\allowbreak Finite\allowbreak Window\allowbreak Spectral},
\texttt{Quantum\allowbreak Action\allowbreak Realizability},
\texttt{Quantum\allowbreak Reachable\allowbreak Action\allowbreak Realizability},
\texttt{Quantum\allowbreak Postselection\allowbreak Error}, and
\texttt{Quantum\allowbreak Initialization\allowbreak Cost}, together with explicit action and
Pauli witnesses. The current Lean tree uses uniform finite carriers for
the system, outcomes, and Kraus indices, and one fixed arbitrary
complete fallback instrument \(J\). If the root contains an action, its
original instrument can itself be chosen as the fallback. Equality of
the joint states has already been proved for two fixed fallbacks with
different hidden carriers. The node-dependent \(J_h\) allowed in the
text is a paper generalization of the same local transport identity.
Zero padding between different node carriers still requires an explicit
adapter.

The full initialization classification, the impossibility argument for
arbitrary references, the sufficient Kraus operator-norm error bound,
and the monotone-convergence interpretation of the expected cost of
independent retries until the first success are paper proofs. Their
verification status does not change because neighbouring modules pass
Lean. Actual builds, named axioms, execution records, and source
identities are documented in this edition's verification ledger.

Finiteness, an explicitly chosen complex quantum domain, a fixed
original controller, randomization over a finite action library, and the
resource assumptions of independent resetting remain substantive
premises. This chapter does not derive the quantum domain uniquely from
Core, prove a general quantum Perron theorem or an unbounded conditioned
path construction, or provide new hardware data. Future experiments can
test whether a specified implementation meets the full-record, cost, and
error contracts established here.

\chapter{67. Operational identity, microscopic truncation, and robust
implementation with shared parameters}\label{ch-67}

This chapter continues the finite conditioning chain already closed in
Chapter 66. The new responsibility is to connect the statistical
identity of actual transformations, oscillator dynamics, finite
computation, and one legal deployment under unknown parameters to the
same complete output. Generic Core, selection of the complex sector, a
physical device, and empirical support remain separate layers. No frozen
Core statement is changed, and internal consistency of standard quantum
structure is not presented as an independent selection of nature's
theory.

\section{67.1 Inheritance, quantifiers, and evidence}

Purification, local tomography, effect separation, and trace distance in
Chapters 54--64, together with Chapter 66's complete conditional CQ
chain, are inherited with their original sources and receipts. The
research baseline is \texttt{16a89154183c\allowbreak 22030d452641\allowbreak 3167d095d94e\allowbreak 6890};
at entry, all 155 historical proof/configuration inputs and nine runtime
inputs matched their historical receipts byte for byte; the new public
imports and audit configuration have separate current identities.
Analytic proof, Lean, exact finite arithmetic, floating diagnostics,
synthetic records, and real measurements are registered separately.

The dependency classes here are B, P, E, and O. The selected complex
quantum structure is an explicit physical input; this chapter has not
reduced that independent premise. Implementation error concerns complete
CQ outputs with the same record alphabet and terminal space; §67.10
separately studies a discarded-record marginal distance. Records include
failure, stopping, and preparation-filter failure; the reference is
untouched. Universal state/effect quantifiers do not assert that an
experimentalist has acquired all those measurement capabilities.

\section{67.2 Operational identity of actual events}

\phantomsection\label{p-opt-02}{}

\textbf{P-OPT-02 (Four equivalent identities for events in
positive-dimensional finite complex systems).} Let the input and output
be positive-dimensional finite complex Hilbert spaces, and let events be
complex-linear, completely positive, trace-nonincreasing maps. The
following are equivalent: equality of raw statistics in every finite
positive-dimensional auxiliary context; equality on every single-system
density and effect; equality of the actual complex-linear maps; and
equality of their fixed-basis unnormalized Choi matrices.

\[
\begin{gathered}
\forall R,\rho_{AR},\ 0\le E_{BR}\le I:\quad
\operatorname{Tr}\!\left[E_{BR}(\Phi\otimes\mathrm{id}_R)(\rho_{AR})\right]
=\operatorname{Tr}\!\left[E_{BR}(\Psi\otimes\mathrm{id}_R)(\rho_{AR})\right],\\
J(\Phi)_{(i,a),(j,b)}=\Phi(|i\rangle\langle j|)_{a,b}.
\end{gathered}
\]

\textbf{Proof.} A one-dimensional auxiliary gives the single-system
condition. For each fixed input, the two event outputs are Hermitian and
may have unequal traces. A nonzero difference is distinguished by its
positive or negative spectral projector, an allowed effect. Hence the
outputs coincide on every density. Every nonzero PSD input is its trace
times a density, and a zero-trace PSD input is zero. Every Hermitian
matrix is a difference of PSD matrices; every complex matrix is a
Hermitian real part plus the imaginary unit times a Hermitian imaginary
part. Complex linearity therefore determines the whole map. Equality
survives tensoring with an identity, including entangled inputs. Each
input block of the Choi matrix recovers one matrix unit's image, and
matrix units form a complex basis, proving the remaining equivalence.
\(\square\)

The zero event has no posterior but is included. Normalizing an event
first loses its probability: one quarter of the identity event and one
half of it have identical positive-branch posteriors but are different
events. Restricted effects or inputs do not satisfy the universal
premise. Approximate Choi entries cannot simply be renamed a worst-case
operational error.

\section{67.3 Visible labels and the operational quotient}

\phantomsection\label{p-opt-03}{}

\textbf{P-OPT-03 (Composition respects identity at every visible
label).} Fix a visible outcome alphabet. Two complete instruments have
identical labelled auxiliary statistics exactly when their actual CP
maps agree at every visible label. Sequential and parallel composition,
specified outcome coarse graining, classical randomization, conditional
tests, and bounded classical adaptive trees preserve this equivalence.

The state carrier used for composition must contain all persistent
memory, or auxiliaries must be independently reset as in this chapter's
example. Equality of single-slot channels does not identify multitime
hardware with different reused hidden environments.

\textbf{Proof.} Each diagonal record block of the complete CQ output is
its unnormalized event output; apply P-OPT-02 to each block.
Composition, tensor product, and finite sums preserve map equality.
Induction on finite tree height proves adaptive substitution.
Conditional completeness follows from trace preservation of each
successor test's sum and completeness of its parent instrument. Control
weights must agree within one observation fibre; hidden Kraus indices
are unavailable. \(\square\)

One identity Kraus operator and two identity Kraus operators scaled by
three fifths and four fifths represent the same CP map. Swapping visible
projective outcomes does not represent the same instrument: on a
suitable input its complete CQ distance is one even though the forgotten
channels agree. A rectangular event can be completed by mapping the
probability of its missing effect to a fixed output density. Positive
output dimension ensures that such a density exists. Empty matrix
carriers remain mathematical boundary tests, rather than physical
objects without normalized states.

This supplies operational semantics inside the chosen complex model. The
Chiribella--D'Ariano--Perinotti reconstruction still requires the
complete finite causal OPT background and its selection principles; see
the \href{https://arxiv.org/abs/1011.6451v3}{2011 original} and
\href{https://arxiv.org/abs/0908.1583v5}{2010 framework}. Neutral
principles of parent-state existence, complete access, and additive
resources also hold in full classical probability theory: use a product
parent distribution, all stochastic maps, and additive nonnegative
costs. These principles do not exclude the classical candidate, and the
complete reconstruction has not been machine-checked here.

\section{67.4 Deriving a uniform truncation from microscopic dynamics}

The system is a two-level probe and one oscillator; its reference is any
finite positive-dimensional system. Excitation number and its budget are
dimensionless. Physical energy requires a separate operator relation to
excitation.

\[
\mathcal H_S=\mathbb C^2\otimes\ell^2(\mathbb N_0),\qquad
G=a^\dagger a+P_e,\qquad
\Pi_N=1_{G\le N},\qquad \operatorname{rank}\Pi_N=2N+1.
\]

Jaynes--Cummings exchange, probe phase, and excitation-preserving free
terms form finite Hermitian blocks in each total-excitation sector.
Their Hilbert direct sum is self-adjoint on vectors for which the
squared norms of all block actions sum finitely: block resolvents are
uniformly bounded, yielding self-adjointness and the unitary group.
These controls strongly preserve the spectral projections. When this
Hamiltonian is used as a device model, rotating-wave and other
approximation errors remain separate obligations.

A vacuum auxiliary with nonnegative excitation ensures that an
excitation-conserving joint unitary, started below the cutoff, leaves
every hidden branch of a complete readout within the system cutoff. The
readout need not be diagonal in auxiliary excitation: its component
along any fixed readout vector cannot create higher system excitation.

\phantomsection\label{p-micro-01}{}

\textbf{P-MICRO-01 (A microscopic bridge to complete adaptive CQ
truncation).} For this actual action family, finite outcome/hidden
refinements, fixed observation permissions, and a maximum finite height,
every legal classical controller, finite reference, and
budget-admissible joint input satisfies:

\[
\begin{gathered}
K_{aok}\Pi_N=\Pi_NK_{aok}\Pi_N,\qquad
s_N=\operatorname{Tr}[(\Pi_N\otimes I_R)\rho_{SR}]>0,\\
\rho_N=\frac{(V_N^\dagger\otimes I_R)\rho_{SR}(V_N\otimes I_R)}{s_N},\qquad
\operatorname{Tr}(G\rho_S)\le\bar n<N+1,\\
D\!\left(\Omega^\pi(\rho_{SR}),\iota_N\Omega_N^\pi(\rho_N)\right)
\le\sqrt{\frac{\bar n}{N+1}}.
\end{gathered}
\]

\textbf{Proof.} The first identity is derived from microscopic
conservation, auxiliary preparation, and readout. The spectral
inequality bounds the tail mass, and the purification argument of
P-QINF-02 bounds the initial normalized projected state. Define finite
branches by compressing the actual Kraus operators. Invariance derives
branch intertwining and finite completeness; intertwining is not
supplied as a certificate field. Multiplying and summing along the tree,
while retaining all terminal labels, derives complete-output
intertwining for the same controller. A complete finite tree is normal
CPTP and contracts Hermitian trace-class differences, proving the
estimate. \(\square\)

There is no step-count multiplier on the tail error because the subspace
remains invariant. Excitation injection, leakage, and unbounded stopping
are outside the proof. Normalized projection depends on the input; the
expression is not an energy-constrained diamond norm of a difference
between two fixed channels.

\section{67.5 Necessary premises, normal trace class, and actual
filtering}

\phantomsection\label{c-micro-01}{}

\textbf{C-MICRO-01 (Mean-energy preservation does not imply branchwise
truncation).} On three levels with energies zero, one, and two, a
measure-and-prepare channel fixes levels zero and two, sends level one
to an equal mixture of zero and two, and removes off-diagonal entries.
It preserves mean excitation for every input. One Kraus operator
nevertheless raises level one to level two, escaping the first cutoff.
Mean-energy nonincrease cannot replace P-MICRO-01's branchwise premise.

A separate actual-operation counterexample exchanges one probe
excitation into the oscillator and then completely damps it, producing
vacuum. Reversing those operations produces an excited oscillator. The
CP compositions have distance one on the same input. This checks actual
channels; anticommutation of two Kraus matrices would not suffice.

\phantomsection\label{p-micro-02}{}

\textbf{P-MICRO-02 (Normal trace-class postselection and complete
preparation records).} For the same normal CP trace-nonincreasing
acceptance extractor, complete normal outputs, and positive acceptance
probabilities, the sharp finite postselection constant extends to
trace-class outputs. For a physical projection filter, the ideal
always-pass output and actual complete pass/fail output also satisfy the
following preparation bound:

\[
\begin{gathered}
D(\omega,\sigma)\le\varepsilon,\quad
A=\mathcal E(\omega),\ B=\mathcal E(\sigma),\quad
p=\operatorname{Tr}A>0,\ q=\operatorname{Tr}B>0,\\
|p-q|\le\varepsilon,\qquad
D(A/p,B/q)\le\min\{1,\varepsilon/\max(p,q)\},\\
\delta_N=1-s_N,\qquad
D_{\mathrm{prep}}\le
\frac{\sqrt{\delta_N(4-3\delta_N)}+\delta_N}{2}
\le\min\{1,\sqrt{\delta_N}+\delta_N\}.
\end{gathered}
\]

\textbf{Proof.} Complete the extractor using its missing effect and an
orthogonal failure flag. Contractivity bounds the accepted-block
trace-norm difference plus the success-probability difference by twice
the original error. If the first probability is larger, write the
normalized difference as the accepted-block difference divided by that
probability plus the second positive block times the difference in
normalization factors. The triangle inequality proves the sharp
estimate; exchange the outputs in the other case. Only trace-class norms
and normal maps are used.

For preparation, purify the input and separate its orthogonal pass and
fail components. In the pass block, ideal minus filtered input is a
two-dimensional Hermitian matrix with diagonal entries zero and the tail
mass, and off-diagonal entries equal to the product of the two
amplitudes. Its trace norm is the displayed square root. Add the
orthogonal fail block's tail mass and divide by two. Partial trace and
the common subsequent channel contract the result. \(\square\)

The common record space must include a preparation-failure label and
pass-prefixed original records. For the low-spectral filter used here,
the passing state does not increase the mean excitation budget, as
proved in Section 68.4; general postselection can increase a conditional
mean. Failed filtering of an unknown state supplies no fresh copy. The
trace-class extension and direct-sum self-adjointness are analytic
proofs; a nearby finite-matrix Lean compilation does not extend their
machine-checked domain.

\section{67.6 One source with shared parameters}

The frozen unknown parameter consists of the source flag probability,
phase pulse, and exchange pulse. The last two may be uncertain, but
source, target, and action library must share the same actual gates. Two
original steps stop at the first failure.

\[
\begin{gathered}
\theta=(q,\phi,\alpha),\quad q\in[12/25,13/25],\quad s=1-q/2,\\
U_0=e^{-i\phi P_e},\qquad
U_1=e^{-i\alpha(a\sigma_++a^\dagger\sigma_-)},\\
K_0=\sqrt{s(1-q)}U_0,\quad
K_1=\sqrt{sq}U_1,\quad K_f=\sqrt{1-s}I,\\
\mathcal R=\{f,0f,1f,00,01,10,11\}.
\end{gathered}
\]

The microscopic realization uses a fresh independent three-state
degenerate logical flag with zero excitation operator. Two real plane
rotations prepare the three amplitudes; flag-controlled gates apply
either actual gate or the identity, and the flag is completely read out.
The joint process preserves system excitation. Logical degeneracy does
not make preparation, readout, or reset free. The original two-success
probability and backward effects are derived from this source:

\[
H_t=s^{2-t}I,\qquad p_{\mathrm{src}}=s^2\ge1369/2500,
\qquad t\in\{0,1,2\}.
\]

The backward-effect formula applies only to successful prefixes; a
failure leaf has zero acceptance effect. The conditional target has
independent Bernoulli probabilities on its four successful histories,
zero blocks on its three original failure histories, and terminal states
obtained by the respective ordered gate products on the original joint
input. Its normalization is input-independent, making it CPTP.
Consequently this example's cutoff comparison contracts directly,
without division by the source success probability. The complete raw
source still fails and must not be identified with its conditioned
target.

\section{67.7 Joint calibration and robust certificates}

\phantomsection\label{p-rob-01}{}

\textbf{P-ROB-01 (Deployment error and separation on one parameter
set).} Fix complete outputs, input/reference scope, and observation
permissions. Calibration produces a joint set. Given a deployment
candidate \(w(D_{\mathrm{cal}})\) depending only on calibration and
legal visible history, require \(w_a\ge0\) and \(\sum_a w_a=1\), with
its nominal residual certified. If target and actions have uniform
errors relative to nominal objects on that common set, then for this
given candidate:

\[
\begin{gathered}
d_{\mathcal X}(T_\theta,\widehat T)\le\eta_T,\quad
d_{\mathcal X}(A_{a,\theta},\widehat A_a)\le\eta_a,\quad
d_{\mathcal X}(\widehat T,\sum_a w_a\widehat A_a)\le r_{\mathrm{nom}},\\
\forall\theta\in C(D_{\mathrm{cal}}),\\
d_{\mathcal X}\!\left(T_\theta,\sum_a w_aA_{a,\theta}\right)
\le r_{\mathrm{nom}}+\eta_T+\sum_a w_a\eta_a.
\end{gathered}
\]

Deployment design requires
\(\exists w(D_{\mathrm{cal}})\,\forall\theta\in C(D_{\mathrm{cal}})\).
This sufficient bound does not itself guarantee a small-residual
candidate.

\textbf{Proof.} For the same input and parameter, insert the nominal
target and nominal mixture between the actual target and mixture. Apply
the triangle inequality and convexity of the norm, then take the
declared input and parameter suprema. Convex mixing preserves instrument
completeness. History-dependent policies may be treated as members of a
complete policy library, or certified fibre by fibre and then composed
using a separate finite error-propagation theorem. That composition step
cannot be omitted. \(\square\)

For a fixed legal input and a complete-record effect, a certified target
expectation lower bound exceeding all action expectation upper bounds
plus the allowed error yields a separation certificate for every
candidate and every mixture. A general linear functional requires its
actual dual norm. If neither upper nor lower bounds decide the requested
precision, report insufficient certification, rather than identifying
solver failure with physical impossibility.

\phantomsection\label{c-rob-01}{}

\textbf{C-ROB-01 (Pointwise realizability does not imply one robust
deployment).} For this chapter's conditional target, each known
parameter permits exact realization by setting each step's label-one
weight to that parameter. When the parameter is unknown and visible
records are generated by the chosen actions, with no additional
parameter-sensitive query, the deployed first-label distribution is
fixed. Targets at two different parameters have different first-label
probabilities and cannot both be reproduced exactly. This establishes
absence of one uniform policy, not impossibility for each physical
candidate.

Forgetting labels creates a different false positive: two instruments
can both have the identity as forgotten channel while always emitting
zero and always emitting one, respectively. Their complete record
distance is one. Changing observation permissions changes the feasible
set; hidden nodes must not be disclosed to the controller. Certified
support and spectral-gap premises likewise cannot be replaced by
floating thresholds.

\section{67.8 Exact benefit of visible history and the complete error}

\phantomsection\label{p-rob-02}{}

\textbf{P-ROB-02 (An exact design-region certificate for the same
source).} At each step the library implements an actual gate and emits
its label. The random coin result is that visible label; no additional
action record is deleted. On a design interval of half-width one
fiftieth, constant half weights have two-step worst complete error
fifty-one over twenty-five hundred. A legal policy using the first
visible label improves it to one fiftieth, attaining the minimax value
on the full design region for this labelled-action library.

\[
\begin{gathered}
\Delta=1/50,\qquad
w_1=1/2,\quad w_{2|0}=312/625,\quad w_{2|1}=313/625,\\
Q_{00}=Q_{11}=1/4+\Delta^2,\qquad
Q_{01}=Q_{10}=1/4-\Delta^2,\\
t=q-1/2,\quad u=\Delta^2-t^2,\qquad
D_{\mathrm{CQ}}=\max\{|t|,u\}+u\le\Delta.
\end{gathered}
\]

\textbf{Proof.} The target and implementation use the same actual gates
along each history, so the corresponding terminal states coincide. Each
block's trace-norm difference is therefore exactly its probability
difference, for every input and finite reference. Substitution gives the
formula. If the absolute offset is at least the remainder, the
difference between the half-width and the error factors into two
nonnegative terms; otherwise the error is at most twice the half-width
squared. The endpoint target record laws have TV distance twice the
half-width. By the triangle inequality, every parameter-independent
record law has worst error at least the half-width. This policy attains
it. \(\square\)

Optimality concerns the entire prefrozen design interval and the
specified library without additional parameter-sensitive in-run queries.
A smaller realized calibration set can admit a better policy. At the
central parameter this policy sacrifices an error of \(1/1250\) to
reduce the worst-case error; it is not better than constant weights at
every parameter value.

For a general input of excitation budget one and cutoff four thousand
ninety-five, retaining preparation-failure records gives the following
analytic sufficient bound within the previously frozen precision one
twentieth:

\[
\begin{gathered}
\delta_N\le1/4096,\\
D_{\mathrm{full}}\le65/4096+1/50
=3673/102400\\
=0.035869140625<1/20,\\
\varepsilon_{\mathrm{remaining}}=1/20-3673/102400=1447/102400.
\end{gathered}
\]

The remaining allowance can accommodate independently proved
common-output bounds for actual gate mismatch or unmodelled dynamics.
Without such measurements or model bounds, there is no device-precision
certificate. No thousands-dimensional matrix allocation is offered as a
substitute for the universal proof. The actual finite computation uses a
state known to lie in the single-excitation sector, with an entangled
reference, and checks source, backward effect, every record, and ordered
gates; cutoff and preparation-failure errors vanish in that calculation.

\section{67.9 Calibration, holdout, costs, and mechanism identification}

The model calibration retains the order of eighty thousand single-step
source flags and estimates the source parameter from failure frequency.
Independence, identical distribution, a fixed parameter, and the correct
source relation are coverage premises.

\[
\begin{gathered}
\widehat q=2k_f/n,\qquad n=80000,\\
C_q=[\widehat q-1/100,\widehat q+1/100]\cap[0,1],\\
\Pr(q\notin C_q)\le2e^{-4}<1/20.
\end{gathered}
\]

The executed synthetic run yielded 19844 failures, estimate 0.4961, and
interval {[}0.4861,0.5061{]}, inside the design region. Flag data do not
calibrate gate mismatch. The full joint set also contains the allowed
phase/exchange nuisance domain and an independently justified
shared-gate contract. Two separate held-out sets of 4096 trials test the
seven-label source law and compiled-policy records; neither is rejected
by the frozen finite record testers. These data are model-generated,
without terminal-state tomography, and are neither empirical closure nor
a measured full trace distance.

A negative control repeats one latent flag through all trials. It has
the same one-shot marginal but destroys the stated coverage, showing
that marginals alone do not support iid concentration. Real analysis
requires independent calibration, ordered raw records, all
failures/leakage/stopping, input/reference preparations, gate-matching
evidence, terminal tester settings, and complete costs. Missing items
retain their unverified responsibilities.

If the source is genuinely independently re-preparable, one attempt has
expected action count one plus the one-step success probability. The
action cost per successful output divides by two-step success. With a
common comprehensive per-step cost, the full expectation is:

\[
C_{\mathrm{src}}=
\frac{c_{\mathrm{prep}}+(1+s)c_{\mathrm{step}}+(1-s^2)c_{\mathrm{reset}}+c_{\mathrm{other}}}{s^2}.
\]

Here the additional term is an expected cost per original-source
attempt. It includes auxiliary preparation, readout, recording, and
classical computation not already counted in a step or preparation term.
Unequal leaf costs require the actual source-leaf weighting. A one-time
calibration batch cost \(C_{\mathrm{cal}}\) instead contributes
\(C_{\mathrm{cal}}/M\) when amortized over \(M\) delivered successful
outputs; a cost already amortized per output must not enter the
numerator and be divided by the success probability again. The compiled
policy likewise pays for filter failures, two gates, randomization,
reset, and storage. Unknown quantum inputs do not permit free
independent retries. No measured advantage in seconds, joules, or
currency is claimed. The microscopic process already has an ordinary
quantum-environment realization and does not itself distinguish
objective actualization from standard environmental mechanisms.

\section{67.10 Long-time marginal mixing versus complete memory}

\phantomsection\label{p-micro-03}{}

\textbf{P-MICRO-03 (Single-excitation marginal mixing without
full-record forgetting).} Further fix the single-excitation basis order
as oscillator excitation followed by probe excitation, and choose the
actual gates below. Uniformly over the source parameter's design
interval, the record-forgotten conditional channel contracts every two
steps and has the sector's maximally mixed state as its unique invariant
density. When the complete successful history is retained, the CQ
distance between any two inputs is instead exactly preserved.

\[
\begin{gathered}
U_0=\operatorname{diag}(1,(3+4i)/5),\qquad U_1=(3I-4iX)/5,\\
M(q)=(1-q)
\begin{pmatrix}3/5&-4/5&0\\4/5&3/5&0\\0&0&1\end{pmatrix}
+q\begin{pmatrix}1&0&0\\0&-7/25&-24/25\\0&24/25&-7/25\end{pmatrix},\\
\|M(q)^2\|_2\le49/50,\qquad
D(\mathcal T_q^n(\rho),I_2/2)
\le(49/50)^{\lfloor n/2\rfloor}D(\rho,I_2/2),\\
D(\Omega_n(\rho_{SR}),\Omega_n(\sigma_{SR}))
=D(\rho_{SR},\sigma_{SR}).
\end{gathered}
\]

\textbf{Proof.} Pauli conjugation by the actual two-dimensional gates
gives the displayed real matrices. At the central parameter, nine tenths
squared times the identity minus the Gram matrix of the two-step matrix
has leading principal minors 4941/12500, 422733249/3906250000, and
29552508909/390625000000, all strictly positive. Thus the central
two-step norm is less than nine tenths. A convex combination of
rotations has norm at most one; changing the parameter changes the
squared matrix by at most four times its offset. This yields the uniform
bound. Bloch-vector distance proves marginal contraction and uniqueness.

In the full record, each block is a unitary conjugate of the input
difference times the same history probability. Trace norm is unitarily
invariant, and the probabilities sum to one, proving the last identity
with finite references. The marginal contraction factor does not thereby
become a bound in an arbitrary-reference channel norm. \(\square\)

The entire oscillator retains invariant excitation weights: vacuum and
the single-excitation maximally mixed state are distinct invariant
states. There is no corresponding global oscillator uniqueness or
mixing. The original source's infinite-success probability tends to
zero; a Doob stationary state, a time-conditioned state limit, and a
path law conditioned on an original zero-probability infinite event are
distinct. This result identifies precisely what a complete record
retains as memory and how discarding that record allows finite-system
mixing.

\section{67.11 Physical arrows that remain open}

The continuing task graph remains in the current ROADMAP. The next
microscopic problem is a full-record error for weak nonconserving terms
or deletion of finite bath modes under the same intervention slots.
Gaussian/regulated free fields, complete same-order EFT observables, and
continuum limits follow only with their additional controls. The
particle line retains both three-colour and five-colour candidates.
Rank-one FCNC structure must connect to a common-parameter low-energy
relation not used for fitting, with the full operator running, matrix
elements, and covariance.

The original direct Fisher--Lorentzian identification has a precise
obstruction: a regular real-parameter Fisher matrix is an expected score
outer product, hence PSD, and its real coordinate pullback remains PSD.
Reparameterization alone cannot produce a nondegenerate Lorentzian
metric. This rejects only direct identification or real pullback; other
information-gravity bridges must supply independent causal or signed
structure. A globally hyperbolic background, free field, Hadamard state,
and local probe response form a separately testable next target;
gravitational feedback additionally requires renormalization and
stress-energy conservation.

\section{67.12 Actual evidence limits}

The new operational-equivalence source, compilation, and named axiom
audit are bound individually by the v31.11 receipts. For P-OPT-02 the
machine endpoints concern equality of actual maps, single-system and
ancillary statistics, and Choi matrices. For P-OPT-03 machine coverage
is limited to labelled sequential composition and sum congruence.
Complete parallel-test bundling, quotient-tree induction, all abstract
OPT coherence laws, and reconstruction retain their paper-level
responsibilities. Infinite-dimensional dynamics, normal trace class,
robust bounds, and mixing retain their separate analytic scopes. Exact
finite checkers and the root agent's independent calculation complement
one another. Synthetic calibration and finite-tester holdout do not
replace actual data. Three native research subagents performed internal
AI cross-review, not external human peer review. Complete source
identities, corrections after failed checks, and uncovered domains
remain in this iteration's audit and execution records.

\chapter{68. Segmented excitation growth and controlled blue-pulse
omission}\label{ch-68}

Chapter 67 established finite truncation for excitation-preserving
operations. This chapter permits one specific form of excitation growth:
an independently applied blue-sideband pulse follows the original
red-sideband pulse. The essential step is to derive a support increase
of at most two from the actual generator. Finite carriers then grow with
the steps and preserve complete predictions. A separate small-pulse
premise supports a bound for omitting the blue pulse. This is a
specified nonconserving extension; it does not cover general Rabi
dynamics, infinite time, or all continuous fields.

\section{68.1 The common model and additional inputs}

Retain the probe, oscillator, excitation observable, any finite
positive-dimensional reference, and a bounded initial mean. The
reference is untouched and may initially be entangled with the system.
The parameter is now a four-tuple. Phase, red, and blue areas are
dimensionless integrated couplings. The order below is red first and
blue second; separated operations are not replaced by simultaneous
evolution.

\[
\begin{gathered}
\mathcal H_S=\mathbb C^2\otimes\ell^2(\mathbb N_0),\quad
G=a^\dagger a+P_e,\quad \Pi_M=1_{G\le M},\quad
\theta=(q,\phi,\alpha,\beta),\\
A=\sigma_+a+\sigma_-a^\dagger,\quad
B=\sigma_+a^\dagger+\sigma_-a,\quad
U_0=e^{-i\phi P_e},\quad U_1=e^{-i\alpha A},\\
B_\beta=e^{-i\beta B},\quad V_0=U_0,\quad V_1=B_\beta U_1,\quad
q\in[12/25,13/25],\quad s=1-q/2.
\end{gathered}
\]

Complex Hilbert structure, the Born rule, tensor composition, the Fock
representation, and CP instruments remain explicit P/S inputs. A chapter
using these inputs cannot establish that Core uniquely selects complex
quantum theory. An external drive supplies work to the blue pulse; a
degenerate logical flag does not turn the operation into an
excitation-conserving closed system.

Red and blue sidebands and their number-state couplings are standard
physics. Wineland and colleagues derive these operators and square-root
couplings in the interaction picture under Lamb-Dicke and rotating-wave
assumptions {[}\emph{Quantum information processing with trapped ions},
DOI 10.1098/rsta.2003.1205, pp.1350-1351{]}. Here the separated
effective Hamiltonians define an exact mathematical model. A device
application separately needs bounds for detuning, pulse overlap,
readout, filtering, and effective-model errors. A large analytic cutoff
does not establish Lamb-Dicke validity throughout its carrier.

\section{68.2 A growth window derived from a self-adjoint generator}

\phantomsection\label{p-micro-04}{}

\textbf{P-MICRO-04 (Finite excitation growth under segmented blue
pulses).} For every nonnegative integer cutoff and every real pulse
area, the blue unitary maps the cutoff carrier into the carrier whose
upper bound is two larger. Any legal finite classical control tree with
a bounded number of blue pulses satisfies the accumulated branchwise
growth bound. This is a consequence of the source operator, not an input
certificate field.

\textbf{Proof.} The red generator vanishes on the vacuum and acts as the
square root of total excitation times the Pauli exchange matrix on each
fixed-excitation block. The blue generator instead decomposes into an
isolated vector and the following orthogonal two-dimensional blocks:

\[
\begin{gathered}
E_n=\operatorname{span}\{|g,n\rangle,|e,n+1\rangle\},\quad
B|_{E_n}=\sqrt{n+1}\,\sigma_x,\quad B|e,0\rangle=0,\\
D(B)=\left\{\psi:\sum_{n\ge0}(n+1)
\bigl(|c_{g,n}|^2+|c_{e,n+1}|^2\bigr)<\infty\right\},\\
B_\beta|g,n\rangle=
\cos(\beta\sqrt{n+1})|g,n\rangle
-i\sin(\beta\sqrt{n+1})|e,n+1\rangle,\\
B_\beta\Pi_M=\Pi_{M+2}B_\beta\Pi_M,\qquad
V_1\Pi_M=\Pi_{M+2}V_1\Pi_M.
\end{gathered}
\]

On the stated graph domain the blue generator is self-adjoint. Each
finite block is symmetric, and the two shifts \(B± iI\) have inverses of
norm at most one. The direct sums are bounded, everywhere-defined
inverses whose ranges lie in the graph domain. Both shifted ranges are
therefore the whole space, establishing self-adjointness. Finite-block
vectors are dense in graph norm. The spectral theorem defines the
unitary for every real area, and exponentiation in each block gives the
displayed formula and its reverse.

The two total excitations in a blue block differ by two, proving the
support statement. Red, phase, and failure-identity operations do not
raise the bound; induction along the tree gives cumulative growth. Small
blue area is unnecessary. A hidden Kraus refinement must represent the
same already established CP event or have an independent microscopic
support proof. Arbitrary additional operations do not inherit the
statement. \(\square\)

\section{68.3 Rectangular instruments, stopping records, and
composition}

\phantomsection\label{p-opt-04}{}

\textbf{P-OPT-04 (Complete rectangular operations on growing carriers).}
For fixed finite actions, outcomes, observation permissions, and maximum
height, the microscopic source defines complete rectangular finite
instruments at every step. After embedding all stopping branches into
one final carrier, their complete CQ output equals that of the infinite
system on the projected input. The equality preserves every visible
label and any finite untouched reference.

Explicitly let \(L\) be the maximum number of operation slots, with at
most one blue pulse per slot, and assume
\(\operatorname{Tr}G\rho_S\le\bar n<N+1\). Set
\(s_N=\operatorname{Tr}[(\Pi_N\otimes I_R)\rho]>0\) and the
finite-coordinate state
\(\rho_N=(J_N^\dagger\otimes I_R)\rho(J_N\otimes I_R)/s_N\). On each CQ
label block, \(\iota_{N+2L}\) conjugates by \(J_{N+2L}\otimes I_R\) into
the common physical Fock space and preserves the label. The ideal output
for an untruncated input may still have an infinite tail; a common
finite-carrier calculation here applies only to projected inputs or
examples with known support. Let each embedding have the corresponding
spectral carrier as its range. The original three-state flag preparation
and complete readout give the source branches; finite branches are
restrictions of the actual maps:

\[
\begin{gathered}
K_0=\sqrt{s(1-q)}V_0,\quad K_1=\sqrt{sq}V_1,\quad K_f=\sqrt{1-s}I,\\
J_M:\mathbb C^{2M+1}\longrightarrow\mathcal H_S,\quad
K_{o,M}^{\rm fin}=J_{M+2}^\dagger K_oJ_M,\\
K_oJ_M=J_{M+2}K_{o,M}^{\rm fin},\qquad
\sum_o(K_{o,M}^{\rm fin})^\dagger K_{o,M}^{\rm fin}=I,\\
\Omega_\beta^\pi\bigl((J_N\otimes I_R)\rho_N(J_N^\dagger\otimes I_R)\bigr)
=\iota_{N+2L}\Omega_{N,\beta}^\pi(\rho_N),\\
D\bigl(\Omega_\beta^\pi(\rho_{SR}),
\iota_{N+2L}\Omega_{N,\beta}^\pi(\rho_N)\bigr)
\le\sqrt{\bar n/(N+1)}.
\end{gathered}
\]

\textbf{Proof.} The preceding support theorem shows that inserting the
output projection leaves each source branch unchanged. Substituting the
derived intertwining identities into source completeness proves
completeness on the entire finite input carrier. A rectangular
restriction of the unitary is an isometry and need not be onto. Multiply
the intertwining identities along each path and take the direct sum over
orthogonal terminal records to obtain complete-output equality. Nodes
within an observation fiber have the same control weights;
distinguishable histories may have different weights. Composition on the
operational quotient remains legitimate.

Every early stopping block is embedded into the common final carrier
before comparison. The completed finite tree is normal CPTP. The
spectral Markov bound and Chapter 67's projected-state estimate then
give the final inequality. The finite-coordinate state is obtained by
restricting the normalized projected joint state through the embedding.
No extra height factor multiplies the initial tail term, because growing
carriers represent its subsequent dynamics exactly. This is an ordinary
trace-distance comparison on a declared input family, not a renamed
energy-constrained diamond norm. \(\square\)

Compressing the blue pulse into the old fixed carrier loses probability:
a boundary ground state transfers squared-sine mass into its excluded
partner. This is a failure of the actual CP output mass, not a
conclusion drawn merely from anticommuting Kraus matrices. The program
compares the rectangular construction with Chapter 66's original
interpreter using a unitary completion on the common carrier. Its action
at unreachable boundary columns is not asserted to equal the full
oscillator unitary.

\section{68.4 Low-spectral filtering, the reference, and a correction}

The previous wording of section 67.5 placed the general possibility of
increased conditional energy after the particular low-spectral filter.
That scope was too broad. This edition makes a minimal bilingual
correction there, while retaining old execution records and result
objects. In this model the retained spectrum is at most the cutoff and
the rejected spectrum is at least one greater. When rejection has
positive probability, the full mean is a convex combination of the two
conditional means, so the retained mean is no larger. If rejection
probability is zero, the passing state is the original state and the
failure reference block is zero; no normalized zero-probability failure
state is defined.

\[
\begin{gathered}
s_N=\operatorname{Tr}[(\Pi_N\otimes I_R)\rho_{SR}],\quad
\delta_N=1-s_N\le e_N:=\bar n/(N+1)<1,\\
\rho_N^{\rm emb}=\frac{(\Pi_N\otimes I_R)\rho_{SR}(\Pi_N\otimes I_R)}{s_N},
\qquad \operatorname{Tr}G(\rho_N^{\rm emb})_S\le\operatorname{Tr}G\rho_S,\\
\rho_R^{\rm fail}=\operatorname{Tr}_S[(Q_N\otimes I_R)\rho_{SR}(Q_N\otimes I_R)],
\quad Q_N=I-\Pi_N,\quad \tau_S=|g,0\rangle\langle g,0|,\\
\mathcal F_N(\rho)=
|{\rm pass}\rangle\langle{\rm pass}|\otimes s_N\rho_N^{\rm emb}
+|{\rm prep\_\allowbreak fail}\rangle\langle{\rm prep\_\allowbreak fail}|\otimes\tau_S\otimes\rho_R^{\rm fail},\\
f(e)=\frac{\sqrt{4e-3e^2}+e}{2},\quad
D\bigl(|{\rm pass}\rangle\langle{\rm pass}|\otimes\rho,\mathcal F_N(\rho)\bigr)
\le f(e_N).
\end{gathered}
\]

The last bound reuses P-MICRO-02 and is not counted as a new theorem. It
follows from at most two nonzero eigenvalues of the purified pass-block
difference and the failure mass, followed by contractivity under
discarding the purification. Failure resets only the system; the
reference retains its actual subnormalized failure-conditioned marginal.
Replacing it by the failure probability times the unconditional
reference changes the complete output, as an entangled-input control
shows. The filter, reset, and record are physical operation assumptions,
not free consequences of writing a projection.

\section{68.5 Energy-weighted omission of the blue pulse}

\phantomsection\label{p-micro-05}{}

\textbf{P-MICRO-05 (Complete-record omission bound using a no-blue
reference).} Fix the same legal record policy and the same red and phase
gates. Compare complete outputs with and without the blue pulses. If the
normalized states along no-blue reference paths preserve the stated
mean-excitation bound, local errors add with their subnormalized prefix
masses. For this chapter's two-step visible policy, the mean selected
blue-pulse count is exactly one, yielding the following uniform bound.

\[
\begin{gathered}
B^2=G+I-2P_e\le G+I,\qquad
\|(e^{-i\beta B}-I)\psi\|^2\le\beta^2\|B\psi\|^2,\\
D\bigl((B_\beta\otimes I_R)\sigma(B_\beta^\dagger\otimes I_R),\sigma\bigr)
\le|\beta|\sqrt{\operatorname{Tr}G\sigma_S+1},\\
Q_{00}=Q_{11}=313/1250,\qquad Q_{01}=Q_{10}=312/1250,\\
\mathbb E_Q[\#1]=1,\qquad
D(C_\beta(\rho),C_0(\rho))\le|\beta|\sqrt{\bar n+1}.
\end{gathered}
\]

\textbf{Proof.} Squaring on the blue blocks gives the first identity.
The spectral theorem and the scalar exponential inequality give the
vector bound. Purify a mixed state without acting on either reference or
purification, bound pure-state trace distance by vector distance, and
then discard the purification. Finite mean excitation places the
purified vector in the domain of \(B\otimes I\). The displayed
expression uses \(\|B\psi\|^2\) and does not require the operator domain
of \(B^2\).

Telescope the difference between the two complete processes slot by
slot. After each local difference, the same actual future channel is
applied and cannot increase trace norm. Before that difference, use the
no-blue reference, so red and phase gates preserve the input energy
bound. Zero-mass prefixes contribute zero without normalizing their
states. Orthogonal record blocks give each contribution the mass of
reaching its prefix and selecting label one. Both label-one marginals of
the displayed policy are one half, so their mass sum is one. The
preceding section ensures that the same mean budget holds after passing
the low-spectral filter. \(\square\)

This proof does not assume that actual blue-containing paths retain
their initial energy, and it does not let the controller use hidden true
parameters. The source's conditional selected-pulse mean is twice its
probability parameter; the raw source including failures has yet another
count. Neither replaces the deployed policy's count in this comparison.
State-dependent readout, additional energy-injecting gates, or altered
observation permissions require a new prefix budget.

\section{68.6 Simultaneous interaction does not preserve a hard window}

\phantomsection\label{c-micro-02}{}

\textbf{C-MICRO-02 (An infinite-tail counterexample for simultaneous red
and blue interaction).} The segmented result does not imply finite
excitation support after finite-time simultaneous red-plus-blue
evolution. Let \(X\) be the canonical self-adjoint closure of
\(a+a^\dagger\) on finite Fock vectors, and choose
\(H_{\rm sim}=\sigma_x\otimes X\) with domain \(D(X)\oplus D(X)\) in the
exchange eigenbasis. It agrees with the formal sum \(A+B\) on the common
finite core. The parameter \(t\) below is dimensionless pulse area.

\[
\begin{gathered}
H_{\rm sim}=\sigma_x\otimes X,\qquad
e^{-itH_{\rm sim}}(|+_x\rangle\otimes|0\rangle)
=|+_x\rangle\otimes|-it\rangle,\\
\Pr(n=k)=e^{-t^2}\frac{t^{2k}}{k!}>0\quad(t\ne0,\ k\ge0).
\end{gathered}
\]

\textbf{Proof.} On the positive eigenstate of the probe exchange
operator, the generator reduces to the oscillator position quadrature.
The Weyl exponential of its self-adjoint closure is a displacement that
sends vacuum to the displayed coherent state. Every number-state mass is
strictly positive, leaving nonzero mass outside every finite cutoff. The
initial probe-vacuum state nevertheless has total excitation support at
most one. \(\square\)

The counterexample refutes exact hard-support extrapolation. It does not
rule out controlled tails, Trotter approximation, or study of the full
Rabi model. General simultaneous interactions, weak leakage, and mode
deletion need their own budgets. A refuted implication is not relabeled
as a proved positive statement.

\section{68.7 A certificate for the complete common-source chain}

\phantomsection\label{p-rob-03}{}

\textbf{P-ROB-03 (A joint certificate for growth, physical filtering,
and pulse omission).} Use the source, two-step acceptance event, and
legal policy above. On the declared design interval, initial mean
budget, and ideal operation contract, both the matched-blue and
omitted-blue implementations have sufficient complete-CQ error bounds.
The latter additionally requires an independent small-pulse guarantee.

\[
\begin{gathered}
H=s^2I,\quad p_{\rm source}=s^2\ge1369/2500,\quad
p_{\rm actual}=s_Ns^2,\\
\bar n=1,\quad N=4095,\quad |\beta|\le1/100,\quad
\varepsilon_{\rm required}=1/20,\\
\varepsilon_{\rm matched}\le\frac1{50}+\frac{1+\sqrt{16381}}{8192},\\
\varepsilon_{\rm omitted}\le\frac1{50}+\frac{1+\sqrt{16381}}{8192}
+\frac{\sqrt2}{100}
<0.049887775359266728<\frac1{20}.
\end{gathered}
\]

Here \(p_{\rm actual}\) denotes an original-source attempt including the
filter. The compiler has no original-source failure after filter passage
and delivers with probability \(s_N\).

\textbf{Proof.} The original seven-leaf instrument gives the scalar
backward effect in the first line. Its complete conditional target is
therefore itself CPTP, and preparation error contracts through it. For a
common input, shared actual gates make the conditional quantum state
identical on each target/compiler path. Complete-CQ distance is exactly
path TV, bounded by one fiftieth on the design interval by P-ROB-02.
Insert the actual filter and add its sharp error by the triangle
inequality. The omitted implementation adds the preceding blue-omission
bound. Failure records, system, and reference all belong to the same
output space.

An independent rational checker encloses both radicals using integer
square roots, proving the strict precision inequality without
floating-point rounding. Its upper endpoint is 4988777535926672699
divided by ten to the twentieth power; the remaining allowed error is
approximately 0.000112224640733273. The looser filtering bound used for
Chapter 67's simpler example would instead give approximately
0.05001127625 and would not certify this precision. Failure of a weaker
upper bound is not physical impossibility. \(\square\)

A real device must additionally combine gate-matching, filtering,
effective-model, and readout errors on this same complete output. None
is measured here. Floating calculation tolerance, calibration coverage
probability, and physical trace-distance error are distinct quantities.
The general sharp postselection extension remains available, but this
example requires no additional division by acceptance: source
construction has already established linearity and trace preservation of
its conditional target.

\section{68.8 Flag counts do not identify blue pulse area}

\phantomsection\label{c-rob-02}{}

\textbf{C-ROB-02 (Identical flag probabilities do not certify the
blue-pulse terminal state).} All original flag-record probabilities in
this model depend only on the probability component of the parameter.
They cannot identify blue pulse area or certify complete terminal-state
accuracy after omission.

\textbf{Proof.} Every labelled event is a parameter scalar times a
unitary, so path masses are products of scalars. Use a vacuum input and
compare zero blue area with an area of half pi. The original red and
phase gates fix the vacuum. On deployed word zero-one, the latter model
sends it to an orthogonal excitation state. That visible path alone
contributes the lower bound

\[
D(C_{\pi/2}(|g,0\rangle\langle g,0|),C_0(|g,0\rangle\langle g,0|))
\ge Q_{01}=156/625>1/20.
\]

All flag distributions remain identical. This control deliberately
violates the independent small-area premise to show why flag calibration
cannot establish that premise. It does not refute the sufficient bound
when the premise holds. \(\square\)

Deployment therefore still chooses one set of legal weights from
independent calibration and requires validity for every parameter in the
resulting joint set. The inherited flag calibration can constrain the
probability parameter, but cannot certify blue area, shared gate
identity, or terminal testers. Independent pulse, spectral, or process
calibration is needed, preserving common nuisance correlations. Without
it, the device conclusion is insufficient evidence, not certified
success or certified impossibility.

\section{68.9 Complete costs and two physical schedules}

One instrument may have different physical schedules. Reading an
independent flag before applying its selected gate gives a selected-call
count derived from all successful and failed leaves. Coherently
controlling gates before flag readout may schedule exactly one blue
control in every active slot. The scheduled equalities below refer
specifically to that implementation, not to every coherent-control
device. The compiler formulas concern the matched-blue implementation;
the omitted implementation has zero blue calls. Labels alone do not
decide which hardware contract applies.

\[
\begin{gathered}
\mathbb E[\#B]_{\rm source,feedforward}=s_Nsq(1+s),\quad
\mathbb E[\#B]_{\rm compiled,feedforward}=s_N,\\
\mathbb E[\#B]_{\rm source,scheduled}=s_N(1+s),\quad
\mathbb E[\#B]_{\rm compiled,scheduled}=2s_N,\\
\mathbb E[C_{\rm per\ delivery}]
=\frac{\mu_{\rm complete\ attempt}}{s_Ns^2}
+\frac{C_{\rm calibration}}{m}.
\end{gathered}
\]

The last formula is restricted to genuinely re-preparable, fresh-reset,
independent identically distributed source attempts; calibration is
amortized over delivered outputs. The attempt mean includes preparation,
filtering, failed resets, auxiliaries, control, randomness, readout,
records, and classical computation, charged on each item's actual
occurrence event. A cost conditional on filter passage is multiplied by
passage probability before entering the numerator. Unknown quantum
inputs are not freely repeatable. Omitting a pulse does not itself
establish measured time, energy, or monetary advantage.

\section{68.10 Evidence, integration, and continuing responsibilities}

The six additions are analytic proofs or counterexamples in their stated
domains, supported by independent exact arithmetic and finite floating
checks. There are no new Lean declarations. The unchanged 158
proof/configuration inputs and 13 runtime inputs are reused by source
identity; their original receipts are not rewritten as chapter-68
executions. New source calculations, rectangular composition, reference
and failure blocks, noncommuting CP maps, and invalid-boundary controls
have new actual execution identities. Three native research tasks
performed cross-review. This is internal AI review, not external human
peer review.

The main numerical witness uses initial cutoff \(N=1\), a
two-dimensional reference, and final carrier \(G\le5\). The
nonzero-failure control uses input support \(G\le2\), filter cutoff
\(N=1\), and common output \(G\le6\), with a deliberately chosen tail
mass \(1/4096\). It checks the physical filtering interface, not
numerical execution of the analytic \(N=4095\) carrier.

No protocol-matched raw device data were acquired. The inherited
synthetic calibration pipeline remains synthetic evidence and cannot
certify the new blue-area premise. The next verifiable theoretical
responsibility is a common-intervention tail bound for simultaneous
evolution or finite-bath mode deletion, together with an assessment of
whether actual control errors fit the remaining precision margin.
Complete OPT, independent complex-quantum selection, mechanism
identification, long-time limits, QFT and EFT, particle and flavour
predictions, gravity, and cosmology retain separate responsibilities in
the current roadmap. Completing this finite growth example does not
complete physics across those domains.

\chapter{69. Simultaneous dynamics, genuine truncation error, and
calibrated complete predictions}\label{ch-69}

\section{69.1 Continuing from the infinite-tail counterexample}

Chapter 68 established growing finite carriers for separated red and
blue pulses, while simultaneous controlled force can produce an infinite
tail at every nonzero time. This chapter makes the latter model
computable with controlled error. The comparison is between the true
unitary and evolution obtained by \textbf{compressing the Hamiltonian
before exponentiating}. Projecting the exact evolved state, normalizing
its surviving block, or observing agreement between two large matrices
does not discharge this obligation. The six results have B/P/E/O
dependency classes. Complex quantum structure, the Fock representation,
Born readout and effective control Hamiltonians remain independent
inputs. Argument, verification and dependency categories remain
distinct.

The system is a complex qubit probe and one oscillator, with
\(G=n+P_e\), \(P_M=1_{G\le M}\) and its canonical rank-\(2M+1\)
embedding \(J_M\). Any positive-dimensional finite complex reference is
untouched and may initially be entangled. The general dynamics and the
calibrated subfamily are

\[
\begin{aligned}
A&=a\sigma_++a^\dagger\sigma_-,&B&=a^\dagger\sigma_++a\sigma_-,\\
H_0&=\omega n+\Delta P_e+gA,&H&=H_0+hB,\quad\omega>0,\\
H_x&=x\sigma_xX=x(A+B),&X&=a+a^\dagger,\\
V_1(x)&=e^{-iH_x},&V_0&=e^{-i(\pi/3)P_e},\\
H_K&=J_K^\dagger HJ_K,&U_K(t)&=e^{-itH_K}.
\end{aligned}
\]

All action durations are \(t\ge0\), with signed couplings allowed. The
general generator uses frequency units, with physical Hamiltonian
\(\hbar H\); the exposure \(\int|h|dt\) is dimensionless. The calibrated
subfamily has unit slot duration and signed area \(x\). Its
zero-frequency quadrature realization is specified independently rather
than obtained by applying the positive-frequency domain theorem at its
boundary. The gate \(V_1\) is simultaneous evolution, not the previous
product of separated exponentials. External driving supplies energy.

The original source remains \(K_0=\sqrt{s(1-q)}V_0\),
\(K_1=\sqrt{sq}V_1\), \(K_f=\sqrt{1-s}I\), with \(s=1-q/2\). Two slots
stop at \(f\), retaining the seven terminal words
\(f,0f,1f,00,01,10,11\); the last four are accepted. These original
events directly give the backward acceptance effect \(s^2I\), so this
model's normalized accepted target \(T_\infty(q,x)\) is CPTP. Failure
words retain zero mass in its common record alphabet. Source and
available physical actions share the same unknown parameter; deployment
reads only permitted visible histories.

\section{69.2 Approximate operational stability of a complete tree}

\phantomsection\label{p-opt-05}{}

\textbf{P-OPT-05 (Complete-CQ error from unnormalized reference
prefixes).} Fix a finite-height classical control tree, actions, visible
outcomes, stopping rule and observation permissions. Two normal CP
instrument families correspond label by label, with each complete test
trace preserving. Persistent memory is included in the system, or the
auxiliary is genuinely reset. Let \(\sigma_h^F\) be the unnormalized
joint state reaching prefix \(h\) when every preceding instrument is the
finite/reference family \(F\). Then their complete terminal CQ outputs
obey

\[
\begin{aligned}
d_h(\sigma)&=\frac12\sum_o\|(E_{ho}-F_{ho})(\sigma)\|_1,\\
D(\Omega_E(\rho),\Omega_F(\rho))
&\le\sum_{h\ \mathrm{internal}}d_h(\sigma_h^F).
\end{aligned}
\]

Comparison uses the same terminal embedding; finite reference states are
embedded into the common system. Every label, early failure, stop,
hidden Kraus sum and untouched reference is retained. A zero-mass prefix
contributes zero without normalization.

\textbf{Proof.} At each prefix insert the local difference
\((E_h-F_h)\sigma_h^F\), followed on the left by the complete true
suffix. The right prefix always remains from the same family \(F\).
Finite recursive expansion cancels the intermediate terms and expresses
the complete output difference as the sum of these suffix-applied local
differences. Each suffix is a normal CPTP map from all current visible
outcomes to all terminal records. For a Hermitian trace-class operator
\(Z=Z_+-Z_-\), positivity and trace preservation give
\(\|\Lambda Z\|_1\le\operatorname{Tr}Z_++\operatorname{Tr}Z_-=\|Z\|_1\).
The trace norm of the record direct sum equals the sum of block norms.
The triangle inequality proves the claim, also after identity extension
on the reference. \(\square\)

One cannot replace an individual finite prefix by an actual prefix while
retaining the former's support or energy budget. Controllers constrained
by observation fibres still use common weights on each fibre; this proof
grants no hidden-node access. Exact operational-quotient compatibility
is a different statement. Additional state-dependent readout uses this
general prefix bound rather than automatically inheriting the
state-independent path probabilities below.

\section{69.3 Energy and boundary flow from the actual generator}

\phantomsection\label{p-micro-06}{}

\textbf{P-MICRO-06 (Self-adjoint simultaneous source, energy budget, and
boundary defect).} For the real parameters of Section 69.1 with
\(\omega>0\), \(H,H_0\) are self-adjoint on \(D(n)\). The force
subfamily instead uses the self-adjoint realization \(xX\oplus(-xX)\) in
the \(\sigma_x\) eigenbasis. The difference between finite compression
and true dynamics follows from the two actual blue boundary edges;
finite-energy inputs satisfy the following budget and local
complete-output bound.

The existing block realizations give \(A^2=G\) and \(B^2=n+P_g\le n+I\).
For \(\psi\in D(n)\), the coupling norm is at most
\((|g|+|h|)\|(n+I)^{1/2}\psi\|\), whose relative bound with respect to
\(n\) can be arbitrarily small:
\(\|(n+I)^{1/2}\psi\|\le\epsilon\|n\psi\|+(1+1/(4\epsilon))\|\psi\|\).
The standard Kato--Rellich theorem supplies self-adjointness on the
stated domain. That theorem is an explicit standard mathematical
dependency, not a new UEOT discovery or machine proof. Every block of
\(H_0\) preserves \(G\). The zero-frequency subfamily is defined
directly by self-adjoint quadrature spectral calculus, and its finite
Fock vectors satisfy the domain requirements of the finite Duhamel
expansion below.

Set \(C=i(a^\dagger\sigma_+-a\sigma_-)\). First work in finite matrices:
\(i[B,G]=-2C\) and \(C_K^2\le J_K^\dagger C^2J_K\le G_K+I\). Integration
gives the following homogeneous budget for an unnormalized input of mass
\(m\) and excitation \(e\):

\[
\begin{aligned}
e_{\mathrm{out}}+m
&\le\bigl(\sqrt{e+m}+z\sqrt m\bigr)^2,
&z&=\int|h(t)|\,dt,\\
\sqrt{E(t)+1}&\le\sqrt{E(0)+1}+z
&&\text{when }m=1.
\end{aligned}
\]

Extension to the infinite model has a noncircular order. The finite
Duhamel theorem in Section 69.4 does not consume this energy bound; it
first provides Galerkin strong convergence for finitely supported
inputs. Lower semicontinuity of the positive \(G\) expectation transfers
the finite-matrix budget. Low-spectral initial-state approximation and a
second lower limit extend it to arbitrary normal finite-mean-excitation
inputs. There is no assumption that every finite-energy vector lies in
\(D(n)\) or has a pointwise differentiable energy expectation. A rare
normalized measurement branch does not automatically inherit an
aggregate mean budget.

The actual embedding defect is \(HJ_K-J_KH_K=h(I-P_K)BJ_K\). Calculating
the blue edges gives

\[
\begin{aligned}
W_K&=J_K^\dagger B(I-P_K)BJ_K\\
&=(K+1)|g,K\rangle\langle g,K|
 +K|g,K-1\rangle\langle g,K-1|,\\
d_K(\sigma)&\le |h|\int_0^{|t|}
\sqrt{m\operatorname{Tr}[(W_K\otimes I_R)\sigma_K(u)]}\,du.
\end{aligned}
\]

For \(K=0\) omit the second term. Here \(\sigma_K(u)\) is the actual
finite-generator evolution of the unnormalized prefix and
\(m=\operatorname{Tr}\sigma\). Apply
\(U(t)J_K-J_KU_K(t)=-i\int_0^tU(t-u)(HJ_K-J_KH_K)U_K(u)du\) to a
purification, followed by unitarity, Cauchy--Schwarz and trace-distance
contraction. The purification has norm \(\sqrt m\), so the mass inside
the square root cannot be omitted. Combining this defect with P-OPT-05
gives a bound for the same complete tree. Numerical certification of the
integral would additionally require validated quadrature; the principal
certificate instead uses the integral-free prior bound next.

\section{69.4 A finite remainder for genuine Galerkin evolution}

\phantomsection\label{p-micro-07}{}

\textbf{P-MICRO-07 (Genuine truncation bound from finite Duhamel
expansion).} Let \(N,L\ge0\) and \(K=N+2L\). In either specified
self-adjoint model of Section 69.1, start from support in \(P_N\).
Finitely many parameter segments and intervening \(G\)-preserving
unitaries are permitted, with total absolute blue exposure \(z\) along a
path. For every untouched finite reference and normal joint input, the
trace distance between the true evolution and embedded \(e^{-itH_K}\)
evolution is bounded by

\[
\begin{aligned}
c_j(N,z)&=\frac{z^j}{j!}\prod_{r=1}^{j}\sqrt{N+2r-1},\qquad c_0=1,\\
\|(U-J_KU_KJ_K^\dagger)P_N\|&\le2c_{L+1}(N,z),\\
d_{N,L}(z)&=\min\{1,2c_{L+1}(N,z)\}.
\end{aligned}
\]

For unnormalized mass \(m\), multiply the last bound by \(m\). The same
cutoff is used throughout the path; the true evolution is not assumed to
remain inside it.

\textbf{Proof.} Begin with \(U(t)=U_0(t)-ih\int_0^tU(t-u)BU_0(u)du\) and
expand only its leftmost true \(U\), for exactly \(L+1\) iterations. The
true unitary remains on the far left as a suffix. Each right-hand
\(U_0\) preserves excitation support, and each \(B\) increases its upper
limit by at most two. The input of the \(r\)th \(B\) therefore lies in
\(P_{N+2r-2}\), where its norm gain is at most \(\sqrt{N+2r-1}\). The
ordered simplex integral of the absolute coupling product is \(z^j/j!\),
and the true suffix has norm one. Hence the order-\(L+1\) remainder is
at most \(c_{L+1}\).

Repeat the same finite expansion for the actual compressed Hamiltonian.
Every intermediate support through order \(L\) lies within \(P_K\), so
those terms agree exactly with their infinite counterparts. The finite
remainder has the same bound because inserted projections cannot
increase norm. Subtraction proves the operator estimate; purification
and partial trace prove the joint-state estimate. Finite segments and
conserving gates can be included in the conserving reference propagator,
with the ordered integrals bounded by total absolute exposure.
\(\square\)

Only a finite expansion is needed; no unaudited infinite Dyson
convergence is invoked. At fixed \(N,z\), the ratio
\(c_{j+1}/c_j=z\sqrt{N+2j+1}/(j+1)\) tends to zero, so a sufficiently
large \(L\) makes the error vanish. This does not authorize unbounded
times, unlimited stopping or coherent controllers. An unfiltered
energy-bounded input additionally requires its initial projection error
\(\sqrt{\bar n/(N+1)}\). Actual physical filtering instead uses the
established complete-preparation bound, retaining its failures.

Here record probabilities are independent of the quantum state, allowing
complete-word weighting. The legal Chapter 67 policy has
\(Q_{00}=Q_{11}=313/1250\) and \(Q_{01}=Q_{10}=312/1250\), with mean
force-call count exactly one. At nominal area \(\widehat x\),

\[
\epsilon_{\mathrm{Gal}}(\widehat x)
\le\frac{624}{1250}d_{N,L}(|\widehat x|)
 +\frac{313}{1250}d_{N,L}(2|\widehat x|).
\]

Failed words of the original source use their actual probabilities, not
\(Q\). This is the interface between the general prefix theorem and the
special path formula.

\section{69.5 Compression need not inherit the infinite invariant}

\phantomsection\label{c-micro-03}{}

\textbf{C-MICRO-03 (Galerkin compression does not preserve this
quadrature invariant).} The infinite force and probe phase strongly
commute with the spectral projections of \(X\), preserving its \(X^2\)
quadratic-form expectation. At \(K=1\), in basis \((g0,g1,e0)\), the
compressed unit-area generator couples only \(g1\leftrightarrow e0\),
whereas \(J_1^\dagger X^2J_1=\operatorname{diag}(1,3,1)\). Starting from
\(e0\), finite evolution has expectation \(1+2\sin^2x\), while true
infinite evolution retains one. They differ unless \(x\) is an integer
multiple of \(\pi\).

The proof is the explicit map \(e0\mapsto\cos x\,e0-i\sin x\,g1\). The
two-quantum same-probe matrix elements in \(X^2=a^2+(a^\dagger)^2+2n+1\)
have no corresponding endpoints within this carrier. This refutes
automatic inheritance of the invariant, without refuting the
trace-distance approximation in P-MICRO-07. Parameter comparison must
therefore first use two exact infinite models and pay the Galerkin error
separately.

\section{69.6 Identifiable calibration and a nominal parameter without
hidden truth}

\phantomsection\label{p-rob-04}{}

\textbf{P-ROB-04 (Joint calibration, nominal parameters, and
shared-force stability).} Freeze independent fresh preparations, fixed
parameters, ideal readout and the shared-pulse model. Estimate \(q\)
from original-source failures. Separately apply the same force to
\(|+_x\rangle\otimes|0\rangle\) and read the binary event \(p\le0\),
where \(p=(a-a^\dagger)/(i\sqrt2)\). Its probability is the strictly
monotone \(\Phi(2x)\), identifying signed area. Complete raw calibration
records define one joint set; deterministic inversion and rounding
select rational \(\widehat x\) and a whole-set displacement radius
\(\Delta_x\). One lawful deployment must work for all parameters in that
set.

Original-source failure probability is \(q/2\), so
\(\widehat q=2F/n_q\). A failure-frequency radius cannot be used as a
\(q\) radius without scaling. Choose \(n_q=120000\) with \(q\) radius
\(0.01\), and \(n_x=1000000\) with drive-probability radius \(0.002\).
Hoeffding and a union bound give

\[
\begin{aligned}
\Pr\{(q,x)\notin\mathcal C(D_{\mathrm{cal}})\}
&\le 2e^{-6}+2e^{-8}<0.01,\\
\mathcal C&=\{q:\ |q-\widehat q|\le0.01\}
\times\{x:\ |\Phi(2x)-\widehat p|\le0.002\},\\
s&=1-q/2,\qquad
\Delta_x\ge\sup_{(q,x)\in\mathcal C}|x-\widehat x|.
\end{aligned}
\]

The product set corresponds to the two declared parameter calibrations,
while the shared relation between \(s\) and \(q\) is preserved. The
analyzer requires the entire untrimmed set to lie in the design window
\(q\in[0.48,0.52]\), \(|x|\le0.1\). It cannot discard adverse parameters
and claim coverage of the original set. The coverage statement bounds
the unconditional probability of issuing an invalid certificate in this
one fixed-sample protocol. It is not conditional coverage after
admission and does not allow uncorrected resampling until admission.
Rational normal-integral enclosures and monotone inversion supply
independently checkable endpoints. The constants, radicals and
exponential bounds also receive outward arithmetic enclosures.

For a passed input with mean excitation at most one, the form inequality
\(X^2\le4n+2I\le4G+2I\) gives expectation at most six. Exact infinite
forces and all probe phases strongly preserve this form domain, so
replacing one actual area by its nominal value has vector error at most
\(|x-\widehat x|\sqrt6\). Telescope each complete word and weight by
\(Q\), using mean call count one:

\[
D(C_\infty(x),C_\infty(\widehat x))
\le\sqrt6\,\Delta_x.
\]

The inputs and complete preparation records agree in this comparison.
Finite mean \(n\) places a purification in \(D(X)\), the quadratic-form
domain of the positive operator \(X^2\); it need not lie in \(D(X^2)\).
This invariant cannot be applied to compressed prefixes, as C-MICRO-03
shows.

Flags remain independent of \(x\), so flags alone do not produce an area
set. Nor does the momentum calibration identify arbitrary devices. An
unknown additive readout bias \(p_{\rm obs}=p+b\) can change its
probability to \(\Phi(2x-\sqrt2b)\); the transformation
\((x,b)\mapsto(x+c,b+\sqrt2c)\) preserves the data. This
nonidentifiability concerns single-dose calibration access. A second
dose with the same fixed bias has probability \(\Phi(4x-\sqrt2b)\), and
the two doses can jointly identify \(x,b\). This analyzer does not refit
bias using holdout data and retains the declared zero-bias model.
Independent zero/bias, preparation and shared-gate contracts are not
validated by the simulation. Correlated trials also cannot impersonate
iid data: repeating one latent binary bit a million times does not
acquire Hoeffding precision.

\section{69.7 One certificate for the complete output}

\phantomsection\label{p-rob-05}{}

\textbf{P-ROB-05 (Composing physical preparation, robust deployment,
nominal calibration, and finite prediction).} Under the shared model,
legal visible policy, physical low-spectral filtering, shared gates and
valid calibration-set contract, every normal joint input with
\(\operatorname{Tr}G\rho_S\le1\) and every finite untouched reference
satisfy

\[
\begin{aligned}
D(T_\infty(q,x),C_K(\widehat x))
&\le f(1/(N+1))+\frac1{50}+\sqrt6\Delta_x
       +\epsilon_{\mathrm{Gal}}(\widehat x),\\
f(e)&=\frac{\sqrt{4e-3e^2}+e}{2}.
\end{aligned}
\]

The complete output retains preparation failure, every path, terminal
states and the reference. Filter failure is
\(\tau_{\rm vac}\otimes\operatorname{Tr}_S[(I-P_N)\rho(I-P_N)]\), with
the actual failure-conditioned reference; the ideal target has zero mass
in this block. Low-spectral passing does not increase mean \(G\), unlike
unrestricted postselection.

\textbf{Proof.} Pay the inherited sharp physical-filter error once. At
the same parameter, source and action terminal blocks agree, so the
legal-policy difference from the conditioned target is the existing
path-TV bound \(1/50\). Next pay P-ROB-04's nominal-parameter difference
only between exact infinite models. Apply P-MICRO-07's complete-path
weighted error to the nominal model. The triangle inequality gives the
result. All four terms concern the same output, reference and acceptance
contract; input tails, dynamics and statistical coverage are not
conflated. \(\square\)

The physical compiler \(C_\infty(x)\) calls the actual shared gates, and
its target error needs only \(f+1/50\). The nominal infinite process is
a comparison object and \(C_K(\widehat x)\) is a predictor; hardware is
not assumed to implement a compressed Hamiltonian. Calibration coverage
concerns whether the set contains the parameter, not measured validity
of the device model. Additional device errors require their own
contributions in the same complete metric.

Freeze \(N=4095,L=40,K=4175\). Over the entire design window, the
Galerkin term is at most approximately \(4.54172\times10^{-5}\); the
calibrated nominal value permits a tighter bound. The actual synthetic
records have \(F=30036\) and \(547267\) drive events, giving
\(\widehat q=0.5006\), \(\widehat p=0.547267\) and

\[
\begin{aligned}
q&\in[0.4906,0.5106],\\
x&\in[0.0568560627871\ldots,0.0619048193330\ldots],\\
\widehat x&=2969/50000,\qquad
\Delta_x\le0.002524819334,\\
\epsilon_{\mathrm{finite}}&<0.041930158795<0.05.
\end{aligned}
\]

Strict endpoints and upper bounds are the rational numbers in the
receipt; the displayed decimals are for reading. Widening the
probability calibration radius to \(0.01\) gives approximately
\(0.06671\) and the analyzer returns insufficient certification, not
physical impossibility. The retained policy is the established
whole-design policy; no new minimax optimality is claimed for every
narrower calibrated set.

\section{69.8 Actual calculation, held-out protocols, and full
resources}

Ordinary small-matrix checks use initial \(G\le1\), \(K=7\) and a
two-dimensional reference, with the existing Chapter 66 interpreter
independently computing the complete source, backward effect and
conditioned target. A larger finite matrix is diagnostic and has its own
analytic error. Independent force checks also use actual
infinite-displacement matrix elements rather than calling a large cutoff
an infinite oracle. The positive-frequency Rabi stress case fixes
\(\omega=1,\Delta=3/5,g=2/5,h=1/10,t=1/5\); these finite numerical
samples do not prove all-parameter claims.

A separate executed integer-interval predictor uses \(K=4175\) and fixed
input \((|g,1\rangle|0_R\rangle+|e,0\rangle|1_R\rangle)/\sqrt2\), which
passes filtering with probability one. A degree-100 Taylor polynomial
uses outward-rounded radicals, products and divisions. The bound
\(\|H_K\|\le13\) supplies an explicit remainder, while sparsity
identifies all unpropagated coordinates. Each word's integer amplitude
centre is normalized by its exact norm, producing a lawful CQ
prediction. Its actually certified numerical error is below \(10^{-8}\),
so this specified output has total bound below \(0.041930168795\). This
numerical certificate concerns that input only. The all-input theorem
above concerns the exact finite map; one interval example does not
certify numerical evaluation on every input.

The synthetic analysis separates calibration, truth provenance and
held-out data; the analyzer never reads generator truth. Independent
drive holdout applies the same gate twice in succession and predicts
\(\Phi(4x)\), testing a different shared-parameter response from the
single-pulse calibration. Complete source and compiler records are held
out separately. A mismatched-drive control uses the same frozen
threshold. These are synthetic finite-tester checks, not an empirical
measurement of full trace distance. The final sample contract contains
120000 source-flag calibration trials, 1000000 drive calibration trials
and 150000 ordinary holdout trials, with intentional mismatch controls
separately registered. Time, energy, monetary cost and real instrument
errors remain unmeasured.

A real source attempt that filters first succeeds with probability
\(s_Ns^2\). Reading flags before applying gates gives mean force calls
\(s_Nsq(1+s)\) for the source and \(s_N\) for the compiler. Coherent
scheduling may instead require \(s_N(1+s)\) and \(2s_N\) gate slots;
label counts do not make scheduling free. Only genuine fresh-reset
independent preparation permits delivery cost
\(\mu/(s_Ns^2)+C_{\rm cal}/m\). One-time calibration is amortized over
deliveries; unknown quantum inputs and references cannot be freely
recreated. Computational prediction cost and physical gate cost are also
recorded separately.

\section{69.9 Sources, verification axes, and the next arrow}

This chapter uses standard CP contraction, Kato--Rellich, Duhamel and
Hamiltonian-truncation methods. Tong et al.,
\href{https://arxiv.org/abs/2110.06942v2}{Quantum 6, 816 (2022)},
especially Appendix D, distinguish state truncation from Hamiltonian
truncation. Woods, Cramer and Plenio,
\href{https://arxiv.org/abs/1504.01531}{PRL 115, 130401 (2015)}, derive
bosonic-bath and Fock-truncation errors. These precedents are not
renamed UEOT discoveries. This increment supplies explicit coefficients
for the declared probe model, complete record/reference/permission
interfaces and shared-source integration with calibrated nominal
predictions. {[}QM3113{]}

The universal arguments of the six results are paper proofs.
Rational/integer intervals, finite floating comparisons, synthetic raw
streams and internal AI cross-review have separate evidence identities.
There are no new Lean declarations. The unchanged 158
proof/configuration inputs and 13 earlier runtime inputs are checked
before reusing historical machine execution. The original 188 result
objects and 111 ledger claims retain their identities. No
protocol-matched real device calibration, complete terminal-state data
or costs were obtained, so the empirical axis remains untested.

The closed arrow is specified simultaneous dynamics to a controlled
finite predictor of the same intervention process. The next admissible
extension is mode deletion or persistent-memory comparison in a
specified finite Gaussian/free-field bath: derive prefix defects from
the actual coupling, bound the same intervention slots and complete
records, and retain a control showing that deleting a mode is not
resetting a bath. Full OPT/independent complex selection, long-time
limits, coherent controls and infinite references retain distinct
obligations. Same-order EFT observables, the two high-energy generation
candidates, flavour relations, new structure between Fisher and
Lorentzian geometry, and the cosmological targets remain on the
continuous roadmap.

\chapter{70. Complete persistent-environment processes, finite
certification, and selection limits}\label{ch-70}

\section{70.1 Shared object and proof obligations}

The actual Chapter 69 workspace, candidate commit, and historical
execution evidence have been recovered. PR465 merged the original
candidate into main \texttt{ee7fb2793a1d\allowbreak bb519703bba0\allowbreak 5f535f0d962d\allowbreak 9e56};
its original computations have not been relabeled as new executions.
This chapter continues from that verified basis with a three-slot
process that retains the same environment throughout. The canonical
contract, frozen before parallel research, is
\texttt{persistent\_\allowbreak bath/\allowbreak SHARED\_\allowbreak MODEL.\allowbreak json}, with SHA256
\texttt{685d373bb5a7\allowbreak 0a677fd4d110\allowbreak cb81c6624260\allowbreak de95ea39cc60\allowbreak 070a8f4bd9f1\allowbreak 9668}.
Dynamics, instruments, complete records, calibration, finite prediction,
and held-out decisions below refer to this one object. {[}QM3114{]}

Two separate questions organize the chapter: whether the available
operational principles uniquely select complex quantum theory, and
whether a specified source yields controlled, reproducible complete
predictions after complex quantum physical inputs are explicitly
accepted. Resolving the latter does not remove the premises of the
former. Core objects, composition, and persistence are read-only
mathematical interfaces. Complex Hilbert spaces, the Born rule, tensor
products, CP instruments, the Fock representation, vacuum preparation,
and the declared readout structure remain independent physical inputs.
Formulas, program receipts, and synthetic samples carry their respective
proof or computational obligations.

The common error object is the CQ output containing every stopped
record. The reference \(R\) may have any finite dimension, may be
entangled with input \(S\), and remains untouched during source
evolution, instruments, and feedback; a declared terminal tester may
read the output \(S/R\). Every all-input statement concerns the external
channel of this fixed policy. Changing intermediate instruments,
feedback, or stopping requires a separately declared process. Numerical
and statistical certification holds inside the stipulated two-mode
vacuum model; it does not certify a real device or complete the
unification of physics.

\section{70.2 Reachable sectors and the boundary of independent quantum
selection}

Complex Hilbert spaces, the Born pairing, tensor products, CP
instruments and Fock dynamics remain explicit physical inputs in this
chapter. A complete process calculation can test a model assembled from
those inputs; its success does not establish that Core has independently
selected them. We distinguish the \textbf{reachable operational theory}
generated by the registered preparations, operations and effects from a
\textbf{global theory} concerning every physically allowed system and
sector. The following construction states why these targets cannot be
silently exchanged.

\phantomsection\label{c-qsel-02}{} \textbf{C-QSEL-02 (Operational
equivalence of hidden completions of a reachable sector and the boundary
of global quantum selection).} Fix a reachable complex-quantum operation
library, including complete finite tests, parallel composition,
conditional composition, coarse graining and discarding. It has a
completion with an additional central superselection sector that
reproduces every finite adaptive circuit from reachable preparations,
with unrestricted finite length and with references, failures, stops and
terminal effects retained. Nevertheless, the completion's full state
space contains a state that has no purification using any finite
direct-sum matrix auxiliary system.

\textbf{Proof.} Complete the algebra of a reachable system to

\[
\widetilde{\mathcal A}=M_d(\mathbb C)\oplus\mathbb C.
\]

States are \((\rho,h)\), with
\(\rho\succeq0,h\ge0,\operatorname{tr}\rho+h=1\); effects are \((E,e)\),
with pairing \(\operatorname{tr}(E\rho)+eh\). Embed reachable
preparations as \((\rho,0)\). Extend each branch of a complete
instrument by

\[
\widetilde\Phi_y(\rho,h)=(\Phi_y(\rho),w_yh),
\qquad w_y\ge0,\qquad\sum_yw_y=1.
\]

Each branch is completely positive and their sum preserves trace. Along
every reachable path the hidden component remains zero and the visible
component remains the original branch. Induction on length, followed by
the direct sum of terminal labels, therefore preserves every record.
Parallel composition, conditional choices and coarse graining preserve
the same conclusion.

This extension concerns the registered elementary registers and the
circuits they generate; it does not identify the systemwise direct-sum
construction with a strong monoidal functor. Complete each nontrivial
elementary register as above, while the empty register, or tensor unit,
remains \(\mathbb C\). Joint carriers use the \textbf{actual algebraic
tensor product} of the completed algebras, for example
\((M_d(\mathbb C)\oplus\mathbb C)\otimes(M_e(\mathbb C)\oplus\mathbb C)\).
Its fully visible central block is
\(M_d(\mathbb C)\otimes M_e(\mathbb C)\); the direct sum of all other
central blocks forms the hidden part and must not be silently replaced
by one copy of \(\mathbb C\). In general, write the input and output as
\(\widetilde{\mathcal A}=\mathcal A_{\rm vis}\oplus\mathcal A_{\rm hid}\)
and
\(\widetilde{\mathcal B}=\mathcal B_{\rm vis}\oplus\mathcal B_{\rm hid}\).
Discarding is the sum of the ordinary matrix traces over their central
blocks. With no input or output register, the corresponding visible
algebra is \(\mathbb C\) and its hidden part is zero.

Any original complete CP instrument \(\{\Phi_y\}\) with multiple input
and output registers can act as originally specified on the visible
joint block, with the extension

\[
\begin{aligned}
\widetilde\Phi_y(\rho\oplus\xi)
&=\iota_{\mathcal B}\Phi_y(\rho)+\Xi_y(\xi),\\
\Xi_y(\xi)&=w_y\operatorname{Tr}(\xi)\tau_{\mathcal B},
\quad w_y\ge0,\quad\sum_yw_y=1.
\end{aligned}
\]

where \(\iota_{\mathcal B}\) embeds the visible block and
\(\tau_{\mathcal B}\) is any fixed normalized state of the actual output
algebra; for a tensor-unit output choose \(\tau_{\mathcal B}=1\). Every
branch is CP, and their sum preserves the full input trace. Identity
operations may instead use the identity extension on the hidden part,
and discarding uses the full total trace. Extensions are chosen only for
the registered generating operations; sequential, parallel and
conditional circuits are then built by actual CP composition and actual
algebraic tensor products. From visible preparations, \(\xi=0\), so
every such circuit retains its original visible branches and
probabilities. The result establishes operational equivalence of
reachable circuits. It does not assert that every original operational
equation also holds on the added unreachable sectors, or that
\(\widehat{\mathcal A\otimes\mathcal B}=\widehat{\mathcal A}\otimes\widehat{\mathcal B}\).
The completed theory may additionally allow preparations of hidden
states, but they are outside the reachable preparation library fixed by
this proposition.

In contrast, the completion admits \((\rho/2,1/2)\), which occupies two
central sectors. A pure joint state must occupy one central block:
occupation of multiple blocks would give a nontrivial convex
decomposition by central projections. Discarding any finite direct-sum
matrix auxiliary system preserves the definite left central label and
cannot produce that mixture. Universal purification therefore fails in
the full completed theory. \(\square\)

This differs from chapter 54's finite Core classical countermodel.
Starting from any specified reachable complex-quantum experiment, it
preserves every circuit response generated by that library while
changing a global property of the full theory. If a UEOT object is
already defined by the \textbf{minimal reachable quotient} of all
allowed future responses, the additional sector may simply be discarded
representation redundancy. The result accepts that minimization. It
neither establishes hidden sectors nor rules out a minimal
complex-quantum realization. It excludes an inference from the present
reachable data directly to the global structure.

\begin{longtable}[]{@{}
  >{\raggedright\arraybackslash}p{(\columnwidth - 6\tabcolsep) * \real{0.2100}}
  >{\raggedright\arraybackslash}p{(\columnwidth - 6\tabcolsep) * \real{0.2000}}
  >{\raggedright\arraybackslash}p{(\columnwidth - 6\tabcolsep) * \real{0.2800}}
  >{\raggedright\arraybackslash}p{(\columnwidth - 6\tabcolsep) * \real{0.3100}}@{}}
\toprule\noalign{}
\begin{minipage}[b]{\linewidth}\raggedright
Candidate
\end{minipage} & \begin{minipage}[b]{\linewidth}\raggedright
Complete basic tests and composition
\end{minipage} & \begin{minipage}[b]{\linewidth}\raggedright
Universal purification existence
\end{minipage} & \begin{minipage}[b]{\linewidth}\raggedright
Local tomography
\end{minipage} \\
\midrule\noalign{}
\endhead
\bottomrule\noalign{}
\endlastfoot
Finite classical simplex & Satisfied & Fails for nontrivial mixtures &
Satisfied \\
Standard real quantum theory & Satisfied & Satisfied & Excluded by
chapter 55's witness \\
Standard finite complex quantum theory & Satisfied & Satisfied &
Satisfied \\
Nontrivial direct sums of complex matrix algebras & Satisfied & Fails
for cross-central mixtures & Satisfied \\
\end{longtable}

Completeness of access to all physical sectors, purification for all
systems, and the full reconstruction background still need independent
grounds. Local tomography alone does not eliminate the final candidate.
The CDP reconstruction also uses the full background of causality,
perfect distinguishability, ideal compression, local distinguishability
and pure conditioning; this chapter's finite interface cannot replace
all of those hypotheses.
\href{https://arxiv.org/abs/1011.6451}{Chiribella, D'Ariano and
Perinotti, 2011}. The next foundational route must first decide whether
selection is up to operational equivalence or concerns global structure,
and then identify an independent source for each new principle. The
present selection verdict is complete; Core's unique selection of
complex quantum structure remains open.

\section{70.3 From finite legal tests to complete external-channel
error}

Once the three durations, instruments, feedback and stopping rules are
fixed, the persistent-bath experiment is a complete recorded channel
whose input and output systems are both two-level systems. Its
inaccessible internal bath is not the carrier that must then be
tomographed. This distinction permits finite ordinary preparations and
readouts to control every untouched reference. It does not permit reuse
of the same certification after changing an intermediate intervention.

\phantomsection\label{p-tom-01}{} \textbf{P-TOM-01 (An explicit
finite-tester stability bound and its preparation, readout and dimension
obligations).} For \(d\ge2\), choose the following \(d^2\) legal pure
preparations and binary-effect settings:

\[
P_j=|j\rangle\langle j|,\quad
P^x_{jk}=\frac{(|j\rangle+|k\rangle)(\langle j|+\langle k|)}2,\quad
P^y_{jk}=\frac{(|j\rangle+i|k\rangle)(\langle j|-i\langle k|)}2
\quad(j<k).
\]

Each readout setting is \(\{P_q,I-P_q\}\). The projectors are not the
outcomes of one common POVM. Define

\[
X_{jk}=|j\rangle\langle k|+|k\rangle\langle j|,\qquad
Y_{jk}=-i|j\rangle\langle k|+i|k\rangle\langle j|.
\]

Their Hermitian duals are

\[
F^x_{jk}=X_{jk},\quad F^y_{jk}=Y_{jk},\quad
F_j=P_j-\frac12\sum_{k\ne j}
\bigl(X_{\min(j,k),\max(j,k)}+Y_{\min(j,k),\max(j,k)}\bigr).
\]

Direct calculation gives \(\operatorname{tr}(P_qF_r)=\delta_{qr}\). For
two instruments with the same input/output dimensions and complete label
set, define

\[
\delta p_{yqr}=\operatorname{tr}
\bigl[P_q(\Phi_y-\Psi_y)(P_r)\bigr].
\]

Their complete CQ channels obey

\[
\frac12\|\Phi_{\rm CQ}-\Psi_{\rm CQ}\|_\diamond
\le\frac12\sum_r\|F_r^{(A)}\|_\infty
\sum_{y,q}\|F_q^{(B)}\|_1|\delta p_{yqr}|.
\tag{70-T}
\]

\textbf{Proof.} Apply the dual expansion once on the input and once on
the output. For every finite reference \(R\) and every matrix \(Z\),

\[
(\Delta_y\otimes\mathrm{id}_R)(Z)
=\sum_{r,q}\delta p_{yqr}F_q^{(B)}\otimes
\operatorname{tr}_A[(F_r^{(A)}\otimes I_R)Z].
\]

Trace-norm duality yields

\[
\|\operatorname{tr}_A[(F\otimes I_R)Z]\|_1
\le\|F\|_\infty\|Z\|_1.
\]

The triangle inequality, multiplicativity of the tensor trace norm and
additivity over orthogonal record blocks then give the claim after
taking the supremum. The quantifiers include every untouched reference
and joint input, rather than a single Bell input. \(\square\)

Set

\[
\begin{aligned}
A_d&=\frac d2(1+\sqrt{2d-1})+d(d-1),\\
B_d&=d\sqrt{2d-1}+2d(d-1),\\
\delta&=\max_{q,r}\sum_y|\delta p_{yqr}|.
\end{aligned}
\]

The simpler upper bound is \(A_{d_A}B_{d_B}\delta/2\), with sufficient
constants:

\begin{longtable}[]{@{}
  >{\raggedright\arraybackslash}p{(\columnwidth - 4\tabcolsep) * \real{0.2500}}
  >{\raggedleft\arraybackslash}p{(\columnwidth - 4\tabcolsep) * \real{0.3400}}
  >{\raggedleft\arraybackslash}p{(\columnwidth - 4\tabcolsep) * \real{0.4100}}@{}}
\toprule\noalign{}
\begin{minipage}[b]{\linewidth}\raggedright
\(d_A=d_B\)
\end{minipage} & \begin{minipage}[b]{\linewidth}\raggedleft
Preparation × readout settings
\end{minipage} & \begin{minipage}[b]{\linewidth}\raggedleft
Half-diamond coefficient
\end{minipage} \\
\midrule\noalign{}
\endhead
\bottomrule\noalign{}
\endlastfoot
2 & 16 & \(9+5\sqrt3\approx17.660254\) \\
3 & 81 & \((225+81\sqrt5)/4\approx101.530377\) \\
4 & 256 & \(196+52\sqrt7\approx333.579068\) \\
\end{longtable}

These constants follow from the actual spectra of the duals and are not
claimed optimal. The source code checks the dual identities and spectral
polynomials in exact arithmetic for dimensions 2--4. The fixed
three-slot protocol has external dimension two, so four calibrated
preparations and four binary effects, with all four terminal records
retained, can use the first row. Internal-bath tomography and
preparation of a reference-entangled experimental state are not
additional requirements for this bound. Dimension and trusted
preparation/readout ranges remain independent premises.

If preparation error has trace distance at most \(\tau_r\), and each
actual effect differs from its nominal effect in operator norm by at
most \(\xi_q\), then

\[
\sum_y|\widetilde p_{yqr}-p_{yqr}|\le2\tau_r+\xi_q.
\tag{70-SPAM}
\]

The preparation term follows from contraction by the complete POVM. The
readout term uses the unit total mass of the subnormalized label blocks.
A simultaneous statistical radius, the actual fit residual after
projection onto physical maps, and these SPAM terms must all enter
(70-T). A common gauge transformation of unknown preparations and
readouts can preserve every probability, so an absolute channel distance
in specified coordinates needs an independent physical frame. This is
consistent with the gauge issue in standard gate-set tomography.
\href{https://arxiv.org/abs/2009.07301}{Nielsen et al., 2021}.

The finite-dimensional assumption is equally essential: coherent leakage
probability \(\ell\) can cause state trace distance \(\sqrt\ell\), so
adding only \(\ell\) is invalid. Whether a leakage flag is physically
measured must be declared. Postselection also requires its own
acceptance-probability bound: two processes at full-record distance
\(q\) can have accepted states at distance one on an event of
probability \(q\).

Complete tomography is not automatically economical. With four labels,
99\% coverage and the whole error budget 0.05 assigned to entrywise
Hoeffding/union bounds, one conservative sufficient budget for this
frame is 9,438,576 shots per setting, or 151,017,216 in total. This is
not a sample lower bound or an experiment executed in this chapter. The
present structured two-parameter calibration does not replace complete
channel tomography; likewise, fixed-protocol tomography does not replace
process-tensor certification for arbitrary intermediate interventions.

\section{70.4 The common multitime source with a persistent bath}

\subsection{Source, records and quantifiers}

The same two-mode environment is retained until each branch actually
stops. The system S is a qubit; the initial state \(\rho_{SR}\) is
arbitrary, and the finite-dimensional reference R is never acted on
during the dynamics. Set \(\hbar=1\) and \(H_S=0\), and fix the source
below:

\[
\begin{aligned}
H_\lambda&=I_S\otimes H_E+\lambda Z\otimes B,
&H_E&=n_1+\frac32n_2,\\
B&=(a_1+a_1^\dagger)+\frac35(a_2+a_2^\dagger),
&\lambda&\in[1/5,2/5],\\
\rho_{SRE}(0)&=\rho_{SR}\otimes|00\rangle\langle00|_E,
&(t_1,t_2,t_3)&=(2/3,3/4,1/2).
\end{aligned}
\]

After every joint evolution, apply the complete two-branch instrument
\(F=\operatorname{diag}(\sqrt{1/8},\sqrt{1/4})\) and
\(P=\operatorname{diag}(\sqrt{7/8},\sqrt{3/4})\). Failure stops
immediately and retains its actual unnormalized SR state. Pass branches
receive \(R_y(\pi/2)\), \(R_x(\pi/3)\) and I, respectively. Every
terminal record f,pf,ppf,ppp is retained. Distance is half the sum of
the trace norms of the record-block differences; E is discarded only at
the stipulated terminal output. Complex Hilbert space, the Born rule, CP
operations, oscillator dynamics and vacuum preparation remain
independent physical inputs, not structures uniquely derived from Core.

\phantomsection\label{p-micro-08}{} \textbf{P-MICRO-08 (A complete
persistent Gaussian-bath process on a common domain).} We first
construct a legitimate common source and then derive its full multitime
response. General finite Gaussian states specify the scope of the
analytic formula only. Every executed, calibrated and numerically
certified source in this section starts in the frozen vacuum above.

\subsection{Field bounds and a common self-adjoint domain}

For finitely many positive-frequency modes, define
\(a(g)=\sum_k g_k^*a_k\) and \(b_g^2=\sum_k|g_k|^2/\omega_k\).
Cauchy--Schwarz and the CCR give, initially on finite Fock vectors,

\[
\begin{aligned}
\|a(g)\psi\|&\le b_g\|H_E^{1/2}\psi\|,\\
\|a(g)^\dagger\psi\|^2
&=\|a(g)\psi\|^2+\|g\|_2^2\|\psi\|^2,\\
\|(a(g)+a(g)^\dagger)\psi\|
&\le2b_g\|H_E^{1/2}\psi\|+\|g\|_2\|\psi\|.
\end{aligned}
\]

The spectral inequality \(\sqrt{x}\le\epsilon x+(4\epsilon)^{-1}\) turns
the last line into an arbitrarily small relative \(H_E\) bound plus a
bounded term. The field therefore has relative bound zero. The system
operator is bounded, so Kato--Rellich gives a self-adjoint realization
on \(\mathbb C^2\otimes D(H_E)\). A compact finite-coupling set and a
positive frequency lower bound make these estimates uniform. Here
completing the square also constructs the unitary explicitly. With
\(\omega=(1,3/2)\) and \(c=(1,3/5)\), the \(Z\)-sector \(z=\pm1\)
satisfies

\[
\begin{aligned}
U_z(t)&=e^{i\phi(t)}D(z\alpha(t))e^{-itH_E},\\
\alpha_k(t)&=\frac{\lambda c_k}{\omega_k}(e^{-i\omega_kt}-1),\\
\phi(t)&=\sum_k\left(\frac{\lambda c_k}{\omega_k}\right)^2
[\omega_kt-\sin(\omega_kt)],\\
D(\zeta)&=\exp\!\left[\sum_k(\zeta_ka_k^\dagger-\zeta_k^*a_k)\right].
\end{aligned}
\]

The identity \(a_kD(\zeta)=D(\zeta)(a_k+\zeta_k)\) shows that finite
displacement preserves \(D(H_E)\); the additional field terms are
controlled by the preceding bound. This explicit unitary agrees with the
self-adjoint source, rather than defining dynamics through an unproved
formal series.

\subsection{Energy with complete stopping records}

Dilate the instruments into a finite record register. After a branch
stops, let its bath evolve freely until the common horizon \(T=23/12\).
This padding changes neither its terminal record nor its SR output.
System and record operations preserve total number \(N=n_1+n_2\), and
the live-branch coupling operator has norm at most one. Put
\(\gamma=\|c\|_2=\sqrt{34/25}\). Vacuum evolves into finitely many
displacement histories with the required number moments. Differentiation
on this actual domain and Cauchy--Schwarz give

\[
\left|\frac{d}{dt}\langle N\rangle\right|
\le2|\lambda|\gamma\sqrt{\langle N\rangle},
\qquad
\sqrt{\langle N(t)\rangle}\le|\lambda|\gamma t.
\]

The zero initial value is handled by first regularizing with
\(\sqrt{\langle N\rangle+\epsilon}\) and then taking \(\epsilon\) to
zero. This is a complete-record, unconditioned energy bound. Dividing a
branch by a small acceptance probability does not preserve it without a
cost. For a separately declared finite-energy Gaussian preparation, the
right-hand side becomes
\(\sqrt{\langle N(0)\rangle}+|\lambda|\gamma t\); this does not extend
the present truncation or numerical certificate to nonvacuum input.

\subsection{Complete history pairs for a correlated Gaussian
preparation}

For a system eigenvalue history \(s=(z_1,\ldots,z_n)\), let
\(R_j=\operatorname{diag}(e^{-i\omega_kt_j})\). Starting at
\(\beta_0=0,\theta_0=0\), recurse

\[
\begin{aligned}
\beta_j&=R_j\beta_{j-1}+z_j\alpha(t_j),\\
\theta_j&=\theta_{j-1}+\phi(t_j)
+\operatorname{Im}\langle R_j\beta_{j-1},z_j\alpha(t_j)\rangle,\\
W_s&=U_{z_n}(t_n)\cdots U_{z_1}(t_1)
=e^{i\theta_s}D(\beta_s)e^{-iH_ET_n}.
\end{aligned}
\]

The second line retains the exchange phase from Weyl multiplication.
After noncommuting intermediate operations, this phase generally cannot
be discarded.

Use quadratures \(r=(q_1,p_1,q_2,p_2)\), with
\(q_k=(a_k+a_k^\dagger)/\sqrt2\), \(p_k=(a_k-a_k^\dagger)/(i\sqrt2)\)
and \(\Omega=\bigoplus_{k=1}^2\begin{pmatrix}0&1\\-1&0\end{pmatrix}\).
Let d be the mean and V the symmetric covariance, with
\(V+i\Omega/2\succeq0\). V is the full matrix and may contain intermode
correlations. The Gaussian characteristic function and history-pair
factor are

\[
\begin{aligned}
\xi(\zeta)&=\sqrt2(\operatorname{Im}\zeta_1,-\operatorname{Re}\zeta_1,
\operatorname{Im}\zeta_2,-\operatorname{Re}\zeta_2)^T,\\
\chi_\sigma(\zeta)&=\operatorname{Tr}[\sigma D(\zeta)]
=e^{i\xi(\zeta)^Td-\frac12\xi(\zeta)^TV\xi(\zeta)},\\
F_{s,r}&=\operatorname{Tr}(W_s\sigma_EW_r^\dagger)\\
&=e^{i[\theta_s-\theta_r+\operatorname{Im}\langle\beta_r,\beta_s\rangle]}
\chi_\sigma\!\left(e^{i\omega T_n}(\beta_s-\beta_r)\right).
\end{aligned}
\]

This follows from cyclicity of the trace, the Weyl phase in
\(D(-\beta_r)D(\beta_s)\), and free rotation. No product decomposition
of a correlated initial bath is made, and no conditional bath is assumed
Gaussian.

For terminal record h, let \(M_{h,j}\) denote the actual failure
operator or rotated pass operator, and define
\(A_{h,s}=M_{h,n}\Pi_{z_n}\cdots M_{h,1}\Pi_{z_1}\). The full output is

\[
\rho_h^{SR}=\sum_{s,r}F_{s,r}
(A_{h,s}\otimes I_R)\rho_{SR}(A_{h,r}^\dagger\otimes I_R).
\]

This is the expansion of the actual joint operator
\(\sum_sA_{h,s}\otimes W_s\) followed by a partial trace, so it is CP;
\((F_{s,r})\) is also a positive Gram matrix. Summing outcomes
successively and using \(F^\dagger F+P^\dagger P=I\) proves unit total
trace and causal prefix constraints. The reference is retained
throughout. Iterating \((3/4)I\preceq P^\dagger P\preceq(7/8)I\) bounds
surviving mass and yields \(27/64\le\Pr(ppp)\le343/512\). This neither
makes the accepted effect scalar nor makes normalized postselection
linear.

Even vacuum does not support a single-Gaussian conditional closure.
After the first pass with Bell input, the bath marginal mixes opposite
coherent states with weights 7/13 and 6/13. With
\(\mu=\sqrt2|\alpha_1(t_1)|\), its displacement-aligned fourth cumulant
is \(\mu^4(-2+8/169-6/28561)\), strictly nonzero for nonzero
displacement. Gaussian preparation and fully Gaussian conditional
dynamics are therefore different assumptions.

\section{70.5 Genuine finite propagation error and parameter stability}

\subsection{Finite Duhamel expansion and the all-input lift}

\phantomsection\label{p-micro-09}{} \textbf{P-MICRO-09 (An all-input
finite approximation of the same persistent source).} Let \(P_K\)
project onto occupation at most K in each mode, let \(J_K\) be its
inclusion, and use the actual compression
\(H_{\lambda,K}=J_K^\dagger H_\lambda J_K\). Free bath evolution and
every dilated record operation preserve total number. In the interaction
picture including these operations, write the field as
\(\widetilde V(s)\). On total number at most \(k-1\), its norm is
bounded by \(\gamma(\sqrt{k-1}+\sqrt{k})\).

The essential Duhamel orientation keeps true propagation on the left:

\[
W(T,0)=I-i\lambda\int_0^T W(T,s)\widetilde V(s)\,ds.
\]

Expand the left suffix only finitely many times. In the order-j
remainder, the true unitary remains to the left of every field factor,
while the right-hand input has finite support originating in vacuum.
Each field increases total number by at most one. Thus every term
through order K stays inside \(N\le K\), itself contained in the
rectangular carrier \(P_K\); the first K terms of the infinite and
finite sources agree exactly. Projections in the finite remainder cannot
increase norms. With \(\lambda_{\max}=2/5\), the ordered time-simplex
volume bounds the two order-\(K+1\) remainders separately:

\[
\begin{aligned}
x&=\lambda_{\max}\gamma T,
&c_j&=\frac{x^j}{j!}\prod_{k=1}^j(\sqrt{k-1}+\sqrt{k}),\\
\|(W-J_KW_KJ_K^\dagger)\psi_{\rm vac}\|
&\le c_{K+1}+c_{K+1},
&d_K&=\min\{1,2c_{K+1}\}.
\end{aligned}
\]

No convergence of an infinite Dyson series is assumed, and no
uncontrolled infinitely evolved state is placed to the right of a field
operator. Purify an arbitrary input and retain every instrument record
in the dilation. The operator estimate is unchanged by tensoring the
identity on any reference. Pure-state trace distance is at most the
corresponding vector distance; discarding E and unused dilation
registers then gives

\[
\begin{gathered}
\sup_{\rho_{SR},\,\dim R<\infty}
\frac12\sum_h\left\|
(\Phi_{\lambda,h}\otimes\operatorname{id}_R)(\rho_{SR})-
(\Phi_{\lambda,K,h}\otimes\operatorname{id}_R)(\rho_{SR})
\right\|_1\le d_K,\\
d_{20}<0.00001344.
\end{gathered}
\]

This is an analytic statement for arbitrary SR input, not an
extrapolation from numerical Bell-input agreement. The same bound
applies when the first two durations are interchanged, with the same
total exposure and legal system operations.

\subsection{Boundary residuals on actual finite prefixes}

A small initial tail does not establish accurate propagation by a
compressed generator. Here only creation operators cross the rectangular
boundary. The two outgoing faces are orthogonal, hence

\[
\begin{aligned}
\mathcal R_K&=H_\lambda J_K-J_KH_{\lambda,K},\\
\mathcal R_K^\dagger\mathcal R_K
&=\lambda^2(K+1)I_S\otimes
\left[\Pi_{n_1=K}+\frac9{25}\Pi_{n_2=K}\right].
\end{aligned}
\]

Let \(\sigma_h\) be the actual unnormalized live prefix of the same
finite model, with mass \(m_h\), and let \(\sigma_h(s)\) be its local
finite evolution. After purifying a subnormalized state, the Duhamel
vector-defect integral must be multiplied by its norm \(\sqrt{m_h}\) to
bound trace distance. Contractivity of the true complete suffix and
telescoping over the tree give

\[
r_h(s)=\operatorname{Tr}\!\left[
(\mathcal R_K^\dagger\mathcal R_K\otimes I_R)\sigma_h(s)\right],
\qquad
D_{\rm full}\le\sum_{h\,\mathrm{live}}\int_0^{t_h}\sqrt{m_h r_h(s)}\,ds.
\]

The occupations and masses must belong to the actual finite prefixes,
not be replaced by infinite-source occupations. This formula can sharpen
the result, but ordinary numerical quadrature is not called an interval
certificate here. The uniform certificate actually used is \(d_K\)
above.

The frozen vacuum has zero initial tail; no mode is deleted and no
Trotter split is used. Changing the preparation requires separate
initial-tail and propagation costs; correlated modes cannot be replaced
by a freshly prepared product vacuum. Mode deletion should switch off
its coupling on the same Hilbert space and bound the coupling difference
on the actual retained initial state and finite domain. Device
preparation, control, leakage, effective-model and readout errors still
require independent evidence.

\subsection{Parameter stability is first proved in the infinite source}

For \(\lambda,\widehat\lambda\in[1/5,2/5]\), compare the same controller
and stopping rules. The parameter Duhamel formula inserts
\((\lambda-\widehat\lambda)Z\otimes B\) with a true infinite-source
prefix on its right. Its energy estimate must therefore come from the
infinite source, rather than an assumed invariant of the finite
compression. For a normalized purified prefix,

\[
\begin{aligned}
\|B\psi(t)\|&\le\gamma\left(\sqrt{\langle N\rangle}
+\sqrt{\langle N\rangle+1}\right)
\le\gamma(1+2\lambda_{\max}\gamma t),\\
D_{\rm full}(\Phi_\lambda,\Phi_{\widehat\lambda})
&\le|\lambda-\widehat\lambda|
\int_0^T\gamma(1+2\lambda_{\max}\gamma t)\,dt\\
&=|\lambda-\widehat\lambda|
[\gamma T+\lambda_{\max}\gamma^2T^2]
\le\frac{51}{10}|\lambda-\widehat\lambda|.
\end{aligned}
\]

The bracket is approximately \(4.2336426708\). The rational estimate
\(\gamma<7/6\) bounds it by \(38111/9000<51/10\). We retain the
previously frozen \(51/10\). Parameter stability is established first;
finite truncation and numerical-representation errors are added
afterwards. It does not authorize silent changes of frequencies,
relative couplings, preparation or control permissions.

\section{70.6 Memory controls and physical numerical representation}

\subsection{A strict memory witness for the fixed three-slot source}

\phantomsection\label{c-qmem-01}{} \textbf{C-QMEM-01 (Equal single-slot
channels do not determine a multitime process).} Repreparing vacuum
after every slot defines a different process. Its isolated single-slot
channel equals the true source for every input and reference, because
that slot still starts from the same vacuum and undergoes the same
evolution. Yet for the same three-slot Bell input at \(\lambda=3/10\),
the event ``record ppp and \(X_S\otimes X_R=+1\)'' has the strict
interval witness

\[
p_{\rm reset}-p_{\rm persistent}
>0.03164037538529123569>\frac3{100},
\qquad
D_{\rm full}(\Phi_{\rm reset},\Phi_{\rm persistent})>\frac3{100}.
\]

This uses 80-digit interval history pairs and an independently
propagated sequence of reset Schur channels, with exact dyadic endpoints
retained. An event gap bounds trace distance directly; interval
diagonalization of a trace norm is unnecessary, and no different
causal-break protocol is substituted. The complete distance
approximately \(0.1012518562\) remains a floating diagnostic. Omitting
the exchange phase causes approximately \(0.0502254699\) error. Vacuum
has exactly zero initial tail above K=0, but its propagation error at
\(\lambda=0.4\) is approximately \(0.2389578616\). These controls
identify the consequences of separate omissions; they do not replace the
strong standard baseline.

\subsection{Why interval Choi bounds control arbitrary reference inputs}

\phantomsection\label{p-num-01}{} \textbf{P-NUM-01 (From interval Choi
bounds to a physical terminal instrument).} With normalized Bell vector
\(|\Omega_2\rangle=(|00\rangle+|11\rangle)/\sqrt2\), define the four
normalized Choi blocks in output-S followed by input-reference order:

\[
J_h=(\Phi_h\otimes\operatorname{id}_2)
(|\Omega_2\rangle\langle\Omega_2|),
\qquad
J_h\succeq0,
\qquad \sum_h\operatorname{Tr}_{\rm out}J_h=I_2/2.
\]

For generator parameter 3/10 and calibration-only nominal parameter
0.299314119192, independent 80-digit interval formulas enclose every
entry for both timing orders. Each timing separately certifies the
actually stored binary64 matrices from the Weyl and K20 Galerkin
implementations, giving eight certificates. Convert the stored matrices
to exact rationals and Hermitianize them, calling them \(Q_h\).
Entrywise intervals bound Frobenius error, and each block has dimension
four, so

Because the target \(J_h\) is Hermitian, exact Hermitianization is an
orthogonal projection in the Frobenius inner product and does not
increase the Frobenius error relative to \(J_h\).

\[
\begin{aligned}
e&:=\sum_h\|Q_h-J_h\|_1
\le2\sum_h\|Q_h-J_h\|_F<10^{-12},\\
\frac12\|\mathcal Q-\Phi\|_\diamond
&\le\frac{d_{\rm in}}2\sum_h\|Q_h-J_h\|_1=e,
\qquad d_{\rm in}=2.
\end{aligned}
\]

The second line follows directly by vectorization. Any unit rank-one
input \(|u\rangle\langle v|\) with an arbitrary reference can be written
as \((I\otimes A)|\Omega_2\rangle\langle\Omega_2|(I\otimes B^\dagger)\),
with \(\|A\|,\|B\|\le\sqrt2\). Left and right multiplication therefore
increase each block trace norm by at most two. A singular-value
decomposition extends the bound to any trace-class input, proving the
diamond inequality. This certifies the stored linear map for every
reference input; it does not extrapolate one floating Bell-input
agreement.

The intervals compare the exact infinite source directly with the actual
stored matrices. For the Galerkin backend, this fixed-nominal-parameter
difference already contains truncation and floating matrix-exponential
errors; it is not merely a rounding estimate. In contrast, \(d_{20}\) is
an independent analytic bound uniform over the entire frozen parameter
interval.

\subsection{Exact CP/TP repair and its term-by-term cost}

The raw \(Q_h\) need not be positive or trace preserving. Choose
rational \(\delta=10^{-10}\) and \(a=(1+20\delta)^{-1}\), and define

\[
\widetilde Q_h=Q_h+\delta I_4,
\qquad
R=I_2/2-a\sum_h\operatorname{Tr}_{\rm out}\widetilde Q_h.
\]

Each block has operator error at most \(e<\delta\), hence
\(\widetilde Q_h\succeq0\). Trace-norm contractivity of partial trace on
Hermitian inputs bounds the aggregate marginal error by \(e<2\delta\).
The four added \(\delta I_4\) blocks contribute \(8\delta I_2\) after
partial trace. Therefore

\[
a\sum_h\operatorname{Tr}_{\rm out}\widetilde Q_h
\preceq a(1/2+10\delta)I_2=I_2/2,
\qquad R\succeq0.
\]

Set \(J_h^{\rm rep}=a\widetilde Q_h\) for every block, then add
\(|0\rangle\langle0|_{\rm out}\otimes R\) only to the existing f block.
Every block is positive and the total output marginal is exactly
\(I_2/2\), defining a CP/TP terminal instrument. Matrices, a, delta and
R are all stored as exact Fractions. Trace preservation and Hermiticity
are verified entrywise with rational arithmetic. Separate review
additionally checked all 504 principal minors of the eight certificates'
blocks and residuals exactly.

To expose the repair cost, let \(\Delta_h=Q_h-J_h\) and
\(s=\sum_h\operatorname{Tr}\Delta_h\), so \(|s|\le e\). Since
\(R\succeq0\) and \(\sum_h\operatorname{Tr}J_h=1\), term-by-term
triangle bounds give

\[
\begin{aligned}
\sum_h\|J_h^{\rm rep}-J_h\|_1
&\le ae+(1-a)+16a\delta+\operatorname{Tr}R,\\
\operatorname{Tr}R&=1-a(1+s+16\delta),\\
\sum_h\|J_h^{\rm rep}-J_h\|_1
&\le ae+2(1-a)-as\\
&\le2ae+40a\delta
=a(40\delta+2e)<10^{-8},\\
\frac12\|\Phi^{\rm rep}-\Phi\|_\diamond&<10^{-8}.
\end{aligned}
\]

The term \(16a\delta\) comes from four four-dimensional blocks;
\(2(1-a)=40a\delta\) comes from normalization scaling. No
input-dimension factor is missing. This is explicitly paid
physicalization, rather than treating raw ``nearly CP'' matrices as CP.

The repaired object is a terminal prediction instrument for one fixed
policy. It is not claimed implementable at each original early stopping
time, nor to repair a complete process tensor for all counterfactual
interventions. Its changed failure block is included in the error bound.
A later floating contraction for arbitrary input requires its own
rounding budget; the fixed Bell-input Pauli readout here can instead be
evaluated exactly from the rational matrices.

The frozen calibration certificate retains the conservative budget

\[
\varepsilon_{\rm CQ}
=\frac{51}{10}r_\lambda+d_{20}+\varepsilon_{\rm phys},
\qquad
\varepsilon_{\rm tester}
=\varepsilon_{\rm CQ}+\frac12r_\eta,
\qquad \varepsilon_{\rm phys}=10^{-8}.
\]

At the current nominal parameter the interval certificate already
directly bounds the infinite-source-to-repaired-predictor difference.
Paying \(d_{20}\) again is a conservative reserve, not a second
independent discovery of the same error. The term \(r_\eta/2\) applies
only to the declared complete binary terminal detector model, not to
arbitrary device SPAM.

Standard persistent joint evolution and the same Gaussian influence
functional predict the same physics. This section connects explicit
premises, complete records, actual finite propagation and auditable
numerical representation, without claiming additional empirical
advantage over standard quantum theory. Infinite modes, local field
theory, energy exchange and thermalization retain separate UV/IR,
representation, causal-composition and limit obligations.

\section{70.7 Persistent-memory lower bounds after a causal break}

\phantomsection\label{p-qmem-02}{} \textbf{P-QMEM-02 (A necessary memory
dimension from causal-break response rank).} At a fixed temporal cut, if
the Hilbert dimension of every history-carrying register crossing the
cut is jointly \(D\), the response matrix for histories \(h\) and future
events \(t\) satisfies

\[
H_{ht}=\operatorname{tr}(F_t\sigma_h),\qquad
\operatorname{rank}_{\mathbb R}H\le D^2.
\tag{70-M}
\]

\textbf{Proof.} Expand Hermitian states and effects in the same real
basis of dimension \(D^2\). The matrix \(H\) factors through this real
vector space and therefore has rank at most \(D^2\). \(\square\)

Future operations must have the same implementation for every history.
If a controller or simulator can read the past label, that log is part
of the memory crossing the cut. A varying accessible reference also
belongs in \(D\). When \(S\) is genuinely discarded and independently
prepared in the same state, with no varying reference crossing the cut,
only \(M\) carries history and
\(d_M\ge\lceil\sqrt{\operatorname{rank}H}\rceil\). This is a necessary
bound, not a theorem implementing CP compression.

For the same two-mode vacuum source, declare a separate diagnostic:
prepare \(S\) in \(Z=s\), \(s=\pm1\), evolve for \(t=2/3\); discard and
prepare \(|+\rangle\), evolve for \(\tau=3/4\); finally measure \(Y\).
Future operations do not read \(s\), the reference is unused, and the
bath is not reset. With \(\kappa_k=\lambda c_k/\omega_k\), the exact
coherence is

\[
f_s=v e^{is\theta},\qquad
v=\exp\!\left[-4\sum_k\kappa_k^2(1-\cos\omega_k\tau)\right],
\] \[
\theta=4\sum_k\kappa_k^2
\left[\sin\omega_k t+\sin\omega_k\tau-\sin\omega_k(t+\tau)\right].
\]

The first slot writes opposite coherent displacements into the bath; the
second reads that history through the Weyl exchange phase. Hence

\[
H=\frac12
\begin{pmatrix}1-b&1+b\\1+b&1-b\end{pmatrix},\qquad
b=v\sin\theta.
\]

Using \(\sin x+\sin y-\sin(x+y)=4\sin(x/2)\sin(y/2)\sin((x+y)/2)\), the
full interval \(\lambda\in[1/5,2/5]\) obeys \(0<\theta<2.2272<\pi\).
Thus \(b>0\), the response rank is two, and every finite quantum-memory
realization requires dimension at least two. This does not show that
two-dimensional memory suffices for all bath histories.

At \(\lambda=0.2,0.3,0.4\), the second singular values are approximately
0.0686433, 0.1432226 and 0.2280272. Independent \(K=14\) mode-matrix
calculations differ from the analytic coherences by less than
\(3.35\times10^{-16}\) in floating arithmetic. Nonvanishing throughout
the interval follows from the analytic inequality, not extrapolation
from those three grid points. A device-level memory-rank claim still
needs a simultaneous confidence region and SPAM bounds excluding every
rank-one response; uncontrolled history-dependent drift can imitate this
signal.

\section{70.8 Joint identification, finite-sample coverage, and
same-source certification}

\phantomsection\label{p-rob-06}{} \textbf{P-ROB-06 (joint identification
and coverage from two trusted settings).} For the declared two-mode
vacuum source, let \(\lambda\in[1/5,2/5]\) be the common coupling scale
and \(\eta\in[9/10,1]\) a fixed symmetric visibility shared by the
declared binary readouts. A fresh trial prepares \(|+X\rangle\), evolves
for t, and measures X. Its observed mean is

\[
\mu_t=\eta e^{-a(t)\lambda^2},\qquad
a(t)=4\left[1-\cos t+\frac4{25}(1-\cos(3t/2))\right].
\]

The zero-time setting and \(t=3/2\) therefore determine

\[
\eta=\mu_0,\qquad
\lambda=\sqrt{\frac{\log\mu_0-\log\mu_{3/2}}{a(3/2)}}.
\]

The identification domain requires positive visibility, the nonnegative
coupling convention, known frequencies and relative couplings, vacuum
preparation, and trusted preparation/readout structure. A single nonzero
time admits the exact equivalence curve
\(\eta'=\eta\exp[a(t)(\lambda'^2-\lambda^2)]\); more shots at that
setting cannot remove the degeneracy. With an unknown common thermal
occupation n, all single-slot coherences identify only
\((2n+1)\lambda^2\). An unknown detector bias or unrestricted
setting-dependent drift requires additional independent anchors and an
appropriate error model. These quantities are not silently treated as
identified by the present calibration.

Each calibration setting generated 200,000 synthetic independent trials.
For the plus-outcome frequency, use the Hoeffding radius

\[
r_p=\sqrt{\frac{\log(4/0.01)}{2\cdot200000}}+10^{-9},
\]

and transform it into expectation intervals \(I_0,I_1\). The joint set
is

\[
C(D)=\{(\lambda,\eta)\in[.2,.4]\times[.9,1]:
\eta\in I_0,\;\eta e^{-a(3/2)\lambda^2}\in I_1\}.
\]

Under the declared independent stationary sampling assumptions, coverage
is at least 99\%. The \(10^{-9}\) term is a numerical
sampling-probability allowance frozen in advance. An independent
interval implementation checked all 98 probabilities actually used by
the generator; the largest discrepancy was below
\(1.011\times10^{-16}\). Coverage is a statistical theorem: one
fixed-seed synthetic run does not empirically establish the coverage
frequency.

The realized joint set is nonempty. Calibration alone selected

\[
\widehat\lambda=0.299314119192,\qquad
\widehat\eta=0.959480000000.
\]

Outward projections satisfy

\[
\begin{aligned}
\lambda&\in[0.292143974948,0.306484263435],\\
|\lambda-\widehat\lambda|&\le0.007170144244,\\
|\eta-\widehat\eta|&\le0.007740457121.
\end{aligned}
\]

Interval arithmetic verifies that the nominal point belongs to the joint
set. The displayed rectangle is only an outer enclosure. The generator's
truth file is read solely by generation and a separate generator audit;
it is not an input to estimation or held-out decisions.

\phantomsection\label{p-rob-07}{} \textbf{P-ROB-07 (a complete CQ
certificate and independent holdout).} Combining the same-source
all-input parameter bound, the K=20 Galerkin bound, and the
interval-validated exact rational CP/TP repair gives

\[
\begin{aligned}
D_{\rm full}
&\le\frac{51}{10}(0.007170144244)+0.00001344+10^{-8}\\
&=0.0365811856444\le0.036581185645<0.05.
\end{aligned}
\]

This inequality retains an arbitrary finite untouched reference and
every \texttt{f,pf,ppf,ppp} record. The pointwise nominal Choi
certificate already compares directly with the infinite source. The
displayed \(K=20\) term is a retained conservative uniform analytic
allowance, not an omitted truncation error of that pointwise
certificate. Preparation, control, and effective-model errors are zero
here as assumptions of the selected synthetic source; corresponding
device certificates have not been established. Normalized successful
conditional states do not inherit this precision without accounting for
the acceptance probability.

For each fixed binary terminal tester, uncertainty in the shared
visibility adds at most \(|\eta-\widehat\eta|/2\). Thus the complete
readout distribution obeys

\[
D_{\rm measured}\le0.0404514142049\le0.040451414205<0.05.
\]

The quantum-output and readout-distribution metrics are recorded
separately. The four terminal observables actually used,
\texttt{XX,XY,YZ,ZZ}, are contracted with the repaired rational Choi
matrices in exact rational arithmetic; an unbudgeted later floating
contraction is not included in this certificate.

The nominal model, statistical rule, and program were hash-bound before
holdout generation. Each family used 100,000 trials for each of the four
observables, retaining all eight stopped-record/sign events:

\begin{longtable}[]{@{}
  >{\raggedright\arraybackslash}p{(\columnwidth - 4\tabcolsep) * \real{0.2500}}
  >{\raggedright\arraybackslash}p{(\columnwidth - 4\tabcolsep) * \real{0.3400}}
  >{\raggedright\arraybackslash}p{(\columnwidth - 4\tabcolsep) * \real{0.4100}}@{}}
\toprule\noalign{}
\begin{minipage}[b]{\linewidth}\raggedright
Holdout family
\end{minipage} & \begin{minipage}[b]{\linewidth}\raggedright
Timing/control change
\end{minipage} & \begin{minipage}[b]{\linewidth}\raggedright
Result
\end{minipage} \\
\midrule\noalign{}
\endhead
\bottomrule\noalign{}
\endlastfoot
Original three-slot deployment & \texttt{(2/\allowbreak 3,3/\allowbreak 4,1/\allowbreak 2)} with
noncommuting intermediate rotations & Frozen contract not rejected \\
Changed intervals & \texttt{(3/\allowbreak 4,2/\allowbreak 3,1/\allowbreak 2)}, same total duration and
controls & Frozen contract not rejected \\
Shared-detector violation & Original timing, generator visibility
eta=.70 & Frozen model contract rejected \\
\end{longtable}

An independent interval history-pair implementation encloses the
continuous parameter set using 288 lambda grid points, the eta
endpoints, and a \(5.1\,\mathrm{mesh}/2\) covering allowance.
Simultaneous confidence intervals cover the 96 events. When the frozen
model holds, the combined calibration/holdout false-rejection
probability is at most \texttt{.\allowbreak 02}. Non-rejection establishes neither
model uniqueness nor device validation.

\section{70.9 Strong baseline, cost, and real-data qualification}

Standard persistent joint evolution, the Gaussian influence functional,
and the process used here describe the same conditional quantum map.
With the same calibration data, parameter set, controls, and robust
rule, their predictions agree. The independent K=20 Galerkin and Weyl
backends differ in the full CQ floating calculation by approximately
\texttt{8.\allowbreak 93e-\allowbreak 16} and \texttt{9.\allowbreak 32e-\allowbreak 16} for the two normal timings;
interval enclosures and exact physical repair separately supply the
numerical guarantee. The method verdict is
\textbf{\texttt{EQUIVALENT\_\allowbreak NO\_\allowbreak ADDED\_\allowbreak GAIN}}. This iteration does not
establish an additional physical prediction or an advantage over a
strong standard quantum method under equal resources.

Resetting the bath between slots changes the full CQ nominal output by
approximately \texttt{.\allowbreak 100895} and \texttt{.\allowbreak 104952}. This confirms that
the ablation loses memory; it is not a substitute for the
strong-baseline comparison and does not reject every possible Markov
model. The retained synthetic data contain 1.6 million trials: 400,000
calibration and 1.2 million holdout trials. No initial projection filter
is used; all failure records of the source itself are retained. The
holdout trials consumed 2,966,811 interaction slots when charged through
their actual stopping prefixes. These are simulated protocol resources,
not hardware calls or laboratory wall time.

A separate, predeclared reproduction used one author-published processed
observable from Zhu et al.,
\texttt{O=X\ tensor\ (2\allowbreak Z+Y)/\allowbreak sqrt(5)\ tensor\ I\ tensor\ .\allowbreak .\allowbreak .\allowbreak }. All
5-by-11 values in the original HDF5 member were retained. The mean-curve
RMSE is \texttt{0.\allowbreak 01860965} against the authors' noisy simulation and
\texttt{0.\allowbreak 03351389} against the ideal simulation on matching grid
points. The original HDF5 member, archive SHA256/MD5, source, and
visible license status are preserved. These are not raw trials for the
oscillator source: the protocol-specific failure, timing, vacuum,
reference, reset, and calibration/holdout records are absent. They
cannot provide the device certificate required by this chapter.

Three conclusions therefore coexist: the declared source has a
mathematical and synthetic certificate; an external processed
experimental observable has been reproduced; the matched device contract
remains \textbf{\texttt{NOT\_\allowbreak CERTIFIED}}. Independent physical selection
and the broader unification task remain open.

\section{70.10 Decisions and the next research boundary}

The three parallel research roles performed substantive cross-review. A
independently read B's eight interval certificates, checked 504
principal minors exactly, and verified Choi normalization and tensor
order. B reviewed C's statistics, error composition, and strong
baseline. C reviewed A's selection counterexample, operational
interface, and all-input quantifiers. Review identified two corrected
issues: variable shadowing had mislabeled an actual \(K=20\) execution
as 24, and one decimal display in the old nominal file was
\(4\times10^{-13}\) below its exact upper bound. Additive corrections
and new outward bounds preserve the old files and execution identities.
They do not conceal refitting or alter frozen held-out decisions. This
is independent-role review within the project, not external peer review.

Nine new stable results are registered with separate scopes for
conditional paper proofs, explicit counterexamples, exact finite
algebra, and executed interval calculations. The previous 194 result
objects and 117 claims retain their historical identities; this edition
has 203 results and 126 claims. No new Lean proof is claimed. The
integrated replay separately records execution time, source hashes,
reuse of original samples, commands, resources, and outputs; old
synthetic samples are not counted as new experiments.

The completed research loop has a declared scope: genuine finite
propagation and parameter bounds connect one persistent source to legal
complete instruments, independent holdout, and auditable numerical
certificates. Available reachable data still do not select a global
quantum theory. Comparison with the strong standard method is equivalent
and yields no additional predictive advantage. The next physical
question can reduce independent physical inputs or exclude competing
theories through a noncircular principle, specify a prediction that
differs from that baseline, or establish a rigorous resource improvement
under the same controls, data, and cost. An independently justified
reduction of premises can advance unification while reproducing standard
predictions. Device validation requires independent preparation,
control, readout, and complete stopped-record evidence matching this
protocol.

Further obligations remain explicit: infinite modes and continuum
limits, UV/IR and renormalization, relativistic causality, gauge and
matter selection, thermalization and energy exchange, genuine gravity
and spacetime dynamics, semiclassical limits, cosmology, and
observational constraints. This finite source does not prove those
interfaces. The release audit retains their premise classes,
distinguishable exit conditions, and full physics responsibility map.
Completed chapter counts do not measure the fraction of physics solved.
{[}QM3114{]}

\chapter{Appendix A. Notation, units, and responsibilities of standard
theorems}\label{appendix-a}

\section{A.1 Common notation}

\begin{longtable}[]{@{}
  >{\raggedright\arraybackslash}p{(\columnwidth - 2\tabcolsep) * \real{0.3000}}
  >{\raggedright\arraybackslash}p{(\columnwidth - 2\tabcolsep) * \real{0.7000}}@{}}
\toprule\noalign{}
\begin{minipage}[b]{\linewidth}\raggedright
Symbol
\end{minipage} & \begin{minipage}[b]{\linewidth}\raggedright
Meaning and convention
\end{minipage} \\
\midrule\noalign{}
\endhead
\bottomrule\noalign{}
\endlastfoot
\(\rho,\sigma\) & Density operators; subnormalization is stated
explicitly when allowed. \\
\(\mathcal E,\mathcal M_a\) & A CPTP channel and a CP branch map of an
instrument. \\
\(E_a,P_a\) & A POVM effect and an orthogonal branch projector;
generally different objects. \\
\(J_{\mathcal E}\) & Unnormalized input-first Choi operator, with trace
equal to the input dimension. \\
\(D(\rho,\sigma)\) & Trace distance where so defined in the text;
relative entropy is written \(D(\rho\|\sigma)\). \\
\(\delta_\diamond\) & Half the diamond-norm difference, unless otherwise
specified. \\
\(C(h)\) & Canonical response for a fixed future target and protocol. \\
FFIPS & Future/intervention predictive equivalence and its canonical
sufficient representation. \\
ORC & Objective Record Certificate: an operational certificate for a
public record. \\
SRN & Strong record nondisturbance: preservation of every state within
each branch. \\
QDRR & The project's record-writing and inverse-writing/recovery
protocol, used with the precise definition in Chapter 8. \\
\(K,\Lambda\) & An SRN correlation matrix and an entrywise-exponential
generator; not a general Kraus operator. \\
\(\mathsf G\) or \(G\) & A Gram matrix, Gaussian quadratic kernel, or
group, according to context; these types cannot be mixed across
chapters. \\
\(\beta_T,\beta_C,\beta_8\) & Inverse temperature, control discount, and
the 331 charge-embedding parameter: three distinct quantities. \\
\(\hbar\) & Calibration of the unit of action; set to one in natural
units in the field-theory chapters and restored explicitly where
needed. \\
\(\sigma,\Omega\) & The phase-space symplectic bilinear form and its
matrix representation. \\
\(\Phi,\Pi\) & Core structure/resource symbols whose specific function
must be declared; not a fixed physical entropy. \\
\(\Pi_{\rm cut}\) & Energy-cutoff projector, kept distinct from a
resource cost. \\
\(\Gamma\) & A field-theory effective action or vertex; not an undefined
general Core complexity. \\
\(\lambda,\lambda_H\) & Scalar quartic and Higgs quartic couplings. \\
\(\mu,\Lambda_*\) & Renormalization scale and EFT high-energy scale in
the stated context; chemical potential is separately identified. \\
\(N_c,N_f,k,t\) & Number of colours, number of generations, number of
quark generations in the horizontal representation \(R\), and the
anomaly-solution multiplier. \\
\(\chi\) & A framing-completion response class, not an ordinary
four-dimensional chiral index. \\
HGP, FAD & The separate physical premise packages in Chapters 31 and
32. \\
GOIT & The UEOT interpretive interface for gauge-orbit identity
transport. \\
D/P/K/H/S/E/N/O & The argument-status classes of Chapter 0; P does not
imply originality or machine verification. \\
\end{longtable}

\section{A.2 Units}

Chapters 1--20 and the actualization and global-phase chapters retain
\(\hbar\). A Hamiltonian has units of energy and a Lindblad generator
has units of inverse time. In
\(\mathcal L\rho=-i[H,\rho]/\hbar+\sum L\rho L^\dagger-\frac12\{L^\dagger L,\rho\}\)
the operator \(L\) has units of inverse square root of time. Position
Fisher information has units of inverse length squared, and an action
has units of action.

In four-dimensional natural units,
\([\phi]=[A_\mu]=1,[\psi]=3/2,[\partial_\mu]=1\) and \([\mathcal L]=4\).
Gauge and quartic couplings have mass dimension zero; a cross section
has mass dimension minus two. In an EFT, an operator, its Wilson
coefficient, and any explicit high-energy scale must combine to total
mass dimension four in the Lagrangian density.

Entropy defined with natural logarithms is measured in nats. Conversion
to bits divides by \(\log2\). For thermodynamic entropy, inclusion of
\(k_B\) must be specified. Here \(\beta_T=(k_BT)^{-1}\), and information
entropy is dimensionless by default.

\section{A.3 Major results invoked without full reproving}

\begin{longtable}[]{@{}
  >{\raggedright\arraybackslash}p{(\columnwidth - 4\tabcolsep) * \real{0.2500}}
  >{\raggedright\arraybackslash}p{(\columnwidth - 4\tabcolsep) * \real{0.3400}}
  >{\raggedright\arraybackslash}p{(\columnwidth - 4\tabcolsep) * \real{0.4100}}@{}}
\toprule\noalign{}
\begin{minipage}[b]{\linewidth}\raggedright
Standard tool
\end{minipage} & \begin{minipage}[b]{\linewidth}\raggedright
Required assumptions and use here
\end{minipage} & \begin{minipage}[b]{\linewidth}\raggedright
Outside the proof supplied here
\end{minipage} \\
\midrule\noalign{}
\endhead
\bottomrule\noalign{}
\endlastfoot
Two-principle finite OPT reconstruction & The original preprint's
complete background, LE/CIT, and ESP. & Treating the abstract as a
complete axiom list or an established final theorem. \\
Complex polarization and quadratic-form characterization & Global
homogeneity, parallelogram relations, and continuity. & Replacing global
quantifiers with finite samples. \\
Relative-entropy data processing & CPTP maps and appropriate support;
primarily finite-dimensional use. & Arbitrary non-CP approximations or
undefined entropy differences. \\
Knill--Laflamme & A fixed code and noise linear span, with perfect
recovery. & Universal hardware fault tolerance without a specified noise
model. \\
Spectral theorem and Stone's theorem & Self-adjointness, strong
continuity, and appropriate operator domains. & Automatic
self-adjointness of formal differential expressions. \\
Kato--Rellich/Friedrichs & Symmetry, relative bounds, semiboundedness,
and closed-form assumptions as applicable. & Validity of a Hamiltonian
without domain conditions. \\
Stone--von Neumann & Finitely many degrees of freedom, regularity,
irreducibility, and nonzero central character. & A globally unique
representation for infinitely many QFT degrees of freedom. \\
Williamson theorem and Gaussian CP conditions & Finitely many modes,
positive-definite covariance, and a declared symplectic convention. &
Determination of arbitrary non-Gaussian operations by second moments. \\
Wigner's theorem & A transition-probability-preserving ray bijection and
applicable continuity assumptions. & Identifying every noisy channel
with a reversible symmetry. \\
Lieb--Robinson bounds & Suitable local bounded interactions and spatial
summability. & One unconditional light cone for all long-range or
unbounded fields. \\
GKSL & The conditions of the chosen version, such as a norm-continuous
finite-dimensional CPTP semigroup. & Assuming every reduced dynamics is
a semigroup. \\
Davies weak-coupling limit & Equilibrium bath, correlation decay, and
the stated scaling limit. & Automatic thermal Lindblad dynamics at
strong coupling or with long memory. \\
Process tensors/combs & Complete positivity, causal recursion, and
specified intervention slots. & Recovering a full process from
correlations without interventions. \\
Spin--statistics and LSZ & Relativistic locality, spectrum and
positivity, or stable-particle asymptotic-state assumptions. & Arbitrary
nonrelativistic models or an uncomplicated charged infrared limit. \\
Perturbative renormalization and Ward/ST identities & Specified
regulator, scheme, and absence of the relevant gauge anomaly. &
Nonperturbative existence and global gauge fixing. \\
Atiyah--Singer/Dai--Freed & Specified manifolds, bundles, operators,
boundaries, and anomaly types. & Substituting Betti numbers for indices
or arbitrary framing maps. \\
\end{longtable}

Invoking a substantial theorem under its conditions records a
dependency; it does not remove a proof obligation. The book supplies
explicit calculable realizations so that basic formulas can be used
without first studying every source. Extending a conclusion beyond its
stated domain requires returning to the relevant rigorous theorem.

\section{A.4 Notation for the global and topological models}

\(\theta\) is the boundary twist on a ring, and \(\eta\) its
gauge-invariant twist. In Chapter 44, \(\mathcal A,\mathcal F\) denote
the Berry connection and curvature, and \(t_1,t_2\) the SSH hoppings. In
the lattice chapter, \(a\) is the lattice spacing and \(U_l\) a
group-valued link variable. The symbol \(p_0\) is a calibrated lower
bound on a conditional event probability; \(L_n\) is a statistical
supermartingale, not a Lindblad operator or a line bundle. Reused
symbols always take their types from the chapter in which they occur.

\section{A.5 Notation for the verified kernel}

In Chapters 49--50, \(R_\rho(w)\) denotes the response to all finite
words, \(W\) the real span of actual word effects, and \(x(\rho)\)
finite response coordinates. The quantity \(d_E\) is the supremum of
probability differences over all effects \(0\preceq E\preceq I\), while
\(\Delta_E\) is the supremum of actual channel-output discrepancies.
Neither is a maximum over a finite sample. The quantities \(p,q\) are
Born probabilities for the same outcome from two inputs, and \(\eta\) is
a specified positive probability lower bound. Sections 50.3 and 50.5
distinguish the original distance interfaces from stronger norm claims;
Chapters 60 and 63 establish state separation and the
half-trace-absolute-value identity. An arbitrary-ancilla diamond-norm
identification remains a separate obligation.

\chapter{Appendix B. Registry of numbered results}\label{appendix-b}

The 3.7 bilingual revision preserves all 173 stable result IDs of
research edition 3.6, without adding or withdrawing another result.
These comprise the 161 results of edition 3.5 and the 11 propositions
and one counterexample added in Chapter 66 by the 3.6 research revision.
The earlier composition is 78 from edition 2.0, 10 in Chapters 49--50,
13 in Chapters 51--53, 16 in Chapters 54--56, 14 in Chapters 57--59, 16
in Chapters 60--62, and 14 in Chapters 63--65. The registry counts
analytic statements in the text, not original discoveries, Lean
declarations, or experimental confirmations. Each chapter supplies the
full hypotheses, proof, and open obligations. The bilingual map
separately binds the titles and anchors in both languages. Edition 3.8
appends the nine Chapter 67 results, giving 182 in total; all 173
previous result objects and their historical verification identities are
preserved. Edition 3.9 adds six Chapter 68 results for a total of 188,
preserving all 182 prior result objects. Edition 3.10 appends six
Chapter 69 results for a total of 194, preserving all 188 previous
result objects. Edition 3.11 adds nine Chapter 70 results for a total of
203; all 194 previous result objects are preserved item by item.

\begin{longtable}[]{@{}
  >{\raggedright\arraybackslash}p{(\columnwidth - 4\tabcolsep) * \real{0.2400}}
  >{\raggedleft\arraybackslash}p{(\columnwidth - 4\tabcolsep) * \real{0.1000}}
  >{\raggedright\arraybackslash}p{(\columnwidth - 4\tabcolsep) * \real{0.6600}}@{}}
\toprule\noalign{}
\begin{minipage}[b]{\linewidth}\raggedright
ID
\end{minipage} & \begin{minipage}[b]{\linewidth}\raggedleft
Chapter
\end{minipage} & \begin{minipage}[b]{\linewidth}\raggedright
Result
\end{minipage} \\
\midrule\noalign{}
\endhead
\bottomrule\noalign{}
\endlastfoot
\hyperref[p-qbr-01]{P-QBR-01} & 1 & P-QBR-01: A classical realization
countermodel. \\
\hyperref[p-qsel-01]{P-QSEL-01} & 2 & P-QSEL-01: Classical and
real-quantum deletion witnesses. \\
\hyperref[p-qact-01]{P-QACT-01} & 3 & P-QACT-01: Conditional
Madelung--Schrödinger equivalence. \\
\hyperref[p-qact-02]{P-QACT-02} & 3 & P-QACT-02. \\
\hyperref[p-qborn-02]{P-QBORN-02} & 3 & P-QBORN-02: Probabilities of an
unbiased final branch. \\
\hyperref[p-qop-01]{P-QOP-01} & 4 & P-QOP-01: Finite Choi--Kraus
equivalence. \\
\hyperref[p-qop-02]{P-QOP-02} & 4 & P-QOP-02: No signalling. \\
\hyperref[p-qop-03]{P-QOP-03} & 4 & P-QOP-03: Operational error and
propagation. \\
\hyperref[p-qff-01]{P-QFF-01} & 5 & P-QFF-01: Finite quantum FFIPS. \\
\hyperref[p-qff-02]{P-QFF-02} & 5 & P-QFF-02: Path preservation by a
quantum quotient. \\
\hyperref[p-qmeas-01]{P-QMEAS-01} & 6 & P-QMEAS-01: Parent-system
completion. \\
\hyperref[p-qmeas-02]{P-QMEAS-02} & 6 & P-QMEAS-02. \\
\hyperref[p-qmeas-03]{P-QMEAS-03} & 6 & P-QMEAS-03: Public blindness. \\
\hyperref[p-qmeas-04]{P-QMEAS-04} & 6 & P-QMEAS-04. \\
\hyperref[p-qsrn-01]{P-QSRN-01} & 7 & P-QSRN-01: Complete SRN
classification. \\
\hyperref[p-qsrn-02]{P-QSRN-02} & 7 & P-QSRN-02. \\
\hyperref[p-qsrn-03]{P-QSRN-03} & 7 & P-QSRN-03: Necessary and
sufficient generator-cone conditions. \\
\hyperref[p-qqdr-01]{P-QQDR-01} & 8 & P-QQDR-01: Coherent pullback. \\
\hyperref[p-qqdr-02]{P-QQDR-02} & 8 & P-QQDR-02. \\
\hyperref[p-qqdr-03]{P-QQDR-03} & 8 & P-QQDR-03: Simultaneous error. \\
\hyperref[p-qinfo-01]{P-QINFO-01} & 9 & P-QINFO-01. \\
\hyperref[p-qinfo-02]{P-QINFO-02} & 9 & P-QINFO-02. \\
\hyperref[p-qinf-01]{P-QINF-01} & 10 & P-QINF-01: Trace-norm continuity
of normal states. \\
\hyperref[p-qinf-02]{P-QINF-02} & 10 & P-QINF-02: Energy-truncation
error. \\
\hyperref[p-qccr-01]{P-QCCR-01} & 11 & P-QCCR-01: A single scale from
symplectic covariance. \\
\hyperref[p-qccr-02]{P-QCCR-02} & 11 & P-QCCR-02. \\
\hyperref[p-qccr-03]{P-QCCR-03} & 11 & P-QCCR-03. \\
\hyperref[p-qgauss-01]{P-QGAUSS-01} & 13 & P-QGAUSS-01: A necessary
covariance condition. \\
\hyperref[p-qmem-01]{P-QMEM-01} & 18 & P-QMEM-01: Irreducible error from
omitted memory. \\
\hyperref[p-qtherm-01]{P-QTHERM-01} & 19 & P-QTHERM-01: Free energy and
relative entropy. \\
\hyperref[p-qtherm-02]{P-QTHERM-02} & 19 & P-QTHERM-02: KMS for finite
Gibbs systems. \\
\hyperref[p-qtherm-03]{P-QTHERM-03} & 19 & P-QTHERM-03. \\
\hyperref[p-qctl-01]{P-QCTL-01} & 20 & P-QCTL-01. \\
\hyperref[p-qkg-01]{P-QKG-01} & 22 & P-QKG-01: Canonical relations and
propagation. \\
\hyperref[p-qkg-02]{P-QKG-02} & 22 & P-QKG-02: Spacelike causality. \\
\hyperref[p-qdir-01]{P-QDIR-01} & 23 & P-QDIR-01. \\
\hyperref[p-qunit-01]{P-QUNIT-01} & 24 & P-QUNIT-01: Optical theorem. \\
\hyperref[p-qren-01]{P-QREN-01} & 25 & P-QREN-01. \\
\hyperref[p-qren-02]{P-QREN-02} & 25 & P-QREN-02: Subtracted bubble. \\
\hyperref[p-qren-03]{P-QREN-03} & 25 & P-QREN-03: Two-particle cut. \\
\hyperref[p-qeft-01]{P-QEFT-01} & 26 & P-QEFT-01: Quadratic Schur
complement. \\
\hyperref[p-qeft-02]{P-QEFT-02} & 26 & P-QEFT-02: A restricted real-time
bridge. \\
\hyperref[p-qqed-01]{P-QQED-01} & 28 & P-QQED-01. \\
\hyperref[p-qqed-02]{P-QQED-02} & 28 & P-QQED-02: Threshold beta
function. \\
\hyperref[p-qsm-01]{P-QSM-01} & 29 & P-QSM-01: Anomaly solutions and
nonuniqueness. \\
\hyperref[p-qhig-01]{P-QHIG-01} & 30 & P-QHIG-01: Mass matrices. \\
\hyperref[p-qflav-01]{P-QFLAV-01} & 30 & P-QFLAV-01: Generic
nondegenerate \(N\)-generation quark counting. \\
\hyperref[p-qnu-01]{P-QNU-01} & 30 & P-QNU-01: Type-I seesaw. \\
\hyperref[p-qg3-01]{P-QG3-01} & 31 & P-QG3-01. \\
\hyperref[p-qg3-02]{P-QG3-02} & 31 & P-QG3-02: All integer solutions. \\
\hyperref[p-qg3-03]{P-QG3-03} & 31 & P-QG3-03: Selection without
primitivity. \\
\hyperref[p-qg3-04]{P-QG3-04} & 31 & P-QG3-04. \\
\hyperref[p-qfr-01]{P-QFR-01} & 32 & P-QFR-01. \\
\hyperref[p-qfr-02]{P-QFR-02} & 32 & P-QFR-02. \\
\hyperref[g-01]{G-01} & 36 & G-01: All integer solutions for the
three-scalar content. \\
\hyperref[g-02]{G-02} & 36 & G-02: Local completion for general colour
number. \\
\hyperref[f-01]{F-01} & 37 & F-01. \\
\hyperref[f-02]{F-02} & 37 & F-02. \\
\hyperref[f-03]{F-03} & 37 & F-03. \\
\hyperref[m-01]{M-01} & 38 & M-01. \\
\hyperref[m-02]{M-02} & 38 & M-02. \\
\hyperref[m-03]{M-03} & 39 & M-03: Four is the minimum number of
branches for single-time separation. \\
\hyperref[m-04]{M-04} & 39 & M-04: A sharp classical lower bound. \\
\hyperref[m-05]{M-05} & 39 & M-05. \\
\hyperref[e-01]{E-01} & 39 & E-01: A finite-sample rejection test. \\
\hyperref[e-02]{E-02} & 39 & E-02: A robust bound for approximate
population preservation. \\
\hyperref[c-01]{C-01} & 40 & C-01: An explicit contradiction from remote
preparation. \\
\hyperref[c-02]{C-02} & 40 & C-02: Norm preservation and a linear
ensemble average. \\
\hyperref[c-03]{C-03} & 40 & C-03: A single spectral branch with Born
weights. \\
\hyperref[c-04]{C-04} & 40 & C-04: Decay of pairwise uncertainty. \\
\hyperref[c-05]{C-05} & 41 & C-05: Equivalence with an unrestricted
quantum environment. \\
\hyperref[p-ugl-01]{P-UGL-01} & 43 & P-UGL-01: Ring spectrum and gauge
invariance. \\
\hyperref[p-utp-01]{P-UTP-01} & 44 & P-UTP-01: Chern number of this
eigenline. \\
\hyperref[p-utp-02]{P-UTP-02} & 44 & P-UTP-02: A semi-infinite zero
mode. \\
\hyperref[p-ulg-01]{P-ULG-01} & 45 & P-ULG-01: Existence and gauge
invariance of the finite model. \\
\hyperref[p-ulg-02]{P-ULG-02} & 45 & P-ULG-02: Exact partition function
and response. \\
\hyperref[p-ust-01]{P-UST-01} & 46 & P-UST-01: A fixed-sample test for
correlated trials. \\
\hyperref[p-ust-02]{P-UST-02} & 46 & P-UST-02: An optional-stopping
certificate. \\
\hyperref[p-qkr-01]{P-QKR-01} & 49 & P-QKR-01 (Posterior closure of the
response object). \\
\hyperref[p-qkr-02]{P-QKR-02} & 49 & P-QKR-02 (A local finite
representation determines every future response). \\
\hyperref[p-qkr-03]{P-QKR-03} & 49 & P-QKR-03 (All-word coordinates
determine a complete finite adaptive tree). \\
\hyperref[p-qkr-04]{P-QKR-04} & 49 & P-QKR-04 (Soundness of an exact
certificate). \\
\hyperref[p-qst-01]{P-QST-01} & 50 & P-QST-01 (Positive pairing and
effect bounds). \\
\hyperref[p-qst-02]{P-QST-02} & 50 & P-QST-02 (Complete pullback of
continuation effects). \\
\hyperref[p-qst-03]{P-QST-03} & 50 & P-QST-03 (Basic contract for the
observable distance). \\
\hyperref[p-qst-04]{P-QST-04} & 50 & P-QST-04 (One-step stability of a
complete Kraus channel). \\
\hyperref[p-qst-05]{P-QST-05} & 50 & P-QST-05 (Composable error for
quantum trajectories). \\
\hyperref[p-qst-06]{P-QST-06} & 50 & P-QST-06 (Postselection error and
preservation of positive support). \\
\hyperref[p-qmx-01]{P-QMX-01} & 51 & P-QMX-01 (An actual realization of
a probability mixture). \\
\hyperref[p-qmx-02]{P-QMX-02} & 51 & P-QMX-02 (A uniform bound for
changed weights in a common library). \\
\hyperref[p-qmx-03]{P-QMX-03} & 51 & P-QMX-03 (A joint bound for weight
and member errors). \\
\hyperref[p-qmx-04]{P-QMX-04} & 51 & P-QMX-04 (Interpretation of an
exact rational budget). \\
\hyperref[p-qce-01]{P-QCE-01} & 52 & P-QCE-01 (Exact residual and
all-effect budget for an average state). \\
\hyperref[p-qce-02]{P-QCE-02} & 52 & P-QCE-02 (Compactness of the actual
density-operator set). \\
\hyperref[p-qce-03]{P-QCE-03} & 52 & P-QCE-03 (A stationary density
operator for a finite complete channel). \\
\hyperref[n-qce-01]{N-QCE-01} & 52 & N-QCE-01 (Average invariance does
not imply orbit convergence or uniqueness). \\
\hyperref[p-qsv-01]{P-QSV-01} & 53 & P-QSV-01 (A complete quantum
instrument for a substochastic kernel). \\
\hyperref[p-qsv-02]{P-QSV-02} & 53 & P-QSV-02 (Intertwining diagonal
dynamics with finite success words). \\
\hyperref[p-qsv-03]{P-QSV-03} & 53 & P-QSV-03 (The geometric survival
law of the actual all-success record). \\
\hyperref[p-qsv-04]{P-QSV-04} & 53 & P-QSV-04 (A complete realization of
the Doob kernel and invariant success records). \\
\hyperref[n-qsv-01]{N-QSV-01} & 53 & N-QSV-01 (The QSD, Doob invariant
law, and an arbitrary initial state are distinct). \\
\hyperref[p-qcl-01]{P-QCL-01} & 54 & P-QCL-01 (States, responses, and
transport on the full simplex). \\
\hyperref[n-qcl-01]{N-QCL-01} & 54 & N-QCL-01 (A two-point mixed state
has no finite classical purification). \\
\hyperref[p-qcl-02]{P-QCL-02} & 54 & P-QCL-02 (The actual terminal
quotient retains two-point support). \\
\hyperref[p-qcl-03]{P-QCL-03} & 54 & P-QCL-03 (A common-source safe
trajectory and a purification-separation certificate). \\
\hyperref[n-qcl-02]{N-QCL-02} & 54 & N-QCL-02 (An operational quotient
can remove microscopic mixing). \\
\hyperref[p-qrt-01]{P-QRT-01} & 55 & P-QRT-01 (Two valid, distinct real
joint states). \\
\hyperref[n-qrt-01]{N-QRT-01} & 55 & N-QRT-01 (Real local product tests
do not determine the joint state). \\
\hyperref[p-qrt-02]{P-QRT-02} & 55 & P-QRT-02 (Deterministic
discrimination by a joint effect). \\
\hyperref[n-qrt-02]{N-QRT-02} & 55 & N-QRT-02 (Discrimination when
complex local effects are allowed). \\
\hyperref[n-qrt-03]{N-QRT-03} & 55 & N-QRT-03 (The distinguishing effect
is outside the finite real span of product effects). \\
\hyperref[p-qqd-01]{P-QQD-01} & 56 & P-QQD-01 (Validity of filter
normalization). \\
\hyperref[p-qqd-02]{P-QQD-02} & 56 & P-QQD-02 (An actual quantum Doob
instrument). \\
\hyperref[p-qqd-03]{P-QQD-03} & 56 & P-QQD-03 (Finite intertwining on
all operators and the actual conditional state). \\
\hyperref[p-qqd-04]{P-QQD-04} & 56 & P-QQD-04 (Quasi-stationarity,
geometric survival, and stationarity of the actual channel). \\
\hyperref[p-qqd-05]{P-QQD-05} & 56 & P-QQD-05 (A nondiagonal Doob
witness computable on all matrices). \\
\hyperref[n-qqd-01]{N-QQD-01} & 56 & N-QQD-01 (Three distinctions that
must be retained). \\
\hyperref[p-pur-01]{P-PUR-01} & 57 & P-PUR-01 (Vector purity in the
complete density-state space). \\
\hyperref[p-pur-02]{P-PUR-02} & 57 & P-PUR-02 (Existence of a
purification for every finite complex density state). \\
\hyperref[p-opt-01]{P-OPT-01} & 57 & P-OPT-01 (Three concrete
presentations in one operational language). \\
\hyperref[n-opt-01]{N-OPT-01} & 57 & N-OPT-01 (Basic operational laws do
not automatically select a quantum theory). \\
\hyperref[p-cq-01]{P-CQ-01} & 58 & P-CQ-01 (Complete positivity, trace
preservation, and effect contraction for rectangular channels). \\
\hyperref[p-cq-02]{P-CQ-02} & 58 & P-CQ-02 (Joint block state and
marginals of an actual instrument). \\
\hyperref[p-cq-03]{P-CQ-03} & 58 & P-CQ-03 (Record-and-memory
realization of every finite protocol). \\
\hyperref[p-cq-04]{P-CQ-04} & 58 & P-CQ-04 (Joint-distance and
perturbation control of record distributions). \\
\hyperref[n-cq-01]{N-CQ-01} & 58 & N-CQ-01 (Equal unconditional channels
do not control record error). \\
\hyperref[p-qdd-01]{P-QDD-01} & 59 & P-QDD-01 (Exact dynamics of the
present Doob channel). \\
\hyperref[p-qdd-02]{P-QDD-02} & 59 & P-QDD-02 (Closed forms for the
actual source success operator and its mass). \\
\hyperref[p-qdd-03]{P-QDD-03} & 59 & P-QDD-03 (The actual conditional
state and its quantitative convergence). \\
\hyperref[p-qdd-04]{P-QDD-04} & 59 & P-QDD-04 (Decay of finite-success
probabilities and the role of the initial state). \\
\hyperref[n-qdd-01]{N-QDD-01} & 59 & N-QDD-01 (A valid identity-Doob
negative control). \\
\hyperref[p-sep-01]{P-SEP-01} & 60 & P-SEP-01 (Representation of a
Hermitian operator by valid effects). \\
\hyperref[p-sep-02]{P-SEP-02} & 60 & P-SEP-02 (Separation by all effects
and the matrix meaning of zero distance). \\
\hyperref[p-clt-01]{P-CLT-01} & 60 & P-CLT-01 (Universal local
tomography of the complete finite complex model). \\
\hyperref[n-sep-01]{N-SEP-01} & 60 & N-SEP-01 (Even all diagonal effects
fail to identify coherent states). \\
\hyperref[p-puc-01]{P-PUC-01} & 61 & P-PUC-01 (Purity characterization
in the complete finite complex density space). \\
\hyperref[p-pun-01]{P-PUN-01} & 61 & P-PUN-01 (Unitary equivalence of
all purifications on the same auxiliary carrier). \\
\hyperref[p-aux-01]{P-AUX-01} & 61 & P-AUX-01 (Actual channel, marginal,
and purity of an auxiliary isometric embedding). \\
\hyperref[p-pun-02]{P-PUN-02} & 61 & P-PUN-02 (Unitary equivalence after
common extension for different finite auxiliary carriers). \\
\hyperref[n-pun-01]{N-PUN-01} & 61 & N-PUN-01 (The adjoint, transpose,
and entrywise conjugate cannot replace the actual auxiliary action). \\
\hyperref[n-aux-01]{N-AUX-01} & 61 & N-AUX-01 (Strict auxiliary
enlargement is an isometric embedding, not a two-sided unitary
equivalence). \\
\hyperref[p-arb-01]{P-ARB-01} & 62 & P-ARB-01 (History blocks,
normalized weights, and valid block effects). \\
\hyperref[p-arb-02]{P-ARB-02} & 62 & P-ARB-02 (Action-weighted bound for
the actual history-reading channel). \\
\hyperref[p-arb-03]{P-ARB-03} & 62 & P-ARB-03 (Matrix agreement between
finite-history evolution and the original complete protocol). \\
\hyperref[p-arb-04]{P-ARB-04} & 62 & P-ARB-04 (Actual per-action sum
budget and complete-record TV). \\
\hyperref[p-arb-05]{P-ARB-05} & 62 & P-ARB-05 (An exact rational witness
saturating the stepwise budget). \\
\hyperref[n-arb-01]{N-ARB-01} & 62 & N-ARB-01 (Zero unlabelled action
defects can coexist with maximum complete-record error). \\
\hyperref[p-tdist-01]{P-TDIST-01} & 63 & P-TDIST-01 (Positive and
negative spectral parts and the actual absolute value of a finite
Hermitian matrix). \\
\hyperref[p-tdist-02]{P-TDIST-02} & 63 & P-TDIST-02
(Half-trace-absolute-value formula and an optimal effect for full-effect
distance). \\
\hyperref[p-dec-01]{P-DEC-01} & 63 & P-DEC-01 (Actual optimal success
probability for equal-prior binary discrimination). \\
\hyperref[p-tcert-01]{P-TCERT-01} & 63 & P-TCERT-01 (Distance
certificate from a PSD support decomposition). \\
\hyperref[p-pdm-01]{P-PDM-01} & 64 & P-PDM-01 (Exact dimension ordering
for a rectangular isometry). \\
\hyperref[p-pdm-02]{P-PDM-02} & 64 & P-PDM-02 (A dimension-directed
direct isometry for arbitrary actual purifications). \\
\hyperref[p-pdm-03]{P-PDM-03} & 64 & P-PDM-03 (Actual rank lower bound
for all convex-pure extensions). \\
\hyperref[p-pdm-04]{P-PDM-04} & 64 & P-PDM-04 (Exact minimum auxiliary
dimension of finite complex purification). \\
\hyperref[n-pdm-01]{N-PDM-01} & 64 & N-PDM-01 (Equal marginals do not
remove the dimension direction required by a direct isometry). \\
\hyperref[p-stop-01]{P-STOP-01} & 65 & P-STOP-01 (Actual absorbing
padding of an original finite stopping tree). \\
\hyperref[p-stop-02]{P-STOP-02} & 65 & P-STOP-02 (Full-effect distance
preservation and the connection to the original joint leaf state). \\
\hyperref[p-stop-03]{P-STOP-03} & 65 & P-STOP-03 (Per-action error
budget under a common finite stopping rule). \\
\hyperref[p-stop-04]{P-STOP-04} & 65 & P-STOP-04 (Exact \(1/16\) record
difference in a stopping protocol). \\
\hyperref[p-stop-05]{P-STOP-05} & 65 & P-STOP-05 (Identical quantum
actions do not control all records under different stopping rules). \\
\hyperref[p-fcond-01]{P-FCOND-01} & 66 & P-FCOND-01 (Acceptance effects
of the same finite original experiment). \\
\hyperref[p-fcond-02]{P-FCOND-02} & 66 & P-FCOND-02 (Kernel annihilation
from the backward Gram sum). \\
\hyperref[p-fcond-03]{P-FCOND-03} & 66 & P-FCOND-03 (Support
transformation and a valid completion on the full space). \\
\hyperref[p-fcond-04]{P-FCOND-04} & 66 & P-FCOND-04 (Conditioning the
full leaf record and terminal state from the same original
experiment). \\
\hyperref[p-fcond-05]{P-FCOND-05} & 66 & P-FCOND-05 (A convex-hull
equivalence preserving outcomes and observation permissions). \\
\hyperref[p-fcond-06]{P-FCOND-06} & 66 & P-FCOND-06 (A finite negative
separation certificate for the allowed class). \\
\hyperref[p-fcond-07]{P-FCOND-07} & 66 & P-FCOND-07 (Classification of
deterministic initialization on all unknown inputs). \\
\hyperref[p-fcond-08]{P-FCOND-08} & 66 & P-FCOND-08 (Actual success
probability and posterior state of the scaled filter). \\
\hyperref[p-fcond-09]{P-FCOND-09} & 66 & P-FCOND-09 (Full cost of
independent resets until success). \\
\hyperref[p-fcond-10]{P-FCOND-10} & 66 & P-FCOND-10 (Propagation of
error from full outputs to conditioned outputs). \\
\hyperref[p-fcond-11]{P-FCOND-11} & 66 & P-FCOND-11 (A usable local
error bound for valid Kraus approximations). \\
\hyperref[c-fcond-01]{C-FCOND-01} & 66 & C-FCOND-01 (Small positive
eigenvalues cannot be treated as zero). \\
\hyperref[p-opt-02]{P-OPT-02} & 67 & P-OPT-02 (Four equivalent
identities for events in positive-dimensional finite complex
systems). \\
\hyperref[p-opt-03]{P-OPT-03} & 67 & P-OPT-03 (Composition respects
identity at every visible label). \\
\hyperref[p-micro-01]{P-MICRO-01} & 67 & P-MICRO-01 (A microscopic
bridge to complete adaptive CQ truncation). \\
\hyperref[c-micro-01]{C-MICRO-01} & 67 & C-MICRO-01 (Mean-energy
preservation does not imply branchwise truncation). \\
\hyperref[p-micro-02]{P-MICRO-02} & 67 & P-MICRO-02 (Normal trace-class
postselection and complete preparation records). \\
\hyperref[p-rob-01]{P-ROB-01} & 67 & P-ROB-01 (Deployment error and
separation on one parameter set). \\
\hyperref[c-rob-01]{C-ROB-01} & 67 & C-ROB-01 (Pointwise realizability
does not imply one robust deployment). \\
\hyperref[p-rob-02]{P-ROB-02} & 67 & P-ROB-02 (An exact design-region
certificate for the same source). \\
\hyperref[p-micro-03]{P-MICRO-03} & 67 & P-MICRO-03 (Single-excitation
marginal mixing without full-record forgetting). \\
\hyperref[p-micro-04]{P-MICRO-04} & 68 & P-MICRO-04 (Finite excitation
growth under segmented blue pulses). \\
\hyperref[p-opt-04]{P-OPT-04} & 68 & P-OPT-04 (Complete rectangular
operations on growing carriers). \\
\hyperref[p-micro-05]{P-MICRO-05} & 68 & P-MICRO-05 (Complete-record
omission bound using a no-blue reference). \\
\hyperref[c-micro-02]{C-MICRO-02} & 68 & C-MICRO-02 (An infinite-tail
counterexample for simultaneous red and blue interaction). \\
\hyperref[p-rob-03]{P-ROB-03} & 68 & P-ROB-03 (A joint certificate for
growth, physical filtering, and pulse omission). \\
\hyperref[c-rob-02]{C-ROB-02} & 68 & C-ROB-02 (Identical flag
probabilities do not certify the blue-pulse terminal state). \\
\hyperref[p-opt-05]{P-OPT-05} & 69 & P-OPT-05 (Complete-CQ error from
unnormalized reference prefixes). \\
\hyperref[p-micro-06]{P-MICRO-06} & 69 & P-MICRO-06 (Self-adjoint
simultaneous source, energy budget, and boundary defect). \\
\hyperref[p-micro-07]{P-MICRO-07} & 69 & P-MICRO-07 (Genuine truncation
bound from finite Duhamel expansion). \\
\hyperref[c-micro-03]{C-MICRO-03} & 69 & C-MICRO-03 (Galerkin
compression does not preserve this quadrature invariant). \\
\hyperref[p-rob-04]{P-ROB-04} & 69 & P-ROB-04 (Joint calibration,
nominal parameters, and shared-force stability). \\
\hyperref[p-rob-05]{P-ROB-05} & 69 & P-ROB-05 (Composing physical
preparation, robust deployment, nominal calibration, and finite
prediction). \\
\hyperref[c-qsel-02]{C-QSEL-02} & 70 & C-QSEL-02 (Operational
equivalence of hidden completions of a reachable sector and the boundary
of global quantum selection). \\
\hyperref[p-tom-01]{P-TOM-01} & 70 & P-TOM-01 (An explicit finite-tester
stability bound and its preparation, readout and dimension
obligations). \\
\hyperref[p-micro-08]{P-MICRO-08} & 70 & P-MICRO-08 (A complete
persistent Gaussian-bath process on a common domain). \\
\hyperref[p-micro-09]{P-MICRO-09} & 70 & P-MICRO-09 (An all-input finite
approximation of the same persistent source). \\
\hyperref[c-qmem-01]{C-QMEM-01} & 70 & C-QMEM-01 (Equal single-slot
channels do not determine a multitime process). \\
\hyperref[p-num-01]{P-NUM-01} & 70 & P-NUM-01 (From interval Choi bounds
to a physical terminal instrument). \\
\hyperref[p-qmem-02]{P-QMEM-02} & 70 & P-QMEM-02 (A necessary memory
dimension from causal-break response rank). \\
\hyperref[p-rob-06]{P-ROB-06} & 70 & P-ROB-06 (joint identification and
coverage from two trusted settings). \\
\hyperref[p-rob-07]{P-ROB-07} & 70 & P-ROB-07 (a complete CQ certificate
and independent holdout). \\
\end{longtable}

\chapter{Appendix C. Sources, inheritance, and research
integrity}\label{appendix-c}

\section{C.1 The source boundary of this edition}

The inherited editorial record documents direct checks of five
original/Core sources: the original UEOT PDF, the complete Core v3
mathematical text, the open-port matrix, the post-Core
scientific-closure account, and the theory-completion roadmap. It also
records the recovery and inspection of the unified version 2.0 Markdown,
LaTeX, PDF, and complete source/verification package. The Markdown and
LaTeX files were byte-identical to the corresponding package members.
The standalone PDF and packaged PDF differed in bytes and retain
separate identities; they are not presented as the same rendered
artifact.

The mathematics and runtime code incorporated from the v31.10 research
revision are identified by that revision's file-level source manifest
and actual verification records. Its research baseline was main commit
\texttt{3d443ec29eb1\allowbreak 7af1fc00699a\allowbreak 74e2d42d7351\allowbreak 6d81}. The earlier v31.9
scientific source \texttt{27f6ac776c56\allowbreak 2512451000d5\allowbreak 957087d3ecdc\allowbreak d086} had
been merged through PR \#356 as
\texttt{e9b61ca9c359\allowbreak 55028c4bfa42\allowbreak 5dfa278fd186\allowbreak 4029}; PR CI
\texttt{39569411258} and main CI \texttt{39570687643} both succeeded.
The previously registered pending status was outdated and must not be
interpreted as evidence that the source remained unpublished. Work in
that research revision's isolated checkout was confined to physics, with
Core retained as a read-only dependency. Earlier scientific sources,
literature, and evidence retain their version identities, and newly
executed verification is registered separately. This does not claim that
every historical attachment was reread or that the unchanged inherited
monograph programs were rerun.

The historical 3.7 publication baseline and the verified v31.9/v31.10
remote outcomes are stated separately in §C.8. Historical dates and
acceptance states are not silently promoted to current facts.

\section{C.2 Project-source map}

\begin{longtable}[]{@{}
  >{\raggedright\arraybackslash}p{(\columnwidth - 2\tabcolsep) * \real{0.3000}}
  >{\raggedright\arraybackslash}p{(\columnwidth - 2\tabcolsep) * \real{0.7000}}@{}}
\toprule\noalign{}
\begin{minipage}[b]{\linewidth}\raggedright
Source
\end{minipage} & \begin{minipage}[b]{\linewidth}\raggedright
Use in this book
\end{minipage} \\
\midrule\noalign{}
\endhead
\bottomrule\noalign{}
\endlastfoot
UEOT\_\allowbreak v3\_\allowbreak top.pdf & Original motivation concerning objects, phase,
resources, and physical unification; physical claims require subsequent
rigorous treatment \\
UEOT\_\allowbreak v1.0\_\allowbreak 核心理论正式版.md & Original definitions, Core, and
directions for physical research \\
UEOT\_\allowbreak 数学基础\_\allowbreak 统一正式版\_\allowbreak v1.1(1).md & Early mathematical support,
compared with later precise boundaries \\
UEOT\_\allowbreak Core\_\allowbreak Mathematics\_\allowbreak v3.0\_\allowbreak Complete.md & Generic bridges between
protocols, prediction, carriers, and control; quantum assumptions are
not inserted into the generic axioms \\
UEOT-QOD\_\allowbreak Interface\_\allowbreak v1.0; Core v2.2 dual track & Separation of physical
realization from generic object mathematics \\
UEOT-QM v8.1 and the v9--20 sequence & Historical sources for
purification, predictive objects, access carriers, process RG, and blind
theory competition; the generic part is developed in the mathematics
volume \\
UEOT-QM v21, v22, v23.0 & Coherent measurement parent layer, public
blindness, SRN classification/generators, and QDRR; Chapters 4--9 \\
P2\_\allowbreak Quantum\_\allowbreak Selection\_\allowbreak Minimality\_\allowbreak Frozen.md & Two-principle selection
and relative independence within a fixed OPT background; Chapter 2 \\
UEOT-QM v24.4 P1-CrossStage-Harmonized & Self-adjoint domains, Weyl
relations, UCA, finite canonical degrees of freedom, and a common
action-scale condition; Chapters 10--13 \\
UEOT-QM v25.2 P3-CrossStage-Harmonized & Symmetry, spin, identical
particles, and cross-stage boundaries; Chapters 14--15 \\
UEOT-QM v26.3 P4-Final-Closure & Many-body systems, open dynamics,
process memory, and statistical physics; Chapters 16--20 \\
UEOT-QM v27.1 Relativistic-Free-Fields & Relativistic assumptions, free
fields, and locality; Chapters 21--23 \\
UEOT-QM v28.1 Interacting-QFT-RG-EFT & Interactions, loop calculations,
renormalization, and bridges among three kinds of RG; Chapters 24--26 \\
UEOT-QM v29.2 Gauge-QED-Final-Closure & QED, Ward identities, UV poles,
scattering, thresholds, and magnetic moments; Chapters 27--28 \\
UEOT-QM v30.0 P8A, v30.2 P8C, v30.3 P8D & Non-Abelian/SM candidates,
anomalies, Higgs/flavour, and the limits of earlier selectors; Chapters
29--30 \\
UEOT-QM v30.4--v30.9 and the v30.10 synthesis & Assessment of CP,
common-q, framing, autonomy, and completion routes; Chapter 32 \\
UEOT-QM v30.11 Gauge-Orbit-Horizontal-Anomaly-Atomicity & GOIT, the
replication obstruction, HGP, 331 models, and conditional algebra;
Chapters 29 and 31 \\
\end{longtable}

Some early materials in this table serve as historical sources through
later syntheses and interface documents. Their presence does not mean
that their entire historical texts were reread for this revision.
Earlier versions' reported PASS counts are not carried over as newly
executed verification results.

\section{C.3 External original research and author resources}

The following references locate mathematical and physical precedents.
Their use does not make established theory an original UEOT result.

\begin{enumerate}
\def\labelenumi{\arabic{enumi}.}
\tightlist
\item
  Kenji Nakahira, \emph{Two Operational Principles Single Out Quantum
  Theory}, 2026 preprint. Chapter 2 depends on its complete OPT theorem
  background. A preprint and an experimentally validated theory have
  different evidential status.
  \href{https://arxiv.org/abs/2605.23217}{Original}.
\item
  John Watrous, \emph{The Theory of Quantum Information}, the 2018
  version supplied by the author. The standard tools concerning
  channels, norms, and information in Chapters 4 and 9 can be checked
  against this source.
  \href{https://cs.uwaterloo.ca/~watrous/TQI/}{Author's page}.
\item
  F. A. Pollock et al., \emph{Operational Markov Condition for Quantum
  Processes}, Physical Review Letters \textbf{120}, 040405 (2018). The
  causal-intervention criterion in Chapter 18 belongs to established
  process theory.
  \href{https://doi.org/10.1103/PhysRevLett.120.040405}{Original paper}.
\item
  F. A. Pollock et al., \emph{Non-Markovian quantum processes: Complete
  framework and efficient characterization}, Physical Review A
  \textbf{97}, 012127 (2018). This is an original framework for
  multi-time process tensors and is distinct from reference 3.
  \href{https://doi.org/10.1103/PhysRevA.97.012127}{DOI},
  \href{https://arxiv.org/abs/1512.00589}{author preprint}.
\item
  David Tong, \emph{Quantum Field Theory}, author lecture notes. These
  supply standard background on free fields, scattering, and gauge
  fields; this book's derivations and numerical benchmarks are stated
  separately.
  \href{https://www.damtp.cam.ac.uk/user/tong/qft.htm}{Author's page},
  \href{https://davidtong.org/teaching/quantum-field-theory/}{current
  teaching locator}.
\item
  Kazuya Yonekura, \emph{Dai-Freed theorem and topological phases of
  matter}. A mathematical-physics source for the determinant-line
  interpretation in Chapter 29.
  \href{https://arxiv.org/abs/1607.01873}{Original}.
\item
  I. García-Etxebarria and M. Montero, \emph{Dai-Freed anomalies in
  particle physics}. Used to distinguish local from global anomaly
  checks. \href{https://arxiv.org/abs/1808.00009}{Original}.
\item
  F. Pisano and V. Pleitez, \emph{SU(3)⊗U(1) Model for Electroweak
  Interactions} (1992). The organization of 331 representations has
  established precedents.
  \href{https://arxiv.org/abs/hep-ph/9206242}{Original}.
\item
  R. Foot et al., \emph{SU(3)L⊗U(1)N and SU(4)L⊗U(1)N gauge models with
  right-handed neutrinos} (1994). A precedent for the 331RHN family.
  \href{https://doi.org/10.1103/PhysRevD.50.R34}{Original paper}.
\item
  R. M. Lutchyn, J. D. Sau, and S. Das Sarma, \emph{Majorana Fermions
  and a Topological Phase Transition in Semiconductor-Superconductor
  Heterostructures}. A theoretical source for the effective device model
  in Chapter 34, not a new experimental claim.
  \href{https://arxiv.org/abs/1002.4033}{Original}.
\item
  Y. Oreg, G. Refael, and F. von Oppen, \emph{Helical liquids and
  Majorana bound states in quantum wires}. An independent theoretical
  precedent for the same application.
  \href{https://arxiv.org/abs/1003.1145}{Original}.
\item
  D. Burgarth, G. Chiribella, V. Giovannetti, P. Perinotti, and K.
  Yuasa, \emph{Ergodic and Mixing Quantum Channels in Finite
  Dimensions}. A reference for the distinction between stationarity,
  ergodicity, and mixing in Chapter 52. The actual scope of this
  project's machine proofs is stated separately.
  \href{https://arxiv.org/abs/1210.5625}{Original}.
\item
  F. Carollo, J. P. Garrahan, I. Lesanovsky, and C. Pérez-Espigares,
  \emph{Making rare events typical in Markovian open quantum systems}.
  Broader background on quantum conditioned dynamics beyond Chapter 53;
  it does not certify that this project has implemented a general
  infinite-conditioning limit.
  \href{https://arxiv.org/abs/1711.10951}{Original}.
\item
  {[}CDP2010{]} G. Chiribella, G. M. D'Ariano, and P. Perinotti,
  \emph{Probabilistic theories with purification}, Physical Review A
  \textbf{81}, 062348 (2010). Chapter 57 uses it to distinguish
  purification existence from the full principle including essential
  uniqueness. \href{https://arxiv.org/abs/0908.1583}{Original}.
\item
  {[}CDP2011{]} The same authors, \emph{Informational derivation of
  Quantum Theory}, Physical Review A \textbf{84}, 012311 (2011). Its
  reconstruction invokes explicit background axioms and is not replaced
  by this book's finite purification-existence result.
  \href{https://arxiv.org/abs/1011.6451}{Original}.
\item
  W. H. Zurek, \emph{Quantum Darwinism}, Nature Physics \textbf{5},
  181--188 (2009). The established environmental-redundancy framework
  used in the discussion of records.
  \href{https://doi.org/10.1038/nphys1202}{DOI}.
\item
  T. P. Le and A. Olaya-Castro, \emph{Objectivity (or lack thereof):
  Comparison between predictions of quantum Darwinism and spectrum
  broadcast structure}, Physical Review A \textbf{98}, 032103 (2018). A
  total-mutual-information plateau alone need not certify accessible
  classical records.
  \href{https://doi.org/10.1103/PhysRevA.98.032103}{DOI}.
\item
  T. P. Le and A. Olaya-Castro, \emph{Strong Quantum Darwinism and
  Strong Independence are Equivalent to Spectrum Broadcast Structure},
  Physical Review Letters \textbf{122}, 010403 (2019); A. Feller et al.,
  Comment, \textbf{126}, 188901 (2021); Le and Olaya-Castro, Reply,
  \textbf{126}, 188902 (2021). The corrected independence requirement
  concerns the full conditional joint state, not merely pairwise
  independence.
  \href{https://doi.org/10.1103/PhysRevLett.122.010403}{Original},
  \href{https://doi.org/10.1103/PhysRevLett.126.188901}{Comment},
  \href{https://doi.org/10.1103/PhysRevLett.126.188902}{Reply}.
\item
  E. Tiesinga et al., \emph{CODATA recommended values of the fundamental
  physical constants: 2018}, Reviews of Modern Physics \textbf{93},
  025010 (2021). This identifies the frozen constant used in §28.1; it
  is not a claim to use the most recent adjustment.
  \href{https://doi.org/10.1103/RevModPhys.93.025010}{DOI},
  \href{https://physics.nist.gov/cuu/Constants/archive2018.html}{NIST
  archive}.
\end{enumerate}

The five laboratory studies relevant to Chapter 8 are cited with their
actual systems and evidential limits in §8.5. The accompanying
\texttt{PRIOR\_\allowbreak ART\_\allowbreak AND\_\allowbreak EXPERIMENTAL\_\allowbreak STATUS.\allowbreak md} gives the
source-review depth and further original references. A bibliography
entry does not mean that raw data were obtained or that all derivations
in that work were independently reproduced.

\section{C.4 The status of the book's synthesis and logical
improvements}

This book is not an unedited collection of historical articles. Its
publication lineage reorganizes units, Choi leg order, process
interfaces, physical/public layers, conditional proofs, failure
conditions, and responsibility for application errors. Chapter 31 proves
that the primitive-package premise is redundant under the stated weak
upper bound. The general 331 presentation explicitly includes the
charged \(E^c\) fields so that their anomaly contributions are not
omitted.

These are contributions to organization, calculation, and conditional
derivation. They do not claim the discovery of new 331 arithmetic or new
particle physics. A scholarly priority claim would require an
independent prior-art search, an exact statement of the new theorem,
peer scrutiny, and any claimed distinguishable physical consequences. A
self-contained exposition can still depend explicitly on established
theorems.

\section{C.5 Inherited advanced modules and sources added along the way}

All numbered results from the quantum-foundations text and the frontier
derivations are incorporated into the main text. The global-cycle model,
Berry/SSH analysis, finite lattice calculations, and conditional-trial
tests were subsequently added. The standard physical sources for
Chapters 44--45 are Berry (1984), Su--Schrieffer--Heeger (1979), and
Wilson (1974), linked next to the relevant derivations. A newly written
analytic proof or a new combination of uses is not evidence of first
discovery. Standard literature and project materials carry different
source responsibilities.

No new complete historical conversation or new hardware dataset was
obtained for this publication revision. Original attachments, accessible
evolving manuscripts, and the two earlier complete texts document the
research history. Inaccessible material is not used as an undisclosed
proof premise.

\section{C.6 Third-edition code and original-source checks}

Chapters 49--59 retain the finite kernels, stability, probability
mixtures and long-time averages, survival, selection separation,
purification, actual protocols, and conditioned dynamics of Kernel
1.1--1.7. Chapters 60--62 incorporate Kernel 1.8: all-effect separation,
complex local tomography, full purity, ancilla comparison, and action
budgets under a common controller. Chapters 63--65 incorporate Kernel
1.9: trace distance, optimal equal-prior discrimination, minimum ancilla
dimension, directed isometries, and finite stopping protocols. Chapter
66 incorporates Kernel 1.10: finite conditioned processes of the
original instrument, realizability under allowed actions, and complete
resource and error accounting. The relevant entry points are
\texttt{formalization/\allowbreak ueot-\allowbreak qm} and
\texttt{physics/\allowbreak qm/\allowbreak kernel/\allowbreak evidence/\allowbreak v31.\allowbreak 10}. Historical receipts keep
their original inputs and execution scopes. Frozen Core and domain
interfaces determine classification; a property within the quantum
sector is not inserted into the generic axioms.

The edition 3.0 revision directly checked Theorem 1 and Appendix
A.2--A.3 of Nakahira's preprint. Its complete OPT background is required
for invoking the theorem, and the conclusion does not require a physical
system of every natural-number dimension. Chapter 2 retains the
external-theorem status K and does not count this theorem among the
project's generic Lean results.
\href{https://arxiv.org/html/2605.23217v1}{Original HTML}. Standard
channel and norm theory can be consulted in
\href{https://cs.uwaterloo.ca/~watrous/TQI/}{Watrous's author
resources}. This project's Lean proofs proceed from their actual matrix
constructions; they neither reproduce his book nor certify that all of
its theory has been re-proved here.

The stronger inferences on pages 30--33 of the original paper---Fisher
information to uncertainty relations, closed action to CCR, Betti
numbers to an index and three generations, and Massey products to a
gauge group---require the bridges, domain inputs, or open obligations
distinguished throughout this book. Retaining an idea does not reinstate
a withdrawn inference. The precise premises, countermodels, and
inspectable proofs in the rigorous version determine what follows.

Source label QM316 identifies the classical sector, the
Core/purification boundary, real tomography, and general quantum Doob
dynamics. QM317 identifies full-state purification, operational
presentations of the classical, real, and complex models, rectangular
channels, joint record memory, and specified conditioned dynamics. QM318
identifies all-effect separation, complex local tomography, full purity,
ancilla-unitary comparison, history channels, and action-record budgets.
QM319 identifies \texttt{Quantum\allowbreak Trace\allowbreak Distance},
\texttt{Binary\allowbreak Discrimination}, \texttt{Purification\allowbreak Isometry},
\texttt{Purification\allowbreak Dimension}, \texttt{Stopping\allowbreak Protocols}, and
\texttt{Stopping\allowbreak Record\allowbreak Budget}, together with three new exact witnesses,
cross-review, and the unified verification gate. Mathlib lemmas have the
proof responsibility of their pinned library source; the text states the
conditions actually invoked and the conclusions actually exported.
Substantial library theorems are not relabelled as original physical
discoveries by this project. Earlier citations retain their registered
scope.

\section{C.7 Original research on finite conditioned processes}

{[}CFW-R1{]} Federico Carollo, Juan P. Garrahan, Igor Lesanovsky, and
Carlos Pérez-Espigares. \emph{Making rare events typical in Markovian
open quantum systems}. Physical Review A \textbf{98}, 010103(R) (2018).
\href{https://arxiv.org/abs/1711.10951}{Original paper},
\href{https://doi.org/10.1103/PhysRevA.98.010103}{journal DOI}. This
reference identifies prior work on quantum Doob transformations of
rare-trajectory dynamics, including finite-time, time-dependent
constructions. The finite-tree, support, action-realizability, and cost
relations in this book are proved separately in Chapter 66; the external
article's abstract is not a substitute for proofs of the complete chain
presented here.

\section{C.8 The 3.7 bilingual publication and the observed remote
record}

This bilingual project began with edition 3.5 at main commit
\texttt{3d443ec29eb1\allowbreak 7af1fc00699a\allowbreak 74e2d42d7351\allowbreak 6d81}. A refresh before
delivery found that main had advanced to
\texttt{925e3a8a411c\allowbreak 0b1bab99dd91\allowbreak 50bf00589365\allowbreak 6140}: PR \#450 incorporated
the v31.10 research and Chapter 66, and PR \#451 completed its observed
remote metadata. Publication revision 3.7 incorporates that advance in
full, retaining all 173 results of edition 3.6, including the original
161 and the 12 added in Chapter 66. It does not replace the new main
manuscript with the older edition.

\begin{longtable}[]{@{}
  >{\raggedright\arraybackslash}p{(\columnwidth - 4\tabcolsep) * \real{0.2500}}
  >{\raggedright\arraybackslash}p{(\columnwidth - 4\tabcolsep) * \real{0.3400}}
  >{\raggedright\arraybackslash}p{(\columnwidth - 4\tabcolsep) * \real{0.4100}}@{}}
\toprule\noalign{}
\begin{minipage}[b]{\linewidth}\raggedright
Research release
\end{minipage} & \begin{minipage}[b]{\linewidth}\raggedright
Actual source and merge identities
\end{minipage} & \begin{minipage}[b]{\linewidth}\raggedright
Actual CI
\end{minipage} \\
\midrule\noalign{}
\endhead
\bottomrule\noalign{}
\endlastfoot
v31.9 & PR \#356; source
\texttt{27f6ac776c56\allowbreak 2512451000d5\allowbreak 957087d3ecdc\allowbreak d086}; tested main
\texttt{e9b61ca9c359\allowbreak 55028c4bfa42\allowbreak 5dfa278fd186\allowbreak 4029} & PR
\texttt{39569411258}; independent main push \texttt{39570687643}, both
successful \\
v31.10 & PR \#450; source
\texttt{4e5de046ad59\allowbreak eaa37856027a\allowbreak f6547ca26bde\allowbreak a1c2}; tested main
\texttt{dee9bfa68c5f\allowbreak c60bed285335\allowbreak 6c28011b0d3c\allowbreak 53a0} & PR
\texttt{42966223125}; independent main push \texttt{42967294067}, both
successful \\
\end{longtable}

The PR, run, job, and step statuses above were actually read from the
remote service during publication work. The tested v31.10 source and
tested main have identical Git trees. All 155 proof and configuration
inputs, including 97 transitively imported Core sources, and all nine
runtime inputs match their recorded SHA-256 values at the tested source,
tested main, publication baseline, and current working copy. All 5033
tracked Core blobs are unchanged between the initial and current
research baselines. The exact inventories and comparisons are in
\texttt{evidence/\allowbreak v3.\allowbreak 7/\allowbreak source-\allowbreak reuse.\allowbreak json}.

These identities permit reuse of existing heavy scientific receipts
without relabelling them as newly executed publication checks. The
existing v31.10 remote-acceptance receipts record inspection of original
artifact contents. This publication separately checks remote metadata
and source identities; it does not claim to have downloaded and
byte-checked the original ZIP archives again. Newly executed
book-structure checks, bilingual mathematics checks, the teaching
example, PDF builds, and page inspection have separate receipts. If the
subsequent publication PR automatically triggers repository CI, its
outcome is recorded only after it actually occurs.

The current mathematics volume is consulted to locate the maintained
Core and its boundaries; the actual Core-to-Physics contract remains
\texttt{physics/\allowbreak interfaces/\allowbreak CORE\_\allowbreak PHYSICS\_\allowbreak CONTRACT.\allowbreak md}. The seven
evidence levels in §0.10 are a reading aid over the existing chapter
dossier and source manifest, not a competing scientific ledger.
Acceptance of source code, completion of a publication, and support from
nature are distinct outcomes.

{[}QM3114{]} Release kernel and execution reports:
\texttt{physics/\allowbreak qm/\allowbreak UEOT\_\allowbreak QM\_\allowbreak KERNEL\_\allowbreak v31.\allowbreak 14.\allowbreak md} and
\texttt{UEOT\_\allowbreak QM\_\allowbreak EXECUTION\_\allowbreak v31.\allowbreak 14.\allowbreak md}; complete proofs, original
synthetic records, and cross-review:
\texttt{physics/\allowbreak qm/\allowbreak kernel/\allowbreak evidence/\allowbreak v31.\allowbreak 14/\allowbreak }; current premises and full
physics obligations:
\texttt{physics/\allowbreak qm/\allowbreak audits/\allowbreak PERSISTENT\_\allowbreak BATH\_\allowbreak v31.\allowbreak 14/\allowbreak PREMISES\_\allowbreak AND\_\allowbreak RESPONSIBILITIES.\allowbreak md}.

\chapter{Appendix D. Complete inherited verification code and versioned
receipts}\label{appendix-d}

\section{D.1 Execution contract}

Edition 3.0 actually reran the three inherited programs at
2026-10-19T23:13:31.030592+00:00, passing 757 checks in groups of 237,
423, and 97. Editions 3.1, 3.2, 3.3, 3.4, and 3.5 preserved their bytes
and checked their hashes. Edition 3.7 likewise reuses this versioned
execution evidence without relabeling it as a new run. The checks
consist of finite exact arithmetic, deterministic numerical regressions,
and fixed-seed synthetic sampling. They are not 757 universal theorems
or experimental results. The complete code and the actual historical
environment/output are retained below; itemized results are in
\texttt{evidence/\allowbreak v3.\allowbreak 0}.

The Lean formalization, 100 runtime tests, four demonstrations, and
concrete proof witnesses associated with Chapters 49--65 are described
separately in Appendix F and are not combined with these counts. The
following scripts do not access the network or experimental data. They
require Python 3, NumPy, and SciPy, and can be saved and run
individually. Their literal code is preserved, including comments and
historical output strings.

\section{D.2 verify\_\allowbreak ueot\_\allowbreak quantum.py}

Code SHA-256:
79dc78035c239957\allowbreak a8ac4f5b2f6a9d1c\allowbreak 1a6b2a84d2cee094\allowbreak f831bad531f983c0

\begin{Verbatim}[breaklines,breakanywhere,fontsize=\scriptsize]
"""UEOT quantum master: deterministic finite/numerical regression, not formal proof."""
import json
import math
import platform
from fractions import Fraction as F
import numpy as np
import scipy
from scipy.linalg import expm, logm
from scipy.integrate import quad, dblquad
from numpy.polynomial.legendre import leggauss

rng = np.random.default_rng(31011)
checks = []
metrics = {}
I = np.eye(2, dtype=complex)
X = np.array([[0, 1], [1, 0]], complex)
Y = np.array([[0, -1j], [1j, 0]], complex)
Z = np.diag([1, -1]).astype(complex)
paulis = [I, X, Y, Z]

def check(name, value, detail=None):
    ok = bool(value)
    checks.append({"name": name, "passed": ok})
    if not ok:
        raise AssertionError((name, detail))

def close(name, a, b, atol=1e-10, rtol=1e-9):
    check(name, np.allclose(a, b, atol=atol, rtol=rtol),
          float(np.max(np.abs(np.asarray(a)-np.asarray(b)))))

def state(d):
    a = rng.normal(size=(d,d)) + 1j*rng.normal(size=(d,d))
    a = a@a.conj().T
    return a/np.trace(a)

def pure(v):
    return np.outer(v, np.conj(v))

def apply(ks, rho):
    return sum(k@rho@k.conj().T for k in ks)

def trace_out_second(rho, da, db):
    return np.einsum("abcb->ac",rho.reshape(da,db,da,db))

def td(a,b):
    return float(np.sum(np.abs(np.linalg.eigvalsh((a-b+(a-b).conj().T)/2)))/2)

def choi(ks):
    d = ks[0].shape[1]
    blocks = []
    for i in range(d):
        row = []
        for j in range(d):
            e = np.zeros((d,d),complex); e[i,j] = 1
            row.append(apply(ks,e))
        blocks.append(row)
    return np.block(blocks)

def herm_min(a):
    return float(np.linalg.eigvalsh((a+a.conj().T)/2).min())

# Complex channel construction and input-first Choi reconstruction.
for rep in range(12):
    z = rng.normal(size=(6,2))+1j*rng.normal(size=(6,2))
    v,_ = np.linalg.qr(z)
    ks = [v[i*2:(i+1)*2] for i in range(3)]
    j = choi(ks); rho = state(2)
    close(f"channel_TP_{rep}",sum(k.conj().T@k for k in ks),I)
    check(f"channel_Choi_PSD_{rep}",herm_min(j)>-1e-12)
    close(f"channel_Choi_partial_trace_{rep}",trace_out_second(j,2,2),I)
    product = np.kron(rho.T,I)@j
    recovered = np.einsum("iaib->ab",product.reshape(2,2,2,2))
    close(f"channel_reconstruction_{rep}",recovered,apply(ks,rho))
    sig=state(2)
    check(f"channel_trace_contraction_{rep}",
          td(apply(ks,rho),apply(ks,sig)) <= td(rho,sig)+1e-12)

bell=np.array([1,0,0,1],complex)/np.sqrt(2)
bellrho=pure(bell)
pt=bellrho.reshape(2,2,2,2).transpose(0,3,2,1).reshape(4,4)
close("transpose_nonCP_eigenvalue",np.linalg.eigvalsh(pt).min(),-0.5)
rho=state(2)
r=[np.trace(rho@s).real for s in [X,Y,Z]]
close("Pauli_state_tomography",rho,(I+sum(a*s for a,s in zip(r,[X,Y,Z])))/2)

# Rebit deletion witness.
rplus=(np.eye(4)+np.kron(Y,Y))/4
rminus=(np.eye(4)-np.kron(Y,Y))/4
check("rebit_positive",min(herm_min(rplus),herm_min(rminus))>-1e-12)
check("rebit_distinct_real",td(rplus,rminus)>0.99 and np.max(abs(rplus.imag))<1e-12)
for aidx,a in enumerate([I,X,Z]):
    for bidx,b in enumerate([I,X,Z]):
        close(f"rebit_local_blind_{aidx}_{bidx}",np.trace((rplus-rminus)@np.kron(a,b)),0)

# Coherent write / public records / inverse write.
cnot=np.array([[1,0,0,0],[0,1,0,0],[0,0,0,1],[0,0,1,0]],complex)
initial=pure(np.array([1,0,1,0],complex)/np.sqrt(2))
written=cnot@initial@cnot
public=sum(p@written@p for p in [np.diag([1,1,0,0]),np.diag([0,0,1,1])])
close("write_isometry",written,bellrho)
close("public_record_probabilities",np.diag(written),np.diag(public))
close("coherent_parent_distinguishability",td(written,public),0.5)
close("inverse_write_exact",cnot@written@cnot,initial)
close("inverse_write_dephased_X",np.trace((cnot@public@cnot)@np.kron(X,I)),0)

# SRN Gram kernel, block-scalar Kraus, phase-sensitive tomography.
m=3
env=rng.normal(size=(m,4))+1j*rng.normal(size=(m,4))
env/=np.linalg.norm(env,axis=1)[:,None]
K=env@env.conj().T
dims=[2,1,2]; labels=np.repeat(np.arange(m),dims)
ks=[np.diag(env[labels,mu]) for mu in range(4)]
rho=state(5)
expected=K[labels[:,None],labels[None,:]]*rho
close("SRN_block_Kraus_Schur",apply(ks,rho),expected)
for a in range(m):
    idx=np.where(labels==a)[0]
    rr=np.zeros((5,5),complex)
    rr[np.ix_(idx,idx)]=state(len(idx))
    close(f"SRN_fixes_branch_internal_{a}",apply(ks,rr),rr)
close("SRN_composition",apply(ks,apply(ks,rho)),
      (K*K)[labels[:,None],labels[None,:]]*rho)
for a in range(m):
    for b in range(a+1,m):
        for phase in [0,np.pi/2]:
            v=np.zeros(m,complex);v[a]=1/np.sqrt(2);v[b]=np.exp(1j*phase)/np.sqrt(2)
            out=K*pure(v)
            xx=np.zeros((m,m),complex);xx[a,b]=xx[b,a]=1
            close(f"QDRR_phase_{a}_{b}_{phase:.3f}",
                  np.trace(xx@out).real,(np.exp(-1j*phase)*K[a,b]).real)
n=math.ceil(2*math.log(2*m*(m-1)/0.05)/0.12**2)
metrics["QDRR_shots_per_pair_phase"]=n
check("QDRR_sample_count",n==762)

for rep in range(12):
    u=rng.normal(size=(4,3))+1j*rng.normal(size=(4,3))
    G=u@u.conj().T; diag=G.diagonal().real
    theta=rng.normal(size=4)
    lam=G-(diag[:,None]+diag[None,:])/2+1j*(theta[:,None]-theta[None,:])
    C=np.eye(4)-np.ones((4,4))/4
    check(f"SRN_generator_CPSD_{rep}",herm_min(C@lam@C)>-1e-10)
    kt=np.exp(0.4*lam)
    check(f"SRN_generator_exponential_PSD_{rep}",herm_min(kt)>-1e-10)
    close(f"SRN_generator_diagonal_{rep}",kt.diagonal(),np.ones(4))
bad=np.array([[0,-1,-1],[-1,0,-10],[-1,-10,0]],float)
z=np.array([2,-1,-1])
close("invalid_generator_quadratic",z@bad@z,-12)
check("invalid_generator_fails_CP",herm_min(np.exp(0.001*bad))<0)

# Finite CV boundary, Gaussian, spin, Bell and control.
N=9
a=np.diag(np.sqrt(np.arange(1,N)),1)
close("truncated_CCR_boundary",a@a.T-a.T@a,
      np.eye(N)-N*np.diag([0]*(N-1)+[1]))
q=(a+a.T)/np.sqrt(2);p=(a-a.T)/(1j*np.sqrt(2))
close("oscillator_vacuum_uncertainty",(q@q)[0,0]*(p@p)[0,0],0.25)
Omega=np.array([[0,1],[-1,0]],float)
check("Gaussian_vacuum_covariance",herm_min(np.eye(2)/2+1j*Omega/2)>-1e-12)
check("Gaussian_subvacuum_rejected",herm_min(np.eye(2)*0.2+1j*Omega/2)<0)
eta=.7
check("Gaussian_loss_CP",herm_min((1-eta)*np.eye(2)/2+
      1j*(Omega-eta*Omega)/2)>-1e-12)
for t in [.1,.7,1.5]:
    H=(.8*X+.6*Z)/2
    u=expm(-1j*H*t)
    close(f"Rabi_{t}",abs(u[1,0])**2,.8**2*np.sin(t/2)**2)
close("spin_commutator",(X@Y-Y@X)/4,1j*Z/2)
close("spin_time_reversal_square",(1j*Y)@(1j*Y).conj(),-I)
B=np.kron(Z,(Z+X)/np.sqrt(2))+np.kron(Z,(Z-X)/np.sqrt(2))+np.kron(X,(Z+X)/np.sqrt(2))-np.kron(X,(Z-X)/np.sqrt(2))
close("CHSH_Tsirelson",np.vdot(bell,B@bell),2*np.sqrt(2))
Hs=(np.kron(X,X)+np.kron(Y,Y)+np.kron(Z,Z))/4
close("Heisenberg_spectrum",np.linalg.eigvalsh(Hs),[-.75,.25,.25,.25])
H=.3*X+.4*Z; O=Z; psi=np.array([1,0],complex); T=.7
U=expm(-1j*H*T); psit=U@psi
chi=U.conj().T@O@psit
# integrated derivative for constant pulse, compared to central finite difference
def grad_integrand(t):
    pt=expm(-1j*H*t)@psi
    ct=expm(-1j*H*(T-t)).conj().T@O@psit
    return 2*np.imag(np.vdot(ct,X@pt))
g=quad(grad_integrand,0,T)[0]
def objective(e):
    v=expm(-1j*(H+e*X)*T)@psi
    return np.vdot(v,O@v).real
close("control_gradient_sign",g,(objective(1e-6)-objective(-1e-6))/2e-6,atol=1e-8)

# Open dynamics, KMS, causal breaks, thermodynamics.
eta=np.exp(-.7)
damp=[np.diag([1,np.sqrt(eta)]),np.array([[0,np.sqrt(1-eta)],[0,0]])]
close("amplitude_damping_TP",sum(k.conj().T@k for k in damp),I)
close("amplitude_damping_population",apply(damp,np.diag([0,1]))[1,1],eta)
beta=.8; H=np.diag([0,1.3]); th=expm(-beta*H);th/=np.trace(th)
A=X+.2*Z; BB=Y+.3*X; t=.4
def tau(t):
    return expm(1j*H*t)@BB@expm(-1j*H*t)
close("Gibbs_KMS_complex_boundary",np.trace(th@A@tau(t+1j*beta)),
      np.trace(th@tau(t)@A))
def relent(r,s):
    return np.trace(r@(logm(r)-logm(s))).real
def entropy(r):
    ev=np.linalg.eigvalsh(r);ev=ev[ev>1e-14]
    return -np.sum(ev*np.log(ev))
rho=state(2)
free=lambda r:np.trace(r@H).real-entropy(r)/beta
close("Gibbs_free_energy_identity",free(rho)-free(th),relent(rho,th)/beta)
relax=lambda r:.63*r+.37*th*np.trace(r)
check("relative_entropy_stationary_contraction",relent(relax(rho),th)<=relent(rho,th)+1e-12)
# causal break resets S, leaves previous record in E; then E controls S
reverse=np.array([[1,0,0,0],[0,0,0,1],[0,0,1,0],[0,1,0,0]],complex)
outs=[]
for bit in [0,1]:
    inp=np.zeros(4);inp[2*bit]=1
    joint=cnot@pure(inp)@cnot
    env=np.einsum("abad->bd",joint.reshape(2,2,2,2))
    reset=np.kron(np.diag([1,0]),env)
    final=reverse@reset@reverse
    outs.append(trace_out_second(final,2,2))
close("causal_break_memory_witness",td(*outs),1)

# Clifford and Dirac mass shell.
zero=np.zeros((2,2),complex)
gam=[np.block([[I,zero],[zero,-I]])]
gam += [np.block([[zero,s],[-s,zero]]) for s in [X,Y,Z]]
metric=[1,-1,-1,-1]
for mu in range(4):
    for nu in range(4):
        close(f"Clifford_{mu}_{nu}",gam[mu]@gam[nu]+gam[nu]@gam[mu],
              2*metric[mu]*np.eye(4) if mu==nu else np.zeros((4,4)))
def slash(v):
    return sum(metric[i]*v[i]*gam[i] for i in range(4))
pp=np.array([np.sqrt(1+.2**2+.3**2),.2,.3,0])
close("Dirac_mass_shell_factorization",(slash(pp)-np.eye(4))@(slash(pp)+np.eye(4)),np.zeros((4,4)))
qq=np.array([.4,.2,-.1,.3])
close("QED_tree_Ward",slash(qq),(slash(pp+qq)-np.eye(4))-(slash(pp)-np.eye(4)))

# Independent numerical spin traces for e+e- -> mu+mu-.
for costh in [-.7,0,.4,.9]:
    E=2.0; st=np.sqrt(1-costh**2)
    p1=np.array([E,0,0,E]);p2=np.array([E,0,0,-E])
    k1=np.array([E,E*st,0,E*costh]);k2=np.array([E,-E*st,0,-E*costh])
    traces=0j
    for mu in range(4):
        for nu in range(4):
            L=np.trace(slash(p2)@gam[mu]@slash(p1)@gam[nu])/4
            R=np.trace(slash(k1)@gam[mu]@slash(k2)@gam[nu])
            traces+=metric[mu]*metric[nu]*L*R
    close(f"QED_scattering_trace_{costh}",traces/(4*E*E)**2,1+costh**2)
alpha=1/137.035999084
sigma=quad(lambda c:2*np.pi*alpha**2/(4*100)*(1+c*c),-1,1)[0]
close("QED_integrated_cross_section",sigma,4*np.pi*alpha**2/300)
metrics["QED_tree_cross_section_inverse_energy_squared"]=sigma

# phi4 renormalization and EFT.
mass=1.3; cutoff=5.0
numeric=quad(lambda q:q**3/(q*q+mass*mass),0,cutoff)[0]/(8*np.pi**2)
closed=(cutoff**2-mass**2*np.log1p(cutoff**2/mass**2))/(16*np.pi**2)
close("phi4_tadpole_radial_integral",numeric,closed)
f=lambda x:np.log((1+x*(1-x)*4)/(1+x*(1-x)*25))/(16*np.pi**2)
bubble=quad(f,0,1,epsabs=1e-13)[0]
gx,gw=leggauss(100)
close("phi4_bubble_two_quadratures",bubble,np.dot(gw,[f((x+1)/2) for x in gx])/2)
close("phi4_bubble_benchmark",bubble,-.006683040905266517)
metrics["phi4_one_loop_F"]=.4-1.5*.4**2*bubble
s=5.;lam=.2;vel=np.sqrt(1-4/s)
im=lam**2*vel/(32*np.pi)
phase=quad(lambda c:vel/(16*np.pi),-1,1)[0]
close("phi4_optical_cut",2*im,lam**2*phase/2)
mass=.7;M=10.;kap=3.;p2=4.
kernel=np.array([[p2-mass**2,-kap],[-kap,p2-M*M]])
schur=kernel[0,0]-kernel[0,1]*kernel[1,0]/kernel[1,1]
close("EFT_Schur_inverse",1/schur,np.linalg.inv(kernel)[0,0])
errors=[]
for order in range(3):
    approx=kap**2/M**2*sum((p2/M**2)**j for j in range(order+1))
    err=abs(kap**2/(M*M-p2)-approx)
    bound=kap**2/M**2*(p2/M**2)**(order+1)/(1-p2/M**2)
    close(f"EFT_remainder_{order}",err,bound)
    errors.append(err)
metrics["EFT_errors_LO_NLO_NNLO"]=errors
def pi_hat(aq):
    return 2*alpha/np.pi*quad(lambda x:x*(1-x)*np.log1p(aq*x*(1-x)),0,1,epsabs=1e-13)[0]
aq=1e-5
close("QED_polarization_IR",pi_hat(aq)/(alpha*aq/(15*np.pi)),1,atol=2e-6)
aq=1e7
close("QED_polarization_UV",pi_hat(aq),alpha/(3*np.pi)*(np.log(aq)-5/3),atol=1e-8)
def threshold(aq):
    return 4*alpha**2/np.pi*quad(lambda x:aq*x*x*(1-x)**2/(1+aq*x*(1-x)),0,1)[0]
close("QED_threshold_UV",threshold(1e7)/(2*alpha**2/(3*np.pi)),1,atol=2e-6)
aq=3.;step=1e-4
close("QED_threshold_derivative",threshold(aq),
      alpha*(pi_hat(aq*np.exp(2*step))-pi_hat(aq*np.exp(-2*step)))/(2*step),
      atol=1e-11)
schwinger=dblquad(lambda y,z:z/(1-z),0,1,lambda z:0,lambda z:1-z)[0]*alpha/np.pi
close("Schwinger_parameter_integral",schwinger,alpha/(2*np.pi))
metrics["Schwinger_one_loop"]=schwinger

# Higgs and seesaw.
g=.65;gp=.36;v=246.
mat=v*v/4*np.array([[g*g,-g*gp],[-g*gp,gp*gp]])
close("Higgs_neutral_eigenvalues",np.linalg.eigvalsh(mat),[0,v*v*(g*g+gp*gp)/4],atol=1e-9)
md=.1;mr=20.;ev=np.linalg.eigvalsh([[0,md],[md,mr]])
check("seesaw_controlled_error",abs(ev[0]+md*md/mr)<md**4/mr**3*1.01)

# Exact local anomaly functions; rows = multiplicity, C dimension/sign, L dimension/sign, X.
def anomalies(rows):
    out=[F(0) for _ in range(6)]
    for mult,dc,sc,dl,sl,x in rows:
        x=F(x)
        out[0]+=mult*dl*sc
        out[1]+=mult*dc*sl
        out[2]+=mult*dl*(F(1,2) if dc>1 else 0)*x
        out[3]+=mult*dc*(F(1,2) if dl>1 else 0)*x
        out[4]+=mult*dc*dl*x
        out[5]+=mult*dc*dl*x**3
    return out
rows331=[(3,1,0,3,1,F(-1,3)),(3,1,0,1,0,1),
         (1,3,1,3,1,F(1,3)),(2,3,1,3,-1,0),
         (4,3,-1,1,0,F(-2,3)),(5,3,-1,1,0,F(1,3))]
check("331RHN_all_six_local_anomalies",anomalies(rows331)==[0]*6)
metrics["331RHN_local_anomalies"]=[str(x) for x in anomalies(rows331)]
for nc in range(3,16):
    q=F(1,2*nc); u=-F(nc+1,2*nc);d=F(nc-1,2*nc);l=F(-1,2);e=F(1)
    vals=[2*q+u+d,nc*q+l,2*nc*q+nc*u+nc*d+2*l+e,
          2*nc*q**3+nc*u**3+nc*d**3+2*l**3+e**3]
    check(f"generalized_SM_local_anomalies_{nc}",vals==[0]*4)
# Generic beta8=sqrt(3)*b, exact rational grid, including E^c.
for nc in range(3,9):
    nf=2*nc; k=nc-1
    for b in [F(-1),F(-1,3),F(0),F(2,3),F(1)]:
        xl=-F(1,2)-b/2;xp=F(1,6)-b/2;xm=F(1,6)+b/2
        qe=-F(1,2)-3*b/2;jp=F(1,6)-3*b/2;jm=F(1,6)+3*b/2
        rows=[(nf,1,0,3,1,xl),(nf,1,0,1,0,1),(nf,1,0,1,0,-qe),
              (k,nc,1,3,1,xp),(nf-k,nc,1,3,-1,xm),
              (nf,nc,-1,1,0,F(-2,3)),(nf,nc,-1,1,0,F(1,3)),
              (k,nc,-1,1,0,-jp),(nf-k,nc,-1,1,0,-jm)]
        pred=[0,0,0,F(nf*(nc-3),12),0,-F(nf*(nc-3),4)*(1+3*b*b)]
        check(f"generalized_331_full_anomaly_{nc}_{b}",anomalies(rows)==pred)
survivors=[];loose=[]
for nc in range(2,31):
    for nf in range(1,101):
        for k in range(nf+1):
            if 2*nc*k==(nc-1)*nf and nf*(nc+1)<22:
                loose.append((nc,nf,k))
                if nc>=3:survivors.append((nc,nf,k))
check("G3_no_primitive_filter",survivors==[(3,3,1)])
check("G3_complex_color_needed",(2,4,1) in loose)
check("G3_AF_needed",2*3*2==(3-1)*6)
metrics["G3_survivors_without_atomicity"]=survivors
candidates=[(nc,nf) for nc in range(3,23) for nf in range(3,23) if nf*(nc+1)<22]
close("framing_candidate_count",len(candidates),7)
framing={str(chi):[(nc,nf) for nc,nf in candidates if nf*(4*nc+3+chi)%48==0] for chi in [0,1,17,33]}
check("framing_zero_empty",framing["0"]==[])
check("framing_one_unique",framing["1"]==[(3,3)])
check("framing_seventeen_competitor",(4,4) in framing["17"])
check("framing_thirtythree_competitors",all(x in framing["33"] for x in [(3,3),(3,4),(3,5)]))
metrics["framing_survivors"]=framing

result={"summary":{"status":"all_passed","checks":len(checks),"seed":31011,
                   "formal_proof_assistant_checked":False,"experimental_data_used":False},
        "environment":{"python":platform.python_version(),"numpy":np.__version__,"scipy":scipy.__version__},
        "metrics":metrics,"checks":checks}
print(json.dumps(result,ensure_ascii=False,indent=2))
\end{Verbatim}

Actual edition 3.0 result summary; the code remains unchanged in
editions 3.1--3.7:

\begin{Verbatim}[breaklines,breakanywhere,fontsize=\scriptsize]
{
  "summary": {
    "status": "all_passed",
    "checks": 237,
    "seed": 31011,
    "formal_proof_assistant_checked": false,
    "experimental_data_used": false
  },
  "environment": {
    "python": "3.12.14",
    "numpy": "2.3.5",
    "scipy": "1.17.0"
  },
  "metrics": {
    "QDRR_shots_per_pair_phase": 762,
    "QED_tree_cross_section_inverse_energy_squared": 2.230587522067968e-06,
    "phi4_one_loop_F": 0.40160392981726395,
    "EFT_errors_LO_NLO_NNLO": [
      0.0037500000000000033,
      0.00014999999999999736,
      5.999999999992123e-06
    ],
    "Schwinger_one_loop": 0.0011614097328858596,
    "331RHN_local_anomalies": [
      "0",
      "0",
      "0",
      "0",
      "0",
      "0"
    ],
    "G3_survivors_without_atomicity": [
      [
        3,
        3,
        1
      ]
    ],
    "framing_survivors": {
      "0": [],
      "1": [
        [
          3,
          3
        ]
      ],
      "17": [
        [
          3,
          3
        ],
        [
          4,
          4
        ]
      ],
      "33": [
        [
          3,
          3
        ],
        [
          3,
          4
        ],
        [
          3,
          5
        ]
      ]
    }
  }
}
\end{Verbatim}

\section{D.3 verify\_\allowbreak frontier.py}

Code SHA-256:
f0eb63e7d743bb9e\allowbreak 807526a47708fc36\allowbreak b73e6762ab4fbd27\allowbreak 510927f3876094d2

\begin{Verbatim}[breaklines,breakanywhere,fontsize=\scriptsize]
"""Independent checks for UEOT physics frontier. No experimental data."""
import json, math, platform
from fractions import Fraction as F
import numpy as np
import scipy
from scipy.linalg import expm
from scipy.special import logsumexp
from scipy.stats import binom

rng=np.random.default_rng(9062026)
checks=[];metrics={}
def test(name,truth,detail=None):
    if not bool(truth):raise AssertionError((name,detail))
    checks.append({"name":name,"passed":True})
def close(name,a,b,tol=1e-10):
    test(name,np.allclose(a,b,atol=tol,rtol=tol),float(np.max(abs(np.asarray(a)-np.asarray(b)))))
def pure(v):return np.outer(v,v.conj())
def td(a,b):return np.linalg.eigvalsh((a-b+(a-b).conj().T)/2).__abs__().sum()/2
def off(a):return a-np.diag(np.diag(a))
def haar(n):
    z=rng.normal(size=(n,n))+1j*rng.normal(size=(n,n))
    q,r=np.linalg.qr(z)
    return q@np.diag(np.diag(r)/abs(np.diag(r)))
def psd(a):return np.linalg.eigvalsh((a+a.conj().T)/2).min()
def anomalies(rows):
    val=[F(0) for _ in range(6)]
    for n,dc,sc,dl,sl,x in rows:
        x=F(x)
        val[0]+=n*dl*sc
        val[1]+=n*dc*sl
        val[2]+=n*dl*(F(1,2) if dc>1 else 0)*x
        val[3]+=n*dc*(F(1,2) if dl>1 else 0)*x
        val[4]+=n*dc*dl*x
        val[5]+=n*dc*dl*x**3
    return val
def spectrum(nc,nf,b):
    k=F((nc-1)*nf,2*nc)
    assert k.denominator==1
    k=int(k)
    qu=F(nc+1,2*nc);qd=-F(nc-1,2*nc)
    xl=-F(1,2)-b/2;xp=F(1,2*nc)-b/2;xm=F(1,2*nc)+b/2
    qe=-F(1,2)-3*b/2;jp=F(1,2*nc)-3*b/2;jm=F(1,2*nc)+3*b/2
    return [(nf,1,0,3,1,xl),(nf,1,0,1,0,1),(nf,1,0,1,0,-qe),
            (k,nc,1,3,1,xp),(nf-k,nc,1,3,-1,xm),
            (nf,nc,-1,1,0,-qu),(nf,nc,-1,1,0,-qd),
            (k,nc,-1,1,0,-jp),(nf-k,nc,-1,1,0,-jm)]
def survivors(ns):
    ans=[]
    for nc in range(3,31):
        for nf in range(1,121):
            if (nc-1)*nf%(2*nc):continue
            bc=F(11*nc,3)-2*nf
            bl=11-F(nf*(nc+1),3)-F(ns,6)
            if bc>0 and bl>0:ans.append((nc,nf,(nc-1)*nf//(2*nc)))
    return ans
for ns in range(46):
    expected=[(3,3,1),(5,5,2)] if ns<=5 else ([(3,3,1)] if ns<=41 else [])
    test(f"UV_integer_classification_ns_{ns}",survivors(ns)==expected)
metrics["UV_three_scalar_survivors"]=survivors(3)
for nc in range(3,11):
    nf=2*nc
    for b in [F(-1),F(-1,3),F(0),F(1,2)]:
        test(f"generalized_complete_anomalies_{nc}_{b}",anomalies(spectrum(nc,nf,b))==[0]*6)
test("five_color_five_family_all_local_anomalies",anomalies(spectrum(5,5,F(-1,3)))==[0]*6)
metrics["five_color_coefficients"]={"bC":str(F(25,3)),"bL":str(F(1,2)),"low_b2":str(F(-17,6))}
# The zero-U(1) anomaly identities fail if generic charged exotic E^c is omitted.
rows=spectrum(5,5,F(0))
test("omitted_charged_Ec_detected",anomalies(rows[:2]+rows[3:])!=[0]*6)

# Flavor, computed through independent mass-basis rotations.
for rep in range(30):
    uu=haar(3);ud=haar(3);P=np.diag([1,0,0]);V=uu.conj().T@ud
    Pu=uu.conj().T@P@uu;Pd=ud.conj().T@P@ud
    close(f"flavor_CKM_transport_{rep}",Pd,V.conj().T@Pu@V)
    close(f"flavor_rank_one_minor_{rep}",Pu[0,1]*Pu[1,2],Pu[0,2]*Pu[1,1])
    close(f"flavor_cycle_real_{rep}",np.imag(Pu[0,1]*Pu[1,2]*Pu[2,0]),0)
    strength=np.linalg.norm(off(Pu))**2+np.linalg.norm(off(Pd))**2
    test(f"flavor_universal_lower_bound_{rep}",strength>=1-abs(V).max()-1e-12)
    C_u=-Pu**2/2;C_d=-Pd**2/2
    test(f"flavor_EFT_coefficient_bound_{rep}",
         np.sum(abs(off(C_u)))+np.sum(abs(off(C_d)))>=(1-abs(V).max())/2-1e-12)
# Illustrative, frozen angles; these are not a fresh fit to data.
t12,t23,t13,delta=.227,.041,.0035,1.2
c12,c23,c13=np.cos([t12,t23,t13]);s12,s23,s13=np.sin([t12,t23,t13])
V=np.array([[c12*c13,s12*c13,s13*np.exp(-1j*delta)],
            [-s12*c23-c12*s23*s13*np.exp(1j*delta),
             c12*c23-s12*s23*s13*np.exp(1j*delta),s23*c13],
            [s12*s23-c12*c23*s13*np.exp(1j*delta),
             -c12*s23-s12*c23*s13*np.exp(1j*delta),c23*c13]])
close("illustrative_CKM_unitarity",V.conj().T@V,np.eye(3))
metrics["illustrative_flavor_lower_bound_over_delta_g_squared"]=float(1-abs(V).max())

# Three-cycle complete positivity, phase loss, and the classical Poisson counterexample.
for rep in range(24):
    L=rng.normal(size=(3,2))+1j*rng.normal(size=(3,2))
    G=L@L.conj().T;di=np.diag(G).real;h=rng.normal(size=3)
    lam=G-(di[:,None]+di[None,:])/2-1j*(h[:,None]-h[None,:])
    rates=-lam.real
    omega=(lam[0,1]+lam[1,2]+lam[2,0]).imag
    bound=4*rates[0,2]*rates[1,2]-(rates[0,2]+rates[1,2]-rates[0,1])**2
    test(f"cycle_sharp_CP_{rep}",omega**2<=bound+1e-10)
    total=rates[0,1]+rates[1,2]+rates[2,0]
    test(f"cycle_phase_loss_{rep}",total>=np.sqrt(3)*abs(omega)-1e-10)
    K=np.exp(.3*lam)
    test(f"cycle_finite_CP_{rep}",psd(K)>-1e-10)
    determinant=1-abs(K[0,1])**2-abs(K[1,2])**2-abs(K[2,0])**2+2*np.real(K[0,1]*K[1,2]*K[2,0])
    close(f"cycle_determinant_identity_{rep}",determinant,np.linalg.det(K))
theta=np.array([0,2*np.pi/3,4*np.pi/3])
lam=np.exp(1j*(theta[:,None]-theta[None,:]))-1
omega=(lam[0,1]+lam[1,2]+lam[2,0]).imag
close("Poisson_saturates_phase_loss",sum(-lam.real[i,(i+1)%3] for i in range(3)),np.sqrt(3)*abs(omega))
test("Poisson_nonzero_cycle",abs(omega)>2)
metrics["classical_Poisson_cycle"]=float(omega)
bad=np.array([[0,-1+2j,-1],[-1-2j,0,-1],[-1,-1,0]])
test("excess_cycle_not_CP",psd(np.exp(.001*bad))<-1e-6)

# Four-level channel and witness.
V=np.array([[1,0],[0,1],[1,1],[1,1j]],complex);V[2:]/=np.sqrt(2)
K=V@V.conj().T
B=np.array([[-1/np.sqrt(2),-1/np.sqrt(2),1,0],
            [-1/np.sqrt(2),-1j/np.sqrt(2),0,1]],complex)
W=B.conj().T@B;kappa=3-2*np.sqrt(2);wmax=2+1/np.sqrt(2);E=W/wmax
close("four_level_nullspace",B@V,np.zeros((2,2)))
close("four_level_witness_zero",np.trace(W@K),0)
close("witness_frobenius",np.linalg.norm(W),3)
close("witness_operator_norm",np.linalg.eigvalsh(W).max(),wmax)
test("effect_valid",psd(E)>-1e-12 and psd(np.eye(4)-E)>-1e-12)
ks=[np.diag(V[:,mu]) for mu in range(2)]
close("four_level_Kraus_TP",sum(a.conj().T@a for a in ks),np.eye(4))
psi=np.ones(4)/2;rho=pure(psi)
out=sum(a@rho@a.conj().T for a in ks)
close("four_level_output",out,K/4)
joint=np.zeros((8,4),complex)
for a in range(4):joint[a*2:(a+1)*2,a]=V[a]
close("joint_isometry",joint.conj().T@joint,np.eye(4))
jointstate=joint@rho@joint.conj().T
reduced=np.einsum("abcb->ac",jointstate.reshape(4,2,4,2))
close("environment_dilation",reduced,out)
zsat=np.array([1,1,1,(1+1j)/np.sqrt(2)])
close("classical_bound_attained",np.vdot(zsat,W@zsat),kappa)
min_random=1e9
for rep in range(20):
    zz=np.exp(1j*rng.uniform(-np.pi,np.pi,size=(5000,4)))
    vals=np.einsum("ni,ij,nj->n",zz.conj(),W,zz).real
    min_random=min(min_random,float(vals.min()))
    test(f"classical_witness_random_batch_{rep}",vals.min()>=kappa-1e-12)
for eps in [0,.01,.04,kappa/4,.08,1]:
    Ke=(1-eps)*K+eps*np.eye(4)
    close(f"mixture_witness_{eps}",np.trace(W@Ke),4*eps)
pcl=kappa/(4*wmax)
n=math.ceil(math.log(.05)/math.log(1-pcl))
test("zero_event_exact_n",n==188)
test("zero_event_size",binom.cdf(0,n,pcl)<=.05)
test("one_less_shot_insufficient",binom.cdf(0,n-1,pcl)>.05)
metrics["four_level"]={"kappa":kappa,"p_classical_min":pcl,"ideal_zero_event_n95":n,
                      "noise_mixture_threshold":kappa/4,"random_min_witness":min_random,
                      "classical_pvalue_zero_188":float(binom.cdf(0,n,pcl)),
                      "population_leakage_positive_bound_limit":pcl**2/8}
# Verify the approximate population-preserving random-unitary bound on genuine unitary samples.
for rep in range(12):
    us=[]
    for j in range(4):
        h=rng.normal(size=(4,4))+1j*rng.normal(size=(4,4));h=(h+h.conj().T)/2
        us.append(expm(-1j*.002*h))
    weights=rng.dirichlet(np.ones(4))
    surv=sum(p*abs(np.diag(u))**2 for p,u in zip(weights,us))
    ell=1-surv.min()
    rout=sum(p*u@rho@u.conj().T for p,u in zip(weights,us))
    ds=[np.diag(np.diag(u)/abs(np.diag(u))) for u in us]
    rdiag=sum(p*d@rho@d.conj().T for p,d in zip(weights,ds))
    test(f"approximate_RU_distance_{rep}",td(rout,rdiag)<=np.sqrt(8*ell)+1e-12)
    test(f"approximate_RU_effect_{rep}",np.trace(E@rout).real>=pcl-np.sqrt(8*ell)-1e-12)

# Nonlinear ensemble inconsistency.
x1=.8*np.exp(-1)+.2*(-np.exp(-1))
x2=.5*.6*np.exp(-.36)+.5*.6*np.exp(-.36)
close("initial_ensemble_same_bloch",.8-.2,.6)
close("signalling_probability_difference",(x2-x1)/2,.0989390654698766)
metrics["nonlinear_signal_probability_difference"]=(x2-x1)/2

# Exact physical filtering sampler, no Euler step or negative probabilities.
p0=np.array([.2,.3,.5]);aval=np.array([-1.,.2,1.3]);gamma=.7;nsamp=120000
v0=np.dot(p0,aval**2)-np.dot(p0,aval)**2
filter_metrics=[]
for t in [.2,1.,4.]:
    labels=rng.choice(3,size=nsamp,p=p0)
    yt=2*np.sqrt(gamma)*aval[labels]*t+np.sqrt(t)*rng.normal(size=nsamp)
    logs=np.log(p0)[None,:]+2*np.sqrt(gamma)*yt[:,None]*aval-2*gamma*aval**2*t
    pp=np.exp(logs-logsumexp(logs,axis=1)[:,None])
    close(f"filter_normalization_{t}",pp.sum(axis=1),np.ones(nsamp))
    for a in range(3):
        stderr=pp[:,a].std(ddof=1)/np.sqrt(nsamp)
        test(f"filter_Born_martingale_{t}_{a}",abs(pp[:,a].mean()-p0[a])<=6*stderr+1e-4)
    bound=0.
    for a in range(3):
        for b in range(a+1,3):
            samples=np.sqrt(pp[:,a]*pp[:,b])
            theory=np.sqrt(p0[a]*p0[b])*np.exp(-gamma*(aval[a]-aval[b])**2*t/2)
            stderr=samples.std(ddof=1)/np.sqrt(nsamp)
            test(f"filter_pair_decay_{t}_{a}_{b}",abs(samples.mean()-theory)<=6*stderr+1e-4)
            bound+=theory
    variance=pp@(aval**2)-(pp@aval)**2
    test(f"filter_variance_bound_{t}",variance.mean()<=v0/(1+4*gamma*v0*t)+.004)
    err=1-pp.max(axis=1)
    test(f"filter_localization_bound_{t}",err.mean()<=bound+.002)
    filter_metrics.append({"t":t,"mean_p":pp.mean(axis=0).tolist(),
        "mean_variance":float(variance.mean()),"mean_unresolved":float(err.mean()),
        "proved_unresolved_bound":float(bound)})
metrics["exact_filtering"]=filter_metrics
# Heating commutator checked on an interior state, not as false full finite CCR.
a=np.diag(np.sqrt(np.arange(1,12)),1);q=(a+a.T)/np.sqrt(2);p=(a-a.T)/(1j*np.sqrt(2))
H=(p@p+q@q)/2
comm=lambda x,y:x@y-y@x
close("position_noise_heating_interior",-comm(q,comm(q,H))[0,0]/2,.5)
result={"summary":{"status":"all_passed","checks":len(checks),"seed":9062026,
        "experimental_data":False,"formal_proof_assistant":False},
        "environment":{"python":platform.python_version(),"numpy":np.__version__,"scipy":scipy.__version__},
        "metrics":metrics,"checks":checks}
print(json.dumps(result,ensure_ascii=False,indent=2))
\end{Verbatim}

Actual edition 3.0 result summary; the code remains unchanged in
editions 3.1--3.7:

\begin{Verbatim}[breaklines,breakanywhere,fontsize=\scriptsize]
{
  "summary": {
    "status": "all_passed",
    "checks": 423,
    "seed": 9062026,
    "experimental_data": false,
    "formal_proof_assistant": false
  },
  "environment": {
    "python": "3.12.14",
    "numpy": "2.3.5",
    "scipy": "1.17.0"
  },
  "metrics": {
    "UV_three_scalar_survivors": [
      [
        3,
        3,
        1
      ],
      [
        5,
        5,
        2
      ]
    ],
    "five_color_coefficients": {
      "bC": "25/3",
      "bL": "1/2",
      "low_b2": "-17/6"
    },
    "illustrative_flavor_lower_bound_over_delta_g_squared": 0.0008465021129667738,
    "classical_Poisson_cycle": -2.598076211353316,
    "four_level": {
      "kappa": 0.1715728752538097,
      "p_classical_min": 0.0158446719248555,
      "ideal_zero_event_n95": 188,
      "noise_mixture_threshold": 0.04289321881345243,
      "random_min_witness": 0.17513824113339804,
      "classical_pvalue_zero_188": 0.04965532503989115,
      "population_leakage_positive_bound_limit": 3.1381703550788005e-05
    },
    "nonlinear_signal_probability_difference": 0.0989390654698766,
    "exact_filtering": [
      {
        "t": 0.2,
        "mean_p": [
          0.19907099901956785,
          0.29966622147876876,
          0.5012627795016538
        ],
        "mean_variance": 0.5373589168712367,
        "mean_unresolved": 0.40514124557832,
        "proved_unresolved_bound": 0.7956706280194334
      },
      {
        "t": 1.0,
        "mean_p": [
          0.19923581748426195,
          0.30010875188535324,
          0.5006554306303944
        ],
        "mean_variance": 0.2081865339186897,
        "mean_unresolved": 0.21236481946950928,
        "proved_unresolved_bound": 0.4512082432205973
      },
      {
        "t": 4.0,
        "mean_p": [
          0.19889313068851436,
          0.2987103279415504,
          0.5023965413698956
        ],
        "mean_variance": 0.034615822280499135,
        "mean_unresolved": 0.03598012354677628,
        "proved_unresolved_bound": 0.10399501336004843
      }
    ]
  }
}
\end{Verbatim}

\section{D.4 verify\_\allowbreak completion.py}

Code SHA-256:
c2107ac235f0d87f\allowbreak aee246412c6478a0\allowbreak 405113cb0113b647\allowbreak 027820106d95ad32

\begin{Verbatim}[breaklines,breakanywhere,fontsize=\scriptsize]
"""Independent finite checks for the unified additions; no experimental data."""
import itertools, json, platform
import numpy as np
import scipy
from scipy.integrate import quad
from scipy.linalg import expm
from scipy.stats import binom

rng=np.random.default_rng(9062048)
checks=[];metrics={}
def test(name,ok):
    if not bool(ok): raise AssertionError(name)
    checks.append(name)
def close(name,a,b,tol=1e-10):
    test(name,np.max(np.abs(np.asarray(a)-np.asarray(b)))<tol)

# Spectral gauge transport and flux periodicity, not finite exact CCR.
n=np.arange(-8,9);hbar=1.3;L=2.1;m=.7;q=.8
for j in range(5):
    theta,A,chi=rng.normal(size=3)
    energy=lambda th,aa: (hbar*(2*np.pi*n+th)/L-q*aa)**2/(2*m)
    close(f"ring_gauge_transport_{j}",energy(theta,A),energy(theta+q*chi*L/hbar,A+chi))
    close(f"ring_spectral_relabel_{j}",energy(theta+2*np.pi,A)[:-1],energy(theta,A)[1:])
close("Berry_Chern_integral",quad(lambda t:-.5*np.sin(t),0,np.pi)[0],-1)
for theta in [.2,.9,2.5]:
    close(f"Berry_patch_difference_{theta}",.5*(1+np.cos(theta))+.5*(1-np.cos(theta)),1)

for ratio in [.2,.7,1.4,2.0]:
    k=np.linspace(-np.pi,np.pi,10001)
    z=ratio+np.exp(-1j*k)
    winding=(np.unwrap(np.angle(z))[-1]-np.unwrap(np.angle(z))[0])/(2*np.pi)
    close(f"SSH_winding_{ratio}",winding,-1 if ratio<1 else 0)
for N,r in [(8,.4),(12,.7),(24,.8)]:
    H=np.zeros((2*N,2*N))
    for j in range(N):
        H[2*j,2*j+1]=H[2*j+1,2*j]=r
        if j+1<N:H[2*j+1,2*j+2]=H[2*j+2,2*j+1]=1.
    norm=np.sqrt((1-r*r)/(1-r**(2*N)))
    psi=np.zeros(2*N);psi[::2]=norm*(-r)**np.arange(N)
    residual=norm*r**N
    close(f"SSH_normalization_{N}",np.linalg.norm(psi),1)
    close(f"SSH_finite_residual_{N}",np.linalg.norm(H@psi),residual)
    test(f"SSH_near_zero_spectrum_{N}",min(abs(np.linalg.eigvalsh(H)))<=residual+1e-12)

# Direct noncommutative plaquette gauge covariance using SU(2) matrices.
pauli=[np.array([[0,1],[1,0]]),np.array([[0,-1j],[1j,0]]),np.diag([1,-1])]
def su2():
    v=rng.normal(size=3)
    return expm(1j*sum(v[i]*pauli[i] for i in range(3)))
for j in range(12):
    U=[su2() for _ in range(4)];G=[su2() for _ in range(4)]
    transformed=[G[k]@U[k]@G[(k+1)%4].conj().T for k in range(4)]
    prod=lambda a:a[0]@a[1]@a[2]@a[3]
    close(f"plaquette_gauge_covariance_{j}",prod(transformed),G[0]@prod(U)@G[0].conj().T)
    close(f"plaquette_trace_invariance_{j}",np.trace(prod(transformed)),np.trace(prod(U)))
F=(.3*pauli[0]+.7*pauli[2])/2;g=.8
for a in [.1,.05,.025]:
    exact=(1-np.trace(expm(1j*g*a*a*F)).real/2)*4/g**2
    predicted=a**4*np.trace(F@F).real
    test(f"Wilson_continuum_coefficient_{a}",abs(exact-predicted)/predicted<2e-6)

configs=np.array(list(itertools.product([-1,1],repeat=4)))
P=configs.prod(axis=1)
for beta in [0,.2,.7,1.4,3.]:
    weights=np.exp(beta*P);Z=weights.sum()
    close(f"Z2_partition_{beta}",Z,16*np.cosh(beta))
    close(f"Z2_plaquette_{beta}",weights@P/Z,np.tanh(beta))
    close(f"Z2_link_expectation_{beta}",weights@configs/Z,np.zeros(4))
for vertex in range(4):
    gauged=configs.copy();gauged[:,vertex]*=-1;gauged[:,(vertex-1)%4]*=-1
    close(f"Z2_vertex_gauge_{vertex}",gauged.prod(axis=1),P)

# Enumerate all histories in a genuinely history-dependent binary model.
p0=.13;ntrial=10;qalt=.02
histories=[((),1.)]
for t in range(ntrial):
    new=[]
    for hist,prob in histories:
        pt=p0+.3*((sum(hist)+t)%3)/2
        new += [(hist+(0,),prob*(1-pt)),(hist+(1,),prob*pt)]
    histories=new
distribution=np.zeros(ntrial+1)
for hist,prob in histories:distribution[sum(hist)]+=prob
close("history_tree_normalization",distribution.sum(),1)
for k in range(ntrial+1):
    test(f"dependent_binomial_domination_{k}",distribution[:k+1].sum()<=binom.cdf(k,ntrial,p0)+1e-12)
for pt in [p0,.3,.7,1.]:
    factor=pt*qalt/p0+(1-pt)*(1-qalt)/(1-p0)
    close(f"supermartingale_conditional_identity_{pt}",factor,1+(pt-p0)*(qalt-p0)/(p0*(1-p0)))
    test(f"supermartingale_factor_{pt}",factor<=1+1e-12)
likelihood=lambda k:(qalt/p0)**k*((1-qalt)/(1-p0))**(ntrial-k)
test("dependent_likelihood_mean",sum(distribution[k]*likelihood(k) for k in range(ntrial+1))<=1+1e-12)
test("marginal_bound_counterexample",1-p0>(1-p0)**ntrial)
pstar=(3-2*np.sqrt(2))/(4*(2+1/np.sqrt(2)))
test("ideal_zero_188_valid",(1-pstar)**188<.05)
test("ideal_zero_187_insufficient",(1-pstar)**187>.05)
metrics.update({"ideal_pstar":pstar,"zero_event_188_pvalue":(1-pstar)**188,
    "Berry_upper_band_Chern":-1,"history_model_zero_probability":distribution[0],
    "history_null_zero_bound":(1-p0)**ntrial})
result={"summary":{"status":"all_passed","checks":len(checks),"seed":9062048,
    "experimental_data":False,"formal_proof_assistant":False},
    "environment":{"python":platform.python_version(),"numpy":np.__version__,"scipy":scipy.__version__},
    "metrics":metrics,"checks":checks}
print(json.dumps(result,ensure_ascii=False,indent=2))
\end{Verbatim}

Actual edition 3.0 result summary; the code remains unchanged in
editions 3.1--3.7:

\begin{Verbatim}[breaklines,breakanywhere,fontsize=\scriptsize]
{
  "summary": {
    "status": "all_passed",
    "checks": 97,
    "seed": 9062048,
    "experimental_data": false,
    "formal_proof_assistant": false
  },
  "environment": {
    "python": "3.12.14",
    "numpy": "2.3.5",
    "scipy": "1.17.0"
  },
  "metrics": {
    "ideal_pstar": 0.0158446719248555,
    "zero_event_188_pvalue": 0.04965532503989112,
    "Berry_upper_band_Chern": -1,
    "history_model_zero_probability": 0.039600353867614305,
    "history_null_zero_bound": 0.24842341419143568
  }
}
\end{Verbatim}

\chapter{Appendix E. Reading editions, compilation, and revision
contract}\label{appendix-e}

\section{E.1 Complete semantic sources and derived reading editions}

The maintained Chinese source \texttt{UEOT\_\allowbreak PHYSICS.\allowbreak md} and its complete
English counterpart \texttt{UEOT\_\allowbreak PHYSICS.\allowbreak en.\allowbreak md} each contain Chapters
0--70 and Appendices A--F. Their corresponding \texttt{.\allowbreak tex} and
\texttt{.\allowbreak pdf} files are actual derivatives of those sources, with
mathematical typesetting, a full chapter contents list, stable result
IDs, original derivations, and embedded code. Edition 3.11 adds the
Chapter 70 persistent-environment processes, finite certification, and
selection limits to the continuous bilingual manuscript while retaining
Chapter 69 and all earlier results. The edition number identifies a
publication; the text and evidence determine which physical questions
are resolved or open.

Both languages carry the same scientific commitments. Their prose is
adapted for natural reading, while assumptions, formulas, constants,
result IDs, source identities, and empirical qualifications are checked
together. Mathematical symbols are not silently changed in translation.
Purely linguistic labels inside formulas and recorded editorial
corrections are listed in the fidelity report.

\section{E.2 Reproducible compilation}

The existing builder \texttt{physics/\allowbreak qm/\allowbreak book/\allowbreak build/\allowbreak build\_\allowbreak pdf.\allowbreak py}
compiles one complete current source at a time. It does not invoke an
old assembler or silently fetch another manuscript. It requires Pandoc,
XeLaTeX, ordinary TeX packages, and suitable fonts. The default
publication uses Noto Serif CJK SC for the main text and Noto Sans CJK
SC for supplementary roles, including the inherited source-code
comments. A local CJK font file can be supplied through
\texttt{-\allowbreak \/\allowbreak -\allowbreak font}. Shell escape is disabled. Font identity is part of
the build receipt, not part of the physical theory.

\begin{Verbatim}[breaklines,breakanywhere,fontsize=\scriptsize]
python3 physics/qm/book/build/build_pdf.py \
  --source physics/qm/book/UEOT_PHYSICS.md \
  --output-dir /tmp/ueot-physics-zh-build

python3 physics/qm/book/build/build_pdf.py \
  --source physics/qm/book/UEOT_PHYSICS.en.md \
  --output-dir /tmp/ueot-physics-en-build
\end{Verbatim}

The builder first checks consecutive chapters, matching appendices,
stable anchors, internal links, and incomplete authoring placeholders.
It produces a print source and LaTeX, then runs XeLaTeX three times to
stabilize the contents list. Source, font, TeX, PDF, and logs are
identified separately. The resulting pages must also be inspected for
text, equations, tables, and code wrapping; a successful exit code alone
does not complete the publication review.

\section{E.3 Maintenance rules}

New theoretical work belongs in the relevant chapters and result
registry. Code changes must be reflected in the formalization scope and
actual receipts. Changes to cross-chapter dependencies belong in
\texttt{CHAPTER\_\allowbreak DOSSIER.\allowbreak json} and the review record. The publication's
status depends on continuous content, traceable evidence, and completed
checks; it does not require unresolved physical questions to be written
as solved.

PDF hashes are stored in external receipts rather than embedded in the
PDFs themselves. Remote source commits, PRs, and main workflows have
their own observed records. A local build is not called a remote
execution. The current revision reuses the verified, unchanged
scientific-source identities described in §C.8 and distinguishes them
from the newly executed publication checks.

\section{E.4 Rights and the scope of this revision}

The theoretical author is Degui Qian. The English translation,
scientific editing, and three specialist reviews were produced with AI
assistance under the author's project instructions; they do not
establish outside endorsement. Historical expression already validly
licensed under CC BY 4.0 retains that license, including its lawful
adaptations. New edition-specific expression and translation are
identified by this revision's source manifest and change history. This
edition makes no additional blanket license grant for those new
contributions; rights, to the extent they subsist, remain with the
respective rightsholders pending an express grant. It makes no claim of
exclusive ownership over physical laws, mathematical results, or
third-party material.

\texttt{PUBLICATION\_\allowbreak RIGHTS.\allowbreak md} provides the edition-specific inventory
and the preserved historical baseline. Referenced articles, figures,
datasets, software, and dependencies retain their own terms. This
revision does not redistribute new third-party figures or raw
experimental data and does not change the licenses of Lean or numerical
code. The same notice appears in the Chinese and English sources and
their generated reading editions.

\chapter{Appendix F. Formalized scope and the complete named axiom
audit}\label{appendix-f}

\section{F.1 Verified identity and counts}

The inherited v31.9 proof evidence uses Lean 4.33.1 (compiler commit
\texttt{819816b2e0a3\allowbreak bf405af45ae5\allowbreak c7af2491d8f5\allowbreak bee6}) and Mathlib
\texttt{0df444a360ea\allowbreak a60ab8c11dca\allowbreak 51a86af69295\allowbreak 5474}; the identities of
nine Git dependencies matched the lock file. The unified build of 41
modules exited with code zero after 8805 jobs. The evidence checks 854
explicit theorems and 1123 unique fully named axiom reports. The only
allowed foundations are \texttt{propext}, \texttt{Classical.\allowbreak choice}, and
\texttt{Quot.\allowbreak sound}. Ordinary linter warnings are retained; a
warning-free build is not claimed. These counts describe historical
v31.9 acceptance; F.4--F.5 give the current v31.10 increment and totals.
Publication revision 3.7 reuses the source-bound heavy proof evidence
locally; any actual automatic CI has its own remote receipt.

The six modules added in v31.9 contain 107 explicit theorems and 138
named audit targets; the 985 inherited requests were preserved. Explicit
declarations, named nonlocal instances, and actual structure
constructors are registered from the source. The 1123 targets consist of
854 explicit theorems and 269 other listed targets. They do not
enumerate every compiler-generated internal constant and are not a count
of new physical laws.

\begin{longtable}[]{@{}
  >{\raggedright\arraybackslash}p{(\columnwidth - 4\tabcolsep) * \real{0.3000}}
  >{\raggedleft\arraybackslash}p{(\columnwidth - 4\tabcolsep) * \real{0.1200}}
  >{\raggedright\arraybackslash}p{(\columnwidth - 4\tabcolsep) * \real{0.5800}}@{}}
\toprule\noalign{}
\begin{minipage}[b]{\linewidth}\raggedright
Module
\end{minipage} & \begin{minipage}[b]{\linewidth}\raggedleft
Explicit theorems
\end{minipage} & \begin{minipage}[b]{\linewidth}\raggedright
Exact responsibility
\end{minipage} \\
\midrule\noalign{}
\endhead
\bottomrule\noalign{}
\endlastfoot
Basic & 7 & Density operators, positive trace, and normalization. \\
Instrument & 14 & Actual Kraus branches, probabilities, and
posteriors. \\
Operational & 20 & All-record responses, response quotients, and
reachable updates. \\
Observables & 25 & Actual word effects, minimal closure, and finite
generation. \\
FiniteCoordinates & 21 & Local presentations, all-word coordinates, and
closed posterior updates. \\
Adaptive & 11 & Complete finite adaptive trees, leaf laws, and
coordinates. \\
Weighted & 16 & Nonnegative weights converted to ordinary complex Kraus
operations. \\
Ragged & 19 & Outcome-dependent finite indices and preservation under
zero-padding. \\
Exact & 24 & Injectivity and algebraic preservation of complex rational
interpretation. \\
ExactCertificate & 9 & Soundness of a finite computable checker and
instance interfaces. \\
Transport & 4 & Scalar instances of existing Core error transport. \\
Selector & 11 & Conditional generation selection over all integers and
premise-deletion counterexamples. \\
Spectrum & 15 & Actual finite spectrum rows mapped to anomaly and
effective-window budgets. \\
Effects & 30 & All actual effects, branch/channel pullbacks, and
word-effect bounds. \\
Stability & 30 & All-effect pseudodistance, channel defects,
trajectories, and posteriors. \\
Mixture & 23 & Actual mixtures, exact TV bounds, and rational
budgets. \\
Cesaro & 23 & Matrix-average residuals, compactness, continuity, and
stationary subsequences. \\
Survival & 32 & Actual diagonal survival records and use of Core
QSD/Doob results. \\
ClassicalSector & 24 & The full classical simplex, all effects, and
purification obstruction for every finite ancilla. \\
CorePurificationBoundary & 32 & One actual Core kernel, terminal
quotient, safe trajectory, and purification separation. \\
RealLocalTomography & 37 & Indistinguishability under all real local
effects and joint/complex-local deletion witnesses. \\
QuantumDoob & 32 & Actual Doob completeness, finite conditioning, and
stationarity with additional spectral premises. \\
RecordProbability & 9 & TV attained by actual finite record events and
equivalence over all events. \\
RectangularChannels & 29 & Rectangular Kraus channels, arbitrary finite
ancillary amplification, partial trace, and effect contraction. \\
QuantumPurification & 25 & Full convex extremality, Gram factors derived
from PSD, and purification of every finite density operator. \\
OperationalSystems & 25 & Three complete state-space presentations,
seven laws, and independent candidate principles. \\
ClassicalQuantumMemory & 39 & Actual joint records, finite-tree memory,
all-event TV, and label defects. \\
QuantumDoobModel & 17 & Extraction of the previous nondiagonal model and
explicit public states; no independent new physics. \\
QuantumDoobDynamics & 43 & Closed forms for specified success records,
exact conditional-state distance decay, and probability limits. \\
QuantumSeparation & 9 & The Hermitian real span of valid effects, full
density separation, and the zero-distance criterion. \\
ComplexLocalTomography & 11 & Actual product effects, real/imaginary
matrix units, and local tomography for arbitrary joint complex
states. \\
PureStateCharacterization & 5 & Full convex-pure-state vector
characterization from a PSD positive-column decomposition. \\
PurificationUniqueness & 22 & Equal-Gram support isometry,
common-ancilla unitaries, rectangular embeddings, and actual tensor
actions. \\
HistoryChannels & 23 & Actual history-reading channels, common PSD block
weights, zero blocks, and valid effect bounds. \\
ActionRecordBudget & 31 & Original fixed-depth protocol semantics,
step/action error sums, and complete-record TV. \\
QuantumTraceDistance & 25 & Positive/negative Hermitian spectral parts,
valid optimal effects, and the half-trace-absolute-value formula. \\
BinaryDiscrimination & 9 & Optimal equal-prior binary discrimination,
PSD-support certificates, and actual unitary spectral data. \\
PurificationIsometry & 9 & Necessary and sufficient ancillary dimension
order, directed isometries, and actual purification maps. \\
PurificationDimension & 14 & Positive spectral-support Gram
construction, rank lower bounds, and the exact minimum-ancilla
criterion. \\
StoppingProtocols & 21 & Actual Option instruments, causal absorbing
control, original-leaf encoding, and all-matrix record identities. \\
StoppingRecordBudget & 29 & Isometric extension of valid effects,
transport to original-leaf joint states, and active-weight action
budgets. \\
\end{longtable}

\section{F.2 Historical execution and evidence reuse}

The v31.9 local acceptance actually built the main package on the frozen
98-file source set, checked all 1123 public named audits, and compiled
the three newly added witnesses. Those three witnesses contain 105
concrete theorems, 39 definitions, three abbreviations, and 147 named
audit targets.

The ten older witnesses retain their original actual receipts for 193
theorems, 71 definitions, five abbreviations, and 269 audits. The v31.9
reuse review checked 86 unchanged scientific inputs, 56 Core sources,
nine dependency versions, and the actual transitive import closures of
the older witnesses and certificates. New umbrella imports, the public
audit list, and witness registration/selection interfaces were outside
those older mathematical import closures. Six exact certificates with 42
audits, four classes of rejected proofs asserting a true checker result
with successful proofs that the checker returns false on the same
corrupted data, 100 runtime methods, and four demonstrations therefore
retained the verified v31.8 local and main evidence. They were not
reported as rerun in that local acceptance. \texttt{legacy-\allowbreak reuse.\allowbreak json}
records the exact scope and identities. The 757 inherited text-program
checks separately retain their edition 3.0 execution identity.

The complete witness registry has thirteen groups: 298 concrete
theorems, 110 definitions, eight abbreviations, and 416 named targets.
The ten inherited and three newly executed groups are distinguished;
their counts are not added to the main package's 854/1123. The new
witnesses cover a noncommuting rational distance, optimal binary
decisions, a complex-phase restriction negative control, matrix
direction and minimum rank under different ancilla dimensions, an exact
stopped-tree budget of 1/16, and record TV equal to one for different
stopping rules. These are formal mathematical instances, not new
experiments.

The v31.9 source-set SHA-256 is
\texttt{bf026d2a1f7d\allowbreak 8c8bd1f9f6c1\allowbreak ad9015aa20bd\allowbreak ae6bd36b9d43\allowbreak 2c7ce6f0581e\allowbreak ffc0};
its 98 scientific input files retained their bytes through that
acceptance. The local source/olean binding and the 56 recursively
imported Core sources have their own receipts. This is not reacceptance
of all 106 Core responsibilities. Local execution, historical reuse,
PR/main CI, and later metadata are distinct evidence events.

For this publication, the repository was refreshed to
\texttt{3d443ec29eb1\allowbreak 7af1fc00699a\allowbreak 74e2d42d7351\allowbreak 6d81}.
\href{https://github.com/MurphyHoops/UEOT/pull/356}{PR \#356} had
already merged source \texttt{27f6ac776c56\allowbreak 2512451000d5\allowbreak 957087d3ecdc\allowbreak d086}
as main \texttt{e9b61ca9c359\allowbreak 55028c4bfa42\allowbreak 5dfa278fd186\allowbreak 4029}. The actual
\href{https://github.com/MurphyHoops/UEOT/actions/runs/39569411258}{PR
run 39569411258} and independent
\href{https://github.com/MurphyHoops/UEOT/actions/runs/39570687643}{main
push run 39570687643} both report successful runtime and Lean jobs. All
98 scientific files and all 56 imported Core files match their recorded
hashes in the source commit, merged commit, and refreshed main. The
older \texttt{LOCAL\_\allowbreak VERIFIED\_\allowbreak REMOTE\_\allowbreak PENDING} metadata is therefore a
historical checkpoint, not the observed publication state. This review
checked repository bytes and workflow/job metadata; it did not inspect
the archived artifact ZIPs or rerun the heavy proof suite.

\section{F.3 All 1123 named axiom requests}

The following reproduces every request from the actual v31.9
\texttt{Axiom\allowbreak Audit.\allowbreak lean}, in order. Each request must have a unique
identically named actual output. Missing or duplicate outputs,
unreviewed axioms, or build errors are failures. The declaration types
and premises in Chapters 49--65 determine scientific scope. The literal
request blocks remain unchanged in this bilingual publication.

\subsection{F.3.1 Requests 1--60}

\begin{Verbatim}[breaklines,breakanywhere,fontsize=\scriptsize]
UEOTQM.Density.ext
UEOTQM.mass_nonneg
UEOTQM.trace_im_eq_zero
UEOTQM.mass_smul
UEOTQM.Density.mass_one
UEOTQM.mass_eq_zero_iff
UEOTQM.normalize_matrix
UEOTQM.krausMap_posSemidef
UEOTQM.krausMap_smul
UEOTQM.trace_krausMap
UEOTQM.Instrument.branch_posSemidef
UEOTQM.Instrument.branch_smul
UEOTQM.Instrument.sum_trace_branch
UEOTQM.Instrument.sum_mass_branch
UEOTQM.Instrument.probability_nonneg
UEOTQM.Instrument.sum_probability
UEOTQM.Instrument.probability_le_one
UEOTQM.Instrument.posterior_matrix
UEOTQM.Instrument.zero_probability_iff_zero_branch
UEOTQM.ancillaKrausMap_posSemidef
UEOTQM.ancillaKrausMap_entry
UEOTQM.recordOperator_posSemidef
UEOTQM.recordOperator_smul
UEOTQM.response_nonneg
UEOTQM.response_nil
UEOTQM.response_singleton
UEOTQM.totalRecordMass_eq_mass
UEOTQM.totalRecordMass_density
UEOTQM.responseClass_surjective
UEOTQM.existsUnique_responseDecoder
UEOTQM.probability_eq_of_response_eq
UEOTQM.response_posterior
UEOTQM.posterior_response_eq
UEOTQM.nextResponse_eq_none_iff
UEOTQM.nextResponse_eq_of_response_eq
UEOTQM.nextResponse_inputFiberCompatible
UEOTQM.existsUnique_reachableResponseUpdate
UEOTQM.Selector.selector_finite_bounds
UEOTQM.Selector.selects_three_without_atomicity
UEOTQM.Selector.dropping_color_bound_admits_four_families
UEOTQM.Selector.dropping_weak_budget_admits_six_families
UEOTQM.Selector.dropping_cubic_equation_admits_one_family
UEOTQM.Selector.dropping_nonempty_family_count_admits_empty_spectrum
UEOTQM.Selector.homogeneous_leptons_recover_cubic_equation
UEOTQM.Selector.mixed_lepton_template_admits_two_families
UEOTQM.Selector.cubic_cancellation_replicates
UEOTQM.Selector.replicated_weak_screening
UEOTQM.Selector.two_or_more_three_family_packets_fail_weak_budget
UEOTQM.Transport.scalarBudget_zero
UEOTQM.Transport.scalarBudget_succ
UEOTQM.Transport.scalarBudget_nonneg
UEOTQM.Transport.scalar_transport_certificate
UEOTQM.normalize
UEOTQM.Instrument.posterior
UEOTQM.nextReachableResponse_map_val
UEOTQM.nextReachableResponse_eq_of_response_eq
UEOTQM.nextReachableResponse_inputFiberCompatible
UEOTQM.existsUnique_closedReachableUpdate
UEOTQM.nextReachableResponse
UEOTQM.pairing_one
\end{Verbatim}

\subsection{F.3.2 Requests 61--120}

\begin{Verbatim}[breaklines,breakanywhere,fontsize=\scriptsize]
UEOTQM.pairing_zero_left
UEOTQM.pairing_add_left
UEOTQM.pairing_smul_left
UEOTQM.pairing_smul_right
UEOTQM.pairing_sum_left
UEOTQM.heisenberg_zero
UEOTQM.heisenberg_add
UEOTQM.heisenberg_smul
UEOTQM.trace_heisenberg
UEOTQM.pairing_heisenberg
UEOTQM.heisenberg_posSemidef
UEOTQM.recordEffect_posSemidef
UEOTQM.recordEffect_isHermitian
UEOTQM.mass_recordOperator_eq_pairing
UEOTQM.response_eq_pairing
UEOTQM.recordEffect_mem
UEOTQM.one_mem_observableSpace
UEOTQM.recordEffect_mem_of_closed
UEOTQM.observableSpace_le_of_closed
UEOTQM.observableSpace_pullbackClosed
UEOTQM.isHermitian_of_mem_observableSpace
UEOTQM.response_eq_iff_pairing_observableSpace
UEOTQM.observableSpace_finiteDimensional
UEOTQM.exists_finite_wordEffect_generators
UEOTQM.FiniteObservablePresentation.observable_mem_span
UEOTQM.FiniteObservablePresentation.observable_unit
UEOTQM.FiniteObservablePresentation.one_mem_span
UEOTQM.FiniteObservablePresentation.pullback_observable
UEOTQM.FiniteObservablePresentation.span_pullbackClosed
UEOTQM.FiniteObservablePresentation.span_eq_observableSpace
UEOTQM.FiniteObservablePresentation.recordEffect_mem_span
UEOTQM.FiniteObservablePresentation.operatorCoordinates_unit
UEOTQM.FiniteObservablePresentation.coordinates_unit
UEOTQM.FiniteObservablePresentation.coordinates_apply
UEOTQM.FiniteObservablePresentation.coordinates_eq_iff_response_eq
UEOTQM.FiniteObservablePresentation.operatorCoordinates_branch
UEOTQM.FiniteObservablePresentation.operatorCoordinates_record
UEOTQM.FiniteObservablePresentation.response_eq_recordCoordinates_unit
UEOTQM.FiniteObservablePresentation.probability_eq_applyBranch_unit
UEOTQM.FiniteObservablePresentation.coordinates_posterior
UEOTQM.FiniteObservablePresentation.coordinates_posterior_eq_normalizedBranch
UEOTQM.FiniteObservablePresentation.normalizedBranch_mem_range
UEOTQM.FiniteObservablePresentation.closedCoordinateUpdate_map_val
UEOTQM.FiniteObservablePresentation.coordinateUpdate_eq_posterior
UEOTQM.FiniteObservablePresentation.coordinateUpdate_eq_none_iff
UEOTQM.Adaptive.Protocol.word_injective
UEOTQM.Adaptive.sum_leafWeight
UEOTQM.Adaptive.leafWeight_nonneg
UEOTQM.Adaptive.leafProbability_eq_response
UEOTQM.Adaptive.leafProbability_nonneg
UEOTQM.Adaptive.sum_leafProbability
UEOTQM.Adaptive.leafProbability_le_one
UEOTQM.Adaptive.leafProbability_eq_of_response_eq
UEOTQM.Adaptive.leafProbability_eq_recordCoordinates_unit
UEOTQM.Adaptive.leafProbability_eq_of_coordinates_eq
UEOTQM.Adaptive.leafProbability_posterior
UEOTQM.Spectrum.quarkCount_le_length
UEOTQM.Spectrum.familyCubic_eq_counts
UEOTQM.Spectrum.leptonCount_eq_length_of_homogeneous
UEOTQM.Spectrum.ExtraMatter.screening_nonneg
\end{Verbatim}

\subsection{F.3.3 Requests 121--180}

\begin{Verbatim}[breaklines,breakanywhere,fontsize=\scriptsize]
UEOTQM.Spectrum.extraScreening_nonneg
UEOTQM.Spectrum.familyScreening_eq_of_active
UEOTQM.Spectrum.selector_anomaly_of_spectrum
UEOTQM.Spectrum.selector_budget_of_spectrum
UEOTQM.Spectrum.spectrum_selects_three
UEOTQM.Spectrum.canonical_rows_and_higgs
UEOTQM.Spectrum.mixed_two_family_rows
UEOTQM.Spectrum.replicated_six_family_rows
UEOTQM.Spectrum.compensated_one_family_rows
UEOTQM.Spectrum.zero_anomaly_unit_one_family_rows
UEOTQM.Spectrum.extra_screening_closes_positive_window
UEOTQM.pairingLinear
UEOTQM.heisenbergLinear
UEOTQM.FiniteObservablePresentation.closedCoordinateUpdate
UEOTQM.Adaptive.Protocol.leafFintype
UEOTQM.Spectrum.oneHiggs
UEOTQM.Spectrum.twentyHiggs
UEOTQM.sqrt_smul_gram
UEOTQM.sqrt_smul_krausTerm
UEOTQM.sqrt_smul_heisenbergTerm
UEOTQM.krausMap_sqrtWeightedKraus
UEOTQM.heisenberg_sqrtWeightedKraus
UEOTQM.weightedKrausMap_posSemidef
UEOTQM.weightedHeisenberg_posSemidef
UEOTQM.pairing_weightedKrausMap
UEOTQM.WeightedInstrument.toInstrument
UEOTQM.WeightedInstrument.toInstrument_kraus
UEOTQM.WeightedInstrument.toInstrument_branch
UEOTQM.WeightedInstrument.toInstrument_heisenberg
UEOTQM.WeightedInstrument.toInstrument_complete
UEOTQM.WeightedInstrument.branch_posSemidef
UEOTQM.WeightedInstrument.toInstrument_probability
UEOTQM.WeightedInstrument.sum_mass_branch
UEOTQM.WeightedInstrument.toInstrument_posterior_matrix
UEOTQM.zeroExtend_apply
UEOTQM.zeroExtend_of_not_mem_range
UEOTQM.sum_zeroExtend
UEOTQM.zeroExtend_nonneg
UEOTQM.weightedKrausMap_zeroExtend
UEOTQM.weightedHeisenberg_zeroExtend
UEOTQM.weightedGram_zeroExtend
UEOTQM.RaggedInstrument.toWeighted
UEOTQM.RaggedInstrument.toWeighted_weight_apply
UEOTQM.RaggedInstrument.toWeighted_kraus_apply
UEOTQM.RaggedInstrument.toWeighted_branch
UEOTQM.RaggedInstrument.toWeighted_toInstrument_branch
UEOTQM.RaggedInstrument.toWeighted_heisenberg
UEOTQM.RaggedInstrument.toWeighted_toInstrument_kraus_of_not_mem_range
UEOTQM.WeightedInstrument.padOutcomes
UEOTQM.WeightedInstrument.padOutcomes_branch_apply
UEOTQM.WeightedInstrument.padOutcomes_branch_of_not_mem_range
UEOTQM.WeightedInstrument.padOutcomes_heisenberg_apply
UEOTQM.WeightedInstrument.padOutcomes_heisenberg_of_not_mem_range
UEOTQM.WeightedInstrument.padOutcomes_probability_apply
UEOTQM.WeightedInstrument.padOutcomes_probability_of_not_mem_range
UEOTQM.Exact.MatrixData.ext
UEOTQM.Exact.MatrixData.toMatrix_re
UEOTQM.Exact.MatrixData.toMatrix_im
UEOTQM.Exact.MatrixData.toMatrix_injective
UEOTQM.Exact.MatrixData.toMatrix_eq_iff
\end{Verbatim}

\subsection{F.3.4 Requests 181--240}

\begin{Verbatim}[breaklines,breakanywhere,fontsize=\scriptsize]
UEOTQM.Exact.MatrixData.toMatrix_zero
UEOTQM.Exact.MatrixData.toMatrix_one
UEOTQM.Exact.MatrixData.toMatrix_add
UEOTQM.Exact.MatrixData.toMatrix_neg
UEOTQM.Exact.MatrixData.toMatrix_sub
UEOTQM.Exact.MatrixData.toMatrix_mul
UEOTQM.Exact.MatrixData.toMatrix_adjoint
UEOTQM.Exact.MatrixData.toMatrix_scale
UEOTQM.Exact.MatrixData.toMatrix_sum
UEOTQM.Exact.MatrixData.instDecidableEq
UEOTQM.Exact.toMatrix_weightedBranch
UEOTQM.Exact.toMatrix_weightedPullback
UEOTQM.Exact.toMatrix_weightedCompleteness
UEOTQM.Exact.toMatrix_weightedBranch_univ
UEOTQM.Exact.toMatrix_weightedPullback_univ
UEOTQM.Exact.Examples.imaginaryUnit_square_exact
UEOTQM.Exact.Examples.imaginaryUnit_adjoint_exact
UEOTQM.Exact.Examples.imaginaryUnit_square_complex
UEOTQM.Exact.Examples.hadamard_complete_exact
UEOTQM.Exact.Examples.hadamard_complete_complex
UEOTQM.ExactCertificate.check_eq_true_iff
UEOTQM.ExactCertificate.check_sound
UEOTQM.ExactCertificate.branch_eq_exact
UEOTQM.ExactCertificate.heisenberg_eq_exact
UEOTQM.ExactCertificate.exactRecordEffect_toMatrix
UEOTQM.ExactCertificate.basis_toMatrix
UEOTQM.ExactCertificate.coordinates_eq_iff_response_eq
UEOTQM.ExactCertificate.response_eq_recordCoordinates_unit
UEOTQM.ExactCertificate.adaptive_leafProbability_eq
UEOTQM.ExactCertificate.conditionsDecidable
UEOTQM.ExactCertificate.weightedInstrument
UEOTQM.ExactCertificate.instrumentFamily
UEOTQM.ExactCertificate.presentation
UEOTQM.pairing_sub_left
UEOTQM.pairing_add_right
UEOTQM.pairing_sum_right
UEOTQM.pairing_nonneg
UEOTQM.Effect.ext
UEOTQM.Effect.zero
UEOTQM.Effect.one
UEOTQM.Effect.complement
UEOTQM.Effect.zero_matrix
UEOTQM.Effect.one_matrix
UEOTQM.Effect.complement_matrix
UEOTQM.Effect.pairing_nonneg
UEOTQM.Effect.pairing_le_mass
UEOTQM.Effect.probability
UEOTQM.Effect.probability_nonneg
UEOTQM.Effect.probability_le_one
UEOTQM.Effect.probability_zero
UEOTQM.Effect.probability_one
UEOTQM.Effect.probability_complement
UEOTQM.Instrument.channel
UEOTQM.Instrument.channel_posSemidef
UEOTQM.Instrument.trace_channel
UEOTQM.Instrument.channelState
UEOTQM.Instrument.channelState_matrix
UEOTQM.Instrument.pullbackEffects
UEOTQM.Instrument.pullbackEffects_matrix
UEOTQM.Instrument.pairing_pullbackEffects
\end{Verbatim}

\subsection{F.3.5 Requests 241--300}

\begin{Verbatim}[breaklines,breakanywhere,fontsize=\scriptsize]
UEOTQM.Instrument.pullbackEffect
UEOTQM.Instrument.pairing_pullbackEffect
UEOTQM.Instrument.probability_channelState
UEOTQM.Instrument.branchEffect
UEOTQM.Instrument.branchEffect_matrix
UEOTQM.Instrument.pairing_branchEffect
UEOTQM.Instrument.probability_branchEffect_one
UEOTQM.Instrument.probability_posterior
UEOTQM.recordTest
UEOTQM.recordTest_matrix
UEOTQM.recordEffect_complement_posSemidef
UEOTQM.response_eq_effect_probability
UEOTQM.response_le_one
UEOTQM.effectGapSet
UEOTQM.Effect.probability_abs_sub_le_one
UEOTQM.effectGapSet_nonempty
UEOTQM.effectGapSet_bddAbove
UEOTQM.effectDistance
UEOTQM.Effect.probability_sub_le_effectDistance
UEOTQM.effectDistance_nonneg
UEOTQM.effectDistance_le_iff
UEOTQM.effectDistance_le_one
UEOTQM.effectDistance_self
UEOTQM.effectDistance_comm
UEOTQM.effectDistance_triangle
UEOTQM.effectDistance_eq_zero_iff
UEOTQM.response_sub_le_effectDistance
UEOTQM.channelDefect
UEOTQM.channelDefect_bddAbove
UEOTQM.channelDefect_nonneg
UEOTQM.channelDefect_le_one
UEOTQM.effectDistance_le_channelDefect
UEOTQM.channelDefect_self
UEOTQM.Instrument.channelState_effectDistance_le
UEOTQM.Instrument.probability_sub_le_effectDistance
UEOTQM.Instrument.probability_pos_of_effectDistance_lt
UEOTQM.Instrument.channelState_perturbation_bound
UEOTQM.normalized_probability_difference_le
UEOTQM.Instrument.posterior_effectDistance_le
UEOTQM.Instrument.posterior_effectDistance_le_of_probability_lower_bound
UEOTQM.channelTrajectory
UEOTQM.channelTrajectory_zero
UEOTQM.channelTrajectory_succ
UEOTQM.channelTrajectory_transport_certificate
UEOTQM.scalarBudget_one_eq_sum
UEOTQM.channelTrajectory_effectDistance_le_sum
UEOTQM.channelTrajectory_response_sub_le_sum
UEOTQM.Mixture.tv
UEOTQM.Mixture.tv_nonneg
UEOTQM.Mixture.tv_self
UEOTQM.Mixture.tv_comm
UEOTQM.Mixture.tv_le_one
UEOTQM.Mixture.expectation_sub_le_tv
UEOTQM.Mixture.density
UEOTQM.Mixture.density_matrix
UEOTQM.Mixture.probability_density
UEOTQM.Mixture.density_weight_change_le_tv
UEOTQM.Mixture.density_fixed_weights_le_sum
UEOTQM.Mixture.density_joint_change_le
UEOTQM.Mixture.weightedInstrument
\end{Verbatim}

\subsection{F.3.6 Requests 301--360}

\begin{Verbatim}[breaklines,breakanywhere,fontsize=\scriptsize]
UEOTQM.Mixture.instrument
UEOTQM.Mixture.instrument_kraus
UEOTQM.Mixture.instrument_branch
UEOTQM.Mixture.instrument_heisenberg
UEOTQM.Mixture.instrument_channel
UEOTQM.Mixture.instrument_channelState
UEOTQM.Mixture.channelDefect_weight_change_le_tv
UEOTQM.Mixture.channelDefect_fixed_weights_le_sum
UEOTQM.Mixture.channelDefect_joint_change_le
UEOTQM.Mixture.RationalWeights.toWeights
UEOTQM.Mixture.rationalTV
UEOTQM.Mixture.tv_toWeights
UEOTQM.Mixture.density_weight_change_le_rational_budget
UEOTQM.Mixture.channelDefect_weight_change_le_rational_budget
UEOTQM.Mixture.channelDefect_fixed_weights_le_rational_sum
UEOTQM.Mixture.rationalJointBudget
UEOTQM.Mixture.channelDefect_joint_change_le_rational_budget
UEOTQM.Cesaro.iterateState
UEOTQM.Cesaro.iterateState_zero
UEOTQM.Cesaro.iterateState_succ
UEOTQM.Cesaro.channelLinear
UEOTQM.Cesaro.channelLinear_apply
UEOTQM.Cesaro.orbitSum
UEOTQM.Cesaro.orbitSum_positive
UEOTQM.Cesaro.orbitSum_trace
UEOTQM.Cesaro.averageState
UEOTQM.Cesaro.averageState_matrix
UEOTQM.Cesaro.orbitSum_telescope
UEOTQM.Cesaro.matrix_residual
UEOTQM.Cesaro.pairing_sub_right
UEOTQM.Cesaro.probability_residual
UEOTQM.Cesaro.probability_residual_le
UEOTQM.Cesaro.effectDistance_residual_le
UEOTQM.Cesaro.density_operatorNorm_le_one
UEOTQM.Cesaro.densityMatrices
UEOTQM.Cesaro.densityMatrices_closed
UEOTQM.Cesaro.densityMatrices_compact
UEOTQM.Cesaro.continuous_channel
UEOTQM.Cesaro.matrix_residual_norm_le
UEOTQM.Cesaro.averageSequence
UEOTQM.Cesaro.matrix_residual_tendsto_zero
UEOTQM.Cesaro.subsequence_limit_invariant
UEOTQM.Cesaro.subsequence_limit_density
UEOTQM.Cesaro.exists_stationary_subsequence
UEOTQM.Cesaro.exists_stationary_of_density
UEOTQM.Cesaro.normalizedIdentity
UEOTQM.Cesaro.exists_stationary
UEOTQM.Survival.diagonalOperator
UEOTQM.Survival.diagonalOperator_posSemidef
UEOTQM.Survival.mass_diagonalOperator
UEOTQM.Survival.diagonalOperator_injective
UEOTQM.Survival.diagonalOperator_smul
UEOTQM.Survival.diagonalDensity
UEOTQM.Survival.diagonalDensity_matrix
UEOTQM.Survival.matrixUnit
UEOTQM.Survival.matrixUnit_gram
UEOTQM.Survival.matrixUnit_sandwich
UEOTQM.Survival.measurePrepare
UEOTQM.Survival.measurePrepare_apply
UEOTQM.Survival.measurePrepare_diagonal
\end{Verbatim}

\subsection{F.3.7 Requests 361--420}

\begin{Verbatim}[breaklines,breakanywhere,fontsize=\scriptsize]
UEOTQM.Survival.measurePrepare_gram
UEOTQM.Survival.failureKernel
UEOTQM.Survival.failureKernel_nonneg
UEOTQM.Survival.outcomeKernel
UEOTQM.Survival.outcomeKernel_complete
UEOTQM.Survival.weightedInstrument
UEOTQM.Survival.instrument
UEOTQM.Survival.instrument_branch
UEOTQM.Survival.success_branch
UEOTQM.Survival.rowApply_nonneg
UEOTQM.Survival.success_probability
UEOTQM.Survival.success_posterior_matrix
UEOTQM.Survival.killedIterate
UEOTQM.Survival.killedIterate_posSemidef
UEOTQM.Survival.killedIterate_diagonal
UEOTQM.Survival.survivalMass
UEOTQM.Survival.survivalMass_eq_response
UEOTQM.Survival.survivalMass_diagonal
UEOTQM.Survival.conditionalState
UEOTQM.Survival.qsd_killedIterate
UEOTQM.Survival.qsd_survival_pow
UEOTQM.Survival.qsd_survival_pos
UEOTQM.Survival.qsd_conditionalState
UEOTQM.Survival.qsd_rate_le_one
UEOTQM.Survival.doobTransform_nonneg
UEOTQM.Survival.doobInstrument
UEOTQM.Survival.doobDensity
UEOTQM.Survival.doob_success_branch
UEOTQM.Survival.doob_success_probability
UEOTQM.Survival.doob_killedIterate
UEOTQM.Survival.doob_survivalMass
UEOTQM.Survival.doob_conditionalState
UEOTQM.Classical.State
UEOTQM.Classical.Effect
UEOTQM.Classical.Effect.mk
UEOTQM.Classical.state_ext
UEOTQM.Classical.evaluate
UEOTQM.Classical.evaluate_nonneg
UEOTQM.Classical.evaluate_le_one
UEOTQM.Classical.indicator
UEOTQM.Classical.evaluate_indicator
UEOTQM.Classical.all_effects_separate
UEOTQM.Classical.dirac
UEOTQM.Classical.dirac_weight
UEOTQM.Classical.evaluate_dirac
UEOTQM.Classical.mix
UEOTQM.Classical.mix_weight
UEOTQM.Classical.halfMix
UEOTQM.Classical.halfMix_weight
UEOTQM.Classical.pushforward
UEOTQM.Classical.pushforward_weight
UEOTQM.Classical.pushforward_dirac
UEOTQM.Classical.pushforward_halfMix
UEOTQM.Classical.marginalLeft
UEOTQM.Classical.marginalLeft_weight
UEOTQM.Classical.marginalLeft_eq_pushforward
UEOTQM.Classical.marginalLeft_dirac
UEOTQM.Classical.marginalLeft_halfMix
UEOTQM.Classical.IsPure
UEOTQM.Classical.dirac_isPure
\end{Verbatim}

\subsection{F.3.8 Requests 421--480}

\begin{Verbatim}[breaklines,breakanywhere,fontsize=\scriptsize]
UEOTQM.Classical.dirac_ne
UEOTQM.Classical.halfMix_not_pure
UEOTQM.Classical.halfRow
UEOTQM.Classical.halfRow_weight
UEOTQM.Classical.marginalLeft_halfRow
UEOTQM.Classical.no_pure_extension_halfMix
UEOTQM.Classical.attach
UEOTQM.Classical.marginalLeft_attach
UEOTQM.Classical.dirac_has_pure_extension
UEOTQM.Classical.toPMF
UEOTQM.Classical.toPMF_apply
UEOTQM.Classical.toPMF_apply_toReal
UEOTQM.Classical.toPMF_dirac
UEOTQM.Classical.toPMF_pushforward
UEOTQM.Classical.rowState
UEOTQM.Classical.rowState_weight
UEOTQM.Classical.stepState
UEOTQM.Classical.stepState_dirac
UEOTQM.Classical.toPMF_stepState
UEOTQM.CorePurificationBoundary.Micro
UEOTQM.CorePurificationBoundary.fullyObservedSurvivalModel
UEOTQM.CorePurificationBoundary.terminal_eq_iff
UEOTQM.CorePurificationBoundary.Terminal
UEOTQM.CorePurificationBoundary.terminalFintype
UEOTQM.CorePurificationBoundary.terminalMeasurableSpace
UEOTQM.CorePurificationBoundary.terminalControl
UEOTQM.CorePurificationBoundary.terminalGood
UEOTQM.CorePurificationBoundary.terminalControl_source
UEOTQM.CorePurificationBoundary.terminalControl_reward
UEOTQM.CorePurificationBoundary.terminalControl_optimalValue
UEOTQM.CorePurificationBoundary.terminalControl_injective
UEOTQM.CorePurificationBoundary.terminalGood_source
UEOTQM.CorePurificationBoundary.viable_all_horizons
UEOTQM.CorePurificationBoundary.initial_support_viable
UEOTQM.CorePurificationBoundary.firstRow
UEOTQM.CorePurificationBoundary.firstRow_halfMix
UEOTQM.CorePurificationBoundary.firstQuotientState
UEOTQM.CorePurificationBoundary.quotient_support_distinct
UEOTQM.CorePurificationBoundary.firstQuotientState_halfMix
UEOTQM.CorePurificationBoundary.firstQuotientState_not_pure
UEOTQM.CorePurificationBoundary.firstQuotientState_no_pure_extension
UEOTQM.CorePurificationBoundary.sourcePMF
UEOTQM.CorePurificationBoundary.piStar
UEOTQM.CorePurificationBoundary.actualPathLaw
UEOTQM.CorePurificationBoundary.sourcePMF_toReal
UEOTQM.CorePurificationBoundary.sourcePMF_firstRow
UEOTQM.CorePurificationBoundary.sourcePMF_rowState
UEOTQM.CorePurificationBoundary.piStar_preserves
UEOTQM.CorePurificationBoundary.actualPath_all_times_safe
UEOTQM.CorePurificationBoundary.actualPath_first_quotient_marginal
UEOTQM.CorePurificationBoundary.false_action_destructive
UEOTQM.CorePurificationBoundary.FiniteCoreSlice
UEOTQM.CorePurificationBoundary.finiteCoreSlice_holds
UEOTQM.CorePurificationBoundary.classical_purification_separation
UEOTQM.CorePurificationBoundary.RestrictedTerminal
UEOTQM.CorePurificationBoundary.restrictedTerminalFintype
UEOTQM.CorePurificationBoundary.restrictedFirstQuotientState
UEOTQM.CorePurificationBoundary.restricted_support_merge
UEOTQM.CorePurificationBoundary.restrictedFirstQuotientState_dirac
UEOTQM.CorePurificationBoundary.restrictedFirstQuotientState_pure
\end{Verbatim}

\subsection{F.3.9 Requests 481--540}

\begin{Verbatim}[breaklines,breakanywhere,fontsize=\scriptsize]
UEOTQM.RealLocalTomography.IsReal
UEOTQM.RealLocalTomography.symmetric_of_real_hermitian
UEOTQM.RealLocalTomography.productEffect
UEOTQM.RealLocalTomography.productEffect_matrix
UEOTQM.RealLocalTomography.plusProjector
UEOTQM.RealLocalTomography.minusProjector
UEOTQM.RealLocalTomography.plusProjector_hermitian
UEOTQM.RealLocalTomography.minusProjector_hermitian
UEOTQM.RealLocalTomography.plusProjector_square
UEOTQM.RealLocalTomography.minusProjector_square
UEOTQM.RealLocalTomography.plusProjector_mul_minusProjector
UEOTQM.RealLocalTomography.one_sub_plusProjector
UEOTQM.RealLocalTomography.plusProjector_posSemidef
UEOTQM.RealLocalTomography.minusProjector_posSemidef
UEOTQM.RealLocalTomography.involutionEffect
UEOTQM.RealLocalTomography.Qubit
UEOTQM.RealLocalTomography.TwoQubit
UEOTQM.RealLocalTomography.pauliY
UEOTQM.RealLocalTomography.parity
UEOTQM.RealLocalTomography.pauliY_hermitian
UEOTQM.RealLocalTomography.pauliY_square
UEOTQM.RealLocalTomography.parity_hermitian
UEOTQM.RealLocalTomography.parity_square
UEOTQM.RealLocalTomography.pauliY_trace
UEOTQM.RealLocalTomography.parity_trace
UEOTQM.RealLocalTomography.parity_isReal
UEOTQM.RealLocalTomography.rhoPlus
UEOTQM.RealLocalTomography.rhoMinus
UEOTQM.RealLocalTomography.rhoPlus_matrix
UEOTQM.RealLocalTomography.rhoMinus_matrix
UEOTQM.RealLocalTomography.rhoPlus_isReal
UEOTQM.RealLocalTomography.rhoMinus_isReal
UEOTQM.RealLocalTomography.rhoPlus_symmetric
UEOTQM.RealLocalTomography.rhoMinus_symmetric
UEOTQM.RealLocalTomography.density_difference
UEOTQM.RealLocalTomography.trace_mul_pauliY_zero
UEOTQM.RealLocalTomography.trace_product_density_difference_zero
UEOTQM.RealLocalTomography.all_real_product_probabilities_equal
UEOTQM.RealLocalTomography.jointTest
UEOTQM.RealLocalTomography.jointTest_matrix
UEOTQM.RealLocalTomography.jointTest_isReal
UEOTQM.RealLocalTomography.jointTest_probability_plus
UEOTQM.RealLocalTomography.jointTest_probability_minus
UEOTQM.RealLocalTomography.rhoPlus_ne_rhoMinus
UEOTQM.RealLocalTomography.real_local_tomography_fails
UEOTQM.RealLocalTomography.complexLocalTest
UEOTQM.RealLocalTomography.complexLocalTest_not_real
UEOTQM.RealLocalTomography.complex_product_probability_plus
UEOTQM.RealLocalTomography.complex_product_probability_minus
UEOTQM.RealLocalTomography.complex_product_distinguishes
UEOTQM.QuantumDoob.successOperator
UEOTQM.QuantumDoob.successOperator_zero
UEOTQM.QuantumDoob.successOperator_succ
UEOTQM.QuantumDoob.successOperator_smul
UEOTQM.QuantumDoob.successOperator_posSemidef
UEOTQM.QuantumDoob.eigen_successOperator
UEOTQM.QuantumDoob.successMass
UEOTQM.QuantumDoob.successMass_eq_response
UEOTQM.QuantumDoob.eigen_successMass
UEOTQM.QuantumDoob.eigen_successMass_pos
\end{Verbatim}

\subsection{F.3.10 Requests 541--600}

\begin{Verbatim}[breaklines,breakanywhere,fontsize=\scriptsize]
UEOTQM.QuantumDoob.conditionalState
UEOTQM.QuantumDoob.eigen_conditionalState
UEOTQM.QuantumDoob.eigen_rate_nonneg
UEOTQM.QuantumDoob.eigen_rate_le_one
UEOTQM.QuantumDoob.InvertibleFilter
UEOTQM.QuantumDoob.InvertibleFilter.mk
UEOTQM.QuantumDoob.InvertibleFilter.apply
UEOTQM.QuantumDoob.InvertibleFilter.unapply
UEOTQM.QuantumDoob.InvertibleFilter.metric
UEOTQM.QuantumDoob.InvertibleFilter.apply_smul
UEOTQM.QuantumDoob.InvertibleFilter.unapply_apply
UEOTQM.QuantumDoob.InvertibleFilter.apply_unapply
UEOTQM.QuantumDoob.InvertibleFilter.apply_posSemidef
UEOTQM.QuantumDoob.InvertibleFilter.mass_apply_pos
UEOTQM.QuantumDoob.InvertibleFilter.state
UEOTQM.QuantumDoob.InvertibleFilter.state_matrix
UEOTQM.QuantumDoob.InvertibleFilter.mass_apply_eq_pairing
UEOTQM.QuantumDoob.InvertibleFilter.metric_conjugate
UEOTQM.QuantumDoob.transformed_gram
UEOTQM.QuantumDoob.transformed_map
UEOTQM.QuantumDoob.doob_complete
UEOTQM.QuantumDoob.doobWeighted
UEOTQM.QuantumDoob.doobInstrument
UEOTQM.QuantumDoob.doob_branch
UEOTQM.QuantumDoob.doob_branch_apply
UEOTQM.QuantumDoob.doob_branch_state
UEOTQM.QuantumDoob.doob_channelState
UEOTQM.QuantumDoob.doob_successOperator_apply
UEOTQM.QuantumDoob.eigenmetric_successMass
UEOTQM.QuantumDoob.eigenmetric_successMass_pos
UEOTQM.QuantumDoob.doob_finite_conditioning
UEOTQM.QuantumDoob.doob_successOperator_state
UEOTQM.QuantumDoob.doob_successMass_state
UEOTQM.RecordProbability.eventProbability
UEOTQM.RecordProbability.positiveEvent
UEOTQM.RecordProbability.eventProbability_empty
UEOTQM.RecordProbability.eventProbability_univ
UEOTQM.RecordProbability.eventProbability_nonneg
UEOTQM.RecordProbability.eventProbability_le_one
UEOTQM.RecordProbability.eventProbability_eq_indicator
UEOTQM.RecordProbability.eventProbability_abs_sub_le_tv
UEOTQM.RecordProbability.positiveEvent_sub_eq_tv
UEOTQM.RecordProbability.positiveEvent_abs_sub_eq_tv
UEOTQM.RecordProbability.tv_le_iff_events
UEOTQM.RectangularChannel
UEOTQM.RectangularChannel.operator
UEOTQM.RectangularChannel.operator_posSemidef
UEOTQM.RectangularChannel.operator_zero
UEOTQM.RectangularChannel.operator_add
UEOTQM.RectangularChannel.operator_smul
UEOTQM.RectangularChannel.operator_real_smul
UEOTQM.RectangularChannel.operatorLinear
UEOTQM.RectangularChannel.operator_sum
UEOTQM.RectangularChannel.trace_operator
UEOTQM.RectangularChannel.mass_operator
UEOTQM.RectangularChannel.state
UEOTQM.RectangularChannel.state_matrix
UEOTQM.RectangularChannel.adjoint
UEOTQM.RectangularChannel.adjoint_posSemidef
UEOTQM.RectangularChannel.adjoint_one
\end{Verbatim}

\subsection{F.3.11 Requests 601--660}

\begin{Verbatim}[breaklines,breakanywhere,fontsize=\scriptsize]
UEOTQM.RectangularChannel.adjoint_sub
UEOTQM.RectangularChannel.trace_adjoint
UEOTQM.RectangularChannel.pullbackEffect
UEOTQM.RectangularChannel.pullbackEffect_matrix
UEOTQM.RectangularChannel.probability_state
UEOTQM.RectangularChannel.state_effectDistance_le
UEOTQM.RectangularChannel.identity
UEOTQM.RectangularChannel.identity_operator
UEOTQM.RectangularChannel.identity_state
UEOTQM.RectangularChannel.comp
UEOTQM.RectangularChannel.comp_operator
UEOTQM.RectangularChannel.comp_state
UEOTQM.RectangularChannel.amplify
UEOTQM.RectangularChannel.amplify_posSemidef
UEOTQM.RectangularChannel.amplify_operator_apply
UEOTQM.RectangularChannel.discardRightKraus
UEOTQM.RectangularChannel.discardRight
UEOTQM.RectangularChannel.discardRight_operator_apply
UEOTQM.RectangularChannel.discardLeftKraus
UEOTQM.RectangularChannel.discardLeft
UEOTQM.RectangularChannel.discardLeft_operator_apply
UEOTQM.RectangularChannel.defect
UEOTQM.RectangularChannel.defect_bddAbove
UEOTQM.RectangularChannel.defect_nonneg
UEOTQM.RectangularChannel.defect_le_one
UEOTQM.RectangularChannel.state_distance_le_defect
UEOTQM.RectangularChannel.perturbation_bound
UEOTQM.QuantumPurification.mix
UEOTQM.QuantumPurification.mix_matrix
UEOTQM.QuantumPurification.IsPure
UEOTQM.QuantumPurification.outer
UEOTQM.QuantumPurification.outer_apply
UEOTQM.QuantumPurification.outer_positive
UEOTQM.QuantumPurification.outer_hermitian
UEOTQM.QuantumPurification.outer_trace
UEOTQM.QuantumPurification.outer_square
UEOTQM.QuantumPurification.outer_sandwich
UEOTQM.QuantumPurification.vectorState
UEOTQM.QuantumPurification.vectorState_matrix
UEOTQM.QuantumPurification.exists_gram
UEOTQM.QuantumPurification.positive_compression_zero
UEOTQM.QuantumPurification.density_eq_vector_of_compression_zero
UEOTQM.QuantumPurification.vectorState_isPure
UEOTQM.QuantumPurification.vectorize
UEOTQM.QuantumPurification.vectorize_outer_trace
UEOTQM.QuantumPurification.discard_vectorize_outer
UEOTQM.QuantumPurification.gram
UEOTQM.QuantumPurification.gram_spec
UEOTQM.QuantumPurification.purificationVector
UEOTQM.QuantumPurification.purificationVector_normalized
UEOTQM.QuantumPurification.purification
UEOTQM.QuantumPurification.purification_matrix
UEOTQM.QuantumPurification.purification_isPure
UEOTQM.QuantumPurification.purification_marginal
UEOTQM.QuantumPurification.exists_purification
UEOTQM.QuantumPurification.density_empty_elim
UEOTQM.QuantumPurification.product
UEOTQM.QuantumPurification.product_matrix
UEOTQM.QuantumPurification.discard_product
UEOTQM.QuantumPurification.probability_mix
\end{Verbatim}

\subsection{F.3.12 Requests 661--720}

\begin{Verbatim}[breaklines,breakanywhere,fontsize=\scriptsize]
UEOTQM.QuantumPurification.state_mix
UEOTQM.OperationalSystems.FiniteSystem
UEOTQM.OperationalSystems.FiniteSystem.of
UEOTQM.OperationalSystems.FiniteSystem.product
UEOTQM.OperationalSystems.Presentation
UEOTQM.OperationalSystems.Presentation.discardRight
UEOTQM.OperationalSystems.Laws
UEOTQM.OperationalSystems.IsPure
UEOTQM.OperationalSystems.Purifies
UEOTQM.OperationalSystems.PurificationExistence
UEOTQM.OperationalSystems.LocalTomography
UEOTQM.OperationalSystems.ClassicalModel.product
UEOTQM.OperationalSystems.ClassicalModel.productEffect
UEOTQM.OperationalSystems.ClassicalModel.act
UEOTQM.OperationalSystems.ClassicalModel.probability_mix
UEOTQM.OperationalSystems.ClassicalModel.product_probability
UEOTQM.OperationalSystems.ClassicalModel.marginal_product
UEOTQM.OperationalSystems.ClassicalModel.act_discard
UEOTQM.OperationalSystems.ClassicalModel.act_mix
UEOTQM.OperationalSystems.ClassicalModel.act_identity
UEOTQM.OperationalSystems.ClassicalModel.act_comp
UEOTQM.OperationalSystems.classicalModel
UEOTQM.OperationalSystems.classical_laws
UEOTQM.OperationalSystems.classical_pure_iff
UEOTQM.OperationalSystems.classical_discard_eq
UEOTQM.OperationalSystems.classical_not_purification
UEOTQM.OperationalSystems.coreSystem
UEOTQM.OperationalSystems.core_first_marginal_no_purifies
UEOTQM.OperationalSystems.QuantumProcess
UEOTQM.OperationalSystems.QuantumModel.trace_product_im_zero
UEOTQM.OperationalSystems.QuantumModel.product_probability
UEOTQM.OperationalSystems.complexModel
UEOTQM.OperationalSystems.complex_laws
UEOTQM.OperationalSystems.complex_pure_iff
UEOTQM.OperationalSystems.complex_purification
UEOTQM.OperationalSystems.RealModel.IsRealMatrix
UEOTQM.OperationalSystems.RealModel.real_add
UEOTQM.OperationalSystems.RealModel.real_smul
UEOTQM.OperationalSystems.RealModel.real_mul
UEOTQM.OperationalSystems.RealModel.real_conjTranspose
UEOTQM.OperationalSystems.RealModel.real_kronecker
UEOTQM.OperationalSystems.RealModel.State
UEOTQM.OperationalSystems.RealModel.Effect
UEOTQM.OperationalSystems.RealModel.mix
UEOTQM.OperationalSystems.RealModel.product
UEOTQM.OperationalSystems.RealModel.productEffect
UEOTQM.OperationalSystems.RealModel.Process
UEOTQM.OperationalSystems.RealModel.operator_real
UEOTQM.OperationalSystems.RealModel.act
UEOTQM.OperationalSystems.RealModel.identity
UEOTQM.OperationalSystems.RealModel.compose
UEOTQM.OperationalSystems.RealModel.discardProcess
UEOTQM.OperationalSystems.realModel
UEOTQM.OperationalSystems.real_laws
UEOTQM.OperationalSystems.real_not_localTomography
UEOTQM.ClassicalQuantumMemory.taggedKraus
UEOTQM.ClassicalQuantumMemory.taggedKraus_gram
UEOTQM.ClassicalQuantumMemory.taggedKraus_conjugation
UEOTQM.ClassicalQuantumMemory.block
UEOTQM.ClassicalQuantumMemory.block_apply_same
\end{Verbatim}

\subsection{F.3.13 Requests 721--780}

\begin{Verbatim}[breaklines,breakanywhere,fontsize=\scriptsize]
UEOTQM.ClassicalQuantumMemory.block_apply_ne
UEOTQM.ClassicalQuantumMemory.trace_block
UEOTQM.ClassicalQuantumMemory.channel
UEOTQM.ClassicalQuantumMemory.channel_operator_apply
UEOTQM.ClassicalQuantumMemory.channel_amplify_operator_apply
UEOTQM.ClassicalQuantumMemory.channel_operator_eq_block
UEOTQM.ClassicalQuantumMemory.state
UEOTQM.ClassicalQuantumMemory.state_matrix
UEOTQM.ClassicalQuantumMemory.discard_record_operator
UEOTQM.ClassicalQuantumMemory.discard_record_state
UEOTQM.ClassicalQuantumMemory.discard_memory_operator
UEOTQM.ClassicalQuantumMemory.discard_memory_state_matrix
UEOTQM.ClassicalQuantumMemory.state_block_mass
UEOTQM.ClassicalQuantumMemory.positive_block_posterior
UEOTQM.ClassicalQuantumMemory.zero_probability_iff_zero_block
UEOTQM.ClassicalQuantumMemory.eventEffect
UEOTQM.ClassicalQuantumMemory.eventEffect_matrix
UEOTQM.ClassicalQuantumMemory.eventEffect_pairing_block
UEOTQM.ClassicalQuantumMemory.weights
UEOTQM.ClassicalQuantumMemory.weights_weight
UEOTQM.ClassicalQuantumMemory.event_probability
UEOTQM.ClassicalQuantumMemory.record_tv_le_joint_distance
UEOTQM.ClassicalQuantumMemory.defect
UEOTQM.ClassicalQuantumMemory.perturbation_bound
UEOTQM.ClassicalQuantumMemory.record_tv_perturbation_bound
UEOTQM.ClassicalQuantumMemory.ProtocolMemory.Leaf
UEOTQM.ClassicalQuantumMemory.ProtocolMemory.KrausPath
UEOTQM.ClassicalQuantumMemory.ProtocolMemory.pathKraus
UEOTQM.ClassicalQuantumMemory.ProtocolMemory.path_complete
UEOTQM.ClassicalQuantumMemory.ProtocolMemory.path_branch_eq_recordOperator
UEOTQM.ClassicalQuantumMemory.ProtocolMemory.channel
UEOTQM.ClassicalQuantumMemory.ProtocolMemory.channel_operator_apply
UEOTQM.ClassicalQuantumMemory.ProtocolMemory.channel_amplify_operator_apply
UEOTQM.ClassicalQuantumMemory.ProtocolMemory.channel_operator_eq_block
UEOTQM.ClassicalQuantumMemory.ProtocolMemory.state
UEOTQM.ClassicalQuantumMemory.ProtocolMemory.state_matrix
UEOTQM.ClassicalQuantumMemory.ProtocolMemory.discard_record_operator
UEOTQM.ClassicalQuantumMemory.ProtocolMemory.discard_memory_operator
UEOTQM.ClassicalQuantumMemory.ProtocolMemory.discard_memory_state_matrix
UEOTQM.ClassicalQuantumMemory.ProtocolMemory.state_block_mass
UEOTQM.ClassicalQuantumMemory.ProtocolMemory.zero_probability_iff_zero_block
UEOTQM.ClassicalQuantumMemory.ProtocolMemory.weights
UEOTQM.ClassicalQuantumMemory.ProtocolMemory.weights_weight
UEOTQM.ClassicalQuantumMemory.ProtocolMemory.event_probability
UEOTQM.ClassicalQuantumMemory.ProtocolMemory.record_tv_le_joint_distance
UEOTQM.ClassicalQuantumMemory.ProtocolMemory.defect
UEOTQM.ClassicalQuantumMemory.ProtocolMemory.perturbation_bound
UEOTQM.ClassicalQuantumMemory.ProtocolMemory.record_tv_perturbation_bound
UEOTQM.QuantumDoobModel.successK
UEOTQM.QuantumDoobModel.failureK
UEOTQM.QuantumDoobModel.success_gram
UEOTQM.QuantumDoobModel.failure_gram
UEOTQM.QuantumDoobModel.sourceK
UEOTQM.QuantumDoobModel.source_complete
UEOTQM.QuantumDoobModel.sourceInstrument
UEOTQM.QuantumDoobModel.source_branch_formula
UEOTQM.QuantumDoobModel.rhoMatrix
UEOTQM.QuantumDoobModel.rhoMatrix_positive
UEOTQM.QuantumDoobModel.rhoMatrix_trace
UEOTQM.QuantumDoobModel.rho
\end{Verbatim}

\subsection{F.3.14 Requests 781--840}

\begin{Verbatim}[breaklines,breakanywhere,fontsize=\scriptsize]
UEOTQM.QuantumDoobModel.S
UEOTQM.QuantumDoobModel.T
UEOTQM.QuantumDoobModel.inverse_mul
UEOTQM.QuantumDoobModel.mul_inverse
UEOTQM.QuantumDoobModel.filter
UEOTQM.QuantumDoobModel.metric_exact
UEOTQM.QuantumDoobModel.source_right_eigen
UEOTQM.QuantumDoobModel.source_left_eigen
UEOTQM.QuantumDoobModel.filtered_mass_exact
UEOTQM.QuantumDoobModel.plusMatrix
UEOTQM.QuantumDoobModel.filtered_state_exact
UEOTQM.QuantumDoobModel.plus
UEOTQM.QuantumDoobModel.plus_matrix
UEOTQM.QuantumDoobModel.transformedInstrument
UEOTQM.QuantumDoobModel.transformed_branch_formula
UEOTQM.QuantumDoobModel.success_kraus_noncommuting
UEOTQM.QuantumDoobModel.input_coherence_retained
UEOTQM.QuantumDoobDynamics.decay
UEOTQM.QuantumDoobDynamics.decay_zero
UEOTQM.QuantumDoobDynamics.decay_pos
UEOTQM.QuantumDoobDynamics.decay_le_one
UEOTQM.QuantumDoobDynamics.decay_succ
UEOTQM.QuantumDoobDynamics.decay_tendsto
UEOTQM.QuantumDoobDynamics.effectDistance_eq_of_matrix_sub_smul
UEOTQM.QuantumDoobDynamics.doob_branch_formula
UEOTQM.QuantumDoobDynamics.doob_branch_trace
UEOTQM.QuantumDoobDynamics.doob_successOperator_formula
UEOTQM.QuantumDoobDynamics.doob_iterateState_matrix
UEOTQM.QuantumDoobDynamics.doob_stationary_plus
UEOTQM.QuantumDoobDynamics.doob_stationary_unique
UEOTQM.QuantumDoobDynamics.doob_effectDistance_exact
UEOTQM.QuantumDoobDynamics.doob_matrix_tendsto
UEOTQM.QuantumDoobDynamics.filterMass
UEOTQM.QuantumDoobDynamics.pullbackWeight
UEOTQM.QuantumDoobDynamics.normalizer
UEOTQM.QuantumDoobDynamics.filterMass_diag
UEOTQM.QuantumDoobDynamics.filterMass_trace
UEOTQM.QuantumDoobDynamics.filterMass_metric_trace
UEOTQM.QuantumDoobDynamics.filterMass_bounds
UEOTQM.QuantumDoobDynamics.pullbackWeight_bounds
UEOTQM.QuantumDoobDynamics.normalizer_bounds
UEOTQM.QuantumDoobDynamics.normalizer_pos
UEOTQM.QuantumDoobDynamics.unapply_plus
UEOTQM.QuantumDoobDynamics.source_successOperator_formula
UEOTQM.QuantumDoobDynamics.source_successOperator_density
UEOTQM.QuantumDoobDynamics.source_successMass_formula
UEOTQM.QuantumDoobDynamics.source_response_formula
UEOTQM.QuantumDoobDynamics.source_successMass_pos
UEOTQM.QuantumDoobDynamics.source_successMass_bounds
UEOTQM.QuantumDoobDynamics.source_conditionalState_matrix
UEOTQM.QuantumDoobDynamics.conditioned
UEOTQM.QuantumDoobDynamics.conditioned_matrix
UEOTQM.QuantumDoobDynamics.conditioned_matrix_affine
UEOTQM.QuantumDoobDynamics.conditioned_matrix_sub
UEOTQM.QuantumDoobDynamics.conditioned_effectDistance_exact
UEOTQM.QuantumDoobDynamics.conditioned_effectDistance_le
UEOTQM.QuantumDoobDynamics.normalizer_tendsto
UEOTQM.QuantumDoobDynamics.conditionalWeight_tendsto
UEOTQM.QuantumDoobDynamics.conditioned_matrix_tendsto
UEOTQM.QuantumDoobDynamics.conditioned_effectDistance_tendsto
\end{Verbatim}

\subsection{F.3.15 Requests 841--900}

\begin{Verbatim}[breaklines,breakanywhere,fontsize=\scriptsize]
UEOTQM.QuantumDoobDynamics.source_successMass_tendsto_zero
UEOTQM.QuantumDoobDynamics.filterMass_rho
UEOTQM.QuantumDoobDynamics.pullbackWeight_rho
UEOTQM.QuantumDoobDynamics.normalizer_rho
UEOTQM.QuantumDoobDynamics.source_successMass_rho
UEOTQM.OperationalSystems.finiteSystemCoeSort
UEOTQM.ClassicalQuantumMemory.ProtocolMemory.krausPathFintype
UEOTQM.OperationalSystems.FiniteSystem.finite
UEOTQM.OperationalSystems.FiniteSystem.decidableEq
UEOTQM.RectangularChannel.mk
UEOTQM.OperationalSystems.FiniteSystem.mk
UEOTQM.OperationalSystems.Presentation.mk
UEOTQM.OperationalSystems.Laws.mk
UEOTQM.OperationalSystems.QuantumProcess.mk
UEOTQM.OperationalSystems.RealModel.Process.mk
UEOTQM.QuantumSeparation.boundedHermitianEffect
UEOTQM.QuantumSeparation.boundedHermitianEffect_matrix
UEOTQM.QuantumSeparation.boundedHermitianEffect_affine
UEOTQM.QuantumSeparation.exists_effect_affine_of_hermitian
UEOTQM.QuantumSeparation.hermitian_mem_real_span_effects
UEOTQM.QuantumSeparation.pairing_eq_of_all_effect_pairings_eq
UEOTQM.QuantumSeparation.pairing_sub_right
UEOTQM.QuantumSeparation.density_eq_of_all_effect_probabilities_eq
UEOTQM.QuantumSeparation.density_eq_iff_all_effect_probabilities_eq
UEOTQM.QuantumSeparation.effectDistance_eq_zero_iff_density_eq
UEOTQM.ComplexLocalTomography.symmetricUnit
UEOTQM.ComplexLocalTomography.imaginaryUnit
UEOTQM.ComplexLocalTomography.symmetricUnit_hermitian
UEOTQM.ComplexLocalTomography.imaginaryUnit_hermitian
UEOTQM.ComplexLocalTomography.sub_kronecker
UEOTQM.ComplexLocalTomography.kronecker_sub
UEOTQM.ComplexLocalTomography.pairing_symmetricUnit
UEOTQM.ComplexLocalTomography.pairing_imaginaryUnit
UEOTQM.ComplexLocalTomography.symmetricUnit_product
UEOTQM.ComplexLocalTomography.imaginaryUnit_product
UEOTQM.ComplexLocalTomography.hermitian_product_pairing_eq
UEOTQM.ComplexLocalTomography.density_eq_of_all_product_effect_probabilities_eq
UEOTQM.ComplexLocalTomography.complex_localTomography
UEOTQM.QuantumPurification.outer_real_smul
UEOTQM.QuantumPurification.gram_eq_sum_outer
UEOTQM.QuantumPurification.IsPure.normalize_eq_of_positive_summand
UEOTQM.QuantumPurification.IsPure.exists_vector
UEOTQM.QuantumPurification.isPure_iff_exists_vector
UEOTQM.QuantumPurification.inner_matrix_image
UEOTQM.QuantumPurification.norm_matrix_image_eq_of_gram
UEOTQM.QuantumPurification.exists_unitary_of_column_gram
UEOTQM.QuantumPurification.exists_auxiliary_unitary_of_gram
UEOTQM.QuantumPurification.auxiliary_action_apply
UEOTQM.QuantumPurification.auxiliary_action_vectorize
UEOTQM.QuantumPurification.outer_mulVec
UEOTQM.QuantumPurification.exists_auxiliary_unitary_of_equal_marginal
UEOTQM.QuantumPurification.pure_states_unitarily_equivalent
UEOTQM.QuantumPurification.auxiliary_tensor_isometry
UEOTQM.QuantumPurification.auxiliaryEmbeddingChannel
UEOTQM.QuantumPurification.auxiliaryEmbeddingChannel_operator
UEOTQM.QuantumPurification.auxiliaryEmbeddingChannel_operator_apply
UEOTQM.QuantumPurification.auxiliaryEmbeddingChannel_discard_operator
UEOTQM.QuantumPurification.auxiliaryEmbeddingChannel_discard_state
UEOTQM.QuantumPurification.auxiliaryEmbeddingChannel_preserves_purity
UEOTQM.QuantumPurification.auxiliaryInl
\end{Verbatim}

\subsection{F.3.16 Requests 901--960}

\begin{Verbatim}[breaklines,breakanywhere,fontsize=\scriptsize]
UEOTQM.QuantumPurification.auxiliaryInr
UEOTQM.QuantumPurification.auxiliaryInl_isometry
UEOTQM.QuantumPurification.auxiliaryInr_isometry
UEOTQM.QuantumPurification.extendAuxiliaryLeft
UEOTQM.QuantumPurification.extendAuxiliaryRight
UEOTQM.QuantumPurification.extendAuxiliaryLeft_pure
UEOTQM.QuantumPurification.extendAuxiliaryRight_pure
UEOTQM.QuantumPurification.extendAuxiliaryLeft_marginal
UEOTQM.QuantumPurification.extendAuxiliaryRight_marginal
UEOTQM.QuantumPurification.pure_states_common_extension_unitary
UEOTQM.HistoryChannels.principal
UEOTQM.HistoryChannels.principal_apply
UEOTQM.HistoryChannels.principal_posSemidef
UEOTQM.HistoryChannels.sum_mass_principal
UEOTQM.HistoryChannels.weight
UEOTQM.HistoryChannels.weight_nonneg
UEOTQM.HistoryChannels.sum_weight
UEOTQM.HistoryChannels.zero_weight_iff_zero_principal
UEOTQM.HistoryChannels.restrictEffect
UEOTQM.HistoryChannels.restrictEffect_matrix
UEOTQM.HistoryChannels.pairing_block
UEOTQM.HistoryChannels.readKraus
UEOTQM.HistoryChannels.readKraus_gram
UEOTQM.HistoryChannels.readKraus_conjugation
UEOTQM.HistoryChannels.controlled
UEOTQM.HistoryChannels.controlled_operator_apply
UEOTQM.HistoryChannels.controlled_operator_eq_block
UEOTQM.HistoryChannels.pairing_operator_difference_le_mass_mul_defect
UEOTQM.HistoryChannels.controlled_state_distance_le_weighted_defect
UEOTQM.HistoryChannels.controlled_state_distance_le_uniform
UEOTQM.HistoryChannels.controlled_defect_le_uniform
UEOTQM.HistoryChannels.weights
UEOTQM.HistoryChannels.weights_weight
UEOTQM.HistoryChannels.event_probability
UEOTQM.HistoryChannels.record_tv_le_joint_distance
UEOTQM.HistoryChannels.relabelKraus
UEOTQM.HistoryChannels.relabelKraus_gram
UEOTQM.HistoryChannels.relabel
UEOTQM.HistoryChannels.relabel_operator_apply
UEOTQM.HistoryChannels.relabel_state_matrix
UEOTQM.HistoryChannels.relabel_state_distance
UEOTQM.ActionRecordBudget.History
UEOTQM.ActionRecordBudget.Controller
UEOTQM.ActionRecordBudget.Controller.dropLast
UEOTQM.ActionRecordBudget.Controller.shift
UEOTQM.ActionRecordBudget.Controller.dropLast_shift
UEOTQM.ActionRecordBudget.Controller.word
UEOTQM.ActionRecordBudget.Controller.word_length
UEOTQM.ActionRecordBudget.Controller.protocol
UEOTQM.ActionRecordBudget.Controller.leafEquiv
UEOTQM.ActionRecordBudget.Controller.protocol_word_eq
UEOTQM.ActionRecordBudget.Controller.snoc_head
UEOTQM.ActionRecordBudget.Controller.tail_snoc
UEOTQM.ActionRecordBudget.Controller.word_snoc
UEOTQM.ActionRecordBudget.Controller.truncate
UEOTQM.ActionRecordBudget.Controller.truncate_last
UEOTQM.ActionRecordBudget.Controller.truncate_dropLast
UEOTQM.ActionRecordBudget.initialEquiv
UEOTQM.ActionRecordBudget.extendEquiv
UEOTQM.ActionRecordBudget.step
\end{Verbatim}

\subsection{F.3.17 Requests 961--1020}

\begin{Verbatim}[breaklines,breakanywhere,fontsize=\scriptsize]
UEOTQM.ActionRecordBudget.step_defect_le_uniform
UEOTQM.ActionRecordBudget.record_defect_self
UEOTQM.ActionRecordBudget.step_state_distance_le_weighted
UEOTQM.ActionRecordBudget.step_state_distance_le_uniform
UEOTQM.ActionRecordBudget.step_operator_apply
UEOTQM.ActionRecordBudget.step_operator_eq_block
UEOTQM.ActionRecordBudget.run
UEOTQM.ActionRecordBudget.recordOperator_append
UEOTQM.ActionRecordBudget.run_matrix
UEOTQM.ActionRecordBudget.run_eq_protocol_state
UEOTQM.ActionRecordBudget.terminal_memory_matrix
UEOTQM.ActionRecordBudget.weight_run
UEOTQM.ActionRecordBudget.weight_run_leaf
UEOTQM.ActionRecordBudget.actionCost
UEOTQM.ActionRecordBudget.actionCost_nonneg
UEOTQM.ActionRecordBudget.actionCost_castSucc
UEOTQM.ActionRecordBudget.actionCost_last
UEOTQM.ActionRecordBudget.run_succ_distance_le
UEOTQM.ActionRecordBudget.run_distance_le_actionCost_sum
UEOTQM.ActionRecordBudget.actionCost_le
UEOTQM.ActionRecordBudget.run_distance_le_budget_sum
UEOTQM.ActionRecordBudget.record_tv_le_actionCost_sum
UEOTQM.ActionRecordBudget.record_tv_le_budget_sum
UEOTQM.ActionRecordBudget.protocol_record_tv_eq
UEOTQM.ActionRecordBudget.protocol_record_tv_le_budget_sum
UEOTQM.QuantumTraceDistance.unitaryDiagonal
UEOTQM.QuantumTraceDistance.unitaryDiagonal_add
UEOTQM.QuantumTraceDistance.unitaryDiagonal_sub
UEOTQM.QuantumTraceDistance.unitaryDiagonal_mul
UEOTQM.QuantumTraceDistance.unitaryDiagonal_zero
UEOTQM.QuantumTraceDistance.unitaryDiagonal_one
UEOTQM.QuantumTraceDistance.unitaryDiagonal_posSemidef
UEOTQM.QuantumTraceDistance.mass_unitaryDiagonal
UEOTQM.QuantumTraceDistance.positivePart
UEOTQM.QuantumTraceDistance.negativePart
UEOTQM.QuantumTraceDistance.positivePart_posSemidef
UEOTQM.QuantumTraceDistance.negativePart_posSemidef
UEOTQM.QuantumTraceDistance.positivePart_sub_negativePart
UEOTQM.QuantumTraceDistance.max_add_max_neg
UEOTQM.QuantumTraceDistance.abs_eq_positivePart_add_negativePart
UEOTQM.QuantumTraceDistance.traceAbsolute
UEOTQM.QuantumTraceDistance.traceAbsolute_nonneg
UEOTQM.QuantumTraceDistance.traceAbsolute_eq_sum_abs_eigenvalues
UEOTQM.QuantumTraceDistance.positiveSpectrumEffect
UEOTQM.QuantumTraceDistance.positiveSpectrumEffect_matrix
UEOTQM.QuantumTraceDistance.positiveSpectrumEffect_idempotent
UEOTQM.QuantumTraceDistance.positiveSpectrumEffect_mul_positivePart
UEOTQM.QuantumTraceDistance.positiveSpectrumEffect_mul_negativePart
UEOTQM.QuantumTraceDistance.positiveSpectrumEffect_pairing
UEOTQM.QuantumTraceDistance.difference_parts_mass_eq
UEOTQM.QuantumTraceDistance.effectDistance_le_mass_of_psd_split
UEOTQM.QuantumTraceDistance.effectDistance_eq_positivePart_mass
UEOTQM.QuantumTraceDistance.positiveSpectrumEffect_attains
UEOTQM.QuantumTraceDistance.effectDistance_eq_half_traceAbsolute
UEOTQM.QuantumTraceDistance.effectDistance_eq_half_sum_abs_eigenvalues
UEOTQM.BinaryDiscrimination.equalPriorSuccess
UEOTQM.BinaryDiscrimination.equalPriorSuccess_eq
UEOTQM.BinaryDiscrimination.equalPriorSuccess_le
UEOTQM.BinaryDiscrimination.positiveSpectrumEffect_success
UEOTQM.BinaryDiscrimination.exists_optimal_equalPriorDecision
\end{Verbatim}

\subsection{F.3.18 Requests 1021--1080}

\begin{Verbatim}[breaklines,breakanywhere,fontsize=\scriptsize]
UEOTQM.BinaryDiscrimination.optimal_equalPriorSuccess_traceAbsolute
UEOTQM.BinaryDiscrimination.effectDistance_eq_mass_of_support_certificate
UEOTQM.BinaryDiscrimination.success_eq_optimum_of_support_certificate
UEOTQM.BinaryDiscrimination.effectDistance_eq_sum_positive_of_unitary_data
UEOTQM.BinaryDiscrimination.effectDistance_eq_half_sum_abs_of_unitary_data
UEOTQM.QuantumPurification.coordinateIsometry
UEOTQM.QuantumPurification.coordinateIsometry_adjoint_mul
UEOTQM.QuantumPurification.exists_matrix_isometry_of_card_le
UEOTQM.QuantumPurification.card_le_of_matrix_isometry
UEOTQM.QuantumPurification.exists_matrix_isometry_iff_card_le
UEOTQM.QuantumPurification.matrix_isometry_not_coisometry_of_card_lt
UEOTQM.QuantumPurification.auxiliary_tensor_mul
UEOTQM.QuantumPurification.exists_auxiliary_isometry_of_equal_marginal
UEOTQM.QuantumPurification.exists_auxiliary_isometry_action
UEOTQM.QuantumPurification.pure_states_isometrically_equivalent
UEOTQM.QuantumPurification.SpectralSupport
UEOTQM.QuantumPurification.spectralSupport_card
UEOTQM.QuantumPurification.supportedGram
UEOTQM.QuantumPurification.density_spectral_expansion
UEOTQM.QuantumPurification.supportedGram_spec
UEOTQM.QuantumPurification.supportedGram_vector_normalized
UEOTQM.QuantumPurification.supportPurification
UEOTQM.QuantumPurification.supportPurification_matrix
UEOTQM.QuantumPurification.supportPurification_isPure
UEOTQM.QuantumPurification.supportPurification_marginal
UEOTQM.QuantumPurification.rank_marginal_le_auxiliary_card
UEOTQM.QuantumPurification.exists_purification_of_rank_le
UEOTQM.QuantumPurification.exists_purification_iff_rank_le
UEOTQM.QuantumPurification.exists_minimal_purification
UEOTQM.QuantumPurification.no_purification_of_card_lt_rank
UEOTQM.QuantumPurification.density_rank_pos
UEOTQM.QuantumPurification.no_purification_empty_auxiliary
UEOTQM.StoppingProtocols.Leaf
UEOTQM.StoppingProtocols.height
UEOTQM.StoppingProtocols.height_next_le
UEOTQM.StoppingProtocols.action
UEOTQM.StoppingProtocols.advance
UEOTQM.StoppingProtocols.follow
UEOTQM.StoppingProtocols.controller
UEOTQM.StoppingProtocols.controller_truncate
UEOTQM.StoppingProtocols.controller_dropLast
UEOTQM.StoppingProtocols.controller_zero
UEOTQM.StoppingProtocols.controller_shift
UEOTQM.StoppingProtocols.controller_word_succ
UEOTQM.StoppingProtocols.encode
UEOTQM.StoppingProtocols.decode
UEOTQM.StoppingProtocols.decode_encode
UEOTQM.StoppingProtocols.encode_of_decode
UEOTQM.StoppingProtocols.encode_injective
UEOTQM.StoppingProtocols.decode_eq_some_iff
UEOTQM.StoppingProtocols.word_length_le_height
UEOTQM.StoppingProtocols.stopTime
UEOTQM.StoppingProtocols.stopTime_le
UEOTQM.StoppingProtocols.encode_record
UEOTQM.StoppingProtocols.liftKraus
UEOTQM.StoppingProtocols.lift
UEOTQM.StoppingProtocols.lift_branch_active
UEOTQM.StoppingProtocols.lift_branch_active_none
UEOTQM.StoppingProtocols.lift_branch_idle
UEOTQM.StoppingProtocols.lift_branch_idle_some
\end{Verbatim}

\subsection{F.3.19 Requests 1081--1123}

\begin{Verbatim}[breaklines,breakanywhere,fontsize=\scriptsize]
UEOTQM.StoppingProtocols.record_zero
UEOTQM.StoppingProtocols.record_eq_decode
UEOTQM.StoppingProtocols.record_encode
UEOTQM.StoppingProtocols.run_matrix
UEOTQM.StoppingRecordBudget.isometryChannel
UEOTQM.StoppingRecordBudget.isometryChannel_operator
UEOTQM.StoppingRecordBudget.isometry_complement_positive
UEOTQM.StoppingRecordBudget.extendEffect
UEOTQM.StoppingRecordBudget.pullback_extendEffect
UEOTQM.StoppingRecordBudget.isometry_state_distance
UEOTQM.StoppingRecordBudget.embeddingChannel
UEOTQM.StoppingRecordBudget.embedding_operator_image
UEOTQM.StoppingRecordBudget.embedding_operator_outside_left
UEOTQM.StoppingRecordBudget.embedding_operator_outside_right
UEOTQM.StoppingRecordBudget.embedding_state_distance
UEOTQM.StoppingRecordBudget.someMemoryEmbedding
UEOTQM.StoppingRecordBudget.active_state_eq_embedding
UEOTQM.StoppingRecordBudget.active_state_distance
UEOTQM.StoppingRecordBudget.active_defect
UEOTQM.StoppingRecordBudget.idle_state_eq
UEOTQM.StoppingRecordBudget.idle_defect
UEOTQM.StoppingRecordBudget.leafEmbedding
UEOTQM.StoppingRecordBudget.outside_leafEmbedding
UEOTQM.StoppingRecordBudget.run_eq_embedding
UEOTQM.StoppingRecordBudget.original_joint_distance_eq
UEOTQM.StoppingRecordBudget.weight_run_encode
UEOTQM.StoppingRecordBudget.weight_run_decode_none
UEOTQM.StoppingRecordBudget.original_record_tv_eq
UEOTQM.StoppingRecordBudget.activeDefect
UEOTQM.StoppingRecordBudget.lift_defect_eq
UEOTQM.StoppingRecordBudget.activeCost
UEOTQM.StoppingRecordBudget.actionCost_eq_activeCost
UEOTQM.StoppingRecordBudget.activeCost_eq_sum_active
UEOTQM.StoppingRecordBudget.original_joint_distance_le_activeCost_sum
UEOTQM.StoppingRecordBudget.original_record_tv_le_activeCost_sum
UEOTQM.StoppingRecordBudget.terminal_memory_distance_le_activeCost_sum
UEOTQM.StoppingRecordBudget.activeProbability
UEOTQM.StoppingRecordBudget.activeProbability_nonneg
UEOTQM.StoppingRecordBudget.activeProbability_le_one
UEOTQM.StoppingRecordBudget.activeCost_le_probability_mul
UEOTQM.StoppingRecordBudget.original_joint_distance_le_probability_budget
UEOTQM.QuantumPurification.spectralSupportFintype
UEOTQM.QuantumPurification.spectralSupportDecidableEq
\end{Verbatim}

\section{F.4 Actual increment from the v31.10 finite-conditioning
research revision}

The research revision published as edition 3.6 added eight public
modules. For the same original finite experiment, they supply a spectral
construction for singular support, the full conditioned CQ state,
allowed-action and observation-fibre interfaces, criteria on specified
reachable inputs, actual filtering and retry costs based on the original
leaves, and sharp conditioning-error bounds. Concrete action
counterexamples and the positive Pauli certificate retain separate
names. The actual integrated build, axiom whitelist, source manifest, 19
semantic/numerical controls, and 95 exact finite checks are recorded in
\texttt{physics/\allowbreak qm/\allowbreak kernel/\allowbreak evidence/\allowbreak v31.\allowbreak 10} and
\texttt{physics/\allowbreak qm/\allowbreak conditioning}, with their respective evidence scopes.
Successful compilation is not a new law of nature, and floating-point
diagnostics are not rigorous numerical certificates. Actual execution
status is determined by the receipts bound to the corresponding source
bytes.

The spectral construction for arbitrary original experimental trees and
the full CQ endpoint theorem are formalized in Lean. The classification
of initialization over all inputs, the impossibility argument when
arbitrary reference correlations must be preserved, the sufficient
Kraus-norm error bound, and the monotone-convergence justification of
the expected cost of independent retries until success are handwritten
proofs in Chapter 66. The eight records and costs of the concrete
three-step model are checked by an exact finite program. No separate
Lean theorem verifying all of that runtime program is claimed. The
earlier 161 numbered results, 11 original-source entries, and three
inherited programs retain their original identities.

The current public package contains 49 modules, 959 explicit theorems,
and 1285 named axiom requests. The v31.10 increment consists of eight
modules, 105 explicit theorems, and 162 requests. The publication
revision adds no mathematical declaration and does not rerun these local
scientific gates. Section C.8 records the actual PR/main identities and
exact reuse conditions.

\section{F.5 The 162 named axiom requests added in v31.10}

The names below follow the order of the current
\texttt{Axiom\allowbreak Audit.\allowbreak lean}. Appending them to the 1123 requests in F.3
gives exactly the current 1285 unique requests. They retain the actual
v31.10 acceptance identity and are not new proofs produced by
translation.

\begin{Verbatim}[breaklines,breakanywhere,fontsize=\scriptsize]
UEOTQM.QuantumActionRealizability.labelInstrument
UEOTQM.QuantumActionRealizability.labelInstrument_branch
UEOTQM.QuantumActionRealizability.labelInstrument_channel
UEOTQM.QuantumActionRealizability.fairWeights
UEOTQM.QuantumActionRealizability.fair_label_mixture_branch
UEOTQM.QuantumActionRealizability.swapped_labels_same_channel
UEOTQM.QuantumActionRealizability.swapped_labels_distinct_visible_branch
UEOTQM.QuantumActionRealizability.swapped_labels_not_realizable
UEOTQM.QuantumActionRealizability.hidden_nodes_individually_realizable
UEOTQM.QuantumActionRealizability.shared_observation_not_realizable
UEOTQM.QuantumActionRealizability.visibleTrueChoiMass
UEOTQM.QuantumActionRealizability.visibleTrueChoiMass_label
UEOTQM.QuantumActionRealizability.swapped_labels_choi_separation
UEOTQM.QuantumActionRealizability.fiberDifferenceChoiMass
UEOTQM.QuantumActionRealizability.shared_observation_choi_separation
UEOTQM.QuantumActionPauliWitness.pauliZ
UEOTQM.QuantumActionPauliWitness.pauliX
UEOTQM.QuantumActionPauliWitness.pauliZ_gram
UEOTQM.QuantumActionPauliWitness.pauliX_gram
UEOTQM.QuantumActionPauliWitness.labelUnitary
UEOTQM.QuantumActionPauliWitness.labelUnitary_branch
UEOTQM.QuantumActionPauliWitness.allowed
UEOTQM.QuantumActionPauliWitness.targetWeighted
UEOTQM.QuantumActionPauliWitness.target
UEOTQM.QuantumActionPauliWitness.target_branch
UEOTQM.QuantumActionPauliWitness.fair_allowed_realizes_target
UEOTQM.QuantumActionPauliWitness.target_choi_positive_certificate
UEOTQM.QuantumActionPauliWitness.sourceWeighted
UEOTQM.QuantumActionPauliWitness.source
UEOTQM.QuantumActionPauliWitness.source_success_eq_half_target
UEOTQM.QuantumActionPauliWitness.source_failure_branch
UEOTQM.QuantumActionRealizability.matrixUnit
UEOTQM.QuantumActionRealizability.matrix_expansion
UEOTQM.QuantumActionRealizability.branchLinear
UEOTQM.QuantumActionRealizability.branchLinear_apply
UEOTQM.QuantumActionRealizability.choi
UEOTQM.QuantumActionRealizability.choi_injective
UEOTQM.QuantumActionRealizability.conjugationLinear
UEOTQM.QuantumActionRealizability.instrumentChoiAfter
UEOTQM.QuantumActionRealizability.targetChoiAfter
UEOTQM.QuantumActionRealizability.instrumentChoiAfter_eq_iff
UEOTQM.QuantumActionRealizability.mixture_instrumentChoiAfter
UEOTQM.QuantumActionRealizability.weightsOfPMF
UEOTQM.QuantumActionRealizability.realizes_of_weights
UEOTQM.QuantumActionRealizability.realizable_after_iff_mem_convexHull
UEOTQM.QuantumActionRealizability.full_instrument_realizable_iff
UEOTQM.QuantumActionRealizability.support_realizable_iff
UEOTQM.QuantumActionRealizability.observation_realizable_iff_fiber_choi
UEOTQM.QuantumActionRealizability.separator_excludes_convexHull
UEOTQM.QuantumActionRealizability.support_agreement_on_fixed_input
UEOTQM.QuantumFiniteWindow.SupportFilter
UEOTQM.QuantumFiniteWindow.SupportFilter.identity
UEOTQM.QuantumFiniteWindow.SupportFilter.zero
UEOTQM.QuantumFiniteWindow.SupportFilter.apply
UEOTQM.QuantumFiniteWindow.SupportFilter.metric
UEOTQM.QuantumFiniteWindow.SupportFilter.complement
UEOTQM.QuantumFiniteWindow.SupportFilter.apply_posSemidef
UEOTQM.QuantumFiniteWindow.SupportFilter.apply_smul
UEOTQM.QuantumFiniteWindow.SupportFilter.apply_zero
UEOTQM.QuantumFiniteWindow.SupportFilter.identity_apply
UEOTQM.QuantumFiniteWindow.SupportFilter.zero_apply
UEOTQM.QuantumFiniteWindow.SupportFilter.complement_adjoint
UEOTQM.QuantumFiniteWindow.SupportFilter.complement_sq
UEOTQM.QuantumFiniteWindow.SupportFilter.complement_matrix
UEOTQM.QuantumFiniteWindow.SupportFilter.inverse_metric
UEOTQM.QuantumFiniteWindow.StepCertificate
UEOTQM.QuantumFiniteWindow.tiltedKraus
UEOTQM.QuantumFiniteWindow.tilted_main_gram
UEOTQM.QuantumFiniteWindow.tilted_fallback_gram
UEOTQM.QuantumFiniteWindow.tiltedInstrument
UEOTQM.QuantumFiniteWindow.tilted_main_transport
UEOTQM.QuantumFiniteWindow.tilted_fallback_zero
UEOTQM.QuantumFiniteWindow.tilted_branch_transport
UEOTQM.QuantumFiniteWindow.FilterTree
UEOTQM.QuantumFiniteWindow.FilterTree.rootFilter
UEOTQM.QuantumFiniteWindow.FilterTree.Valid
UEOTQM.QuantumFiniteWindow.FilterTree.Leaf
UEOTQM.QuantumFiniteWindow.FilterTree.leafFintype
UEOTQM.QuantumFiniteWindow.FilterTree.accepts
UEOTQM.QuantumFiniteWindow.FilterTree.sourceWord
UEOTQM.QuantumFiniteWindow.FilterTree.sourceWord_injective
UEOTQM.QuantumFiniteWindow.FilterTree.sourceLeafOperator
UEOTQM.QuantumFiniteWindow.FilterTree.sourceLeafOperator_eq_recordOperator
UEOTQM.QuantumFiniteWindow.FilterTree.acceptedLeafOperator
UEOTQM.QuantumFiniteWindow.FilterTree.tiltedLeafOperator
UEOTQM.QuantumFiniteWindow.FilterTree.leaf_transport
UEOTQM.QuantumFiniteWindow.FilterTree.tiltedLeafOperator_smul
UEOTQM.QuantumFiniteWindow.FilterTree.tiltedLeafOperator_posSemidef
UEOTQM.QuantumFiniteWindow.FilterTree.sum_tiltedLeafOperator_trace
UEOTQM.QuantumFiniteWindow.FilterTree.successMass
UEOTQM.QuantumFiniteWindow.FilterTree.successMass_eq_filtered_mass
UEOTQM.QuantumFiniteWindow.FilterTree.successMass_nonneg
UEOTQM.QuantumFiniteWindow.FilterTree.initialState
UEOTQM.QuantumFiniteWindow.FilterTree.initialState_matrix
UEOTQM.QuantumFiniteWindow.FilterTree.conditional_leaf
UEOTQM.QuantumFiniteWindow.FilterTree.zero_success_filter
UEOTQM.QuantumFiniteWindow.FilterTree.zero_success_no_positive
UEOTQM.QuantumFiniteWindow.FilterTree.tiltedLeafOperator_zero
UEOTQM.QuantumFiniteWindow.FilterTree.zero_success_all_accepted_blocks_zero
UEOTQM.QuantumFiniteWindow.FilterTree.sourceLeafOperator_posSemidef
UEOTQM.QuantumFiniteWindow.FilterTree.sum_sourceLeafOperator_trace
UEOTQM.QuantumFiniteWindow.FilterTree.sourceWeights
UEOTQM.QuantumFiniteWindow.FilterTree.sourceWeights_accepted_sum
UEOTQM.QuantumFiniteWindow.FilterTree.successMass_le_one
UEOTQM.QuantumFiniteWindow.FilterTree.metric_complement_posSemidef
UEOTQM.QuantumFiniteWindow.FilterTree.leafDecidableEq
UEOTQM.QuantumFiniteWindow.block_eq_sum_lifts
UEOTQM.QuantumFiniteWindow.block_posSemidef
UEOTQM.QuantumFiniteWindow.block_smul
UEOTQM.QuantumFiniteWindow.FilterTree.jointState
UEOTQM.QuantumFiniteWindow.FilterTree.acceptedJointOperator
UEOTQM.QuantumFiniteWindow.FilterTree.conditional_joint
UEOTQM.QuantumFiniteWindow.FilterTree.conditional_record_probability
UEOTQM.QuantumFiniteWindow.FilterTree.fallback_independent
UEOTQM.QuantumFiniteWindow.spectralFilter
UEOTQM.QuantumFiniteWindow.spectralFilter_matrix_posSemidef
UEOTQM.QuantumFiniteWindow.spectralFilter_metric
UEOTQM.QuantumFiniteWindow.SupportFilter.matrix_projection
UEOTQM.QuantumFiniteWindow.SupportFilter.matrix_complement
UEOTQM.QuantumFiniteWindow.gram_sum_zero_component
UEOTQM.QuantumFiniteWindow.support_of_backward
UEOTQM.QuantumFiniteWindow.StepCertificate.of_backward
UEOTQM.QuantumFiniteWindow.SourceTree
UEOTQM.QuantumFiniteWindow.backwardGram
UEOTQM.QuantumFiniteWindow.backwardGram_posSemidef
UEOTQM.QuantumFiniteWindow.SourceTree.canonical
UEOTQM.QuantumFiniteWindow.SourceTree.canonical_valid
UEOTQM.QuantumFiniteWindow.FilterTree.erase
UEOTQM.QuantumFiniteWindow.canonical_erases_to_source
UEOTQM.QuantumFiniteWindow.source_node_conditional_joint
UEOTQM.QuantumInitializationCost.gramFactor
UEOTQM.QuantumInitializationCost.gramFactor_gram
UEOTQM.QuantumInitializationCost.filterInstrument
UEOTQM.QuantumInitializationCost.filter_success_branch
UEOTQM.QuantumInitializationCost.gram_filter_mass
UEOTQM.QuantumInitializationCost.filter_success_probability
UEOTQM.QuantumInitializationCost.filter_success_probability_pos
UEOTQM.QuantumInitializationCost.filter_posterior_matrix
UEOTQM.QuantumInitializationCost.scaled_filter_mean_attempts
UEOTQM.QuantumInitializationCost.successProbability
UEOTQM.QuantumInitializationCost.meanTrialCost
UEOTQM.QuantumInitializationCost.retryTreeCost
UEOTQM.QuantumInitializationCost.successProbability_nonneg
UEOTQM.QuantumInitializationCost.failureWeight_eq
UEOTQM.QuantumInitializationCost.successProbability_le_one
UEOTQM.QuantumInitializationCost.retryTreeCost_succ
UEOTQM.QuantumInitializationCost.retryTreeCost_closed
UEOTQM.QuantumInitializationCost.retryTreeCost_tendsto
UEOTQM.QuantumInitializationCost.sourceInitialFilter
UEOTQM.QuantumInitializationCost.sourceInitialFilter_probability
UEOTQM.QuantumInitializationCost.source_successProbability
UEOTQM.QuantumInitializationCost.source_retryTreeCost_tendsto
UEOTQM.QuantumPostselectionError.binary_gap_le_div_of_le
UEOTQM.QuantumPostselectionError.binary_gap_le_div_max
UEOTQM.QuantumPostselectionError.probability_mul_posterior
UEOTQM.QuantumPostselectionError.posterior_effectDistance_le_div_max
UEOTQM.QuantumPostselectionError.posterior_effectDistance_le_min
UEOTQM.QuantumPostselectionError.posterior_half_traceAbsolute_le_min
UEOTQM.QuantumActionRealizability.branchOutputFamily
UEOTQM.QuantumActionRealizability.mixture_branchOutputFamily
UEOTQM.QuantumActionRealizability.fixed_input_branch_realizable_iff
UEOTQM.QuantumActionRealizability.fixed_input_observation_realizable_iff
\end{Verbatim}

\section{F.6 The v31.11 operational-identity increment}

This increment adds one public module,
\texttt{Operational\allowbreak Equivalence.\allowbreak lean}, with 28 explicit theorems and 43
explicitly named declarations. The actual namespace audit covers 80
exported constants, including generated constants. The target module
compiled successfully; only \texttt{propext}, \texttt{Classical.\allowbreak choice},
and \texttt{Quot.\allowbreak sound} are accepted as foundations. Exact execution and
source hashes are in
\texttt{.\allowbreak .\allowbreak /\allowbreak kernel/\allowbreak evidence/\allowbreak v31.\allowbreak 11/\allowbreak lean-\allowbreak verification.\allowbreak json}. The 8730
build jobs are dependency-graph work items, not a count of new theorems.
Ordinary linter warnings remain.

The current public source inventory is 50 modules, 987 explicit theorems
and 1365 named axiom requests. These inventory totals do not claim a
fresh local rerun of every historical gate. The old 49 source modules
and 1285 requests retain their historical identities; the new module and
public entrypoints have separately bound current evidence. Chapter 67's
ledger gives P-OPT-02's actual-map, single-system, complete
ancillary-statistics and Choi endpoints and rectangular-event coverage.
P-OPT-03's Lean subset covers sequential composition and sum congruence
only. The complete finite Test/parallel/quotient-tree background and the
infinite-dimensional, trace-class, robust and mixing paper arguments are
not made machine theorems by a neighboring build.

The following is the complete list of constants actually audited in the
new namespace, matching the added public audit requests:

\begin{Verbatim}[breaklines,breakanywhere,fontsize=\scriptsize]
UEOTQM.OperationalEquivalence.ContextStatistics
UEOTQM.OperationalEquivalence.Event
UEOTQM.OperationalEquivalence.Event._sizeOf_1
UEOTQM.OperationalEquivalence.Event._sizeOf_inst
UEOTQM.OperationalEquivalence.Event.amplify
UEOTQM.OperationalEquivalence.Event.amplify._proof_1
UEOTQM.OperationalEquivalence.Event.amplify._simp_1
UEOTQM.OperationalEquivalence.Event.amplify.eq_1
UEOTQM.OperationalEquivalence.Event.amplify_operator_apply
UEOTQM.OperationalEquivalence.Event.amplify_operator_eq
UEOTQM.OperationalEquivalence.Event.amplify_posSemidef
UEOTQM.OperationalEquivalence.Event.casesOn
UEOTQM.OperationalEquivalence.Event.ctorIdx
UEOTQM.OperationalEquivalence.Event.defect_positive
UEOTQM.OperationalEquivalence.Event.kraus
UEOTQM.OperationalEquivalence.Event.mass_operator_le
UEOTQM.OperationalEquivalence.Event.mk
UEOTQM.OperationalEquivalence.Event.mk._flat_ctor
UEOTQM.OperationalEquivalence.Event.mk.congr_simp
UEOTQM.OperationalEquivalence.Event.mk.inj
UEOTQM.OperationalEquivalence.Event.mk.injEq
UEOTQM.OperationalEquivalence.Event.mk.noConfusion
UEOTQM.OperationalEquivalence.Event.mk.sizeOf_spec
UEOTQM.OperationalEquivalence.Event.noConfusion
UEOTQM.OperationalEquivalence.Event.noConfusionType
UEOTQM.OperationalEquivalence.Event.ofChannel
UEOTQM.OperationalEquivalence.Event.ofChannel._proof_1
UEOTQM.OperationalEquivalence.Event.ofChannel_operator
UEOTQM.OperationalEquivalence.Event.operatorLinear
UEOTQM.OperationalEquivalence.Event.operatorLinear._proof_1
UEOTQM.OperationalEquivalence.Event.operatorLinear._proof_2
UEOTQM.OperationalEquivalence.Event.operatorLinear.eq_1
UEOTQM.OperationalEquivalence.Event.operator_posSemidef
UEOTQM.OperationalEquivalence.Event.rec
UEOTQM.OperationalEquivalence.Event.recOn
UEOTQM.OperationalEquivalence.Event.statistics_iff_map_eq
UEOTQM.OperationalEquivalence.Event.trace_operator
UEOTQM.OperationalEquivalence.PhysicalSystem
UEOTQM.OperationalEquivalence.SingleStatistics
UEOTQM.OperationalEquivalence.amplify
UEOTQM.OperationalEquivalence.amplify._proof_1
UEOTQM.OperationalEquivalence.amplify._proof_2
UEOTQM.OperationalEquivalence.amplify_congr
UEOTQM.OperationalEquivalence.choi
UEOTQM.OperationalEquivalence.choi_injective
UEOTQM.OperationalEquivalence.complex_linear_eq_of_eq_on_densities
UEOTQM.OperationalEquivalence.contextStatistics_iff_choi_eq
UEOTQM.OperationalEquivalence.contextStatistics_iff_map_eq
UEOTQM.OperationalEquivalence.contextStatistics_iff_singleStatistics
UEOTQM.OperationalEquivalence.context_statistics_of_singleStatistics
UEOTQM.OperationalEquivalence.hermitian_eq_of_all_effect_pairings_eq
UEOTQM.OperationalEquivalence.imaginaryUnit
UEOTQM.OperationalEquivalence.imaginaryUnit.eq_1
UEOTQM.OperationalEquivalence.imaginaryUnit_hermitian
UEOTQM.OperationalEquivalence.imaginaryUnit_hermitian._simp_1_2
UEOTQM.OperationalEquivalence.labelled_comp_congr
UEOTQM.OperationalEquivalence.labelled_sum_congr
UEOTQM.OperationalEquivalence.map_eq_iff_choi_eq
UEOTQM.OperationalEquivalence.matrixUnit
UEOTQM.OperationalEquivalence.matrixUnit.congr_simp
UEOTQM.OperationalEquivalence.matrixUnit.eq_1
UEOTQM.OperationalEquivalence.matrixUnit_hermitian_expansion
UEOTQM.OperationalEquivalence.matrixUnit_hermitian_expansion._simp_1_4
UEOTQM.OperationalEquivalence.matrix_expansion
UEOTQM.OperationalEquivalence.physicalState
UEOTQM.OperationalEquivalence.physicalState._proof_1
UEOTQM.OperationalEquivalence.real_maps_eq_on_hermitian_of_eq_on_densities
UEOTQM.OperationalEquivalence.real_maps_eq_on_positive_of_eq_on_densities
UEOTQM.OperationalEquivalence.scalarDensity
UEOTQM.OperationalEquivalence.scalarDensity._proof_1
UEOTQM.OperationalEquivalence.scalarDensity._proof_2
UEOTQM.OperationalEquivalence.scalarEffect
UEOTQM.OperationalEquivalence.scalarEffect._proof_1
UEOTQM.OperationalEquivalence.scalarEffect._proof_2
UEOTQM.OperationalEquivalence.scalar_context_pairing
UEOTQM.OperationalEquivalence.singleStatistics_iff_map_eq
UEOTQM.OperationalEquivalence.sub_kronecker
UEOTQM.OperationalEquivalence.symmetricUnit
UEOTQM.OperationalEquivalence.symmetricUnit.eq_1
UEOTQM.OperationalEquivalence.symmetricUnit_hermitian
\end{Verbatim}

\section{F.7 Proof and execution scope of the v31.12 growing-carrier
increment}

The infinite-dimensional, domain, complete-process, and
error-composition arguments for the six Chapter 68 results are paper
proofs. This increment adds no Lean source, theorem, or new axiom-audit
execution. \texttt{growing\_\allowbreak cutoff/\allowbreak verify\_\allowbreak inherited\_\allowbreak evidence.\allowbreak py}
checks 158 proof/configuration inputs and 13 earlier runtime inputs byte
for byte against pinned historical receipts, and checks the complete
proof/tool inventory for additions or changes. When these conditions
hold, it reuses the actual v31.11 CI on commit
\texttt{06ca0dce8b03\allowbreak cd34e86785e9\allowbreak 7fdc3c7aa2c5\allowbreak 82bd}; the 1365 historical
public axiom reports keep their original identity. Any proof-input or
inventory change selects full machine verification rather than extending
this reuse conclusion.

Four new runtime groups in \texttt{physics/\allowbreak qm/\allowbreak growing\_\allowbreak cutoff/\allowbreak } cover
the complete rectangular source, references, stopping records, the
original Chapter 66 interpreter, independent growth checks, and exact
rational certificates. Floating finite examples, exact algebraic
certificates, synthetic observations, and general paper quantifiers are
registered separately in \texttt{kernel/\allowbreak evidence/\allowbreak v31.\allowbreak 12/\allowbreak }. The analytic
cutoff 4095 does not assert construction of that matrix dimension:
full-record numerical examples use smaller known support explicitly
stated in their receipts. Every execution records the actual environment
and checkout identity; checker authors are not relabelled as remote
executors. The three-agent cross-review is internal AI review, not
external human peer review. No real device data were obtained.

\section{F.8 Verification scope of the v31.13 simultaneous-dynamics
increment}

The Chapter 69 trace-class prefix, self-adjointness, finite-Duhamel,
calibration-coverage and error-composition arguments are paper proofs.
None of the six stable results has a new Lean binding. The runtime
directory \texttt{physics/\allowbreak qm/\allowbreak simultaneous\_\allowbreak cutoff/\allowbreak } distinguishes
floating source interpretation, independent prefix and
infinite-displacement controls, rational calibration with an independent
checker, integer-interval prediction and raw synthetic streams. The
K=4175 interval algorithm actually evaluates a sparse polynomial for one
specified \(G\le1\) entangled input; it does not numerically certify
every input or the entire 8351-dimensional matrix. The unchanged 158
proof inputs retain the original 1365 public axiom reports; new runs do
not extend their machine-verified quantifiers. Three-agent internal AI
cross-review is not external human peer review, and the empirical axis
remains untested.

\section{F.9 Verification scope of the v31.14 persistent-bath increment}

The nine Chapter 70 stable results have paper proofs, exact finite
algebra, directed intervals, and synthetic statistical evidence, with
their respective quantifiers stated in the text and
\texttt{chapter70\_\allowbreak formal\_\allowbreak bindings.\allowbreak json}. They are not new Lean
declarations; Chapter 69 likewise adds no Lean coverage. After the
current 158 unchanged proof inputs were hash-checked, their original
execution was reused without relabeling the historical 1365 axiom
reports.

The integrated replay actually reads the original 1.6 million synthetic
rows and separately records its execution time and input/output hashes.
Reused data are not new shots. Eight Choi interval certificates, 504
exact principal minors, frozen calibration/holdout, external
processed-data reproduction, and the first missing-h5py failure attempt
retain separate identities. Complete evidence is in
\texttt{physics/\allowbreak qm/\allowbreak kernel/\allowbreak evidence/\allowbreak v31.\allowbreak 14/\allowbreak }. Publication identity and
scoped visual review are in \texttt{physics/\allowbreak qm/\allowbreak book/\allowbreak evidence/\allowbreak v3.\allowbreak 11/\allowbreak };
reviewing the new and changed pages does not claim a fresh page-by-page
review of all inherited content.

\backmatter
\end{document}
